\documentclass[12pt]{article}

\usepackage[utf8]{inputenc}
\usepackage{CJKutf8}
\usepackage{amsmath,amssymb}
\usepackage{amsthm}
\usepackage{geometry}
\geometry{
  a4paper,
  top=25mm,
  bottom=25mm,
  left=20mm,
  right=20mm
}
\usepackage{xcolor}
\usepackage{url}
\usepackage[unicode,colorlinks=true,linkcolor=blue,urlcolor=blue]{hyperref}
\usepackage{xurl}
\Urlmuskip=0mu plus 2mu
\sloppy
\usepackage{tocloft}
\usepackage{array}
\usepackage{longtable}

% 目录中 section 条目使用点线填充到页码
\renewcommand{\cftsecleader}{\cftdotfill{\cftdotsep}}
\renewcommand{\contentsname}{目录}
\renewcommand{\refname}{参考文献}

% 基本定理环境（工作笔记：形式系统风格；参考 tex21）
\newtheorem{definition}{定义}[section]
\newtheorem{assumption}{假设}[section]
\newtheorem{theorem}{定理}[section]
\newtheorem{proposition}{命题}[section]
\newtheorem{corollary}{推论}[section]
\newtheorem{lemma}{引理}[section]
\newtheorem{remark}{备注}[section]
\newtheorem{Certificate}{证书}[section]

% 记号（仅用于本笔记自洽引用，避免依赖主稿宏定义）
\newcommand{\code}[1]{\path{#1}}
\newcommand{\RR}{\mathbb{R}}
\newcommand{\NN}{\mathbb{N}}

% 纯数卷符号体系：对齐 docs/PureMath/Pub_GFramework_PureMath.tex
\newcommand{\GFramework}{\text{G-Framework}}
\newcommand{\GAlgebra}{\text{G-Algebra}}

\newcommand{\GZLStxt}{\textsf{GZ-LS}}
\newcommand{\GZLTtxt}{\textsf{GZ-LT}}
\newcommand{\GZLOCtxt}{\textsf{GZ-LOC}}
\newcommand{\GZLCurvtxt}{\textsf{GZ-LCurv}}

\newcommand{\GZLS}{\ensuremath{\mathfrak{L}_{\mathrm{GZ}}}}
\newcommand{\GZLT}{\ensuremath{M_{!w}}}
\newcommand{\GZLOC}{\ensuremath{A_M}}
\newcommand{\GZLCurv}{\ensuremath{\mathcal{F}_{\mathrm{law}}}}

\newcommand{\LawSpace}{\ensuremath{\mathcal{L}}}
\newcommand{\Base}{\ensuremath{M}}
\newcommand{\Bundle}{\ensuremath{P}}
\newcommand{\pr}{\operatorname{pr}}
\newcommand{\Group}{\ensuremath{G}}
\newcommand{\Conn}{\ensuremath{\omega}}
\newcommand{\Curv}{\ensuremath{\Omega}}
\newcommand{\Homotopy}{\ensuremath{\mathcal{H}}}
\newcommand{\Jac}{\ensuremath{\mathsf{Jac}}}
\newcommand{\JacGap}{\ensuremath{\mathrm{JacGap}}} % Jacobiator-gap（非导数 Jacobian）
\newcommand{\Layer}{\ensuremath{\ell}}

% 本笔记使用的工程符号
\newcommand{\PTOPB}{\ensuremath{\mathrm{PTOPB}}}
\newcommand{\ETOPB}{\ensuremath{\mathrm{ETOPB}}}
\newcommand{\TOPB}{\ensuremath{\mathrm{TOPB}}}
\newcommand{\Aug}{\ensuremath{\mathrm{Aug}}}
\newcommand{\DeltaRT}{\ensuremath{\Delta_{\mathrm{rt}}}}

% 避免少量不可控的换行溢出（尤其是混排 CJK 与等宽字体时）
\setlength{\emergencystretch}{6em}

\begin{document}
\begin{CJK*}{UTF8}{gkai}

\title{PTOPB/ETOPB 的算子二分：外界实时干扰与非交换构造\\（工作笔记：形式系统版）}
\author{Gao Zheng（高政）}
\date{2026-01-18}
\maketitle

\section*{前言}
\addcontentsline{toc}{section}{前言}

本文的目标不是单独定义 \PTOPB{} 与 \ETOPB{} 两个术语，而是在 \GFramework{} 的
算子工程链条中，把外界实时干扰、非交换构造、纤维内优化、包线迁移、同伦扭曲
实时控制与动态算法选型组织成一个可审计的算子二分系统。若这些对象分散出现，
外界干扰容易被误写为可控构造，工程增强容易被误写为真实环境改造；本文首先
建立来源/可控性二分，再讨论二者如何在同伦扭曲中协同。

本文的第一层立场是 PTOPB/ETOPB 的来源区分。\PTOPB{} 关注被动承受或可观测的
外部实时扰动，\ETOPB{} 关注系统内部可设计、可组合、可调度的构造算子。二者
都影响轨迹，但语义不同：前者提供环境压强，后者提供工程动作。只有把二者分开，
后续的可消/不可消、增强/对抗、相内/相间搜索才有清晰判据。

本文的第二层立场是增强的同伦化。工程增强不是简单增加算子，而是改变纤维内
优化、包线迁移和同伦扭曲的可达结构。相内搜索处理局部结构细化，相间搜索处理
不同结构盆地之间的迁移，实时控制则在动态环境下决定何时调整算法、何时切换
路径、何时承认某些外界扰动不可消。

本文的第三层立场是算子二分的治理意义。PTOPB/ETOPB 的划分并不只服务于记号清晰，
而是服务于工程责任边界：哪些变化来自外界，哪些变化来自构造，哪些变化可被同伦
修复，哪些变化只能被纳入鲁棒控制。由此，系统获得了对干扰、构造与修复行动的
可审计分账。

本文的第四层立场是非交换构造的必要性。工程构造算子通常不能交换：先改变材料
路线再调环境，与先调环境再改变路线，可能得到完全不同的轨迹。本文把这种非交换性
纳入 \ETOPB{} 的核心，而不是将其视作实现细节。由此，算子词、执行顺序和同伦路径
共同成为可审计对象。

本文的第五层立场是两层增强结构。纤维内优化处理同一相或同一局部盆地内的细化，
包线迁移处理跨相、跨盆地或跨策略族的切换。二者的任务不同：前者提高局部质量，
后者改变全局可达结构。本文把二者放在同一同伦扭曲框架中，以避免把局部优化误认为
全局迁移。

本文的第六层立场是实时控制与动态算法选型。外界扰动具有时间性，构造算子也具有
执行窗口，因此算法不能预先固定为单一路线。本文通过相空间穿梭、动态算法选型和
实时控制，把算子二分转化为一套随环境和证书状态更新的控制系统。

本文的第七层立场是为开放系统和命中范式提供来源审计。后续同伦命中若要判断一次
穿透或修复来自何处，必须知道它是环境推动、工程构造、同伦迁移还是鲁棒吸收。
本文的 PTOPB/ETOPB 二分，正是这种来源审计的算子基础。

因此，本文在整体 \GFramework{} 中承担的是“外界扰动与工程构造的算子分解”作用。
它向前连接 law-space 与主丛联络，向中间连接纤维优化、包线迁移和实时控制，向后
为开放系统治理、统计命中与 PDE 残差范式提供可控/不可控来源的判定接口。概括地说，
本文把“环境影响与工程动作混在一起”的问题，推进为 PTOPB/ETOPB 二分下的同伦
扭曲控制系统。

更具体地说，本文把“扰动”与“构造”分账。外界实时扰动可能改变目标、边界或可行域；
工程构造算子则改变系统自身的结构、路径或控制策略。二者若被混为一谈，系统就
无法判断失败是外部压力导致，还是内部构造不足导致。本文的二分正是为这种判断
提供形式接口。

从方法论上看，本文使实时控制具备可追溯来源。每一次算法切换、相空间穿梭或包线
迁移，都可以追问它是在响应 \PTOPB{}，还是在施加 \ETOPB{}，还是二者的组合。这种
来源标记为后续审计、治理和命中泛函分析提供了必要的日志结构。

从整体谱系看，本文位于开放系统治理和策略命中之间。开放系统需要识别外界干扰，
统计命中需要选择可控路径，PDE 残差范式需要分清残差来自环境还是构造不足。本文
给出的算子二分，使这些后续问题有了统一的来源语义。

\setcounter{tocdepth}{2}
\setcounter{secnumdepth}{2}
\tableofcontents
\clearpage

\section*{形式系统写作约定（公理降级、证书与谓词因果链）}
\addcontentsline{toc}{section}{形式系统写作约定（公理降级、证书与谓词因果链）}
本笔记采用“形式逻辑 + 证书”的工作笔记体例，并统一使用\textbf{基于数学符号推演逻辑的谓词因果链}来组织证明要点。
为避免口头叙述与工程接口脱钩，正文默认遵循如下约定：
\begin{itemize}
  \item \textbf{公理降级}：不单列“公理”环境；凡需作为基础假设固定者，均以“公理（降级为定理）”写成\textbf{定理}，
  并配套可复算的\textbf{证书}与\textbf{证明}。若断言包含“按步/按层/按回合/按长度”的全称迭代结构，默认给出\textbf{数学归纳法证书}。
  \item \textbf{定理/引理/命题/推论}：每条编号断言均同时给出（i）证书（可审计见证/日志/不变量/基步-归纳步等），
  （ii）证明（以谓词链闭合方式把推理写成可复核的符号推演步骤）。
  \item \textbf{备注}：只用于动机、解释、示例与工程语境对齐，不承担证明义务。
  \item \textbf{证书调用}：对外复用时建议成对引用“断言 + 证书”（如 \texttt{定理~\ref{thm:short-exact-implies-long-exact}} 与 \texttt{证书~\ref{cert:short-exact-implies-long-exact}}），
  以确保他处调用时具备可追溯的接口；全文证书调用约定与索引见附录B。
\end{itemize}
\clearpage

\section{来源/可控性二分：PTOPB/ETOPB 的干扰与构造算子}
\label{sec:ptopb-etopb-operator-decomposition}

\begin{remark}[核心陈述（来源/可控性二分）]
在“总算子丛（\TOPB）”的表达中，可以按\textbf{来源/可控性}把 \PTOPB{}（物理总算子丛）与 \ETOPB{}（工程总算子丛）各自再切成两大类：
\begin{enumerate}
  \item 外界实时干扰算子族（$\Delta_{\mathrm{P}},\Delta_{\mathrm{E}}$）：来自系统之外，随时间变化，通常不可控或仅能间接影响；
  \item 系统内“功能 + 保障”的非交换构造算子族（$\mathcal{C}_{\mathrm{P}},\mathcal{C}_{\mathrm{E}}$）：来自系统之内，可控，并且一般满足顺序敏感的非交换性。
\end{enumerate}
该二分的意义在于：把工程对抗中的“外生输入”与“内生可控结构”分离出来，并把顺序敏感性提升为一等结构，
从而为后续的博弈对抗分析与证书化审计提供统一的算子语言。
\end{remark}

\subsection{统一形式化框架：总体状态与总算子}
\label{sec:unified-framework}

\begin{definition}[总体状态分解]
令 $\mathcal{X}$ 为物理态空间，$\mathcal{Z}$ 为工程态空间，定义总体状态空间
\[
  \mathcal{S} := \mathcal{X}\times \mathcal{Z},
\]
并将任意状态写为
\[
  s := (x,z),\qquad x\in\mathcal{X},\ z\in\mathcal{Z}.
\]
其中 $x$ 表示物理态（位置/速度/能量/温度/电磁等效量等），$z$ 表示工程态（算法参数、网络拓扑、任务编排、密钥与访问控制、冗余切换、运维策略与保障链路状态等）。
\end{definition}

\begin{definition}[总算子族与阶段更新]
称一族映射
\[
  \mathcal{O} := \{\,O_{\theta}:\mathcal{S}\to\mathcal{S}\,\}
\]
为总算子族（总算子丛的算子侧抽象），其中 $\theta$ 汇总了时间、空间、工况与对抗态势等参数。
在离散阶段口径下，系统更新可写为
\[
  s_{t+1} = O_t(s_t),
\]
其中 $O_t\in\mathcal{O}$。
\end{definition}

\begin{definition}[\PTOPB/\ETOPB{} 的两类主导口径（投影分层）]
令 $\pi_{\mathrm{P}}:\mathcal{S}\to\mathcal{X}$ 与 $\pi_{\mathrm{E}}:\mathcal{S}\to\mathcal{Z}$ 为自然投影。
本章以如下口径抽象区分 \PTOPB{} 与 \ETOPB{}：
\begin{enumerate}
  \item $\PTOPB$-侧：强调算子对 $\pi_{\mathrm{P}}(s)=x$ 的直接塑形与约束（允许通过物理效应间接影响 $z$）；
  \item $\ETOPB$-侧：强调算子对 $\pi_{\mathrm{E}}(s)=z$ 的结构与策略塑形，并通过闭环反馈影响 $x$。
\end{enumerate}
形式上，可将 \PTOPB{} 与 \ETOPB{} 视为 $\mathcal{O}$ 的两个可识别子族：
\[
  \PTOPB \subseteq \mathcal{O},\qquad
  \ETOPB \subseteq \mathcal{O},
\]
其中“可识别”由工程证书（来源、触发与审计边界）给出，而不是由口头解释给出。
\end{definition}

\begin{definition}[外界实时干扰算子与系统构造算子（阶段口径）]
在每个阶段 $t$，设存在一对算子
\[
  \delta_t:\mathcal{S}\to\mathcal{S},\qquad
  \pi_t:\mathcal{S}\to\mathcal{S},
\]
满足 $\delta_t$ 由系统外部因素（环境/对手/基础设施/人因等）触发并参数化，$\pi_t$ 由系统内部结构与策略选择触发并参数化。
我们称 $\delta_t$ 为外界实时干扰算子，称 $\pi_t$ 为系统构造算子，并约定阶段更新算子可写为
\[
  O_t = \pi_t\circ \delta_t.
\]
其中复合顺序体现了“先发生外生输入，再由系统执行内部策略”的一种常用抽象；若场景需要亦可将顺序调换为 $\delta_t\circ \pi_t$。
\end{definition}

\begin{theorem}[公理降级：来源/可控性交错分解与二分成立]
\label{thm:source-controllability-splitting}
对任意有限阶段数 $T\in\mathbb{N}$，若每个阶段 $t\in\{0,1,\dots,T-1\}$ 的更新算子满足 $O_t=\pi_t\circ\delta_t$，
则整体更新算子满足交错分解
\[
  O_{0:T-1}
  :=
  O_{T-1}\circ\cdots\circ O_0
  =
  (\pi_{T-1}\circ\delta_{T-1})\circ\cdots\circ(\pi_0\circ\delta_0).
\]
并且在该口径下，\PTOPB{} 与 \ETOPB{} 均可按来源/可控性作二分：
\[
  \PTOPB = \Delta_{\mathrm{P}}\uplus \mathcal{C}_{\mathrm{P}},
  \qquad
  \ETOPB = \Delta_{\mathrm{E}}\uplus \mathcal{C}_{\mathrm{E}},
  \qquad
  \Delta_{\bullet}\cap \mathcal{C}_{\bullet}=\varnothing.
\]
\end{theorem}

\begin{Certificate}[数学归纳法证书：按阶段数 $T$ 归纳的交错分解]
\label{cert:splitting-by-time-horizon}
定义谓词 $P(T)$ 为：“任意长度为 $T$ 的阶段更新算子复合 $O_{0:T-1}$ 可写为交错复合 $(\pi_{T-1}\circ\delta_{T-1})\circ\cdots\circ(\pi_0\circ\delta_0)
  $。”
证书化要求为：
\[
  P(1)\ \text{成立},\qquad
  (\forall T\in\NN_{\ge 1})\ \bigl(P(T)\Rightarrow P(T+1)\bigr).
\]
并将 $\Delta_{\mathrm{P}},\mathcal{C}_{\mathrm{P}},\Delta_{\mathrm{E}},\mathcal{C}_{\mathrm{E}}$ 的“来源标签”视为每个阶段算子携带的可审计元数据：
外生来源 $\Rightarrow$ 进入 $\Delta_{\bullet}$，内生来源 $\Rightarrow$ 进入 $\mathcal{C}_{\bullet}$，从而由标签给出不交性。
\end{Certificate}

\begin{proof}[证明（数学归纳法证书；谓词因果链）]
由证书~\ref{cert:splitting-by-time-horizon}：
\begin{itemize}
  \item 基步：当 $T=1$ 时，$O_{0:0}=O_0=\pi_0\circ\delta_0$，故 $P(1)$ 成立；
  \item 归纳步：假设 $P(T)$ 成立，则
  \[
    O_{0:T}
    = O_T\circ O_{0:T-1}
    = (\pi_T\circ\delta_T)\circ(\pi_{T-1}\circ\delta_{T-1})\circ\cdots\circ(\pi_0\circ\delta_0),
  \]
  从而 $P(T+1)$ 成立。
\end{itemize}
由数学归纳法得对任意有限 $T$ 交错分解成立。

对二分式：按定义外生来源的阶段算子进入 $\Delta_{\bullet}$，内生来源的阶段算子进入 $\mathcal{C}_{\bullet}$。
令谓词链为
\[
\begin{aligned}
  Q_1 &: \mathrm{Src}(A)=\text{外生}\Rightarrow A\in\Delta_{\bullet},\\
  Q_2 &: \mathrm{Src}(A)=\text{内生}\Rightarrow A\in\mathcal{C}_{\bullet},\\
  Q_3 &: \mathrm{Src}(A)\ \text{外生与内生互斥}.
\end{aligned}
\]
由 $Q_1,Q_2,Q_3$ 得 $\Delta_{\bullet}\cap \mathcal{C}_{\bullet}=\varnothing$ 且 $\Delta_{\bullet}\uplus\mathcal{C}_{\bullet}
  $ 覆盖对应算子子族，故二分成立。证毕。
\end{proof}

\subsection{外界实时干扰算子：判据与两类实例}
\label{sec:disturbance-operators}

\begin{definition}[外界实时干扰算子（判据化定义）]
称算子 $\delta:\mathcal{S}\to\mathcal{S}$ 为外界实时干扰算子，若满足以下三条判据：
\begin{enumerate}
  \item 外生性 (Exogeneity): $\delta$ 的触发与参数主要由系统外部决定；
  \item 实时性 (Real-time): $\delta$ 在运行过程中随时间变化；
  \item 弱可控或不可控 (Weak controllability): 系统通常无法直接决定其取值，只能通过鲁棒性设计与在线适配抑制其影响。
\end{enumerate}
记所有满足判据的算子集合为 $\Delta$。
\end{definition}

\begin{lemma}[统一表示：加性/乘性/切换扰动可归入同一外界算子口径]
\label{lem:disturbance-unified-representation}
若物理/工程系统中的外界扰动以如下三类形式出现：
\begin{enumerate}
  \item 加性扰动：$s' = F(s) + d(t)$；
  \item 乘性/参数扰动：$s' = F(s;p(t))$；
  \item 模式切换扰动：$s' = F_{m(t)}(s)$，
\end{enumerate}
则它们均可被表示为某个外界实时干扰算子 $\delta_t\in\Delta$ 对状态的作用：$s_{t+1}=\delta_t(s_t)$。
\end{lemma}

\begin{Certificate}[证书：三类扰动的等价封装]
\label{cert:disturbance-unified-representation}
证书只登记一个事实：上述三类扰动都可通过扩充外界参数空间 $\Xi$ 统一为
\[
  \delta_{\xi}:\mathcal{S}\to\mathcal{S},\qquad
  s\mapsto \delta_{\xi}(s),
\]
其中 $\xi\in\Xi$ 记录 $d(t)$、$p(t)$ 或 $m(t)$ 的当期取值。
\end{Certificate}

\begin{proof}[证明（谓词因果链）]
由证书~\ref{cert:disturbance-unified-representation}，令谓词链为：
\[
\begin{aligned}
  R_1 &: \text{存在统一参数空间 }\Xi,\\
  R_2 &: \xi\ \text{可分别承载 }d(t),p(t),m(t),\\
  R_3 &: \text{据此定义 }\delta_{\xi}\ \text{使其复现三类更新式}.
\end{aligned}
\]
由 $R_1\wedge R_2\Rightarrow R_3$，得三类扰动均可写为外界算子作用，从而引理成立。证毕。
\end{proof}

\begin{proposition}[\PTOPB{} 与 \ETOPB{} 中的干扰算子分化（投影主导口径）]
\label{prop:disturbance-ptopb-etopb}
在总体状态 $s=(x,z)$ 的分解下：
\begin{enumerate}
  \item \PTOPB{} 的外界实时干扰算子 $\delta_{\mathrm{P}}\in\Delta_{\mathrm{P}}$ 以作用于物理态 $x$ 为主，并可通过物理效应诱导工程态 $z$ 的变化；
  \item \ETOPB{} 的外界实时干扰算子 $\delta_{\mathrm{E}}\in\Delta_{\mathrm{E}}$ 以作用于工程态 $z$ 为主，并通过闭环削弱对物理态 $x$ 的控制能力。
\end{enumerate}
\end{proposition}

\begin{Certificate}[证书：投影主导的判定规则]
\label{cert:disturbance-ptopb-etopb}
证书登记如下判定口径：
\[
  \delta\in\Delta_{\mathrm{P}}
  \Longleftrightarrow
  \pi_{\mathrm{P}}\circ \delta\ \text{承载主要外生变化},
  \qquad
  \delta\in\Delta_{\mathrm{E}}
  \Longleftrightarrow
  \pi_{\mathrm{E}}\circ \delta\ \text{承载主要外生变化}.
\]
其中“主要”由工程侧的可观测量与审计字段定义（例如传感器/网络/保障链路的事件流证据）。
\end{Certificate}

\begin{proof}[证明（谓词因果链）]
由证书~\ref{cert:disturbance-ptopb-etopb}，令谓词链为：
\[
  S_1:\ \pi_{\mathrm{P}}\circ\delta\ \text{承载主要外生变化}\Rightarrow \delta\in\Delta_{\mathrm{P}},
  \qquad
  S_2:\ \pi_{\mathrm{E}}\circ\delta\ \text{承载主要外生变化}\Rightarrow \delta\in\Delta_{\mathrm{E}}.
\]
由 $S_1,S_2$ 得命题两条结论。证毕。
\end{proof}

\begin{remark}[示例（非穷举）]
\textbf{\PTOPB{} 的外界实时干扰}可包括：风场突变/湍流/雨雪海况、噪声干扰与遮挡、多径与频谱拥挤、GNSS 受扰/欺骗、振动冲击/老化热漂移、对手物理施加的压制/诱偏/阻断/毁伤等。

\textbf{\ETOPB{} 的外界实时干扰}可包括：链路拥塞、时延抖动、丢包与带宽骤降、边缘算力被抢占、扫描/注入/拒绝服务/协议降级、备件到货延迟与补给受限、维护窗口压缩、误操作与协同方响应延迟等。
\end{remark}

\subsection{功能与保障的非交换构造算子：内生可控与顺序敏感}
\label{sec:construction-operators}

\begin{definition}[系统构造算子（内生可控）]
称算子 $\pi:\mathcal{S}\to\mathcal{S}$ 为系统构造算子，若其触发与参数主要由系统内部选择决定（设计/配置/运行策略），并在证书边界内可被审计与复核。
记系统构造算子集合为 $\mathcal{C}$。
\end{definition}

\begin{definition}[功能/保障二分（构造算子族内部结构）]
将构造算子族分为两类子族：
\[
  \mathcal{C} = \mathcal{C}^{\mathrm{fun}}\ \cup\ \mathcal{C}^{\mathrm{ass}},
\]
其中 $\mathcal{C}^{\mathrm{fun}}$ 表示功能类算子（使系统“能做事”的链路），$\mathcal{C}^{\mathrm{ass}}$ 表示保障类算子（使系统在扰动与对抗下“持续能做事”的链路）。
\end{definition}

\begin{definition}[非交换构造（顺序敏感）]
称构造算子族具有非交换构造特征，若存在 $\pi_A,\pi_B\in\mathcal{C}$ 使得
\[
  \pi_A\circ \pi_B \neq \pi_B\circ \pi_A.
\]
此时我们称 $\pi_A$ 与 $\pi_B$ 为顺序敏感的一对构造算子。
\end{definition}

\begin{theorem}[公理降级：信息依赖与资源竞争导致非交换构造是常态]
\label{thm:noncommutativity-is-default}
若系统构造链中存在以下任一结构：
\begin{enumerate}
  \item 信息依赖：后一步算子的定义域/参数依赖前一步的估计、标定或认证输出；
  \item 资源竞争：算力、带宽、电能、热预算等共享资源被不同模块争用且存在时序约束；
  \item 不可逆/不可回滚：某些动作执行后导致信息不可恢复或回滚代价不可忽略；
  \item 安全边界：认证/授权/审计的先后顺序改变系统风险边界；
  \item 时序窗口：对抗态势使延迟本身成为状态变量，
\end{enumerate}
则在证书边界内可构造一对构造算子 $\pi_A,\pi_B\in\mathcal{C}$ 使得 $\pi_A\circ \pi_B \neq \pi_B\circ \pi_A$。
\end{theorem}

\begin{Certificate}[证书：非交换性的可观测见证（交换子非零）]
\label{cert:noncommutativity-witness}
证书登记如下可观测见证口径：
存在状态 $s\in\mathcal{S}$ 使得
\[
  (\pi_A\circ \pi_B)(s)\neq (\pi_B\circ \pi_A)(s),
\]
从而交换子在该点的作用非零。该见证可由运行时日志/审计事件流或仿真对比复验给出。
\end{Certificate}

\begin{proof}[证明（谓词因果链）]
令谓词链为：
\[
\begin{aligned}
  T_1 &: \text{存在信息依赖/资源竞争/不可逆/安全边界/时序窗口之一},\\
  T_2 &: \Rightarrow\ \exists s\in\mathcal{S}\ \text{使得顺序改变}\\
      &\qquad \text{导致中间状态或资源余量不同},\\
  T_3 &: \Rightarrow\ \exists(\pi_A,\pi_B)\ \text{满足}\\
      &\qquad (\pi_A\circ \pi_B)(s)\neq(\pi_B\circ \pi_A)(s).
\end{aligned}
\]
由 $T_1\Rightarrow T_2\Rightarrow T_3$，并结合证书~\ref{cert:noncommutativity-witness} 的可复验见证口径，得定理结论。证毕。
\end{proof}

\begin{lemma}[非交换性与可达集合的顺序敏感]
\label{lem:reachable-set-order-sensitive}
若构造算子族中存在顺序敏感对 $(\pi_A,\pi_B)$，则对某些初始状态 $s_0$，由两步构造产生的可达集合
\[
  \mathcal{R}_2(s_0)
  :=
  \{(\pi_2\circ\pi_1)(s_0):\ \pi_1,\pi_2\in\mathcal{C}\}
\]
对“动作顺序”敏感：存在同一元素多重集合 $\{\pi_A,\pi_B\}$ 的两种排列，导致终态不同。
\end{lemma}

\begin{Certificate}[证书：顺序敏感对的可达差异]
\label{cert:reachable-set-order-sensitive}
证书登记：存在 $s_0$ 使得
\[
  (\pi_A\circ\pi_B)(s_0)\neq (\pi_B\circ\pi_A)(s_0),
\]
从而两条长度为 $2$ 的可行构造词产生不同终态。
\end{Certificate}

\begin{proof}[证明（谓词因果链）]
由证书~\ref{cert:reachable-set-order-sensitive}，存在 $s_0$ 使得两种排列的复合结果不同。
取谓词链：
\[
\begin{aligned}
  U_1 &: (\pi_A\circ\pi_B)(s_0)\neq (\pi_B\circ\pi_A)(s_0),\\
  U_2 &: \Rightarrow\ \text{同一多重集合的不同排列给出不同终态},\\
  U_3 &: \Rightarrow\ \mathcal{R}_2(s_0)\ \text{对顺序敏感}.
\end{aligned}
\]
由 $U_1\Rightarrow U_2\Rightarrow U_3$ 得结论。证毕。
\end{proof}

\begin{remark}[\PTOPB/\ETOPB{} 中的构造算子示例（非穷举）]
\textbf{\PTOPB{} 的构造算子}可分为：
\begin{itemize}
  \item 功能类：传感器布置与选择、执行机构配置、结构与材料方案、能源链路设计等；
  \item 保障类：冗余与隔离、故障检测/隔离/重构（FDIR）、抗扰裕度设计、降级曲线与安全模式等。
\end{itemize}

\textbf{\ETOPB{} 的构造算子}可分为：
\begin{itemize}
  \item 功能类：信息链（采集-同步-清洗-融合-识别-决策）、控制链（建模-估计-预测-控制律-调度执行）、通信链（编码/路由/拥塞控制/重传）、任务链（分解-资源分配-编排-一致性）；
  \item 保障类：安全算子（认证/授权/密钥/完整性/审计）、鲁棒/自适应（在线辨识/自整定/策略切换）、弹性（熔断/隔离/限流/快速恢复/灰度回退）、运维保障（监测告警/健康管理/补丁策略/备份恢复演练）等。
\end{itemize}
\end{remark}

\subsection{博弈对抗语境：外生干扰与内生构造的夹逼形式}
\label{sec:game-adversarial}

\begin{definition}[对抗口径：干扰序列与构造序列]
设 $\delta_{0:T-1}$ 为外界干扰序列（由环境或对手选择），$\pi_{0:T-1}$ 为系统构造序列（由我方选择）。
给定目标泛函 $J$（效能、损失、资源消耗、暴露风险等），可将对抗问题抽象为
\[
  \max_{\pi_{0:T-1}\in\mathcal{C}^{T}}
  \ \min_{\delta_{0:T-1}\in\Delta^{T}}
  J(\pi_{0:T-1},\delta_{0:T-1}).
\]
\end{definition}

\begin{proposition}[非交换构造使“顺序”成为对抗中的结构变量]
\label{prop:order-is-structural-in-adversarial}
若存在顺序敏感对 $(\pi_A,\pi_B)\in\mathcal{C}^2$，则对抗优化中必须把顺序作为结构变量登记；
忽略顺序（将动作集误当作交换集合）将导致可达集合与风险边界的系统性误判。
\end{proposition}

\begin{Certificate}[证书：顺序敏感导致风险边界差异]
\label{cert:order-is-structural}
证书登记：存在扰动序列 $\delta_{0:T-1}$ 与初始状态 $s_0$，使得在相同的多重集合动作下，不同顺序产生不同终态，
从而导致 $J$ 的取值不同：
\[
  J(\cdots,\pi_A\circ\pi_B,\cdots;\delta_{0:T-1})
  \neq
  J(\cdots,\pi_B\circ\pi_A,\cdots;\delta_{0:T-1}).
\]
\end{Certificate}

\begin{proof}[证明（谓词因果链）]
由证书~\ref{cert:order-is-structural}，存在同一动作多重集合的两种排列导致 $J$ 值不同。
令谓词链为：
\[
\begin{aligned}
  V_1 &: \exists s_0,\delta_{0:T-1}\ \text{使得不同顺序给出不同终态},\\
  V_2 &: \Rightarrow\ J\ \text{取值不同},\\
  V_3 &: \Rightarrow\ \text{顺序必须作为结构变量进入优化与审计}.
\end{aligned}
\]
由 $V_1\Rightarrow V_2\Rightarrow V_3$ 得命题结论。证毕。
\end{proof}

\begin{corollary}[顺序敏感是“非交换构造”建模的必要入口]
\label{cor:noncommutative-modeling-necessary}
在存在外界干扰与内生构造的对抗系统中，若构造链条满足顺序敏感性，则“非交换构造”不是修辞而是必要的最弱建模入口：
它保证对抗分析在最弱形式下仍能区分不同策略链条。
\end{corollary}

\begin{Certificate}[证书：最弱入口的区分性]
\label{cert:noncommutative-modeling-necessary}
证书登记：存在两条策略链条（同一动作多重集合、不同排列）在可复验样例上产生不同结果；
因此若将算子交换化，则该差异会在模型中被抹平。
\end{Certificate}

\begin{proof}[证明（谓词因果链）]
由证书~\ref{cert:noncommutative-modeling-necessary}，存在两条策略链条在可复验样例上产生不同结果。
令谓词链为：
\[
\begin{aligned}
  W_1 &: \exists\ \text{两条同多重集合、不同排列的策略链条结果不同},\\
  W_2 &: \Rightarrow\ \text{交换化会抹平该差异},\\
  W_3 &: \Rightarrow\ \text{非交换构造是保持区分性的最弱入口}.
\end{aligned}
\]
由 $W_1\Rightarrow W_2\Rightarrow W_3$ 得推论结论。证毕。
\end{proof}

\subsection{算子清单模板：证书化登记与可复用算子库}
\label{sec:operator-library-template}

\begin{definition}[算子条目（可审计记录模板）]
为将 $\Delta$ 与 $\mathcal{C}$ 落到可维护的算子库中，对每个算子条目建议登记如下字段：
\begin{enumerate}
  \item 作用域：对 $x$ 为主 / 对 $z$ 为主 / 强耦合；
  \item 输入输出：$(x,z)\mapsto(x',z')$ 的哪些分量被改变；
  \item 触发方式：外生（环境/对手）或内生（策略/调度）；
  \item 时效：连续/离散，更新频率，延迟窗口；
  \item 资源代价：功耗、算力、带宽、时延、热预算等；
  \item 可观测性：能否检测其发生，检测置信度与证据来源；
  \item 可逆性：是否可回滚，回滚代价与风险；
  \item 交换性标注：与哪些算子存在强非交换关系（顺序敏感清单或交换子见证）。
\end{enumerate}
该记录既可用于工程文档，也可作为后续证书系统的最小条目载体。
\end{definition}

\begin{proposition}[算子库条目可作为 \ETOPB{} 的证书化边界接口]
\label{prop:operator-library-as-certificate-interface}
若对每个算子条目均保留其触发方式、证据来源与顺序敏感标注，
则算子库可作为 \ETOPB{} 的证书化边界接口：它把“哪些算子可用、何时可用、如何复验”从口头解释降级为可追溯记录。
\end{proposition}

\begin{Certificate}[证书：条目字段 $\Rightarrow$ 可追溯调用边界]
\label{cert:operator-library-as-certificate-interface}
证书登记如下蕴含链：
\[
  (\text{触发方式}+\text{证据来源}+\text{顺序敏感标注})
  \Rightarrow
  (\text{可复验调用条件}+\text{审计边界}+\text{回归基线}).
\]
其中“证据来源”包括运行时日志、仿真对比、试验记录与供应链/运维事件流等。
\end{Certificate}

\begin{proof}[证明（谓词因果链）]
由证书~\ref{cert:operator-library-as-certificate-interface}，令谓词链为：
\[
\begin{aligned}
  X_1 &: \text{触发方式与证据来源可定位责任域与可复验数据},\\
  X_2 &: \text{顺序敏感标注可定位非交换风险边界},\\
  X_3 &: X_1\wedge X_2\Rightarrow \text{形成可复验调用条件与审计边界}.
\end{aligned}
\]
由 $X_1,X_2,X_3$ 得命题结论。证毕。
\end{proof}

\subsection{进一步细化：稳定性分档与功能/保障二级目录}
\label{sec:finer-taxonomy}

\begin{remark}[干扰算子可按可预测性/统计稳定性再分层]
在不改变 $\Delta$ 的外生定义前提下，可将外界干扰按统计稳定性分为稳态扰动（可标定、可建模）与非稳态扰动（突变、对手自适应）。
该分层的工程意义在于：前者对应“提高鲁棒裕度与在线估计”，后者对应“引入对抗检测与策略切换证书”。
\end{remark}

\begin{remark}[构造算子建议固化为“功能/保障”二级目录]
在不改变 $\mathcal{C}$ 的内生定义前提下，把构造算子固化为功能/保障二级目录，有利于方案评审时回答两个关键问题：
\begin{enumerate}
  \item 系统在哪些外界扰动下会掉出可控域（$\Delta$ 覆盖性）？
  \item 哪个保障算子能把系统拉回可控域，代价是多少，顺序是否敏感（$\mathcal{C}^{\mathrm{ass}}$ 的有效性与非交换性）？
\end{enumerate}
该目录化同时也为后续“算子最小生成元表”与“治理证书链条”提供可维护入口。
\end{remark}

\section{工程增强与对抗可消/不可消：上同调/同伦形式系统}
\label{sec:etopb-augmentation-cohomology-homotopy}

\begin{remark}[定位（把“工程增强”与“可消/不可消”统一到同一套不变量语言）]
本节把 \ETOPB{} 视为 \PTOPB{} 的工程增强，并把“对抗博弈中的可消/不可消干扰”解释为上同调/同伦意义下的
\textbf{上边界（可消）}与\textbf{非平凡上同调类（不可消）}。
其目的不是引入玄学比喻，而是给出一套可被\textbf{证书化登记、可复验压缩、可回归审计}的形式化接口。
\end{remark}

\subsection{将 \ETOPB{} 视为 \PTOPB{} 的增强：扩张与长正合列}
\label{sec:etopb-as-ptopb-augmentation}

\begin{definition}[上同调接口（记号约定）]
为将 \PTOPB/\ETOPB{} 置入同一套“不变量”口径，本节把它们视为某类可计算对象（例如共链复形、带微分的模、或层/系的截面复形），
并用 $H^k(-)$ 表示其第 $k$ 阶上同调（或同调型不变量）：
\[
  H^k(\mathcal{E})
  := H^k\bigl(C^\bullet(\mathcal{E})\bigr).
\]
其中 $C^\bullet(\mathcal{E})$ 的选取、微分与可计算实现由工程侧证书给出；本节仅使用其满足的函子性质（短正合列诱导长正合列）。
\end{definition}

\begin{definition}[\ETOPB{} 对 \PTOPB{} 的增强扩张（短正合列口径）]
称 \ETOPB{} 是 \PTOPB{} 的增强（augmentation），若存在增广项 \Aug{} 与态射 $i,\pi$ 使得在上述对象范畴中成立短正合列：
\[
  0
  \longrightarrow
  \PTOPB
  \xrightarrow{i}
  \ETOPB
  \xrightarrow{\pi}
  \Aug
  \longrightarrow
  0.
\]
其中 $i$ 表示“增强不丢物理”（物理算子侧嵌入），$\pi$ 表示“抽取纯增广部分”（信息/控制/网络/保障/运维等相对纯物理新增结构）。
\end{definition}

\begin{theorem}[定理：短正合列诱导长正合列（增强的不变量审计接口）]
\label{thm:short-exact-implies-long-exact}
若 \ETOPB{} 对 \PTOPB{} 的增强扩张满足短正合列
\[
  0\to\PTOPB\xrightarrow{i}\ETOPB\xrightarrow{\pi}\Aug\to 0,
\]
则存在一族连接同态 $\partial^k:H^k(\Aug)\to H^{k+1}(\PTOPB)$ 使得诱导长正合列：
\[
  \cdots
  \to H^k(\PTOPB)
  \xrightarrow{i_*}
  H^k(\ETOPB)
  \xrightarrow{\pi_*}
  H^k(\Aug)
  \xrightarrow{\partial^k}
  H^{k+1}(\PTOPB)
  \to \cdots .
\]
因此：
\begin{enumerate}
  \item \ETOPB{} 相对 \PTOPB{} 的新增不可忽略结构，集中体现在 $H^k(\Aug)$（及其与 $\partial^k$ 的耦合）中；
  \item 若在关心的阶数 $k$ 上 $i_*:H^k(\PTOPB)\to H^k(\ETOPB)$ 为同构（或准同构意义下的同构），则增强在该阶不改变对应不变量。
\end{enumerate}
\end{theorem}

\begin{Certificate}[证书：短正合列 $\Rightarrow$ 长正合列]
\label{cert:short-exact-implies-long-exact}
证书登记如下蕴含链（同调代数标准接口）：
\[
\begin{aligned}
  Z_1 &: 0\to \PTOPB \xrightarrow{i} \ETOPB \xrightarrow{\pi} \Aug \to 0\ \text{短正合列}\\
  Z_2 &: \Rightarrow\ \exists\,\partial^k:H^k(\Aug)\to H^{k+1}(\PTOPB)\\
  Z_3 &: \Rightarrow\ \cdots\to H^k(\PTOPB)\to H^k(\ETOPB)\to H^k(\Aug)\\
      &\qquad \to H^{k+1}(\PTOPB)\to\cdots\ \text{长正合列}.
\end{aligned}
\]
\end{Certificate}

\begin{proof}[证明（谓词因果链）]
由证书~\ref{cert:short-exact-implies-long-exact}：$Z_1$ 给出对象级短正合列。
在可计算实现层（共链复形口径）上，$Z_1$ 可被落实为共链短正合列；
对其应用标准连接同态构造（蛇形引理/连接态射口径）得到 $Z_2$；
并由连接同态的定义与函子性得到 $Z_3$ 的精确性（$\operatorname{im}=\ker$）。
故定理成立。证毕。
\end{proof}

\begin{definition}[分层增强（有限扩张过滤）]
\label{def:layered-augmentation-filtration}
称 \ETOPB{} 是 \PTOPB{} 的 $N$ 层增强，若存在对象链
\[
  \mathcal{E}_0 \subseteq \mathcal{E}_1 \subseteq \cdots \subseteq \mathcal{E}_N,
  \qquad
  \mathcal{E}_0=\PTOPB,\ \mathcal{E}_N=\ETOPB,
\]
以及每一层的短正合列
\[
  0\longrightarrow \mathcal{E}_{j-1}\longrightarrow \mathcal{E}_j \longrightarrow \Aug_j \longrightarrow 0,
  \qquad j=1,\dots,N.
\]
其中每个 $\Aug_j$ 表示第 $j$ 层工程增广（例如通信层、控制层、保障层、运维层等可审计分层）。
\end{definition}

\begin{theorem}[公理降级：分层增强可归约为逐层证书链（按层数 $N$ 归纳）]
\label{thm:layered-augmentation-by-induction}
若存在分层增强过滤（定义~\ref{def:layered-augmentation-filtration}），
则对任意阶数 $k$ 与任意层 $j$，均存在可复验的逐层长正合列片段：
\[
  \cdots
  \to H^k(\mathcal{E}_{j-1})
  \to H^k(\mathcal{E}_j)
  \to H^k(\Aug_j)
  \to H^{k+1}(\mathcal{E}_{j-1})
  \to \cdots .
\]
从而“\ETOPB{} 相对 \PTOPB{} 的新增结构”可被拆解为 $N$ 个可追溯的增广层条目，而不是一次性黑盒差分。
\end{theorem}

\begin{Certificate}[数学归纳法证书：按增强层数 $N$ 归纳]
\label{cert:layered-augmentation-by-induction}
证书登记谓词 $P(N)$：
\[
  P(N):\ \text{任意 $N$ 层增强过滤都能对每层给出长正合列片段}.
\]
证书化要求为：
\[
  P(1)\ \text{成立},\qquad
  (\forall N\in\NN_{\ge 1})\ \bigl(P(N)\Rightarrow P(N+1)\bigr).
\]
\end{Certificate}

\begin{proof}[证明（数学归纳法证书；谓词因果链）]
由证书~\ref{cert:layered-augmentation-by-induction}：
\begin{itemize}
  \item 基步（$N=1$）：此时只有一条短正合列 $0\to\PTOPB\to\ETOPB\to\Aug_1\to 0$；
  由定理~\ref{thm:short-exact-implies-long-exact} 得到对应长正合列片段，故 $P(1)$ 成立。
  \item 归纳步：假设 $P(N)$ 成立。对任意 $(N+1)$ 层过滤，
  前 $N$ 层由归纳假设分别给出逐层长正合列片段；
  最后一层满足短正合列 $0\to \mathcal{E}_N\to \mathcal{E}_{N+1}\to \Aug_{N+1}\to 0$，
  再次应用定理~\ref{thm:short-exact-implies-long-exact} 得到第 $(N+1)$ 层的长正合列片段。
  因此 $P(N+1)$ 成立。
\end{itemize}
由数学归纳法得对任意有限层数 $N$ 命题成立。证毕。
\end{proof}

\subsection{实时干扰算子同伦类：内外统一的负算子口径}
\label{sec:disturbance-homotopy-class}

\begin{definition}[实时干扰算子（负算子口径）]
令 $J:\mathcal{S}\to\RR$ 为损失/风险泛函，并令 $\mathcal{K}\subseteq\mathcal{S}$ 为关键约束域（稳定性、安全性、资源、可审计性等约束共同定义）。
称算子 $d:\mathcal{S}\to\mathcal{S}$ 为实时干扰算子，若满足：
\begin{enumerate}
  \item 负效应：存在 $s\in\mathcal{K}$ 使 $J(d(s))>J(s)$；
  \item 约束风险：或存在 $s\in\mathcal{K}$ 使 $d(s)\notin\mathcal{K}$；
  \item 实时性：$d$ 的参数可随时间变化（可视为 $d_t$ 的族）。
\end{enumerate}
记所有实时干扰算子的集合为 $\DeltaRT$。
其中前文的“外界实时干扰”集合 $\Delta$ 可视为 $\DeltaRT$ 的外生子集：$\Delta\subseteq \DeltaRT$。
\end{definition}

\begin{definition}[干扰同伦等价（不穿越关键约束边界）]
对 $d_0,d_1\in\DeltaRT$，称 $d_0\sim d_1$，若存在连续族 $\{d_\lambda\}_{\lambda\in[0,1]}\subseteq\DeltaRT$ 使得：
\[
  d_{\lambda=0}=d_0,\qquad d_{\lambda=1}=d_1,
\]
并且该变形在整个过程中不穿越关键约束边界（即对 $\lambda\in[0,1]$ 的每一步均保持在允许域内）。
记等价类为 $[d]\in \DeltaRT/\!\sim$。
\end{definition}

\begin{lemma}[干扰同伦等价 $\sim$ 是等价关系]
\label{lem:disturbance-homotopy-is-equivalence}
关系 $\sim$ 在 $\DeltaRT$ 上满足自反、对称、传递。
\end{lemma}

\begin{Certificate}[证书：等价关系三性]
\label{cert:disturbance-homotopy-is-equivalence}
证书登记如下谓词链：
\[
\begin{aligned}
  E_1 &: d\sim d\ \text{（取常值同伦）},\\
  E_2 &: d_0\sim d_1\Rightarrow d_1\sim d_0\ \text{（反向参数）},\\
  E_3 &: (d_0\sim d_1)\wedge(d_1\sim d_2)\Rightarrow d_0\sim d_2\ \text{（拼接同伦）}.
\end{aligned}
\]
\end{Certificate}

\begin{proof}[证明（谓词因果链）]
由证书~\ref{cert:disturbance-homotopy-is-equivalence}：
$E_1$ 由 $d_\lambda\equiv d$ 给出；
$E_2$ 由 $d'_\lambda:=d_{1-\lambda}$ 给出；
$E_3$ 由把两段同伦在 $\lambda=1/2$ 处拼接给出。
故 $\sim$ 为等价关系，引理成立。证毕。
\end{proof}

\begin{proposition}[换招不换势：同伦类压缩是对抗分析的最小充分口径]
\label{prop:homotopy-class-compression}
令 $I:\DeltaRT\to \mathcal{V}$ 为任意“同伦不变量”（满足 $d_0\sim d_1\Rightarrow I(d_0)=I(d_1)$）。
则 $I$ 可通过商映射因子化为
\[
  I = \overline{I}\circ q,\qquad
  q:\DeltaRT\to \DeltaRT/\!\sim,
\]
从而对抗分析无需枚举每个具体干扰手段 $d$，只需在同伦类 $[d]$ 上登记与推演。
\end{proposition}

\begin{Certificate}[证书：同伦不变量因子化]
\label{cert:homotopy-class-compression}
证书登记：若 $d_0\sim d_1\Rightarrow I(d_0)=I(d_1)$，则存在唯一映射 $\overline{I}$ 使 $I=\overline{I}\circ q$。
\end{Certificate}

\begin{proof}[证明（谓词因果链）]
由证书~\ref{cert:homotopy-class-compression}，令谓词链为：
\[
\begin{aligned}
  F_1 &: d_0\sim d_1\Rightarrow I(d_0)=I(d_1),\\
  F_2 &: \Rightarrow\ I\ \text{在每个等价类上取常值},\\
  F_3 &: \Rightarrow\ \exists\,\overline{I}\ \text{使}\ I=\overline{I}\circ q.
\end{aligned}
\]
由 $F_1\Rightarrow F_2\Rightarrow F_3$ 得结论。证毕。
\end{proof}

\subsection{功能+保障构造的同伦类与代表元最优化}
\label{sec:construction-homotopy-and-representative-optimization}

\begin{definition}[构造同伦类（功能+保障的结构等价）]
令 $\mathcal{C}$ 为系统构造算子集合（包含功能链与保障链及其编排/调度/切换），并令 $\mathcal{K}$ 为关键约束域。
对 $c_0,c_1\in\mathcal{C}$，称 $c_0\simeq c_1$，若存在连续族 $\{c_\lambda\}_{\lambda\in[0,1]}\subseteq\mathcal{C}$ 使
$c_{\lambda=0}=c_0,c_{\lambda=1}=c_1$ 且变形过程不穿越关键约束边界（始终保持在允许域内）。
记构造同伦类为 $[c]\in \mathcal{C}/\!\simeq$。
\end{definition}

\begin{definition}[代表元最优化（在同伦类内选最优实现）]
给定干扰同伦类 $[d]\in\DeltaRT/\!\sim$ 与构造同伦类 $[c]\in\mathcal{C}/\!\simeq$，
“在线调优/重构”的形式化目标可记为在同伦类内选取最优代表元：
\[
  c^\star \in [c],
  \qquad
  c^\star \in \arg\min_{c\in[c]} J(c,[d]).
\]
其中 $J(c,[d])$ 表示在干扰类 $[d]$ 下、采用构造实现 $c$ 的等效损失/风险指标（由证书化评测流程定义）。
\end{definition}

\begin{proposition}[同伦类代表元选择是“可调参/可重构/可编排”的统一抽象]
\label{prop:representative-optimization-unified-abstract}
在不穿越关键约束边界的前提下，常见工程操作（连续调参、阈值移动、策略切换、资源重分配、编排拓扑等价改写）均可被视为
在固定的构造同伦类 $[c]$ 中选取不同代表元 $c$；
因此最优化选择问题发生在“同伦类的代表元空间”而非“单点实现”。
\end{proposition}

\begin{Certificate}[证书：代表元选择的审计边界]
\label{cert:representative-optimization-audit}
证书登记：只要变更路径 $\{c_\lambda\}$ 保持在允许域内（不穿越关键约束边界），则其均属于同一同伦类 $[c]$，
并可被同一套审计字段与回归基线复验；在线选择仅改变代表元而不改变结构类标签。
\end{Certificate}

\begin{proof}[证明（谓词因果链）]
由证书~\ref{cert:representative-optimization-audit}，令谓词链为：
\[
\begin{aligned}
  G_1 &: \{c_\lambda\}\subseteq\mathcal{C}\ \text{且不穿越边界}\Rightarrow c_0\simeq c_1,\\
  G_2 &: \Rightarrow\ c_0,c_1\ \text{同属}\ [c],\\
  G_3 &: \Rightarrow\ \text{可用同一审计字段/回归基线复验}.
\end{aligned}
\]
由 $G_1\Rightarrow G_2\Rightarrow G_3$ 得结论。证毕。
\end{proof}

\subsection{上同调口径：可消干扰是上边界，不可消干扰落在上同调类中}
\label{sec:cohomology-clear-cancel-noncancel}

\begin{definition}[两端复形：构造调整 $\to$ 干扰 $\to$ 越界见证]
为刻画“可消/不可消”，引入一个最小共链框架：
\begin{enumerate}
  \item $C^0$：构造调整空间（功能+保障的允许改动方向/改动原语）；
  \item $C^1$：干扰空间（实时干扰算子/干扰等效项）；
  \item $C^2$：约束越界见证空间（稳定/安全/资源/可审计性越界的证书化见证）。
\end{enumerate}
并给定共边界算子
\[
  \partial^0:C^0\to C^1,\qquad
  \partial^1:C^1\to C^2,
\]
满足复形条件 $\partial^1\circ \partial^0 = 0$。
直观上：$\partial^0$ 把“构造改动”映射为其在等效层面对干扰的补偿/抵消项；
$\partial^1$ 把“干扰项”映射为其触发越界的见证。
\end{definition}

\begin{definition}[可消/不可消（上同调判据）]
定义第一上同调
\[
  H^1 := \ker(\partial^1)\big/\operatorname{im}(\partial^0).
\]
对任意闭合干扰 $d\in\ker(\partial^1)$：
\begin{enumerate}
  \item 若存在 $c\in C^0$ 使 $d=\partial^0(c)$，则称 $d$ \textbf{可消}（$d$ 为上边界）；
  \item 若 $[d]\neq 0$ 在 $H^1$ 中，则称 $d$ \textbf{不可消}（结构性剩余干扰）。
\end{enumerate}
\end{definition}

\begin{theorem}[定理：可消干扰 $\Leftrightarrow$ 上同调类平凡]
\label{thm:cancellable-iff-coboundary}
对任意 $d\in\ker(\partial^1)$，以下命题等价：
\[
  d\ \text{可消}
  \quad\Longleftrightarrow\quad
  [d]=0\ \text{于}\ H^1.
\]
\end{theorem}

\begin{Certificate}[证书：上边界判据]
\label{cert:cancellable-iff-coboundary}
证书登记：$[d]=0$ 当且仅当 $d\in \operatorname{im}(\partial^0)$；
而“可消”的定义恰为 $d=\partial^0(c)$ 的存在性。
\end{Certificate}

\begin{proof}[证明（谓词因果链）]
由证书~\ref{cert:cancellable-iff-coboundary}，令谓词链为：
\[
  H_1:\ d\ \text{可消}\Rightarrow (\exists c\in C^0)\ d=\partial^0(c)\Rightarrow d\in\operatorname{im}(\partial^0)
    \Rightarrow [d]=0,
\]
\[
  H_2:\ [d]=0\Rightarrow d\in\operatorname{im}(\partial^0)\Rightarrow (\exists c\in C^0)\ d=\partial^0(c)\Rightarrow d\
    \text{可消}.
\]
由 $H_1,H_2$ 得等价性成立。证毕。
\end{proof}

\begin{lemma}[引理：非交换结构通过交换子进入共边界算子]
\label{lem:commutator-enters-coboundary}
在顺序敏感/非交换构造场景中，共边界算子 $\partial^0$ 可包含交换子项作为可观测见证：
\[
  \partial^0(c)\ \sim\ \partial_t c + [A,c] + \cdots,
  \qquad
  [A,c] := A\circ c - c\circ A.
\]
因此 $[A,c]\neq 0$ 可作为“顺序敏感耦合导致补偿项改变”的最小见证。
\end{lemma}

\begin{Certificate}[证书：交换子见证非交换耦合]
\label{cert:commutator-enters-coboundary}
证书登记：若 $A\circ c\neq c\circ A$，则交换子 $[A,c]\neq 0$，从而 $\partial^0(c)$ 的值随顺序结构改变而改变。
\end{Certificate}

\begin{proof}[证明（谓词因果链）]
由证书~\ref{cert:commutator-enters-coboundary}，令谓词链为：
\[
\begin{aligned}
  K_1 &: A\circ c\neq c\circ A\Rightarrow [A,c]\neq 0,\\
  K_2 &: [A,c]\neq 0\Rightarrow \partial^0(c)\ \text{含非交换耦合贡献},\\
  K_3 &: \Rightarrow\ \text{补偿/抵消项对顺序敏感}.
\end{aligned}
\]
由 $K_1\Rightarrow K_2\Rightarrow K_3$ 得结论。证毕。
\end{proof}

\begin{corollary}[推论：非平凡上同调类刻画结构性短板（同伦内调优无法根除）]
\label{cor:nontrivial-cohomology-is-structural}
若某类闭合干扰 $d\in\ker(\partial^1)$ 在 $H^1$ 中满足 $[d]\neq 0$，
则在给定的构造调整空间 $C^0$（即给定的“功能+保障结构类”）内不存在 $c$ 使其被完全消去；
此时若要消除该类压力，必须改变增强项（\Aug{} 的结构/边界）或改变物理约束（\PTOPB{} 的底层可达域），
对应于定理~\ref{thm:short-exact-implies-long-exact} 所刻画的“跨阶映射/连接同态”处的结构变更。
\end{corollary}

\begin{Certificate}[证书：不可消 $\Rightarrow$ 需要改变结构类而非仅调参]
\label{cert:nontrivial-cohomology-is-structural}
证书登记：$[d]\neq 0\Rightarrow d\notin\operatorname{im}(\partial^0)$（由定义），
因此同伦内的代表元选择无法产生 $d=\partial^0(c)$；要使该类消失只能改变 $C^0$ 或改变干扰闭合条件/约束域（即改变结构类边界）。
\end{Certificate}

\begin{proof}[证明（谓词因果链）]
由证书~\ref{cert:nontrivial-cohomology-is-structural}，令谓词链为：
\[
\begin{aligned}
  L_1 &: [d]\neq 0\Rightarrow d\notin\operatorname{im}(\partial^0),\\
  L_2 &: \Rightarrow\ \nexists c\in C^0\ \text{使}\ d=\partial^0(c),\\
  L_3 &: \Rightarrow\ \text{仅在同一结构类内调参/改写无法根除该压力}.
\end{aligned}
\]
由 $L_1\Rightarrow L_2\Rightarrow L_3$ 得推论。证毕。
\end{proof}

\subsection{工程对抗的直接产出：归并、选代表元、识别短板}
\label{sec:engineering-outcomes}

\begin{remark}[三类直接产出（可落地到算子库与证书系统）]
在上述口径下，可直接得到三类工程化产出：
\begin{enumerate}
  \item \textbf{干扰归并（等价类压缩）}：用 $[d]\in\DeltaRT/\!\sim$ 代替无穷多具体手段 $d$，把“换招不换势”升级为可审计标签；
  \item \textbf{构造归并（方案族而非单点）}：把可调参/可重构/可编排统一成在 $[c]\in\mathcal{C}/\!\simeq$ 内选代表元 $c^\star$；
  \item \textbf{短板识别（上同调非平凡项）}：用 $H^1$ 的非平凡类定位“结构性不可消压力”，区分“调参可改进”与“必须改结构”。
\end{enumerate}
\end{remark}

\subsection{相对上同调：用差分不变量刻画增强项的本质贡献}
\label{sec:relative-cohomology-augmentation}

\begin{definition}[相对上同调（增强差分口径）]
在共链实现口径下，定义相对上同调为商对象的上同调：
\[
  H^k(\ETOPB,\PTOPB) := H^k(\ETOPB/\PTOPB).
\]
其工程含义是：只度量“相对纯物理新增的那部分结构”，把共同部分（\PTOPB{}）从指标中消去。
\end{definition}

\begin{theorem}[定理：相对上同调在增强短正合列下等价于增广项上同调]
\label{thm:relative-cohomology-is-aug}
若存在增强短正合列 $0\to\PTOPB\to\ETOPB\to\Aug\to 0$，则存在自然同构
\[
  H^k(\ETOPB,\PTOPB) \cong H^k(\Aug).
\]
\end{theorem}

\begin{Certificate}[证书：商对象同构 $\Rightarrow$ 上同调同构]
\label{cert:relative-cohomology-is-aug}
证书登记：由短正合列得 $\ETOPB/\PTOPB \cong \Aug$；对同构应用 $H^k(-)$ 得 $H^k(\ETOPB/\PTOPB)\cong H^k(\Aug)$。
\end{Certificate}

\begin{proof}[证明（谓词因果链）]
由证书~\ref{cert:relative-cohomology-is-aug}，令谓词链为：
\[
\begin{aligned}
  R_1 &: 0\to\PTOPB\to\ETOPB\to\Aug\to 0\Rightarrow \ETOPB/\PTOPB \cong \Aug,\\
  R_2 &: \Rightarrow\ H^k(\ETOPB/\PTOPB)\cong H^k(\Aug),\\
  R_3 &: \Rightarrow\ H^k(\ETOPB,\PTOPB)\cong H^k(\Aug).
\end{aligned}
\]
由 $R_1\Rightarrow R_2\Rightarrow R_3$ 得结论。证毕。
\end{proof}

\begin{remark}[下一步（从抽象到可计算）]
若要把本节框架进一步落到可计算/可推演层面，关键是固定一个具体对抗域并给出证书化的 $C^\bullet$ 与 $\partial$：
例如通信压制、导航欺骗、分布式协同破坏、供应链/保障链受扰、网络与控制耦合攻击等。
在该锚点下，$\DeltaRT/\!\sim$ 的压缩、$[c]$ 的代表元选择、以及 $H^1$ 的非平凡类定位
即可被固化为“干扰库/构造库/证书链条”的可回归工具。
\end{remark}

\section{两层增强：纤维内优化与包线迁移（\texorpdfstring{$\PTOPB(\lambda)$}{PTOPB(lambda)} 族口径）}
\label{sec:two-level-augmentation}

\begin{remark}[核心陈述（两层增强的精确定义）]
更精确地说，“\ETOPB{} 与 \PTOPB{} 可视为 \PTOPB{} 增强”可以同时包含两种不同层级的增强：
\begin{enumerate}
  \item \textbf{不突破 \PTOPB{} 边界的增强}：\PTOPB{} 的定义域/可达域/工作包线不变，\ETOPB{} 只是在同一物理包线内，通过工程构造把性能、鲁棒性与可验证性推到更优；
  \item \textbf{突破 \PTOPB{} 边界的增强}：\ETOPB{} 不仅在既有 \PTOPB{} 的纤维内做优化，还通过“材料—工艺—结构—环境”等构造，把系统迁移到新的物理包线里，从而得到一个新的物理总算子丛
    $\PTOPB^{\prime}$（更大的定义域或不同的约束结构）。
\end{enumerate}
第二类增强并非“违背物理定律”，而是改变了被建模的 \PTOPB{} 的参数域与可达域：把原先只在某一参数区间可行的机制迁移到新的参数区间可稳定/可复验。
\end{remark}

\subsection{形式化：\texorpdfstring{$\PTOPB(\lambda)$}{PTOPB(lambda)} 族与两类增强算子}
\label{sec:ptopb-family-and-two-operators}

\begin{definition}[$\PTOPB(\lambda)$ 族（物理包线的参数化）]
令 $\Lambda$ 为“物理工作点/材料相图参数”空间。把物理总算子丛视为一族对象：
\[
  \PTOPB(\lambda),\qquad \lambda\in\Lambda.
\]
其中 $\lambda$ 可包含（非穷举）：材料成分、缺陷与氧含量、应变、维度（块材/薄膜）、压力、温度、制备路径等。
固定 $\lambda_0\in\Lambda$，得到当前物理包线下的 $\PTOPB\bigl(\lambda_0\bigr)$。
\end{definition}

\begin{definition}[纤维内增强算子（不突破包线）]
给定 $\lambda_0\in\Lambda$。称算子 $c$ 为纤维内增强算子，若
\[
  c:\PTOPB(\lambda_0)\to \PTOPB(\lambda_0),
\]
并且其语义为：在保持物理包线（$\lambda_0$）不变的前提下，通过功能链与保障链的构造/编排/调参，使目标泛函更优，且尽可能把实时干扰压入上边界（可消项）。
\end{definition}

\begin{definition}[基底迁移算子（突破包线）]
称映射 $g:\Lambda\to\Lambda$ 为基底迁移算子，若它把参数点从 $\lambda_0$ 迁移到 $\lambda_1=g(\lambda_0)$，并诱导物理包线发生改变：
\[
  \PTOPB(\lambda_0)\ \rightsquigarrow\ \PTOPB(\lambda_1)=:\PTOPB^{\prime}.
\]
直观含义是：通过材料/工艺/结构/环境的可执行构造，把系统搬迁到新的相区、新的稳定域或新的有效自由度结构里。
\end{definition}

\begin{remark}[同伦/模空间语言的对应（与前文同伦口径对齐）]
在前文“构造同伦类 $[c]$ 与代表元最优化”的口径下，两类增强可重述为：
\begin{enumerate}
  \item 纤维内增强：固定 $\lambda_0$ 不变，在同一构造同伦类 $[c]$ 内选择更优代表元 $c^\star$；
  \item 基底迁移：在参数空间 $\Lambda$（材料/相图/工艺模空间）上选取路径或映射 $g$，把 $\lambda_0$ 迁移到 $\lambda_1$，从而更换对象 $\PTOPB(\lambda)$ 本身（包线搬家）。
\end{enumerate}
因此 A 类是“同伦类内选代表元”，B 类是“在模空间上走路并改变对象”，二者层级不同但可在同一增强序列中交错出现。
\end{remark}

\begin{theorem}[公理降级：任意增强序列可分解为“包线迁移 + 纤维内优化”（按步数归纳）]
\label{thm:two-level-augmentation-decomposition}
考虑一个有限增强序列 $u_{0:T-1}$，其中每一步 $u_t$ 要么为纤维内增强 $c_t$，要么为基底迁移 $g_t$。
则存在一个由迁移步组成的复合 $G_T$ 与一个由纤维步组成的复合 $C_T$ 使得：
\begin{enumerate}
  \item 参数演化只由迁移步决定：$\lambda_T = G_T(\lambda_0)$；
  \item 在最终包线 $\PTOPB(\lambda_T)$ 上，系统的“纤维内实现”可被视为在该包线内的可执行构造复合（即纤维内优化代表元选择）。
\end{enumerate}
特别地，若序列中不含任何迁移步，则 $\lambda_T=\lambda_0$，增强不突破 $\PTOPB\bigl(\lambda_0\bigr)$ 的边界。
\end{theorem}

\begin{Certificate}[数学归纳法证书：按增强步数 $T$ 归纳]
\label{cert:two-level-augmentation-decomposition}
定义谓词 $P(T)$ 为：“任意长度为 $T$ 的增强序列，其最终参数 $\lambda_T$ 仅由其中迁移步的复合 $G_T$ 决定；且纤维步可被视为最终包线内的构造复合。”
证书化要求为：
\[
  P(1)\ \text{成立},\qquad
  (\forall T\in\NN_{\ge 1})\ \bigl(P(T)\Rightarrow P(T+1)\bigr).
\]
\end{Certificate}

\begin{proof}[证明（数学归纳法证书；谓词因果链）]
由证书~\ref{cert:two-level-augmentation-decomposition}：
\begin{itemize}
  \item 基步：$T=1$ 时，若 $u_0=c_0$ 则 $\lambda_1=\lambda_0$；若 $u_0=g_0$ 则 $\lambda_1=g_0(\lambda_0)$。两种情形均满足“参数仅由迁移步决定”，
    且纤维步显然是最终包线内的构造复合。故 $P(1)$ 成立。
  \item 归纳步：假设 $P(T)$ 成立。考虑长度为 $T+1$ 的序列 $u_{0:T}$。其前 $T$ 步由归纳假设给出
  $\lambda_T = G_T(\lambda_0)$。第 $T$ 步：
  \begin{itemize}
    \item 若 $u_T=c_T$，则参数不变，$\lambda_{T+1}=\lambda_T$，仍由 $G_T$ 决定；
    \item 若 $u_T=g_T$，则参数更新为 $\lambda_{T+1}=g_T(\lambda_T)=(g_T\circ G_T)(\lambda_0)$，仍仅由迁移步复合决定。
  \end{itemize}
  纤维内实现的复合性由“纤维步在对应包线内闭合”给出（可执行构造复合仍落在最终包线内）。因此 $P(T+1)$ 成立。
\end{itemize}
由数学归纳法得对任意有限 $T$ 结论成立。证毕。
\end{proof}

\subsection{上同调解释：突破边界 = 改变复形，而非只换代表元}
\label{sec:cohomology-interpretation-boundary-shift}

\begin{lemma}[纤维内增强不改变复形，只改变代表元]
\label{lem:fiber-augmentation-does-not-change-complex}
固定 $\lambda_0$。若增强序列仅由纤维内增强步组成，则其仅在同一对象 $\PTOPB(\lambda_0)$ 上选取不同实现代表元；
因此可视为在固定复形 $C^\bullet\bigl(\PTOPB(\lambda_0)\bigr)$ 内做代表元调整，而不改变复形本身。
\end{lemma}

\begin{Certificate}[证书：对象不变 $\Rightarrow$ 复形不变]
\label{cert:fiber-augmentation-does-not-change-complex}
证书登记：当 $\lambda$ 固定时，对象 $\PTOPB(\lambda)$ 固定；
而共链复形 $C^\bullet(\PTOPB(\lambda))$ 是对对象的函子构造，因此对象不变 $\Rightarrow$ 复形不变。
\end{Certificate}

\begin{proof}[证明（谓词因果链）]
由证书~\ref{cert:fiber-augmentation-does-not-change-complex}，令谓词链为：
\[
\begin{aligned}
  Y_1 &: \text{仅含纤维内增强}\Rightarrow \lambda\ \text{不变},\\
  Y_2 &: \Rightarrow\ \PTOPB(\lambda)\ \text{不变},\\
  Y_3 &: \Rightarrow\ C^\bullet(\PTOPB(\lambda))\ \text{不变}.
\end{aligned}
\]
由 $Y_1\Rightarrow Y_2\Rightarrow Y_3$ 得引理结论。证毕。
\end{proof}

\begin{proposition}[突破边界的增强允许“障碍类”翻转为可消项（条件命题）]
\label{prop:boundary-shift-can-trivialize-obstruction}
给定 $\lambda_0\in\Lambda$ 与某一闭合干扰类 $[d]\in H^1\bigl(\PTOPB(\lambda_0)\bigr)$。
若存在基底迁移 $g$ 使 $\lambda_1=g(\lambda_0)$ 并满足在新包线中
\[
  [d] = 0\ \text{于}\ H^1\bigl(\PTOPB(\lambda_1)\bigr),
\]
则原先在 $\lambda_0$ 下不可消的结构性压力，在 $\lambda_1$ 的新包线下可被视为可消项（上边界）并可由构造算子抵消。
\end{proposition}

\begin{Certificate}[证书：障碍类平凡化 $\Rightarrow$ 存在上边界补偿]
\label{cert:boundary-shift-can-trivialize-obstruction}
证书登记：$[d]=0$ 在 $H^1$ 中等价于 $d\in\operatorname{im}(\partial^0)$；
因此一旦迁移使得 $[d]$ 平凡化，则存在 $c\in C^0$ 使 $d=\partial^0(c)$，从而 $d$ 可消。
\end{Certificate}

\begin{proof}[证明（谓词因果链）]
由证书~\ref{cert:boundary-shift-can-trivialize-obstruction}，令谓词链为：
\[
\begin{aligned}
  Z_1 &: [d]=0\ \text{于}\ H^1\bigl(\PTOPB(\lambda_1)\bigr),\\
  Z_2 &: \Rightarrow\ d\in\operatorname{im}(\partial^0),\\
  Z_3 &: \Rightarrow\ (\exists c)\ d=\partial^0(c)\Rightarrow d\ \text{可消}.
\end{aligned}
\]
由 $Z_1\Rightarrow Z_2\Rightarrow Z_3$ 得命题结论。证毕。
\end{proof}

\begin{corollary}[目标牵引下的必要性：若目标区间在当前包线外，则仅纤维内增强不足]
\label{cor:need-boundary-shift-for-out-of-envelope-goal}
令 $\mathcal{T}\subseteq\Lambda$ 表示目标区间（例如同时要求 $T\approx 300\text{K}$ 且 $P\approx 1\text{atm}$ 且可工程化的工作点集合）。
若 $\lambda_0\notin\mathcal{T}$ 且在仅含纤维内增强的口径下无法把系统迁移到 $\mathcal{T}$，
则必须引入基底迁移算子 $g$（突破边界）以获得 $\lambda_1=g(\lambda_0)\in\mathcal{T}$ 的新包线。
\end{corollary}

\begin{Certificate}[证书：目标区间外 $\Rightarrow$ 需要参数迁移]
\label{cert:need-boundary-shift-for-out-of-envelope-goal}
证书登记：仅含纤维内增强 $\Rightarrow$ 参数 $\lambda$ 不变（由定理~\ref{thm:two-level-augmentation-decomposition}）。
因此若目标区间以参数域刻画且 $\lambda_0\notin\mathcal{T}$，则要进入 $\mathcal{T}$ 必须引入改变 $\lambda$ 的迁移步。
\end{Certificate}

\begin{proof}[证明（谓词因果链）]
由证书~\ref{cert:need-boundary-shift-for-out-of-envelope-goal}，令谓词链为：
\[
\begin{aligned}
  S_1 &: \text{仅纤维内增强}\Rightarrow \lambda_T=\lambda_0,\\
  S_2 &: \lambda_0\notin\mathcal{T}\Rightarrow \lambda_T\notin\mathcal{T},\\
  S_3 &: \Rightarrow\ \text{要进入}\ \mathcal{T}\ \text{必须引入基底迁移}\ g.
\end{aligned}
\]
由 $S_1\Rightarrow S_2\Rightarrow S_3$ 得推论结论。证毕。
\end{proof}

\subsection{备注：以“室温超导”类目标为例的工程直觉}
\label{sec:room-temperature-superconductivity-remark}

\begin{remark}[为何“突破边界增强”在该类目标中更突出]
在“常压常温超导”这类目标牵引工程中，目标本身要求把既有“高压可行”的物理区间迁移到“常压常温可用”的区间。
因此若把目标写为“$T\approx 300\text{K}$ 且 $P\approx 1\text{atm}$ 且可工程化”，则常见情形是：单纯的纤维内增强（不改变包线）不足以覆盖目标区间，必须引入基底迁移（改变 $\lambda$）。
例如一类综述指出：富氢体系（如 $\mathrm{H_3S}$、$\mathrm{LaH_{10}}$ 等）在约 $140\text{--}200\,\mathrm{GPa}$ 的高压区间存在可重复的高 $T_c$ 证据；
而若把目标进一步推到 $T\approx 300\,\mathrm{K}$，讨论中常伴随更高压力尺度（如 $350\text{--}400\,\mathrm{GPa}$ 等量级）的需求。
在本章口径下，这些事实并不意味着“违背定律”，而是提示：当目标区间与现有包线不相交时，必须通过改变 $\lambda$ 的迁移步去重塑可达域\cite{ref_nsr_nwae047}。
\end{remark}

\begin{remark}[两类“基底迁移算子”的雏形（仅作方法学对应）]
从工程方法学角度看，以下路线可被视为“突破边界”的迁移算子雏形（用于说明口径，而非对具体结论作背书）：
\begin{enumerate}
  \item \textbf{用外延应变替代体压力（把压力维度部分等效为应变/界面条件维度）}：
	  以某些原先需要高压才能稳定结构相区/相态的体系为例，通过薄膜外延与基底引入的面内压缩应变来稳定相区，
	  从而在\textbf{常压}下获得原先需要高压才能进入的结构窗口。
	  同时也提示：需要把“起始转变温度”“零电阻窗口”“抗磁响应窗口”等不同层级证据严格区分并证书化登记，以避免把不同物理含义的测量现象混为一谈\cite{ref_slac_2025_room_pressure,
	    ref_nature_2025_08755_z}；
	  \item \textbf{压力淬火把高压相带回常压（pressure-quench protocol, PQP；亚稳保留）}：
	  先在高压诱导目标相态/机制，再按特定工艺路径淬火并回收到常压，使其以亚稳态形式保留并开展后续表征。
	  该路线的关键不只在于“能否回收”，更在于“能否可重复、可持续、可审计”，
	    因此天然要求把“加压$\to$诱导$\to$淬火$\to$回收$\to$稳定性测试$\to$测量链”作为顺序敏感的非交换构造词来登记\cite{ref_uh_2025_pqp}。
\end{enumerate}
这些路线的共性是：迁移步本身通常是非交换的构造序列（顺序敏感），并且其可复验性强依赖“保障链/证据链”的证书化登记。
\end{remark}

\begin{remark}[为何必须把“干扰”扩展为实时干扰算子（内外负算子）]
在“突破边界增强”的语境下，把干扰仅限定为外部环境往往会失真，因为内部负效应同样会形成决定性障碍，例如：
\begin{itemize}
  \item 样品非均匀与相分离导致的“看起来像转变”的伪影；
  \item 数据处理流程（背景扣除、滤波、选择性展示等）引入的不可审计性；
  \item 表征链路不完备（只看电阻不看磁响应，或反之）导致的证据不闭合；
  \item 工艺窗口极窄导致的低成功率/不可复现。
\end{itemize}
这些都属于 $\DeltaRT$ 应捕捉的负算子族。
相关争议与撤稿事件也提示：证据链与分析流程本身必须进入“保障算子”与“干扰算子”的统一审计口径\cite{ref_nature_2022_05294_9,ref_sciam_retraction_note}。
\end{remark}

\begin{remark}[媒体报道参考（非学术，便于追溯争议脉络）]
见\cite{ref_wsj_61288727,ref_wsj_aa2f9fd4,ref_time_ambient_superconductor}。
\end{remark}

\section{同伦扭曲两层搜索：相内/相间的工程化架构}
\label{sec:twisted-homotopy-two-level-search}

\begin{remark}[定位（把“两层增强”落成可执行搜索）]
本章把前文“两层增强”（定理~\ref{thm:two-level-augmentation-decomposition}）进一步落成一套可执行的搜索架构：
把搜索显式分成\textbf{相内（不跨相变边界）}与\textbf{相间（通过工程算子改变相空间/相图参数域）}两层。
并且在每一层里分别进行“算法选型$\to$代表元更新$\to$退出条件”的证书化流程，从而把“增强/迁移”从叙述层提升为可回归审计的工程流程。
\end{remark}

\subsection{对象与目标：相空间参数、构造代表元与鲁棒泛函}
\label{sec:twisted-search-objects-and-objective}

\begin{definition}[搜索态（相空间参数 + 构造代表元）]
令 $\Lambda$ 为相空间/相图参数域（材料成分、缺陷与氧含量、应变、维度、压力、温度、制备路径、界面条件等）。
对任意 $\lambda\in\Lambda$，记 $\mathcal{C}(\lambda)\subseteq\mathcal{C}$ 为在包线 $\PTOPB(\lambda)$ 下\textbf{可执行}且\textbf{可审计}
  的构造集合（功能链 + 保障链 + 编排顺序/阈值/资源分配/验证链等）。
称
\[
  \sigma := (\lambda,c),\qquad \lambda\in\Lambda,\ c\in\mathcal{C}(\lambda)
\]
为搜索态；其中 $c$ 可理解为单个构造算子或构造词（非交换复合）的等价类代表元。
\end{definition}

\begin{definition}[鲁棒目标（实时干扰同伦类下）]
令 $J(c,\lambda;d)$ 为在构造 $c$、相参数 $\lambda$、实时干扰 $d$ 下的目标泛函（可包含性能、成本、可复现性等）。
给定干扰同伦类 $[d]\in\DeltaRT/\!\sim$（见第~\ref{sec:disturbance-homotopy-class} 节），定义鲁棒/风险口径的目标为（两种口径择一即可）：
\[
  \tilde J_{\max}(c,\lambda)
  := \max_{d\in[d]} J(c,\lambda;d),
\]
或
\[
  \tilde J_{\mathrm{risk}}(c,\lambda)
  := \mathbb{E}\bigl[J(c,\lambda;d)\bigr]+\beta\cdot \mathrm{Risk}\bigl(J(c,\lambda;d)\bigr),
\]
其中 $\beta\ge 0$ 为风险权重，期望与风险算子由证书化数据口径给出。
本章记 $\tilde J$ 为选定口径下的鲁棒目标。
\end{definition}

\begin{remark}[两类同伦类（与前文定义对齐）]
\textbf{干扰同伦类} $[d]\in\DeltaRT/\!\sim$ 用于把“内外负算子”按不穿越关键约束边界的连续变形归并（第~\ref{sec:disturbance-homotopy-class} 节）。
\textbf{构造同伦类} $[c]\in\mathcal{C}/\!\simeq$ 用于把“功能+保障”的结构等价类归并，并在类内选取最优代表元（第~\ref{sec:construction-homotopy-and-representative-optimization} 节）。
因此鲁棒优化可以表述为：在给定（或可迁移的）$\lambda$ 下，从某个构造同伦类 $[c]$ 中选取代表元 $c^\star$ 使 $\tilde J(c^\star,\lambda)$ 最优。
\end{remark}

\subsection{两层搜索算子：相内局部搜索与相间迁移}
\label{sec:twisted-search-two-level-operators}

\begin{definition}[相内区域与相内局部搜索算子（不跨相变）]
称 $R\subseteq\Lambda$ 为相内区域，若在工程容差与观测分辨率下，$R$ 内不穿越相变边界（序参量跃迁/强滞回/统计结构突变等报警信号不触发）。
给定 $\lambda\in R$，称映射
\[
  \mathsf{LS}_{\lambda,R}:\mathcal{C}(\lambda)\to \mathcal{C}(\lambda)
\]
为相内局部搜索算子，若它仅在固定包线 $\PTOPB(\lambda)$ 内更新构造代表元（例如调参、重构、编排微调），并以降低 $\tilde J(\cdot,\lambda)$ 为目标。
\end{definition}

\begin{definition}[相间迁移算子（改变相图参数域）]
称映射 $g:\Lambda\to\Lambda$ 为相间迁移算子（工程算子），若它把参数点从 $\lambda$ 迁移到 $\lambda' = g(\lambda)$，并诱导物理包线改变：
\[
  \PTOPB(\lambda)\ \rightsquigarrow\ \PTOPB(\lambda').
\]
工程上 $g$ 可由材料/成分/缺陷工程、应变与界面工程、压力--温度路径与淬火路径、制备路线改变等可执行动作实现；
并且当多个动作按序复合时，一般满足非交换性（顺序敏感）。
\end{definition}

\begin{definition}[两层搜索步（阶段更新口径）]
称搜索态更新 $\sigma\mapsto\sigma'$ 为：
\begin{enumerate}
  \item 相内步（A 类）：$\sigma'=(\lambda,\,\mathsf{LS}_{\lambda,R}(c))$，其中 $\lambda$ 固定；
  \item 相间步（B 类）：$\sigma'=(g(\lambda),\,c')$，其中 $g$ 为迁移算子，$c'$ 由跨相搬运/重整定给出（见第~\ref{sec:twisted-search-transport} 节）。
\end{enumerate}
因此“两层搜索”是对 A/B 两类步的交错编排，而不是单一连续优化过程。
\end{definition}

\begin{theorem}[公理降级：两层搜索轨迹的结构分解（按外层步数归纳）]
\label{thm:twisted-two-level-search-decomposition}
考虑长度为 $T$ 的两层搜索序列 $\{u_t\}_{t=0}^{T-1}$，其中每一步 $u_t$ 要么为相内步，要么为相间步。
记搜索态序列 $\sigma_t=(\lambda_t,c_t)$ 由 $\sigma_{t+1}=u_t(\sigma_t)$ 生成。
则存在由相间步复合得到的迁移算子 $G_T$，使得参数演化仅由相间步决定：
\[
  \lambda_T = G_T(\lambda_0).
\]
并且在任意固定的相内区域中，所有相内步的构造更新可视为同一纤维 $\mathcal{C}(\lambda)$ 内的代表元迭代（即在固定包线内做同伦类内搜索）。
\end{theorem}

\begin{Certificate}[数学归纳法证书：按外层步数 $T$ 归纳的轨迹分解]
\label{cert:twisted-two-level-search-decomposition}
定义谓词 $P(T)$ 为：
“任意长度为 $T$ 的两层搜索序列，其最终参数 $\lambda_T$ 仅由其中相间步的复合决定；且相内步不改变 $\lambda$ 并在对应纤维内闭合。”
证书化要求为：
\[
  P(1)\ \text{成立},\qquad
  (\forall T\in\NN_{\ge 1})\ \bigl(P(T)\Rightarrow P(T+1)\bigr).
\]
\end{Certificate}

\begin{proof}[证明（数学归纳法证书；谓词因果链）]
由证书~\ref{cert:twisted-two-level-search-decomposition}：
\begin{itemize}
  \item 基步：$T=1$ 时，若 $u_0$ 为相内步，则 $\lambda_1=\lambda_0$；若 $u_0$ 为相间步，则 $\lambda_1=g_0(\lambda_0)$。两种情形均满足“参数仅由相间步决定”。
    同时相内步显然在固定 $\mathcal{C}(\lambda_0)$ 内闭合。故 $P(1)$ 成立。
  \item 归纳步：假设 $P(T)$ 成立。考虑长度为 $T+1$ 的序列 $u_{0:T}$。其前 $T$ 步由归纳假设给出
  $\lambda_T = G_T(\lambda_0)$。第 $T$ 步：
  \begin{itemize}
    \item 若 $u_T$ 为相内步，则参数不变，$\lambda_{T+1}=\lambda_T$，仍由 $G_T$ 决定；
    \item 若 $u_T$ 为相间步，则参数更新为 $\lambda_{T+1}=g_T(\lambda_T)=(g_T\circ G_T)(\lambda_0)$，仍仅由相间步复合决定。
  \end{itemize}
  相内闭合性由相内步定义给出（$\mathsf{LS}_{\lambda,R}:\mathcal{C}(\lambda)\to\mathcal{C}(\lambda)$）。因此 $P(T+1)$ 成立。
\end{itemize}
由数学归纳法得对任意有限 $T$ 结论成立。证毕。
\end{proof}

\begin{corollary}[推论：不含迁移步则退化为相内代表元搜索]
\label{cor:twisted-search-reduces-to-intra-phase}
若两层搜索序列中不含任何相间步，则对任意 $T$ 有 $\lambda_T=\lambda_0$，搜索仅在固定纤维 $\mathcal{C}(\lambda_0)$ 内进行代表元更新。
\end{corollary}

\begin{Certificate}[证书：无相间步 $\Rightarrow$ $\lambda$ 常值]
\label{cert:twisted-search-reduces-to-intra-phase}
证书登记：相内步不改变 $\lambda$；若全序列仅由相内步组成，则对任意步数 $T$ 有 $\lambda_T=\lambda_0$。
\end{Certificate}

\begin{proof}[证明（谓词因果链）]
由证书~\ref{cert:twisted-search-reduces-to-intra-phase}，令谓词链为：
\[
\begin{aligned}
  W_1 &: \text{相内步}\Rightarrow \lambda\ \text{不变},\\
  W_2 &: \text{全序列仅含相内步}\Rightarrow \lambda_T=\lambda_0,\\
  W_3 &: \Rightarrow\ \text{搜索仅在}\ \mathcal{C}(\lambda_0)\ \text{内更新}.
\end{aligned}
\]
由 $W_1\Rightarrow W_2\Rightarrow W_3$ 得推论结论。证毕。
\end{proof}

\subsection{同伦扭曲：跨相搬运算子与热启动一致性}
\label{sec:twisted-search-transport}

\begin{definition}[扭曲搬运算子（平行搬运/扭曲映射）]
给定相间迁移 $g:\lambda\mapsto\lambda'$。称映射
\[
  T_g:\mathcal{C}(\lambda)\to\mathcal{C}(\lambda')
\]
为扭曲搬运算子，若它把旧包线下的构造代表元 $c$ 迁移到新包线下的“等效构造”$T_g(c)$，
用于跨相热启动与局部搜索初始化。
工程上 $T_g$ 可由规则驱动（编排重排、参数缩放、重整定规则）、模型驱动（等效参数变换）、数据驱动（迁移历史学习）等方式实现。
\end{definition}

\begin{definition}[扭曲一致性（同伦类保持）]
称 $T_g$ 满足扭曲一致性，若对任意 $c_0,c_1\in\mathcal{C}(\lambda)$ 有
\[
  c_0\simeq c_1\ \Rightarrow\ T_g(c_0)\simeq T_g(c_1),
\]
并且 $T_g$ 保持可执行性：$c\in\mathcal{C}(\lambda)\Rightarrow T_g(c)\in\mathcal{C}(\lambda')$。
直观含义是：跨相搬运不引入“结构性偷换”，使同伦类代表元的传递在审计口径下不失真。
\end{definition}

\begin{lemma}[扭曲一致性 $\Rightarrow$ 搬运在同伦类上良定义]
\label{lem:transport-induces-map-on-homotopy-classes}
若 $T_g$ 满足扭曲一致性，则 $T_g$ 在构造同伦类上诱导良定义映射
\[
  \bar T_g:\mathcal{C}(\lambda)/\!\simeq\ \to\ \mathcal{C}(\lambda')/\!\simeq,
  \qquad
  [c]\mapsto [T_g(c)].
\]
因此跨相热启动可以在“同伦类水平”进行：把旧类 $[c]$ 的代表元搬运到新类 $\bar T_g([c])$ 再做相内代表元优化。
\end{lemma}

\begin{Certificate}[证书：等价关系因子化]
\label{cert:transport-induces-map-on-homotopy-classes}
证书登记：若 $c_0\simeq c_1\Rightarrow T_g(c_0)\simeq T_g(c_1)$，则映射 $c\mapsto T_g(c)$ 可因子化为商映射上的映射 $[c]\mapsto [T_g(c)]$。
\end{Certificate}

\begin{proof}[证明（谓词因果链）]
由证书~\ref{cert:transport-induces-map-on-homotopy-classes}，令谓词链为：
\[
\begin{aligned}
  V_1 &: c_0\simeq c_1\Rightarrow T_g(c_0)\simeq T_g(c_1),\\
  V_2 &: \Rightarrow\ [c]\ \text{的取值与代表元选择无关},\\
  V_3 &: \Rightarrow\ \bar T_g([c]):=[T_g(c)]\ \text{良定义}.
\end{aligned}
\]
由 $V_1\Rightarrow V_2\Rightarrow V_3$ 得引理结论。证毕。
\end{proof}

\begin{remark}[非交换性提示：跨相搬运通常是序列算子]
当构造以“算子词”形式出现且顺序敏感时，$T_g$ 往往不是单个闭式变换，而是一个\textbf{非交换的重整定/重排序列}：
例如“先重标定$\to$再改阈值$\to$再换验证链”的顺序与“先换验证链$\to$再重标定”的风险边界一般不同。
因此把 $T_g$ 纳入证书边界并记录其顺序，是“同伦扭曲”可复验性的关键。
\end{remark}

\subsection{流程与退出条件：相内 $\to$ 相间 $\to$ 再相内}
\label{sec:twisted-search-workflow}

\begin{definition}[相内区域分层与相变报警（工程判据口径）]
设 $\{R_i\}_{i\in I}$ 为对 $\Lambda$ 的一组分层/分区（可迭代修正），并记判别器
\[
  \mathsf{Region}:\Lambda\to I,\qquad \lambda\mapsto i\ \text{使}\ \lambda\in R_i.
\]
设 $\mathsf{PhaseAlarm}$ 为相变报警信号（可由序参量跃迁、滞回、响应折点稳定复现、噪声统计结构突变等触发），其输出进入“是否继续相内搜索/是否需要换相”的退出条件证书。
\end{definition}

\begin{definition}[同伦扭曲两层搜索（骨架流程）]
给定初始搜索态 $\sigma_0=(\lambda_0,c_0)$、工程算子库 $\mathcal{G}$ 与搬运算子族 $\{T_g\}_{g\in\mathcal{G}}$。
两层搜索流程可写为：
\begin{enumerate}
  \item 相内阶段（A）：在 $R_i\ni\lambda$ 内，令 $c\leftarrow \mathsf{LS}_{\lambda,R_i}(c)$ 迭代，直到触发退出条件；
  \item 相间阶段（B）：从候选迁移算子序列中选择 $g^\star$，令 $\lambda'\leftarrow g^\star(\lambda)$，
  并用 $c'\leftarrow T_{g^\star}(c)$ 热启动；
  \item 回到相内阶段，在新包线 $\PTOPB(\lambda')$ 下继续代表元搜索。
\end{enumerate}
\end{definition}

\begin{remark}[退出/切换条件（非穷举）]
相内阶段的常用退出条件可包括：
\begin{enumerate}
  \item 改善幅度低于阈值（平台期）；
  \item 继续相内优化的边际收益低于换相的期望收益（成本--收益证书）；
  \item 相变报警升高（接近相边界：响应不稳定、滞回出现、方差放大、可重复性下降等）。
\end{enumerate}
相间阶段则需要“候选生成$\to$廉价筛选$\to$风险约束下选择”的证书链，以控制算子序列空间的指数增长。
\end{remark}

\begin{remark}[伪代码骨架（工程实现模板）]
\begin{verbatim}
Input: lambda0, c0
       Region(lambda) -> Ri
       PhaseAlarm(...)
       operator library G, transport Tg
       budget Budget

(lambda, c) <- (lambda0, c0)
Best <- (lambda, c, Jtilde(c,lambda))

while Budget not exhausted:
  i <- Region(lambda)

  # A: intra-phase search
  c <- LocalSearch(c; lambda, Ri, rule=FeatureBased)
  if Jtilde(c,lambda) better than Best: update Best

	  if SwitchCondition(plateau / marginal_gain_low / PhaseAlarm_high):
	      # B: inter-phase search
	      Cand <- GenerateCandidates(G, lambda, Best,
	                                 constraints/risk)
	      Cand <- CheapScreen(Cand)               # surrogate/multifidelity/prior
	      g*   <- Select(Cand,
	                     criteria=E[gain]-cost-risk)

	      lambda' <- g*(lambda)
	      c'      <- Tg*(c)
	      # warm start via twisted transport
	      (lambda, c) <- (lambda', c')
	      continue
	  else:
	      lambda <- SmallPerturbWithin(Ri, lambda)

Output: Best
\end{verbatim}
\end{remark}

\subsection{算法选型：相内局部优化与相间策略搜索}
\label{sec:twisted-search-algorithm-selection}

\begin{remark}[相内局部搜索的选型规则（按可观测特征）]
相内区域几何相对稳定，优先追求效率与可解释收敛。可按观测特征选型：
\begin{enumerate}
  \item 响应近似光滑、噪声低、可得梯度或近似梯度：选基于梯度的局部优化（步长控制 + 约束处理）；
  \item 光滑性不确定或含离散编排决策（顺序敏感强）：选导数无关局部搜索，并把“算子词/顺序”作为显式离散变量处理；
  \item 评估代价高、维度不高、噪声中等：选代理模型的局部优化（多保真/显式噪声建模）；
  \item 接近相边界（滞回、突变、方差放大）：选信赖域/小步探索 + 多起点，并降低对光滑性的假设。
\end{enumerate}
\end{remark}

\begin{remark}[相间探索的选型规则（把换相当作策略/图搜索）]
相间阶段的核心是选择 $g$ 或非交换算子序列 $g_k\circ\cdots\circ g_1$，更接近策略/图搜索而非连续优化：
\begin{enumerate}
  \item 算子库规模较小且代价/收益可估：用启发式图搜索（以相内区域 $R_i$ 为节点、$g$ 为边）；
  \item 序列很长且强非交换：用分层策略搜索 + 失败快速剪枝（中途出现不可逆负面指标即终止）；
  \item 跨相成功率低、路径依赖强（如淬火/界面诱导等）：用路径搜索 + 扭曲搬运热启动（$T_g$）以降低重启成本。
\end{enumerate}
\end{remark}

\begin{remark}[实施关键建议（把“可复验”纳入优化变量）]
\begin{enumerate}
  \item 把每个相内区域 $R_i$ 做成数据对象：记录典型响应形态、噪声水平、滞回风险、适用局部算法、以及高概率有效的迁移边（哪些 $g$ 更可能可行）。
  \item 把“验证链/保障链”纳入 $c$ 一并优化：在目标牵引工程中，验证链决定假阳性/不可复现是否会被当成改进信号；将其纳入构造变量可避免搜索被测量伪影牵引到错误同伦类。
\end{enumerate}
\end{remark}

\section{同伦扭曲实时控制：相空间穿梭与动态算法选型}
\label{sec:twisted-homotopy-real-time-control}

\begin{remark}[定位（从离线搜索到“实验--表征--工艺--控制”闭环）]
当“搜索”从离线优化推进到“实验--表征--工艺--控制”的在线闭环时，
实时控制本身就演化为一种\textbf{相空间穿梭}问题：控制的首要对象不再只是某个物理状态 $x_{\mathrm{phys}}(t)$ 的稳定，
而是相图参数轨迹 $\lambda(t)$ 的规划与跟踪，使系统进入并尽量停留在目标相域或目标亚稳盆地。
当系统存在强路径依赖（滞回、多稳态、不可逆微结构演化、顺序敏感的工艺链）时，
该实时控制天然带有\textbf{同伦扭曲}特征：同一终点的不同路径不等价，需要用“搬运/续接/扭曲映射”保持连续可行。
\end{remark}

\subsection{相空间穿梭对象：相图参数轨迹与控制旋钮}
\label{sec:phase-space-shuttling-object}

\begin{definition}[相图参数轨迹（相空间控制旋钮）]
令 $\Lambda$ 为相图参数空间（基底空间）。实时控制的相空间轨迹记为
\[
  \lambda(t)=\bigl(P(t),\,T(t),\,\varepsilon(t),\,x(t),\,\mu(t),\,\dots\bigr)\in\Lambda,
\]
其中 $P$ 表示压力/应力口径，$T$ 表示温度与路径段落，$\varepsilon$ 表示应变（含外延应变/界面失配等），
$x$ 表示成分/掺杂/缺陷浓度/氧含量等，$\mu$ 表示化学势相关变量（氢化程度、气氛分压、离子迁移环境等）。
\end{definition}

\begin{definition}[相空间穿梭控制问题（目标集与可复验窗口）]
令 $\mathcal{T}\subseteq\Lambda$ 为目标相域/目标工作点集合（例如“常压常温且可工程化”的区间），
令 $\mathcal{F}\subseteq\Lambda$ 为可控/安全运行域（由设备边界、样品不可逆损伤阈值、风险边界等定义并证书化）。
相空间穿梭控制的基本目标是：在保持 $\lambda(t)\in\mathcal{F}$ 的约束下，使 $\lambda(t)$ 尽快进入 $\mathcal{T}$ 并最大化其停留时间（可复验窗口）。
\end{definition}

\subsection{同伦扭曲控制：路径依赖与非平凡回归}
\label{sec:twisted-homotopy-control}

\begin{definition}[路径诱导搬运（纤维态的平行搬运口径）]
为刻画路径依赖，引入“纤维态”集合族 $\{\mathcal{Y}(\lambda)\}_{\lambda\in\Lambda}$，
其中 $\mathcal{Y}(\lambda)$ 记录在相参数为 $\lambda$ 时的内部慢变量与工程配置（例如微结构/相分数/缺陷谱/残余应力 + 表征链与保障链配置等）。
对任意可执行路径 $\gamma:[0,1]\to\Lambda$，记其诱导的搬运算子为
\[
  T_{\gamma}:\mathcal{Y}\bigl(\gamma(0)\bigr)\to \mathcal{Y}\bigl(\gamma(1)\bigr).
\]
\end{definition}

\begin{definition}[同伦扭曲（twisted homotopy）控制判据]
称系统的相空间控制具有同伦扭曲特征，若存在两条端点相同的可执行路径 $\gamma_1,\gamma_2$，
满足 $\gamma_1(0)=\gamma_2(0)$ 且 $\gamma_1(1)=\gamma_2(1)$，但存在 $y\in\mathcal{Y}(\gamma_1(0))$ 使
\[
  T_{\gamma_1}(y)\neq T_{\gamma_2}(y).
\]
等价地，存在回路 $\ell$ 使得 $T_{\ell}$ 在纤维上作用非平凡（非平凡回归/holonomy）。
\end{definition}

\begin{proposition}[命题：非交换工艺序列给出同伦扭曲控制的可观测见证]
\label{prop:noncommutative-process-implies-twist}
若存在两条由可执行工艺动作构成的控制词（顺序敏感序列）$\mathbf{u}_A,\mathbf{u}_B$，
它们在相空间上诱导端点相同的路径 $\gamma_A,\gamma_B$，但在某些可观测输出/纤维态上产生不同结果，
则控制具有同伦扭曲特征。
\end{proposition}

\begin{Certificate}[证书：端点相同但纤维态不同的见证]
\label{cert:noncommutative-process-implies-twist}
证书登记：存在 $y\in\mathcal{Y}(\gamma_A(0))$ 使
$T_{\gamma_A}(y)\neq T_{\gamma_B}(y)$，并且 $\gamma_A(0)=\gamma_B(0)$、$\gamma_A(1)=\gamma_B(1)$。
该见证可由滞回回线、亚稳保留时间差异、不可逆指标、或可复现实验窗口差异等审计字段给出。
\end{Certificate}

\begin{proof}[证明（谓词因果链）]
由证书~\ref{cert:noncommutative-process-implies-twist}，令谓词链为：
\[
\begin{aligned}
  A_1 &: \gamma_A(0)=\gamma_B(0)\ \wedge\ \gamma_A(1)=\gamma_B(1),\\
  A_2 &: (\exists y)\ T_{\gamma_A}(y)\neq T_{\gamma_B}(y),\\
  A_3 &: \Rightarrow\ \text{存在端点相同但搬运不同的两条路径}.
\end{aligned}
\]
由 $A_1\wedge A_2\Rightarrow A_3$，满足同伦扭曲判据，命题成立。证毕。
\end{proof}

\begin{corollary}[推论：无记忆/可交换情形下可退化为端点控制]
\label{cor:untwisted-reduces-to-endpoint-control}
若对任意端点相同的可执行路径 $\gamma_1,\gamma_2$ 都有 $T_{\gamma_1}=T_{\gamma_2}$（搬运仅依赖端点），
则同伦扭曲消失；实时控制可在相空间上退化为“点到点”的端点规划问题，而无需把路径同伦类作为一等变量。
\end{corollary}

\begin{Certificate}[证书：搬运仅依赖端点 $\Rightarrow$ 无同伦扭曲]
\label{cert:untwisted-reduces-to-endpoint-control}
证书登记：$(\forall \gamma_1,\gamma_2)$ 若端点相同则 $T_{\gamma_1}=T_{\gamma_2}$；
于是对任意回路 $\ell$ 有 $T_{\ell}=\mathrm{id}$，不存在非平凡回归。
\end{Certificate}

\begin{proof}[证明（谓词因果链）]
由证书~\ref{cert:untwisted-reduces-to-endpoint-control}，令谓词链为：
\[
\begin{aligned}
  B_1 &: \text{端点相同}\Rightarrow T_{\gamma_1}=T_{\gamma_2},\\
  B_2 &: \Rightarrow\ \text{任意回路}\ \ell\ \text{满足}\ T_{\ell}=\mathrm{id},\\
  B_3 &: \Rightarrow\ \text{不存在非平凡回归}\Rightarrow \text{无同伦扭曲}.
\end{aligned}
\]
由 $B_1\Rightarrow B_2\Rightarrow B_3$ 得推论结论。证毕。
\end{proof}

\begin{remark}[两类典型穿梭场景（方法学对应）]
\textbf{高压诱导$\to$路径回收常压（回收/保态穿梭）}：
关键不在于“到达高压点”，而在于卸压--温控--等待--表征路径是否把系统带入可保留的亚稳盆地；
同一终点常压常温，不同路径可得到不同的保态时间与可复验窗口\cite{ref_pnas_2501048122}。

\textbf{应变/界面重参数化（等效替代部分压力自由度）}：
通过外延应变、界面效应等把相空间坐标从 $(P,T,\dots)$ 扩展为 $(\varepsilon,T,x,\dots)$，
并引入强离散决策（衬底选择、层数、退火段落），使局部几何与噪声结构发生改变\cite{ref_pubmed_39701131,ref_aps_prmaterials_8_044801}。
\end{remark}

\subsection{续接型滚动优化：实时可行解追踪（NMPC/RTI 口径）}
\label{sec:continuation-rt-optimization}

\begin{definition}[滚动优化子问题（离散采样口径）]
对离散时刻 $k\in\NN$，记观测/估计到的相参数为 $\lambda_k$，纤维态估计为 $y_k\in\mathcal{Y}(\lambda_k)$。
设存在一族滚动优化子问题 $\mathsf{P}_k$（例如 NMPC），其求解输出为一段可执行控制/工艺段落与配置计划
\[
  \pi_k := (u_{k:k+H-1},\,c_{k:k+H-1}),
\]
并且 $\pi_k$ 满足约束（安全边界、设备边界、证据链完整性等）与鲁棒目标降低（使用 $\tilde J$ 口径）。
\end{definition}

\begin{definition}[续接型热启动更新（Real-Time Iteration 抽象）]
称映射
\[
  \mathsf{RTI}:\ (\lambda_k,y_k,\pi_{k-1})\ \mapsto\ \pi_k
\]
为续接型热启动更新算子，若它以 $\pi_{k-1}$ 为初值（热启动），
在严格时限内对 $\mathsf{P}_k$ 做一次或有限次局部更新以得到新的可执行计划 $\pi_k$。
该口径与同伦/续接型 NMPC 实现一致：其目标是“每个周期给出可行且持续改进的解”，而非每次全局收敛\cite{ref_syscop_zanelli2017_homotopy_nmpc,ref_merl_tr2018_082}。
\end{definition}

\begin{theorem}[公理降级：续接型热启动保证滚动可行（按周期归纳）]
\label{thm:continuation-ensures-feasible-rolling}
假设存在常数 $\eta>0$ 与可行修复算子 $\mathsf{Repair}$，使得对任意时刻 $k$：
当相参数变化满足 $\|\lambda_{k+1}-\lambda_k\|\le \eta$ 且上一时刻计划 $\pi_k$ 可行时，
则热启动后的候选计划经 $\mathsf{Repair}$ 可得到下一时刻的可行计划 $\pi_{k+1}$。
则对任意有限步数 $K\in\NN$，若初始计划 $\pi_0$ 可行，则存在一条长度为 $K$ 的可行计划序列 $\{\pi_k\}_{k=0}^{K}$，
从而滚动闭环在前 $K$ 个周期内保持在可行域内。
\end{theorem}

\begin{Certificate}[数学归纳法证书：按周期步数 $K$ 归纳的可行性保持]
\label{cert:continuation-ensures-feasible-rolling}
定义谓词 $P(K)$ 为：“若 $\pi_0$ 可行且对所有 $k<K$ 有 $\|\lambda_{k+1}-\lambda_k\|\le\eta$，
则存在可行计划序列 $\pi_0,\pi_1,\dots,\pi_K$。”
证书化要求为：
\[
  P(1)\ \text{成立},\qquad
  (\forall K\in\NN_{\ge 1})\ \bigl(P(K)\Rightarrow P(K+1)\bigr).
\]
\end{Certificate}

\begin{proof}[证明（数学归纳法证书；谓词因果链）]
由证书~\ref{cert:continuation-ensures-feasible-rolling}：
\begin{itemize}
  \item 基步：$K=1$。已知 $\pi_0$ 可行且 $\|\lambda_1-\lambda_0\|\le\eta$。
  由定理条件中“可行 $\Rightarrow$ 经修复可得下一步可行”，得 $\pi_1$ 可行，故 $P(1)$ 成立。
  \item 归纳步：假设 $P(K)$ 成立，即存在可行序列 $\pi_0,\dots,\pi_K$。
  由条件 $\|\lambda_{K+1}-\lambda_K\|\le\eta$ 与 $\pi_K$ 可行，热启动候选经 $\mathsf{Repair}$ 得到 $\pi_{K+1}$ 可行。
  因此存在 $\pi_0,\dots,\pi_{K+1}$ 可行，故 $P(K+1)$ 成立。
\end{itemize}
由数学归纳法得对任意有限 $K$ 结论成立。证毕。
\end{proof}

\subsection{动态算法选型：风险边界驱动的求解器簇切换}
\label{sec:dynamic-solver-selection}

\begin{definition}[风险报警与求解器选型映射]
设 $\mathsf{Region}(\lambda_k,y_k)$ 给出当前相内区域索引，$\mathsf{Alarm}(y_k)$ 给出风险等级（例如相边界接近、不可逆指标升高、方差放大等）。
给定求解器族 $\{\mathsf{Solver}_{i,r}\}$，定义选型映射
\[
  \mathsf{SelectSolver}:(i,r)\mapsto \mathsf{Solver}_{i,r},
\]
其中 $i$ 为区域索引，$r$ 为风险等级；并由证书化规则登记“何时切换、切换到哪一类算法簇”。
\end{definition}

\begin{lemma}[平滑切换口径：同伦续接保证切换时的连续可行]
\label{lem:smooth-switching-preserves-feasibility}
若从求解器 $\mathsf{Solver}^-$ 切换到 $\mathsf{Solver}^+$ 时，引入同伦参数 $\alpha\in[0,1]$ 构造一族过渡更新
\[
  \mathsf{Update}_{\alpha} := (1-\alpha)\,\mathsf{Update}^- + \alpha\,\mathsf{Update}^+,
\]
并且对所有 $\alpha$ 该过渡更新在证书边界内保持可行性（约束不被破坏），
则切换可以在有限步内完成且不破坏滚动可行解序列的续接性。
\end{lemma}

\begin{Certificate}[证书：同伦参数插值的可行性保持]
\label{cert:smooth-switching-preserves-feasibility}
证书登记：存在 $\alpha_0=0<\alpha_1<\cdots<\alpha_m=1$，
使得对每个 $j$，$\mathsf{Update}_{\alpha_j}$ 将可行计划映射到可行计划。
\end{Certificate}

\begin{proof}[证明（谓词因果链）]
由证书~\ref{cert:smooth-switching-preserves-feasibility}，令谓词链为：
\[
\begin{aligned}
  C_1 &: \pi\ \text{可行}\Rightarrow \mathsf{Update}_{\alpha_j}(\pi)\ \text{可行}\quad(\forall j),\\
  C_2 &: \Rightarrow\ \pi\ \text{沿}\ \alpha_0\to\alpha_m\ \text{的有限链保持可行},\\
  C_3 &: \Rightarrow\ \text{完成切换且不破坏续接可行}.
\end{aligned}
\]
由 $C_1\Rightarrow C_2\Rightarrow C_3$ 得引理结论。证毕。
\end{proof}

\begin{remark}[三层闭环（快/中/慢环）作为可实现架构]
工程实现中可用三层闭环把责任拆开：
\begin{enumerate}
  \item \textbf{快环}：以安全边界为主，保证设备与样品运行在可控域内（不可逆损伤/风险上界）；
  \item \textbf{中环}：续接型滚动优化（第~\ref{sec:continuation-rt-optimization} 节），每周期输出可行且持续改进的动作；
  \item \textbf{慢环}：当相内平台期或结构性障碍出现时，触发相间迁移算子（改变 $\lambda$ 的工程动作），并用扭曲搬运保持热启动一致性（第~\ref{sec:twisted-search-transport} 节）。
\end{enumerate}
在该拆分下，“局部最优算法选型”不是一次性选择，而是随 $(i,r)$ 动态切换，并用同伦续接保证切换时连续可行。
\end{remark}

\section{相空间切换与规范结构：三层级澄清与贴图搬运}
\label{sec:gauge-structure-switching}

\begin{remark}[定位（避免把“换写法”与“换物理相”混淆）]
“切换规范场”这句话可以成立，但需要把它分成不同层级来讲清楚：
否则容易把\textbf{纯粹的规范变换}（只换描述方式）与\textbf{物理相的改变}（有效理论真的改变）混在一起。
本章用“纤维丛/连接”的语言把三层级精确化，并与前文“相空间穿梭 + 同伦扭曲”（第~\ref{sec:twisted-homotopy-real-time-control} 节）对齐。
\end{remark}

\subsection{三层级：规范选择、背景连接与规范结构}
\label{sec:gauge-switching-three-levels}

\begin{definition}[固定相域内的规范结构（主丛 + 连接口径）]
在某一固定相域/固定有效自由度集合中，记规范结构为
\[
  \mathcal{M}:=(\Bundle\xrightarrow{\pr}\Base,\ \Group,\ \Conn),
\]
其中 $\Bundle$ 为主 $\Group$-丛，$\Conn$ 为其上的连接（规范势），$\Curv$ 为曲率（场强）。
“固定相域”表示：有效自由度与约束组织方式不变（仅允许在同一物理包线/同一模型口径内连续变形）。
\end{definition}

\begin{definition}[A 级：规范选择变换（同一相内的 gauge fixing）]
在固定规范结构 $\mathcal{M}$ 下，称“切换规范场”为 A 级，若它仅是规范选择/坐标截面的更换：
存在规范变换 $u:\Base\to \Group$ 使连接变为
\[
  \Conn^{u}:=u^{-1}\Conn\,u + u^{-1}\mathrm{d}u,
\]
并且所有规范不变量（以 $\Curv$ 或 Wilson 回路等为代表）保持不变。
其工程意义对应于：相内搜索时的“坐标图切换/重参数化”，改善数值条件但不改变物理相。
\end{definition}

\begin{proposition}[命题：A 级变换不改变规范不变量（只换写法）]
\label{prop:gauge-fixing-does-not-change-invariants}
在固定规范结构 $\mathcal{M}$ 下，若 $\Conn^{u}$ 由规范变换 $u$ 得到，则曲率满足
\[
  \Curv^{u}=u^{-1}\Curv\,u.
\]
因此任意规范不变量 $I$（满足 $I(u^{-1}\Curv u)=I(\Curv)$）保持不变：
\[
  I(\Curv^{u})=I(\Curv).
\]
\end{proposition}

\begin{Certificate}[证书：曲率协变变换 $\Rightarrow$ 不变量不变]
\label{cert:gauge-fixing-does-not-change-invariants}
证书登记如下谓词链：
\[
\begin{aligned}
  G_1 &: \Conn^{u}=u^{-1}\Conn u + u^{-1}\mathrm{d}u,\\
  G_2 &: \Rightarrow\ \Curv^{u}=u^{-1}\Curv u,\\
  G_3 &: I\ \text{为规范不变量}\Rightarrow I(u^{-1}\Curv u)=I(\Curv),\\
  G_4 &: \Rightarrow\ I(\Curv^{u})=I(\Curv).
\end{aligned}
\]
\end{Certificate}

\begin{proof}[证明（谓词因果链）]
由证书~\ref{cert:gauge-fixing-does-not-change-invariants}：
$G_1\Rightarrow G_2$ 给出曲率的协变变换；
再由 $G_3$ 得 $I(\Curv^{u})=I(u^{-1}\Curv u)=I(\Curv)$；
故命题成立。证毕。
\end{proof}

\begin{definition}[B 级：背景连接随参数变化（同一规范结构内的连接迁移）]
在固定规范结构 $\mathcal{M}$ 下，称“切换规范场”为 B 级，若规范结构不变，但“背景连接/外参场”随相参数变化：
存在一族连接 $\Conn(\lambda)$（例如外加电磁势、应变诱导等效规范势、Berry 连接等）随 $\lambda\in\Lambda$ 改变。
沿参数路径 $\gamma:[0,1]\to\Lambda$ 的平行搬运记为 $\mathrm{PT}_{\gamma}$，并允许出现路径依赖（holonomy）。
其工程意义对应于：仍在“非相变区域”内，但路径依赖开始显著（第~\ref{sec:twisted-homotopy-control} 节）。
\end{definition}

\begin{lemma}[引理：非平凡回归 $\Rightarrow$ 路径同伦类成为决策变量]
\label{lem:holonomy-implies-path-class-relevant}
若存在回路 $\ell$ 使得 $\mathrm{PT}_{\ell}\neq \mathrm{id}$（非平凡回归/holonomy），
则存在两条端点相同的路径 $\gamma_1,\gamma_2$ 满足
\[
  \mathrm{PT}_{\gamma_1}\neq \mathrm{PT}_{\gamma_2}.
\]
因此“只看终点 $\lambda$”不足以刻画系统演化，必须把路径（至少其同伦类）纳入控制/搜索变量。
\end{lemma}

\begin{Certificate}[证书：回路非平凡 $\Rightarrow$ 存在端点相同但搬运不同的两条路径]
\label{cert:holonomy-implies-path-class-relevant}
证书登记谓词链：
\[
\begin{aligned}
  H_1 &: \exists\,\ell\ \text{使}\ \mathrm{PT}_{\ell}\neq \mathrm{id},\\
  H_2 &: \Rightarrow\ \text{取任意}\ \gamma_1,\ \text{令}\ \gamma_2:=\ell\ast\gamma_1,\\
  H_3 &: \gamma_1,\gamma_2\ \text{端点相同},\\
  H_4 &: \Rightarrow\ \mathrm{PT}_{\gamma_2}=\mathrm{PT}_{\gamma_1}\circ \mathrm{PT}_{\ell}\neq \mathrm{PT}_{\gamma_1}.
\end{aligned}
\]
\end{Certificate}

\begin{proof}[证明（谓词因果链）]
由证书~\ref{cert:holonomy-implies-path-class-relevant}：
$H_1$ 给出非平凡回路；按 $H_2$ 用回路拼接构造 $\gamma_2$；
由 $H_3$ 得两条路径端点相同；
由 $H_4$ 得其搬运算子不同。
故引理成立。证毕。
\end{proof}

\begin{definition}[C 级：规范结构切换（跨相/换有效理论）]
称“切换规范场”为 C 级，若不再局限于同一规范结构 $\mathcal{M}$，
而是发生模型对象的切换：
\[
  \mathcal{M}_i=(\Bundle_i\to\Base_i,\ \Group_i,\ \Conn_i)\ \rightsquigarrow\
  \mathcal{M}_j=(\Bundle_j\to\Base_j,\ \Group_j,\ \Conn_j).
\]
其物理含义可表述为：有效自由度--约束--冗余的组织方式发生重组（相变或进入不同亚稳盆地），
因而“规则”改变；这不是 A 级的换 gauge fixing，也不是 B 级的换背景连接，而是换\textbf{规范结构本身}。
\end{definition}

\begin{definition}[贴图与扭曲搬运（模型间映射）]
当发生 C 级切换 $\mathcal{M}_i\to\mathcal{M}_j$ 时，
需要构造贴图/搬运映射
\[
  T_{i\to j}:\ \mathcal{C}_i\to \mathcal{C}_j,
\]
把旧模型下的构造代表元映射为新模型下的可行初值。
该映射至少应保持两类可审计条件：
\begin{enumerate}
  \item \textbf{判据可解释}：关键可观测判据与保障约束在新模型中仍可解释、可检查；
  \item \textbf{可行可续接}：输出落入新模型可执行集合，保证实时性与连续推进（对齐第~\ref{sec:twisted-search-transport} 节的扭曲搬运口径）。
\end{enumerate}
\end{definition}

\begin{remark}[三个层级与“相空间穿梭/同伦扭曲”的对应]
\begin{enumerate}
  \item A 级：同一相内换规范/换坐标图；不改变物理不变量，只改善表述与数值条件；
  \item B 级：同一规范结构内改变背景连接；若出现 holonomy，则路径同伦类进入决策变量（同伦扭曲的来源之一）；
  \item C 级：跨相切换导致规范结构改变；必须做贴图与扭曲搬运，否则相内最优解无法自然延拓到相外。
\end{enumerate}
\end{remark}

\subsection{公理降级：规范切换序列的分解与归纳证书}
\label{sec:gauge-switching-decomposition}

\begin{theorem}[公理降级：任意“切换规范场”序列可分解为 A/B 内部演化与 C 级结构切换（按步数归纳）]
\label{thm:gauge-switching-sequence-decomposition}
考虑一个长度为 $T$ 的“切换规范场”序列 $\{u_t\}_{t=0}^{T-1}$，
其中每一步 $u_t$ 不是 A 级（规范选择变换）、就是 B 级（背景连接迁移）、就是 C 级（规范结构切换）。
记模型序列为 $\{\mathcal{M}_t\}$。
则：
\begin{enumerate}
  \item $\mathcal{M}_t$ 的改变仅发生在 C 级步上；A/B 级步不改变模型类型；
  \item 因此序列可被分段为若干“同一模型内的 A/B 演化段”，并以 C 级步为段间贴图接口。
\end{enumerate}
\end{theorem}

\begin{Certificate}[数学归纳法证书：按步数 $T$ 归纳的分段分解]
\label{cert:gauge-switching-sequence-decomposition}
定义谓词 $P(T)$ 为：
“任意长度为 $T$ 的序列，其模型类型只在 C 级步处发生变化，且可按 C 级步将序列分割为 A/B 演化段。”
证书化要求为：
\[
  P(1)\ \text{成立},\qquad
  (\forall T\in\NN_{\ge 1})\ \bigl(P(T)\Rightarrow P(T+1)\bigr).
\]
\end{Certificate}

\begin{proof}[证明（数学归纳法证书；谓词因果链）]
由证书~\ref{cert:gauge-switching-sequence-decomposition}：
\begin{itemize}
  \item 基步：$T=1$。若 $u_0$ 为 A/B 级，则模型不变；若 $u_0$ 为 C 级，则模型在此步切换。两种情形均满足“模型仅在 C 级处改变”，并可视为只有一个段。故 $P(1)$ 成立。
  \item 归纳步：假设 $P(T)$ 成立。考虑长度为 $T+1$ 的序列 $u_{0:T}$。
  其前 $T$ 步按归纳假设可分段且模型改变只发生在 C 级步上。第 $T$ 步：
  \begin{itemize}
    \item 若 $u_T$ 为 A/B 级，则模型类型不变，仅延长末段；
    \item 若 $u_T$ 为 C 级，则在该处发生模型切换，并在此处新增一个分段边界。
  \end{itemize}
  因此 $P(T+1)$ 成立。
\end{itemize}
由数学归纳法得对任意有限 $T$ 结论成立。证毕。
\end{proof}

\begin{corollary}[推论：若不发生 C 级切换，则全过程属于同一规范结构（A/B 口径）]
\label{cor:no-structure-switch-implies-same-model}
若序列 $\{u_t\}$ 中不含任何 C 级步，则对任意 $T$ 有 $\mathcal{M}_T=\mathcal{M}_0$；
此时所谓“切换规范场”只能是 A 级（换写法）或 B 级（换背景连接），而不等价于跨相改变有效理论。
\end{corollary}

\begin{Certificate}[证书：无 C 级步 $\Rightarrow$ 模型常值]
\label{cert:no-structure-switch-implies-same-model}
证书登记：A/B 级步不改变模型类型；若全序列仅含 A/B 级步，则模型类型在每一步保持不变，故 $\mathcal{M}_T=\mathcal{M}_0$。
\end{Certificate}

\begin{proof}[证明（谓词因果链）]
由证书~\ref{cert:no-structure-switch-implies-same-model}，令谓词链为：
\[
\begin{aligned}
  D_1 &: \text{A/B 级}\Rightarrow \mathcal{M}\ \text{不变},\\
  D_2 &: \text{全序列仅含 A/B}\Rightarrow \mathcal{M}_T=\mathcal{M}_0,\\
  D_3 &: \Rightarrow\ \text{“切换规范场”不等价于跨相}.
\end{aligned}
\]
由 $D_1\Rightarrow D_2\Rightarrow D_3$ 得推论结论。证毕。
\end{proof}

\subsection{工程规则：触发条件、贴图要求与推荐表述}
\label{sec:gauge-switching-engineering-rules}

\begin{remark}[规则 1：相内控制/搜索 = 固定规范结构下的续接与局部优化]
在同一相域内，应优先把“切换规范场”解释为 A/B 级：
通过坐标图重参数化改善数值条件（A 级），并用平行搬运/热启动保持连续可行（B 级）。
这与第~\ref{sec:twisted-homotopy-two-level-search} 节的相内搜索与第~\ref{sec:continuation-rt-optimization} 节的续接型滚动优化一致。
\end{remark}

\begin{remark}[规则 2：出现“不可延拓”迹象时触发 C 级切换]
当观测/审计字段显示仅靠 A/B 级不足以维持可行或持续改进时，应触发 C 级切换（跨相/换模型），典型信号包括：
滞回显著或回线不闭合、局部拟合残差结构突变、噪声分布突变、关键约束变为不可满足或代价急剧增大、序参量跃迁/响应不可微等。
此时“换写法/换连接”不足以解决结构性障碍，需要通过工程算子改变相空间包线并切换模型口径。
\end{remark}

\begin{remark}[规则 3：跨相之后必须做贴图与扭曲搬运，不能直接复用旧解]
跨相切换相当于 $\mathcal{M}_i\to\mathcal{M}_j$，旧模型的代表元解一般\textbf{不}能自然延拓到新模型。
必须构造 $T_{i\to j}$ 把旧解映射为新相域的可行初值，并把“保障链/证据链”的可检查性纳入映射约束，
否则搜索/控制会被伪信号牵引到错误的同伦类。
\end{remark}

\begin{remark}[推荐改写（更严格且不改变核心意思）]
为避免歧义，可将
“切换相空间相当于切换规范场，等价于切换实际粒子的虚拟组织形态从而改变规范场规则”
改写为：
\begin{quote}
跨相切换相空间等价于切换系统的有效自由度与约束组织方式；在纤维丛表述中，这通常体现为连接的改变，甚至规范结构（场内容/约束/转移函数）的改变，而不仅是规范选择的改变。
\end{quote}
该表述保留了“规则会变”的核心意思，同时明确区分 A/B 与 C 两类根本不同的“切换”。
\end{remark}

\section{守恒底座与涌现切换：代数—拓扑—边缘算子三元组}
\label{sec:emergent-switching-algebra-topology-edge}

\begin{remark}[定位（把“有效理论切换”落到可推演的结构对象）]
在“有效理论/涌现描述”的口径下，可以把“相空间切换”的核心意思进一步钉成三层：
\begin{enumerate}
  \item \textbf{守恒底座固定}：能量与相关守恒荷等全局约束作为底层不变量，不随相域切换而被违反；
  \item \textbf{变化的是有效自由度的组织方式}：体现在可观测算子所满足的\textbf{代数规则}与允许的\textbf{拓扑扇区/约束}发生改变；
  \item \textbf{代数与拓扑的改变对应规范结构与上同调微分的改变}：从而“规范场形态/规则”随之改变，并可被用于解释“上同调修复边缘算子”的工程含义。
\end{enumerate}
本章把上述三层落成一个可计算的三元组模型，并给出“边缘修复算子”在上同调口径下的形式化位置。
\end{remark}

\subsection{守恒底座：能量与守恒荷的形式化口径}
\label{sec:conserved-bedrock}

\begin{definition}[守恒底座（不变量集合）]
\label{def:conserved-bedrock}
令 $\mathcal{Q}$ 为一组全局不变量（例如能量与相关守恒荷），其中每个 $Q\in\mathcal{Q}$ 视为函数
\[
  Q:\mathcal{S}\to\RR.
\]
称系统在物理更新算子 $O_t:\mathcal{S}\to\mathcal{S}$ 下满足守恒底座，若对任意 $t$ 与任意 $s\in\mathcal{S}$ 有
\[
  Q\bigl(O_t(s)\bigr)=Q(s)\qquad(\forall Q\in\mathcal{Q}).
\]
\end{definition}

\begin{remark}[关于“粒子数守恒”的更稳妥表述]
在某些有效理论（例如含配对凝聚的表述）中，“粒子数”在有效描述里不一定以严格守恒的形式出现；
但守恒底座更一般地指“能量与相关守恒荷”的不变量集合 $\mathcal{Q}$。
因此本章将“守恒前提”统一表述为 $\mathcal{Q}$ 固定，而不把“粒子数”当作必须守恒的唯一对象。
\end{remark}

\begin{proposition}[命题：守恒底座为跨相/跨模型切换提供全局一致的约束口径]
\label{prop:conserved-bedrock-provides-global-constraints}
若系统满足守恒底座（定义~\ref{def:conserved-bedrock}），则任意由物理更新算子复合得到的轨迹
\[
  s_T := O_{T-1}\circ\cdots\circ O_0(s_0)
\]
满足 $Q(s_T)=Q(s_0)$（$\forall Q\in\mathcal{Q}$）。
因此在“相域切换/模型切换”的上层描述中，可将 $\mathcal{Q}$ 视为\textbf{全局一致的硬约束}，而把“变化”限制在有效自由度组织方式与约束结构上。
\end{proposition}

\begin{Certificate}[证书：不变量沿复合算子保持]
\label{cert:conserved-bedrock-provides-global-constraints}
证书登记如下谓词链：
\[
\begin{aligned}
  Q_1 &: Q(O_t(s))=Q(s)\quad(\forall t,\forall s),\\
  Q_2 &: \Rightarrow\ Q(O_{T-1}\circ\cdots\circ O_0(s_0))=Q(s_0),\\
  Q_3 &: \Rightarrow\ \mathcal{Q}\ \text{可作为跨相/跨模型的全局硬约束口径}.
\end{aligned}
\]
\end{Certificate}

\begin{proof}[证明（谓词因果链）]
由证书~\ref{cert:conserved-bedrock-provides-global-constraints}：
对复合算子逐步应用 $Q_1$ 得 $Q_2$；由 $Q_2$ 可把 $\mathcal{Q}$ 视为与相域无关的不变量约束，从而得到 $Q_3$。
命题成立。证毕。
\end{proof}

\subsection{有效模型三元组：代数、拓扑与上同调微分}
\label{sec:effective-model-triple}

\begin{definition}[有效模型三元组（代数—拓扑—微分）]
对每个相域/模型口径 $i$，引入三元组
\[
  \mathcal{M}_i := \bigl(\mathfrak{A}_i,\ \Bundle_i\xrightarrow{\pr_i} X_i,\ \partial_i\bigr),
\]
其中：
\begin{enumerate}
  \item $\mathfrak{A}_i$ 为可观测量/算子集合在该相域内满足的代数结构（生成元、关系式、交换/非交换规则、约束理想等）；
  \item $\Bundle_i\xrightarrow{\pr_i} X_i$ 为定义在参数/慢变量空间 $X_i$ 上的主丛（或等价几何承载体），其结构群可理解为某个母规范群 $\Group$ 的子结构；
  \item $\partial_i$ 为用于“可修复/不可修复”分类的上同调微分（共边界算子），它决定了何谓闭、何谓边界（可修复）、何谓非平凡上同调类（结构性剩余）。
\end{enumerate}
\end{definition}

\begin{remark}[与前文三层级（A/B/C）的一致性]
在第~\ref{sec:gauge-switching-three-levels} 节的三层级口径下：
\begin{enumerate}
  \item A 级（规范选择）与 B 级（背景连接迁移）均不改变 $\mathcal{M}_i$ 的类型：它们只是在同一三元组内改变表示或改变连接代表；
  \item C 级（规范结构切换）对应至少一项发生改变：$\mathfrak{A}_i$、$\Bundle_i\to X_i$ 的拓扑/约束扇区，或 $\partial_i$ 的“修复规则”发生改变。
\end{enumerate}
因此“相空间切换导致规则变化”的严格落点，是 $\mathcal{M}_i\to\mathcal{M}_j$ 的三元组切换，而不是 A 级意义下的换写法。
\end{remark}

\begin{theorem}[公理降级：相空间切换可分解为守恒底座约束下的三元组切换]
\label{thm:phase-switching-as-model-triple-switching}
在守恒底座 $\mathcal{Q}$ 固定的前提下，任意“相空间切换/有效理论切换”都可被形式化为
\[
  \mathcal{M}_i \ \longrightarrow\ \mathcal{M}_j,
\]
其中变化仅允许发生在三元组分量 $\mathfrak{A}$、$\Bundle\to X$、$\partial$ 上；
而 $\mathcal{Q}$ 作为全局硬约束在切换前后保持不变。
\end{theorem}

\begin{Certificate}[证书：守恒底座固定 + 三元组分量变化的封闭性]
\label{cert:phase-switching-as-model-triple-switching}
证书登记如下谓词链：
\[
\begin{aligned}
  T_1 &: \mathcal{Q}\ \text{固定}\Rightarrow \text{切换前后必须满足同一组全局不变量约束},\\
  T_2 &: \text{有效理论差异}\Rightarrow \text{用}\ (\mathfrak{A},\Bundle\to X,\partial)\ \text{三类结构承载},\\
  T_3 &: \Rightarrow\ \mathcal{M}_i\to\mathcal{M}_j\ \text{给出相空间切换的可推演表述}.
\end{aligned}
\]
\end{Certificate}

\begin{proof}[证明（谓词因果链）]
由证书~\ref{cert:phase-switching-as-model-triple-switching}：
$T_1$ 给出切换的底座约束；$T_2$ 指定“变化”被三元组承载；合并得 $T_3$，即相空间切换可用三元组切换表述。
定理成立。证毕。
\end{proof}

\subsection{边缘修复算子：改变可修复性的微分改造}
\label{sec:edge-repair-operator}

\begin{definition}[边缘修复算子与改造微分]
在某一相域 $i$ 的共链复形 $(C_i^\bullet,\partial_i)$ 上，称 $\partial_{\mathrm{edge},i}$ 为\textbf{边缘修复算子}，若它是次数 $+1$ 的线性算子，
并定义改造微分为
\[
  \partial_i' := \partial_i + \partial_{\mathrm{edge},i}.
\]
其工程语义为：把“只能在边界/界面处起作用的修复机制（补偿自由度/保障通道/边界条件改变）”显式并入上同调判据，从而改变“可修复/不可修复”的分类规则。
\end{definition}

\begin{lemma}[引理：边缘修复算子给出合法微分的充分条件]
\label{lem:edge-repair-yields-valid-differential}
若 $\partial_i$ 已满足 $\partial_i\circ \partial_i=0$，并且 $\partial_{\mathrm{edge},i}$ 满足
\[
  \partial_{\mathrm{edge},i}\circ\partial_{\mathrm{edge},i}=0,\qquad
  \partial_i\circ\partial_{\mathrm{edge},i} + \partial_{\mathrm{edge},i}\circ\partial_i=0,
\]
则 $\partial_i'=\partial_i+\partial_{\mathrm{edge},i}$ 仍满足 $(\partial_i')^2=0$，从而 $(C_i^\bullet,\partial_i')$ 是合法的共链复形。
\end{lemma}

\begin{Certificate}[证书：展开平方并逐项归零]
\label{cert:edge-repair-yields-valid-differential}
证书登记如下谓词链：
\[
\begin{aligned}
  E_1 &: (\partial_i')^2=(\partial_i+\partial_{\mathrm{edge},i})^2\\
      &\Rightarrow (\partial_i')^2=\partial_i^2+\partial_i\partial_{\mathrm{edge},i}+\partial_{\mathrm{edge},i}
      \partial_i+\partial_{\mathrm{edge},i}^2,\\
  E_2 &: \partial_i^2=0\ \wedge\ \partial_{\mathrm{edge},i}^2=0\ \wedge\ (\partial_i\partial_{\mathrm{edge},i}
    +\partial_{\mathrm{edge},i}\partial_i)=0,\\
  E_3 &: \Rightarrow (\partial_i')^2=0.
\end{aligned}
\]
\end{Certificate}

\begin{proof}[证明（谓词因果链）]
由证书~\ref{cert:edge-repair-yields-valid-differential}：
按 $E_1$ 展开 $(\partial_i')^2$；由给定条件得到 $E_2$；代入即得 $E_3$。
引理成立。证毕。
\end{proof}

\begin{proposition}[命题：边缘修复使“不可修复类”转化为“可修复项”的判据]
\label{prop:edge-repair-trivializes-obstruction-class}
令 $d\in C_i^1$ 为某一干扰/负算子见证。
若在原复形上 $[d]\neq 0\in H^1(C_i^\bullet,\partial_i)$，但在改造复形上 $[d]=0\in H^1(C_i^\bullet,\partial_i')$，
则存在 $c\in C_i^0$ 使得
\[
  d=\partial_i'(c),
\]
从而 $d$ 在改造规则下成为可修复项（上边界）。
\end{proposition}

\begin{Certificate}[证书：上同调类平凡 $\Leftrightarrow$ 存在上边界表达]
\label{cert:edge-repair-trivializes-obstruction-class}
证书登记如下谓词链：
\[
\begin{aligned}
  R_1 &: [d]=0\ \text{在}\ H^1(C_i^\bullet,\partial_i')\\
      &\Rightarrow d\in \operatorname{im}(\partial_i':C_i^0\to C_i^1),\\
  R_2 &: \Rightarrow (\exists c\in C_i^0)\ d=\partial_i'(c),\\
  R_3 &: \Rightarrow d\ \text{为上边界}\Rightarrow d\ \text{可修复}.
\end{aligned}
\]
\end{Certificate}

\begin{proof}[证明（谓词因果链）]
由证书~\ref{cert:edge-repair-trivializes-obstruction-class}：
上同调类平凡的定义给出 $R_1$；由像空间定义得 $R_2$；再由“上边界 = 可修复”判据得 $R_3$。
命题成立。证毕。
\end{proof}

\begin{corollary}[推论：调整修复微分等价于改变“可消/不可消”的分类规则]
\label{cor:edge-repair-changes-repairability-classification}
在不改变守恒底座 $\mathcal{Q}$ 的前提下，若通过边缘修复算子把 $\partial_i$ 改造为 $\partial_i'$，
则“哪些干扰属于可修复项、哪些属于结构性剩余”的判据随之改变。
因此“突破边界的增强”（第~\ref{sec:ptopb-family-and-two-operators} 节）可被理解为：
通过改变模型三元组中的 $\partial$（以及必要时 $\mathfrak{A}$ 与拓扑扇区）使障碍类平凡化。
\end{corollary}

\begin{Certificate}[证书：判据依赖于微分 $\partial$]
\label{cert:edge-repair-changes-repairability-classification}
证书登记：上同调 $H^1(C^\bullet,\partial)$ 的定义依赖于微分 $\partial$；
因此 $\partial\to\partial'$ 会改变闭/边界/上同调类，从而改变“可修复/不可修复”的分类规则。
\end{Certificate}

\begin{proof}[证明（谓词因果链）]
由证书~\ref{cert:edge-repair-changes-repairability-classification}：
“可修复/不可修复”由上同调类是否平凡给出，而上同调由微分定义，故 $\partial$ 的改变会改变分类规则。
推论成立。证毕。
\end{proof}

\begin{remark}[相对上同调与贴锥口径（对“边缘修复”的结构化表达）]
当系统存在边界/界面 $\partial X$ 时，很多“边缘修复”更自然地对应：
\begin{enumerate}
  \item 从绝对上同调 $H^\ast(X)$ 转到相对上同调 $H^\ast(X,\partial X)$；
  \item 或把“补偿自由度/保障机制”并入扩张复形（例如贴锥 \emph{mapping cone}）以显式改变长正合列中的连接映射。
\end{enumerate}
本笔记的工程口径不要求固定采用某一种实现方式；其共同点是：通过改变复形/微分，把原先闭但非边界的障碍类转化为边界（可修复项）。
与第~\ref{sec:relative-cohomology-augmentation} 节的“相对上同调（增强差分口径）”一致。
\end{remark}

\subsection{对接同伦扭曲搜索/控制：模型切换与贴图接口的证据链}
\label{sec:triple-switching-aligns-with-twisted-homotopy}

\begin{remark}[对接规则（算法实现口径）]
把同伦扭曲搜索/控制（第~\ref{sec:twisted-homotopy-two-level-search} 节与第~\ref{sec:twisted-homotopy-real-time-control} 节）落到本章三元组口径，
  可得到一条可直接用于实现的规则：
\begin{enumerate}
  \item \textbf{相内续接}：保持 $\mathcal{M}_i$ 不变（尤其是 $\partial_i$ 与代数/拓扑口径不变），仅在同一模型内做热启动与局部优化；
  \item \textbf{跨相切换}：当观测与审计字段指示存在结构性障碍类（不可修复）时，触发 $\mathcal{M}_i\to\mathcal{M}_j$；
  \item \textbf{贴图/搬运接口}：跨相后必须给出 $T_{i\to j}$（第~\ref{sec:twisted-search-transport} 节）把旧构造代表元映射为新模型下的可行初值，并保持证据链可解释、可检查。
\end{enumerate}
在该规则下，“规范群内切换”可被理解为：在某个母群框架内，不同相域对应不同的代数/拓扑扇区与微分规则分支；跨相是分支切换而非单纯的 gauge fixing。
\end{remark}

\begin{remark}[落地清单（把“规则变化”压缩为最小可登记集合）]
若要把“涌现切换/规范群内切换”进一步落实为可计算对象，建议把“规则变化”的最小集合登记为三项：
\begin{enumerate}
  \item 代数规则变化：交换/非交换结构变化、约束理想变化，或生成元集合变化；
  \item 拓扑约束变化：允许的缺陷类型变化、边界通量规则变化，或主丛拓扑类变化；
  \item 边缘修复算子：对应哪类工程机制（新增边界自由度、边界补偿项、边界测量与反馈通道、界面结构重构等）。
\end{enumerate}
这三项一旦证书化登记，即可直接作为“模型切换 $\mathcal{M}_i\to\mathcal{M}_j$”与“贴图接口 $T_{i\to j}$”的实现输入。
\end{remark}

\section{全局守恒与局部破缺：边缘算子闭合子系统创生/湮灭}
\label{sec:global-conservation-local-breaking}

\begin{remark}[核心断言（闭合总系统口径）]
在把系统视为\textbf{闭合总系统}（或把外部库显式并入系统的边界自由度）的前提下，可以同时成立两句话：
\begin{enumerate}
  \item \textbf{守恒量固定}：总系统的能量与相关守恒荷满足全局守恒（见定义~\ref{def:conserved-bedrock}）；
  \item \textbf{创生/湮灭破缺成立}：在子系统或有效准粒子层面，“粒子数/数算符”不再守恒（出现创生/湮灭项或等价的数守恒破缺），但其缺口由边界交换/补偿通道在全局上闭合。
\end{enumerate}
二者不矛盾；把它们粘合起来的关键结构是：\textbf{边界交换项}与\textbf{边缘修复算子}。
\end{remark}

\subsection{兼容条件：全局守恒与子系统非守恒的同时成立}
\label{sec:compatibility-global-vs-local}

\begin{definition}[闭合总系统与边界库（交换通道）]
\label{def:closed-total-system-reservoir}
将总系统分解为“体内子系统 $X$”与“边界库/外部库 $R$”（可显式建模为边界自由度），并令总哈密顿量（或等价生成元）分解为
\[
  H_{\mathrm{tot}} = H_X + H_R + H_{\mathrm{edge}},
\]
其中 $H_{\mathrm{edge}}$ 为体与库之间的边界交换/耦合项（顺序敏感时可视为非交换构造词的一段）。
\end{definition}

\begin{remark}[工程化兼容条件（账本口径）]
把“守恒量固定”解释为全局守恒，需要显式给出两类账本检查：
\begin{enumerate}
  \item \textbf{能量账本}：若 $H_{\mathrm{tot}}$ 对时间不显含（或将驱动源也纳入 $H_R$ 的能量账本），则总能量守恒；若边界耦合显式随时间变化，则除非把驱动源计入，否则不能声称能量守恒。
  \item \textbf{守恒荷账本}：明确声称守恒的“相关守恒荷”是哪一个（电荷/某个全局量子数/拓扑荷等），并证书化登记其在边界交换项上的补偿通道（流入/流出/异常流入等）。
\end{enumerate}
\end{remark}

\begin{theorem}[公理降级：全局守恒下的局部非守恒闭合（按子系统数归纳）]
\label{thm:global-conservation-compat-local-nonconservation}
考虑由 $n$ 个部件组成的闭合总系统，记其“局部守恒荷”算符为 $Q_1,\dots,Q_n$，并令
\[
  Q_{\mathrm{tot}}:=\sum_{k=1}^{n} Q_k.
\]
若总生成元满足 $[H_{\mathrm{tot}},Q_{\mathrm{tot}}]=0$，则全局守恒成立：
\[
  \frac{\mathrm{d}}{\mathrm{d}t}Q_{\mathrm{tot}}(t)=0,
\]
同时允许局部非守恒：
\[
  \frac{\mathrm{d}}{\mathrm{d}t}Q_k(t)=\mathrm{i}[H_{\mathrm{tot}},Q_k(t)]\ \text{可不为零},
\]
但其总和必闭合：
\[
  \sum_{k=1}^{n}\frac{\mathrm{d}}{\mathrm{d}t}Q_k(t)=0.
\]
\end{theorem}

\begin{Certificate}[数学归纳法证书：按部件数 $n$ 归纳的闭合性]
\label{cert:global-conservation-compat-local-nonconservation}
定义谓词 $P(n)$ 为：
“对任意由 $n$ 个部件组成的闭合系统，只要 $[H_{\mathrm{tot}},Q_{\mathrm{tot}}]=0$，则 $\sum_{k=1}^{n}\frac{\mathrm{d}}{\mathrm{d}t}Q_k(t)=0$
  成立。”
证书化要求为：
\[
  P(1)\ \text{成立},\qquad
  (\forall n\in\NN_{\ge 1})\ \bigl(P(n)\Rightarrow P(n+1)\bigr).
\]
\end{Certificate}

\begin{proof}[证明（数学归纳法证书；谓词因果链）]
由证书~\ref{cert:global-conservation-compat-local-nonconservation}：
\begin{itemize}
  \item 基步：$n=1$ 时，$Q_{\mathrm{tot}}=Q_1$，由 $[H_{\mathrm{tot}},Q_{\mathrm{tot}}]=0$ 得 $\frac{\mathrm{d}}{\mathrm{d}t}
    Q_1(t)=0$，故闭合性成立。
  \item 归纳步：假设 $P(n)$ 成立。对 $n+1$ 个部件，令
  \[
    Q_{\mathrm{blk}}:=\sum_{k=1}^{n}Q_k,\qquad Q_{\mathrm{tot}}=Q_{\mathrm{blk}}+Q_{n+1}.
  \]
  由 $[H_{\mathrm{tot}},Q_{\mathrm{tot}}]=0$ 得
  \[
    0=\frac{\mathrm{d}}{\mathrm{d}t}Q_{\mathrm{tot}}(t)=\frac{\mathrm{d}}{\mathrm{d}t}Q_{\mathrm{blk}}(t)
      +\frac{\mathrm{d}}{\mathrm{d}t}Q_{n+1}(t).
  \]
  而 $\frac{\mathrm{d}}{\mathrm{d}t}Q_{\mathrm{blk}}(t)=\sum_{k=1}^{n}\frac{\mathrm{d}}{\mathrm{d}t}Q_k(t)$；
  故 $\sum_{k=1}^{n+1}\frac{\mathrm{d}}{\mathrm{d}t}Q_k(t)=0$，即 $P(n+1)$ 成立。
\end{itemize}
由数学归纳法得对任意 $n$ 结论成立。证毕。
\end{proof}

\subsection{边界通量与对易关系：创生/湮灭破缺的可观测见证}
\label{sec:boundary-flux-and-commutator-witness}

\begin{definition}[体内守恒荷与边界通量公式（连续性口径）]
设子系统区域为 $X$，密度为 $\rho(\mathbf{x},t)$，流密度为 $\mathbf{j}(\mathbf{x},t)$，定义体内荷为
\[
  Q_X(t):=\int_{X}\rho(\mathbf{x},t)\,\mathrm{d}^d x.
\]
对连续性方程积分可得形式恒等式：
\[
  \frac{\mathrm{d}}{\mathrm{d}t}Q_X(t)
  =-\int_{\partial X}\mathbf{j}\cdot \mathrm{d}\mathbf{S}
  + \int_X s(\mathbf{x},t)\,\mathrm{d}^d x,
\]
其中 $s$ 可解释为体内源项；当把源项显式并入边界库/边界自由度后，可将其视为边界交换通道的等效记账。
\end{definition}

\begin{proposition}[命题：子系统数算符非守恒由边界交换项见证]
\label{prop:subsystem-number-nonconservation-witnessed-by-edge}
在定义~\ref{def:closed-total-system-reservoir} 的分解下，令 $N_X$ 为子系统的数算符（或有效准粒子数算符），并假设 $N_X$ 无显式时间依赖。
若 $[H_X,N_X]=0$ 且 $[H_R,N_X]=0$，则
\[
  \frac{\mathrm{d}}{\mathrm{d}t}N_X(t)=\mathrm{i}[H_{\mathrm{edge}},N_X(t)].
\]
因此 $[H_{\mathrm{edge}},N_X]\neq 0$ 就是“创生/湮灭破缺”（子系统数守恒破缺）的可观测见证；
而在全局守恒前提下（例如 $[H_{\mathrm{tot}},N_X+N_R]=0$），该非守恒缺口由库侧 $N_R$ 的变化补偿并闭合。
\end{proposition}

\begin{Certificate}[证书：海森堡方程与对易分解]
\label{cert:subsystem-number-nonconservation-witnessed-by-edge}
证书登记如下谓词链：
\[
\begin{aligned}
  N_1 &: \frac{\mathrm{d}}{\mathrm{d}t}N_X(t)=\mathrm{i}[H_{\mathrm{tot}},N_X(t)],\\
  N_2 &: H_{\mathrm{tot}}=H_X+H_R+H_{\mathrm{edge}},\\
  N_3 &: [H_X,N_X]=[H_R,N_X]=0\Rightarrow [H_{\mathrm{tot}},N_X]=[H_{\mathrm{edge}},N_X],\\
  N_4 &: \Rightarrow \frac{\mathrm{d}}{\mathrm{d}t}N_X(t)=\mathrm{i}[H_{\mathrm{edge}},N_X(t)].
\end{aligned}
\]
\end{Certificate}

\begin{proof}[证明（谓词因果链）]
由证书~\ref{cert:subsystem-number-nonconservation-witnessed-by-edge}：
海森堡方程给出 $N_1$；代入分解得 $N_2$；
由对易条件得 $N_3$；合并得 $N_4$。命题成立。证毕。
\end{proof}

\begin{remark}[两种常见语义：子系统非守恒 vs 有效准粒子非守恒]
“创生/湮灭破缺成立”通常至少对应两种语义之一：
\begin{enumerate}
  \item \textbf{子系统非守恒}：$X$ 与 $R$ 存在交换通道，$N_X$ 不守恒但 $N_X+N_R$ 守恒；
  \item \textbf{有效准粒子非守恒}：微观守恒荷仍成立，但低能自由度重组（例如配对凝聚/安德烈夫过程），使“准粒子数”不再是好量子数。
\end{enumerate}
两者都与“全局守恒底座固定”（定义~\ref{def:conserved-bedrock}）兼容，但需要在证书中明确是哪个层面发生非守恒。
\end{remark}

\subsection{上同调修复：把“破缺项”纳入可修复结构}
\label{sec:cohomology-repair-closes-breaking}

\begin{remark}[从“源项”到“上边界”：边缘修复的上同调解释]
在本笔记的同调代数口径下，“局部看起来像源项/破缺项”的对象可被登记为某个共链 $d\in C^1$。
若在原微分 $\partial$ 下它是闭但非边界（$[d]\neq 0$），则属于结构性不可修复类；
而引入边缘修复算子把微分改造为 $\partial'=\partial+\partial_{\mathrm{edge}}$ 后，
若障碍类在新结构下平凡化，则存在 $c\in C^0$ 使
\[
  d=\partial'(c),
\]
从而 $d$ 变为上边界并被视为可修复项。
该机制已在命题~\ref{prop:edge-repair-trivializes-obstruction-class} 与推论~\ref{cor:edge-repair-changes-repairability-classification} 中证书化登记。
\end{remark}

\begin{remark}[工程含义：边缘算子把“局部破缺”闭合为“全局守恒”]
把“边界库/补偿自由度”并入系统边界，相当于扩张了可用的修复机制与约束结构；
在守恒账本层面，表现为把局部非守恒的缺口转化为边界交换通量；
在上同调层面，表现为把原先非平凡的障碍类转化为边界（可修复项）。
这正是“边缘算子把全局守恒与局部破缺粘合起来”的形式化落点。
\end{remark}

\subsection{相变触发与算法接口：作为跨相算子进入同伦扭曲框架}
\label{sec:edge-trigger-phase-change-and-algorithm-interface}

\begin{remark}[相变触发：改变拓扑扇区与有效代数规则]
当“引入/调整边缘算子”同时改变允许的拓扑扇区（边界通量规则、缺陷进入/退出、holonomy 约束等）与有效代数规则（生成元与关系式）时，
系统的低能自由度可能发生重组，表现为相变/相域切换。
在“母规范群框架”下，这更像是在不同扇区/分支之间切换，而不是对同一物理态做纯 A 级规范变换（见第~\ref{sec:gauge-switching-three-levels} 节）。
\end{remark}

\begin{remark}[推荐改写（更严格且不改变含义）]
为避免把“全局守恒”与“局部破缺”混为矛盾，可将
\begin{quote}
通过附加或调整边缘算子调整规范场产生相变，使粒子创生或湮灭破缺成立属于“守恒量（能量与相关守恒荷）固定”
\end{quote}
改写为：
\begin{quote}
在总系统能量与相关守恒荷保持守恒的前提下，引入/调整边缘算子等价于改变边界补偿自由度与拓扑约束（即改变上同调“可修复结构”）；这可导致有效自由度重组与相变，并使子系统或有效准粒子层面出现创生/湮灭项（数守恒破缺），
  而其守恒缺口由边界交换/补偿通道在全局上闭合。
\end{quote}
\end{remark}

\begin{remark}[对接同伦扭曲搜索/控制（可执行接口）]
把“边缘算子调整”纳入第~\ref{sec:twisted-homotopy-two-level-search} 节与第~\ref{sec:twisted-homotopy-real-time-control} 节的框架，可按如下口径实现：
\begin{enumerate}
  \item 将“边缘算子调整”登记为一类\textbf{跨相触发算子}：它改变拓扑扇区与修复微分（$\partial\to\partial'$），因此改变相空间本身；
  \item 将“守恒底座固定”（定义~\ref{def:conserved-bedrock}）写成算法的硬约束（能量预算、净荷预算等）；
  \item 将“创生/湮灭项”写成子系统开放项，并要求由边界补偿通道在证书边界内闭合；
  \item 跨相后用贴图/搬运算子 $T_{i\to j}$（第~\ref{sec:twisted-search-transport} 节）保持热启动可行与证据链可审计。
\end{enumerate}
\end{remark}

\section{扇区化三层统一：公理层—扇区层—算法层的对象—映射—切换规则}
\label{sec:axiom-sector-algorithm-unified-framework}

\begin{remark}[定位（把三条主张压成可推演、可实现的接口）]
可将“公理层—扇区层—算法层”的三条主张压缩为一组\textbf{对象—映射—切换规则}：
\begin{enumerate}
  \item \textbf{公理层}：守恒量固定，定义全局可行母空间（控制—搜索的母空间）；
  \item \textbf{扇区层}：通过边缘算子调节改变修复结构与拓扑扇区，从而在同一母规范群框架下切换到不同的非规范等价扇区；
  \item \textbf{算法层}：每个扇区对应一套“同伦扭曲搜索/控制配置栈”，并在切换时用贴图/搬运算子保持可行与证据链一致。
\end{enumerate}
本章将上述三条固化为可计算对象，并给出最小切换逻辑与实现输入表。
\end{remark}

\subsection{公理层：守恒量固定与全局可行母空间}
\label{sec:axiom-layer-global-feasible-space}

\begin{definition}[全局不变量约束集与可行母空间]
令 $\mathcal{I}=\{I_1,\dots,I_m\}$ 为一组全局不变量约束（可取为守恒底座 $\mathcal{Q}$ 的具体化，见定义~\ref{def:conserved-bedrock}）。
令 $\lambda$ 为相空间基底变量，$c$ 为工程构造配置，$\Conn$ 为连接/规范势（亦可用 $A$ 表示），并允许附加其他可审计变量（以 $\ldots$ 表示）。
定义全局可行母空间为
\[
  \Omega_{\mathrm{glob}}
  :=
  \Bigl\{\,(\lambda,c,\Conn,\ldots)\ \Big|\ I_k(\lambda,c,\Conn,\ldots)=I_k^\star,\ \forall k\,\Bigr\},
\]
其中 $I_k^\star$ 为固定常数（预算/守恒值）。
\end{definition}

\begin{theorem}[公理降级：母空间不变性（按步数归纳）]
\label{thm:global-feasible-space-invariant}
考虑一个长度为 $T$ 的离散更新序列 $\{U_t\}_{t=0}^{T-1}$，其中每一步 $U_t$ 作用于 $(\lambda,c,\Conn,\ldots)$。
若对任意 $t$ 与任意 $x$ 有
\[
  I_k\bigl(U_t(x)\bigr)=I_k(x)\qquad(\forall k),
\]
则对任意初值 $x_0\in\Omega_{\mathrm{glob}}$ 有
\[
  x_T:=U_{T-1}\circ\cdots\circ U_0(x_0)\in\Omega_{\mathrm{glob}}.
\]
因此无论后续如何切换扇区或切换算法，所有动作必须落在同一母空间 $\Omega_{\mathrm{glob}}$ 内。
\end{theorem}

\begin{Certificate}[数学归纳法证书：按步数 $T$ 归纳的母空间闭合性]
\label{cert:global-feasible-space-invariant}
定义谓词 $P(T)$ 为：“对任意满足不变量保持的 $T$ 步更新序列，若 $x_0\in\Omega_{\mathrm{glob}}$，则 $x_T\in\Omega_{\mathrm{glob}}$。”
证书化要求为：
\[
  P(1)\ \text{成立},\qquad
  (\forall T\in\NN_{\ge 1})\ \bigl(P(T)\Rightarrow P(T+1)\bigr).
\]
\end{Certificate}

\begin{proof}[证明（数学归纳法证书；谓词因果链）]
由证书~\ref{cert:global-feasible-space-invariant}：
\begin{itemize}
  \item 基步：$T=1$。若 $x_0\in\Omega_{\mathrm{glob}}$，则 $I_k(x_0)=I_k^\star$。由不变量保持得 $I_k(U_0(x_0))=I_k(x_0)=I_k^\star$，故
    $x_1\in\Omega_{\mathrm{glob}}$。
  \item 归纳步：假设 $P(T)$ 成立。对 $T+1$ 步序列，有 $x_T\in\Omega_{\mathrm{glob}}$；
  再由不变量保持得 $I_k(U_T(x_T))=I_k(x_T)=I_k^\star$，故 $x_{T+1}\in\Omega_{\mathrm{glob}}$。
\end{itemize}
由数学归纳法得对任意有限 $T$ 结论成立。证毕。
\end{proof}

\begin{remark}[解释：为什么它是“基础公理”]
母空间 $\Omega_{\mathrm{glob}}$ 的作用是裁剪全局控制—搜索空间：
扇区切换与算法切换只改变“如何在 $\Omega_{\mathrm{glob}}$ 内走路”，而不允许突破 $\Omega_{\mathrm{glob}}$ 的硬约束。
\end{remark}

\subsection{扇区层：规范场扇区三元组与边缘算子切换}
\label{sec:sector-layer-sector-triple}

\begin{definition}[规范场扇区（相空间片区/扇区模型）]
将每个可识别扇区定义为三元组
\[
  S_i := \bigl(\mathfrak{A}_i,\ \mathcal{T}_i,\ \partial_i\bigr),
\]
其中：
\begin{enumerate}
  \item $\mathfrak{A}_i$ 为有效算子/可观测代数（生成元、关系式、交换/非交换规则、约束理想等）；
  \item $\mathcal{T}_i$ 为拓扑约束与允许扇区（允许的缺陷类型、通量/缠绕、边界可通行规则、holonomy 约束等）；
  \item $\partial_i$ 为“可修复/不可修复”分类所依赖的上同调微分（本笔记统一用 $\partial$，以避免与外界实时干扰算子 $\delta_t$ 混淆）。
\end{enumerate}
\end{definition}

\begin{remark}[与三元组模型 $\mathcal{M}_i$ 的对应]
第~\ref{sec:effective-model-triple} 节采用 $\mathcal{M}_i=(\mathfrak{A}_i,\Bundle_i\to X_i,\partial_i)$。
本节的 $\mathcal{T}_i$ 可视为从 $\Bundle_i\to X_i$ 提取的“拓扑扇区与边界允许性”数据，
从而 $S_i$ 是对 $\mathcal{M}_i$ 的扇区化压缩表示，便于直接驱动切换逻辑。
\end{remark}

\begin{definition}[边缘算子调节与扇区切换映射]
令 $\alpha$ 为可控的边缘调节参数（边界自由度、界面耦合、边界补偿规则、边界测量与反馈通道等）。
称映射
\[
  \mathsf{EdgeSwitch}_\alpha:\ S_i\mapsto S_j
\]
为边缘算子切换映射，若其作用至少包含对修复结构的改造：
\[
  \partial_j = \partial_i + \partial_{\mathrm{edge}}(\alpha),
\]
并允许同时改变 $\mathcal{T}$（边界通量/缺陷通行规则）甚至改变有效代数 $\mathfrak{A}$（低能自由度重组）以反映相变/扇区重组。
\end{definition}

\begin{proposition}[命题：扇区切换在母空间约束下进行（不突破守恒底座）]
\label{prop:sector-switching-respects-global-feasible-space}
若边缘算子切换 $\mathsf{EdgeSwitch}_\alpha$ 被建模为闭合总系统内部的结构调整（或把外部库显式并入系统边界自由度），
并且满足定理~\ref{thm:global-feasible-space-invariant} 的不变量保持条件，
则任意扇区切换前后仍有 $(\lambda,c,\Conn,\ldots)\in\Omega_{\mathrm{glob}}$。
因此扇区层的“切换规则”不能被用来绕开公理层的全局预算/守恒约束。
\end{proposition}

\begin{Certificate}[证书：母空间不变性直接推出切换闭合性]
\label{cert:sector-switching-respects-global-feasible-space}
证书登记：将扇区切换视为某一步更新 $U_t$；若该步保持所有 $I_k$，则由定理~\ref{thm:global-feasible-space-invariant} 得结论成立。
\end{Certificate}

\begin{proof}[证明（谓词因果链）]
由证书~\ref{cert:sector-switching-respects-global-feasible-space}：
把 $\mathsf{EdgeSwitch}_\alpha$ 视为更新序列中的一步 $U_t$，若其保持不变量，则由定理~\ref{thm:global-feasible-space-invariant} 得切换前后仍在
  $\Omega_{\mathrm{glob}}$。
命题成立。证毕。
\end{proof}

\begin{remark}[扇区切换与“规范群内切换”的严格含义]
扇区切换并非 A 级意义下的纯规范变换；
在“母规范群框架”下，它更接近于在不同的\textbf{非规范等价}扇区/分支之间切换（通常伴随 $\mathcal{T}$ 与 $\partial$ 的改变，甚至 $\mathfrak{A}$ 的改变），
属于第~\ref{sec:gauge-switching-three-levels} 节的 C 级口径。
\end{remark}

\subsection{算法层：扇区到“同伦扭曲搜索/控制配置栈”的映射}
\label{sec:algorithm-layer-sector-to-stack}

\begin{definition}[扇区算法配置栈（最低配置）]
对每个扇区 $S_i$，定义其算法配置栈为
\[
  \mathcal{A}_i
  :=
  \bigl(
    M_i,\ J_i,\ \mathsf{Constraints}_i,\ \mathsf{Estimator}_i,\ \mathsf{Controller}_i,\ \mathsf{LocalSolver}_i,\
      \mathsf{Verifier}_i,\ T_i
  \bigr),
\]
其中各分量含义如下：
\begin{enumerate}
  \item $M_i$：扇区内有效模型（状态、演化、观测口径）；
  \item $J_i$：目标函数（可含鲁棒/风险项）；
  \item $\mathsf{Constraints}_i$：扇区内约束（含拓扑/边界允许性与资源约束）；
  \item $\mathsf{Estimator}_i$：估计器（由观测得到扇区内状态/参数）；
  \item $\mathsf{Controller}_i$：控制器或滚动优化控制器（相内追踪/稳定/安全）；
  \item $\mathsf{LocalSolver}_i$：局部求解器族（按相域特征动态选型，见第~\ref{sec:twisted-homotopy-real-time-control} 节）；
  \item $\mathsf{Verifier}_i$：验证链/保障链（审计闭环、防伪信号、防不可复现牵引）；
  \item $T_i$：扭曲搬运结构（含扇区内沿路径的搬运 $U_\gamma$，以及跨扇区贴图 $T_{i\to j}$ 的接口，见第~\ref{sec:twisted-search-transport} 节）。
\end{enumerate}
\end{definition}

\begin{definition}[扇区到算法的映射]
定义映射
\[
  \Phi:\ S_i\longmapsto \mathcal{A}_i,
\]
将扇区结构对象映射到其对应的算法配置栈。
该映射的工程意义是：\textbf{扇区（规范场相空间）变了，算法必须切到对应的配置栈}，否则约束语义、估计口径与验证链将不再一致。
\end{definition}

\begin{lemma}[切换一致性：扇区切换必须伴随算法栈切换与贴图搬运]
\label{lem:sector-switch-implies-stack-switch-and-transport}
若在某一步发生扇区切换 $S_i\to S_j$，则必须同时执行：
\begin{enumerate}
  \item 算法栈切换：$\mathcal{A}_i\to\mathcal{A}_j=\Phi(S_j)$；
  \item 贴图/搬运：用 $T_{i\to j}$ 将旧扇区的构造配置与证据链参数映射为新扇区的可行初值。
\end{enumerate}
否则即使母空间约束仍满足，也可能出现“可行性断裂/验证链失配/不可复现牵引”的工程失败模式。
\end{lemma}

\begin{Certificate}[证书：语义一致性链（扇区约束 $\Rightarrow$ 算法口径）]
\label{cert:sector-switch-implies-stack-switch-and-transport}
证书登记如下谓词链：
\[
\begin{aligned}
  A_1 &: S_i\to S_j\Rightarrow (\mathfrak{A},\mathcal{T},\partial)\ \text{口径发生改变},\\
  A_2 &: \Rightarrow \mathsf{Constraints},\ \mathsf{Estimator},\ \mathsf{Verifier}\ \text{的语义随扇区改变},\\
  A_3 &: \Rightarrow \mathcal{A}_i\ \text{不再匹配}\Rightarrow \text{需要}\ \mathcal{A}_j=\Phi(S_j),\\
  A_4 &: \Rightarrow \text{需用}\ T_{i\to j}\ \text{保持热启动可行与证据链一致}.
\end{aligned}
\]
\end{Certificate}

\begin{proof}[证明（谓词因果链）]
由证书~\ref{cert:sector-switch-implies-stack-switch-and-transport}：
扇区切换改变结构口径得 $A_1$；
由结构口径变化推出约束/估计/验证语义变化得 $A_2$；
语义变化导致原算法栈失配并要求切换到 $\Phi(S_j)$ 得 $A_3$；
为保持连续可行与证据链一致需贴图搬运得 $A_4$。引理成立。证毕。
\end{proof}

\subsection{统一闭环：扇区化同伦扭曲搜索—控制的形式系统}
\label{sec:sectorized-twisted-search-control-loop}

\begin{definition}[扇区化闭环状态与两类步]
定义闭环状态为
\[
  \Sigma_t := \bigl(S_{i_t},\ x_t,\ \mathcal{A}_{i_t}\bigr),
  \qquad x_t\in\Omega_{\mathrm{glob}},\ \mathcal{A}_{i_t}=\Phi(S_{i_t}).
\]
称一步更新 $\Sigma_t\to\Sigma_{t+1}$ 为：
\begin{enumerate}
  \item \textbf{相内步}：保持扇区 $S_{i_{t+1}}=S_{i_t}$，在 $\mathcal{A}_{i_t}$ 下做续接/局部搜索/控制更新（含扇区内搬运 $U_\gamma$）；
  \item \textbf{跨扇区步}：选择边缘调节参数 $\alpha_t$ 触发 $S_{i_{t+1}}=\mathsf{EdgeSwitch}_{\alpha_t}(S_{i_t})$，并执行算法栈切换与贴图搬运
    $T_{i_t\to i_{t+1}}$。
\end{enumerate}
\end{definition}

\begin{theorem}[公理降级：扇区化对象—映射—切换闭环的自洽性（按步数归纳）]
\label{thm:sectorized-loop-consistency}
考虑任意有限步数 $T$ 的闭环序列 $\{\Sigma_t\}_{t=0}^{T}$。
若每一步更新满足：
\begin{enumerate}
  \item 母空间不变性：对 $x_t$ 的更新保持 $\mathcal{I}$（定理~\ref{thm:global-feasible-space-invariant}）；
  \item 跨扇区一致性：发生扇区切换时满足引理~\ref{lem:sector-switch-implies-stack-switch-and-transport}；
\end{enumerate}
则对任意初值 $\Sigma_0=(S_{i_0},x_0,\mathcal{A}_{i_0})$（其中 $x_0\in\Omega_{\mathrm{glob}}$ 且 $\mathcal{A}_{i_0}=\Phi(S_{i_0})$），
  对所有 $t\le T$ 有：
\[
  x_t\in\Omega_{\mathrm{glob}},\qquad \mathcal{A}_{i_t}=\Phi(S_{i_t}).
\]
即闭环在公理层约束下保持可行，并在每一步保持“扇区—算法”语义一致。
\end{theorem}

\begin{Certificate}[数学归纳法证书：按步数 $T$ 归纳的一致性保持]
\label{cert:sectorized-loop-consistency}
定义谓词 $P(T)$ 为：
“任意长度为 $T$ 的闭环序列，只要每一步满足母空间不变性与跨扇区一致性，则对所有 $t\le T$ 有 $x_t\in\Omega_{\mathrm{glob}}$ 且 $\mathcal{A}_{i_t}=\Phi(S_{i_t})$。
  ”
证书化要求为：
\[
  P(1)\ \text{成立},\qquad
  (\forall T\in\NN_{\ge 1})\ \bigl(P(T)\Rightarrow P(T+1)\bigr).
\]
\end{Certificate}

\begin{proof}[证明（数学归纳法证书；谓词因果链）]
由证书~\ref{cert:sectorized-loop-consistency}：
\begin{itemize}
  \item 基步：$T=1$。母空间不变性保证 $x_1\in\Omega_{\mathrm{glob}}$；
  若无扇区切换，则 $\mathcal{A}_{i_1}=\mathcal{A}_{i_0}=\Phi(S_{i_0})=\Phi(S_{i_1})$；
  若有扇区切换，则由引理~\ref{lem:sector-switch-implies-stack-switch-and-transport} 得 $\mathcal{A}_{i_1}=\Phi(S_{i_1})$。故 $P(1)$ 成立。

  \item 归纳步：假设 $P(T)$ 成立。对 $T+1$ 步序列，由归纳假设得 $x_T\in\Omega_{\mathrm{glob}}$ 且 $\mathcal{A}_{i_T}=\Phi(S_{i_T})$。
  第 $T$ 步由母空间不变性得 $x_{T+1}\in\Omega_{\mathrm{glob}}$；
  又由“是否扇区切换”的两分与引理~\ref{lem:sector-switch-implies-stack-switch-and-transport} 得 $\mathcal{A}_{i_{T+1}}=\Phi(S_{i_{T+1}})
    $。
  故 $P(T+1)$ 成立。
\end{itemize}
由数学归纳法得对任意有限 $T$ 结论成立。证毕。
\end{proof}

\begin{corollary}[推论：若不发生扇区切换，则退化为相内续接 + 局部算法动态选型]
\label{cor:no-sector-switch-reduces-to-intra-sector}
若闭环序列中不发生任何扇区切换，则 $S_{i_t}=S_{i_0}$（$\forall t$），算法栈保持为 $\mathcal{A}_{i_0}$，
系统退化为在固定扇区内的续接型滚动优化与局部求解器动态选型（第~\ref{sec:twisted-homotopy-real-time-control} 节），无需跨扇区贴图 $T_{i\to j}$。
\end{corollary}

\begin{Certificate}[证书：无切换 $\Rightarrow$ 扇区常值 $\Rightarrow$ 算法栈常值]
\label{cert:no-sector-switch-reduces-to-intra-sector}
证书登记：无扇区切换 $\Rightarrow S_{i_t}$ 常值；由 $\mathcal{A}_{i_t}=\Phi(S_{i_t})$ 得算法栈常值；因此仅剩相内续接与局部选型。
\end{Certificate}

\begin{proof}[证明（谓词因果链）]
由证书~\ref{cert:no-sector-switch-reduces-to-intra-sector}：
扇区常值推出算法栈常值；此时系统只在固定扇区内迭代，结论成立。证毕。
\end{proof}

\begin{remark}[最小切换逻辑（可直接工程落地）]
一个可直接落地的最小触发器可由三类信号组成（各信号均可证书化登记）：
\begin{enumerate}
  \item \textbf{相内平台信号}：局部最优化边际收益持续低于阈值；
  \item \textbf{结构性障碍信号}：出现不可修复类迹象（某类负项始终无法被现有 $\partial_i$ 归入上边界）；
  \item \textbf{边界风险信号}：接近相边界/滞回增强/噪声结构突变/不可逆指标上升。
\end{enumerate}
触发后执行：选择候选 $\alpha$ 完成 $S_i\to S_j$，切换算法栈 $\mathcal{A}_i\to\mathcal{A}_j$，并用 $T_{i\to j}$ 贴图搬运热启动与证据链参数。
\end{remark}

\begin{remark}[两张输入表（把三层框架变成可运行系统）]
要把本章框架推进到可实现系统，最小需要补齐两张表（可逐步迭代完善）：
\begin{enumerate}
  \item \textbf{扇区目录表}：每个 $S_i=(\mathfrak{A}_i,\mathcal{T}_i,\partial_i)$ 的最小描述、可观测判据、允许的边缘调节参数集；
  \item \textbf{算法映射表}：$\Phi(S_i)=\mathcal{A}_i$ 的求解器选型、验证链口径、扇区内搬运 $U_\gamma$ 与跨扇区贴图 $T_{i\to j}$ 的定义与校验条件。
\end{enumerate}
\end{remark}

\section{生命周期规划与逆向设计：扇区切换—算法切换—算子序列的有向路径问题}
\label{sec:lifecycle-planning-inverse-design}

\begin{remark}[核心推论（生命周期 = 受门槛约束的有向路径）]
在公理层母空间 $\Omega_{\mathrm{glob}}$（第~\ref{sec:axiom-layer-global-feasible-space} 节）固定、扇区对象 $S_i$ 与算法映射 $\Phi$（第~\ref{sec:sector-layer-sector-triple} 节与第~\ref{sec:algorithm-layer-sector-to-stack} 节）已给定的前提下，
一个产品/解决方案从概念到交付再到退役的全生命周期，可被抽象为：
\begin{enumerate}
  \item 在\textbf{扩展状态空间}（成熟度等级 + 扇区 + 算法栈 + 证据包 + 格点索引）上的一个\textbf{有向路径}；
  \item 路径的边由\textbf{边缘算子库}及其\textbf{非交换序列}生成（必要时包含算法切换算子与相内调参算子）；
  \item 每一步都受\textbf{门槛（Gate）与成熟度等级（TRL/GRL/\,$\cdots$）}的证据谓词约束。
\end{enumerate}
本章将上述推论固化为可推演、可实现的形式系统：对象、算子、图模型、门槛谓词与逆向（反向）规划接口。
\end{remark}

\subsection{扩展状态：成熟度、扇区、算法栈、证据包与格点索引}
\label{sec:lifecycle-extended-state}

\begin{definition}[成熟度等级与门槛方向]
令 $R\in\NN$ 为目标成熟度上界，定义离散成熟度集合
\[
  \mathcal{R}:=\{0,1,\dots,R\}.
\]
其中 $r=0$ 表示概念/探索起点，$r=R$ 表示目标交付成熟度。
门槛方向记为 $r\to r+1$。
\end{definition}

\begin{remark}[TRL/GRL/Stage-Gate 的统一口径]
成熟度等级可承载不同体系的标签（例如技术成熟度 \textsf{TRL}\cite{ref_nasa_trl}、面向系统集成的 \textsf{GRL} 讨论\cite{ref_energycentral_grl}、
以及治理/路线类成熟度框架\cite{ref_oecd_advanced_materials}，或阶段门槛 \textsf{Stage-Gate}\cite{ref_stage_gate_dtl}）。
本笔记不强制固定某一体系，而是把它们统一为 $r\in\mathcal{R}$ 与门槛谓词 $\mathsf{Gate}(r\to r+1)$ 的可审计接口。
\end{remark}

\begin{definition}[证据包（Evidence Package）]
令 $\mathcal{E}$ 为证据空间（试验/仿真/运行数据、审计记录、可靠性与可重复性证明、制造与运维可行性材料等）。
称 $E\in\mathcal{E}$ 为证据包，若 $E$ 可被解析为一组可审计原子证据及其元数据（来源、时间、口径、置信度与可复验条件）。
\end{definition}

\begin{definition}[格点索引（几何/材料/环境的离散化占位符）]
令 $\mathcal{K}$ 为离散设计格点集合，用于承载几何形态、材料选择与包络环境的离散化索引：
\[
  \kappa\in\mathcal{K}.
\]
$\kappa$ 的工程语义是：把连续高维设计域压缩为“可枚举的格点块/候选族”，以支撑边缘算子的枚举、前置条件检查与自适应加密（见备注~\ref{rem:adaptive-gridding}）。
\end{definition}

\begin{remark}[自适应格点而非均匀格点]
\label{rem:adaptive-gridding}
对高维几何/材料/环境空间做全局均匀格点会快速导致维数灾难。
更可执行的口径是：粗网格覆盖全域，并在“灵敏度高/接近扇区边界/证据不确定性高”的区域自适应加密；
这与两层搜索（相内 $\to$ 相间 $\to$ 再相内；第~\ref{sec:twisted-homotopy-two-level-search} 节）一致。
\end{remark}

\begin{definition}[生命周期节点（扩展组合状态）]
定义生命周期节点为
\[
  v := (r,\ S_i,\ x,\ \mathcal{A}_i,\ \kappa,\ E),
\]
其中：
\begin{enumerate}
  \item $r\in\mathcal{R}$ 为成熟度；
  \item $S_i$ 为扇区（第~\ref{sec:sector-layer-sector-triple} 节）；
  \item $x\in\Omega_{\mathrm{glob}}$ 为母空间内的连续状态（第~\ref{sec:axiom-layer-global-feasible-space} 节）；
  \item $\mathcal{A}_i=\Phi(S_i)$ 为扇区算法栈（第~\ref{sec:algorithm-layer-sector-to-stack} 节）；
  \item $\kappa\in\mathcal{K}$ 为离散化索引；
  \item $E\in\mathcal{E}$ 为证据包。
\end{enumerate}
\end{definition}

\subsection{算子库：边缘算子、算法切换与相内调参（非交换生成元）}
\label{sec:lifecycle-operator-library}

\begin{definition}[三类算子库（最小闭环覆盖）]
为使生命周期规划能闭环运行，引入三类算子库：
\[
  \mathcal{G}_{\mathrm{edge}},\qquad
  \mathcal{G}_{\mathrm{alg}},\qquad
  \mathcal{G}_{\mathrm{intra}}.
\]
其语义分别为：
\begin{enumerate}
  \item $\mathcal{G}_{\mathrm{edge}}$：边缘算子（触发扇区切换、改变拓扑扇区与修复结构；见第~\ref{sec:sector-layer-sector-triple} 节的
    $\mathsf{EdgeSwitch}_\alpha$ 口径）；
  \item $\mathcal{G}_{\mathrm{alg}}$：算法切换算子（在同一扇区内更换/重配算法栈组件，例如估计器、局部求解器族、验证链策略）；
  \item $\mathcal{G}_{\mathrm{intra}}$：相内调参算子（在固定扇区与算法栈下做连续或离散的小步调整，用于局部搜索与在线控制）。
\end{enumerate}
\end{definition}

\begin{definition}[算子签名（前置/后置条件与证据钩子）]
对任意算子 $g$（可属于上述任一库），登记其签名为四元组
\[
  \mathsf{Sig}(g):=(\mathsf{Pre}(g),\ \mathsf{Post}(g),\ \mathsf{Cost}(g),\ \mathsf{Verify}(g)).
\]
其中：
\begin{enumerate}
  \item $\mathsf{Pre}(g)$ 为前置条件（适用扇区/适用格点域/先决证据/资源预算）；
  \item $\mathsf{Post}(g)$ 为后置条件（可能进入的扇区集合、母空间不变量保持声明、证据增量与不可逆标记）；
  \item $\mathsf{Cost}(g)$ 为成本/周期/风险等代价向量；
  \item $\mathsf{Verify}(g)$ 为验证链钩子（执行该算子必须触发的证据采集、审计与复验动作）。
\end{enumerate}
\end{definition}

\begin{definition}[非交换算子序列（构造词）]
称有限序列 $w=g_{T-1}\circ\cdots\circ g_0$ 为算子词（构造词），其中每个 $g_t$ 来自 $\mathcal{G}_{\mathrm{edge}}\cup\mathcal{G}_{\mathrm{alg}}
  \cup\mathcal{G}_{\mathrm{intra}}$。
一般情形下满足非交换性：
\[
  g_a\circ g_b \neq g_b\circ g_a,
\]
因此生命周期必须把\textbf{顺序}作为一等对象进行证书化登记与搜索（与第~\ref{sec:ptopb-etopb-operator-decomposition} 节的非交换构造口径一致）。
\end{definition}

\subsection{门槛谓词：成熟度提升的证据化条件}
\label{sec:lifecycle-gate-policy}

\begin{definition}[门槛谓词（Gate Policy）]
称谓词
\[
  \mathsf{Gate}_{r\to r+1}:\ \mathcal{E}\times \Omega_{\mathrm{glob}}\times \{S_i\}\to\{\text{True},\text{False}\}
\]
为门槛谓词，若它用证据包 $E$（以及与之对应的口径信息）判定在给定扇区与母空间状态下是否允许从 $r$ 进入 $r+1$：
\[
  \mathsf{Gate}_{r\to r+1}(E,x,S_i)=\text{True}\ \Rightarrow\ r\ \text{可提升}.
\]
\end{definition}

\begin{remark}[门槛的工程含义]
门槛不只评价“性能”，还应包含可制造、可运维、可验证、可扩展与治理合规等证据维度；
阶段门槛方法（\textsf{Stage-Gate}\cite{ref_stage_gate_dtl}）在工程管理上正是以“证据驱动的 Go/Kill 决策点”为核心思想之一。
\end{remark}

\subsection{有向图模型：生命周期 = 受门槛约束的路径可行性与最优性}
\label{sec:lifecycle-directed-graph}

\begin{definition}[生命周期有向图]
定义生命周期有向图为 $\mathcal{G}=(V,E)$，其中：
\begin{enumerate}
  \item 节点集 $V$ 为所有满足 $x\in\Omega_{\mathrm{glob}}$ 且 $\mathcal{A}_i=\Phi(S_i)$ 的生命周期节点 $v=(r,S_i,x,\mathcal{A}_i,\kappa,
    E)$；
  \item 有向边 $v\to v'$ 属于 $E$ 当且仅当存在算子 $g$ 使得：
  \begin{enumerate}
    \item $\mathsf{Pre}(g)$ 在 $v$ 上满足；
    \item 施加 $g$ 得到新节点 $v'$（允许扇区变化、算法栈切换与证据包更新）；
    \item 若 $r' = r+1$，则满足门槛谓词 $\mathsf{Gate}_{r\to r+1}(E',x',S_{i'})=\text{True}$；若 $r'=r$，则不要求通过门槛。
  \end{enumerate}
\end{enumerate}
\end{definition}

\begin{theorem}[公理降级：任意有限生命周期可表示为有向图上的一条路径（按步数归纳）]
\label{thm:lifecycle-is-a-directed-path}
给定初始节点 $v_0$ 与目标成熟度 $R$。
若存在一段有限生命周期执行记录（由算子序列与证据更新组成）使成熟度从 $r_0$ 推进到 $R$，则存在有向路径
\[
  v_0\to v_1\to\cdots\to v_T
\]
满足 $v_T$ 的成熟度分量为 $R$，且每条边对应一条可审计的算子作用并满足相应门槛谓词。
\end{theorem}

\begin{Certificate}[数学归纳法证书：按生命周期步数 $T$ 归纳的路径表示]
\label{cert:lifecycle-is-a-directed-path}
定义谓词 $P(T)$ 为：
“任意长度为 $T$ 的生命周期执行记录，都可被嵌入为有向图 $\mathcal{G}$ 上从 $v_0$ 出发的长度为 $T$ 的路径，并保持门槛谓词与前置条件的一致性。”
证书化要求为：
\[
  P(1)\ \text{成立},\qquad
  (\forall T\in\NN_{\ge 1})\ \bigl(P(T)\Rightarrow P(T+1)\bigr).
\]
\end{Certificate}

\begin{proof}[证明（数学归纳法证书；谓词因果链）]
由证书~\ref{cert:lifecycle-is-a-directed-path}：
\begin{itemize}
  \item 基步：$T=1$。执行记录只含一步算子 $g_0$ 与证据更新。由边的定义，存在边 $v_0\to v_1$；若成熟度提升则同时满足对应门槛谓词。故 $P(1)$ 成立。
  \item 归纳步：假设 $P(T)$ 成立。对长度为 $T+1$ 的执行记录，其前 $T$ 步由归纳假设对应到路径 $v_0\to\cdots\to v_T$。
  第 $T$ 步算子 $g_T$ 在 $v_T$ 上满足前置条件并产生 $v_{T+1}$；若发生成熟度提升则满足门槛谓词。故 $v_T\to v_{T+1}$ 是有向边，从而得到长度为 $T+1$ 的路径。即 $P(T+1)$ 成立。
\end{itemize}
由数学归纳法得对任意有限步数 $T$ 结论成立。证毕。
\end{proof}

\begin{remark}[最优路径问题（多目标口径）]
当为每条边附加代价向量 $\mathsf{Cost}(g)$（成本/周期/风险/资源消耗等）后，
生命周期规划可被表达为“在门槛约束下的最短路/最优路径”问题（可为多目标帕累托或加权和）。
工程上常用做法是：把风险与不可逆性作为硬约束或强惩罚项，以避免在证据不足时走入不可回退路径。
\end{remark}

\subsection{非交换序列压缩：重写等价与偏序约束}
\label{sec:lifecycle-rewriting-and-partial-order}

\begin{definition}[重写规则与策略类]
令 $\Rightarrow$ 为算子词上的重写关系，表示在给定容差与审计口径下允许的序列变形（例如某些算子对在特定扇区与格点域内近似可交换）。
称两个算子词 $w_1,w_2$ 属于同一策略类，若存在有限重写链
\[
  w_1 \Rightarrow \cdots \Rightarrow w_2.
\]
\end{definition}

\begin{proposition}[命题：重写等价可用于压缩非交换序列搜索空间（条件命题）]
\label{prop:rewriting-reduces-search-space}
若重写关系 $\Rightarrow$ 满足：对任意起点节点 $v$，$w_1\Rightarrow^\ast w_2$ 蕴含
\[
  w_1(v)\ \text{与}\ w_2(v)\ \text{落入同一扇区与同一门槛口径下（并满足等价证据判据）},
\]
则可用策略类代替具体序列进行搜索，从而把“全排列枚举”压缩为“等价类枚举”。
\end{proposition}

\begin{Certificate}[证书：等价类枚举的正确性前提]
\label{cert:rewriting-reduces-search-space}
证书登记：若重写保持扇区/门槛/证据判据不变，则同一策略类内任一代表词均可替代；因此搜索可在等价类上进行而不丢失可行解。
\end{Certificate}

\begin{proof}[证明（谓词因果链）]
由证书~\ref{cert:rewriting-reduces-search-space}：
重写保持判据不变 $\Rightarrow$ 同类代表词可互换 $\Rightarrow$ 搜索可在等价类上进行且不丢失可行路径。
命题成立。证毕。
\end{proof}

\begin{definition}[偏序先决条件（Partial Order）]
对算子集合引入先决条件偏序 $\prec$：若 $g_a\prec g_b$，表示 $g_b$ 的可行性依赖 $g_a$ 先行满足的证据或结构条件。
则允许的算子序列必须是该偏序的线性扩展，从而把“全排列”约束为“受约束调度”。
\end{definition}

\begin{remark}[工程含义：避免组合爆炸]
重写等价（可交换近似/可重排关系）与偏序约束（先决条件）是抑制非交换序列组合爆炸的两类关键压缩机制。
它们本质上是对算子库的“可审计元数据建模”，而不是对优化器做玄学调参。
\end{remark}

\subsection{逆向设计：反向规划与双向搜索接口}
\label{sec:lifecycle-backward-planning}

\begin{definition}[前驱生成器（逆向设计接口）]
对任意节点 $v$，称映射 $\mathsf{Pred}(v)$ 为前驱生成器，若它返回所有可能的前驱节点集合：
\[
  \mathsf{Pred}(v) := \{\,u\in V\mid (u\to v)\in E\,\}.
\]
\end{definition}

\begin{lemma}[引理：存在路径 $\Leftrightarrow$ 正向可达与反向可达相交]
\label{lem:bidirectional-search-correctness}
令 $\mathsf{Reach}_F(v_0)$ 为从起点 $v_0$ 经边关系正向可达的节点集，$\mathsf{Reach}_B(v_{\mathrm{goal}})$ 为从目标集合（例如成熟度为 $R$ 的节点集）按
  $\mathsf{Pred}$ 反向可达的节点集。
若 $\mathsf{Pred}$ 正确刻画边关系，则存在从 $v_0$ 到某个目标节点的路径，当且仅当
\[
  \mathsf{Reach}_F(v_0)\cap \mathsf{Reach}_B(v_{\mathrm{goal}})\neq\varnothing.
\]
\end{lemma}

\begin{Certificate}[证书：边关系的双向可达性等价]
\label{cert:bidirectional-search-correctness}
证书登记：正向可达定义给出“路径 $\Rightarrow$ 交集非空”；反向前驱定义给出“交集非空 $\Rightarrow$ 拼接出路径”。
\end{Certificate}

\begin{proof}[证明（谓词因果链）]
由证书~\ref{cert:bidirectional-search-correctness}：
若存在路径，则其中任一点同时是正向可达与反向可达，故交集非空；
若交集非空，取交点并拼接正向路径与反向路径（取反向路径的逆序），得从起点到目标的路径。
引理成立。证毕。
\end{proof}

\begin{remark}[不可逆边与“近逆”]
生命周期中的边缘算子通常不可逆；逆向设计不要求真实逆算子存在，而是要求可生成满足前置条件的候选前驱集合（$\mathsf{Pred}$）。
这与材料/工艺/组织治理等不可逆动作的工程现实一致。
\end{remark}

\subsection{工程交付件：目录表、算子库、映射表与门槛矩阵}
\label{sec:lifecycle-deliverables}

\begin{remark}[四个最低交付件（保证“能跑起来”）]
要将“生命周期 = 扇区化有向路径”落成可运行系统，建议最少形成四个交付件：
\begin{enumerate}
  \item \textbf{扇区目录（Sector Catalog）}：每个 $S_i$ 的边界判据、可观测判据、适用格点域、典型风险；
  \item \textbf{边缘算子库（Edge Operator Library）}：$\mathcal{G}_{\mathrm{edge}}$ 的算子签名 $\mathsf{Sig}(g)$，含依赖/强非交换对、验证链钩子；
  \item \textbf{算法映射表（Sector $\to$ Stack）}：$\Phi(S_i)=\mathcal{A}_i$ 的具体化（估计/求解/控制/验证/搬运）；
  \item \textbf{门槛矩阵（Gate Policy）}：$\mathsf{Gate}_{r\to r+1}$ 的证据谓词清单，对齐 \textsf{TRL}\cite{ref_nasa_trl} 或
    \textsf{Stage-Gate}\cite{ref_stage_gate_dtl} 等成熟度体系。
\end{enumerate}
这四件缺一不可：缺扇区目录则无法判边界；缺算子库则无法生成边；缺映射表则无法保证语义一致；缺门槛矩阵则无法推进成熟度。
\end{remark}

\section{OpenRA 即时战略：子丛解耦与非因果组织基准的可执行抽象}
\label{sec:openra-subbundle-noncausal-basis}

\begin{remark}[定位（用沙盘验证“扇区/边缘算子/非交换序列”）]
\textsf{OpenRA} 这类即时战略游戏非常适合作为“算子丛—扇区—边缘算子—非交换序列”的可执行沙盘：
\begin{enumerate}
  \item 系统状态可观测但信息不完全（战争迷雾），扰动强且实时（对手行为）；
  \item 算子库离散且明确（造建筑/出兵/移动/交战/侦察），并且强非交换（顺序/时机/位置决定结果）；
  \item 既能做相内局部优化（相同态势内调参），也能触发“边缘算子”切换态势扇区（扩张、切科技线、前推防线等）。
\end{enumerate}
本章聚焦你提出的关键句：
\textbf{“对子丛解耦等价于算子空间搜索的非因果组织基准”}——给出其形式化定义、为何有用、以及在 \textsf{OpenRA} 中如何落地为可执行规划器。
\end{remark}

\subsection{对象与算子：状态、扰动与非交换构造}
\label{sec:openra-state-and-operators}

\begin{definition}[OpenRA 状态的三分：布局—能力—资源/节奏]
将我方状态抽象为三类信息的组合：
\[
  s := (b,k,\rho),
\]
其中：
\begin{enumerate}
  \item $b$：基底布局（建筑位置、通道/瓶颈占用、扩张点与采矿线）；
  \item $k$：能力结构（科技阶段、生产能力、供电能力、支援设施能力）；
  \item $\rho$：资源与节奏（资金流、电力盈余/赤字、建造队列占用、部队数量与机动性、地图控制）。
\end{enumerate}
\end{definition}

\begin{definition}[实时干扰算子与我方构造算子（OpenRA 口径）]
沿用第~\ref{sec:ptopb-etopb-operator-decomposition} 节的阶段更新口径，
将对手攻击/侦察不完全/地形路径约束统一归并为外界实时干扰算子 $\delta_t$；
将我方建造/生产/机动/交战/侦察等操作归并为系统构造算子 $\pi_t$，
并写作
\[
  s_{t+1} = (\pi_t\circ \delta_t)(s_t).
\]
\end{definition}

\begin{remark}[强非交换性（顺序/时机/位置决定结果）]
在 \textsf{OpenRA} 中，同一组建筑与单位集合只要建造顺序、时机或位置不同，就可能产生完全不同的资源曲线、防御覆盖与作战窗口，
因此构造算子一般满足非交换性：$\pi_a\circ\pi_b\neq \pi_b\circ\pi_a$。
这使得“序列”必须作为一等对象进入审计与搜索（与本笔记对非交换构造的总口径一致）。
\end{remark}

\subsection{子丛解耦：功能域分解与接口变量}
\label{sec:openra-subbundle-decoupling}

\begin{definition}[功能域子丛（子算子丛）]
把我方构造算子族按功能域拆为若干子丛：
\[
  \Pi
  =
  \Pi_{\mathrm{eco}}
  \uplus
  \Pi_{\mathrm{pow}}
  \uplus
  \Pi_{\mathrm{prod}}
  \uplus
  \Pi_{\mathrm{def}}
  \uplus
  \Pi_{\mathrm{mob}}
  \uplus
  \Pi_{\mathrm{info}},
\]
分别对应经济、供电、生产/科技、防御、机动打击、信息/侦察等算子子库。
\end{definition}

\begin{definition}[接口变量（子丛耦合的最小载体）]
称向量 $y$ 为接口变量，若它由少量可审计字段组成，足以刻画子丛之间的主要耦合：
\[
  y := (\text{资金预算},\ \text{电力预算},\ \text{建造队列时间},\ \text{空间/瓶颈占用},\ \text{侦察可信度}).
\]
\end{definition}

\begin{theorem}[公理降级：子丛解耦的接口闭合分解（按步数归纳）]
\label{thm:openra-subbundle-interface-decomposition}
考虑一段长度为 $T$ 的构造序列 $\pi_{0:T-1}$。若每一步 $\pi_t$ 都属于某一子丛 $\Pi_\bullet$，
则存在一条接口变量轨迹 $y_{0:T}$ 与一组子丛内序列 $\pi^{(\bullet)}$，
使得全局序列可被表示为“子丛序列 + 接口一致性约束”的组合：
\begin{enumerate}
  \item $y_{t+1}$ 由 $(y_t,\pi_t,\delta_t)$ 按审计规则更新；
  \item 每个子丛内序列只在其作用域内更新局部对象，并通过 $y$ 与其他子丛耦合。
\end{enumerate}
因此“对子丛解耦”不是声称子丛完全独立，而是把耦合压缩到少量接口变量轨迹上。
\end{theorem}

\begin{Certificate}[数学归纳法证书：按步数 $T$ 归纳的接口闭合性]
\label{cert:openra-subbundle-interface-decomposition}
定义谓词 $P(T)$ 为：
“任意长度为 $T$ 的构造序列，都可被表示为子丛内序列与接口变量轨迹 $y_{0:T}$ 的组合，并且 $y$ 的更新只依赖于当前步的 $(y_t,\pi_t,\delta_t)$。”
证书化要求为：
\[
  P(1)\ \text{成立},\qquad
  (\forall T\in\NN_{\ge 1})\ \bigl(P(T)\Rightarrow P(T+1)\bigr).
\]
\end{Certificate}

\begin{proof}[证明（数学归纳法证书；谓词因果链）]
由证书~\ref{cert:openra-subbundle-interface-decomposition}：
\begin{itemize}
  \item 基步：$T=1$。$\pi_0$ 属于某一子丛，按定义可登记 $y_0$ 并由 $(y_0,\pi_0,\delta_0)$ 更新得到 $y_1$，从而 $P(1)$ 成立。
  \item 归纳步：假设 $P(T)$ 成立。对长度为 $T+1$ 的序列，前 $T$ 步给出 $y_{0:T}$ 与子丛序列登记。
  第 $T$ 步 $\pi_T$ 仍属于某一子丛，按同一规则由 $(y_T,\pi_T,\delta_T)$ 更新得到 $y_{T+1}$，并把该步追加到对应子丛序列中，从而仍得到“子丛序列 + 接口轨迹”的表示。故 $P(T+1)$
    成立。
\end{itemize}
由数学归纳法得结论成立。证毕。
\end{proof}

\begin{proposition}[命题：给定接口轨迹，子丛可并行做局部搜索并可拼接为全局可行解（条件命题）]
\label{prop:openra-subbundle-parallel-local-search}
若接口变量轨迹 $y_{0:T}$ 已给定，并且每个子丛都有一个局部可行更新器
\[
  \mathsf{LS}_{\bullet}(\,\cdot\,;y_{0:T}):\Pi_{\bullet}^{T}\to \Pi_{\bullet}^{T},
\]
使其更新不破坏接口一致性（预算/队列/空间/信息约束不被破坏），
则可对各子丛并行优化并将结果拼接为一条全局可行的构造序列（仍落在同一 $y_{0:T}$ 约束下）。
\end{proposition}

\begin{Certificate}[证书：接口一致性 $\Rightarrow$ 可拼接]
\label{cert:openra-subbundle-parallel-local-search}
证书登记如下谓词链：
\[
\begin{aligned}
  S_1 &: \forall \bullet,\ \mathsf{LS}_{\bullet}\ \text{更新不破坏}\ y_{0:T},\\
  S_2 &: \Rightarrow\ \text{各子丛更新后的序列共享同一接口轨迹},\\
  S_3 &: \Rightarrow\ \text{可按时间步拼接为全局序列且保持可行}.
\end{aligned}
\]
\end{Certificate}

\begin{proof}[证明（谓词因果链）]
由证书~\ref{cert:openra-subbundle-parallel-local-search}：
若每个子丛更新都保持同一接口轨迹，则子丛之间的耦合条件已被 $y_{0:T}$ 同步对齐；
因此把各子丛在每个时间步的动作拼接即可得到全局序列，且不会违反预算/队列/空间/信息等接口约束。
命题成立。证毕。
\end{proof}

\subsection{非因果组织基准：从模块目标到序列化实现}
\label{sec:openra-noncausal-organization-basis}

\begin{definition}[非因果组织基准（结构坐标而非动作时间线）]
称“子丛解耦”提供了算子空间搜索的\textbf{非因果组织基准}，若搜索不以“下一步动作”为主坐标，
而以“能力模块目标 + 接口约束”为主坐标：先在结构空间中确定目标轮廓
\[
  g := (g_{\mathrm{eco}},g_{\mathrm{pow}},g_{\mathrm{prod}},g_{\mathrm{def}},g_{\mathrm{mob}},g_{\mathrm{info}}),
\]
再在满足 $y_{0:T}$ 的条件下，将 $g$ 序列化为可执行的建造/生产/机动/交战动作序列。
\end{definition}

\begin{lemma}[引理：任意动作序列诱导一个结构轮廓；固定结构轮廓则搜索退化为受约束排程]
\label{lem:openra-structure-profile-and-scheduling}
给定一条可执行的构造序列 $\pi_{0:T-1}$，总存在一个结构轮廓 $g$ 使该序列在每个子丛上达到（或近似达到）$g$ 所规定的模块水平。
反之，在固定 $g$ 与接口轨迹 $y_{0:T}$ 的条件下，寻找满足 $g$ 的构造序列可被表述为一个受约束排程问题（队列/依赖/预算/空间/信息）。
\end{lemma}

\begin{Certificate}[证书：投影到子丛 $\Rightarrow$ 结构轮廓；固定轮廓 $\Rightarrow$ 排程]
\label{cert:openra-structure-profile-and-scheduling}
证书登记如下谓词链：
\[
\begin{aligned}
  N_1 &: \pi_{0:T-1}\ \text{投影到各子丛}\Rightarrow \text{得到子丛内达成水平},\\
  N_2 &: \Rightarrow\ \exists g\ \text{使}\ \pi_{0:T-1}\ \text{诱导}\ g,\\
  N_3 &: (g,y_{0:T})\ \text{固定}\Rightarrow \text{约束只剩队列/依赖/预算/空间/信息},\\
  N_4 &: \Rightarrow\ \text{求序列}=\text{受约束排程问题}.
\end{aligned}
\]
\end{Certificate}

\begin{proof}[证明（谓词因果链）]
由证书~\ref{cert:openra-structure-profile-and-scheduling}：
把序列按子丛投影即可得到各功能域在 $T$ 步内的达成水平，从而存在结构轮廓 $g$（可取为这些水平的汇总或其容差包络），得 $N_1\Rightarrow N_2$。
当固定 $g$ 与接口轨迹 $y_{0:T}$ 后，剩余自由度主要是把动作放入队列并满足依赖与预算等约束，故 $N_3\Rightarrow N_4$。
引理成立。证毕。
\end{proof}

\begin{corollary}[推论：结构化搜索（非因果）可压缩动作序列搜索并提升鲁棒性（工程推论）]
\label{cor:openra-noncausal-search-advantage}
在满足接口闭合分解（定理~\ref{thm:openra-subbundle-interface-decomposition}）的前提下，
采用“先结构轮廓 $g$、再排程序列”的非因果组织基准，可以：
\begin{enumerate}
  \item 把指数爆炸的动作时间线搜索压缩为“结构空间搜索 + 受约束排程”两阶段；
  \item 在对手扰动变化时，通过调整少数模块目标与接口阈值实现鲁棒策略切换，而非从头枚举动作序列。
\end{enumerate}
\end{corollary}

\begin{Certificate}[证书：由引理与接口压缩推出两阶段优势]
\label{cert:openra-noncausal-search-advantage}
证书登记如下谓词链：
\[
\begin{aligned}
  P_1 &: \text{定理~\ref{thm:openra-subbundle-interface-decomposition}}\Rightarrow \text{耦合压缩到}\ y_{0:T},\\
  P_2 &: \text{引理~\ref{lem:openra-structure-profile-and-scheduling}}\Rightarrow \text{序列搜索}\Rightarrow (g\text{搜索}
    +\text{排程}),\\
  P_3 &: \Rightarrow\ \text{两阶段结构化搜索压缩动作空间并利于鲁棒切换}.
\end{aligned}
\]
\end{Certificate}

\begin{proof}[证明（谓词因果链）]
由证书~\ref{cert:openra-noncausal-search-advantage}：
接口闭合分解把跨域耦合压缩到 $y_{0:T}$；结构轮廓引理把序列生成退化为排程。
因此搜索可先在结构空间中确定 $g$，再在接口约束下排程生成序列；当扰动变化时调整 $g$ 与阈值即可快速切换。
推论成立。证毕。
\end{proof}

\subsection{扇区切换与落地：OpenRA 的四扇区目录与边缘算子候选}
\label{sec:openra-sector-switching-and-deployment}

\begin{definition}[OpenRA 态势扇区标签（避免与抽象扇区三元组混淆）]
为避免与第~\ref{sec:sector-layer-sector-triple} 节的抽象扇区三元组 $S_i=(\mathfrak{A}_i,\mathcal{T}_i,\partial_i)$ 混淆，
本节用 $\Sigma$ 表示 \textsf{OpenRA} 的态势扇区标签，典型取四扇区：
\[
  \Sigma\in\{\Sigma_{\mathrm{A}},\Sigma_{\mathrm{B}},\Sigma_{\mathrm{C}},\Sigma_{\mathrm{D}}\},
\]
分别对应开局经济扩张、被压制/被骚扰、中期对抗、后期决胜。
\end{definition}

\begin{definition}[边缘算子候选（扇区切换生成元）]
称算子 $g$ 为 \textsf{OpenRA} 的边缘算子，若它能显著改变约束结构或有效自由度组织方式，从而触发态势扇区切换：
\[
  \Sigma \xrightarrow{g} \Sigma'.
\]
典型候选包括：扩张/转基地、切科技线/切兵种线、前推防线/封锁瓶颈、拆对手电力/矿线/科技等。
\end{definition}

\begin{proposition}[命题：扇区切换要求算法栈切换（验证链与阈值随态势改变）]
\label{prop:openra-sector-switch-requires-algorithm-switch}
若扇区从 $\Sigma$ 切换到 $\Sigma'$，则对应的规划/控制算法栈（风险阈值、局部求解器、验证链）也必须切换；
否则会出现“过度贪心导致崩盘”或“过度保守导致错失窗口”等失配风险。
\end{proposition}

\begin{Certificate}[证书：态势改变 $\Rightarrow$ 约束与风险口径改变 $\Rightarrow$ 栈需切换]
\label{cert:openra-sector-switch-requires-algorithm-switch}
证书登记如下谓词链：
\[
\begin{aligned}
  Q_1 &: \Sigma\to\Sigma'\Rightarrow \text{资源/时间/威胁结构改变},\\
  Q_2 &: \Rightarrow \text{风险阈值与验证链优先级改变},\\
  Q_3 &: \Rightarrow \text{局部求解器/控制器需切换以保持可行与鲁棒}.
\end{aligned}
\]
\end{Certificate}

\begin{proof}[证明（谓词因果链）]
由证书~\ref{cert:openra-sector-switch-requires-algorithm-switch}：
态势扇区切换意味着威胁结构与资源窗口发生改变，从而风险阈值、验证链与局部求解器的适用条件也随之改变；
若不切换算法栈，则将以错误的约束口径做决策，导致系统性失配风险。
命题成立。证毕。
\end{proof}

\begin{remark}[可执行落地（最小三层规划器规格）]
将上述结构落成可执行规划器，可按三层实现：
\begin{enumerate}
  \item \textbf{扇区判别器}：由侦察与风险信号判别 $\Sigma$（骚扰强度、矿线受损、电力赤字、兵种克制等）；
  \item \textbf{子丛目标器}：在给定 $\Sigma$ 下设定结构轮廓 $g$ 与接口阈值 $y$（先结构后排程）；
  \item \textbf{序列生成器}：在 $\mathcal{G}_{\mathrm{edge}}\cup\mathcal{G}_{\mathrm{intra}}$ 上生成候选非交换序列，并用验证链过滤不可复验/高风险序列。
\end{enumerate}
该实现把“动作时间线盲搜”替换为“结构坐标搜索 + 受约束排程”，并允许在 $\Sigma$ 变化时快速切换算法栈与结构目标。
\end{remark}

\subsection{开放系统扇区切换守卫条件：对手信念—资源边界—布局脆弱性}
\label{sec:openra-guard-three-criteria}

\begin{remark}[核心点（扇区切换不能只依赖我方内部状态）]
一旦把 \textsf{OpenRA}（或一般开放系统博弈）纳入“扇区—算子—算法栈”框架，
扇区切换的触发条件（守卫条件，guard）就不可能只由我方内部状态决定：
至少需要同时引入三类要件作为守卫条件的组成部分：
\begin{enumerate}
  \item \textbf{对手策略态势要件}：在战争迷雾下通常不是“确定事实”，而是对手策略类型的\textbf{信念/概率分布}；
  \item \textbf{资源匮乏要件}：资金流/电力/队列占用/损耗补充能力等构成可行性边界；
  \item \textbf{基底布局脆弱性要件}：尤其是“防御核心攻击不利”的几何拓扑条件（路径短、瓶颈缺失、关键建筑集中、矿线暴露、覆盖缺口等）。
\end{enumerate}
其工程含义是：扇区切换属于开放系统的混合系统模式切换，必须以“外生态势 + 内生可行性 + 结构脆弱性”共同守卫。
\end{remark}

\begin{definition}[对手策略类型与信念（belief）]
令 $\Theta$ 为对手策略类型集合（例如快空、快车、双矿贪经济、科技直冲、偷矿骚扰等）。
记 $\theta_t\in\Theta$ 为时刻 $t$ 的对手策略类型（潜变量），$\mathcal{Y}_{0:t}$ 为截至 $t$ 的情报证据（侦察所见、交战信息、时间点特征等）。
定义信念谓词为：
\[
  \mathsf{Opp}_{\Sigma\to\Sigma'}(t)
  :=
  \Pr\bigl(\theta_t\in\Theta_{\Sigma\to\Sigma'}\mid \mathcal{Y}_{0:t}\bigr)\ge \tau_{\theta,\Sigma\to\Sigma'},
\]
其中 $\Theta_{\Sigma\to\Sigma'}\subseteq\Theta$ 为触发该切换所关注的策略子集，$\tau_{\theta,\Sigma\to\Sigma'}\in[0,1]$ 为阈值（可做风险敏感调节）。
\end{definition}

\begin{definition}[资源状态向量与可行性谓词]
令 $\mathbf{R}_t$ 为资源状态向量（资金、电力、建造队列占用、产能、补给线安全、损耗率等）；
定义资源谓词为：
\[
  \mathsf{Res}_{\Sigma\to\Sigma'}(t) := \mathbf{R}_t\in\mathcal{U}_{\Sigma\to\Sigma'},
\]
其中 $\mathcal{U}_{\Sigma\to\Sigma'}$ 表示“允许执行该切换及其最小修复序列”的资源可行域（可显式区分硬不可行/软不可行/机会成本过高三层口径）。
\end{definition}

\begin{definition}[布局脆弱性向量与脆弱性谓词]
令 $\mathbf{V}_t$ 为布局脆弱性向量，至少包含三类可量化字段：
\[
  \mathbf{V}_t := (\Delta t_{\mathrm{core}},\ \mathrm{Gap}_{\mathrm{AA}},\ \mathrm{Gap}_{\mathrm{AG}},\ \mathrm{SPOF},
    \ldots),
\]
分别表示到达时间差、反空覆盖缺口、反地覆盖缺口、单点失效风险（关键设施集中度）等。
定义脆弱性谓词为：
\[
  \mathsf{Vul}_{\Sigma\to\Sigma'}(t) := \mathbf{V}_t\in\mathcal{W}_{\Sigma\to\Sigma'},
\]
其中 $\mathcal{W}_{\Sigma\to\Sigma'}$ 表示触发该切换的脆弱性区域（例如“暴露/缺口/时间差超阈值”的集合）。
\end{definition}

\begin{definition}[三要件守卫条件（guard）]
定义从态势扇区 $\Sigma$ 到 $\Sigma'$ 的三要件守卫条件为：
\[
  G_{\Sigma\to\Sigma'}(t)
  :=
  \mathsf{Opp}_{\Sigma\to\Sigma'}(t)
  \ \wedge\
  \mathsf{Res}_{\Sigma\to\Sigma'}(t)
  \ \wedge\
  \mathsf{Vul}_{\Sigma\to\Sigma'}(t).
\]
\end{definition}

\begin{theorem}[公理降级：三要件守卫条件给出开放系统模式切换的可审计口径]
\label{thm:openra-guard-three-criteria}
在开放系统博弈口径下，若一次扇区切换 $\Sigma\to\Sigma'$ 被定义为“响应外生威胁并保持内生可行”，
则采用三要件守卫条件 $G_{\Sigma\to\Sigma'}$ 作为触发标准，可以把切换判据分解为：
\textbf{对手信念（外生）}、\textbf{资源可行（内生）}、\textbf{结构脆弱（几何/拓扑）}三类可审计字段的合取，
从而避免“仅凭内部状态盲切换”导致的系统性失配。
\end{theorem}

\begin{Certificate}[证书：三要件合取作为守卫条件的审计链]
\label{cert:openra-guard-three-criteria}
证书登记如下谓词链：
\[
\begin{aligned}
  G_1 &: \mathsf{Opp}\Rightarrow \text{对手威胁在证据口径下成立（信念阈值）},\\
  G_2 &: \mathsf{Res}\Rightarrow \text{切换与最小修复序列在资源账本下可行},\\
  G_3 &: \mathsf{Vul}\Rightarrow \text{存在结构性脆弱性/覆盖缺口需被修复或规避},\\
  G_4 &: G_1\wedge G_2\wedge G_3 \Rightarrow \text{切换理由可审计且与开放系统要件对齐}.
\end{aligned}
\]
\end{Certificate}

\begin{proof}[证明（谓词因果链）]
由证书~\ref{cert:openra-guard-three-criteria}：
对手信念、资源可行与结构脆弱分别给出切换的外生必要证据、内生可行性与几何/拓扑风险依据；
三者合取即可形成可审计的守卫条件，避免仅靠内部状态做盲切换。
定理成立。证毕。
\end{proof}

\subsection{响应型/结构型切换：必要要件、滞回防抖与最小修复序列模板}
\label{sec:openra-switch-guard-templates}

\begin{definition}[响应型切换与结构型切换]
称扇区切换 $\Sigma\to\Sigma'$ 为：
\begin{enumerate}
  \item \textbf{响应型}：其切换目的主要是响应特定对手威胁类型；在该口径下，$\mathsf{Opp}_{\Sigma\to\Sigma'}$ 通常作为必要要件进入守卫条件；
  \item \textbf{结构型/约束型}：其切换目的主要是修复资源不可行或布局结构性脆弱（即便对手未显式施压也必须切换）；在该口径下，$\mathsf{Res}$ 或 $\mathsf{Vul}$ 可能单独构成必要要件。
\end{enumerate}
\end{definition}

\begin{proposition}[命题：响应型切换中，对手信念是必要的判别要件（可分辨性条件）]
\label{prop:openra-opponent-belief-necessary}
设响应型切换 $\Sigma\to\Sigma'$ 的设计目标要求：
在存在对手威胁类型 $\theta_t\in\Theta_{\Sigma\to\Sigma'}$ 时触发切换，而当 $\theta_t\notin\Theta_{\Sigma\to\Sigma'}$ 时不触发切换。
若切换守卫条件只依赖我方内部状态与内部可行性字段（不依赖任何关于 $\theta_t$ 的情报/信念），
则当存在两段历史在同一时刻具有相同的我方内部状态但对手策略类型不同，守卫条件将无法区分两种情形，从而无法同时满足上述设计目标。
\end{proposition}

\begin{Certificate}[证书：同态不可分辨性（同内态不同外态）]
\label{cert:openra-opponent-belief-necessary}
证书登记如下谓词链：
\[
\begin{aligned}
  B_1 &: \exists t,\ \exists h_1,h_2:\ s_t(h_1)=s_t(h_2)\ \wedge\ \theta_t(h_1)\in\Theta_{\Sigma\to\Sigma'}\ \wedge\
    \theta_t(h_2)\notin\Theta_{\Sigma\to\Sigma'},\\
  B_2 &: G_{\Sigma\to\Sigma'}\ \text{仅依赖}\ s_t \Rightarrow G_{\Sigma\to\Sigma'}(t;h_1)=G_{\Sigma\to\Sigma'}(t;h_2),\\
  B_3 &: \Rightarrow\ \text{无法同时满足“对威胁触发/对非威胁不触发”的设计目标}.
\end{aligned}
\]
\end{Certificate}

\begin{proof}[证明（谓词因果链）]
由证书~\ref{cert:openra-opponent-belief-necessary}：
若守卫条件不依赖对手信念，则在两段历史具有相同内态时其输出必相同；
但响应型切换目标要求对两段历史给出不同决策，矛盾。
因此在响应型切换中，对手信念（或其代理）是必要判别要件。证毕。
\end{proof}

\begin{definition}[防抖与滞回：连续满足与回切阈值]
为避免误判与抖动，定义两类机制：
\begin{enumerate}
  \item \textbf{持续满足（防抖）}：对任意布尔谓词 $G(t)$ 与整数 $L\ge 1$，定义
  \[
    \mathsf{Hold}_{L}(G,t):=\bigwedge_{k=0}^{L-1} G(t-k),
  \]
  即连续 $L$ 步满足才允许触发；
  \item \textbf{滞回}：对某个阈值判据（例如 $\Pr(\cdot)\ge\tau$），登记进入阈值 $\tau_{\mathrm{in}}$ 与退出阈值 $\tau_{\mathrm{out}}$，
  并要求 $\tau_{\mathrm{out}}<\tau_{\mathrm{in}}$，以避免来回震荡。
\end{enumerate}
\end{definition}

\begin{lemma}[引理：防抖与滞回可抑制短暂假信号导致的扇区震荡（条件引理）]
\label{lem:openra-debounce-hysteresis-reduces-chattering}
若扇区切换由 $\mathsf{Hold}_{L}(G_{\Sigma\to\Sigma'},t)$ 触发，并且回切采用更严格的退出阈值（滞回），
则由短暂情报噪声或瞬态资源波动引发的“来回切换”将被抑制到至少 $L$ 步尺度以上。
\end{lemma}

\begin{Certificate}[证书：连续满足 + 退出更严格 $\Rightarrow$ 不会被单步噪声触发回切]
\label{cert:openra-debounce-hysteresis-reduces-chattering}
证书登记：单步噪声不足以满足 $\mathsf{Hold}_{L}$；滞回使回切条件更难满足；二者叠加给出抑制震荡的充分条件。
\end{Certificate}

\begin{proof}[证明（谓词因果链）]
由证书~\ref{cert:openra-debounce-hysteresis-reduces-chattering}：
防抖要求连续 $L$ 步满足，因此单步假信号不会触发切换；
滞回要求退出阈值更严格，因此短暂回落也不致立刻回切。
故可抑制短暂噪声导致的来回震荡。引理成立。证毕。
\end{proof}

\begin{definition}[最小修复序列与不可交换约束（切换模板）]
对每条切换 $\Sigma\to\Sigma'$，登记一份切换模板：
\[
  \mathsf{Tpl}_{\Sigma\to\Sigma'}
  :=
  \bigl(
    \tau,\ L,\ \text{滞回参数},\ w_{\min},\ \mathsf{NonComm}
  \bigr),
\]
其中 $w_{\min}=g_{m-1}\circ\cdots\circ g_0$ 为切换后必须优先执行的最小修复序列（通常由边缘算子与保障算子组成），
$\mathsf{NonComm}$ 为不可交换约束集合（明确哪些相邻算子顺序不可调换，例如“补电”与“上关键建筑”的先后）。
该模板将“非交换序列”从描述语句落实为可执行控制逻辑。
\end{definition}

\begin{remark}[与算法栈切换的关系]
切换模板只规定“何时切、切后先做什么”；
而命题~\ref{prop:openra-sector-switch-requires-algorithm-switch} 规定“切换必须伴随算法栈切换”。
二者合并后，开放系统扇区切换才具备可行性、鲁棒性与可审计性。
\end{remark}

\subsection{生成泛函中心论：行为由变分泛函生成，扇区为派生表达}
\label{sec:openra-generating-functional-centrality}

\begin{remark}[主张（生成泛函决定行为；扇区枚举仅是表达层选择）]
在 \textsf{OpenRA} 这类开放系统博弈中，若采用“非交换算子词（build order/action word）+ 同伦修复塔（L0--L5）+ 字典序正则”的框架，
则“是否显式枚举扇区（经济/防空/科技/反骚扰等）”更多属于\textbf{表达层}选择；
真正决定策略系统行为的是\textbf{生成泛函/变分泛函}：它规定
何为更优、何时触发更强修复、以及修复代价如何计量，
从而在可审计证书（replay 事件流）约束下生成那条非交换算子词及其修复塔跃迁。
\end{remark}

\begin{definition}[回放事件流证书（可回放、可对齐、可审计）]
\label{def:openra-replay-event-stream}
令 $\gamma$ 表示一局对局在固定版本与输入口径下的可回放轨迹。
由回放（replay）可抽取一条事件流证书
\[
  \mathsf{Ev}(\gamma)=\bigl(e_0,e_1,\dots,e_{N-1}\bigr),
\]
其中事件 $e_k$ 可以包含：建筑开始/完成、单位训练开始/完成、供电盈余/赤字变化、采矿循环中断、侦察发现、交火与损失、关键指令切换等。
称 $\mathsf{Ev}(\gamma)$ 为轨迹 $\gamma$ 的审计证书载体。
\end{definition}

\begin{definition}[绩效路径积分泛函（评价/学习生成泛函）]
\label{def:openra-performance-action-functional}
给定任意局部评价泛函 $\mathcal{L}:\mathsf{Ev}\to\RR$（对单个事件给出可计算的局部代价/收益），定义路径积分式绩效泛函
\[
  S_{\mathrm{perf}}[\gamma;\mathcal{L}]
  :=
  \sum_{e\in \mathsf{Ev}(\gamma)} \mathcal{L}(e).
\]
若 $\mathcal{L}$ 专注于“从开局到战备成型”的事件结构（经济上线、供电稳定、生产线成型、覆盖成型等），则 $S_{\mathrm{perf}}$ 可视为一种“Game Readiness Level（GRL）口径”的评价泛函；
本笔记只要求：$\mathcal{L}$ 可由回放证书复算，从而允许对任意策略轨迹做一致对齐与审计。
\end{definition}

\begin{definition}[结构缺口证书与修复烈度计数（同伦修复塔的代价坐标）]
\label{def:openra-gap-and-tau}
对轨迹（或当前态势片段）定义一个非负缺口量
\[
  J(\gamma):=\mathrm{Gap}(\gamma)\in\RR_{\ge 0},
\]
用于证书化“结构性缺口”（例如被硬克制、矿线崩溃、断电停摆、基地拓扑被穿透等）。
并定义一组修复烈度计数向量
\[
  \tau
  :=
  \bigl(\tau_{\mathrm{pred}},\tau_{\mathrm{org}},\tau_{\mathrm{pos}},\tau_{\mathrm{param}},\tau_{\mathrm{algo}},
    \tau_{\mathrm{rev}}\bigr)\in\NN^6,
\]
分别对应：对手信念扭曲强度、队列组织重排强度、几何/布局重构强度、参数调节强度、算法栈切换强度、回退/止损强度。
给定权重 $w\in\RR_{>0}^6$，定义加权烈度
\[
  |\tau|_w := \sum_{i=1}^{6} w_i\,\tau_i.
\]
\end{definition}

\begin{definition}[字典序正则：主损失优先，其次最小烈度]
\label{def:openra-lexicographic-regularization}
对任意二元对 $(a,b),(a',b')\in\RR^2$，定义字典序
\[
  (a,b)\prec_{\mathrm{lex}}(a',b')
  \iff
  \bigl(a<a'\bigr)\ \vee\ \bigl(a=a'\ \wedge\ b<b'\bigr).
\]
称以 $\prec_{\mathrm{lex}}$ 最小化 $(J(\cdot),|\tau(\cdot)|_w)$ 为“字典序正则”：
先最小化缺口 $J$，在缺口同等最优时再最小化修复烈度 $|\tau|_w$。
\end{definition}

\begin{proposition}[命题：给定生成泛函后，扇区枚举可视为派生标签]
\label{prop:openra-sector-derived-from-functional}
固定事件流证书口径（定义~\ref{def:openra-replay-event-stream}）与生成泛函族（定义~\ref{def:openra-performance-action-functional}、定义~\ref{def:openra-gap-and-tau}、定义~\ref{def:openra-lexicographic-regularization}），
设在线决策在每个时刻只需在有限候选集合上比较字典序目标并选取最优候选。
则所谓“扇区”（例如经济/防空/科技/反骚扰）可以定义为如下派生标签：
\[
  \Sigma := \bigl(\text{活跃硬约束集合},\ \text{活跃缺口类型集合},\ \text{允许的修复层级}\bigr),
\]
它仅用于表达与审计（便于人类复盘与规则备案），不会改变由生成泛函诱导的最优候选选择。
\end{proposition}

\begin{Certificate}[证书：派生标签不改变最优解（同一目标与候选集）]
\label{cert:openra-sector-derived-from-functional}
证书登记：在线决策只依赖“候选集 + 比较准则（字典序目标）”；
扇区标签 $\Sigma$ 由硬约束活跃集与缺口类型等派生得到，仅改变描述坐标，不改变比较准则与候选集本身，
因此不会改变最优候选的选择结果。
\end{Certificate}

\begin{proof}[证明（证书；谓词因果链）]
由证书~\ref{cert:openra-sector-derived-from-functional}：
若两种描述方式（显式枚举扇区/不枚举扇区）对应相同候选集合与相同字典序比较准则，
则其在每一步选择的最优候选相同；
而扇区标签作为派生描述不改变候选集合与比较准则，故结论成立。证毕。
\end{proof}

\subsection{同伦修复塔（L0--L5）：受限变分族与字典序正则}
\label{sec:openra-repair-tower-lexicographic}

\begin{definition}[同伦修复塔：分层许可的候选变分族]
\label{def:openra-repair-tower}
对任意时刻的态势片段（或短窗口轨迹）$\gamma$，定义一族分层候选集合
\[
  \mathcal{U}_0(\gamma)\subseteq \mathcal{U}_1(\gamma)\subseteq \cdots \subseteq \mathcal{U}_5(\gamma),
\]
并将层级 $L\in\{0,1,2,3,4,5\}$ 解释为“允许的修复权限/变分类型”：
\begin{enumerate}
  \item L0：不修复（坚持当前算子词与组织）；
  \item L1：\textbf{slicing}（短窗口内选最小可行动作子集，如补电/补关键覆盖）；
  \item L2：\textbf{profile 同伦}（调生产比例/侦察频率/阈值参数与轻量策略参数）；
  \item L3：\textbf{组织同伦}（允许对建造/生产队列与优先级做有限重排，候选交换集合受限）；
  \item L4：\textbf{战略重写}（切科技线/切兵种体系/基地几何大改造）；
  \item L5：\textbf{全局退出/止损}（转场、全军一波、卖建筑止损等）。
\end{enumerate}
其中层级越高，允许的结构性动作越强，通常对应更大的 $\tau_{\mathrm{org}},\tau_{\mathrm{pos}},\tau_{\mathrm{rev}}$ 等计数。
\end{definition}

\begin{theorem}[公理降级：修复塔的逐层扩张与缺口最优值的单调性]
\label{thm:openra-repair-tower-monotone}
在定义~\ref{def:openra-repair-tower} 的修复塔上，令
\[
  J_L^\star(\gamma)
  :=
  \min_{u\in\mathcal{U}_L(\gamma)} J(\gamma\uparrow u),
\]
其中 $\gamma\uparrow u$ 表示对 $\gamma$ 施加候选修复 $u$ 后得到的短窗口轨迹（或态势更新）。
则对任意 $L\in\{0,1,2,3,4\}$ 有
\[
  J_{L+1}^\star(\gamma)\le J_L^\star(\gamma),
\]
即允许更高层修复不会使“可达最优缺口”变差。
\end{theorem}

\begin{Certificate}[数学归纳法证书：按层级 $L$ 归纳的单调性]
\label{cert:openra-repair-tower-monotone}
定义谓词 $P(L)$ 为：“对任意 $\gamma$，有 $\mathcal{U}_L(\gamma)\subseteq\mathcal{U}_{L+1}(\gamma)$ 且 $J_{L+1}^\star(\gamma)\le
  J_L^\star(\gamma)$。”
证书化要求为：
\[
  P(0)\ \text{成立},\qquad
  (\forall L\in\{0,1,2,3,4\})\ \bigl(P(L)\Rightarrow P(L+1)\bigr).
\]
\end{Certificate}

\begin{proof}[证明（数学归纳法证书；谓词因果链）]
由证书~\ref{cert:openra-repair-tower-monotone}：
\begin{itemize}
  \item 基步：$L=0$ 时，$\mathcal{U}_0(\gamma)\subseteq\mathcal{U}_1(\gamma)$ 由定义~\ref{def:openra-repair-tower} 的“逐层许可”给出；
  因为在更大的候选集上取最小值不会变大，故 $J_1^\star(\gamma)\le J_0^\star(\gamma)$，从而 $P(0)$ 成立。
  \item 归纳步：假设 $P(L)$ 成立，则 $\mathcal{U}_{L+1}(\gamma)\subseteq\mathcal{U}_{L+2}(\gamma)$ 同理由逐层许可成立；
  由“最小值在超集上不增”得到 $J_{L+2}^\star(\gamma)\le J_{L+1}^\star(\gamma)$，从而 $P(L+1)$ 成立。
\end{itemize}
由归纳得到对所有 $L\in\{0,1,2,3,4\}$ 单调性成立。证毕。
\end{proof}

\begin{corollary}[推论：字典序正则实现“低层优先”的修复选择]
\label{cor:openra-lexicographic-implies-low-level-first}
固定权重 $w\in\RR_{>0}^6$。
在每个时刻对有限候选集合 $\mathcal{U}_5(\gamma)$ 执行字典序最小化
\[
  u^\star \in \operatorname*{arg\,min}_{u\in\mathcal{U}_5(\gamma)} (J(\gamma\uparrow u),|\tau(u)|_w)_{\mathrm{lex}},
\]
则当存在某个最小层级 $L^\star$ 使得 $J_{L^\star}^\star(\gamma)=J_5^\star(\gamma)$（即更高层不再带来缺口改进）时，
可选择一个字典序最优解满足 $u^\star\in \mathcal{U}_{L^\star}(\gamma)$，
从而在“缺口同等最优”意义下实现低层优先。
\end{corollary}

\begin{Certificate}[证书：若更高层不改进缺口，则字典序会选更小烈度的低层代表]
\label{cert:openra-lexicographic-implies-low-level-first}
证书登记：当 $J_{L^\star}^\star=J_5^\star$ 时，存在 $u_0\in\mathcal{U}_{L^\star}$ 使 $J(\gamma\uparrow u_0)=J_5^\star$；
任取 $u\in\mathcal{U}_5$ 若同样达到 $J_5^\star$ 但具有更大烈度 $|\tau(u)|_w$，
则 $(J(\gamma\uparrow u_0),|\tau(u_0)|_w)\prec_{\mathrm{lex}}(J(\gamma\uparrow u),|\tau(u)|_w)$，
因而字典序最优可取为低层代表元。
\end{Certificate}

\begin{proof}[证明（证书；谓词因果链）]
由证书~\ref{cert:openra-lexicographic-implies-low-level-first}：
当更高层不再改进缺口时，低层已存在达到全局最小缺口的代表元；
字典序比较在缺口相同的候选之间按烈度择小，因此可选取烈度更小的低层代表元作为字典序最优解。
推论成立。证毕。
\end{proof}

\subsection{局部变分节点的证书化：L1 slicing 与 L3 组织同伦}
\label{sec:openra-variational-local-certificates}

\begin{definition}[L1 slicing：紧急动作池、硬约束与字典序选择]
\label{def:openra-l1-slicing}
在短窗口内构造一个有限紧急动作池 $\mathcal{Q}$（例如补电厂、补关键防空、补侦察、补矿线防护等）。
对任意子集 $S\subseteq\mathcal{Q}$，定义硬约束谓词
\[
  \mathsf{Feas}(S):=
  \Bigl(\sum_{o\in S}\mathrm{AA}(o)\ge L_{\mathrm{AA}}\Bigr)
  \wedge
  \Bigl(\sum_{o\in S}\mathrm{Power}(o)\ge L_{\mathrm{Power}}\Bigr)
  \wedge
  (|S|\le N).
\]
并定义三元字典序目标
\[
  \min_{\mathrm{lex}}
  \Bigl(
    \sum_{o\in S}\mathrm{OverBuild}(o),\
    \sum_{o\in S}\mathrm{TimeToOnline}(o),\
    |S|
  \Bigr),
\]
用于在满足生存/上线硬约束的前提下抑制过度建造与过慢上线。
\end{definition}

\begin{lemma}[引理：在有限动作池上，L1 slicing 可给出严格最优证书]
\label{lem:openra-l1-slicing-optimal-certificate}
若 $\mathcal{Q}$ 有限且存在至少一个可行子集满足 $\mathsf{Feas}(S)$，
则存在严格最优的 $S^\star\subseteq\mathcal{Q}$ 使得
$S^\star$ 在所有可行子集中按定义~\ref{def:openra-l1-slicing} 的字典序目标最小；
并且该最优性可由“枚举式比较证书”审计复现。
\end{lemma}

\begin{Certificate}[证书：有限可行集的字典序最小元存在（可枚举）]
\label{cert:openra-l1-slicing-optimal-certificate}
证书登记：可行集 $\{S\subseteq\mathcal{Q}\mid \mathsf{Feas}(S)\}$ 是有限集合；
对每个可行 $S$ 计算字典序目标并取最小元，即得 $S^\star$；
将“所有候选的目标值 + 选择最小者”的记录作为可复算审计证书。
\end{Certificate}

\begin{proof}[证明（证书；谓词因果链）]
由证书~\ref{cert:openra-l1-slicing-optimal-certificate}：
有限非空集合在全序（字典序）下存在最小元；
枚举并比较即给出严格最优解与可审计证书。
引理成立。证毕。
\end{proof}

\begin{definition}[L3 组织同伦：有限重排候选集与 $\varepsilon$-近最优证书]
\label{def:openra-l3-organization-homotopy}
令 $\mathcal{C}(\gamma)$ 为一个有限的“队列重排/优先级重写”候选集合，包含空操作 $\varnothing$。
对任意 $c\in\mathcal{C}(\gamma)$ 定义缺口代价 $J(c):=\mathrm{Gap}(\gamma\uparrow c)$。
给定容差 $\varepsilon\ge 0$，称 $c^\star$ 为 $\varepsilon$-近最优，若
\[
  J(c^\star)\le \min_{c\in\mathcal{C}(\gamma)} J(c) + \varepsilon,
\]
并且满足“确有改进”的防抖判据
\[
  J(c^\star)+\varepsilon < J(\varnothing),
\]
以避免无意义抖动重排。
\end{definition}

\begin{lemma}[引理：有限候选上的 $\varepsilon$-近最优重排可证书化]
\label{lem:openra-l3-epsilon-optimal-certificate}
在定义~\ref{def:openra-l3-organization-homotopy} 的有限候选集上，
若存在 $c$ 使 $J(c)+\varepsilon<J(\varnothing)$，则存在满足条件的 $\varepsilon$-近最优 $c^\star$；
并且其近最优性与“确有改进”均可由回放证书与枚举比较记录复算审计。
\end{lemma}

\begin{Certificate}[证书：有限集取最小元 + 与空操作比较]
\label{cert:openra-l3-epsilon-optimal-certificate}
证书登记：对有限集合 $\mathcal{C}(\gamma)$ 枚举计算 $J(c)$ 并取最小元 $\hat c$；
若 $J(\hat c)+\varepsilon<J(\varnothing)$ 则取 $c^\star=\hat c$；
否则拒绝重排（以 $\varnothing$ 表示），从而满足防抖判据。
\end{Certificate}

\begin{proof}[证明（证书；谓词因果链）]
由证书~\ref{cert:openra-l3-epsilon-optimal-certificate}：
有限集存在最小元 $\hat c$，满足 $J(\hat c)=\min J$；
若其满足“确有改进”判据，则 $c^\star=\hat c$ 同时满足 $\varepsilon$-近最优与防抖；
否则输出空操作避免抖动。
引理成立。证毕。
\end{proof}

\subsection{主泛函合成与在线下降：回放评价、修复正则与塔行为涌现}
\label{sec:openra-master-functional-online-descent}

\begin{definition}[主生成泛函：绩效路径积分 + 缺口聚合 + 修复烈度正则 + 硬约束壁垒]
\label{def:openra-master-functional}
对轨迹 $\gamma$ 定义一个主生成泛函（用于离线对齐/回放复盘/参数学习）：
\[
  \mathcal{S}_{\mathrm{master}}[\gamma]
  :=
  S_{\mathrm{perf}}[\gamma;\mathcal{L}_{\mathrm{perf}}]
  +
  \alpha\sum_{t}\mathrm{Gap}(\gamma_t)
  +
  \beta\,|\tau|_w
  +
  \mathcal{B}_{\mathrm{hard}}(\mathsf{I}),
  \qquad
  \alpha,\beta>0,
\]
其中 $\gamma_t$ 表示按窗口/时刻切片的局部态势，$\mathsf{I}$ 表示硬约束集合（例如关键资产不得长时断电、经济不得归零、侦察覆盖不得为零等），
$\mathcal{B}_{\mathrm{hard}}$ 为硬约束壁垒项（违反即给出不可接受的代价）。
\end{definition}

\begin{proposition}[命题：有限候选的离散变分下降足以诱导“修复塔行为”]
\label{prop:openra-online-descent-induces-tower}
在每个时刻 $t$，若在线决策只在有限候选集合 $\mathcal{U}_5(\gamma_t)$ 上比较两个增量：
\[
  \Delta_u J := J(\gamma_t\uparrow u)-J(\gamma_t),
  \qquad
  \Delta_u|\tau|_w := |\tau(u)|_w,
\]
并按字典序优先选择“缺口下降最大且烈度最小”的候选（或其带阈值/滞回的近似版本），
则在不求解全局最小化的条件下，仍可自然出现：
\begin{enumerate}
  \item 缺口优先下降（先把关键缺口压到可控域）；
  \item 低层优先（在缺口收益相当时选更小烈度）；
  \item 修复塔涌现（高层修复因烈度正则被抑制，除非缺口收益足够大）。
\end{enumerate}
\end{proposition}

\begin{Certificate}[证书：字典序比较 + 有限候选 $\Rightarrow$ 选择规则可执行且可审计]
\label{cert:openra-online-descent-induces-tower}
证书登记：对有限候选 $u$ 逐一复算 $(\Delta_u J,\Delta_u|\tau|_w)$ 并取字典序最小者，
即可得到可执行动作；
将候选列表、增量计算与选择结果记录为审计证书即可复现在线决策。
\end{Certificate}

\begin{proof}[证明（证书；谓词因果链）]
由证书~\ref{cert:openra-online-descent-induces-tower}：
有限候选上字典序最小元可直接枚举求得；
字典序首先压缺口，因而得到“缺口优先下降”；
在缺口增量相同或容差内等价时，选择更小烈度候选，得到“低层优先”；
高层修复一般携带更大烈度计数，若缺口收益不足则不会被字典序选中，故形成“修复塔涌现”。
命题成立。证毕。
\end{proof}

\begin{remark}[与守卫条件的接口：缺口/烈度可作为切换触发量]
第~\ref{sec:openra-guard-three-criteria} 节的守卫条件三要件（对手信念—资源边界—布局脆弱性）可以与本节的缺口 $J$ 与烈度 $\tau$ 对接：
当 $J$ 持续高企且在低层候选族内无法下降（或需不可接受烈度），即可触发跨扇区/升层修复；
这相当于把“扇区切换”内化为“候选变分族与硬约束活跃集的改变”，并保持可审计一致性。
\end{remark}

\subsection{术语映射：从“证书/算子词/GRL/修复塔”到 OpenRA 对象}
\label{sec:openra-generating-functional-terminology}

\begin{remark}[映射表（确保对标可落地）]
\begin{center}
\begin{tabular}{p{0.38\linewidth} p{0.56\linewidth}}
\hline
抽象术语 & OpenRA 对标对象 \\
\hline
证书事件流/可审计轨迹 & replay command log / event stream（定义~\ref{def:openra-replay-event-stream}） \\
非交换算子词 & build order + training order + move/attack command word \\
结构缺口证书 & CounterGap/\allowbreak DefenseGap/\allowbreak EconGap/\allowbreak PowerGap/\allowbreak TopologyGap 的聚合\newline
  $J=\mathrm{Gap}$（定义~\ref{def:openra-gap-and-tau}） \\
同伦修复塔 L0--L5 & 分层许可的候选变分族 $\mathcal{U}_0\subseteq\cdots\subseteq\mathcal{U}_5$\newline （定义~\ref{def:openra-repair-tower}）
  \\
边缘算子代价（$\tau$） & 取消/重排/搬迁/切策略/撤退等结构性代价计数（定义~\ref{def:openra-gap-and-tau}） \\
扇区切换 & 权重/硬约束/候选族激活切换（命题~\ref{prop:openra-sector-derived-from-functional} 与第~\ref{sec:openra-switch-guard-templates} 节）
  \\
\hline
\end{tabular}
\end{center}
\end{remark}

\subsection{可实现 AI 规格：证书快照与一致性复算}
\label{sec:openra-mvp-certificate-snapshot}

\begin{remark}[目标（统一在线对局与离线回放的计算接口）]
为使“同一策略”在在线决策与离线复盘中可对齐、可复算、可审计，
需要把原始数据源（在线游戏状态 / 离线 replay 事件流）统一归约为同一口径的\textbf{证书快照}，
并要求所有缺口指标（Gap）与修复烈度（$\tau$）均仅依赖证书快照计算。
\end{remark}

\begin{definition}[证书快照（每 tick 的统一接口）]
\label{def:openra-certificate-snapshot}
对任意时刻 $t$，定义证书快照为
\[
  \mathsf{Snap}_t
  :=
  \bigl(t,\ s_t,\ \mathsf{Ev}_{t-W:t},\ \widehat{\theta}_t,\ \mathsf{Queue}_t\bigr),
\]
其中：
\begin{enumerate}
  \item $s_t$ 为在线状态摘要（经济、电力、设施状态、单位/建筑清单、已知地图信息等）；
  \item $\mathsf{Ev}_{t-W:t}$ 为长度为 $W$ 的事件流窗口（可由定义~\ref{def:openra-replay-event-stream} 的 $\mathsf{Ev}(\gamma)$ 切片得到）；
  \item $\widehat{\theta}_t$ 为对手策略类型的信念分布（由侦察与交战信息更新）；
  \item $\mathsf{Queue}_t$ 为建造/训练队列与其优先级结构（用于组织同伦的候选生成）。
\end{enumerate}
称 $\mathsf{Snap}_t$ 为“在线/离线统一”的审计证书载体。
\end{definition}

\begin{remark}[事件记录的工程字段（示意）]
事件流窗口中的单条事件可包含如下字段（实现可裁剪）：
\begin{verbatim}
{
  "tick": 12345,
  "type": "BuildingCompleted",
  "player": "self",
  "entity": "PowerPlant",
  "pos": [x, y],
  "meta": { "cost": 300, "buildTime": 450 }
}
\end{verbatim}
常用事件类型最小集可取：建造开始/完成、训练开始/完成、建筑/单位损失、供电变化、资金变化、侦察发现、关键命令下发、扩张点占用/丢失等。
该示意仅用于强调“可回放事件流证书”，不作为硬编码格式要求。
\end{remark}

\begin{theorem}[公理降级：证书快照一致 $\Rightarrow$ 指标与决策可复算一致]
\label{thm:openra-snapshot-consistency-reproducible}
设存在一组确定性计算规则，使得对任意 $t$：
\begin{enumerate}
  \item 缺口指标 $\mathbf{J}_t$ 与烈度计数 $\tau_t$ 仅由 $\mathsf{Snap}_t$ 计算得到；
  \item 在线决策仅由 $\mathsf{Snap}_t$ 生成候选并做字典序选择（第~\ref{sec:openra-online-policy-and-level-switch} 节）。
\end{enumerate}
若离线回放可抽取到与在线一致的证书快照序列 $\{\mathsf{Snap}_t\}$，
则离线复算得到的 $\{\mathbf{J}_t,\tau_t\}$ 与在线一致，且逐步决策结果一致。
\end{theorem}

\begin{Certificate}[数学归纳法证书：按 tick 归纳的可复算一致性]
\label{cert:openra-snapshot-consistency-reproducible}
定义谓词 $P(T)$ 为：“对任意 $t\le T$，若离线与在线的 $\mathsf{Snap}_t$ 相同，则离线复算的 $(\mathbf{J}_t,\tau_t)$ 与在线相同，且本步决策相同。”
证书化要求为：
\[
  P(0)\ \text{成立},\qquad
  (\forall T\in\NN)\ \bigl(P(T)\Rightarrow P(T+1)\bigr).
\]
\end{Certificate}

\begin{proof}[证明（数学归纳法证书；谓词因果链）]
由证书~\ref{cert:openra-snapshot-consistency-reproducible}：
\begin{itemize}
  \item 基步：$t=0$ 时，若 $\mathsf{Snap}_0$ 一致，则由“指标/决策仅依赖 $\mathsf{Snap}_0$ 的确定性规则”可得离线与在线的输出一致，故 $P(0)$ 成立；
  \item 归纳步：假设 $P(T)$ 成立。
  若 $t=T+1$ 时离线与在线的 $\mathsf{Snap}_{T+1}$ 一致，则同理由确定性规则得到本步指标与决策一致；
  因而 $P(T+1)$ 成立。
\end{itemize}
归纳成立，定理成立。证毕。
\end{proof}

\subsection{Gap 证书最小集合：六项核心缺口与字典序总缺口}
\label{sec:openra-gap-certificate-minimal-set}

\begin{remark}[目标（把“为何要修复/升级”证书化）]
本节给出一个最小可用的 Gap 指标集合（6 项），它们共同构成结构缺口证书 $\mathbf{J}$：
每一项都可从证书快照 $\mathsf{Snap}_t$ 计算，并能在 \textsf{OpenRA} 中直接解释为“为何要切换修复层级/动作结构”。
\end{remark}

\begin{definition}[经济续航缺口：$\mathrm{EconRunwayGap}$]
\label{def:openra-gap-econ-runway}
令 $M_t$ 表示当前资金，$\dot M_t$ 表示净现金流估计（收入率减支出率）。
定义续航时间（以秒计）
\[
  \mathrm{Runway}_t
  :=
  \begin{cases}
    \dfrac{M_t}{-\dot M_t}, & \dot M_t < 0,\\[6pt]
    +\infty, & \dot M_t \ge 0,
  \end{cases}
\]
并给定最低续航阈值 $T_{\min}>0$（例如 $30$--$60$ 秒），定义
\[
  \mathrm{EconRunwayGap}_t := \max(0,\ T_{\min}-\mathrm{Runway}_t).
\]
\end{definition}

\begin{definition}[供电缺口：$\mathrm{PowerGap}$]
\label{def:openra-gap-power}
令 $P^{\mathrm{gen}}_t$ 为发电能力，$P^{\mathrm{load}}_t$ 为当前负载，定义裕度
\[
  \mathrm{Margin}_t := P^{\mathrm{gen}}_t - P^{\mathrm{load}}_t.
\]
给定安全裕度 $P_{\mathrm{safe}}>0$ 与未来负载预测 $P^{\mathrm{forecast}}_t$（由即将完成建筑/生产引入的负载估计），定义
\[
  \mathrm{PowerGap}_t
  :=
  \max(0,\ P_{\mathrm{safe}}-\mathrm{Margin}_t)
  +
  \lambda\cdot \max(0,\ P^{\mathrm{forecast}}_t-\mathrm{Margin}_t),
  \qquad \lambda\ge 0.
\]
\end{definition}

\begin{definition}[生产连续性缺口：$\mathrm{ProductionContinuityGap}$]
\label{def:openra-gap-production}
令 $\mathcal{F}$ 为关键生产设施集合（兵营/车厂/机场等），取滑窗长度 $W$（例如 $30$ 秒）。
对每个设施 $f\in\mathcal{F}$，令 $\rho_f$ 为该滑窗内的空闲时间比例：
\[
  \rho_f := \frac{T_{\mathrm{idle}}(f)}{W}.
\]
给定容许上界 $\rho_{\max}\in(0,1)$（例如 $0.25$），定义
\[
  \mathrm{ProdGap}_t := \sum_{f\in\mathcal{F}} \max(0,\ \rho_f-\rho_{\max}).
\]
\end{definition}

\begin{definition}[覆盖缺口：$\mathrm{DefenseCoverageGap}$（反空/反地）]
\label{def:openra-gap-defense-coverage}
选取关键资产区（zones）集合 $\mathcal{Z}$（矿线区、生产区、科技区、电力区、基地核心区等）。
对每个 $\mathrm{zone}\in\mathcal{Z}$，定义我方覆盖能力 $C_{\mathrm{AA}}(\mathrm{zone}),C_{\mathrm{AD}}(\mathrm{zone})$，
以及由对手信念 $\widehat{\theta}_t$ 诱导的威胁强度
\[
  \begin{aligned}
    T_{\mathrm{air}}(\mathrm{zone})
    &:=
    \Pr(\text{空袭类}\mid \widehat{\theta}_t)\cdot \widehat{\mathrm{AirPower}}(\mathrm{zone}),\\
    T_{\mathrm{ground}}(\mathrm{zone})
    &:=
    \Pr(\text{地推类}\mid \widehat{\theta}_t)\cdot \widehat{\mathrm{GroundPower}}(\mathrm{zone}).
  \end{aligned}
\]
定义覆盖缺口
\[
  \mathrm{DefGap}_t
  :=
  \sum_{\mathrm{zone}\in\mathcal{Z}}
  \Bigl[
    \max(0,\ T_{\mathrm{air}}(\mathrm{zone})-C_{\mathrm{AA}}(\mathrm{zone}))
    +
    \max(0,\ T_{\mathrm{ground}}(\mathrm{zone})-C_{\mathrm{AD}}(\mathrm{zone}))
  \Bigr].
\]
\end{definition}

\begin{definition}[情报不确定性缺口：$\mathrm{InfoGap}$]
\label{def:openra-gap-info}
令 $q\in[0,1]$ 表示在滑窗 $W$ 内对关键区域的侦察覆盖率，令 $u\in[0,1]$ 表示关键前置/科技信息的未知程度，
给定系数 $a,b\ge 0$，定义
\[
  \mathrm{InfoGap}_t := a(1-q)+bu.
\]
\end{definition}

\begin{definition}[布局拓扑脆弱缺口：$\mathrm{TopologyBreachGap}$]
\label{def:openra-gap-topology}
选取关键资产集合 $\mathcal{K}$（电力群、生产设施、科技设施、主基地核心等）。
对每个 $k\in\mathcal{K}$，估计对手从可达入口到 $k$ 的穿透代价
\[
  \mathrm{CostToBreach}(k)
  \approx
  \frac{\mathrm{PathLength}(k)}{\mathrm{OurDPSAlongPath}(k)+\varepsilon},
  \qquad \varepsilon>0.
\]
给定阈值 $\theta_{\mathrm{breach}}>0$，定义
\[
  \mathrm{TopoGap}_t
  :=
  \sum_{k\in\mathcal{K}} \max\bigl(0,\ \theta_{\mathrm{breach}}-\mathrm{CostToBreach}(k)\bigr).
\]
\end{definition}

\begin{definition}[字典序总缺口向量（主缺口优先）]
\label{def:openra-gap-vector-lex}
定义总缺口向量
\[
  \mathbf{J}_t
  :=
  \bigl(
    \mathrm{TopoGap}_t,\ \mathrm{DefGap}_t,\ \mathrm{PowerGap}_t,\ \mathrm{EconRunwayGap}_t,\ \mathrm{ProdGap}_t,\
      \mathrm{InfoGap}_t
  \bigr)\in\RR_{\ge 0}^6,
\]
并对 $\RR^6$ 采用字典序（先比第一分量，再比第二分量，依此类推）。
该顺序对应工程直觉：拓扑与覆盖缺口优先级通常高于生产与信息缺口；同时保留可调性（可通过重排分量顺序备案）。
\end{definition}

\begin{remark}[与前文标量缺口 $J$ 的兼容口径]
前文以标量 $J=\mathrm{Gap}$ 讨论“缺口优先下降”等性质（例如命题~\ref{prop:openra-online-descent-induces-tower}）。
工程实现中更推荐使用向量 $\mathbf{J}$ 做主缺口优先；
若需要与标量形式对接，可取任意与字典序单调一致的标量化 $\widehat{J}=\Psi(\mathbf{J})$（例如带大权重的加权和或分层罚项），用于粗粒度统计与日志摘要。
\end{remark}

\begin{proposition}[命题：按优先级排序的缺口比较等价于字典序比较]
\label{prop:openra-gap-lex-ordering}
在定义~\ref{def:openra-gap-vector-lex} 的顺序下，
“先最小化 $\mathrm{TopoGap}$，再最小化 $\mathrm{DefGap}$，依次类推”与对 $\mathbf{J}$ 取字典序最小元等价。
\end{proposition}

\begin{Certificate}[证书：字典序的逐分量判定规则]
\label{cert:openra-gap-lex-ordering}
证书登记：字典序的定义即为“在第一处不同分量上取更小者”；因此其最小化恰等价于按给定优先级逐层最小化分量。
\end{Certificate}

\begin{proof}[证明（证书；谓词因果链）]
由证书~\ref{cert:openra-gap-lex-ordering}：
字典序最小化首先确保第一分量最小；在第一分量相同的候选间，再确保第二分量最小；依此类推。
与“按优先级逐分量最小化”一致。
命题成立。证毕。
\end{proof}

\subsection{修复烈度证书：$\tau$ 六维计数与加权正则}
\label{sec:openra-tau-certificate-implementation}

\begin{remark}[目标（把“修复代价”证书化，避免黑箱大改）]
修复塔允许更高层结构性动作（队列重排、布局重构、策略栈切换、止损回退等），
若不对这些动作计量并正则化，就会导致“为了一次小缺口频繁大改”的策略抖动与不可审计性。
因此需要把修复代价以 $\tau$ 的形式证书化并累积。
\end{remark}

\begin{definition}[候选变分的烈度增量（$\Delta\tau$）]
\label{def:openra-tau-increment}
对任意候选变分 $u$，定义烈度增量
\[
  \Delta\tau(u)
  :=
  \bigl(\Delta\tau_{\mathrm{pred}},\Delta\tau_{\mathrm{org}},\Delta\tau_{\mathrm{pos}},\Delta\tau_{\mathrm{param}},
    \Delta\tau_{\mathrm{algo}},\Delta\tau_{\mathrm{rev}}\bigr)\in\NN^6,
\]
其中各分量的工程口径可以登记为规则：
\begin{enumerate}
  \item $\Delta\tau_{\mathrm{pred}}$：对对手策略分布作强制性假设/扭曲的次数（例如强制进入 anti-air posture）；
  \item $\Delta\tau_{\mathrm{org}}$：队列重排/取消/优先级重写的次数；
  \item $\Delta\tau_{\mathrm{pos}}$：建筑布局迁移/基地形态重构的次数（含关键资产搬迁、前推防线模板触发等）；
  \item $\Delta\tau_{\mathrm{param}}$：阈值/配比/集结点/回防规则等参数修改次数；
  \item $\Delta\tau_{\mathrm{algo}}$：策略栈切换次数（经济$\to$防守$\to$反骚扰$\to$一波等）；
  \item $\Delta\tau_{\mathrm{rev}}$：回退/止损动作次数（取消关键投资、卖建筑、全局撤退/转场等）。
\end{enumerate}
给定权重 $w\in\RR_{>0}^6$，定义烈度代价 $\Delta|\tau|_w:=|\Delta\tau(u)|_w$。
\end{definition}

\begin{lemma}[引理：若 $\Delta\tau$ 可复算，则修复代价可审计追溯]
\label{lem:openra-tau-auditable}
若每次候选变分 $u$ 的类型与参数可由证书快照与回放事件流复算，
并且 $\Delta\tau(u)$ 按定义~\ref{def:openra-tau-increment} 的规则确定，
则累计烈度 $\tau$ 可被离线复算并用于审计“为何发生高层结构性修复”。
\end{lemma}

\begin{Certificate}[证书：逐步累积与可复算性闭合]
\label{cert:openra-tau-auditable}
证书登记：每一步执行的候选变分 $u_t$ 可由回放复现；由规则确定 $\Delta\tau(u_t)$；
累积式 $\tau_{t+1}=\tau_t+\Delta\tau(u_t)$ 仅由可复算量构成，因此可审计追溯。
\end{Certificate}

\begin{proof}[证明（证书；谓词因果链）]
由证书~\ref{cert:openra-tau-auditable}：
每一步 $u_t$ 可复现且 $\Delta\tau(u_t)$ 可复算，故递推累积得到的 $\tau_t$ 可复算；
于是修复代价可审计追溯。
引理成立。证毕。
\end{proof}

\subsection{候选变分生成器：L0--L5 分层模板与受限搜索}
\label{sec:openra-variation-generators}

\begin{remark}[核心原则（只扩大“允许的变分类型”，不扩到全动作空间）]
修复塔的关键不是“层级越高动作越多”，而是：
每层只允许一类新的结构性变分，并保持候选集有限（或可控的有限近似），
从而保证字典序选择可执行、可复算、可审计，并自然呈现“低层优先”的塔行为（推论~\ref{cor:openra-lexicographic-implies-low-level-first}）。
\end{remark}

\begin{definition}[候选变分生成器（分层）]
\label{def:openra-variation-generator}
定义分层候选生成器为一族集合值映射
\[
  \mathsf{Gen}_L:\ \mathsf{Snap}_t \longmapsto \mathcal{U}_L(\mathsf{Snap}_t),
  \qquad L\in\{0,1,2,3,4,5\},
\]
满足 $\mathcal{U}_0\subseteq\cdots\subseteq\mathcal{U}_5$，
并要求每个 $\mathcal{U}_L(\mathsf{Snap}_t)$ 为有限集合（或可控截断下的有限集合），以支撑枚举式证书。
\end{definition}

\begin{definition}[L0：不修复（仅允许微小参数扰动）]
\label{def:openra-gen-l0}
令 $\Theta_{\mathrm{micro}}$ 为一组微调参数（集火优先级、回防阈值、集结半径、侦察频率的微小扰动等）。
定义
\[
  \mathcal{U}_0(\mathsf{Snap}_t)
  :=
  \{\ \Delta\theta\in\Theta_{\mathrm{micro}}\ \},
\]
并约束：不允许队列重排/取消、不允许布局改造、不允许策略栈切换。
\end{definition}

\begin{definition}[L1：slicing（紧急最小补救子集）]
\label{def:openra-gen-l1}
设时间视界 $H$（例如 $30$ 秒），构造一个可在 $H$ 内落地的紧急动作池 $\mathcal{Q}_t$：
补电、补最小防空/防地、补侦察单位、短时矿线保护等。
定义 $\mathcal{U}_1(\mathsf{Snap}_t)$ 为所有满足硬约束的候选子集（见定义~\ref{def:openra-l1-slicing}）或其贪心截断近似，
并要求：候选集在实现层面保持有限（例如限制 $|S|\le N$ 且限制动作池规模）。
\end{definition}

\begin{definition}[L2：profile/param 同伦（配比与参数同伦）]
\label{def:openra-gen-l2}
定义一组可调的配比与策略参数候选集合 $\Theta_{\mathrm{profile}}$：
生产配比（反空/反地占比、矿车数量目标）、侦察频率与路径、回防阈值、预算比例等。
定义
\[
  \mathcal{U}_2(\mathsf{Snap}_t)
  :=
  \mathcal{U}_1(\mathsf{Snap}_t)\ \cup\ \{\ \Delta\theta\in\Theta_{\mathrm{profile}}\ \},
\]
并约束：不允许显式的队列结构重排与布局重构。
\end{definition}

\begin{definition}[L3：组织同伦（有限交换/取消/插入/转产）]
\label{def:openra-gen-l3}
在建造/训练队列 $\mathsf{Queue}_t$ 上，定义有限候选操作集：
\begin{enumerate}
  \item $\mathsf{Swap}$：交换相邻两项或将一项前移 $k$ 位（$k$ 很小）；
  \item $\mathsf{Cancel}$：取消低优先级且可回滚的项；
  \item $\mathsf{Insert}$：在队列头插入一项 L1 紧急动作；
  \item $\mathsf{Retarget}$：令某生产设施切换到另一兵种谱内的单位（受限候选）。
\end{enumerate}
令 $\mathcal{C}_t$ 为“对队列执行至多一次上述操作”的有限候选集合（包含空操作 $\varnothing$，见定义~\ref{def:openra-l3-organization-homotopy}）。
定义
\[
  \mathcal{U}_3(\mathsf{Snap}_t)
  :=
  \mathcal{U}_2(\mathsf{Snap}_t)\ \cup\ \mathcal{C}_t.
\]
\end{definition}

\begin{definition}[L4：强结构重写（布局/科技线/兵种体系模板库）]
\label{def:openra-gen-l4}
为避免“自由组合导致候选爆炸”，将强结构重写登记为有限模板库：
\begin{enumerate}
  \item 几何形态模板：集中$\to$分散、关键资产内缩、瓶颈化入口、前推防线、双基地分区等；
  \item 兵种体系模板：地面为主$\to$反空为主$\to$机动骚扰为主$\to$重装推进为主等；
  \item 科技线模板：停科技保生存、快速上关键防御科技、以反制为核心转线等；
  \item 经济结构模板：主矿线迁移、双矿$\to$三矿、矿线防护优先级翻转等。
\end{enumerate}
令 $\mathcal{T}^{(4)}_t$ 为在 $\mathsf{Snap}_t$ 下可执行的模板候选集合（有限），定义
\[
  \mathcal{U}_4(\mathsf{Snap}_t)
  :=
  \mathcal{U}_3(\mathsf{Snap}_t)\ \cup\ \mathcal{T}^{(4)}_t.
\]
\end{definition}

\begin{definition}[L5：全局止损/转场/一波（全局模板库）]
\label{def:openra-gen-l5}
将全局性动作登记为严格触发的有限模板库（并配置冷却时间）：
全军一波、转场/撤离（若机制允许）、大规模取消/卖出非关键资产止损、强制进入最小可存活模式等待机会等。
令 $\mathcal{T}^{(5)}_t$ 为满足触发条件且未被冷却禁止的模板集合（有限），定义
\[
  \mathcal{U}_5(\mathsf{Snap}_t)
  :=
  \mathcal{U}_4(\mathsf{Snap}_t)\ \cup\ \mathcal{T}^{(5)}_t.
\]
\end{definition}

\begin{theorem}[公理降级：若模板库有限且操作步长有界，则每层候选集有限]
\label{thm:openra-generator-finiteness}
若紧急动作池规模有界、队列操作次数有界（如“至多一次”）、L4/L5 模板库为有限集合，
则对任意 $L\in\{0,1,2,3,4,5\}$ 与任意 $\mathsf{Snap}_t$，候选集合 $\mathcal{U}_L(\mathsf{Snap}_t)$ 有限。
\end{theorem}

\begin{Certificate}[证书：有限并集/有限子集/有界操作的封闭性]
\label{cert:openra-generator-finiteness}
证书登记：L1 的动作池有限 $\Rightarrow$ 可行子集候选有限（或按 $|S|\le N$ 截断有限）；
L3 的“至多一次操作”使候选队列变换数有界；
L4/L5 模板库有限；
并集保持有限性。
\end{Certificate}

\begin{proof}[证明（证书；谓词因果链）]
由证书~\ref{cert:openra-generator-finiteness}：
有限集合在“取子集并施加有界约束”“执行有界次数操作”“取有限并集”下保持有限；
故各层候选集合有限。
定理成立。证毕。
\end{proof}

\subsection{在线决策与层级升降：字典序选择、跳级与滞回}
\label{sec:openra-online-policy-and-level-switch}

\begin{remark}[目标（把“修复塔”落成在线可运行控制律）]
在线对局无法每 tick 求解全局最优；但只要在有限候选集上做字典序选择，并配合层级升降逻辑与防抖滞回，
即可获得：缺口优先、低层优先、必要时才高层重写的“塔行为”，并且全程可审计复现。
\end{remark}

\begin{definition}[在线一步决策：有限候选的字典序选择]
\label{def:openra-online-one-step-policy}
令当前层级为 $L_t$，由 $\mathsf{Gen}_{L_t}$ 生成候选集合 $\mathcal{U}_{L_t}(\mathsf{Snap}_t)$。
对每个候选 $u$，用轻量预测（短视界 $H$ 的近似仿真/估计）得到施加后的缺口向量 $\widehat{\mathbf{J}}(\mathsf{Snap}_t\uparrow u)$，
并计算烈度代价 $\Delta|\tau(u)|_w$。
定义在线选择为
\[
  u_t^\star
  \in
  \operatorname*{arg\,min}_{u\in\mathcal{U}_{L_t}(\mathsf{Snap}_t)}
  \bigl(\widehat{\mathbf{J}}(\mathsf{Snap}_t\uparrow u),\ \Delta|\tau(u)|_w\bigr)_{\mathrm{lex}}.
\]
该规则保证：先压主缺口，再以最小烈度完成修复（在预测口径下）。
\end{definition}

\begin{definition}[升级规则（升层/跳级）：缺口超阈值且低层无显著改进]
\label{def:openra-level-upgrade-rule}
给定“显著改进”阈值谓词 $\mathsf{Improve}(u)$ 与层级阈值 $\Theta_L$。
若满足任一条件则触发升级：
\begin{enumerate}
  \item （常规升层）$\mathbf{J}_t \succ \Theta_{L_t}$ 且对所有 $u\in\mathcal{U}_{L_t}(\mathsf{Snap}_t)$ 均不满足 $\mathsf{Improve}(u)$；

  \item （硬触发跳级）出现硬触发：$\mathrm{PowerGap}_t$ 爆表、$\mathrm{TopoGap}_t$ 爆表、或经济续航极低且敌压强显著等。
\end{enumerate}
升级后令 $L_{t+1}:=\min(L_t+\Delta L,5)$，其中 $\Delta L\ge 1$ 由触发类型决定（硬触发允许 $\Delta L>1$）。
\end{definition}

\begin{definition}[降级规则（降层）：缺口低且稳定 + 滞回/驻留时间]
\label{def:openra-level-downgrade-rule}
给定低缺口阈值 $\Theta_{\mathrm{low}}$ 与驻留时间 $T_{\mathrm{dwell}}$。
若在连续 $T_{\mathrm{dwell}}$ 的时间内均满足 $\mathbf{J}_t \preceq \Theta_{\mathrm{low}}$，
则令 $L_{t+1}:=\max(L_t-1,0)$；
并要求采用滞回与最小驻留时间，以避免层级频繁抖动（对齐第~\ref{sec:openra-switch-guard-templates} 节与引理~\ref{lem:openra-debounce-hysteresis-reduces-chattering}）。
\end{definition}

\begin{theorem}[公理降级：升级/降级与防抖滞回结合可抑制层级抖动]
\label{thm:openra-level-switch-stability}
若层级切换采用定义~\ref{def:openra-level-upgrade-rule} 与定义~\ref{def:openra-level-downgrade-rule}，
并对触发谓词使用持续满足与滞回机制（第~\ref{sec:openra-switch-guard-templates} 节、引理~\ref{lem:openra-debounce-hysteresis-reduces-chattering}）
  ，
则由短暂情报噪声或瞬态资源波动引发的“层级来回切换”将被抑制到至少驻留时间尺度以上。
\end{theorem}

\begin{Certificate}[证书：把扇区防抖/滞回引理平移到层级切换]
\label{cert:openra-level-switch-stability}
证书登记：层级切换的触发谓词与扇区切换同型（均为布尔守卫条件）；
将 $\mathsf{Hold}_L$ 与滞回阈值机制应用到层级触发上，即可复用引理~\ref{lem:openra-debounce-hysteresis-reduces-chattering} 的结论。
\end{Certificate}

\begin{proof}[证明（证书；谓词因果链）]
由证书~\ref{cert:openra-level-switch-stability}：
层级切换触发与扇区切换触发同型；
引理~\ref{lem:openra-debounce-hysteresis-reduces-chattering} 给出“持续满足 + 滞回 $\Rightarrow$ 抑制抖动”的充分条件；
应用于层级切换即得结论。
定理成立。证毕。
\end{proof}

\subsection{实现清单：Gap 计算表与 L0--L5 候选生成表（MVP）}
\label{sec:openra-mvp-gap-and-generator-tables}

\begin{remark}[工程表格（可直接作为实现文档与配置备案）]
为便于直接落地实现，本节将最小 Gap 指标与 L0--L5 候选生成器整理为两张表：
一张描述“如何计算与如何触发”，另一张描述“每层允许什么候选以及如何计量烈度”。
建议将表中阈值、滑窗与模板库版本化，并与回放证书一起归档，形成可审计的策略演化记录。
\end{remark}

\begin{remark}[说明（实现口径约束）]
\begin{enumerate}
  \item tick 到秒的换算可取 \code{sec = tick / ticksPerSecond}（\textsf{OpenRA} 常见为 $25$ tick/s，但以引擎实际为准）；
  \item 所有 Gap 输出均为 $[0,+\infty)$，$0$ 表示无缺口，值越大表示越紧急；
  \item 总缺口建议按字典序比较：$\mathrm{Topo}>\mathrm{Def}>\mathrm{Power}>\mathrm{Econ}>\mathrm{Prod}>\mathrm{Info}$（与定义~\ref{def:openra-gap-vector-lex} 一致）；
  \item “推荐触发层级”指：当该 gap 主导时，最低应允许的修复层级（并不禁止更高层）。
\end{enumerate}
\end{remark}

\begingroup
\scriptsize
\setlength{\tabcolsep}{2pt}
\renewcommand{\arraystretch}{1.2}
\setlength{\LTleft}{0pt}
\setlength{\LTright}{0pt}
\begin{longtable}{|
  >{\raggedright\arraybackslash}p{0.10\linewidth}|
  >{\raggedright\arraybackslash}p{0.11\linewidth}|
  >{\raggedright\arraybackslash}p{0.16\linewidth}|
  >{\raggedright\arraybackslash}p{0.10\linewidth}|
  >{\raggedright\arraybackslash}p{0.24\linewidth}|
  >{\raggedright\arraybackslash}p{0.12\linewidth}|
  >{\raggedright\arraybackslash}p{0.10\linewidth}|}
\hline
\textbf{Gap 名称} & \textbf{工程含义} & \textbf{需要输入（在线状态/事件流）} & \textbf{滑窗/频率} & \textbf{计算（建议实现方式）} & \textbf{阈值建议（默认）} &
  \textbf{推荐触发层级} \\
\hline
\endfirsthead
\hline
\textbf{Gap 名称} & \textbf{工程含义} & \textbf{需要输入（在线状态/事件流）} & \textbf{滑窗/频率} & \textbf{计算（建议实现方式）} & \textbf{阈值建议（默认）} &
  \textbf{推荐触发层级} \\
\hline
\endhead
\textbf{Econ\allowbreak Runway\allowbreak Gap}
& 经济续航不足导致“算子不可行域塌缩”
& 当前资金 \texttt{M}；收入率 \texttt{IncomeRate}；消耗率 \texttt{SpendRate}；矿车被打断/损失事件
& 滑窗 $W=30$--$60\,s$ 更新；每 $1$--$2\,s$ 评估一次
& 估计净流 $\dot M:=\texttt{IncomeRate}-\texttt{SpendRate}$；若 $\dot M<0$，$\mathrm{Runway}:=\texttt{M}/(-\dot M)$；否则
  $\mathrm{Runway}:=+\infty$；
\newline Gap $:=\max(0,T_{\min}-\mathrm{Runway})$；可叠加“矿线被打断惩罚”
& $T_{\min}=30$--$60\,s$；惩罚系数 $k=0.5$--$2.0$
& L1（紧急止损）到 L5（无解时）
\\ \hline
\textbf{PowerGap}
& 供电赤字或即将赤字导致关键能力停摆
& 当前发电 \texttt{Pgen}、负载 \texttt{Pload}；队列预测负载 \texttt{Pforecast}；断电事件
& 每 tick 维护；每 $0.5$--$1\,s$ 输出
& $\mathrm{Margin}:=\texttt{Pgen}-\texttt{Pload}$；
\newline Gap $:=\max(0,P_{\mathrm{safe}}-\mathrm{Margin})+\lambda\max(0,\texttt{Pforecast}-\mathrm{Margin})$；
\newline \texttt{Pforecast} 可由“视界 $T_{\mathrm{horizon}}$ 内即将上线的耗电项”累加估计
& $P_{\mathrm{safe}}$ 取负载的 $10$--$20\%$；
\newline $T_{\mathrm{horizon}}=20$--$40\,s$；$\lambda=0.5$--$1.0$
& L1/L3（补电或重排/取消）
\\ \hline
\textbf{ProdGap}
& 生产连续性断档（产能闲置/队列占用错误）
& 关键生产设施集合 $\mathcal{F}$；滑窗内 Idle 时间；队列占用情况
& 滑窗 $W=30$--$45\,s$；每 $2$--$3\,s$ 评估
& 对每 $f\in\mathcal{F}$：$\rho_f:=T_{\mathrm{idle}}(f)/W$；
\newline Gap $:=\sum_{f\in\mathcal{F}}\max(0,\rho_f-\rho_{\max})$；
\newline 可加入“关键兵种缺口”附加项（如反空占比过低）
& $\rho_{\max}=0.25$--$0.35$
& L2/L3（改配比/重排）
\\ \hline
\textbf{Def\allowbreak Coverage\allowbreak Gap}
& 防空/防地覆盖不足（与对手威胁耦合）
& zone 列表（矿线、生产、科技、电力、核心）；
\newline 覆盖能力：\texttt{CAA}\newline \texttt{CAD}；
\newline 对手威胁信念：$\Pr(\mathrm{air})$\newline $\Pr(\mathrm{ground})$；
\newline 敌方强度估计：\texttt{AirPower}\newline \texttt{GroundPower}
& 每 $1$--$2\,s$；侦察事件触发后立即更新
& \(
\begin{aligned}
  T_a &:= \Pr(a)\cdot \widehat{P}_a,\\
  T_g &:= \Pr(g)\cdot \widehat{P}_g,\\
  \mathrm{Gap}
  &:= \sum_{\mathrm{zone}} \Bigl[\\
    &\qquad \max(0,T_a-\texttt{CAA})\\
    &\qquad +\max(0,T_g-\texttt{CAD})
  \Bigr].
\end{aligned}
\)
\newline \texttt{CAA/CAD} 可先用“防御 DPS 加权 + 距离衰减”粗算
& $\Pr(\mathrm{air})$ 与 $\Pr(\mathrm{ground})$
\newline 触发阈值 $0.55$--$0.70$；
\newline Gap 阈值按地图规模缩放
& L1/\allowbreak L3/\allowbreak L4
\newline 补防$\to$重排$\to$重构
\\ \hline
\textbf{InfoGap}
& 情报缺口导致误切扇区/被诱导
& 关键区域侦察覆盖 $q$（最近 $W$ 秒可见比例）；
\newline 敌方关键前置/科技未知度 $u$；
\newline 侦察单位存活与路径
& 滑窗 $W=30$--$60\,s$；每 $2$--$3\,s$
& Gap $:=a(1-q)+bu$；
\newline $q$ 可用“关键区域最近一次可见时间”计算；
\newline $u$ 可用“关键建筑/科技状态未证实的概率和”近似
& 初始可取 $a=b=1$；
\newline 若被诱导频繁可增大 $a$
& L2（提高侦察/保守化）
\\ \hline
\textbf{Topology\allowbreak Breach\allowbreak Gap}
& 基地拓扑/布局结构性脆弱：存在低成本穿透到关键资产的路径
& 关键资产集合 $\mathcal{K}$（电力/生产/科技/矿线/核心等）；
\newline 入口集合；地形可达图；
\newline 沿路径火力覆盖估计
& 每 $2$--$5\,s$；建筑完工/摧毁后立即更新
& 对每 $k$：
\newline \(
\begin{aligned}
  \ell_k &:= \mathrm{PathLen}(k),\\
  d_k &:= \mathrm{DPS}_{\mathrm{path}}(k),\\
  C_k &:= \ell_k/(d_k+\varepsilon),\\
  \mathrm{Gap} &:= \sum_{k\in\mathcal{K}}\max(0,\theta-C_k).
\end{aligned}
\)
\newline \texttt{DPS\_path} 可用路径采样点附近防御 DPS 近似
& $\theta_{\mathrm{breach}}$ 可先取等效 $6$--$12\,s$ 量级并调参
& L4（几何重构/结构改写），严重时 L5
\\ \hline
\end{longtable}
\endgroup

\begingroup
\scriptsize
\setlength{\tabcolsep}{2pt}
\renewcommand{\arraystretch}{1.2}
\setlength{\LTleft}{0pt}
\setlength{\LTright}{0pt}
\begin{longtable}{|
  >{\raggedright\arraybackslash}p{0.07\linewidth}|
  >{\raggedright\arraybackslash}p{0.09\linewidth}|
  >{\raggedright\arraybackslash}p{0.23\linewidth}|
  >{\raggedright\arraybackslash}p{0.12\linewidth}|
  >{\raggedright\arraybackslash}p{0.12\linewidth}|
  >{\raggedright\arraybackslash}p{0.15\linewidth}|
  >{\raggedright\arraybackslash}p{0.15\linewidth}|}
\hline
\textbf{层级} & \textbf{允许变分类型} & \textbf{候选生成器（模板库，最小集）} & \textbf{$\Delta\tau$ 规则} & \textbf{强非交换约束} & \textbf{升级触发（进入该层）
  } & \textbf{降级条件（退出该层）} \\
\hline
\endfirsthead
\hline
\textbf{层级} & \textbf{允许变分类型} & \textbf{候选生成器（模板库，最小集）} & \textbf{$\Delta\tau$ 规则} & \textbf{强非交换约束} & \textbf{升级触发（进入该层）
  } & \textbf{降级条件（退出该层）} \\
\hline
\endhead
\textbf{L0}
& 仅微小参数微调；不动结构
& 1) 集火优先级微调\newline 2) 回防阈值微调\newline 3) 集结点微调（不迁移建筑）
& $\Delta\tau_{\mathrm{param}}=1$（按“是否修改”计）；其余为 $0$
& 禁止：取消建造、重排队列、切策略、搬迁建筑
& 所有 Gap 低且稳定；或 InfoGap 高但无硬缺口
& 若硬缺口（Power/\allowbreak Def/\allowbreak Topo）越阈值，
\newline 立即升级到 L1/\allowbreak L3/\allowbreak L4
\\ \hline
\textbf{L1}
& slicing：短视界最小动作子集
& 候选池 $\mathcal{Q}$（建议 $20$--$40$ 个以内）：
\newline 补电（电厂/备用电）；
\newline 最小防空/防地；
\newline 最小侦察补齐；
\newline 矿线短时保护；
\newline 关键建筑抢修/替代（若机制支持）
& 若插队/抢占队列：$\Delta\tau_{\mathrm{org}}=1$；
\newline 若修改侦察/回防阈值：$\Delta\tau_{\mathrm{param}}=1$
& 常见约束：\textbf{先补电再上耗电项}；
\newline Runway 极短时：\textbf{先止血再扩张}
& PowerGap/\allowbreak EconRunwayGap/\allowbreak DefGap 超阈值且存在短链显著下降
& 连续 $20$--$40\,s$ 关键 Gap 低位且无新威胁证据，可降到 L0/L2
\\ \hline
\textbf{L2}
& profile/\allowbreak param 同伦：
\newline 改配比/阈值/侦察，不改结构
& 1) 生产配比向量（反空/反地/侦察/矿车）；
\newline 2) 侦察周期与路线模板（$2$--$4$ 条）；
\newline 3) 预算分配（防御/扩张/科技）；
\newline 4) 回防半径与触发阈值
& 参数组更新：$\Delta\tau_{\mathrm{param}}=1$；
\newline 若更新威胁阈值/先验：$\Delta\tau_{\mathrm{pred}}=1$
& 参数更新需冷却（如 $10$--$20\,s$）避免振荡
& InfoGap 高升；或 DefGap 中等且可通过“提高侦察+改配比”下降
& 侦察充分且 Gap 稳定下降可降到 L1/L0；
\newline 若参数不足以压缺口则升到 L3
\\ \hline
\textbf{L3}
& 组织同伦：有限重排/取消/插队/转产
& 1) Swap（相邻交换）\newline 2) Promote（前移 $1$--$2$ 位）
\newline 3) Insert（插入 L1 紧急项）
\newline 4) Cancel（取消非关键可回滚项）
\newline 5) Retarget（生产设施切换兵种）
& $\Delta\tau_{\mathrm{org}}$ 按交换/插队/取消次数累计；
\newline 若取消已投入较多的关键项：$\Delta\tau_{\mathrm{rev}}{+}{=}1$
& 禁止“取消$\to$重建”循环；
\newline 仅在 $\Delta J>\varepsilon$ 时落地（防抖）
& L1/L2 无法在短视界显著降低 Def/\allowbreak Power/\allowbreak Prod；
\newline 或威胁确认需要快速重排
& 连续 $30$--$60\,s$ 不再需要重排且 $\Delta\tau_{\mathrm{org}}$ 趋零，可降到 L2
\\ \hline
\textbf{L4}
& 强结构改写：布局/科技/兵种体系模板
& 模板库（建议固化 $6$--$12$ 个）：
\newline 核心内缩+外围防线；核心分散+分区隔离；
\newline 瓶颈化入口（墙/塔/地形利用）；前推哨点；
\newline 切兵种体系（转重装/转机动/转防空为主）；
\newline 停科技保生存/快科技反制；扩张迁移（主矿线迁移）
& $\Delta\tau_{\mathrm{pos}}$：
\newline 建筑移动数 $+$ 防线节点数；
\newline 策略栈切换：
\newline $\Delta\tau_{\mathrm{algo}}{+}{=}1$；
\newline 必要时 $\Delta\tau_{\mathrm{org}}$ 同步上升
& 必须满足“收益足够大”：
\newline $\Delta J$ 明显且可在 $1$--$2$ 分钟内体现；
\newline 最小驻留 $60$--$120\,s$
& Topology\allowbreak Breach\allowbreak Gap 超阈值；
\newline 或 DefGap 持续高且 L3 无解（结构性缺口）
& Topo/\allowbreak Def 持续低位且经济恢复，
\newline 逐步降到 L2/\allowbreak L3；
\newline 保留布局成果，不做反向重构
\\ \hline
\textbf{L5}
& 全局止损/硬重置：转场/一波/卖资产
& 模板库（建议 $3$--$6$ 个）：
\newline 全军一波拆电/拆矿；
\newline 转场（若可行）；
\newline 大规模卖出/取消止血；
\newline 最小可存活模式等待机会
& 每次止损：
\newline $\Delta\tau_{\mathrm{rev}}{+}{=}1$；
\newline 通常伴随
\newline $\Delta\tau_{\mathrm{algo}}$ 上升
& 必须严格触发与冷却
\newline （如 $120$--$180\,s$）；
\newline 防止“情绪化”反复 L5
& EconRunway\allowbreak Gap 极高且持续、Topo\allowbreak Gap 极高且不可修复；
\newline 或进入不可逆崩溃态
& 若扭转成功（Gap 大幅下降并稳定）可降到 L3/L2；否则维持到终局
\\ \hline
\end{longtable}
\endgroup

\begin{remark}[实施备注（可显著减少策略抖动与候选爆炸）]
\begin{enumerate}
  \item 候选集合必须强截断：每层只给几十个候选，否则在线决策不稳定；
  \item 必须做滞回/最小驻留时间：否则会出现 L2$\leftrightarrow$L3 抖动或“防空$\leftrightarrow$扩张”抖动；
  \item $\mathrm{TopoGap}$ 可先粗后细：先用简单路径与局部 DPS 近似跑通，再逐步增强几何/覆盖估计。
\end{enumerate}
\end{remark}

\section{工程场景移植：从 OpenRA 沙盘到特高压抗地磁暴与量子节点抗退相干}
\label{sec:engineering-migration-openra-to-grid-qpu}

\subsection{移植不变量：证书事件流、两层生成泛函与修复塔}
\label{sec:migration-invariants-event-functional-tower}

\begin{remark}[三根主轴（从“对局”到“运行/实验”的不变量）]
无论是电网节点还是量子节点，\textbf{生成泛函中心论 + 非交换算子词 + 同伦修复塔（L0--L5）+ 字典序正则}迁移后仍只依赖三根主轴：
\begin{enumerate}
  \item \textbf{证书化事件流}：所有决策必须能由可回放、可审计的日志/测量数据重建；
  \item \textbf{两层生成泛函}：第一层为绩效路径积分（用于评价/学习/复盘），第二层为缺口--烈度字典序目标（用于在线控制/修复）；
  \item \textbf{修复塔不是扇区枚举}：L0--L5 是“允许的变分族”逐层扩张；是否显式命名为“扇区”，属于表达层选择。
  工程上所谓“扇区切换”，可内化为\textbf{权重/硬约束/候选变分族的激活条件改变}（由风险与证书触发）。
\end{enumerate}
\end{remark}

\begin{definition}[通用证书事件流与路径积分（场景无关骨架）]
\label{def:migration-generic-eventstream-functional}
对任意场景 $\mathsf{S}$，令 $\gamma$ 表示一段可回放/可重放的运行或实验轨迹，定义证书事件流
\[
  \mathsf{Ev}_{\mathsf{S}}(\gamma)=\bigl(e_0,e_1,\dots,e_{N-1}\bigr).
\]
给定局部评价泛函 $\mathcal{L}_{\mathrm{perf}}^{(\mathsf{S})}:\mathsf{Ev}\to\RR$，定义绩效路径积分泛函
\[
  S_{\mathrm{perf}}^{(\mathsf{S})}[\gamma]
  :=
  \sum_{e\in \mathsf{Ev}_{\mathsf{S}}(\gamma)} \mathcal{L}_{\mathrm{perf}}^{(\mathsf{S})}(e).
\]
\end{definition}

\begin{definition}[通用缺口/烈度与字典序目标（场景无关骨架）]
\label{def:migration-generic-gap-tau-lex}
对场景 $\mathsf{S}$，令 $\mathbf{J}^{(\mathsf{S})}$ 表示结构缺口向量（主缺口优先的字典序），令 $\Delta\tau^{(\mathsf{S})}\in\NN^6$
  表示修复烈度增量（pred/org/pos/param/algo/rev 六维对齐）。
在线修复的字典序目标写为
\[
  \min_{\mathrm{lex}}\bigl(\mathbf{J}^{(\mathsf{S})},\ |\tau^{(\mathsf{S})}|_w\bigr),
\]
即先压主缺口，再在缺口同等可接受时最小化结构代价。
\end{definition}

\begin{proposition}[命题：扇区切换可内化为“泛函耦合改变”，而非手工枚举]
\label{prop:migration-sector-switch-as-functional-coupling}
在定义~\ref{def:migration-generic-eventstream-functional} 与定义~\ref{def:migration-generic-gap-tau-lex} 的骨架下，
若某场景的在线决策由“候选变分生成器 + 字典序比较”决定，
则所谓“扇区切换”可以等价表述为以下三类变化之一（或其组合）：
\begin{enumerate}
  \item 权重/阈值变化：$\mathcal{L}_{\mathrm{perf}}^{(\mathsf{S})}$、$\mathbf{J}^{(\mathsf{S})}$ 的分量优先级或阈值发生变化；
  \item 硬约束激活：$\mathcal{B}_{\mathrm{hard}}$（硬约束壁垒项）从不活跃变为活跃，或反之；
  \item 候选族激活：可用候选变分族 $\mathcal{U}_L$ 的允许层级发生变化（升层/降层）。
\end{enumerate}
因此，“是否显式枚举扇区标签”只影响表达与审计坐标，不改变由生成泛函诱导的最优候选选择。
\end{proposition}

\begin{Certificate}[证书：扇区标签是派生描述，不改变候选集与比较准则]
\label{cert:migration-sector-switch-as-functional-coupling}
证书登记：在线决策的输入为（候选集、比较准则）。
当“扇区切换”被实现为权重/硬约束/候选族的变化时，比较准则与候选集的变化已由上述三类变化显式给出；
而扇区标签若存在，仅是把这些变化以人可读方式命名与归档，不改变决策机制本身。
\end{Certificate}

\begin{proof}[证明（证书；谓词因果链）]
由证书~\ref{cert:migration-sector-switch-as-functional-coupling}：
若两种实现只是在标签层不同，而比较准则与候选集一致，则每一步字典序最优候选相同；
若比较准则与候选集发生变化，则变化已经由“权重/硬约束/候选族激活”显式描述；
因此扇区切换可内化为泛函耦合改变而无需手工枚举，命题成立。证毕。
\end{proof}

\subsection{场景 A：特高压输电节点抵御太阳风暴/地磁暴（GMD/GIC）}
\label{sec:migration-uhv-gmd}

\begin{remark}[场景背景与合规口径（证书化动机）]
太阳风暴引发的地磁扰动可在输电网络中诱发地磁感应电流（GIC），导致变压器半周饱和、谐波与无功需求上升，并可能触发电压失稳或保护误动。
该类机制与风险在 NERC 的综述材料中有系统阐述（例如 \cite{ref_nerc_gmd_report_2012,ref_nerc_gmd_planning_guide_2013}）。
在北美体系中，\textsf{TPL-007} 系列规划标准（如 \textsf{TPL-007-4}）以“评估--纠正行动计划--实施”的方式规定了 GMD 规划与缓解的合规路径 \cite{ref_nerc_tpl_007_4}。
因此，电网场景的“证书事件流”不仅用于工程复盘，也直接服务于合规审计。
\end{remark}

\begin{definition}[电网节点的证书事件流：空间天气--运行--保护--设备健康]
\label{def:grid-eventstream}
对电网节点场景 $\mathsf{S}=\mathsf{Grid}$，定义事件流证书 $\mathsf{Ev}_{\mathsf{Grid}}(\gamma)$ 至少包含以下类别事件（可按工程需要裁剪）：
\begin{enumerate}
  \item 空间天气/地磁输入：预警等级、地电场模型参数、区域缩放与参考波形（规划口径可参考 \textsf{TPL-007-4} 的定义域 \cite{ref_nerc_tpl_007_4}）；
  \item 运行事件：断路器/刀闸状态变化、线路潮流、母线电压、频率偏差、无功装置投切（SVC/STATCOM/电容电抗等）、分接头动作；
  \item 保护事件：继电保护告警、动作/误动记录、保护闭锁与复归；
  \item 设备健康事件：变压器油温/绕组热点、谐波含量、寿命损耗代理指标；
  \item GIC 相关测量/估计：中性点直流分量（若有监测）、或基于潮流/无功响应的等效估计（工程实施策略可参见 \cite{ref_doe_gmd_monitoring_strategies}）。
\end{enumerate}
\end{definition}

\begin{definition}[抗 GMD 的绩效路径积分泛函（电网版）]
\label{def:grid-performance-functional}
给定电网局部评价泛函 $\mathcal{L}_{\mathrm{perf}}^{(\mathsf{Grid})}$，定义
\[
  S_{\mathrm{perf}}^{(\mathsf{Grid})}[\gamma]
  :=
  \sum_{e\in \mathsf{Ev}_{\mathsf{Grid}}(\gamma)} \mathcal{L}_{\mathrm{perf}}^{(\mathsf{Grid})}(e),
  \qquad
  \mathcal{L}_{\mathrm{perf}}^{(\mathsf{Grid})}(e)\in\RR.
\]
典型的评价分量（可字典序或加权合成）包括：未供电负荷/供电质量、电压合格率与裕度、设备热裕度消耗与寿命损耗代理、谐波与保护误动风险、操作代价与恢复时间等。
\end{definition}

\begin{definition}[结构缺口向量（电网版，6 项最小集）]
\label{def:grid-gap-vector}
定义电网场景的结构缺口向量为
\[
  \begin{aligned}
    \mathbf{J}^{(\mathsf{Grid})}
    &:=
    \bigl(
      \mathrm{GICSatGap},\ \mathrm{QMarginGap},\ \mathrm{VoltStabGap},\ \mathrm{ProtCoordGap},\\
      &\qquad\mathrm{InfoGap}_{\mathsf{Grid}},\ \mathrm{TopoCascadeGap}
    \bigr)\in\RR_{\ge 0}^6,
  \end{aligned}
\]
其中各分量分别证书化：
\begin{enumerate}
  \item $\mathrm{GICSatGap}$：GIC/半周饱和与谐波风险（可由中性点直流、谐波与无功异常代理估计）\cite{ref_nerc_gmd_report_2012}；
  \item $\mathrm{QMarginGap}$：无功裕度缺口（饱和吸无功导致的电压下陷风险）\cite{ref_nerc_gmd_planning_guide_2013}；
  \item $\mathrm{VoltStabGap}$：电压稳定裕度缺口（关键母线电压与扰动敏感性）；
  \item $\mathrm{ProtCoordGap}$：保护协调缺口（谐波/异常电流导致误动风险上升；相关工程评估材料可参见 \cite{ref_iso_ne_gmd_2026_na_sow}）；
  \item $\mathrm{InfoGap}_{\mathsf{Grid}}$：态势不确定性缺口（空间天气输入与地电场模型不确定、关键监测缺失）；
  \item $\mathrm{TopoCascadeGap}$：拓扑级联缺口（网络拓扑下存在“低成本路径”导致电压/无功薄弱区级联失效）。
\end{enumerate}
并对 $\RR^6$ 采用字典序比较（主缺口优先级可由工程备案指定）。
\end{definition}

\begin{definition}[修复烈度（电网版，pred/org/pos/param/algo/rev 六维对齐）]
\label{def:grid-tau}
定义电网场景的修复烈度增量 $\Delta\tau^{(\mathsf{Grid})}\in\NN^6$ 的对齐解释：
\begin{enumerate}
  \item $\Delta\tau_{\mathrm{pred}}$：空间天气/GIC 风险先验与模型更新幅度（地电场模型切换、阈值提高等）；
  \item $\Delta\tau_{\mathrm{org}}$：调度组织重排强度（发电计划与无功资源调度重排、操作序列重排）；
  \item $\Delta\tau_{\mathrm{pos}}$：拓扑/结构变更强度（开关操作导致网络结构变化、分区运行、关键通道改构型；更长周期硬件改造亦可计入此维）；
  \item $\Delta\tau_{\mathrm{param}}$：控制参数调整强度（AVR/励磁限制、无功装置设定、分接头策略等）；
  \item $\Delta\tau_{\mathrm{algo}}$：算法/控制模式切换强度（常态 OPF $\to$ SCOPF/应急控制/防 GMD 模式等）；
  \item $\Delta\tau_{\mathrm{rev}}$：回退/止损动作强度（负荷切除、受控解列、设备保护性退出、黑启动/恢复路径执行等）。
\end{enumerate}
\end{definition}

\begin{definition}[修复塔 L0--L5（电网版：从常态控制到受控解列）]
\label{def:grid-repair-tower}
以候选变分族逐层扩张的方式定义电网修复塔（示意模板库口径）：
\begin{enumerate}
  \item L0：常态运行与常态控制（常规调压、常规无功调节）；
  \item L1：最小补救链（快速投切无功装置、优化分接头策略的最短操作链）；
  \item L2：profile/param 同伦（切换到更保守的电压/无功控制参数，提高监测频率以降低 $\mathrm{InfoGap}_{\mathsf{Grid}}$）；
  \item L3：组织同伦（调度重排 + 拓扑微调；仅在有限候选集合上做 $\varepsilon$-近最优交换以避免操作爆炸）；
  \item L4：强结构改写（分区运行、关键设备保护性退出、关键通道改构型；并对齐 \textsf{TPL-007-4} 的“纠正行动计划/实施”口径）\cite{ref_nerc_tpl_007_4}；
  \item L5：全局止损/退出（受控解列、负荷切除、黑启动与恢复路径执行）。
\end{enumerate}
各层动作普遍强非交换：同一组投切/分接头/拓扑操作若顺序不同，会导致不同的电压轨迹、保护风险与设备应力。
\end{definition}

\begin{lemma}[引理：电网操作词的强非交换性（存在性引理）]
\label{lem:grid-operator-word-noncommutative}
存在电网状态 $s$ 与两类操作算子 $a,b$（例如“无功投切”与“分接头动作”或“拓扑重构”与“降载”），使得
\[
  (a\circ b)(s)\neq (b\circ a)(s),
\]
即操作词在一般情况下不可交换。
\end{lemma}

\begin{Certificate}[证书：以电压轨迹/保护风险差异作见证]
\label{cert:grid-operator-word-noncommutative}
证书登记：选取一组使系统处于电压裕度较紧的状态 $s$，并选取两类会改变无功需求与电压支撑路径的操作 $a,b$。
由于分接头与无功投切（或拓扑重构与降载）对电压/潮流的影响具有路径依赖，
两种顺序将产生不同的中间状态，从而在有限时间尺度内诱导不同的电压轨迹与保护风险代理，作为 $(a\circ b)(s)\neq (b\circ a)(s)$ 的可观测见证。
\end{Certificate}

\begin{proof}[证明（证书；谓词因果链）]
由证书~\ref{cert:grid-operator-word-noncommutative}：
在电压裕度紧的状态下，操作 $a,b$ 改变的是不同的“支撑路径/约束活跃集”；
顺序不同导致中间状态不同，从而产生不同的电压轨迹与风险代理输出，因此两种复合不相等。
引理成立。证毕。
\end{proof}

\begin{remark}[三维格点对应（几何--材料--包络环境）]
电网抗 GMD 场景中的“三维格点”可直接对齐为：
\begin{enumerate}
  \item 几何形态格点：网络拓扑（开关状态、分区划分、关键通道选择）、站内母线接线方式；
  \item 材料格点：变压器抗饱和能力、GIC 监测/阻断装置、无功补偿设备能力；
  \item 包络环境格点：地电导模型（区域差异）、地磁扰动强度/频谱、土壤电阻率与纬度暴露。
\end{enumerate}
这些格点决定“候选变分族”的可行域与代价结构，因此自然落入 $\mathbf{J}^{(\mathsf{Grid})}$ 与 $\Delta\tau^{(\mathsf{Grid})}$ 的证书化坐标中。
\end{remark}

\subsection{场景 B：量子计算节点抗退相干（含“室温超导”作为材料格点目标语境）}
\label{sec:migration-qpu-decoherence}

\begin{remark}[现实边界与证书化动机（避免概念跳跃）]
截至目前，以公开可核验文献为准，“常压常温且可工程化的室温超导”并无被广泛接受的定论性实现；并且相关高调主张曾出现撤稿与严重争议（例如 \cite{ref_nature_2022_05294_9}）。
与此同时，主流超导量子计算平台通常依赖毫开尔文级环境与稀释制冷机体系；工程综述与国家科学院材料均明确了低温运行与制冷口径（例如 \cite{ref_aip_apr_superconducting_qubits,
  ref_nationalacademies_25196_ch13}）。
因此，本节采取“两层移植”口径：
\begin{enumerate}
  \item 现实层：以现有超导/封装互连工程为对象，建立抗退相干修复塔；
  \item 目标层：把“室温超导互连/屏蔽/器件”视为材料格点上的高阶候选，用以描述其如何改变缺口与候选变分族（并不预设其已实现）。
\end{enumerate}
\end{remark}

\begin{definition}[量子节点的证书事件流：性能--控制--环境--封装互连]
\label{def:qpu-eventstream}
对量子节点场景 $\mathsf{S}=\mathsf{QPU}$，定义事件流证书 $\mathsf{Ev}_{\mathsf{QPU}}(\gamma)$ 至少包含：
\begin{enumerate}
  \item 性能证书：$T_1/T_2$、门保真度、读出误差、串扰指标、泄漏率、噪声谱估计等（工程综述口径可参见 \cite{ref_aip_apr_superconducting_qubits}）；
  \item 控制证书：脉冲参数（频率/幅度/相位/波形）、偏置点、校准版本号、编译器/调度器版本；
  \item 环境证书：多级温度平台、振动、磁场/EMI 监测、辐照/宇宙线事件（若有）；
  \item 硬件互连/封装证书：线缆数量与热负载、封装/屏蔽结构、近端控制电子学集成方式；
  大规模低温工程的互连与热负载扩展瓶颈可参见综述 \cite{ref_springer_epjqt_2019_0072_0}；
  近端控制电子学的低温集成案例可参见 \cite{ref_nature_2025_09157_x}。
\end{enumerate}
\end{definition}

\begin{definition}[量子就绪度的绩效路径积分泛函（QPU 版）]
\label{def:qpu-performance-functional}
给定量子节点的局部评价泛函 $\mathcal{L}_{\mathrm{perf}}^{(\mathsf{QPU})}$，定义
\[
  S_{\mathrm{perf}}^{(\mathsf{QPU})}[\gamma]
  :=
  \sum_{e\in \mathsf{Ev}_{\mathsf{QPU}}(\gamma)} \mathcal{L}_{\mathrm{perf}}^{(\mathsf{QPU})}(e).
\]
其典型评价分量包括：任务级成功率与吞吐、可用 qubit 数与有效线路深度、门/读出误差与漂移速度、校准时间占比、掉线/重启频率、制冷负载与能耗、故障恢复时间等。
\end{definition}

\begin{definition}[结构缺口向量（QPU 版，6 项最小集）]
\label{def:qpu-gap-vector}
定义量子节点的结构缺口向量为
\[
  \begin{aligned}
    \mathbf{J}^{(\mathsf{QPU})}
    &:=
    \bigl(
      \mathrm{DecoherenceGap},\ \mathrm{ControlGap},\ \mathrm{CrosstalkLeakageGap},\\
      &\qquad\mathrm{ThermalGap},\ \mathrm{InfoGap}_{\mathsf{QPU}},\ \mathrm{PackagingTopoGap}
    \bigr)\in\RR_{\ge 0}^6,
  \end{aligned}
\]
其中各分量分别证书化：
\begin{enumerate}
  \item $\mathrm{DecoherenceGap}$：$T_1/T_2$ 无法满足门深度需求或出现异常衰减机制；
  \item $\mathrm{ControlGap}$：控制误差、漂移、脉冲失配导致保真度下降；
  \item $\mathrm{CrosstalkLeakageGap}$：串扰/泄漏在规模扩大时成为结构性障碍；
  \item $\mathrm{ThermalGap}$：热激发/热噪声导致激发态占据与读出噪声上升（低温必要性与制冷口径可参见 \cite{ref_nationalacademies_25196_ch13}）；
  \item $\mathrm{InfoGap}_{\mathsf{QPU}}$：噪声模型不确定性缺口（噪声谱/漂移机制未被识别导致修复无效或过度）；
  \item $\mathrm{PackagingTopoGap}$：封装/互连/屏蔽几何导致的结构性串扰或热负载瓶颈（扩展工程瓶颈可参见 \cite{ref_springer_epjqt_2019_0072_0}）。
\end{enumerate}
并对 $\RR^6$ 采用字典序比较（主缺口优先级可备案）。
\end{definition}

\begin{definition}[修复烈度（QPU 版，pred/org/pos/param/algo/rev 六维对齐）]
\label{def:qpu-tau}
定义量子节点场景的修复烈度增量 $\Delta\tau^{(\mathsf{QPU})}\in\NN^6$ 的对齐解释：
\begin{enumerate}
  \item $\Delta\tau_{\mathrm{pred}}$：噪声模型/环境扰动先验的切换与更新幅度（噪声谱模型切换等）；
  \item $\Delta\tau_{\mathrm{org}}$：校准与作业调度重排（校准插队、批量重排、测量计划改变）；
  \item $\Delta\tau_{\mathrm{pos}}$：封装/互连/屏蔽/线缆路由等几何变更（模块重排、屏蔽结构调整）；
  \item $\Delta\tau_{\mathrm{param}}$：脉冲/偏置/动态解耦序列参数调整强度；
  \item $\Delta\tau_{\mathrm{algo}}$：算法栈切换（编译映射策略、误差缓解/纠错方案切换）；
  \item $\Delta\tau_{\mathrm{rev}}$：回退/止损（回滚校准版本、降级使用 qubit 子集、维护/重新热循环等）。
\end{enumerate}
\end{definition}

\begin{definition}[修复塔 L0--L5（QPU 版：从脉冲微调到硬件级重构）]
\label{def:qpu-repair-tower}
以候选变分族逐层扩张的方式定义量子节点修复塔（示意模板库口径）：
\begin{enumerate}
  \item L0：小幅参数微调（脉冲幅相、偏置点微调），不触碰系统结构；
  \item L1：slicing 最小补救（短窗口内选择最小校准/补偿动作子集，例如关键门的快速重校准、读出链快速校正）；
  \item L2：profile/param 同伦（系统性调整脉冲族、动态解耦参数，提高噪声表征频率以降低 $\mathrm{InfoGap}_{\mathsf{QPU}}$）；
  \item L3：组织同伦（编译层 qubit mapping 与门调度的有限重排，以降低串扰；校准/测量计划的 $\varepsilon$-近最优交换）；
  \item L4：强结构改写（封装/互连/屏蔽/近端电子学集成策略调整或模块级替换；中大型系统扩展工程可参见 \cite{ref_springer_epjqt_2019_0072_0}，低温近端控制案例可参见
    \cite{ref_nature_2025_09157_x}）；
  \item L5：全局止损/退出（系统降级、长时间维护/重新热循环、模块替换与重验收）。
\end{enumerate}
\end{definition}

\begin{proposition}[命题：材料格点迁移（“室温超导”目标）属于扇区切换而非参数微调]
\label{prop:qpu-roomtemp-sc-as-sector-switch}
在量子节点场景中，若把“室温超导互连/屏蔽/器件”视为材料格点上的高阶候选，
则其工程含义更接近于“改变包络环境与结构约束”的扇区切换：
它可能降低 $\mathrm{ThermalGap}$ 与 $\mathrm{PackagingTopoGap}$ 的主导性，同时也可能引入新的缺口项（材料缺陷谱、界面态、磁通噪声等）。
因此该迁移应被证书化为 L4/L5 级别的结构改写候选，而不应被误表述为 L0--L2 的局部参数微调。
\end{proposition}

\begin{Certificate}[证书：约束结构改变 $\Rightarrow$ 候选族与缺口分解改变]
\label{cert:qpu-roomtemp-sc-as-sector-switch}
证书登记：材料格点迁移将改变可用制冷/封装/互连的约束集合与可达域，从而改变缺口分解（$\mathbf{J}^{(\mathsf{QPU})}$ 的主导项）与候选变分族（允许的 $L$ 层级）。
该变化属于“权重/硬约束/候选族激活”的组合变化（命题~\ref{prop:migration-sector-switch-as-functional-coupling}），因此应按扇区切换备案。
\end{Certificate}

\begin{proof}[证明（证书；谓词因果链）]
由证书~\ref{cert:qpu-roomtemp-sc-as-sector-switch}：
材料格点迁移改变的是约束结构与可达域，而不仅是数值参数；
因此它会引起缺口分解与候选族激活的改变；
由命题~\ref{prop:migration-sector-switch-as-functional-coupling} 的等价描述，可将其归类为扇区切换，命题成立。证毕。
\end{proof}

\subsection{统一模板：把“扇区切换”内化为权重/硬约束/候选族激活}
\label{sec:migration-unified-template}

\begin{theorem}[公理降级：两类工程场景共享同一“证书--泛函--修复塔”闭环接口]
\label{thm:migration-shared-closed-loop-interface}
若某工程系统（电网节点或量子节点）能够：
\begin{enumerate}
  \item 抽取证书事件流并形成可复算证书快照序列；
  \item 在证书快照上计算主缺口向量 $\mathbf{J}$ 与烈度增量 $\Delta\tau$；
  \item 为每个层级 $L$ 生成有限候选集合，并以字典序规则选择候选；
  \item 以防抖/滞回控制层级升降。
\end{enumerate}
则其在线决策过程可复用 OpenRA 章“可复算一致性 + 候选有限性 + 抖动抑制”的证明结构，
从而形成一个可审计的“缺口优先 + 低层优先 + 必要时结构改写”的闭环修复系统。
\end{theorem}

\begin{Certificate}[证书：由三个已证结论拼接闭环接口]
\label{cert:migration-shared-closed-loop-interface}
证书登记以下拼接链：
\begin{enumerate}
  \item 证书快照一致 $\Rightarrow$ 指标与决策可复算一致（参照定理~\ref{thm:openra-snapshot-consistency-reproducible}）；
  \item 模板库有限 + 操作步长有界 $\Rightarrow$ 候选集有限（参照定理~\ref{thm:openra-generator-finiteness}）；
  \item 触发谓词采用持续满足与滞回 $\Rightarrow$ 抑制抖动（参照定理~\ref{thm:openra-level-switch-stability}）。
\end{enumerate}
将三者合取即可得到工程场景下的闭环可运行与可审计接口。
\end{Certificate}

\begin{proof}[证明（证书；谓词因果链）]
由证书~\ref{cert:migration-shared-closed-loop-interface}：
第一条给出“可复算一致性”；第二条给出“候选枚举可执行”；第三条给出“层级切换不抖动”；
三者合取即得到闭环系统可运行、可复算、可审计。
定理成立。证毕。
\end{proof}

\subsection{最小落地建议：先跑证书与 L0--L2，再逐步放开 L3--L5}
\label{sec:migration-minimal-roadmap}

\begin{remark}[最小闭环（两场景通用）]
要从概念快速走向可运行原型，建议采用同一条最小路线（两场景通用）：
\begin{enumerate}
  \item 先把 $\mathsf{Ev}_{\mathsf{S}}$ 与 $\mathbf{J}^{(\mathsf{S})}$ 跑起来（哪怕缺口指标很粗糙）；
  \item 只实现 L0--L2（软件/策略层即可），确保“缺口优先 + 最小代价”的决策可复算；
  \item 再按“低层不可修复证书”逐步放开 L3（有限重排）$\to$ L4（结构改写）$\to$ L5（全局止损），并配置最小驻留与冷却。
\end{enumerate}
该路线与电网合规的“评估$\to$纠正行动计划$\to$实施”思路（\textsf{TPL-007-4}）以及量子系统的“表征$\to$校准$\to$映射$\to$封装迭代”工程节奏在结构上对齐。
\end{remark}

\section{运行期在线控制移植（MVP）：多速率 tick、Gap 证书与修复塔 L0--L5}
\label{sec:online-control-mvp}

\subsection{统一接口：多速率 tick 仍可归约为“证书快照 $\to$ 候选 $\to$ 字典序选择”}
\label{sec:online-control-unified-interface}

\begin{remark}[定位（与规划期离线评估的差异）]
规划期离线评估更关注“长期可达性与纠正行动计划”；而运行期在线控制更关注：
\textbf{在强时限与强约束下，每个 tick 都给出可执行且可审计的动作}，
并通过“缺口优先 + 低层优先 + 必要时结构改写”的分层机制避免误动与抖动。
\end{remark}

\begin{definition}[多速率 tick（运行期常态）]
\label{def:online-multi-rate-tick}
对运行期系统，令 $\mathbb{T}$ 为一组 tick 类型集合（多速率）：
\[
  \mathbb{T}=\{\mathsf{fast},\ \mathsf{op},\ \mathsf{topo},\ \mathsf{emg}\},
\]
分别表示：快速监测 tick、操作/设定值 tick、拓扑/调度 tick、紧急事件触发 tick。
对每种 tick 类型 $\tau\in\mathbb{T}$，系统在触发时刻 $t$ 都必须输出一份证书快照 $\mathsf{Snap}_{t}^{(\tau)}$，
并在该快照上执行同一决策接口（候选生成 + 字典序选择）。
\end{definition}

\begin{definition}[在线候选变分的最小签名（含 time-to-effect）]
\label{def:online-candidate-signature}
称一个在线候选变分 $u$ 的最小签名为
\[
  \mathsf{Cand}(u)
  :=
  \bigl(
    \mathsf{word}(u),\ \mathsf{feas}(u),\ \mathsf{tte}(u),\ \widehat{\Delta\mathbf{J}}(u),\ \Delta\tau(u),\ \mathsf{fail}
      (u)
  \bigr),
\]
其中：
\begin{enumerate}
  \item $\mathsf{word}(u)$ 为动作词（非交换序列或模板）；
  \item $\mathsf{feas}(u)$ 为可行性谓词（硬约束/前置条件）；
  \item $\mathsf{tte}(u)$ 为 time-to-effect（动作生效时间尺度）；
  \item $\widehat{\Delta\mathbf{J}}(u)$ 为对主缺口向量变化的短视界估计；
  \item $\Delta\tau(u)$ 为烈度增量（pred/org/pos/param/algo/rev 六维）；
  \item $\mathsf{fail}(u)$ 为失败条件/禁止条件（例如冷却未到、风险超阈等）。
\end{enumerate}
该签名保证：候选可执行、可评估、可审计。
\end{definition}

\begin{theorem}[公理降级：多速率 tick 可统一为同一在线决策接口]
\label{thm:online-multi-rate-unified-interface}
在定义~\ref{def:online-multi-rate-tick} 的多速率 tick 口径下，
若对每种 tick 类型 $\tau\in\mathbb{T}$：
\begin{enumerate}
  \item 可从监测/日志抽取证书快照 $\mathsf{Snap}_{t}^{(\tau)}$；
  \item 可生成有限候选集合 $\mathcal{U}^{(\tau)}(\mathsf{Snap}_{t}^{(\tau)})$，并为每个候选给出最小签名（定义~\ref{def:online-candidate-signature}
    ）；
  \item 比较准则统一为字典序目标
  \[
    \min_{\mathrm{lex}}\bigl(\widehat{\Delta\mathbf{J}}(u),\ \Delta|\tau(u)|_w\bigr),
  \]
  并配置防抖/滞回与冷却（对齐 OpenRA 章的定义与引理）。
\end{enumerate}
则无论 tick 频率如何变化，运行期控制均可归约为同一“证书快照 $\to$ 候选 $\to$ 字典序选择”的接口，
从而保持决策逻辑一致性与审计一致性。
\end{theorem}

\begin{Certificate}[证书：把 tick 类型差异吸收到“候选族 + time-to-effect”中]
\label{cert:online-multi-rate-unified-interface}
证书登记如下因果链：
\[
\begin{aligned}
  C_1 &: \tau\ \text{只影响}\ \mathsf{Snap}_{t}^{(\tau)}\ \text{的可获得字段与}\ \mathcal{U}^{(\tau)}\ \text{的可用候选族},\\
  C_2 &: \mathsf{tte}(u)\ \text{将“动作生效时间尺度”显式写入候选签名，从而允许跨 tick 统一比较},\\
  C_3 &: \text{统一字典序准则 + 防抖/滞回/冷却}\ \Rightarrow\ \text{决策接口一致且可审计}.
\end{aligned}
\]
\end{Certificate}

\begin{proof}[证明（证书；谓词因果链）]
由证书~\ref{cert:online-multi-rate-unified-interface}：
tick 的多速率只改变“此刻可观测什么、可执行什么”，对应到接口即为快照字段与候选族的差异；
而候选签名显式包含 time-to-effect，使不同时间尺度动作可在统一准则下比较；
统一字典序准则与防抖/滞回/冷却保证选择过程一致且不抖动，故可归约为同一接口。
定理成立。证毕。
\end{proof}

\begin{remark}[统一在线控制器接口（最小 API 口径）]
两个工程场景（电网/量子）均可复用同一套“证书 $\to$ 候选 $\to$ 选择 $\to$ 执行 $\to$ 写证书”的最小接口。
一个可直接用于实现文档/任务书的接口口径示例如下：
\begin{verbatim}
snapshot = get_certificate_snapshot()
cands    = gen_candidates(level, snapshot)
score    = eval_candidate(cand, snapshot)   # -> (ΔJ, Δ|τ|_w) lex
actuate(cand)                               # -> append EventStream
\end{verbatim}
其中每个候选需携带可行性、time-to-effect、烈度增量与失败条件（定义~\ref{def:online-candidate-signature}），并配置冷却/驻留与安全联锁。
\end{remark}

\section{运行期在线控制 MVP：特高压输电关键节点抗地磁暴（GMD/GIC）}
\label{sec:online-grid-gmd-mvp}

\subsection{运行期在线控制的边界与护栏（建议器口径）}
\label{sec:online-grid-boundary-guards}

\begin{remark}[定位：证书化建议器 + 有限候选集选择器]
本节的运行期在线控制器定位为“\textbf{证书化建议器 + 有限候选集选择器}”：
\begin{enumerate}
  \item 它只在已批准的运行期 Operating Plan/Procedures 与站内联锁约束下枚举动作词（并输出建议与证据包），不在运行期做自由组合生成；
  \item 对高风险动作（L3--L5：拓扑/分区/负荷切除），候选必须来自经审查的预案模板库，并配置证书门槛、最小驻留时间与冷却时间；
  \item 外部态势输入（SWPC 的 Watch/Warning/Alert 与 NOAA G1--G5 等级）必须证书化写入事件流，以保证“告警 $\to$ 权重/阈值切换 $\to$ 动作选择”的链路可复算可审计。
\end{enumerate}
其中第三点可参照定义~\ref{def:online-grid-threat-certificates}；第二点的必要性可由运行期安全命题~\ref{prop:online-l3plus-template-only}
  的“组合爆炸与不可预期耦合”证书给出。
同时，电网运行期合规常要求维护并执行 GMD Operating Plan（例如 \cite{ref_nerc_eop_010_1}），因此该定位也为合规审计提供统一接口。
\end{remark}

\subsection{证书事件流与 tick 设定（电网运行期）}
\label{sec:online-grid-event-tick}

\begin{remark}[证书事件流最小集（可回放、可审计）]
电网运行期的在线控制必须以可审计事件流为底座，最小建议覆盖：
\begin{enumerate}
  \item 空间天气态势：\textsf{NOAA SWPC} 的 Alert/Watch/Warning（区分“正在发生/可能发生”的告警口径）\cite{ref_noaa_swpc_alerts}；
  \item 电气运行态：PMU/SCADA（母线电压、频率、线路潮流、无功出力、HVDC 交换功率、开关状态）；
  \item 设备响应/风险证书：变压器告警、热点温度、谐波/THD、保护装置告警；
  \item GIC/GMD 影响代理：中性点电流（若有）、谐波与无功吸收异常代理；GIC 导致半周饱和、谐波与无功吸收上升的风险机制可参见 NERC 材料 \cite{ref_nerc_gmd_report_2012}。
\end{enumerate}
同时，运行期合规常要求维护并执行 GMD Operating Plan（例如 \cite{ref_nerc_eop_010_1}），因此“证书事件流 + 可复算”也是合规与事后分析的统一接口。
\end{remark}

\begin{definition}[外部威胁输入证书：Watch/Warning/Alert 与 NOAA 地磁暴等级（G1--G5）]
\label{def:online-grid-threat-certificates}
在运行期在线控制中，空间天气输入可被视为“对手策略信号”的证书化外生扰动。
建议至少对以下信息证书化并写入事件流：
\begin{enumerate}
  \item \textsf{Watch/Warning/Alert}：用于区分“预期/正在逼近/正在发生”的运行态势口径（可参见 NWS 对 Space Weather Watches, Warnings and Alerts 的说明
    \cite{ref_nws_space_ww}）；
  \item \textsf{订阅/推送通道}：运行期系统可通过 SWPC 的订阅服务将告警输入作为可追溯证书接入（见 \cite{ref_swpc_subscription_services}）；
  \item \textsf{G 级（G1--G5）}：作为阈值级别的统一枚举口径，用于驱动 $\Delta\tau_{\mathrm{pred}}$ 的“先验扭曲”与缺口权重调度；其影响示例性解释可参见 NOAA Scales
    \cite{ref_swpc_noaa_scales_explanation}。
\end{enumerate}
该定义的工程含义是：外部威胁不再是口头叙述，而是进入“证书事件流 $\Rightarrow$ 可复算决策”的闭环输入。
\end{definition}

\begin{remark}[证书事件流的工程落地：append-only + schema 版本化]
建议将运行期在线控制的证书事件流实现为 append-only 日志（并配合不可篡改存储策略，例如 WORM），并对 schema 做版本化备案。
一个可落地的事件记录示意如下（仅强调字段口径，不要求硬编码格式）：
\begin{verbatim}
{
  "ts_utc": "2026-01-19T03:21:10Z",
  "tick_class": "FAST|OP|TOPO|EMERGENCY",
  "event_type": [
    "SWPC_WATCH", "SWPC_WARNING", "SWPC_ALERT",
    "PMU_SNAPSHOT", "SCADA_SWITCH", "FACTS_SETPOINT",
    "RELAY_ALARM", "RELAY_TRIP", "HARMONICS_SNAPSHOT",
    "TX_HEALTH", "LOAD_SHED", "ISLANDING", "MODEL_UPDATE"
  ],
  "asset_id": "SUBSTATION_X/TRANSFORMER_T1",
  "payload": { "...": "..." },
  "certificate": {
    "source": "PMU|SCADA|SWPC|EMS",
    "signature": "optional",
    "version": "schema_v1"
  }
}
\end{verbatim}
该日志结构与定义~\ref{def:openra-replay-event-stream} 的“可回放证书载体”同构，从而使运行期在线控制与离线复盘共享同一审计接口。
\end{remark}

\begin{definition}[电网运行期多速率 tick（示意）]
\label{def:online-grid-ticks}
将电网运行期 tick 设为多速率但统一接口：
\begin{enumerate}
  \item Fast tick（$1$--$2$s）：PMU 驱动缺口更新，用于 L0--L2 的快速候选评估；
  \item Op tick（$10$--$30$s）：无功装置投切/设定值更新（SVC/STATCOM、电容电抗、部分策略切换）；
  \item Topo/Dispatch tick（$1$--$5$min）：拓扑切换、发电/无功重调度、跨区潮流限额调整；
  \item Emergency tick（事件触发）：负荷切除、受控解列/孤岛等 L5 动作。
\end{enumerate}
每类 tick 均输出证书快照并调用统一接口（定理~\ref{thm:online-multi-rate-unified-interface}）。
\end{definition}

\begin{definition}[电网证书快照（MVP 字段口径）]
\label{def:online-grid-snapshot-fields}
对每个运行期 tick 类型 $\tau\in\mathbb{T}$ 与触发时刻 $t$，证书快照 $\mathsf{Snap}_{t}^{(\tau)}$ 至少应包含以下字段（并以 schema 版本化备案）：
\begin{enumerate}
  \item \code{t_utc} 与 \code{tick_class}（$\in\mathbb{T}$），用于对齐多速率决策与 time-to-effect；
  \item \code{threat}：外部威胁证书（Watch/Warning/Alert、G 级别、有效期等；见定义~\ref{def:online-grid-threat-certificates}）；
  \item \code{grid_state}：关键母线 $V$、关键廊道潮流、无功备用、频率/振荡代理等；
  \item \code{asset_state}：关键变压器负载率与健康/告警、谐波/THD、（若有）中性点直流代理；
  \item \code{control_state}：无功装置投切状态、FACTS/HVDC 设定值、LTC/电压计划策略状态；
  \item \code{gaps}：主缺口向量 $\mathbf{J}$ 的六个分量（见第~\ref{sec:online-grid-gap-table} 节）；
  \item \code{tau}：烈度计数器与最近增量窗口（见定义~\ref{def:online-grid-tau}）；
  \item \code{level}：当前修复层级 $L\in\{0,\dots,5\}$ 与驻留/冷却计时器（用于抑制抖动与误触发）。
\end{enumerate}
该定义的工程含义是：在线决策在“快照字段”层面可复算，从而保证事后审计可重放与可归因。
\end{definition}

\subsection{Gap 指标计算表（电网运行期 6 核心）}
\label{sec:online-grid-gap-table}

\begin{remark}[说明（电网口径的“缺口”）]
电网场景的主缺口不以“收益最大化”为核心，而以“避免逼近电压崩溃、误动与设备损伤等不可逆边界”为核心；
并可按字典序组织主缺口（拓扑脆弱/电压稳定/无功裕度/饱和代理/保护风险/不确定性）。
该口径与 NERC 对 GMD 风险叙事（设备损伤 + 无功损失导致电压不稳/崩溃）一致；
并可用基准事件描述把“半周饱和 $\Rightarrow$ 谐波 + 无功吸收”这一机制固定为可审计参考口径 \cite{ref_nerc_gmd_report_2012,ref_nerc_benchmark_gmd_event}。
\end{remark}

\begingroup
\scriptsize
\setlength{\tabcolsep}{2pt}
\renewcommand{\arraystretch}{1.2}
\setlength{\LTleft}{0pt}
\setlength{\LTright}{0pt}
\begin{longtable}{|
  >{\raggedright\arraybackslash}p{0.11\linewidth}|
  >{\raggedright\arraybackslash}p{0.15\linewidth}|
  >{\raggedright\arraybackslash}p{0.17\linewidth}|
  >{\raggedright\arraybackslash}p{0.11\linewidth}|
  >{\raggedright\arraybackslash}p{0.20\linewidth}|
  >{\raggedright\arraybackslash}p{0.11\linewidth}|
  >{\raggedright\arraybackslash}p{0.08\linewidth}|}
\hline
\textbf{Gap 名称} & \textbf{工程含义} & \textbf{在线输入（最小集）} & \textbf{频率/滑窗} & \textbf{计算（可落地近似）} & \textbf{阈值建议（起步）} &
  \textbf{触发层级} \\
\hline
\endfirsthead
\hline
\textbf{Gap 名称} & \textbf{工程含义} & \textbf{在线输入（最小集）} & \textbf{频率/滑窗} & \textbf{计算（可落地近似）} & \textbf{阈值建议（起步）} &
  \textbf{触发层级} \\
\hline
\endhead
\textbf{GIC\allowbreak Saturation\allowbreak Gap}
& 变压器半周饱和风险（谐波 + 无功吸收异常）
& 中性点电流（若有）、THD/偶次谐波、变压器无功吸收偏离、热点温度趋势
& $1$--$2$s 更新；$W=30$--$120$s
& 构造饱和指数：$S=a\cdot \mathrm{Norm}(I_{\mathrm{dc}})+b\cdot\mathrm{Norm}(\mathrm{THD})+c\cdot\mathrm{Norm}(\Delta
  Q_{\mathrm{tx}})+d\cdot\mathrm{Norm}(\dot T)$；
\newline Gap $:=\max(0,S-S_{\mathrm{crit}})$
& $S_{\mathrm{crit}}$ 可先取“健康日”99\% 分位并调参
& L1/L2
\\ \hline
\textbf{Var\allowbreak Margin\allowbreak Gap}
& 无功裕度缺口（电压支撑资源不足）
& 无功备用、SVC/STATCOM 余量、机组无功裕度、$\widehat{\Delta Q}_{\mathrm{GIC}}$ 估计
& $2$--$5$s
& $Q_{\mathrm{margin}}:=Q_{\mathrm{reserve}}-\widehat{\Delta Q}_{\mathrm{GIC}}$；
\newline Gap $:=\max(0,Q_{\min}-Q_{\mathrm{margin}})$
& $Q_{\min}$ 可由关键母线裕度倒推或经验阈值设定
& L1/L3
\\ \hline
\textbf{Voltage\allowbreak Stability\allowbreak Gap}
& 电压稳定缺口（逼近崩溃边界）
& 关键母线 $V$、$\dot V$、V-Q 灵敏度代理、无功装置饱和信息
& $1$--$2$s；$W=10$--$60$s
& Gap $:=\max(0,V_{\mathrm{target}}-V)+k\max(0,-\dot V-\dot V_{\min})$（趋势项提前预警）
& $V_{\mathrm{target}}$ 常取 $0.95$ pu；趋势阈值随系统惯性调参
& L2/L3/L4
\\ \hline
\textbf{Protection\allowbreak Misop\allowbreak Gap}
& 保护误动/协调破坏风险（谐波/异常电流）
& 保护告警、谐波谱、负序/零序异常、断路器动作统计
& $1$--$5$s
& Gap $:=$ 告警计数加权 + “接近动作门槛”的裕度反比（保护整定裕度作代理）
& 用保护整定裕度与历史误动样本定阈
& L2/L4
\\ \hline
\textbf{Situational\allowbreak Uncertainty\allowbreak Gap}
& 态势不确定性（告警 vs 本地观测偏差；监测缺失）
& SWPC 告警、站点磁测/地电场代理、监测覆盖率
& $10$--$60$s
& Gap $:=a\cdot\mathrm{Var}(\widehat{E}_g)$
\newline $+b(1-\text{监测覆盖})$
\newline $+c\cdot \mathrm{Mismatch}(\text{SWPC},\text{Loc})$；
\newline Alert 触发时提高 $a$
& 监测覆盖 $<0.8$ 时 gap 快速上升（起步）
& L2
\\ \hline
\textbf{Topology\allowbreak Fragility\allowbreak Gap}
& 拓扑结构脆弱：存在“低成本级联”路径
& 开关拓扑、关键廊道潮流、离线敏感度表（预计算）、关键变压器负载率
& $10$--$60$s
& Gap $:=\sum_{k\in K}\max(0,\theta-\mathrm{FragilityScore}(k))$；
\newline $\mathrm{FragilityScore}$ 可用“预计算敏感度 $\times$ 当前负载因子 $\times$ 无功/电压紧约束代理”组合
& $\theta$ 可由离线 N-1/N-k + GMD 敏感度排名阈值起步
& L4（严重时 L5）
\\ \hline
\end{longtable}
\endgroup

\begin{remark}[符号约定：$\mathrm{Norm}(\cdot)$ 与 $\mathrm{Mismatch}(\cdot)$（电网在线口径）]
为便于工程实现与审计复算，本节表格中的两个辅助算子建议按如下口径落地：
\begin{enumerate}
  \item $\mathrm{Norm}(\cdot)$：对不同量纲的观测量做归一化，可用“健康运行日”的分位数或均值/方差作归一化基准，并将基准版本化备案；
\item $\mathrm{Mismatch}(\text{SWPC},\text{Loc})$：刻画外部告警级别与本地观测/估计之间的一致性偏差；
  例如“外部宣告强但本地观测弱或缺失”应提升不确定性缺口并优先触发 L2 的监测增强与保守化策略。
\end{enumerate}
该口径的目的不是追求精确物理建模，而是把“告警--观测一致性”证书化，以抑制误判与策略抖动。
\end{remark}

\subsection{修复烈度与 L0--L5 候选变分生成表（电网运行期）}
\label{sec:online-grid-tau-and-tower}

\begin{definition}[电网运行期烈度六维（与 OpenRA 对齐）]
\label{def:online-grid-tau}
电网在线控制的 $\Delta\tau\in\NN^6$ 与 OpenRA 六维对齐如下：
\[
  \Delta\tau
  =
  (\Delta\tau_{\mathrm{pred}},\Delta\tau_{\mathrm{org}},\Delta\tau_{\mathrm{pos}},\Delta\tau_{\mathrm{param}},
    \Delta\tau_{\mathrm{algo}},\Delta\tau_{\mathrm{rev}}),
\]
分别表示：空间天气/模型先验切换、调度/操作序列重排、拓扑/分区变化、设定值/参数修改、控制模式切换、强止损动作（负荷切除/受控解列）次数。
在线控制建议取 $w_{\mathrm{rev}},w_{\mathrm{pos}}$ 显著更大，以体现“少切负荷、少大改拓扑”的结构代价。
\end{definition}

\begingroup
\scriptsize
\setlength{\tabcolsep}{2pt}
\renewcommand{\arraystretch}{1.2}
\setlength{\LTleft}{0pt}
\setlength{\LTright}{0pt}
\begin{longtable}{|
  >{\raggedright\arraybackslash}p{0.06\linewidth}|
  >{\raggedright\arraybackslash}p{0.15\linewidth}|
  >{\raggedright\arraybackslash}p{0.22\linewidth}|
  >{\raggedright\arraybackslash}p{0.13\linewidth}|
  >{\raggedright\arraybackslash}p{0.15\linewidth}|
  >{\raggedright\arraybackslash}p{0.13\linewidth}|
  >{\raggedright\arraybackslash}p{0.09\linewidth}|}
\hline
\textbf{层级} & \textbf{允许变分类型} & \textbf{候选生成器（最小模板库）} & \textbf{$\Delta\tau$ 计数} & \textbf{强非交换约束} & \textbf{升级触发（典型）} &
  \textbf{降级条件} \\
\hline
\endfirsthead
\hline
\textbf{层级} & \textbf{允许变分类型} & \textbf{候选生成器（最小模板库）} & \textbf{$\Delta\tau$ 计数} & \textbf{强非交换约束} & \textbf{升级触发（典型）} &
  \textbf{降级条件} \\
\hline
\endhead
\textbf{L0}
& 常态闭环；不启用 GMD 特殊策略
& 保持当前设定；仅微调电压参考/告警阈值（小步）
& 主要为 $\Delta\tau_{\mathrm{param}}$
& 禁止触发拓扑/负荷动作
& 任一硬缺口（电压/无功/饱和）上升
& 全部 Gap 低且平稳
\\ \hline
\textbf{L1}
& 最小补救：快速无功支撑与局部抑制
& 1) 投切电容/电抗\newline 2) SVC/STATCOM 设定值提升\newline 3) HVDC/FACTS 小幅调节以减轻关键廊道
& $\Delta\tau_{\mathrm{param}}$（投切/设定值）\newline 可伴随 $\Delta\tau_{\mathrm{org}}$（插队）
& \textbf{先保无功裕度再做潮流转移}（多数系统更安全）
& Var\allowbreak Margin\allowbreak Gap/\allowbreak GIC\allowbreak Saturation\allowbreak Gap 上升
\newline 且可用快速手段压下
& 连续 $T$ 秒缺口回落且无新告警
\\ \hline
\textbf{L2}
& 参数/策略同伦：运行保守化 + 监测增强
& 1) 提高电压计划目标\newline 2) 限制/冻结部分 LTC 策略（防错误调压）\newline 3) 提高监测频率与告警权重\newline 4) 启用 GMD Operating Plan 的参数集
  \cite{ref_nerc_eop_010_1}
& $\Delta\tau_{\mathrm{param}}+\Delta\tau_{\mathrm{algo}}$
& 参数集切换需冷却（避免来回振荡）
& Voltage\allowbreak Stability\allowbreak Gap
\newline 或 Situational\allowbreak Uncertainty\allowbreak Gap 主导
& 缺口稳定低位且态势解除
\\ \hline
\textbf{L3}
& 组织同伦：有限候选重调度/拓扑微调
& 1) 无功/发电重调度（提升关键区无功备用）\newline 2) 降低关键廊道潮流上限\newline 3) 有限开关操作集合（预案库）
& $\Delta\tau_{\mathrm{org}},\Delta\tau_{\mathrm{pos}}$
& \textbf{先释放无功/降载，再做拓扑切换}（常见安全顺序）
& L1/L2 无法显著降低缺口或保护风险升高
& $T$ 分钟缺口持续低位可回到 L2
\\ \hline
\textbf{L4}
& 强结构改写：分区/设备退出运行/结构性规避
& 1) 分区运行（sectionalize）\newline 2) 退出高风险变压器/设备\newline 3) 预案级结构改写开关序列（更大范围）
& $\Delta\tau_{\mathrm{pos}}+\Delta\tau_{\mathrm{algo}}$（高）
& 必须满足“显著降低 Topology\allowbreak Fragility/\allowbreak VoltStab”的证书门槛；
\newline 最小驻留
& Topology\allowbreak Fragility\allowbreak Gap 主导
\newline 或结构性无解
& 态势解除后逐步回切（不反向重构）
\\ \hline
\textbf{L5}
& 全局止损：负荷切除/受控解列/恢复路径
& 1) 分级负荷切除\newline 2) 受控孤岛\newline 3) 恢复/黑启动预案触发
& $\Delta\tau_{\mathrm{rev}}{+}{=}1$（强止损）
& 必须严格触发 + 冷却（防误触）
& Voltage\allowbreak Stability\allowbreak Gap 爆表
\newline 或保护连锁已开始
& 恢复后回到 L3/L2
\\ \hline
\end{longtable}
\endgroup

\subsection{在线 $\widehat{\Delta\mathbf{J}}$ 评估器的最小实现（电网版）}
\label{sec:online-grid-evaluator}

\begin{definition}[候选评估器（电网运行期；有限候选集）]
\label{def:online-grid-eval}
对有限候选集合 $\mathcal{U}_L$，定义在线评估器输出
\[
  \mathsf{Eval}(u;\mathsf{Snap})
  :=
  \bigl(
    \widehat{\Delta\mathbf{J}}(u),\ \Delta|\tau(u)|_w,\ \mathsf{tte}(u),\ \mathsf{feas}(u)
  \bigr),
\]
其中 $\widehat{\Delta\mathbf{J}}(u)$ 为短视界缺口变化估计，$\Delta|\tau(u)|_w$ 为结构代价增量，
$\mathsf{tte}(u)$ 与 $\mathsf{feas}(u)$ 为候选签名的 time-to-effect 与可行性谓词（定义~\ref{def:online-candidate-signature}）。
\end{definition}

\begin{remark}[实现建议：设定值类用灵敏度近似；拓扑/预案类用预计算表 + 在线校验]
为保证运行期可用性（强时限）与审计一致性（可复算），建议采用“分层评估”：
\begin{enumerate}
  \item L1/L2 的设定值/投切类候选：用局部线性化灵敏度或离线标定近似其对关键母线电压与无功裕度的影响，
  例如 $\Delta V \approx H_{VQ}\Delta Q$，并据此快速估计 $\widehat{\Delta\mathrm{VarMarginGap}}$ 与
    $\widehat{\Delta\mathrm{VoltageStabilityGap}}$；
  同时可将 $\widehat{\Delta\mathrm{GICSaturationGap}}$ 通过 $\Delta Q_{\mathrm{tx}}$ 与 THD 代理项映射到饱和指数的变化。
  \item L3/L4 的调度/拓扑类候选：对每个预案模板离线准备“场景表/敏感度表”（例如对 TopologyFragility 的方向性影响与安全约束边界），
  运行期只做表查询给出 $\widehat{\Delta\mathbf{J}}$，并再做一次硬约束可行性校验（设备可用、联锁约束、冷却/驻留满足）。
\end{enumerate}
该建议与本节“建议器口径”一致：在线不做自由组合与全局求解，只做\textbf{有限候选}的可执行比较。
\end{remark}

\subsection{tick 级在线决策规则（电网版）}
\label{sec:online-grid-policy}

\begin{definition}[电网 tick 级在线决策（字典序 + 显著改进 + 升降级）]
\label{def:online-grid-policy}
在每个决策 tick：
\begin{enumerate}
  \item （周期建议）可将决策周期取为 $10$--$30$s 一次；其中 Fast tick 指标可在 $1$--$2$s 内刷新并写入证书快照，用于候选评估的快速更新；
  \item 计算主缺口向量（建议字典序）：
  \[
    \begin{aligned}
      \mathbf{J}
      &:=
      \bigl(
        \mathrm{TopologyFragilityGap},\mathrm{VoltageStabilityGap},\mathrm{VarMarginGap},\\
        &\qquad \mathrm{GICSaturationGap},\mathrm{ProtectionMisopGap},\mathrm{SituationalUncertaintyGap}
      \bigr).
    \end{aligned}
  \]
	  \item 在当前层级 $L$ 生成有限候选集合 $\mathcal{U}_L$，并为每个候选附带 $\mathsf{tte}$ 与可行性；
	  \item 估计候选评分 $\bigl(\widehat{\Delta\mathbf{J}}(u),\Delta|\tau(u)|_w\bigr)$，按字典序取最小；
  \item 仅当候选满足“显著改进”（例如 $\widehat{\Delta\mathbf{J}}$ 在主缺口分量上超过阈值，或在容差 $\varepsilon$ 内严格优于空操作）才落地执行，以避免抖动；
  \item 若当前层无显著下降候选或硬缺口跨阈，则升级 $L\leftarrow L+1$（允许跳级）；
  \item 连续 $T$（按层设置，分钟级并带滞回/最小驻留）缺口稳定低位，则降级。
\end{enumerate}
该规则将运行期在线控制落成“可执行、可复算、可审计”的 tick 级闭环。
\end{definition}

\section{运行期在线控制 MVP：量子计算节点抗退相干（中大型；兼容室温超导格点）}
\label{sec:online-qpu-mvp}

\subsection{证书事件流与 tick 设定（量子版）}
\label{sec:online-qpu-event-tick}

\begin{remark}[现实边界与证书事件流动机]
主流超导量子计算系统通常依赖稀释制冷达到毫开尔文级温区（工程公开材料可参见 \cite{ref_nationalacademies_25196_ch13}，以及产业侧对 cryostat 的公开说明
  \cite{ref_ibm_goldeneye_cryostat}）。
中大型系统扩展时，控制/读出通道与线缆规模、低温端热负载与资源拥塞往往成为结构性瓶颈；相关研究与资源估计可参见 \cite{ref_arxiv_2601_03922}。
因此量子节点运行期在线控制需要覆盖“性能--控制--环境--资源”全链路证书事件流（对齐第~\ref{sec:migration-qpu-decoherence} 节的迁移口径）。
\end{remark}

\begin{remark}[定位：这里的“tick 在线控制”是策略层，而非纳秒级物理控制环]
本节的 tick 级在线控制对标 \textsf{OpenRA} 的“策略层 tick 决策”：
它在每 job/每秒/每分钟尺度上，从有限候选变分集合中选择一个可执行动作，
驱动校准、映射、调度、噪声建模与（必要时）结构级改写；
它不替代 AWG/FPGA 等纳秒级物理控制环，而是为其提供“何时改什么、改到什么程度”的证书化上层控制律。
\end{remark}

\begin{remark}[证书事件类型最小集（量子版；可回放）]
为保证“决策可复算、策略可回滚、证据可审计”，量子节点运行期证书事件流建议至少覆盖：
\begin{enumerate}
  \item \texttt{CALIBRATION\_UPDATE}（参数版本、覆盖范围、耗时）；
  \item \texttt{RB\_RESULT}/\texttt{IRB\_RESULT}（门误差证书）；
  \item \texttt{T1T2\_MEAS}（相干证书）与 \texttt{READOUT\_CAL}（读出证书）；
  \item \texttt{JOB\_SCHEDULED}/\texttt{MAPPING\_CHANGED}（调度/映射证书）；
  \item \texttt{ENV\_SNAPSHOT}（温度/热负载代理/EMI 等环境证书）；
  \item \texttt{MODEL\_UPDATE}（噪声模型切换）；
  \item \texttt{ROLLBACK}/\texttt{DEGRADE}/\texttt{RESET}（回退/降级/重置）。
\end{enumerate}
该最小集的目的不是约束实现细节，而是固定“事件类型 $\to$ 指标计算 $\to$ 候选评估”的可追溯链路。
\end{remark}

\begin{remark}[系统扩展的使能方向：cryo 控制与低温/片上脉冲发生（参考）]
当系统规模扩大时，控制/读出栈与线缆/热负载瓶颈会把一部分“结构性缺口”推到 L4（结构改写）层级；
例如低温/片上微波脉冲发生器与 cryogenic CMOS 控制电子学被视为重要的系统级使能方向之一（参见 \cite{ref_naturecomm_cryogenic_pulse_gen,ref_prxq_5_010326}）。
在本框架中，这类能力属于“候选变分族的扩张”与“结构代价坐标的改变”，必须通过证书事件流版本化纳入审计。
\end{remark}

\begin{definition}[量子节点 tick（建议三速率）]
\label{def:online-qpu-ticks}
量子节点在线控制可采用三速率 tick（并保持统一接口）：
\begin{enumerate}
  \item Fast tick（每 job 或 $1$--$10$s）：漂移检测、轻量噪声模型更新、调度决策；
  \item Cal tick（分钟级）：增量校准、局部重标定；
  \item Reset tick（小时级或事件触发）：全量校准、降级/维护/系统重启（必要时热循环）。
\end{enumerate}
不同 tick 类型的差异由候选族与 time-to-effect 吸收（定理~\ref{thm:online-multi-rate-unified-interface}）。
\end{definition}

\subsection{Gap 指标计算表（量子运行期 6 核心）}
\label{sec:online-qpu-gap-table}

\begingroup
\scriptsize
\setlength{\tabcolsep}{2pt}
\renewcommand{\arraystretch}{1.2}
\setlength{\LTleft}{0pt}
\setlength{\LTright}{0pt}
\begin{longtable}{|
  >{\raggedright\arraybackslash}p{0.13\linewidth}|
  >{\raggedright\arraybackslash}p{0.17\linewidth}|
  >{\raggedright\arraybackslash}p{0.17\linewidth}|
  >{\raggedright\arraybackslash}p{0.11\linewidth}|
  >{\raggedright\arraybackslash}p{0.20\linewidth}|
  >{\raggedright\arraybackslash}p{0.09\linewidth}|
  >{\raggedright\arraybackslash}p{0.06\linewidth}|}
\hline
\textbf{Gap 名称} & \textbf{工程含义} & \textbf{在线输入（最小集）} & \textbf{频率/滑窗} & \textbf{计算（可落地近似）} & \textbf{阈值建议} & \textbf{层级}
  \\
\hline
\endfirsthead
\hline
\textbf{Gap 名称} & \textbf{工程含义} & \textbf{在线输入（最小集）} & \textbf{频率/滑窗} & \textbf{计算（可落地近似）} & \textbf{阈值建议} & \textbf{层级}
  \\
\hline
\endhead
\textbf{Coherence\allowbreak Budget\allowbreak Gap}
& 相干预算不足（无法支撑所需深度/时长）
& $T_1/T_2$、门时长、目标电路深度/执行时间
& 每 job 或 $1$--$10$s
& 可达深度 $D_{\mathrm{ach}}\approx \min(T_1,T_2)/t_{\mathrm{gate}}$；
\newline Gap $:=\max(0,D_{\mathrm{req}}-D_{\mathrm{ach}})$
& 以业务目标或错误预算设 $D_{\mathrm{req}}$
& L2/L3
\\ \hline
\textbf{Gate\allowbreak Fidelity\allowbreak Gap}
& 门保真不足（RB/IRB 误差超阈）
& RB/IRB、门误差、漂移代理
& 每分钟或事件触发
& Gap $:=\max(0,\mathrm{err}-\mathrm{err}_{\mathrm{target}})$（可对关键门加权）
& 以目标应用或纠错阈值设定
& L1/L2
\\ \hline
\textbf{Crosstalk\allowbreak Leakage\allowbreak Gap}
& 串扰/泄漏结构性障碍
& 同时 RB、泄漏率、频谱碰撞事件
& 每分钟或每批任务
& Gap $:=$ 串扰指标超阈加权和 $+$ 频谱碰撞计数惩罚
& 随规模扩大逐步收紧
& L3/L4
\\ \hline
\textbf{Thermal\allowbreak Population\allowbreak Gap}
& 热激发/热噪声导致占据异常、读出退化
& 激发态占据 $p_e$、有效温度代理、冷端温漂
& 每秒--每分钟
& Gap $:=\max(0,p_e-p_{e,\mathrm{target}})+k\max(0,\Delta T-\Delta T_{\max})$
& $p_{e,\mathrm{target}}$ 由器件经验设定
& L2/\allowbreak L5（严重时）
\\ \hline
\textbf{Model\allowbreak Uncertainty\allowbreak Gap}
& 噪声模型/漂移机制不确定导致修复无效
& 预测--观测残差、校准年龄、拟合置信度
& 每 job / 每分钟
& Gap $:=$ 残差分位数 $+$ 校准过期惩罚
& 校准窗口按漂移速度设定
& L2
\\ \hline
\textbf{Topology\allowbreak Packaging\allowbreak Gap}
& 封装/互连/读出栈扩展导致瓶颈（热负载/线缆/通道拥塞/反射）
& 通道占用、读出带宽、反射/驻波代理、低温端热负载代理
& 每分钟
& Gap $:=$ 资源拥塞（排队/冲突）$+$ 热负载趋势 $+$ 反射代理超阈
& 随规模设“拥塞红线”
& L4
\\ \hline
\end{longtable}
\endgroup

\subsection{修复烈度与 L0--L5 候选变分生成表（量子运行期）}
\label{sec:online-qpu-tau-and-tower}

\begin{definition}[量子节点烈度六维（pred/org/pos/param/algo/rev）]
\label{def:online-qpu-tau}
量子节点在线控制的 $\Delta\tau\in\NN^6$ 与 OpenRA 六维对齐如下：
\begin{enumerate}
  \item $\Delta\tau_{\mathrm{pred}}$：噪声/漂移模型先验更新幅度；
  \item $\Delta\tau_{\mathrm{org}}$：作业调度/校准计划重排次数；
  \item $\Delta\tau_{\mathrm{pos}}$：有效拓扑/子图选择/频率规划等结构变更次数；
  \item $\Delta\tau_{\mathrm{param}}$：脉冲/偏置/解耦序列参数调整次数；
  \item $\Delta\tau_{\mathrm{algo}}$：编译映射/误差缓解策略栈切换次数；
  \item $\Delta\tau_{\mathrm{rev}}$：回退/止损次数（回滚校准、禁用 qubit 子集、系统重启/维护）。
\end{enumerate}
\end{definition}

\begingroup
\scriptsize
\setlength{\tabcolsep}{2pt}
\renewcommand{\arraystretch}{1.2}
\setlength{\LTleft}{0pt}
\setlength{\LTright}{0pt}
\begin{longtable}{|
  >{\raggedright\arraybackslash}p{0.06\linewidth}|
  >{\raggedright\arraybackslash}p{0.15\linewidth}|
  >{\raggedright\arraybackslash}p{0.22\linewidth}|
  >{\raggedright\arraybackslash}p{0.13\linewidth}|
  >{\raggedright\arraybackslash}p{0.15\linewidth}|
  >{\raggedright\arraybackslash}p{0.13\linewidth}|
  >{\raggedright\arraybackslash}p{0.09\linewidth}|}
\hline
\textbf{层级} & \textbf{允许变分类型} & \textbf{候选生成器（最小模板库）} & \textbf{$\Delta\tau$ 计数} & \textbf{强非交换约束} & \textbf{升级触发} &
  \textbf{降级条件} \\
\hline
\endfirsthead
\hline
\textbf{层级} & \textbf{允许变分类型} & \textbf{候选生成器（最小模板库）} & \textbf{$\Delta\tau$ 计数} & \textbf{强非交换约束} & \textbf{升级触发} &
  \textbf{降级条件} \\
\hline
\endhead
\textbf{L0}
& 常态运行（不改结构）
& 微小参数伺服（频率跟踪微调）
& 主要为 $\Delta\tau_{\mathrm{param}}$
& 禁止大规模重映射/重校准
& Gate/Coherence 缺口上升
& Gap 低且稳定
\\ \hline
\textbf{L1}
& slicing：快速校准子集
& 1) 读出校准刷新\newline 2) 单比特门微校准\newline 3) 局部二比特门校准（仅关键门）
& $\Delta\tau_{\mathrm{org}}$（插队校准）\newline $\Delta\tau_{\mathrm{param}}$（参数更新）
& \textbf{先修读出再判门误差}（避免误判）
& GateFidelityGap 上升但可快速压下
& 连续 $T$ 批任务稳定则降级
\\ \hline
\textbf{L2}
& 参数/策略同伦：系统性参数调整与表征增强
& 1) 动态解耦参数族切换\newline 2) 小幅 detune 频率避碰\newline 3) 提高噪声表征频率\newline 4) 保守脉冲整形模板
& $\Delta\tau_{\mathrm{param}}+\Delta\tau_{\mathrm{pred}}$
& 参数集切换需冷却（避免振荡）
& Model\allowbreak Uncertainty\allowbreak Gap
\newline 或 Coherence\allowbreak Budget\allowbreak Gap 主导
& 缺口稳定下降回到 L1/L0
\\ \hline
\textbf{L3}
& 组织同伦：有限候选映射/调度重排
& 1) qubit mapping 在候选交换集合内重排\newline 2) 作业队列避开热/拥塞窗口\newline 3) 有限规则库的门调度避串扰
& $\Delta\tau_{\mathrm{org}}+\Delta\tau_{\mathrm{algo}}$
& “先重映射还是先重校准”需证书判定（强非交换）
& Crosstalk/拥塞上升且 L2 无解
& 稳定后回落到 L2
\\ \hline
\textbf{L4}
& 强结构改写：有效拓扑/子图运行/模式切换
& 1) 禁用噪声子集 qubit，切换子图运行\newline 2) 可调耦合/频率规划更大幅重写\newline 3) 控制栈模式切换（若具备）
& $\Delta\tau_{\mathrm{pos}}+\Delta\tau_{\mathrm{algo}}$（高）
& 必须满足“显著降低 Packaging/\allowbreak Crosstalk”的证书门槛；
\newline 最小驻留
& Packaging/\allowbreak 拓扑类缺口主导
& 稳定后逐步放开但不反复大改
\\ \hline
\textbf{L5}
& 全局止损：全量校准/回滚/系统重启
& 1) 全量重校准\newline 2) 回滚到已知稳定版本\newline 3) 必要时维护/系统重启（热循环按需）
& $\Delta\tau_{\mathrm{rev}}{+}{=}1$
& 严格触发 + 冷却
& Thermal\allowbreak Population\allowbreak Gap 爆表
\newline 或系统不可用
& 恢复后回到 L3/L2
\\ \hline
\end{longtable}
\endgroup

\begin{remark}[“室温超导”格点的启用/禁用口径（运行期）]
在运行期在线控制中，“室温超导互连/屏蔽/器件”若作为目标格点，只应作为 L4/L5 级结构改写候选进入算子库；
若当前系统不具备该能力，则应在候选库中标记为 \textsf{disabled}（$\mathsf{feas}(u)=\texttt{false}$），避免被误当作 L0--L2 的可执行参数修复。
该口径与命题~\ref{prop:qpu-roomtemp-sc-as-sector-switch} 一致，并由撤稿/争议历史强化了“证书化/可复算/可审计”的必要性 \cite{ref_nature_2022_05294_9}。
\end{remark}

\subsection{tick 级在线决策规则（量子版）}
\label{sec:online-qpu-policy}

\begin{definition}[量子 planner tick：字典序缺口 + 候选有限 + 升降级滞回]
\label{def:online-qpu-policy}
对每个 planner tick（每 job 或 $1$--$10$s）：
\begin{enumerate}
  \item 计算主缺口向量（建议字典序）：
  \[
    \begin{aligned}
      \mathbf{J}
      &:=
      \bigl(
        \mathrm{TopologyPackagingGap},\mathrm{ThermalPopulationGap},\mathrm{CrosstalkLeakageGap},\\
        &\qquad \mathrm{CoherenceBudgetGap},\mathrm{GateFidelityGap},\mathrm{ModelUncertaintyGap}
      \bigr).
    \end{aligned}
  \]
  \item 在当前层级 $L$ 生成有限候选集合 $\mathcal{U}_L$（每个候选带 $\mathsf{tte}$，例如 L1 立刻生效、L3 需一批任务生效、L5 需长时间）；
  \item 估计候选评分 $\bigl(\widehat{\Delta\mathbf{J}}(u),\Delta|\tau(u)|_w\bigr)$，按字典序选择；
  \item 若无显著下降候选或硬缺口越阈，则升级（允许跳级）；若稳定则降级（带滞回与最小驻留）。
\end{enumerate}
\end{definition}

\section{运行期安全声明：护栏、人机协同与回滚（必须项）}
\label{sec:online-safety-guards}

\begin{remark}[电网场景：高风险系统必须人机协同与安全联锁]
电网运行期在线控制属于高风险系统：
\begin{enumerate}
  \item L3+（尤其 L4/L5）的候选必须来自经审查的预案模板库，并建议配置双确认、人机协同与安全联锁；
  \item 算法只做“有限候选集合上的证书化选择”，不做自由组合生成（避免不可预期耦合与操作爆炸）；
  \item 每一次动作必须写入 append-only 事件流以形成可追溯证书（见定义~\ref{def:online-grid-threat-certificates} 与相关 remark）。
\end{enumerate}
\end{remark}

\begin{remark}[量子场景：版本化、回滚与证书一致性是运行期底线]
量子节点运行期在线控制同样需要强护栏：
\begin{enumerate}
  \item 所有校准/编译/映射/控制参数变更必须版本化、可回滚、可复算，并写入证书事件流；
  \item L4/L5 的结构改写候选（子图切换、频率规划大改、控制栈模式切换、系统重启/维护）必须配置最小驻留与冷却；
  \item “室温超导”若作为目标格点，只能以 L4/L5 候选进入算子库，并以证书化门槛控制启用（见命题~\ref{prop:qpu-roomtemp-sc-as-sector-switch}）。
\end{enumerate}
\end{remark}

\begin{proposition}[命题：运行期 L3+ 必须采用“预案模板库 + 有限候选枚举”，禁止自由组合生成]
\label{prop:online-l3plus-template-only}
在运行期在线控制中，若 L3+ 允许对高代价、强非交换的操作词做自由组合生成，
则候选空间将出现组合爆炸并引入不可审计的不可预期耦合；
因此，为保证可执行性、安全性与审计一致性，L3+ 的候选应当限制为经审查的预案模板库，并在有限候选集合上枚举比较。
\end{proposition}

\begin{Certificate}[证书：组合爆炸 + 高风险耦合 $\Rightarrow$ 模板库约束是必要护栏]
\label{cert:online-l3plus-template-only}
证书登记以下谓词链：
\[
\begin{aligned}
  S_1 &: \text{强非交换操作词的自由组合}\ \Rightarrow\ \text{候选数随长度指数增长（组合爆炸）},\\
  S_2 &: \text{运行期高风险系统（电网/量子）}\ \Rightarrow\ \text{不可预期耦合不可接受且难以事后审计},\\
  S_3 &: (S_1\wedge S_2)\ \Rightarrow\ \text{必须将 L3+ 候选限制为预案模板库，并在有限集合上枚举}.
\end{aligned}
\]
\end{Certificate}

\begin{proof}[证明（证书；谓词因果链）]
由证书~\ref{cert:online-l3plus-template-only}：
自由组合导致候选空间指数爆炸，使在线枚举与一致性审计不可行；
而运行期高风险系统无法接受不可预期耦合与不可审计决策链；
因此必须以预案模板库约束 L3+ 候选，并在有限集合上证书化选择。
命题成立。证毕。
\end{proof}

\section{表达校准：公开里程碑、开放系统 ORL 与证书化壁垒}
\label{sec:calibration-public-open-orl}

\begin{remark}[目的（避免“已公开能力”与“稀缺增益”混淆）]
本节用于校准两类容易混淆的表达：
\begin{enumerate}
  \item 事实层：哪些能力已在公开文献中出现过（例如实时对抗、部分可观测、多智能体协作等）；
  \item 壁垒层：在开放高风险系统（电网/制造/量子等）中真正稀缺的增益是什么（例如强安全约束、证书化审计、算力/延迟硬约束、人机协同与回滚护栏等）。
\end{enumerate}
该校准的目标是：避免把“公开已达成的能力”当作壁垒去宣传，反而削弱真正的差异化点。
\end{remark}

\subsection{CPU 约束下的运行期 ORL：作为 GRL 路径积分的在线近似}
\label{sec:calibration-cpu-orl-online-approx}

\begin{remark}[校准点：运行期不重训练，而是有限候选的证书化变分下降]
在 CPU/延迟受限的运行期部署中，“在线 \textsf{ORL}（online reinforcement learning）”更可执行的口径不是运行期反向传播重训练，
而是：离线构建局部代价与可行性约束（含模板库与灵敏度近似），运行期在每个 tick 上执行
\[
  \mathsf{Snap}\ \to\ \mathcal{U}_L\ \to\ \text{字典序选择}\ \to\ \text{写入证书事件流},
\]
从而实现对主泛函的“证书化在线近似”。
该口径与 OpenRA 章的主生成泛函定义~\ref{def:openra-master-functional} 以及在线下降命题~\ref{prop:openra-online-descent-induces-tower} 同构。
\end{remark}

\begin{theorem}[公理降级：CPU 约束下的运行期 ORL 是主泛函的证书化在线近似（按 tick 数归纳）]
\label{thm:cpu-bounded-orl-online-approx}
设系统以证书事件流 $\mathsf{Ev}_{\mathsf{S}}$ 与证书快照 $\mathsf{Snap}_t$ 驱动运行期闭环；
并且在每个 tick $t$，控制器只在有限候选集合 $\mathcal{U}_L(\mathsf{Snap}_t)$ 上比较
\[
  \bigl(\widehat{\Delta\mathbf{J}}(u),\ \Delta|\tau(u)|_w\bigr)
\]
并按字典序选择（或其带阈值/滞回的近似版本），当当前层级无显著下降候选时才允许层级扩张（L0$\to$L5）。
则该运行期过程可视为对主生成泛函 $\mathcal{S}_{\mathrm{master}}$（定义~\ref{def:openra-master-functional}）的一个证书化在线近似：
\begin{enumerate}
  \item 在线决策链可复算：每个 tick 的“快照--候选--选择--执行”均可由证书事件流重放；
  \item 在线近似链可解释：在候选集与 time-to-effect 约束下，字典序选择实现“缺口优先 + 低层优先”的局部变分下降；
  \item 层级扩张可控：仅在低层不可修复证书成立时才扩大可行变分族，从而与运行期安全护栏（第~\ref{sec:online-safety-guards} 节）一致。
\end{enumerate}
\end{theorem}

\begin{Certificate}[数学归纳法证书：按 tick 数归纳的“可复算 + 局部下降 + 可控扩张”链]
\label{cert:cpu-bounded-orl-online-approx}
定义谓词 $P(T)$ 为：对前 $T$ 个 tick 的执行序列 $\mathsf{word}(u_0)\circ\cdots\circ\mathsf{word}(u_{T-1})$，
存在一份可重放的证书链，逐 tick 记录 $\bigl(\mathsf{Snap}_t,\mathcal{U}_L,\widehat{\Delta\mathbf{J}},\Delta\tau,\mathsf{tte},
  \mathsf{feas}\bigr)$，
并满足以下三条：
\begin{enumerate}
  \item 可复算：由记录可重放每个 tick 的候选枚举与字典序选择结果；
  \item 局部下降：若存在可行候选使主缺口分量显著下降，则选择器不会输出严格劣于空操作的候选；
  \item 可控扩张：仅当“当前层无显著下降候选”或“硬缺口跨阈”成立时才允许升级到更高层级。
\end{enumerate}
证书化要求为：$P(1)$ 成立，且 $(\forall T\in\NN_{\ge 1})\bigl(P(T)\Rightarrow P(T+1)\bigr)$。
\end{Certificate}

\begin{proof}[证明（数学归纳法证书；谓词因果链）]
由证书~\ref{cert:cpu-bounded-orl-online-approx}：
\begin{itemize}
  \item 基步：$T=1$ 时，$\mathsf{Snap}_0$ 与有限候选集 $\mathcal{U}_L(\mathsf{Snap}_0)$ 可枚举比较并记录；
  字典序选择输出 $u_0$ 或空操作，满足“可复算/不劣于空操作”的防抖判据，且升级条件可由记录判定，故 $P(1)$ 成立。
  \item 归纳步：假设 $P(T)$ 成立。对第 $T$ 个 tick，证书事件流给出 $\mathsf{Snap}_T$；
  有限候选集与候选签名（定义~\ref{def:online-candidate-signature}）保证可枚举比较并记录；
  选择器按同一字典序准则输出 $u_T$ 或空操作，并仅在升级谓词成立时扩张层级。
  因此 $P(T+1)$ 成立。
\end{itemize}
由数学归纳法得对任意有限 $T$，证书化在线近似链成立；从而定理~\ref{thm:cpu-bounded-orl-online-approx} 成立。证毕。
\end{proof}

\subsection{公开里程碑校准：实时、部分可观测、多智能体对抗并非空白}
\label{sec:calibration-public-realtime-pomdp-marl}

\begin{proposition}[命题：公开文献已包含实时对抗与部分可观测多智能体的强基准]
\label{prop:public-realtime-marl-exists}
若以“公开触达”指“存在公开论文/技术报告可供引用”，则实时对抗、部分可观测、多智能体协作对抗并非空白；
例如在 RTS 与实时团队对抗类基准上，公开工作已展示过达到顶级水平的系统（可参见 \cite{ref_pubmed_31666705,ref_openai_dota2}）。
\end{proposition}

\begin{Certificate}[证书：公开论文/技术报告记录（引用即证书）]
\label{cert:public-realtime-marl-exists}
证书登记以下公开记录作为“公开触达”的证据：
\begin{enumerate}
  \item RTS/部分可观测多智能体强化学习：公开论文/数据库收录记录 \cite{ref_pubmed_31666705}；
  \item 实时团队对抗强化学习：公开技术报告 \cite{ref_openai_dota2}。
\end{enumerate}
\end{Certificate}

\begin{proof}[证明（证书；谓词因果链）]
由证书~\ref{cert:public-realtime-marl-exists}，上述工作以公开论文/公开技术报告的形式存在并可被引用；
因此“实时对抗与部分可观测多智能体对抗完全未公开触达”的表述不成立，命题成立。证毕。
\end{proof}

\subsection{公开里程碑校准：自驱动实验室闭环已公开，但与开放系统博弈仍有差距}
\label{sec:calibration-public-sdl-gap}

\begin{proposition}[命题：自动化科研/制造的闭环系统在公开领域已出现]
\label{prop:public-self-driving-lab-exists}
若以“闭环系统”指“机器人/设备 + 数据/模型 + 主动学习/优化 + 自动执行”的循环结构，
则材料/化学方向的自驱动实验室（self-driving laboratory）已在公开文献中出现并形成综述体系（例如 \cite{ref_nature_2023_06734_w,ref_acs_chemrev_4c00055}）。
\end{proposition}

\begin{Certificate}[证书：闭环自驱动实验室的公开记录（论文/综述）]
\label{cert:public-self-driving-lab-exists}
证书登记以下公开记录作为闭环系统的引用证据：
\begin{enumerate}
  \item 自主实验室案例论文：\cite{ref_nature_2023_06734_w}；
  \item 自驱动实验室综述：\cite{ref_acs_chemrev_4c00055}。
\end{enumerate}
\end{Certificate}

\begin{proof}[证明（证书；谓词因果链）]
由证书~\ref{cert:public-self-driving-lab-exists}，闭环自驱动实验室在公开论文与综述中存在且可被引用；
因此“闭环自动化科研仅停留在封闭回合制”的笼统表述需要校准。命题成立。证毕。
\end{proof}

\begin{remark}[差距口径：公开 SDL $\neq$ 运行期开放高风险 ORL]
上述公开 SDL 的闭环形态并不自动等价于“开放系统博弈下的运行期 ORL”：
在电网/制造/量子等高风险系统中，更关键的差距往往集中在
实时非平稳扰动、强安全约束、证书化审计、人机协同与算力/延迟硬约束等方面；
相关的 safe RL 与 sim-to-real 研究脉络可参见综述 \cite{ref_jmlr_garcia15a_safe_rl,ref_sciencedirect_s0004370222001515,
  ref_arxiv_2502_13187}。
\end{remark}

\subsection{差异化壁垒：公开已覆盖“实时+对抗”，但未普遍覆盖“开放+高风险+证书化+算力约束”}
\label{sec:calibration-diff-barrier}

\begin{definition}[公开触达的三维分解：实时性/开放性/审计性]
\label{def:public-coverage-3axis}
为避免“公开触达”含混不清，建议将其分解为三个维度：
\begin{enumerate}
  \item 实时性维度：回合制 vs 实时（含长时程与高频动作）；
  \item 开放性维度：封闭可模拟环境 vs 真实开放环境（非平稳扰动、外部世界不可完全模拟）；
  \item 审计性维度：黑箱策略输出 vs 证书化可复算控制（事件流、回滚与可归责）。
\end{enumerate}
\end{definition}

\begin{theorem}[公理降级：更精确的结论应为“公开已覆盖实时对抗，但开放高风险证书化 ORL 仍是壁垒”]
\label{thm:public-coverage-corrected-statement}
在定义~\ref{def:public-coverage-3axis} 的三维分解下，更精确且不自相矛盾的公开触达结论应为：
\begin{enumerate}
  \item 在“实时性维度”，公开工作已覆盖实时对抗与多智能体协作对抗的强基准（证据见命题~\ref{prop:public-realtime-marl-exists}）；
  \item 但在“开放性 + 审计性 + 算力约束”叠加的运行期部署维度，仍需依赖证书事件流、修复烈度会计与分层修复塔等护栏结构（第~\ref{sec:online-control-mvp}--\ref{sec:online-safety-guards} 节），
  因而该部分更接近工程壁垒而非单纯算法跑分。
\end{enumerate}
\end{theorem}

\begin{Certificate}[证书：三轴分解下的证据拼接]
\label{cert:public-coverage-corrected-statement}
证书登记以下谓词链：
\[
\begin{aligned}
  U_1 &: \text{实时对抗公开证据存在}\Rightarrow \text{实时性维度已覆盖} \quad (\text{由命题~\ref{prop:public-realtime-marl-exists}}),\\
  U_2 &: \text{开放高风险部署}\Rightarrow \text{必须满足安全/合规/审计/算力约束（见第~\ref{sec:online-safety-guards} 节）},\\
  U_3 &: (U_1\wedge U_2)\Rightarrow \text{“公开已覆盖实时对抗，但开放高风险证书化 ORL 仍是壁垒”}.
\end{aligned}
\]
并将 safe RL、sim-to-real 与可解释性研究作为背景参照 \cite{ref_jmlr_garcia15a_safe_rl,ref_arxiv_2502_13187,ref_acm_3616864}。
\end{Certificate}

\begin{proof}[证明（证书；谓词因果链）]
由证书~\ref{cert:public-coverage-corrected-statement}：
$U_1$ 给出“实时性维度已有公开强基准”；
$U_2$ 给出“开放高风险部署必须满足证书化护栏与算力约束”；
二者合取得到 $U_3$，从而定理~\ref{thm:public-coverage-corrected-statement} 成立。证毕。
\end{proof}

\subsection{“上帝视角”校准：不是全知，而是证书化全链路控制视角}
\label{sec:calibration-godview}

\begin{definition}[工程化“上帝视角”】【证书化全链路控制视角】]
\label{def:godview-auditable}
在本笔记的工程语境中，“上帝视角”不应理解为“全知全能”，而应降级为如下可检验对象：
\[
  \mathsf{GodView}
  :=
  (\mathsf{Ev}_{\mathsf{S}},\ \mathsf{Snap},\ \mathsf{word},\ \mathbf{J},\ \tau,\ \mathsf{Rollback}),
\]
其中：
\begin{enumerate}
  \item $\mathsf{Ev}_{\mathsf{S}}$ 为 append-only 证书事件流；
  \item $\mathsf{Snap}$ 为证书快照（例如定义~\ref{def:online-grid-snapshot-fields}）；
  \item $\mathsf{word}$ 为可复算的非交换算子词（动作模板序列）；
  \item $\mathbf{J}$ 与 $\tau$ 分别为缺口证书与结构代价会计；
  \item $\mathsf{Rollback}$ 为版本化回滚与止损接口（第~\ref{sec:online-safety-guards} 节）。
\end{enumerate}
该定义强调的不是“看见一切”，而是“\textbf{可复算、可审计、可回滚}的全链路控制视角”。
\end{definition}

\begin{corollary}[推论：将“战略价值无限大”替换为三条可验证主张]
\label{cor:replace-infinite-strategic-value}
在公开表达中，为避免不可证伪修辞，建议将“战略价值无限大”替换为以下三条可验证主张：
\begin{enumerate}
  \item 将开放系统控制从经验策略提升为证书化变分下降：可复算、可审计、可回滚（定义~\ref{def:godview-auditable}）；
  \item 将高风险探索从自由探索变为分层可行域扩张：L0--L5 的修复塔与模板库护栏（命题~\ref{prop:online-l3plus-template-only}）；
  \item 将“对手/扰动”从外生噪声内化为泛函耦合项：权重/硬约束/候选族随证书触发切换（第~\ref{sec:online-control-mvp} 节）。
\end{enumerate}
\end{corollary}

\begin{Certificate}[证书：三条主张的可检验落点]
\label{cert:replace-infinite-strategic-value}
证书登记三条主张的检验接口分别为：
\begin{enumerate}
  \item 证书化变分下降：事件流可重放 + 候选比较记录（定理~\ref{thm:cpu-bounded-orl-online-approx}）；
  \item 分层可行域扩张：模板库 + 冷却/驻留 + 升降级滞回（第~\ref{sec:online-control-mvp}--\ref{sec:online-safety-guards} 节）；
  \item 扰动内化：缺口向量与硬约束壁垒进入主泛函与在线选择准则（定义~\ref{def:openra-master-functional}）。
\end{enumerate}
\end{Certificate}

\begin{proof}[证明（证书；谓词因果链）]
由证书~\ref{cert:replace-infinite-strategic-value}：
三条主张分别对应到可重放事件流、分层模板库护栏与主泛函/在线准则三个可检验接口；
因此其可验证性不依赖修辞而依赖证书链本身。推论成立。证毕。
\end{proof}

\section{理论/战略分离：Jacobiator-gap、同伦修复塔与分级发布}
\label{sec:jacgap-theory-strategy}

\begin{remark}[分层动机（两类命题不可混写）]
在对外叙事中容易把两类不同层面的命题混写，从而造成“事实可反驳/壁垒不清晰”的沟通风险。
本节将命题分为两层：
\begin{enumerate}
  \item \textbf{理论层命题}：以“\JacGap{} 证书 + 同伦修复塔 + 变分泛函驱动的同伦扭曲”为核心结构，
  讨论其是否构成开放系统博弈（运行期在线控制 + 对抗 + 非平稳扰动 + 强约束）的一个\textbf{可证伪普适闭包骨架}；
  \item \textbf{战略层命题}：若理论在高风险领域具备能力优势，则公开应用细节可能带来竞争风险与信息危害，
  因此需要分级发布与最小披露策略。
\end{enumerate}
本节分别给出两层命题的形式化口径，并把“可控风险的发布策略”写成可执行的分级规则。
\end{remark}

\subsection{术语校准：Jacobiator-gap（\JacGap）不是 Jacobian（导数矩阵）}
\label{sec:jacgap-terminology}

\begin{definition}[Jacobiator-gap（\JacGap）证书：法则一致性/闭包缺口]
\label{def:jacobiator-gap-certificate}
设 $\mathcal{G}$ 为某个扇区 $\Sigma$ 中的“生成元算子集合”（可理解为当前允许的关键操作词的生成集合）。
对任意 $A,B\in\mathcal{G}$ 定义对易子
\[
  [A,B]:=A\circ B - B\circ A.
\]
对任意三元组 $(A,B,C)\in\mathcal{G}^3$ 定义 Jacobiator
\[
  \mathrm{Jacobiator}(A,B,C)
  :=
  [[A,B],C] + [[B,C],A] + [[C,A],B].
\]
称 $\JacGap(\Sigma)$ 为一个\textbf{Jacobiator-gap 证书}，若它由上述 Jacobiator 的某种可审计聚合得到（例如范数平方和或上确界）：
\[
  \JacGap(\Sigma)
  :=
  \mathrm{Agg}\bigl(\|\mathrm{Jacobiator}(A,B,C)\|:\ (A,B,C)\in\mathcal{G}^3\bigr),
  \qquad
  \JacGap(\Sigma)\ge 0.
\]
\end{definition}

\begin{remark}[校准说明：为何不用“Jacobian-gap”这个词]
“Jacobian”在常见数学语境中指导数矩阵；而本节的 \JacGap{} 指的是 \textbf{Jacobi identity 的缺口（Jacobiator）}。
为避免误读，本笔记统一使用 “Jacobiator-gap（\JacGap）”。
\end{remark}

\subsection{理论层：可证伪普适闭包骨架（桥接命题 A/B/C）}
\label{sec:jacgap-bridge-abc}

\begin{remark}[质变点（对齐本笔记坐标系）]
本笔记关注的不是“在固定规则下赢得某个博弈”，而是开放高风险系统中：
\textbf{规则/可行动作族会随时间漂移}、外界扰动非平稳、且必须可审计可回滚。
因此核心困难不是单次策略最优，而是“法则一致性/闭包性”的运行期证书与恢复路径。
\JacGap{} 在该口径下充当“规则一致性缺口证书”，同伦修复塔则提供“可行域分层扩张”的恢复机制（参见第~\ref{sec:online-control-mvp}--\ref{sec:online-safety-guards} 节）。
\end{remark}

\begin{proposition}[桥接命题 A（表示性）：运行期策略可表示为动作词（或其分布）]
\label{prop:bridge-a-representation}
设运行期系统的执行接口为一个动作词集合 $\mathcal{U}$（可为模板库的有限并），
并且策略 $\pi$ 在每个 tick 将历史观测（或证书快照）映射为一个动作词 $u\in\mathcal{U}$。
则任意有限时域 $T$ 上的运行期执行过程都可表示为一个非交换动作词序列
\[
  \mathsf{word}(\pi;T) := u_{T-1}\circ\cdots\circ u_0,
  \qquad u_t\in\mathcal{U},
\]
或在随机策略情形下表示为动作词序列的分布。
\end{proposition}

\begin{Certificate}[证书：策略定义 $\Rightarrow$ 动作词序列的存在]
\label{cert:bridge-a-representation}
证书登记：对每个 tick，策略输出 $u_t\in\mathcal{U}$；
由复合定义直接得到长度为 $T$ 的动作词序列；
若策略随机，则对输出分布做同样登记即可得到序列分布。
\end{Certificate}

\begin{proof}[证明（证书；谓词因果链）]
由证书~\ref{cert:bridge-a-representation}：
每个 tick 的输出属于 $\mathcal{U}$，按时间顺序复合即可得到动作词序列；
随机情形对输出的分布登记即可。
命题成立。证毕。
\end{proof}

\begin{proposition}[桥接命题 B（必要性）：若缺少“规则一致性证书”，则运行期审计闭环无法自洽]
\label{prop:bridge-b-necessity}
在开放高风险运行期闭环中，若系统允许“动作词语义/约束集/可行动作族”随时间漂移，
但决策链缺少任何形式的“规则一致性证书”（\JacGap{} 或其等价替代），
则至少出现以下一类不可避免的问题：
\begin{enumerate}
  \item 决策不可复算：事后无法判定当时使用的规则/约束版本，从而证书事件流无法闭合；
  \item 决策不可回滚：缺少规则版本与结构代价会计（$\tau$）会导致回滚边界不清；
  \item 安全约束不可保证：在规则漂移未被证书化的情况下，硬约束合规仅能依赖经验假设而无法审计验证。
\end{enumerate}
\end{proposition}

\begin{Certificate}[证书：漂移存在 + 无证书 $\Rightarrow$ 证书链断裂]
\label{cert:bridge-b-necessity}
证书登记谓词链：
\[
\begin{aligned}
  B_1 &: \text{规则/约束/动作族随时间漂移} \Rightarrow \text{需要记录“当时版本”才能复算},\\
  B_2 &: \text{缺少一致性证书} \Rightarrow \text{版本无法在事件流中被可核对地登记},\\
  B_3 &: (B_1\wedge B_2) \Rightarrow \text{证书链断裂（不可复算/不可回滚/不可审计）}.
\end{aligned}
\]
\end{Certificate}

\begin{proof}[证明（证书；谓词因果链）]
由证书~\ref{cert:bridge-b-necessity}：
若存在规则漂移，则复算必需知道当时采用的规则/约束版本；
若缺少一致性证书，则版本无法以可核对方式进入事件流；
因此证书链断裂，至少导致不可复算/不可回滚/不可审计之一成立。
命题成立。证毕。
\end{proof}

\begin{proposition}[桥接命题 C（最小性）：修复塔给出“最小烈度的规则一致性恢复路径”（在本框架内）]
\label{prop:bridge-c-minimality}
设修复塔给出一族分层候选集合
\[
  \mathcal{U}_0 \subseteq \mathcal{U}_1 \subseteq \cdots \subseteq \mathcal{U}_5,
\]
并以字典序目标 $\bigl(\mathbf{J},|\tau|_w\bigr)$ 选择候选。
若存在某层级 $L$ 与候选 $u\in\mathcal{U}_L$ 使得 $\JacGap$ 满足容差约束（例如 $\JacGap\le \varepsilon$）且主缺口显著下降，
则采用“从低到高逐层放开”的策略，总能找到一个\textbf{最小层级} $L^\star$ 使上述条件成立；
并且在固定 $L^\star$ 的候选集合内，字典序选择会给出相对于同集合的\textbf{最小结构烈度}（不劣于空操作的防抖口径）。
\end{proposition}

\begin{Certificate}[证书：有限层级取最小 + 层内字典序最小元]
\label{cert:bridge-c-minimality}
证书登记两步：
\begin{enumerate}
  \item 对 $L=0,1,\dots,5$ 依次检查“是否存在 $u\in\mathcal{U}_L$ 使得 $\JacGap\le\varepsilon$ 且缺口显著下降”；取第一个满足者为 $L^\star$；
  \item 在 $\mathcal{U}_{L^\star}$ 内按字典序比较候选评分，并以“确有改进”防抖判据拒绝劣于空操作的候选。
\end{enumerate}
\end{Certificate}

\begin{proof}[证明（证书；谓词因果链）]
由证书~\ref{cert:bridge-c-minimality}：
层级集合有限，故存在最小满足层级 $L^\star$；
在该层级的候选集合内，字典序比较给出最小元；
防抖判据保证输出不劣于空操作。
命题成立。证毕。
\end{proof}

\subsection{战略层：信息危害分级与发布策略（P/S/R）}
\label{sec:jacgap-release}

\begin{definition}[分级发布策略（P/S/R）]
\label{def:release-psr}
为兼顾“学术坐标/优先权”与“高风险系统的信息危害控制”，建议采用分级发布：
\begin{enumerate}
  \item \textbf{P（公开纯理论）}：公开 \JacGap{} 证书、变分泛函、修复塔语义、以及与传统方法的投影关系；不公开现实高风险领域的具体预案模板库与阈值点表。
  \item \textbf{S（公开沙盘实现）}：公开 OpenRA/仿真环境下的候选生成器、Gap 计算、证书回放与 ablation；仍不公开现实关键基础设施的可直接部署细节。
  \item \textbf{R（受限披露落地）}：在白名单合作或 NDA 等受限条件下，披露行业版的算子库、阈值、预案模板与联锁接口，并配置人机协同与回滚护栏。
\end{enumerate}
\end{definition}

\begin{proposition}[命题：高风险领域的 L3+ 预案模板库应当受限披露]
\label{prop:release-template-restriction}
在高风险系统（例如电网、量子等）的运行期在线控制中，
若将 L3+（拓扑/分区/止损等）预案模板库的具体动作词与阈值细节完全公开，
则会同时放大以下风险：
\begin{enumerate}
  \item 信息危害：误用/滥用导致的安全风险；
  \item 对抗风险：对手可据此推断防护边界并构造对抗性输入；
  \item 竞争风险：可直接复刻的工程壁垒被抢占。
\end{enumerate}
因此分级发布（定义~\ref{def:release-psr}）是更稳健的发布策略。
\end{proposition}

\begin{Certificate}[证书：高风险动作模板与运行期护栏的耦合]
\label{cert:release-template-restriction}
证书登记谓词链：
\[
\begin{aligned}
  R_1 &: \text{L3+ 候选必须来自预案模板库（命题~\ref{prop:online-l3plus-template-only}）},\\
  R_2 &: \text{预案模板库包含高风险可执行动作词与阈值细节},\\
  R_3 &: (R_1\wedge R_2)\Rightarrow \text{公开模板库会直接公开高风险可执行细节},\\
  R_4 &: \text{公开高风险细节}\Rightarrow \text{信息危害/对抗风险/竞争风险至少一项上升}.
\end{aligned}
\]
\end{Certificate}

\begin{proof}[证明（证书；谓词因果链）]
由证书~\ref{cert:release-template-restriction}：
命题~\ref{prop:online-l3plus-template-only} 给出 L3+ 对预案模板库的依赖；
而模板库包含可直接执行的高风险动作词与阈值细节；
若完全公开模板库，则等价于公开高风险可执行细节，从而放大信息危害/对抗风险/竞争风险。
因此采用分级发布更稳健。命题成立。证毕。
\end{proof}

\begin{remark}[对标翻译层：把主流语言映射为投影条件]
为降低沟通摩擦，可显式给出“翻译层”：将 Markov game、robust control、safe RL、分层 RL 等主流语言
映射为本框架的投影条件（例如 $\JacGap=0$、禁用高层烈度、候选集退化为固定策略族等）。
同时应承认：实时对抗与多智能体协作对抗在公开领域已有强里程碑（例如 \cite{ref_pubmed_31666705,ref_openai_dota2}），
本框架的差异化更多来自“开放+高风险+证书化+算力约束”的叠加维度（见第~\ref{sec:calibration-diff-barrier} 节）。
\end{remark}

\section{优先权固化：学术时间戳、prior art 与信息危害分级}
\label{sec:priority-ip-infohazard}

\subsection{学术层：公开固化优先权的四要件}
\label{sec:priority-academic}

\begin{definition}[学术优先权固化谓词（四要件）]
\label{def:priority-four-conditions}
称一个公开发布版本 $D$（例如某份 PDF/预印本/论文集版本）在争议场景下\textbf{固化学术优先权}，若其满足以下四个外部可验证要件：
\begin{enumerate}
  \item \textbf{公共可访问}（Publicly accessible）：共同体成员无需经作者许可即可获取同一版本；
  \item \textbf{第三方可验证时间戳}（Third-party timestamp）：发布时间由第三方平台记录（例如提交时间/DOI 发行时间等）；
  \item \textbf{可检索可引用}（Citable）：他人可稳定引用（URL/DOI/版本号/哈希等）；
  \item \textbf{足够披露/可复核}（Enablement）：包含可复核的定义、命题与构造关系，而不是口号。
\end{enumerate}
我们用谓词
\[
  \mathsf{Priority}(D)
\]
表示“$D$ 满足上述四要件并因此在共同体语境下固化优先权”。
\end{definition}

\begin{theorem}[公理降级：四要件足以固化学术优先权（按要件数归纳）]
\label{thm:priority-four-conditions-sufficient}
若公开版本 $D$ 满足定义~\ref{def:priority-four-conditions} 的四要件，则 $\mathsf{Priority}(D)$ 成立。
该结论的工程含义是：仅有“作者/年份/版本号”的内部元数据并不足以在外部争议场景中固化优先权；
必须具备“可访问 + 时间戳 + 可引用 + 可复核”的外部可验证证据链（参见讨论类参考 \cite{ref_academia_priority}）。
\end{theorem}

\begin{Certificate}[数学归纳法证书：按要件数 $n$ 归纳的优先权证据链]
\label{cert:priority-four-conditions-sufficient}
令四要件依次记为 $C_1,C_2,C_3,C_4$，并定义谓词 $P(n)$ 为：
“满足 $C_1,\dots,C_n$ 时，优先权证据链在共同体语境下\textbf{至少}具备前 $n$ 个环节的可验证性”。
证书化要求为：
\[
  P(1)\ \text{成立},\qquad
  (\forall n\in\{1,2,3\})\ \bigl(P(n)\Rightarrow P(n+1)\bigr).
\]
并规定 $P(4)\Rightarrow \mathsf{Priority}(D)$（四环节闭合即固化）。
\end{Certificate}

\begin{proof}[证明（数学归纳法证书；谓词因果链）]
由证书~\ref{cert:priority-four-conditions-sufficient}：
\begin{itemize}
  \item 基步 $n=1$：满足公共可访问 $C_1$，则共同体成员可以独立获取同一版本，证据链的“可见性”环节成立，故 $P(1)$ 成立。
  \item 归纳步：假设 $P(n)$ 成立。
  \begin{itemize}
    \item 若再满足 $C_{n+1}$，则证据链新增一个外部可验证环节：
    $C_2$ 提供第三方时间戳（使“何时公开”可验证）；
    $C_3$ 提供稳定引用键（使“如何引用”可验证）；
    $C_4$ 提供可复核披露（使“主张内容”可验证）。
  \end{itemize}
  因而 $P(n+1)$ 成立。
\end{itemize}
由数学归纳法得 $P(4)$ 成立，按证书约定 $P(4)\Rightarrow \mathsf{Priority}(D)$，定理成立。证毕。
\end{proof}

\begin{remark}[校准：固化强度取决于外部证据链，而非自述]
若仅在私域发布、或缺少第三方时间戳/稳定引用键，则“优先权固化”只在作者自洽体系内部成立，
在外部争议场景中证据链可能不闭合；因此建议将“发言权”表达为更可验证的两条断言：
\textbf{优先权断言}（可引用时间戳披露）与 \textbf{prior art 断言}（构成现有技术披露）。
\end{remark}

\subsection{知识产权层：防御性公开（prior art）与专利窗口权衡}
\label{sec:priority-ip}

\begin{definition}[防御性公开（prior art）口径]
\label{def:defensive-publication}
称公开版本 $D$ 具备“防御性公开”作用，若其满足“可公开检索 + 可核验时间戳 + 足够披露”的证据链，
从而在专利语境下更可能被认定为现有技术（prior art）并限制第三方对同一核心思想取得排他性控制权
（相关实践讨论可参见 \cite{ref_perkinscoie_prior_art}）。
\end{definition}

\begin{proposition}[命题：规范公开可增强 prior art 证据效力（但也可能压缩自身专利窗口）]
\label{prop:prior-art-and-window-tradeoff}
若某个核心思想以满足定义~\ref{def:priority-four-conditions} 的方式公开，则：
\begin{enumerate}
  \item 从防御性公开角度看，它更可能形成可检索、可核验的 prior art 证据，降低第三方在同一核心思想上取得专利垄断的空间；
  \item 从自身专利布局角度看，公开披露可能压缩后续在部分法域的专利新颖性窗口（不同法域的宽限期制度差异显著，需逐国评估）。
\end{enumerate}
\end{proposition}

\begin{Certificate}[证书：prior art 强化与新颖性窗口压缩的双刃链]
\label{cert:prior-art-and-window-tradeoff}
证书登记两条并行谓词链：
\begin{enumerate}
  \item $I_1$：可检索 + 可核验时间戳 + 足够披露 $\Rightarrow$ prior art 证据更强（参见 \cite{ref_perkinscoie_prior_art}）；
  \item $I_2$：公开披露 $\Rightarrow$ 在部分法域触发新颖性约束；宽限期制度差异需参考制度说明（参见 \cite{ref_wipo_patent_protection,ref_uspto_mpep_2153,
    ref_cnipa_patent_law}）。
\end{enumerate}
\end{Certificate}

\begin{proof}[证明（证书；谓词因果链）]
由证书~\ref{cert:prior-art-and-window-tradeoff}：
$I_1$ 给出“规范公开增强 prior art 证据效力”；
$I_2$ 给出“公开披露可能影响后续新颖性窗口，且制度差异需逐国评估”；
二者并行成立即得到命题的双刃结论。
命题成立。证毕。
\end{proof}

\begin{remark}[免责声明（非法律意见）]
本节只用于工程叙事中的“风险结构校准”，不构成法律意见；
若确有专利/商业秘密布局意图，应由专业律师按目标国家制度评估披露与申请节奏。
\end{remark}

\subsection{战略与信息危害：分级发布与“优先权加固包”}
\label{sec:priority-release-and-pack}

\begin{theorem}[公理降级：P/S/R 分级发布可同时服务优先权固化与信息危害控制（按级别归纳）]
\label{thm:release-psr-serves-both}
在定义~\ref{def:release-psr} 的 P/S/R 分级发布策略下：
\begin{enumerate}
  \item P 级（纯理论）优先满足学术优先权固化所需的可公开可引用披露（定理~\ref{thm:priority-four-conditions-sufficient}）；
  \item S 级（沙盘）提供可复现实验与 ablation，用于增强共同体可复核性与工程可信度；
  \item R 级（受限披露）对高风险系统的 L3+ 预案模板库进行最小披露，从而与命题~\ref{prop:release-template-restriction} 的信息危害控制一致。
\end{enumerate}
\end{theorem}

\begin{Certificate}[数学归纳法证书：按发布级别链 P$\to$S$\to$R 归纳的“固化--可信--控危害”]
\label{cert:release-psr-serves-both}
定义谓词 $Q(k)$ 为：“发布到第 $k$ 个级别后，同时满足：
(i) 优先权证据链增强；(ii) 可复核可信度增强；(iii) 高风险细节仍受控披露。”
其中 $k=1,2,3$ 分别对应 P、P+S、P+S+R。
证书化要求为：
\[
  Q(1)\ \text{成立},\qquad Q(1)\Rightarrow Q(2),\qquad Q(2)\Rightarrow Q(3).
\]
\end{Certificate}

\begin{proof}[证明（数学归纳法证书；谓词因果链）]
由证书~\ref{cert:release-psr-serves-both}：
\begin{itemize}
  \item $Q(1)$：P 级公开纯理论与核心命题，增强“可访问/可引用/可复核”的优先权证据链；
  \item $Q(1)\Rightarrow Q(2)$：在不公开高风险模板库的前提下加入 S 级沙盘复现，使可复核可信度增强；
  \item $Q(2)\Rightarrow Q(3)$：R 级对行业细节做受限披露，从而在增强落地能力的同时保持信息危害可控（命题~\ref{prop:release-template-restriction}）。
\end{itemize}
因此 $Q(3)$ 成立，从而定理~\ref{thm:release-psr-serves-both} 成立。证毕。
\end{proof}

\begin{definition}[优先权加固包（Priority Pack；最小六件套）]
\label{def:priority-pack}
为将“已公开”提升到“外部争议场景中几乎不可争议”，建议形成一个\textbf{优先权加固包}：
\begin{enumerate}
  \item 上链式可验证时间戳：选择公认平台发布并获得 DOI/提交编号；
  \item 版本化与变更日志：明确版本序列与差异点；
  \item 哈希指纹：登记 PDF 的 SHA256（或等价哈希）并与公开页面同步；
  \item 一页核心命题清单：集中列出 3--6 条关键命题（例如 \JacGap{} 必要性、修复塔最小性、投影条件等）；
  \item 对标投影表：将主流方法映射为本框架的投影条件（见第~\ref{sec:jacgap-release} 节与第~\ref{sec:calibration-diff-barrier} 节）；
  \item 沙盘复现实验：以 OpenRA/仿真系统给出可重复的 ablation（避免泄露现实高风险模板库）。
\end{enumerate}
\end{definition}

\begin{proposition}[命题：优先权加固包可显著提升外部争议场景下的证据强度]
\label{prop:priority-pack-strengthens}
若公开版本 $D$ 附带定义~\ref{def:priority-pack} 的优先权加固包，
则其在外部争议场景下满足定义~\ref{def:priority-four-conditions} 的证据强度显著提升：
时间戳更可核验、版本更可定位、内容更可复核、引用更稳定，从而优先权固化更稳健。
\end{proposition}

\begin{Certificate}[证书：六件套分别补强“时间戳/版本/可复核/可引用”链路]
\label{cert:priority-pack-strengthens}
证书登记：六件套分别作用于四要件的不同薄弱环节——
\textbf{(1)(3)} 加固时间戳与版本指纹，
\textbf{(2)} 加固版本序列可定位，
\textbf{(4)(5)(6)} 加固可复核性与共同体可检索引用。
因此总体证据链更稳健。
\end{Certificate}

\begin{proof}[证明（证书；谓词因果链）]
由证书~\ref{cert:priority-pack-strengthens}：
六件套对四要件分别给出可核验补强项，合取后使证据链更难被争议推翻；
从而命题~\ref{prop:priority-pack-strengthens} 成立。证毕。
\end{proof}

\section{对话增益总结：从分散直觉到目录化系统}
\label{sec:dialogue-gains-summary}

\begin{remark}[主张（目录化系统收敛）]
本次对话的主要增益可以概括为：把你原先分散的直觉（\PTOPB/\ETOPB、干扰与保障、同伦与上同调、扇区切换、非交换序列、工程逆向设计、对抗博弈触发条件等）
收敛为一套可被实现为“目录 + 库 + 规划器/控制器”的统一框架。
更具体地，增益体现在以下 8 个方面：
\begin{enumerate}
  \item 概念层统一：\ETOPB{} 作为 \PTOPB{} 的增强，且包含“扩边界增强”（第~\ref{sec:etopb-as-ptopb-augmentation}--\ref{sec:two-level-augmentation} 节）；
  \item 分类层澄清：将“外界实时干扰”提升为“实时干扰”（内外统一），并以同伦类组织（第~\ref{sec:disturbance-homotopy-class} 节）；
  \item 数学结构落点：用上同调刻画“可消/不可消”并形成扇区切换触发判据（第~\ref{sec:cohomology-clear-cancel-noncancel}--\ref{sec:relative-cohomology-augmentation} 节）；
  \item 算法框架：两阶段搜索 + 同伦扭曲搬运 + 动态算法选型，形成可运行闭环（第~\ref{sec:twisted-homotopy-two-level-search}--\ref{sec:dynamic-solver-selection} 节）；
  \item 规范结构视角：扇区切换可解释为规范场扇区/修复结构切换，并在跨相处需要贴图搬运（第~\ref{sec:gauge-structure-switching}--\ref{sec:edge-trigger-phase-change-and-algorithm-interface} 节）；
  \item 生命周期与逆向设计：将工程推进表述为边缘算子非交换序列的门槛约束路径规划，并以证据包审计（第~\ref{sec:lifecycle-planning-inverse-design}--\ref{sec:lifecycle-deliverables} 节；参考 \cite{ref_nasa_trl,ref_stage_gate_dtl}）；
  \item 对抗博弈补齐：开放系统下扇区切换守卫条件必须引入对手策略信念、资源可行域与布局脆弱性（第~\ref{sec:openra-guard-three-criteria}--\ref{sec:openra-switch-guard-templates} 节）；
  \item OpenRA 沙盘：用低成本即时战略环境验证“非交换序列 + 守卫条件 + 算法切换”，并可迁移回工程系统（第~\ref{sec:openra-subbundle-noncausal-basis} 节）。
\end{enumerate}
\end{remark}

\begin{definition}[目录化系统（最小交付件集合）]
\label{def:catalog-system-minimal-deliverables}
称一个实现上述框架的“目录化系统”由以下对象组成：
\begin{enumerate}
  \item 扇区目录 $\mathsf{Catalog}$：对每个扇区 $\Sigma$ 给出可观测判据、允许算子与失败模式；
  \item 边缘算子库 $\mathsf{EdgeLib}$：对每个边缘算子给出前置/后置条件、代价/风险、不可交换约束与验证钩子；
  \item 映射表 $\Phi$：将扇区映射为算法栈 $\Phi(\Sigma)=\mathcal{A}_{\Sigma}$（第~\ref{sec:algorithm-layer-sector-to-stack} 节）；
  \item 守卫条件族 $\mathsf{Guard}$：给出 $\Sigma\to\Sigma'$ 的切换触发（对手信念—资源边界—布局脆弱性）及滞回防抖（第~\ref{sec:openra-guard-three-criteria}
    --\ref{sec:openra-switch-guard-templates} 节）；
  \item 门槛策略 $\mathsf{GatePolicy}$：以证据谓词定义成熟度/阶段门的可审计推进规则（第~\ref{sec:lifecycle-gate-policy} 节）。
\end{enumerate}
\end{definition}

\begin{theorem}[公理降级：对话增益收敛为目录化系统的充要对象]
\label{thm:dialogue-gains-to-catalog-system}
在第~\ref{sec:ptopb-etopb-operator-decomposition}--\ref{sec:openra-switch-guard-templates} 节建立的对象与判据基础上，
存在一个目录化系统 $(\mathsf{Catalog},\mathsf{EdgeLib},\Phi,\mathsf{Guard},\mathsf{GatePolicy})$ 满足：
\begin{enumerate}
  \item （可实现性）其对象集足以生成“相内搜索—跨扇区切换—再搜索”的闭环规划/控制器（第~\ref{sec:sectorized-twisted-search-control-loop} 节）；
  \item （可审计性）其每一次扇区切换与算子序列选择均可由证据包与证书化约束追溯（第~\ref{sec:lifecycle-gate-policy} 节）；
  \item （可迁移性）在 OpenRA 沙盘上的守卫条件与算子序列模板可作为工程系统的最小原型（第~\ref{sec:openra-switch-guard-templates} 节）。
\end{enumerate}
\end{theorem}

\begin{Certificate}[数学归纳法证书：按 8 项增益递增构造目录化对象]
\label{cert:dialogue-gains-to-catalog-system}
定义谓词 $P(k)$ 为：“前 $k$ 项增益已被落实为某个目录化对象/映射/守卫/门槛中的一个可计算分量，并与前述分量兼容。”
证书化要求为：
\[
  P(1)\ \text{成立},\qquad
  (\forall k\in\{1,2,\dots,7\})\ \bigl(P(k)\Rightarrow P(k+1)\bigr).
\]
其中“兼容”指：新增分量不会破坏全局可行母空间与扇区化类型一致性（第~\ref{sec:axiom-sector-algorithm-unified-framework} 节）。
\end{Certificate}

\begin{proof}[证明（数学归纳法证书；谓词因果链）]
由证书~\ref{cert:dialogue-gains-to-catalog-system}，逐项说明：
\begin{itemize}
  \item 基步（$k=1$）：由第~\ref{sec:etopb-as-ptopb-augmentation}--\ref{sec:two-level-augmentation} 节，\ETOPB{} 作为 \PTOPB{}
    增强且含“扩边界增强”的两层结构成立；
  该结构为 $\mathsf{Catalog}$ 提供“相内/相间”最小扇区刻画口径，故 $P(1)$ 成立。
  \item 归纳步：假设 $P(k)$ 成立。
  当加入第 $k+1$ 项增益时：
  \begin{itemize}
    \item 若为“实时干扰同伦类”，则由第~\ref{sec:disturbance-homotopy-class} 节可将其登记为 $\mathsf{EdgeLib}$ 的扰动/负算子标签与风险字段；
    \item 若为“上同调可消判据”，则由第~\ref{sec:cohomology-clear-cancel-noncancel}--\ref{sec:relative-cohomology-augmentation}
      节可将其登记为 $\mathsf{Guard}$ 的结构性障碍触发量与修复目标；
    \item 若为“同伦扭曲搬运/动态算法选型”，则由第~\ref{sec:twisted-search-transport} 与第~\ref{sec:dynamic-solver-selection} 节可扩充
      $\Phi(\Sigma)$ 中的搬运算子与求解器簇；
    \item 若为“规范结构/边缘算子切换”，则由第~\ref{sec:gauge-switching-three-levels} 与第~\ref{sec:edge-repair-operator} 节可将其作为
      $\mathsf{EdgeLib}$ 的扇区切换生成元与不可交换约束来源；
    \item 若为“生命周期门槛与证据包”，则由第~\ref{sec:lifecycle-gate-policy} 节可将其登记为 $\mathsf{GatePolicy}$；
    \item 若为“开放系统守卫条件与模板”，则由第~\ref{sec:openra-guard-three-criteria}--\ref{sec:openra-switch-guard-templates} 节可扩充
      $\mathsf{Guard}$ 与最小修复序列模板。
  \end{itemize}
  以上扩充均保持在同一全局母空间与扇区化类型系统内（第~\ref{sec:axiom-sector-algorithm-unified-framework} 节），故 $P(k+1)$ 成立。
\end{itemize}
由归纳得到 $P(8)$ 成立，从而 8 项增益可收敛为定义~\ref{def:catalog-system-minimal-deliverables} 的对象集合；定理成立。
\end{proof}

\subsection{总结性评价：从概念到可运行目录}
\label{sec:dialogue-evaluation}

\begin{lemma}[目录化系统给出闭环可运行的最小接口]
\label{lem:catalog-system-closed-loop-interface}
若存在定义~\ref{def:catalog-system-minimal-deliverables} 的目录化系统，则可构造一个“扇区内搜索/控制 + 守卫触发切换 + 最小修复序列执行 + 证据门槛推进”的闭环过程，
并使得每一步均可被审计追溯为目录/库/映射/门槛中的一个有限记录。
\end{lemma}

\begin{Certificate}[证书：对象—映射—切换规则闭合性检查清单]
\label{cert:catalog-system-closed-loop-interface}
证书由 4 个谓词组成：
\begin{enumerate}
  \item $\mathsf{Obj}$：扇区、算子、算法栈、证据包的类型均在 $\mathsf{Catalog}$ 与 $\mathsf{EdgeLib}$ 的登记域内；
  \item $\mathsf{Map}$：对任意扇区 $\Sigma$，$\Phi(\Sigma)$ 给出相内求解器与搬运/验证接口；
  \item $\mathsf{Switch}$：对任意切换 $\Sigma\to\Sigma'$，$\mathsf{Guard}$ 给出触发与滞回，且模板给出最小修复序列；
  \item $\mathsf{Gate}$：$\mathsf{GatePolicy}$ 给出证据谓词，判定是否推进阶段/成熟度。
\end{enumerate}
若上述谓词同时成立，则闭环过程在对象层面闭合并可审计。
\end{Certificate}

\begin{proof}[证明（证书；谓词因果链）]
由证书~\ref{cert:catalog-system-closed-loop-interface}：
\begin{itemize}
  \item 由 $\mathsf{Obj}$ 可知闭环中每一步引用的对象均有登记域，故“可追溯的索引键”存在；
  \item 由 $\mathsf{Map}$ 得到相内的搜索/控制与搬运接口，因此闭环在扇区内可执行；
  \item 由 $\mathsf{Switch}$ 得到切换触发与最小修复序列，因此闭环在跨扇区处可连续推进且避免抖动；
  \item 由 $\mathsf{Gate}$ 得到证据门槛判据，因此闭环可将“推进”落实为可审计事件，而不是口头判断。
\end{itemize}
四项合取即得到结论：闭环过程可运行且可审计。
\end{proof}

\subsection{下一步最有价值的落地动作：三件最小交付件}
\label{sec:dialogue-next-actions}

\begin{proposition}[三件交付件的最小可运行性]
\label{prop:next-actions-three-deliverables-minimal}
若要将本笔记框架推进为一个最小可运行原型，则至少需要补齐以下三类交付件：
\begin{enumerate}
  \item 扇区目录最小集（建议 5--9 个扇区）：每扇区的可观测判据、允许算子、失败模式；
  \item 边缘算子库最小集：每类格点维度各给出 10--30 个算子条目（含前置/后置、代价/风险、强非交换对、验证钩子）；
  \item 守卫条件指标化：对手策略信念阈值、资源可行域阈值、布局脆弱性指标（暴露度/到达时间差/最小割等）。
\end{enumerate}
并且，在不改变全局公理约束的前提下，上述三件交付件足以派生出定义~\ref{def:catalog-system-minimal-deliverables} 的目录化系统的一个最小实例。
\end{proposition}

\begin{Certificate}[证书：三件交付件到目录化系统五元组的派生规则]
\label{cert:next-actions-three-deliverables-minimal}
证书给出从三件交付件到五元组的构造规则：
\begin{enumerate}
  \item 由扇区目录最小集直接得到 $\mathsf{Catalog}$；
  \item 由边缘算子库最小集直接得到 $\mathsf{EdgeLib}$；
  \item 将“扇区 $\mapsto$ 估计器/求解器/验证链/搬运接口”的默认装配规则登记为 $\Phi$（可从第~\ref{sec:algorithm-layer-sector-to-stack} 节的最小组件抽取）；
  \item 将守卫条件指标化规则登记为 $\mathsf{Guard}$（可从第~\ref{sec:openra-guard-three-criteria} 节的三要件口径抽取）；
  \item 将阶段门/证据谓词登记为 $\mathsf{GatePolicy}$（可从第~\ref{sec:lifecycle-gate-policy} 节抽取；参考 \cite{ref_nasa_trl,
    ref_stage_gate_dtl}）。
\end{enumerate}
\end{Certificate}

\begin{proof}[证明（证书；谓词因果链）]
由证书~\ref{cert:next-actions-three-deliverables-minimal} 的逐项构造规则：
\begin{itemize}
  \item 前两项给出 $\mathsf{Catalog}$ 与 $\mathsf{EdgeLib}$；
  \item 第三项给出 $\Phi$ 的最小装配口径，使每个扇区具备“可执行的相内算法栈”；
  \item 第四项给出 $\mathsf{Guard}$，使闭环具备可计算切换触发与滞回；
  \item 第五项给出 $\mathsf{GatePolicy}$，使推进/验收具备证据谓词与审计接口。
\end{itemize}
因此五元组存在且满足定义~\ref{def:catalog-system-minimal-deliverables}，命题成立。
\end{proof}

\subsection{迁移推论：OpenRA 沙盘到工程系统}
\label{sec:dialogue-transfer-corollary}

\begin{corollary}[沙盘守卫模板的可迁移性推论]
\label{cor:openra-guard-template-transfer}
若某个工程系统可被抽象为扇区化开放系统，并能在其观测口径下定义“对手/环境策略信念、资源可行域、布局/拓扑脆弱性”三类触发量，
则 OpenRA 章给出的守卫条件与最小修复序列模板（第~\ref{sec:openra-guard-three-criteria}--\ref{sec:openra-switch-guard-templates} 节）
可作为该工程系统扇区切换逻辑的一个可用初始模板。
\end{corollary}

\begin{Certificate}[证书：同构替换与接口保持]
\label{cert:openra-guard-template-transfer}
证书由一个“接口保持替换”谓词给出：存在一组替换映射，将 OpenRA 的三要件触发量与算子模板字段
同构替换为工程系统的对应触发量与字段，并保持以下接口不变：
\[
  (\text{触发阈值},\ \text{持续时间},\ \text{滞回},\ \text{最小修复序列},\ \mathsf{NonComm}).
\]
\end{Certificate}

\begin{proof}[证明（证书；谓词因果链）]
由证书~\ref{cert:openra-guard-template-transfer}，若存在接口保持替换映射，则：
\begin{itemize}
  \item 触发量替换保证守卫条件的输入语义在工程系统中仍可计算；
  \item 字段接口保持保证去抖/滞回与最小修复序列模板在结构上不变；
  \item 因而模板可直接作为工程系统的切换逻辑初始版本，并在后续通过证据包校准阈值与不可交换约束。
\end{itemize}
故推论成立。
\end{proof}

\begin{remark}[后续可选动作（按场景生成最小模板）]
若进一步指定一个更关注的沙盘威胁类型或工程场景（例如 OpenRA 的“快空骚扰/装甲推进/偷矿消耗/科技直冲”，或特高压抗太阳风暴节点/量子退相干节点），
则可在上述三件交付件口径下直接给出“扇区目录 + 边缘算子库 + 守卫阈值表”的最小可用版本，便于启动原型规划器/控制器。
\end{remark}

\section{Zenodo DOI 证据链的形式逻辑：公开优先权与统计口径}
\label{sec:zenodo-priority-formal-logic}

\subsection{公理降级：第三方时间戳与不可篡改的优先权固化}
\label{sec:zenodo-priority-axiom-downgrade}

\begin{theorem}[公理降级：Zenodo 记录给出可核验的优先权固化]
\label{thm:zenodo-priority-fixed}
令记录集 $\mathcal{R}=\{r_{\mathrm{PM}},r_{\mathrm{AM1}}\}$，其中：
\begin{itemize}
  \item $r_{\mathrm{PM}}$：Zenodo 记录 17651584（DOI：\code{10.5281/zenodo.17651584}；\code{publication_date}=2025-11-19；状态：
    \code{published}）；
  \item $r_{\mathrm{AM1}}$：Zenodo 记录 17686133（DOI：\code{10.5281/zenodo.17686133}；版本：v1.2；\code{publication_date}
    =2025-11-23；状态：\code{published}）。
\end{itemize}
参见 \cite{GaoZheng2025GFramework,GaoZhengAppVol1}。
若每个 $r\in\mathcal{R}$ 均满足：
\begin{enumerate}
  \item DOI 已发布并可被第三方检索；
  \item 发布后文件与 DOI 不可被修改，更新需生成新版本\cite{ref_zenodo_about_records}，
\end{enumerate}
则“公开优先权固化”不再是主张，而是可被第三方复算核验的事实；
且该事实对任意有限记录集 $|\mathcal{R}|=N$ 成立。
\end{theorem}

\begin{Certificate}[数学归纳法证书：按记录数 $N$ 归纳的优先权固化]
\label{cert:zenodo-priority-fixed}
定义谓词 $P(N)$：“任意 $N$ 个已发布记录均满足 DOI 可检索且版本不可篡改，则其公开优先权可被第三方复算核验。”
证书化要求为：
\[
  P(1)\ \text{成立},\qquad
  (\forall N\in\NN_{\ge 1})\ \bigl(P(N)\Rightarrow P(N+1)\bigr).
\]
\end{Certificate}

\begin{proof}[证明（数学归纳法证书；谓词因果链）]
由证书~\ref{cert:zenodo-priority-fixed}：
\begin{itemize}
  \item 基步：$N=1$ 时，已发布记录给出 DOI 与 \texttt{publication\_date}，
  且发布后不可修改\cite{ref_zenodo_about_records}，因此可被第三方复算核验，$P(1)$ 成立；
  \item 归纳步：假设 $P(N)$ 成立。新增记录 $r_{N+1}$ 满足同一发布与不可篡改规则，
  则 $P(N)$ 的证据链与 $r_{N+1}$ 的证据链并列存在，故 $P(N+1)$ 成立。
\end{itemize}
由数学归纳法得 $P(N)$ 对任意有限 $N$ 成立，从而定理成立。证毕。
\end{proof}

\subsection{定理：Version DOI（版本 DOI）/Concept DOI（概念 DOI）的双口径引用}
\label{sec:zenodo-doi-versioning}

\begin{theorem}[Version DOI（版本 DOI）与 Concept DOI（概念 DOI）的双口径引用规则]
\label{thm:zenodo-doi-versioning}
令 $d_{\mathrm{v}}$ 表示 Version DOI（版本 DOI），$d_{\mathrm{c}}$ 表示 Concept DOI（概念 DOI）。
对 $r_{\mathrm{PM}}$ 与 $r_{\mathrm{AM1}}$，其概念 DOI 分别为
\texttt{10.5281/zenodo.17651583} 与 \texttt{10.5281/zenodo.17683849}\cite{GaoZheng2025GFramework,GaoZhengAppVol1}。
若 DOI 版本规则满足：“版本 DOI 锁定具体版本，概念 DOI 指向该概念的最新版本”\cite{ref_zenodo_doi_versioning}，
则对外引用口径可标准化为：
\begin{enumerate}
  \item 引用具体定义/定理/符号时使用 $d_{\mathrm{v}}$（锁定版本）；
  \item 引用整卷或持续更新体系时附带 $d_{\mathrm{c}}$（指向最新）。
\end{enumerate}
从而避免“版本漂移”导致的引用歧义。
\end{theorem}

\begin{Certificate}[证书：版本/概念 DOI 的元数据可追溯性]
\label{cert:zenodo-doi-versioning}
证书登记如下谓词链：
\[
\begin{aligned}
  A_1 &: \text{API 元数据含}\ \text{\texttt{conceptdoi}}\ \text{字段},\\
  A_2 &: d_{\mathrm{v}}\ \text{解析到具体版本},\ d_{\mathrm{c}}\ \text{解析到最新版本}\cite{ref_zenodo_doi_versioning},\\
  A_3 &: \text{据此可对“版本锁定/概念更新”分别给出引用}.
\end{aligned}
\]
\end{Certificate}

\begin{proof}[证明（谓词因果链）]
由证书~\ref{cert:zenodo-doi-versioning}：
$A_1$ 给出概念 DOI 的可追溯性，
$A_2$ 给出版本 DOI 与概念 DOI 的解析语义，
$A_1\wedge A_2\Rightarrow A_3$，因此双口径引用规则成立。证毕。
\end{proof}

\subsection{引理：访问/下载统计口径的一致性与去重规则}
\label{sec:zenodo-stats-consistency}

\begin{lemma}[统计口径一致性与去重解释]
\label{lem:zenodo-stats-consistency}
对 $r_{\mathrm{PM}}$ 与 $r_{\mathrm{AM1}}$，其 API 统计字段满足：
\[
\begin{aligned}
  r_{\mathrm{PM}} &: \text{\texttt{downloads}}=57,\ \text{\texttt{unique\_downloads}}=36,\ \text{\texttt{views}}=84,\
    \text{\texttt{unique\_views}}=57;\\
  r_{\mathrm{AM1}} &: \text{\texttt{downloads}}=49,\ \text{\texttt{unique\_downloads}}=23,\ \text{\texttt{views}}=55,\
    \text{\texttt{unique\_views}}=32.
\end{aligned}
\]
且 Zenodo 对 \texttt{unique\_*} 采用 1 小时窗口去重，并对概念 DOI 展示全版本聚合统计\cite{ref_zenodo_statistics}。
因此据 \texttt{downloads}/\texttt{views} 计算的转化率属于“全版本聚合口径”；
若需评估 v1.2 的发布增量，应使用 \texttt{version\_downloads}/\texttt{version\_views} 口径。
\end{lemma}

\begin{Certificate}[证书：统计字段与去重规则]
\label{cert:zenodo-stats-consistency}
证书登记：
\[
\begin{aligned}
  S_1 &: \text{记录的 API }\text{\texttt{stats}}\ \text{字段给出}\ \text{\texttt{downloads/views/unique\_*}},\\
  S_2 &: \text{\texttt{unique\_*}}\ \text{按 1 小时窗口去重，概念 DOI 为聚合口径}\cite{ref_zenodo_statistics},\\
  S_3 &: \text{版本字段}\ \text{\texttt{version\_*}}\ \text{对应当前版本统计}.
\end{aligned}
\]
\end{Certificate}

\begin{proof}[证明（谓词因果链）]
由证书~\ref{cert:zenodo-stats-consistency}：
$S_1$ 给出数值口径，$S_2$ 给出去重与聚合规则，$S_3$ 给出版本口径。
由 $S_1\wedge S_2$ 得“聚合转化率”的解释；
由 $S_3$ 得“版本增量”口径的区分性。证毕。
\end{proof}

\subsection{命题：统计口径双报表消除版本迁移混淆}
\label{sec:zenodo-stats-dual-reporting}

\begin{proposition}[双口径统计报表的可复算性]
\label{prop:zenodo-stats-dual-reporting}
若对外报告同时给出概念聚合统计
（\code{downloads}、\code{views}、\code{unique_*}）
与当前版本统计（\code{version_downloads}、\code{version_views}、\code{version_unique_*}），
则可同时刻画“全历史累计”与“当前版本增量”，从而避免版本迁移导致的口径混淆。
\end{proposition}

\begin{Certificate}[证书：双口径字段的可比性]
\label{cert:zenodo-stats-dual-reporting}
证书登记谓词链：
\[
\begin{aligned}
  Y_1 &: \text{概念聚合统计覆盖所有版本},\\
  Y_2 &: \text{版本统计仅覆盖当前版本},\\
  Y_3 &: Y_1\wedge Y_2\Rightarrow \text{可同时给出累计与增量口径}.
\end{aligned}
\]
\end{Certificate}

\begin{proof}[证明（谓词因果链）]
由证书~\ref{cert:zenodo-stats-dual-reporting}：
$Y_1$ 与 $Y_2$ 分别固定“历史累计/当前版本”的口径边界，
$Y_1\wedge Y_2\Rightarrow Y_3$，故命题成立。证毕。
\end{proof}

\subsection{推论：优先权固化与防御性公开的双重表述}
\label{sec:zenodo-priority-corollary}

\begin{corollary}[优先权与防御性公开表述]
\label{cor:zenodo-priority-defensive}
由定理~\ref{thm:zenodo-priority-fixed} 可得以下可核验表述：
\begin{enumerate}
  \item 公开优先权：基础定义与框架已于 2025-11-19（\texttt{10.5281/zenodo.17651584}）与
  2025-11-23（\texttt{10.5281/zenodo.17686133}，v1.2）公开并被冻结，第三方可复算核验；
  \item 防御性公开：上述 DOI 记录构成公开披露证据链，可压缩第三方“更早/独占”主张的事实空间
  （法律效力需按法域评估）。
\end{enumerate}
\end{corollary}

\begin{Certificate}[证书：公开披露与防御性公开的蕴含链]
\label{cert:zenodo-priority-defensive}
证书登记谓词链：
\[
\begin{aligned}
  C_1 &: \text{第三方时间戳 + DOI 可检索 + 版本冻结},\\
  C_2 &: \Rightarrow\ \text{公开披露可复算核验},\\
  C_3 &: \Rightarrow\ \text{形成防御性公开证据链}.
\end{aligned}
\]
\end{Certificate}

\begin{proof}[证明（谓词因果链）]
由证书~\ref{cert:zenodo-priority-defensive}：
$C_1\Rightarrow C_2$ 给出公开披露的可复算性，
$C_2\Rightarrow C_3$ 给出防御性公开的证据链语义。
因此推论成立。证毕。
\end{proof}

\subsection{备注：五项加固动作（不增加应用泄露风险）}
\label{sec:zenodo-priority-remarks}

\begin{remark}[加固动作清单]
不增加应用泄露风险的前提下，可执行如下加固动作：
\begin{enumerate}
  \item 在每卷首页/README（说明文件）增加 “How to cite（引用方式，版本 DOI + 概念 DOI）”，并注明“引用具体定理请用版本 DOI”；
  \item 在 Zenodo 描述区置顶版本变更摘要（保留 v1.1/v1.2 级别变更说明）；
  \item 在文档首页或 Zenodo 描述补充 PDF/TeX 的 SHA256 指纹；
  \item 为 PureMath 与 Vol1 互相添加 Related identifiers / references（关联标识/参考）交叉引用；
  \item 统计报表同时给出聚合口径与当前版本口径，避免口径混淆。
\end{enumerate}
\end{remark}

\section{开放系统运行期在线控制：JacGap 与修复塔的必要性、最小性与投影性}
\label{sec:open-system-runtime-control}

\subsection{形式化设定与定义}
\label{sec:open-system-runtime-formalization}

\begin{definition}[开放系统博弈的运行期在线控制抽象]
令 $x_t\in\mathcal{X}$ 为物理/工程状态，$\omega_t\in\Omega$ 为对手/扰动输入，$u_t\in\mathcal{U}_t$ 为我方控制输入。
将“规则对象”显式化为 $L_t\in\mathsf{Law}$，允许其随运行期演化：
\[
  x_{t+1}=F(x_t,u_t,\omega_t;L_t),
  \qquad
  L_{t+1}=G(L_t,x_t,u_t,\omega_t).
\]
其中 $L_t$ 约束/参数化转移算子、可行动作与可审计边界；$G$ 描述开放系统的规则漂移。
\end{definition}

\begin{definition}[规则一致性证书：Jacobiator 与 JacGap]
给定 $L\in\mathsf{Law}$，令 $\mathcal{O}_L$ 为其操作空间，定义 Jacobiator
\[
  J_L:\mathcal{O}_L^{\otimes 3}\to\mathcal{O}_L.
\]
取非负连续泛函 $\phi_k:\mathcal{O}_L\to\RR_{\ge 0}$，定义
\[
  \JacGap_L(a,b,c):=\sum_{k=1}^m \phi_k(J_L(a,b,c)).
\]
若存在常数 $C>0$ 使得
\[
  \|J_L(a,b,c)\|\le C\bigl(1+\JacGap_L(a,b,c)\bigr),
\]
则称 $\JacGap_L$ 为可计算的规则一致性证书。
并要求 $\JacGap_L(a,b,c)=0\Rightarrow J_L(a,b,c)=0$（分离性）。
\end{definition}

\begin{definition}[绩效--一致性联合目标]
令 $\mathcal{J}_{\mathrm{perf}}(u)$ 表示绩效目标（效能/风险/成本），
并定义运行期一致性度量 $\JacGap_{\mathrm{tot}}(t;u)$（由 $\JacGap_L$ 在运行期观测点聚合而得）。
称问题为“绩效--一致性联合优化”，若满足以下任一口径：
\[
  \min_{u(\cdot)}\ \mathcal{J}_{\mathrm{perf}}(u)
  \quad\text{s.t.}\quad
  \sup_{t\in[0,T]}\JacGap_{\mathrm{tot}}(t;u)\le \varepsilon,
\]
或
\[
  \min_{u(\cdot)}\ \mathcal{J}_{\mathrm{perf}}(u)
  +\lambda\sup_{t\in[0,T]}\JacGap_{\mathrm{tot}}(t;u),
  \qquad \lambda>0.
\]
\end{definition}

\begin{definition}[分层修复塔与烈度度量]
给定允许变分族的分层嵌套
\[
  \mathcal{U}_0\subset \mathcal{U}_1\subset\cdots\subset\mathcal{U}_N,
\]
定义“烈度证书” $\tau(u)\in\RR_{\ge 0}^d$，其加权强度为
\[
  |\tau(u)|_w:=\sum_{i=1}^d w_i\tau_i(u),\qquad w_i>0.
\]
若采用字典序目标
\[
  \min\Bigl(\sup_{t\in[0,T]}\JacGap_{\mathrm{tot}}(t;u),\ |\tau(u)|_w\Bigr),
\]
则称该优化为“修复塔的最小烈度选择”。
\end{definition}

\subsection{必要性：绩效与规则一致性证书必须并行优化}
\label{sec:open-system-necessity}

\begin{lemma}[括号化敏感性由 Jacobiator 控制]
\label{lem:jacobiator-bracketing-sensitivity}
设存在表示 $\rho_L:\mathcal{O}_L\to \mathrm{End}(\mathcal{X})$，并存在常数 $K>0$ 使
$\|\rho_L(o)\|\le K\|o\|$ 对任意 $o\in\mathcal{O}_L$ 成立。
令
\[
  F_{(ab)c}(x):=\rho_L((a\cdot b)\cdot c)(x),
  \qquad
  F_{a(bc)}(x):=\rho_L(a\cdot(b\cdot c))(x),
\]
则对任意 $x\in\mathcal{X}$ 有
\[
  \|F_{(ab)c}(x)-F_{a(bc)}(x)\|
  \le K\|J_L(a,b,c)\|.
\]
\end{lemma}

\begin{Certificate}[证书：表示有界与 Jacobiator 关联]
\label{cert:jacobiator-bracketing-sensitivity}
证书登记谓词链：
\[
\begin{aligned}
  P_1 &: \exists\,\rho_L\ \text{使}\ \|\rho_L(o)\|\le K\|o\|,\\
  P_2 &: (a\cdot b)\cdot c-a\cdot(b\cdot c)=J_L(a,b,c),\\
  P_3 &: \|F_{(ab)c}(x)-F_{a(bc)}(x)\|=\|\rho_L(J_L(a,b,c))(x)\|.
\end{aligned}
\]
\end{Certificate}

\begin{proof}[证明（谓词因果链）]
由证书~\ref{cert:jacobiator-bracketing-sensitivity}：
由 $P_2$ 与 $P_3$ 得差异项等于 $\rho_L(J_L(a,b,c))$ 的作用；
由 $P_1$ 得 $\|\rho_L(J_L(a,b,c))\|\le K\|J_L(a,b,c)\|$，
故结论成立。证毕。
\end{proof}

\begin{lemma}[硬约束下不控制 JacGap 将导致语义不闭合]
\label{lem:safety-nonclosure-without-jacgap}
设安全约束集 $S\subset\mathcal{X}$，要求 $x_t\in S$ 在运行期保持。
若存在 $a,b,c\in\mathcal{O}_L$ 与 $x\in S$ 使
\[
  \|\rho_L(J_L(a,b,c))(x)\|>\mathrm{dist}(x,\partial S),
\]
则存在两种实现（对应两种括号化/调度路径）使得
一条轨迹保持在 $S$ 内，另一条轨迹跨出 $S$。
\end{lemma}

\begin{Certificate}[证书：安全裕度被括号化差异跨越]
\label{cert:safety-nonclosure-without-jacgap}
证书登记谓词链：
\[
\begin{aligned}
  Q_1 &: \mathrm{dist}(x,\partial S)>0,\ \text{为可容忍裕度},\\
  Q_2 &: \|F_{(ab)c}(x)-F_{a(bc)}(x)\|> \mathrm{dist}(x,\partial S),\\
  Q_3 &: \Rightarrow\ \text{存在一条实现跨出 }S.
\end{aligned}
\]
\end{Certificate}

\begin{proof}[证明（谓词因果链）]
由引理~\ref{lem:jacobiator-bracketing-sensitivity} 得
$F_{(ab)c}(x)-F_{a(bc)}(x)=\rho_L(J_L(a,b,c))(x)$。
若满足引理假设，则 $Q_2$ 成立；结合 $Q_1$ 得 $Q_3$，
从而存在一条实现跨出 $S$。证毕。
\end{proof}

\begin{theorem}[必要性定理：开放系统强约束下必须纳入 JacGap]
\label{thm:necessity-jacgap}
若满足以下条件：
\begin{enumerate}
  \item 开放系统：$L_t$ 可随扰动/对手/内部漂移演化；
  \item 运行期存在不可避免的实现差异（并发/调度/延迟），可抽象为括号化差异；
  \item 要求强安全约束与可审计复算（语义不依赖实现细节），
\end{enumerate}
则任何仅优化绩效而不约束/惩罚 $\JacGap_{\mathrm{tot}}$ 的控制问题在该口径下是不适定的；
必须将 $\JacGap_{\mathrm{tot}}$ 纳入目标或约束。
\end{theorem}

\begin{Certificate}[证书：强约束 + 实现差异 + 开放系统]
\label{cert:necessity-jacgap}
证书登记谓词链：
\[
\begin{aligned}
  R_1 &: \text{开放系统可进入高 JacGap 区域},\\
  R_2 &: \text{实现差异}\Rightarrow\text{括号化差异不可避免},\\
  R_3 &: \text{硬约束与可审计要求} \Rightarrow \text{语义需对实现不变}.
\end{aligned}
\]
\end{Certificate}

\begin{proof}[证明（谓词因果链）]
由证书~\ref{cert:necessity-jacgap}：
若不控制 $\JacGap_{\mathrm{tot}}$，则证书无法排除 $J_L$ 的大偏差；
结合 $R_1$ 与 $R_2$ 得“高 JacGap 情形必然暴露实现差异”；
由引理~\ref{lem:safety-nonclosure-without-jacgap} 得在高 JacGap 下安全性依赖实现细节；
结合 $R_3$，得“仅优化绩效”不足以满足强约束与可审计要求。
因此必须纳入 $\JacGap_{\mathrm{tot}}$。证毕。
\end{proof}

\subsection{最小性：同伦扭曲与修复塔的最小烈度路径}
\label{sec:open-system-minimality}

\begin{definition}[可行层级与最小可行层级]
给定容差 $\varepsilon>0$，定义第 $k$ 层可行集
\[
  \mathcal{F}_k(\varepsilon)
  :=
  \bigl\{u(\cdot)\in\mathcal{U}_k:\ \sup_{t\in[0,T]}\JacGap_{\mathrm{tot}}(t;u)\le\varepsilon,\ x_t\in S\bigr\}.
\]
若 $\mathcal{F}_N(\varepsilon)\neq\varnothing$，定义最小可行层级
\[
  k^\star:=\min\{k:\ \mathcal{F}_k(\varepsilon)\neq\varnothing\}.
\]
\end{definition}

\begin{lemma}[可行集单调性]
\label{lem:feasible-set-monotone}
对任意 $k$，有 $\mathcal{F}_k(\varepsilon)\subseteq\mathcal{F}_{k+1}(\varepsilon)$。
\end{lemma}

\begin{Certificate}[证书：嵌套可行域]
\label{cert:feasible-set-monotone}
证书登记：$\mathcal{U}_k\subseteq\mathcal{U}_{k+1}$ 且约束与容差条件对层级无偏置。
\end{Certificate}

\begin{proof}[证明（谓词因果链）]
由证书~\ref{cert:feasible-set-monotone}，
任意 $u(\cdot)\in\mathcal{F}_k(\varepsilon)$ 同时属于 $\mathcal{U}_{k+1}$，
且满足相同容差与安全约束，故 $\mathcal{F}_k(\varepsilon)\subseteq\mathcal{F}_{k+1}(\varepsilon)$。证毕。
\end{proof}

\begin{theorem}[公理降级：有限层修复塔的最小可行层级存在]
\label{thm:min-feasible-layer-exists}
设修复塔为有限层数 $N$。若 $\mathcal{F}_N(\varepsilon)\neq\varnothing$，则存在最小可行层级 $k^\star$，
且满足 $\mathcal{F}_{k^\star}(\varepsilon)\neq\varnothing$ 与 $\mathcal{F}_{k^\star-1}(\varepsilon)=\varnothing$（约定
  $\mathcal{F}_{-1}(\varepsilon)=\varnothing$）。
\end{theorem}

\begin{Certificate}[数学归纳法证书：按层数 $N$ 归纳]
\label{cert:min-feasible-layer-exists}
定义谓词 $P(N)$：“任意 $N$ 层修复塔若顶层可行，则存在最小可行层级 $k^\star$。”
证书化要求为：
\[
  P(1)\ \text{成立},\qquad
  (\forall N\in\NN_{\ge 1})\ \bigl(P(N)\Rightarrow P(N+1)\bigr).
\]
\end{Certificate}

\begin{proof}[证明（数学归纳法证书；谓词因果链）]
由证书~\ref{cert:min-feasible-layer-exists}：
\begin{itemize}
  \item 基步：$N=1$ 时，若 $\mathcal{F}_1(\varepsilon)\neq\varnothing$，
  则 $k^\star=1$，结论成立；
  \item 归纳步：假设 $P(N)$ 成立。对 $(N+1)$ 层修复塔：
  若 $\mathcal{F}_N(\varepsilon)\neq\varnothing$，由归纳假设存在最小 $k^\star\le N$；
  若 $\mathcal{F}_N(\varepsilon)=\varnothing$ 而 $\mathcal{F}_{N+1}(\varepsilon)\neq\varnothing$，
  则 $k^\star=N+1$。因此 $P(N+1)$ 成立。
\end{itemize}
由数学归纳法得对任意有限 $N$ 命题成立。证毕。
\end{proof}

\begin{theorem}[最小烈度定理：字典序目标的最小层级实现]
\label{thm:lexicographic-min-intensity}
若烈度度量满足可行层级下界单调性：
\[
  \inf_{u\in\mathcal{F}_k(\varepsilon)}|\tau(u)|_w
  \le
  \inf_{u\in\mathcal{F}_{k+1}(\varepsilon)}|\tau(u)|_w,
\]
则在字典序目标
\[
  \min\Bigl(\sup_{t\in[0,T]}\JacGap_{\mathrm{tot}}(t;u),\ |\tau(u)|_w\Bigr)
\]
下，任意最优解均可选择在最小可行层级 $k^\star$ 上实现，从而达到最小烈度。
\end{theorem}

\begin{Certificate}[证书：最小层级与烈度单调]
\label{cert:lexicographic-min-intensity}
证书登记谓词链：
\[
\begin{aligned}
  M_1 &: \mathcal{F}_{k^\star}(\varepsilon)\neq\varnothing,\ \mathcal{F}_{k^\star-1}(\varepsilon)=\varnothing,\\
  M_2 &: \inf_{u\in\mathcal{F}_k(\varepsilon)}|\tau(u)|_w\ \text{随}\ k\ \text{不减},\\
  M_3 &: \text{字典序优先保证 }\JacGap_{\mathrm{tot}}\ \text{容差}.
\end{aligned}
\]
\end{Certificate}

\begin{proof}[证明（谓词因果链）]
由证书~\ref{cert:lexicographic-min-intensity}：
$M_1$ 给出最小可行层级；
由 $M_3$，任意可行最优解必须满足容差要求，因此属于某一 $\mathcal{F}_k(\varepsilon)$；
若取 $k>k^\star$，由 $M_2$ 得其烈度不小于 $k^\star$ 层可行解。
因此字典序最优解可选在 $k^\star$ 层，且烈度最小。证毕。
\end{proof}

\begin{proposition}[同伦扭曲/修复完成的最小修复数据（同伦等价口径）]
\label{prop:homotopy-twist-min-repair}
设存在同伦扭曲族 $\{L_\lambda\}_{\lambda\in[0,1]}$ 使 $L_0=L$ 且 $L_1$ 满足 $\JacGap$ 在证书集上最小，
并且 $L_1$ 具有如下最小性：任何达到同等或更低 $\JacGap$ 的修复系统 $L'$，
在算子同伦意义下都可因子化到 $L_1$。
则 $L_1$ 所对应的修复数据给出“最小烈度路径”的同伦等价代表元。
\end{proposition}

\begin{Certificate}[证书：最小修复数据的同伦因子化]
\label{cert:homotopy-twist-min-repair}
证书登记谓词链：
\[
\begin{aligned}
  N_1 &: \exists\ \{L_\lambda\}\ \text{使}\ \JacGap\ \text{单调不增},\\
  N_2 &: L_1\ \text{为证书集上的 gap-最小元},\\
  N_3 &: L'\ \text{若不劣于}\ L_1,\ \text{则 }L'\ \text{可因子化到}\ L_1.
\end{aligned}
\]
\end{Certificate}

\begin{proof}[证明（谓词因果链）]
由证书~\ref{cert:homotopy-twist-min-repair}：
$N_1$ 给出可行同伦路径，$N_2$ 给出目标最小性，
$N_3$ 给出“任何更优修复皆可因子化到 $L_1$”的结构性约束；
因此 $L_1$ 的修复数据在同伦等价意义下为最小代表元。证毕。
\end{proof}

\subsection{投影性：JacGap=0 或修复塔截断的退化}
\label{sec:open-system-projection}

\begin{theorem}[JacGap=0 的严格化投影]
\label{thm:jacgap-zero-strictification}
若在关注时间窗内 $\JacGap_{\mathrm{tot}}(t)=0$，且对应算子层满足严格化可实施条件，
则存在投影 $P:\mathsf{Law}\to\mathsf{Law}_{\mathrm{strict}}$ 使
\[
  J_{P(L)}=0,
\]
从而运行期问题退化为固定规则或准固定规则的经典控制/对策问题。
\end{theorem}

\begin{Certificate}[证书：JacGap 零与严格化]
\label{cert:jacgap-zero-strictification}
证书登记谓词链：
\[
\begin{aligned}
  S_1 &: \JacGap_{\mathrm{tot}}(t)=0\ \text{且分离性成立}\Rightarrow J_L=0,\\
  S_2 &: J_L=0\Rightarrow \text{二元结构闭合且括号化无差异},\\
  S_3 &: \Rightarrow\ \exists P(L)\ \text{实现严格化}.
\end{aligned}
\]
\end{Certificate}

\begin{proof}[证明（谓词因果链）]
由证书~\ref{cert:jacgap-zero-strictification}：
$S_1$ 给出 Jacobiator 消失，
$S_2$ 得括号化差异消除，
由 $S_3$ 得投影 $P(L)$ 的存在，从而退化为经典固定规则问题。证毕。
\end{proof}

\begin{proposition}[修复塔截断即为可行域投影]
\label{prop:repair-tower-truncation-projection}
若仅允许 $\mathcal{U}_k$ 内动作，则原问题等价于在可行域上作投影/截断：
\[
  \mathcal{U}_N\ \longrightarrow\ \mathcal{U}_k.
\]
因此仅在低层动作空间中求解的算法可视为“修复塔截断下的投影方法”。
\end{proposition}

\begin{Certificate}[证书：截断可行域与求解等价]
\label{cert:repair-tower-truncation-projection}
证书登记谓词链：
\[
\begin{aligned}
  T_1 &: \text{原问题的可行域为 }\mathcal{U}_N,\\
  T_2 &: \text{仅允许}\ \mathcal{U}_k\ \text{等价于}\ \mathcal{U}_N\ \text{投影},\\
  T_3 &: \Rightarrow\ \text{解法退化为低层动作空间的最优化}.
\end{aligned}
\]
\end{Certificate}

\begin{proof}[证明（谓词因果链）]
由证书~\ref{cert:repair-tower-truncation-projection}：
$T_1$ 给出原始可行域，$T_2$ 给出截断等价于投影，
从而 $T_3$ 成立，命题得证。证毕。
\end{proof}

\begin{corollary}[现有方法的投影性归纳]
\label{cor:existing-methods-projection}
若 $\JacGap_{\mathrm{tot}}=0$ 或修复塔被截断至低层动作空间，
则多类传统控制/博弈/在线学习方法可被视为本框架在“严格化/投影”条件下的特例。
\end{corollary}

\begin{Certificate}[证书：严格化与截断的归并]
\label{cert:existing-methods-projection}
证书登记谓词链：
\[
\begin{aligned}
  U_1 &: \JacGap_{\mathrm{tot}}=0\Rightarrow \text{严格化投影成立},\\
  U_2 &: \text{修复塔截断}\Rightarrow \text{可行域投影成立},\\
  U_3 &: U_1\ \text{或}\ U_2\Rightarrow \text{现有方法为投影特例}.
\end{aligned}
\]
\end{Certificate}

\begin{proof}[证明（谓词因果链）]
由证书~\ref{cert:existing-methods-projection}：
$U_1$ 与 $U_2$ 分别给出严格化与截断两类投影，
由 $U_3$ 得现有方法可归入投影特例。证毕。
\end{proof}

\subsection{备注：可检验假设与证书接口}
\label{sec:open-system-remarks}

\begin{remark}[形式化证明所需的可检验前提]
为使上述证书链条在工程上可复验，建议显式记录以下可检验前提：
\begin{enumerate}
  \item 表示假设：$\rho_L$ 的有界性/ Lipschitz 常数 $K$ 的估计来源（实验、仿真或审计日志）；
  \item 烈度单调假设：高层动作在 $\tau$ 指标上不可被低层动作抵消的验证规则；
  \item 证书数据接口：$\JacGap$ 估计所依赖的可观测三元组与采样窗口。
\end{enumerate}
这些前提不改变结论结构，但决定“必要性与最小性”能否在具体系统上被复算与审计。
\end{remark}

\section{回合语义与上同调闭合：曲率/Gap 驱动的窗口判据}
\label{sec:round-cohomology-window}

\subsection{形式化设定：回合窗与同调归约}
\label{sec:round-window-formalization}

\begin{definition}[回合窗（round window）]
\label{def:round-window}
给定时间窗 $I=[t,t+\Delta]$，称其为一个 $(\varepsilon,\kappa)$-回合窗，若满足：
\begin{enumerate}
  \item 曲率可压缩：对任意端点相同的 law-trajectory $\gamma_1,\gamma_2$，
  其 holonomy 差异满足
  \[
    \|H(\gamma_1)-H(\gamma_2)\|
    \le \kappa\cdot \operatorname{Area}(\gamma_1\cup\gamma_2),
  \]
  且 $\kappa\cdot \operatorname{Area}_{\max}(I)\ll 1$；
  \item 代数可压缩：在该窗内发生的操作三元组满足
  \[
    \JacGap_{L_s}(a,b,c)\le \varepsilon,\qquad \forall s\in I;
  \]
  \item 同调归约可闭合：存在归约算子 $\mathcal{R}_I$，对窗内事件链 $E_I$ 满足
  \[
    E_I-E_I'=\partial B\ \Rightarrow\ \mathcal{R}_I(E_I)=\mathcal{R}_I(E_I'),
  \]
  即“只依赖同调类的商化结果”可以在窗末端稳定给出。
\end{enumerate}
该定义不依赖“外部信息是否持续加入”，只依赖曲率与 JacGap 口径下的可闭合性。
\end{definition}

\begin{theorem}[公理降级：有限回合闭合的可拼接性]
\label{thm:round-window-concatenation}
若存在连续的回合窗序列 $\{I_j\}_{j=0}^{N-1}$，每个 $I_j$ 均为 $(\varepsilon,\kappa)$-回合窗，
且其归约算子满足相容性
\[
  \mathcal{R}_{I_{j+1}}\circ \mathcal{R}_{I_j}=\mathcal{R}_{I_{j:j+1}},
\]
则拼接窗口 $I=\bigcup_{j=0}^{N-1} I_j$ 也为 $(\varepsilon,\kappa)$-回合窗，
其归约算子可取为 $\mathcal{R}_I=\mathcal{R}_{I_{N-1}}\circ\cdots\circ\mathcal{R}_{I_0}$。
\end{theorem}

\begin{Certificate}[数学归纳法证书：按回合数 $N$ 归纳]
\label{cert:round-window-concatenation}
定义谓词 $P(N)$：“任意 $N$ 个相容回合窗的拼接仍为回合窗，且归约算子可复合。”
证书化要求为：
\[
  P(1)\ \text{成立},\qquad
  (\forall N\in\NN_{\ge 1})\ \bigl(P(N)\Rightarrow P(N+1)\bigr).
\]
\end{Certificate}

\begin{proof}[证明（数学归纳法证书；谓词因果链）]
由证书~\ref{cert:round-window-concatenation}：
\begin{itemize}
  \item 基步：$N=1$ 时命题显然成立；
  \item 归纳步：假设 $P(N)$ 成立。若第 $N+1$ 个窗与前 $N$ 个窗相容，
  则由相容性得到 $\mathcal{R}_{I_{0:N}}$ 可复合，并保持同调归约稳定性；
  因此 $P(N+1)$ 成立。
\end{itemize}
由数学归纳法得定理成立。证毕。
\end{proof}

\subsection{引理：曲率与 JacGap 给出回合窗上限}
\label{sec:round-window-lemmas}

\begin{lemma}[曲率控制回合窗可压缩性]
\label{lem:curvature-bounds-round-window}
设 $\GZLCurv$ 在时间窗 $I$ 上有界。则对任意小闭环 $C\subset I$，
其 holonomy 满足估计
\[
  \|H(C)-\mathrm{id}\|
  \lesssim
  \sup_{s\in I}\|\GZLCurv(s)\|\cdot \operatorname{Area}(C).
\]
因此若 $\sup_{s\in I}\|\GZLCurv(s)\|\cdot \operatorname{Area}_{\max}(I)\le \delta$，
则 $I$ 上的路径压缩误差可被 $\delta$ 控制。
\end{lemma}

\begin{Certificate}[证书：曲率--holonomy 估计]
\label{cert:curvature-bounds-round-window}
证书登记谓词链：
\[
\begin{aligned}
  H_1 &: \text{存在局部 holonomy--曲率估计},\\
  H_2 &: \|H(C)-\mathrm{id}\|\ \text{由}\ \int_C \GZLCurv\ \text{控制},\\
  H_3 &: \Rightarrow\ \text{可用 }\sup\|\GZLCurv\|\cdot \operatorname{Area}(C)\ \text{界定}.
\end{aligned}
\]
\end{Certificate}

\begin{proof}[证明（谓词因果链）]
由证书~\ref{cert:curvature-bounds-round-window}：
$H_1$ 与 $H_2$ 给出局部曲率控制；
取上界 $\sup_{s\in I}\|\GZLCurv(s)\|$ 并用面积估计得到 $H_3$，
从而引理成立。证毕。
\end{proof}

\begin{lemma}[JacGap 控制括号化语义差异]
\label{lem:jacgap-bounds-bracketing}
设存在表示 $\rho_L:\mathcal{O}_L\to \mathrm{End}(\mathcal{X})$ 与常数 $K>0$，
对任意 $a,b,c\in\mathcal{O}_L$ 与 $x\in\mathcal{X}$ 有
\[
  \|\rho_L(J_L(a,b,c))(x)\|\le K\|J_L(a,b,c)\|.
\]
若在窗内满足 $\JacGap_{L_s}(a,b,c)\le\varepsilon$，
则存在常数 $K'$ 使
\[
  \|\text{Outcome}((ab)c;x)-\text{Outcome}(a(bc);x)\|
  \le K'(1+\varepsilon).
\]
\end{lemma}

\begin{Certificate}[证书：JacGap 与语义不变性]
\label{cert:jacgap-bounds-bracketing}
证书登记谓词链：
\[
\begin{aligned}
  J_1 &: \|\rho_L(J_L(a,b,c))(x)\|\le K\|J_L(a,b,c)\|,\\
  J_2 &: \|J_L(a,b,c)\|\le C(1+\JacGap_{L_s}(a,b,c)),\\
  J_3 &: \Rightarrow\ \text{括号化差异由}\ \varepsilon\ \text{控制}.
\end{aligned}
\]
\end{Certificate}

\begin{proof}[证明（谓词因果链）]
由证书~\ref{cert:jacgap-bounds-bracketing}：
$J_1$ 与 $J_2$ 给出差异项上界；
结合 $\JacGap_{L_s}\le\varepsilon$ 得到 $J_3$，从而引理成立。证毕。
\end{proof}

\subsection{定理：封闭/开放系统的回合尺度}
\label{sec:round-window-theorems}

\begin{theorem}[封闭（strict-flat）$\Rightarrow$ 可大回合]
\label{thm:closed-system-large-round}
若在时间窗 $I$ 上满足：
\begin{enumerate}
  \item $\sup_{s\in I}\|\GZLCurv(s)\|$ 足够小，使曲率压缩误差可忽略；
  \item $\JacGap_{L_s}\le\varepsilon$，并且 $\varepsilon$ 小到允许同调归约闭合；
  \item 对抗交互点由规则冻结为离散时刻，
\end{enumerate}
则存在 $\Delta$ 可大于计算延迟的回合窗，使得
运行期交互可被压缩为“端点 + 一次同调归约”的回合语义。
\end{theorem}

\begin{Certificate}[证书：封闭系统的回合闭合]
\label{cert:closed-system-large-round}
证书登记谓词链：
\[
\begin{aligned}
  C_1 &: \text{曲率压缩满足}\ \delta\ \text{阈值},\\
  C_2 &: \JacGap_{L_s}\le\varepsilon\ \Rightarrow\ \text{括号化语义稳定},\\
  C_3 &: \text{交互冻结}\Rightarrow \text{回合端点可定义}.
\end{aligned}
\]
\end{Certificate}

\begin{proof}[证明（谓词因果链）]
由引理~\ref{lem:curvature-bounds-round-window} 与 $C_1$ 得路径压缩误差可控；
由引理~\ref{lem:jacgap-bounds-bracketing} 与 $C_2$ 得语义不变性可控；
由 $C_3$ 给出离散回合端点，因此回合同调归约可在窗末端闭合。
故定理成立。证毕。
\end{proof}

\begin{theorem}[开放（curved-gapped）$\Rightarrow$ 只能 tick 回合]
\label{thm:open-system-tick-round}
设系统满足以下任一条件：
\begin{enumerate}
  \item $\sup_{s\in I}\|\GZLCurv(s)\|$ 使任意 $\Delta\gg\delta t$ 的窗都无法满足曲率压缩阈值；
  \item $\JacGap_{L_s}$ 在长窗内无法保持在 $\varepsilon$ 阈值内；
  \item 对手可在计算期间插入操作，导致交互点无法冻结。
\end{enumerate}
则回合窗必须收缩到最小决策周期 $\delta t$ 级别，即只能采用 tick 回合。
\end{theorem}

\begin{Certificate}[证书：开放系统的回合收缩]
\label{cert:open-system-tick-round}
证书登记谓词链：
\[
\begin{aligned}
  D_1 &: \text{曲率或 JacGap 阈值在长窗内失效},\\
  D_2 &: \text{交互点不可冻结},\\
  D_3 &: \Rightarrow\ \text{定义~\ref{def:round-window} 失效于大回合}.
\end{aligned}
\]
\end{Certificate}

\begin{proof}[证明（谓词因果链）]
由 $D_1$ 得回合窗条件的曲率/代数可压缩性失败；
由 $D_2$ 得回合端点不可稳定定义；
因此定义~\ref{def:round-window} 对 $\Delta\gg\delta t$ 不成立，
只能取 $\Delta=\delta t$，即 tick 回合。证毕。
\end{proof}

\subsection{命题与推论：外部信息不是回合语义判据}
\label{sec:round-window-proposition-corollary}

\begin{proposition}[外部信息持续加入只改变回合窗上限]
\label{prop:external-info-not-defining}
外部信息的持续加入仅通过改变 $\|\GZLCurv\|$ 与 $\JacGap$ 的上界影响可容许回合窗长度；
它不构成回合语义的定义性判据。
\end{proposition}

\begin{Certificate}[证书：信息注入 $\Rightarrow$ 曲率/Gap 变动]
\label{cert:external-info-not-defining}
证书登记谓词链：
\[
\begin{aligned}
  E_1 &: \text{外部信息改变 law-trajectory 与操作集合},\\
  E_2 &: \Rightarrow\ \|\GZLCurv\|\ \text{与}\ \JacGap\ \text{上界变化},\\
  E_3 &: \text{回合语义仍由定义~\ref{def:round-window} 决定}.
\end{aligned}
\]
\end{Certificate}

\begin{proof}[证明（谓词因果链）]
由 $E_1$ 得 law-trajectory 与操作族变化；
由 $E_2$ 得回合窗长度上限变化；
Definition~\ref{sec:round-window-formalization} 不以信息加入为判据，
因此 $E_3$ 成立。证毕。
\end{proof}

\begin{corollary}[回合语义由同调闭合条件决定]
\label{cor:round-window-cohomology}
回合窗是否成立仅取决于曲率压缩与 JacGap 可控性，以及回合端点是否可冻结；
“信息是否持续加入”仅影响这些量的大小，而不改变回合语义本身。
\end{corollary}

\begin{Certificate}[证书：判据归并]
\label{cert:round-window-cohomology}
证书登记谓词链：
\[
\begin{aligned}
  F_1 &: \text{回合窗判据由曲率/Gap/端点冻结给出},\\
  F_2 &: \text{信息持续加入}\Rightarrow \text{只改变阈值大小},\\
  F_3 &: F_1\wedge F_2\Rightarrow \text{判据不变}.
\end{aligned}
\]
\end{Certificate}

\begin{proof}[证明（谓词因果链）]
由 $F_1$ 给出判据结构，$F_2$ 给出外部信息的影响路径，
由 $F_1\wedge F_2$ 得 $F_3$，从而推论成立。证毕。
\end{proof}

\subsection{备注：可核验的回合窗口口径}
\label{sec:round-window-remarks}

\begin{remark}[回合窗的工程核验要点]
回合窗口径可通过以下三类可核验指标落地：
\begin{enumerate}
  \item 曲率证据：以 $\GZLCurv$ 的上界与可观测路径面积给出 holonomy 误差阈值；
  \item Gap 证据：以 $\JacGap$ 的采样窗口与三元组清单给出语义不变性阈值；
  \item 冻结证据：对抗交互点是否被规则冻结（或被系统资源限制为离散点）。
\end{enumerate}
该核验口径不依赖“信息持续加入”的叙述，而依赖可复算的曲率/Gap/冻结证据链。
\end{remark}

\clearpage

\section{从“魔法棒”比喻到可核验结构等价}
\label{sec:magic-wand-verified-equivalence}

\begin{remark}[动机（把比喻降级为可判定对象）]
要把“等价于魔法棒”从文学比喻提升为“可严格证明”，关键在于：先把“魔法棒”定义为形式系统中可判定、可核验的结构，
并给出与 \GFramework{} 体系（法对象/证书/同伦修复塔/变分泛函/开放博弈在线控制）可比对的最小骨架。
否则“等价”缺少共同的比较对象。
\end{remark}

\subsection{形式逻辑：可验证问题规格与证书化解}
\label{sec:magic-wand-spec-certificate}

\begin{definition}[可验证问题规格]
\label{def:magic-wand-verifiable-spec}
一个任务以规格（specification）表示为四元组
\[
  P := (\Sigma, J, \mathcal{B}, \mathrm{Loc}),
\]
其中：
\begin{itemize}
  \item $\Sigma$：规则/约束集合（“规则一致性”的可检验载体）；
  \item $J$：绩效/目标泛函（需要优化或达成的指标）；
  \item $\mathcal{B}$：预算（算力、时间、能耗、样本、可观测性等资源上界）；
  \item $\mathrm{Loc}$：运行期定位信息与外部扰动语境（开放系统的不可控输入接口）。
\end{itemize}
\end{definition}

\begin{definition}[证书化解（解--证书--验证器三元组）]
\label{def:magic-wand-certified-solution}
给定规格 $P$，称三元组
\[
  \mathbf{s} := (x,\pi,\mathsf{Ver})
\]
为一个证书化解，其中 $x\in\mathcal{X}$ 为候选解（策略、控制律、算子丛切面、程序、设计图等），
$\pi\in\Pi$ 为可独立核验的证书对象，$\mathsf{Ver}$ 为验证谓词（验证器）。
其语义约束为
\[
  \mathsf{Ver}(P,x,\pi)=\top
  \ \Longrightarrow\
  x \text{ 满足 } \Sigma \text{ 且在预算 } \mathcal{B} \text{ 内达到 } J \text{ 在 } \mathrm{Loc} \text{ 下的要求。}
\]
该定义把“相信生成过程”替换为“相信可核验的证书”。
\end{definition}

\begin{lemma}[验证通过的语义充分性]
\label{lem:magic-wand-verifier-sufficiency}
若 $\mathbf{s}=(x,\pi,\mathsf{Ver})$ 是规格 $P$ 的证书化解（定义~\ref{def:magic-wand-certified-solution}），则
\[
  \mathsf{Ver}(P,x,\pi)=\top \ \Longrightarrow\ x \text{ 是 } P \text{ 的可接受解。}
\]
\end{lemma}

\begin{Certificate}[证书：验证通过 $\Rightarrow$ 规格满足]
\label{cert:magic-wand-verifier-sufficiency}
证书登记谓词链：
\[
\begin{aligned}
  V_1 &: \mathbf{s}=(x,\pi,\mathsf{Ver})\ \text{满足定义~\ref{def:magic-wand-certified-solution} 的语义约束},\\
  V_2 &: \mathsf{Ver}(P,x,\pi)=\top,\\
  V_3 &: \Rightarrow\ x\ \text{满足}\ \Sigma\ \text{且在}\ \mathcal{B}\ \text{内达成}\ J\ \text{于}\ \mathrm{Loc}.
\end{aligned}
\]
\end{Certificate}

\begin{proof}[证明（谓词因果链）]
由 $V_1$ 直接推出“$\mathsf{Ver}(P,x,\pi)=\top$ 蕴含规格满足”的含义，
与 $V_2$ 合并得到 $V_3$，故引理成立。证毕。
\end{proof}

\subsection{形式逻辑：魔法棒算子（公理降级为定理）}
\label{sec:magic-wand-operator}

\begin{definition}[魔法棒（带证书的求解算子与在线修复族）]
\label{def:magic-wand-operator}
称 $\mathbb{W}$ 为一根“魔法棒”，若它为每个规格 $P$ 生成一对 $(x,\pi)$，
并配套一个验证器 $\mathsf{Ver}$ 与一族局部修复算子 $\{\mathbb{R}_{\delta P}\}$：
\[
  \mathbb{W}: P \longmapsto (x,\pi),\qquad
  \mathbb{R}_{\delta P}:(x,\pi)\longmapsto (x',\pi').
\]
其中 $\delta P$ 表示开放系统下的语境扰动（$P\mapsto P'=P+\delta P$）。
\end{definition}

\begin{theorem}[公理降级：正确性与多步扰动可修复性（数学归纳法证书）]
\label{thm:magic-wand-axioms-downgrade}
设存在 $\mathbb{W}$、$\mathsf{Ver}$ 与 $\{\mathbb{R}_{\delta P}\}$ 满足：
\begin{enumerate}
  \item （一步正确性）对任意规格 $P$，
  \[
    \mathsf{Ver}\!\bigl(P,\mathbb{W}(P)\bigr)=\top;
  \]
  \item （一步修复性）对任意扰动 $P\mapsto P'=P+\delta P$ 与任意已核验对 $(x,\pi)$，
  \[
    \mathsf{Ver}(P,x,\pi)=\top
    \ \Longrightarrow\
    \mathsf{Ver}\!\bigl(P',\mathbb{R}_{\delta P}(x,\pi)\bigr)=\top.
  \]
\end{enumerate}
则对任意有限扰动链
\[
  P_0 \xrightarrow{\delta P_0} P_1 \xrightarrow{\delta P_1} \cdots \xrightarrow{\delta P_{n-1}} P_n,
\]
定义递推
\[
  (x_0,\pi_0):=\mathbb{W}(P_0),\qquad
  (x_{k+1},\pi_{k+1}):=\mathbb{R}_{\delta P_k}(x_k,\pi_k),
\]
则对所有 $k\in\{0,1,\dots,n\}$ 成立
\[
  \mathsf{Ver}(P_k,x_k,\pi_k)=\top.
\]
\end{theorem}

\begin{Certificate}[数学归纳法证书：按扰动步数 $n$ 的在线可修复性]
\label{cert:magic-wand-axioms-downgrade}
令谓词 $P(n)$ 表示：“对任意长度为 $n$ 的扰动链，按定理~\ref{thm:magic-wand-axioms-downgrade} 的递推生成 $(x_k,\pi_k)$，
均满足 $\forall k\le n,\ \mathsf{Ver}(P_k,x_k,\pi_k)=\top$。”
证书化要求为：
\[
  P(0)\ \text{成立},\qquad
  (\forall n\in\NN)\ \bigl(P(n)\Rightarrow P(n+1)\bigr).
\]
\end{Certificate}

\begin{proof}[证明（数学归纳法证书；谓词因果链）]
由证书~\ref{cert:magic-wand-axioms-downgrade}：
\begin{itemize}
  \item 基步：$n=0$ 时仅有 $P_0$，且 $(x_0,\pi_0)=\mathbb{W}(P_0)$。
  由“一步正确性”得 $\mathsf{Ver}(P_0,x_0,\pi_0)=\top$，故 $P(0)$ 成立；
  \item 归纳步：假设 $P(n)$ 成立。
  对长度为 $n+1$ 的扰动链，前 $n$ 步由归纳假设已得 $\mathsf{Ver}(P_n,x_n,\pi_n)=\top$。
  再由“一步修复性”对扰动 $\delta P_n$ 得
  \[
    \mathsf{Ver}\!\bigl(P_{n+1},\mathbb{R}_{\delta P_n}(x_n,\pi_n)\bigr)=\top,
  \]
  即 $\mathsf{Ver}(P_{n+1},x_{n+1},\pi_{n+1})=\top$，从而 $P(n+1)$ 成立。
\end{itemize}
由数学归纳法，$P(n)$ 对任意 $n$ 成立，定理得证。证毕。
\end{proof}

\subsection{形式逻辑：法对象骨架（缺陷势/一致性谓词/同伦修复塔）}
\label{sec:magic-wand-law-object-skeleton}

\begin{definition}[法对象编码与缺陷势语义]
\label{def:magic-wand-law-object}
将规格 $P$ 编码为法对象（law-object）
\[
  L(P)\in\LawSpace.
\]
并配备缺陷势（defect potential）与一致性谓词
\[
  \Phi:\LawSpace\to\RR_{\ge 0},\qquad
  \mathcal{E}:\LawSpace\to\{\top,\bot\}.
\]
其语义约束为：
\[
  \mathcal{E}(L)=\top\ \wedge\ \Phi(L)=0
  \ \Longrightarrow\
  L\ \text{所编码的}\ (\Sigma,J,\mathcal{B},\mathrm{Loc})\ \text{被满足。}
\]
\end{definition}

\begin{definition}[同伦修复塔（证书保持的缺陷势下降链）]
\label{def:magic-wand-homotopy-repair-tower}
称序列
\[
  L_0 \xrightarrow{h_1} L_1 \xrightarrow{h_2}\cdots\xrightarrow{h_n} L_n
\]
为同伦修复塔，若每一步满足
\[
  \mathcal{E}(L_k)=\top,\qquad
  \Phi(L_{k+1})<\Phi(L_k).
\]
直观上：修复不是“随便改”，而是在强约束下的受控变形链，并以缺陷势单调下降作为可核验进度指标。
\end{definition}

\begin{proposition}[离散缺陷势下的有限步终止性]
\label{prop:magic-wand-discrete-termination}
若存在一类问题族使得其法对象缺陷势满足 $\Phi(\LawSpace)\subseteq\NN$，
且任一同伦修复塔都满足严格下降
\[
  \Phi(L_{k+1})\le \Phi(L_k)-1,
\]
则任意以 $L_0$ 为起点的修复塔长度至多为 $\Phi(L_0)$，并在有限步内到达某个 $L^\star$ 使得 $\Phi(L^\star)=0$。
\end{proposition}

\begin{Certificate}[证书：严格下降 $\wedge$ 离散值域 $\Rightarrow$ 有界步数]
\label{cert:magic-wand-discrete-termination}
证书登记谓词链：
\[
\begin{aligned}
  T_1 &: \Phi(\LawSpace)\subseteq\NN,\\
  T_2 &: \forall k,\ \Phi(L_{k+1})\le \Phi(L_k)-1,\\
  T_3 &: \Rightarrow\ \forall k\le \Phi(L_0),\ \Phi(L_k)\in\NN\ \text{且严格下降},\\
  T_4 &: \Rightarrow\ \exists k^\star\le \Phi(L_0),\ \Phi(L_{k^\star})=0.
\end{aligned}
\]
\end{Certificate}

\begin{proof}[证明（谓词因果链）]
由 $T_1$ 得 $\Phi(L_0)\in\NN$。
由 $T_2$ 得每步缺陷势至少下降 $1$，因此对任意 $k\in\NN$ 有
\[
  \Phi(L_k)\le \Phi(L_0)-k.
\]
取 $k=\Phi(L_0)$ 得 $\Phi(L_{\Phi(L_0)})\le 0$。又因 $\Phi(L_{\Phi(L_0)})\in\RR_{\ge 0}$，
故 $\Phi(L_{\Phi(L_0)})=0$，从而 $T_4$ 成立。证毕。
\end{proof}

\subsection{定理：由法对象体系构造一根魔法棒}
\label{sec:magic-wand-theorem-forward}

\begin{theorem}[从法对象/修复塔构造魔法棒]
\label{thm:magic-wand-from-law-object}
假设对每个规格 $P$，存在 $L_0:=L(P)$ 以及在预算 $\mathcal{B}$ 内可执行的同伦修复塔，
得到某个 $L^\star$ 满足
\[
  \mathcal{E}(L^\star)=\top,\qquad
  \Phi(L^\star)=0.
\]
再假设存在提取算子与证书算子
\[
  \mathrm{Sol}:\LawSpace\to\mathcal{X},\qquad
  \mathrm{Cert}:\LawSpace\to\Pi,
\]
使得 $\mathrm{Cert}(L)$ 可核验 $\mathcal{E}(L)=\top$ 以及 $\Phi(L)=0$（并与 $x=\mathrm{Sol}(L)$ 匹配）。
则可构造一根魔法棒 $\mathbb{W}_{\mathrm{GZ}}$、验证器 $\mathsf{Ver}$ 与在线修复族 $\{\mathbb{R}_{\delta P}\}$，
满足对任意 $P$ 均有
\[
  \mathsf{Ver}\!\bigl(P,\mathbb{W}_{\mathrm{GZ}}(P)\bigr)=\top,
\]
并且当 $P\mapsto P'=P+\delta P$ 时，可由修复塔（原塔续接或重算）给出 $\mathbb{R}_{\delta P}$ 使得
\[
  \mathsf{Ver}\!\bigl(P',\mathbb{R}_{\delta P}(\mathbb{W}_{\mathrm{GZ}}(P))\bigr)=\top.
\]
\end{theorem}

\begin{Certificate}[证书：由缺陷势归零与一致性证书推出可核验解]
\label{cert:magic-wand-from-law-object}
证书登记谓词链：
\[
\begin{aligned}
  G_1 &: \exists L^\star,\ \mathcal{E}(L^\star)=\top\ \wedge\ \Phi(L^\star)=0,\\
  G_2 &: x:=\mathrm{Sol}(L^\star),\ \pi:=\mathrm{Cert}(L^\star),\\
  G_3 &: \mathsf{Ver}(P,x,\pi)=\top\ \Longleftrightarrow\ \pi\ \text{可核验}\ \mathcal{E}(L^\star)=\top\ \wedge\
    \Phi(L^\star)=0\ \text{且匹配}\ x,\\
  G_4 &: \mathcal{E}(L^\star)=\top\ \wedge\ \Phi(L^\star)=0\Rightarrow P\ \text{被满足（定义~\ref{def:magic-wand-law-object}）},
    \\
  G_5 &: P\mapsto P'\Rightarrow \text{修复塔给出}\ L(P')^\star\ \text{并诱导}\ \mathbb{R}_{\delta P}.
\end{aligned}
\]
\end{Certificate}

\begin{proof}[证明（谓词因果链）]
给定 $P$，由 $G_1$ 取到满足 $\mathcal{E}(L^\star)=\top,\ \Phi(L^\star)=0$ 的 $L^\star$。
按 $G_2$ 定义 $x,\pi$，并按 $G_3$ 定义验证器 $\mathsf{Ver}$。
由 $G_1$ 与 $G_3$ 立得 $\mathsf{Ver}(P,x,\pi)=\top$；再由 $G_4$ 得 $x$ 满足 $P$ 的规格语义。
令 $\mathbb{W}_{\mathrm{GZ}}(P):=(x,\pi)$，即得“带证书正确性”。

对扰动 $P\mapsto P'$，由 $G_5$ 可（续接或重算）得到 $L(P')^\star$，从而得到 $(x',\pi')$ 并定义
$\mathbb{R}_{\delta P}(x,\pi):=(x',\pi')$。
同理有 $\mathsf{Ver}(P',x',\pi')=\top$，即在线可修复性成立。证毕。
\end{proof}

\subsection{定理：任意魔法棒都可嵌入为法对象机制}
\label{sec:magic-wand-theorem-backward}

\begin{theorem}[反向嵌入：任何魔法棒都可表达为法对象系统]
\label{thm:magic-wand-into-law-object}
给定任意魔法棒结构 $(\mathbb{W},\mathsf{Ver},\{\mathbb{R}_{\delta P}\})$（定义~\ref{def:magic-wand-operator}），
可构造一个法对象系统 $(\LawSpace,\Phi,\mathcal{E},\text{修复塔})$，使得由该系统按定理~\ref{thm:magic-wand-from-law-object} 构造出的
$\mathbb{W}_{\mathrm{GZ}}$ 与原 $\mathbb{W}$ 在“带证书可核验”的意义下等价：
对每个规格 $P$，二者输出的 $(x,\pi)$ 在同一验证器 $\mathsf{Ver}$ 下均可核验通过（且候选解语义一致）。
\end{theorem}

\begin{Certificate}[证书：把 $\mathbb{W}$ 编码进法对象与缺陷势]
\label{cert:magic-wand-into-law-object}
证书登记谓词链：
\[
\begin{aligned}
  H_1 &: \text{定义法对象}\ L(P):=(P,x,\pi)\ \text{其中}\ (x,\pi):=\mathbb{W}(P),\\
  H_2 &: \mathcal{E}(L(P)):=\mathsf{Ver}(P,x,\pi),\\
  H_3 &: \Phi(L(P)):=\begin{cases}0,&\mathcal{E}(L(P))=\top,\\ 1,&\text{否则},\end{cases}\\
  H_4 &: \text{由定理~\ref{thm:magic-wand-axioms-downgrade} 得}\ \mathcal{E}(L(P))=\top\ \text{从而}\ \Phi(L(P))=0,\\
  H_5 &: \text{修复塔取为“调用 }\mathbb{R}_{\delta P}\text{（或重新调用 }\mathbb{W}\text{）”的迭代}.
\end{aligned}
\]
\end{Certificate}

\begin{proof}[证明（谓词因果链）]
按 $H_1$--$H_3$ 完成法对象与谓词/缺陷势的构造。
由 $H_4$ 得对任意 $P$ 都有 $\mathcal{E}(L(P))=\top$ 且 $\Phi(L(P))=0$，因此“缺陷势归零”语义成立。
由 $H_5$ 可在扰动 $P\mapsto P'$ 下沿同一修复算子族生成修复塔，重现原魔法棒的在线修复行为。
因此由该法对象系统反推出的 $\mathbb{W}_{\mathrm{GZ}}$ 与 $\mathbb{W}$ 在可核验意义下等价。证毕。
\end{proof}

\subsection{推论与备注：结构等价、观感来源与边界}
\label{sec:magic-wand-corollary-remarks}

\begin{corollary}[“等价于魔法棒”的严格含义（互相可构造/可表示）]
\label{cor:magic-wand-equivalence}
在“对可证书化/可验证的问题族给出 $\text{解}+\text{证书}$，并在开放扰动下可证书化修复”的语义层面上：
\begin{enumerate}
  \item 由定理~\ref{thm:magic-wand-from-law-object}，法对象/证书/修复塔机制诱导一根魔法棒 $\mathbb{W}_{\mathrm{GZ}}$；
  \item 由定理~\ref{thm:magic-wand-into-law-object}，任意魔法棒都可嵌入为法对象系统并被该机制表达。
\end{enumerate}
因此二者在该语义层面构成严格的结构等价（可理解为在合适的“可核验解”范畴里自然同构）。
\end{corollary}

\begin{Certificate}[证书：互相构造 $\Rightarrow$ 结构等价]
\label{cert:magic-wand-equivalence}
证书登记谓词链：
\[
\begin{aligned}
  I_1 &: \text{定理~\ref{thm:magic-wand-from-law-object}}\Rightarrow \text{法对象机制}\ \mapsto\ \mathbb{W}_{\mathrm{GZ}},\\
  I_2 &: \text{定理~\ref{thm:magic-wand-into-law-object}}\Rightarrow \mathbb{W}\ \mapsto\ \text{法对象机制},\\
  I_3 &: I_1\wedge I_2\Rightarrow \text{互相可构造/可表示}.
\end{aligned}
\]
\end{Certificate}

\begin{proof}[证明（谓词因果链）]
由 $I_1$、$I_2$ 得二者互相可构造，按“可核验解语义”的等价定义得到 $I_3$，故推论成立。证毕。
\end{proof}

\begin{remark}[边界说明（保持严格）]
若把“魔法棒”理解为“对所有可描述问题都必定给出解”，将触及一般意义的不可判定/不可解边界。
本章采用的口径是：仅对\textbf{可证书化/可验证}的问题族，给出\textbf{解 + 证书 + 可修复}。
这也对应工程与元数学意义下更可落地、且可审计的主战场。
\end{remark}

\clearpage

\section{PureMath + Vol1 记号对齐：魔法棒=带证书的规范形选取器/终对象机制}
\label{sec:magic-wand-normal-form-terminal}

\begin{remark}[目标（把“魔法棒”提升为可核验的规范形机制）]
本章在 \GFramework{} 体系的记号骨架下，把“魔法棒”严格化为一个\textbf{带证书的规范形选取器}（normal-form selector），
并给出其“终对象（terminal object）/普适性”口径的形式化版本。
核心结论是：在 PureMath + Vol1 已给出的法对象入口、$\JacGap$ 缺陷证书与同伦补齐（OHU）原则下，
可以构造一个 $\Phi_\lambda$-最小的规范形代表，并附带可核验的证书链条。
\end{remark}

\subsection{形式逻辑：法对象入口函子与 $\JacGap$ 缺陷证书}
\label{sec:magic-wand-puremath-law-entry}

\begin{definition}[法对象空间与对象层表示映射]
\label{def:magic-wand-puremath-lawspace}
令
\[
  \GZLS=\mathfrak{L}_{\mathrm{GZ}},\qquad L\in\GZLS
\]
为法对象空间与法对象。
记三类对象层语义载体为
\[
  \mathsf{PFB\text{-}GNLA}_{\mathrm{strict}},\qquad
  \mathsf{PFB\text{-}GNLA}_{\mathrm{homo}},\qquad
  \mathsf{PFB\text{-}GNLA}_{\mathrm{vf\text{-}op}}.
\]
假设存在法空间表示映射（可视作函子在对象层面的抽象）
\[
  L(-):
  \mathsf{PFB\text{-}GNLA}_{\mathrm{vf\text{-}op}}
  \cup
  \mathsf{PFB\text{-}GNLA}_{\mathrm{homo}}
  \cup
  \mathsf{PFB\text{-}GNLA}_{\mathrm{strict}}
  \longrightarrow \GZLS.
\]
其含义是：无论输入处于 strict/homo/vf-op 哪一层，均可在统一入口 $L(-)$ 下落到同一法对象空间中执行“缺陷测量--同伦补齐--证书核验”。
\end{definition}

\begin{lemma}[统一入口：一切缺陷测量可在 $\GZLS$ 上表述]
\label{lem:magic-wand-puremath-entry}
对任意对象层输入 $\mathcal{Q}$（处于上述三类之一），其可证书化的缺陷与修复判据
均可等价地表述为 $\GZLS$ 上关于 $L(\mathcal{Q})$ 的判定谓词链。
\end{lemma}

\begin{Certificate}[证书：对象层 $\Rightarrow$ 法对象层（统一入口）]
\label{cert:magic-wand-puremath-entry}
证书登记谓词链：
\[
\begin{aligned}
  U_1 &: \mathcal{Q}\in \mathsf{PFB\text{-}GNLA}_{\mathrm{vf\text{-}op}}\cup \mathsf{PFB\text{-}GNLA}_{\mathrm{homo}}
    \cup \mathsf{PFB\text{-}GNLA}_{\mathrm{strict}},\\
  U_2 &: \text{由定义~\ref{def:magic-wand-puremath-lawspace}，存在 }L(\mathcal{Q})\in\GZLS,\\
  U_3 &: \text{缺陷/证书/修复判据以 }L(\mathcal{Q})\text{ 为参数可被统一编码}.
\end{aligned}
\]
\end{Certificate}

\begin{proof}[证明（谓词因果链）]
由 $U_1$ 与 $U_2$ 得 $L(\mathcal{Q})\in\GZLS$ 可用作统一入口；
将对象层判据一律通过 $L(-)$ 推送到 $\GZLS$ 上，即得 $U_3$。证毕。
\end{proof}

\begin{definition}[$\JacGap$ 的提升（缺陷函数/证书口径）]
\label{def:magic-wand-puremath-jacgap}
设存在缺陷泛函（Jacobiator-gap functional）
\[
  \JacGap:\GZLS\to\RR_{\ge 0}.
\]
对任意对象层输入 $\mathcal{Q}$，定义其缺陷为
\[
  \JacGap(\mathcal{Q}) := \JacGap\bigl(L(\mathcal{Q})\bigr).
\]
该口径把 $\JacGap$ 同时用作：
\begin{enumerate}
  \item 缺陷度量（defect functional）：用于指示“缺失的同伦数据/不严格的 Jacobi 结构”；
  \item 证书接口（certificate interface）：通过 $\JacGap$ 的数值/上界与构造日志，提供可核验的偏差证据链。
\end{enumerate}
\end{definition}

\begin{definition}[Upgrade/Strictify 与 $\JacGap$-free 子范畴]
\label{def:magic-wand-puremath-upgrade-strictify}
设存在
\[
  \mathsf{Upgrade}:\mathsf{PFB\text{-}GNLA}_{\mathrm{strict}}\to \mathsf{PFB\text{-}GNLA}_{\mathrm{homo}},
\]
以及一个仅在 $\JacGap$-free 子范畴上定义的部分算子
\[
  \begin{aligned}
    \mathsf{Strictify}&:\mathsf{PFB\text{-}GNLA}^{(0)}_{\mathrm{homo}}\dashrightarrow \mathsf{PFB\text{-}GNLA}
      _{\mathrm{strict}},\\
    \mathsf{PFB\text{-}GNLA}^{(0)}_{\mathrm{homo}}
    &:=
    \bigl\{\mathcal{Q}\in \mathsf{PFB\text{-}GNLA}_{\mathrm{homo}}:\JacGap(\mathcal{Q})=0\bigr\}.
  \end{aligned}
\]
\end{definition}

\begin{theorem}[$\JacGap=0$ 的严格化刻画（PureMath 证书）]
\label{thm:magic-wand-puremath-jacgap-vanishing}
对任意 $L\in\GZLS$，若 $\JacGap(L)=0$，则存在严格代表 $L_{\mathrm{strict}}$，
使其严格满足 Jacobi（strict Jacobi），并且在严格化过程中不破坏受保护的低阶可观测量（protected observables）。
等价地，对任意 $\mathcal{Q}\in \mathsf{PFB\text{-}GNLA}^{(0)}_{\mathrm{homo}}$，$\mathsf{Strictify}(\mathcal{Q})$ 有定义且给出 strict
  层代表。
\end{theorem}

\begin{Certificate}[证书：JacGap 消失 $\Rightarrow$ strictify 可用且不破坏 protected observables]
\label{cert:magic-wand-puremath-jacgap-vanishing}
证书登记谓词链：
\[
\begin{aligned}
  J_1 &: \JacGap(L)=0,\\
  J_2 &: \text{PureMath：}\JacGap(L)=0\ \Longrightarrow\ \exists L_{\mathrm{strict}}\ \text{严格满足 Jacobi},\\
  J_3 &: \text{PureMath：严格化过程保持 protected observables（不改变受保护观测）},\\
  J_4 &: \mathcal{Q}\in \mathsf{PFB\text{-}GNLA}^{(0)}_{\mathrm{homo}}\Rightarrow \JacGap(\mathcal{Q})=0\Rightarrow
    \mathsf{Strictify}(\mathcal{Q})\ \text{有定义}.
\end{aligned}
\]
\end{Certificate}

\begin{proof}[证明（谓词因果链）]
由 $J_1$ 与 $J_2$ 得存在严格代表 $L_{\mathrm{strict}}$；
由 $J_3$ 得严格化不破坏受保护观测；
若以对象层表述，则由 $J_4$ 得 $\JacGap$-free 时 $\mathsf{Strictify}$ 可用并给出 strict 代表。证毕。
\end{proof}

\subsection{形式逻辑：规格接口、受保护一致性与合并缺陷势}
\label{sec:magic-wand-puremath-spec-potential}

\begin{definition}[规格接口（\texorpdfstring{$(\mathrm{Loc},\Sigma)$}{(Loc,Sigma)}）]
\label{def:magic-wand-puremath-spec}
对给定任务/系统，记运行期语境与硬约束为
\[
  (\mathrm{Loc},\Sigma),
\]
其中 $\mathrm{Loc}$ 表示运行期定位/语境/可观测接口，$\Sigma$ 表示硬约束/规则集合。
\end{definition}

\begin{definition}[受保护一致性谓词（protected 证书接口）]
\label{def:magic-wand-puremath-protected}
把“受保护（protected）可观测量/不变量集合不被破坏”编码为可核验谓词（或能量）：
\[
  \mathcal{E}_{\mathrm{prot}}(L;\mathrm{Loc},\Sigma)=0.
\]
其语义是：在规格接口 $(\mathrm{Loc},\Sigma)$ 下，$L$ 的 protected observables 仍保持一致性（可审计、可复算）。
\end{definition}

\begin{definition}[合并缺陷势（defect potential）]
\label{def:magic-wand-puremath-potential}
令 $\mathcal{E}_{\mathrm{CH}}(L;\mathrm{Loc},\Sigma)\in\RR_{\ge 0}$ 表示 PureMath 口径下的“结构/连接/补齐代价”项（例如 CH-energy），
定义合并缺陷势为
\[
  \Phi_\lambda(L;\mathrm{Loc},\Sigma)
  :=
  \mathcal{E}_{\mathrm{CH}}(L;\mathrm{Loc},\Sigma)
  + \lambda\,\JacGap(L),
  \qquad \lambda>0.
\]
直观上：$\mathcal{E}_{\mathrm{CH}}$ 负责“完成/补齐的几何代价”，$\JacGap$ 负责“代数缺陷证据”，$\lambda$ 用于权衡两者。
\end{definition}

\begin{definition}[可接受类（admissible class）]
\label{def:magic-wand-puremath-admissible}
定义在 $(\mathrm{Loc},\Sigma)$ 下的可接受法对象集合为
\[
  \mathsf{Adm}(\mathrm{Loc},\Sigma)
  :=
  \{\,L\in\GZLS:\mathcal{E}_{\mathrm{prot}}(L;\mathrm{Loc},\Sigma)=0\,\}.
\]
并记 $L\simeq_{\mathrm{OH}}\widetilde L$ 表示二者在 operator-homotopy 意义下等价，且等价变形保持 protected observables。
\end{definition}

\subsection{形式逻辑：规范形选取器（魔法棒）与公理降级}
\label{sec:magic-wand-puremath-wand}

\begin{definition}[规范形选取器（\texorpdfstring{$\mathbb{W}_{\mathrm{GZ}}$}{W\_GZ}）]
\label{def:magic-wand-puremath-wand}
给定固定的 $(\mathrm{Loc},\Sigma)$，称映射
\[
  \mathbb{W}_{\mathrm{GZ}}:\mathsf{Adm}(\mathrm{Loc},\Sigma)\to \mathsf{Adm}(\mathrm{Loc},\Sigma),
  \qquad
  L\longmapsto L^\star
\]
为规范形选取器（亦即“魔法棒”的严格版本），若它输出的 $L^\star$ 满足：
\begin{enumerate}
  \item 证书保持：$\mathcal{E}_{\mathrm{prot}}(L;\mathrm{Loc},\Sigma)=0\Rightarrow \mathcal{E}_{\mathrm{prot}}(L^\star;
    \mathrm{Loc},\Sigma)=0$；
  \item 缺陷势最小：不存在 $\widetilde L\in\mathsf{Adm}(\mathrm{Loc},\Sigma)$ 使 $\Phi_\lambda(\widetilde L;\mathrm{Loc},\Sigma)
    <\Phi_\lambda(L^\star;\mathrm{Loc},\Sigma)$；
  \item 同伦普适性：对任意 $\widetilde L\in\mathsf{Adm}(\mathrm{Loc},\Sigma)$，存在保持 protected observables 的 operator-homotopy 证书链把
    $\widetilde L$ 因子化到 $L^\star$（终对象口径见备注~\ref{rem:magic-wand-puremath-terminal}）。
\end{enumerate}
\end{definition}

\begin{theorem}[公理降级：OHU 迭代的证书保持与缺陷势单调性（数学归纳法证书）]
\label{thm:magic-wand-puremath-monotone}
给定初值 $L^{(0)}\in\mathsf{Adm}(\mathrm{Loc},\Sigma)$。
设存在一条同伦补齐/修复迭代（GZ-OHU 口径）生成序列 $\{L^{(n)}\}_{n\in\NN}$，满足对任意 $n\in\NN$：
\begin{enumerate}
  \item （一步证书保持）$\mathcal{E}_{\mathrm{prot}}(L^{(n)};\mathrm{Loc},\Sigma)=0\Rightarrow \mathcal{E}_{\mathrm{prot}}
    (L^{(n+1)};\mathrm{Loc},\Sigma)=0$；
  \item （一步单调性）$\Phi_\lambda(L^{(n+1)};\mathrm{Loc},\Sigma)\le \Phi_\lambda(L^{(n)};\mathrm{Loc},\Sigma)$。
\end{enumerate}
则对所有 $n\in\NN$ 都成立：
\[
  L^{(n)}\in\mathsf{Adm}(\mathrm{Loc},\Sigma),
  \qquad
  \Phi_\lambda(L^{(n)};\mathrm{Loc},\Sigma)\le \Phi_\lambda(L^{(0)};\mathrm{Loc},\Sigma).
\]
\end{theorem}

\begin{Certificate}[数学归纳法证书：按迭代步数 $n$ 的保持性与单调性]
\label{cert:magic-wand-puremath-monotone}
令谓词 $P(n)$ 为：
“$L^{(n)}\in\mathsf{Adm}(\mathrm{Loc},\Sigma)$ 且 $\Phi_\lambda(L^{(n)};\mathrm{Loc},\Sigma)\le \Phi_\lambda(L^{(0)};
  \mathrm{Loc},\Sigma)$。”
证书化要求为：
\[
  P(0)\ \text{成立},\qquad
  (\forall n\in\NN)\ \bigl(P(n)\Rightarrow P(n+1)\bigr).
\]
\end{Certificate}

\begin{proof}[证明（数学归纳法证书；谓词因果链）]
由证书~\ref{cert:magic-wand-puremath-monotone}：
\begin{itemize}
  \item 基步：$n=0$ 时由 $L^{(0)}\in\mathsf{Adm}(\mathrm{Loc},\Sigma)$ 得 $\mathcal{E}_{\mathrm{prot}}(L^{(0)};\mathrm{Loc},
    \Sigma)=0$，
  且 $\Phi_\lambda(L^{(0)};\mathrm{Loc},\Sigma)\le \Phi_\lambda(L^{(0)};\mathrm{Loc},\Sigma)$，故 $P(0)$ 成立；
  \item 归纳步：假设 $P(n)$ 成立，则
  \[
    \mathcal{E}_{\mathrm{prot}}(L^{(n)};\mathrm{Loc},\Sigma)=0.
  \]
  由“一步证书保持”得
  \[
    \mathcal{E}_{\mathrm{prot}}(L^{(n+1)};\mathrm{Loc},\Sigma)=0,
  \]
  即 $L^{(n+1)}\in\mathsf{Adm}(\mathrm{Loc},\Sigma)$。
  再由“一步单调性”得
  \[
    \Phi_\lambda(L^{(n+1)};\mathrm{Loc},\Sigma)\le \Phi_\lambda(L^{(n)};\mathrm{Loc},\Sigma)\le \Phi_\lambda(L^{(0)};
      \mathrm{Loc},\Sigma),
  \]
  从而 $P(n+1)$ 成立。
\end{itemize}
由数学归纳法，对任意 $n$，$P(n)$ 成立，定理得证。证毕。
\end{proof}

\begin{theorem}[GZ-OHU + 极限最小性 $\Rightarrow$ 存在规范形选取器 $\mathbb{W}_{\mathrm{GZ}}$]
\label{thm:magic-wand-puremath-existence}
在定理~\ref{thm:magic-wand-puremath-monotone} 的条件下，进一步假设：
\begin{enumerate}
  \item （极限存在）序列 $\{L^{(n)}\}$ 在 $\GZLS$ 中存在一个累积点/极限点 $L^\star$；
  \item （不可再降）不存在 $\widetilde L\in\mathsf{Adm}(\mathrm{Loc},\Sigma)$ 使
  $\Phi_\lambda(\widetilde L;\mathrm{Loc},\Sigma)<\Phi_\lambda(L^\star;\mathrm{Loc},\Sigma)$。
\end{enumerate}
则映射 $\mathbb{W}_{\mathrm{GZ}}(L^{(0)}):=L^\star$ 满足定义~\ref{def:magic-wand-puremath-wand} 的前两条公理（证书保持与缺陷势最小）。
并且若再引入终对象口径（备注~\ref{rem:magic-wand-puremath-terminal}），则满足同伦普适性公理。
\end{theorem}

\begin{Certificate}[证书：单调下降 + 极限不可再降 $\Rightarrow$ 规范形最小代表]
\label{cert:magic-wand-puremath-existence}
证书登记谓词链：
\[
\begin{aligned}
  M_1 &: L^{(0)}\in\mathsf{Adm}(\mathrm{Loc},\Sigma)\ \text{且}\ \{L^{(n)}\}\ \text{满足定理~\ref{thm:magic-wand-puremath-monotone} 的两条一步性质},\\
  M_2 &: \exists L^\star,\ L^\star\ \text{为}\{L^{(n)}\}\ \text{的累积点/极限点},\\
  M_3 &: \neg\exists \widetilde L\in\mathsf{Adm}(\mathrm{Loc},\Sigma)\ \text{使}\ \Phi_\lambda(\widetilde L)
    <\Phi_\lambda(L^\star),\\
  M_4 &: \Rightarrow\ L^\star\in\mathsf{Adm}(\mathrm{Loc},\Sigma)\ \text{且}\ L^\star\ \text{为}\Phi_\lambda\
    \text{的不可再降代表},\\
  M_5 &: \Rightarrow\ \mathbb{W}_{\mathrm{GZ}}(L^{(0)}):=L^\star\ \text{满足定义~\ref{def:magic-wand-puremath-wand} 的(1)(2)}
    .
\end{aligned}
\]
\end{Certificate}

\begin{proof}[证明（谓词因果链）]
由 $M_1$ 与定理~\ref{thm:magic-wand-puremath-monotone} 得对任意 $n$，$L^{(n)}\in\mathsf{Adm}(\mathrm{Loc},\Sigma)$ 且
  $\Phi_\lambda(L^{(n)})$ 单调不增并下界为 $0$。
由 $M_2$ 取到极限点 $L^\star$；结合“证书保持”的闭性（在本章口径下作为可核验前提），得到 $L^\star\in\mathsf{Adm}(\mathrm{Loc},\Sigma)$，即 $M_4$ 的可接受性部分。
再由 $M_3$ 得 $L^\star$ 不可被任何可接受对象在 $\Phi_\lambda$ 上严格压低，因此满足“缺陷势最小/不可再降”的判据。
令 $\mathbb{W}_{\mathrm{GZ}}(L^{(0)}):=L^\star$，则定义~\ref{def:magic-wand-puremath-wand} 的(1)(2)成立，即 $M_5$。证毕。
\end{proof}

\begin{corollary}[$\JacGap$-free 分支的 strict 规范形（Strictify 证书化）]
\label{cor:magic-wand-puremath-strict-branch}
在定理~\ref{thm:magic-wand-puremath-existence} 的条件下，若额外满足 $\JacGap(L^\star)=0$，则存在 strict 规范形代表 $L^\star_{\mathrm{strict}}$，

并且 strict 化过程保持 protected observables。
\end{corollary}

\begin{Certificate}[证书：$\JacGap(L^\star)=0\Rightarrow$ strictify 可用]
\label{cert:magic-wand-puremath-strict-branch}
证书登记谓词链：
\[
\begin{aligned}
  S_1 &: L^\star\in\mathsf{Adm}(\mathrm{Loc},\Sigma),\\
  S_2 &: \JacGap(L^\star)=0,\\
  S_3 &: \text{由定理~\ref{thm:magic-wand-puremath-jacgap-vanishing}}\Rightarrow \exists L^\star_{\mathrm{strict}}\
    \text{且保持 protected observables}.
\end{aligned}
\]
\end{Certificate}

\begin{proof}[证明（谓词因果链）]
由 $S_2$ 与定理~\ref{thm:magic-wand-puremath-jacgap-vanishing} 直接得到 strict 代表存在且保持 protected observables；
$S_1$ 保证其仍处在同一规格接口的可接受类内。证毕。
\end{proof}

\subsection{推论与备注：与 Vol1 证书化 AI 接口对齐，以及终对象版本}
\label{sec:magic-wand-puremath-vol1-terminal}

\begin{corollary}[Vol1：证书化 AI 的规范形输出口径（policy certificates + JacGap）]
\label{cor:magic-wand-puremath-vol1}
将 Vol1 的证书化 AI 规格接口抽象为 $(\mathrm{Loc},\Sigma)$，
并把“环境--策略组合”编码为 combined law-object（记为 $\mathcal{L}_{\mathrm{RL}}$）。
若存在可核验的 $\mathsf{Cert}_{\mathrm{pol}}$（策略证书）与 $\JacGap$（缺陷证书），
则 $\mathbb{W}_{\mathrm{GZ}}$ 可输出一个在受保护约束下 $\Phi_\lambda$-最小的规范形代表，
并附带 $(\mathsf{Cert}_{\mathrm{pol}},\JacGap)$ 的证书链条作为审计接口。
\end{corollary}

\begin{Certificate}[证书：Vol1 接口 $\Rightarrow$ 规范形 + 证书输出]
\label{cert:magic-wand-puremath-vol1}
证书登记谓词链：
\[
\begin{aligned}
  A_1 &: (\mathrm{Loc},\Sigma)\ \text{给出运行期接口与安全/合规硬约束},\\
  A_2 &: \mathcal{L}_{\mathrm{RL}}\ \text{为 env/pol 的 combined law-object 表示},\\
  A_3 &: \mathsf{Cert}_{\mathrm{pol}}\ \text{可核验策略侧约束与一致性},\\
  A_4 &: \JacGap\ \text{提供可核验缺陷证据},\\
  A_5 &: \text{由定理~\ref{thm:magic-wand-puremath-existence}}\Rightarrow \mathbb{W}_{\mathrm{GZ}}\ \text{输出}
    \Phi_\lambda\text{-最小规范形并保持受保护约束}.
\end{aligned}
\]
\end{Certificate}

\begin{proof}[证明（谓词因果链）]
由 $A_1$--$A_4$ 给出“规格接口 + 证书接口”的可核验输入口径；
由 $A_5$ 得 $\mathbb{W}_{\mathrm{GZ}}$ 产生规范形代表并保持受保护约束，
且其最小性由 $\Phi_\lambda$ 的定义（定义~\ref{def:magic-wand-puremath-potential}）与不可再降判据保证。
因此可把 $(\mathsf{Cert}_{\mathrm{pol}},\JacGap)$ 作为输出证书链条的一部分给出审计接口。证毕。
\end{proof}

\begin{remark}[终对象口径（更强版本的形式化）]
\label{rem:magic-wand-puremath-terminal}
固定 $(\mathrm{Loc},\Sigma)$，可将定义~\ref{def:magic-wand-puremath-admissible} 的可接受对象组织成一个范畴（或 2-范畴）：
对象为 $\mathsf{Adm}(\mathrm{Loc},\Sigma)$ 中的法对象，
态射为保持 protected observables 的 operator-homotopy（按同伦等价类取商可得到 1-范畴）。
在该口径下，“同伦普适性”可等价表述为：
$L^\star$ 在该范畴中为终对象（terminal object）或“终对象 up to homotopy”，并且同时最小化 $\Phi_\lambda$。
于是 $\mathbb{W}_{\mathrm{GZ}}$ 可视作诱导一个规范化函子（normalization functor）。
\end{remark}

\begin{proposition}[终对象 $\Rightarrow$ 规范化函子（定义层面的可核验证书）]
\label{prop:magic-wand-puremath-terminal-functor}
在备注~\ref{rem:magic-wand-puremath-terminal} 的范畴口径下，
若存在对象 $L^\star\in\mathsf{Adm}(\mathrm{Loc},\Sigma)$ 满足：
\begin{enumerate}
  \item $L^\star$ 为终对象（或终对象 up to homotopy）；
  \item $L^\star$ 在 $\mathsf{Adm}(\mathrm{Loc},\Sigma)$ 内 $\Phi_\lambda$-不可再降；
\end{enumerate}
则映射 $L\mapsto L^\star$ 诱导一个规范化函子（把任意可接受对象送到同一规范形代表），
并满足定义~\ref{def:magic-wand-puremath-wand} 的三条公理。
\end{proposition}

\begin{Certificate}[证书：终对象性质 + 最小性 $\Rightarrow$ 普适规范形]
\label{cert:magic-wand-puremath-terminal-functor}
证书登记谓词链：
\[
\begin{aligned}
  F_1 &: L^\star\ \text{为（up to homotopy 的）终对象}\Rightarrow \forall L\in\mathsf{Adm},\ \exists!\ (L\to L^\star)\
    \text{（同伦意义下）},\\
  F_2 &: L^\star\ \Phi_\lambda\text{-不可再降}\Rightarrow \text{满足缺陷势最小公理},\\
  F_3 &: \forall L\in\mathsf{Adm},\ \exists (L\to L^\star)\Rightarrow \text{满足同伦普适性公理},\\
  F_4 &: \text{对象均在 }\mathsf{Adm}\Rightarrow \text{证书保持公理}.
\end{aligned}
\]
\end{Certificate}

\begin{proof}[证明（谓词因果链）]
由 $F_1$ 得到任意可接受对象到 $L^\star$ 的（同伦意义下）唯一因子化，从而给出同伦普适性（$F_3$）。
由 $F_2$ 得到缺陷势最小性；由 $F_4$ 得证书保持。
合并三者即满足定义~\ref{def:magic-wand-puremath-wand} 的三条公理，命题成立。证毕。
\end{proof}

\clearpage

\section{魔法棒同伦类的双重同属性：文明跃迁与文明失控（同源/同构/同伦不变）}
\label{sec:magic-wand-homotopy-dual-properties}

\begin{remark}[定位（同一结构的正负两面）]
本章把“魔法棒同伦类”的两种社会层面观感写成同一套形式系统中的两个可证明性质：
\textbf{颠覆性（文明跃迁）}与\textbf{毁灭性（文明失控）}。
我们不引入外部材料，只依赖已固定的骨架：法对象 $L$、受保护约束 $\mathcal{E}_{\mathrm{prot}}$、缺陷势 $\Phi$、同伦修复塔（在线修复）。
\end{remark}

\subsection{形式逻辑：文明开放系统与法对象编码}
\label{sec:magic-wand-civilization-law-encoding}

\begin{definition}[文明作为开放系统（最小抽象）]
\label{def:civilization-open-system}
用开放系统描述“文明的可控演化”：
\begin{itemize}
  \item 状态空间 $X$（能源、产业、制度、科技、基础设施等的合成态）；
  \item 行为空间 $A$（可执行的行动/政策/工程/算法/组织动作集合）；
  \item 转移 $T:X\times A\to X$；
  \item 约束 $\Sigma\subseteq A$（法律/伦理/安全/资源/物理可行性等硬约束合成）；
  \item 外部语境 $\mathrm{Loc}$（环境/对手/自然扰动/信息流入等开放性接口）；
  \item 预算 $\mathcal{B}$（计算、时间、资源与可观测性等上界）。
\end{itemize}
\end{definition}

\begin{definition}[状态到法对象的编码]
\label{def:civilization-law-encoding}
用法对象编码状态：
\[
  L:X\to \GZLS,\qquad x\mapsto L(x).
\]
并约定受保护约束以谓词接口表达：
\[
  \mathcal{E}_{\mathrm{prot}}(L(x);\mathrm{Loc},\Sigma)=0.
\]
\end{definition}

\subsection{形式逻辑：魔法棒代表元、证书验证器与同伦等价}
\label{sec:magic-wand-representative-homotopy}

\begin{definition}[魔法棒代表元（带证书的规范化行动算子）]
\label{def:magic-wand-representative}
称 $W$ 为一个“魔法棒代表元”，若它对任意输入
\[
  \bigl(L(x_0),\Sigma,\mathrm{Loc},\mathcal{B}\bigr)
\]
输出一个行动序列与证书对
\[
  W:\ (L(x_0),\Sigma,\mathrm{Loc},\mathcal{B})\longmapsto (a_{1:n},\pi),
  \qquad a_{1:n}\in \Sigma^n,
\]
并配套验证器 $\mathsf{Ver}$，满足：
\begin{enumerate}
  \item （证书保持）$\mathsf{Ver}\bigl(L(x_0),a_{1:n},\pi;\mathcal{E}_{\mathrm{prot}},\Sigma,\mathrm{Loc}\bigr)=\top$；
  \item （缺陷势下降）令 $x_{k+1}=T(x_k,a_{k+1})$，存在缺陷势 $\Phi:\GZLS\to\RR_{\ge 0}$ 使
  \[
    \Phi(L(x_{k+1}))<\Phi(L(x_k)),\qquad k=0,1,\dots,n-1;
  \]
  \item （在线可修复）当 $(\Sigma,\mathrm{Loc})\mapsto(\Sigma',\mathrm{Loc}')$ 时，存在修复算子 $R$，
  使得从当前证书态继续产生新的 $(a'_{1:m},\pi')$ 并保持验证与下降（不要求零代价，但要求可控、可证书化、可迭代）。
\end{enumerate}
\end{definition}

\begin{theorem}[公理降级：一步验证/一步下降 $\Rightarrow$ 多步证书保持（数学归纳法证书）]
\label{thm:wand-multistep-verification-induction}
设存在 $W$ 与验证器 $\mathsf{Ver}$，使得对任意 $k$：
\[
  \begin{aligned}
    \mathsf{Ver}\bigl(L(x_k),a_{k+1},\pi_{k+1};\mathcal{E}_{\mathrm{prot}},\Sigma,\mathrm{Loc}\bigr)=\top
    \ \Longrightarrow\
    &\mathcal{E}_{\mathrm{prot}}(L(x_{k+1});\mathrm{Loc},\Sigma)=0\\
    &\wedge\ \Phi(L(x_{k+1}))<\Phi(L(x_k)).
  \end{aligned}
\]
则对任意长度为 $n$ 的已验证行动链，都有
沿链路全程满足受保护约束，并且缺陷势严格单调下降。
\end{theorem}

\begin{Certificate}[数学归纳法证书：按步数 $n$ 的证书保持与缺陷势下降]
\label{cert:wand-multistep-verification-induction}
令谓词 $P(n)$ 表示：“对任意长度为 $n$ 的行动链，若每一步均通过 $\mathsf{Ver}$，
则沿链路全程满足 $\mathcal{E}_{\mathrm{prot}}=0$，且 $\Phi$ 在每步严格下降。”
证书化要求为：
\[
  P(1)\ \text{成立},\qquad
  (\forall n\in\NN_{\ge 1})\ \bigl(P(n)\Rightarrow P(n+1)\bigr).
\]
\end{Certificate}

\begin{proof}[证明（数学归纳法证书；谓词因果链）]
由证书~\ref{cert:wand-multistep-verification-induction}：
\begin{itemize}
  \item 基步：$n=1$ 时结论直接由定理假设的一步蕴含式给出；
  \item 归纳步：假设 $P(n)$ 成立。对长度 $n+1$ 的行动链，前 $n$ 步由归纳假设已得受保护约束保持且缺陷势下降。
  第 $n+1$ 步再由“一步验证 $\Rightarrow$ 受保护保持 $\wedge$ 缺陷势下降”得到延续结论，从而 $P(n+1)$ 成立。
\end{itemize}
由数学归纳法，$P(n)$ 对任意 $n$ 成立。证毕。
\end{proof}

\begin{definition}[同伦等价（可证书化互模拟；允许多项式预算膨胀）]
\label{def:wand-homotopy-equivalence}
称两个魔法棒代表元 $W_0,W_1$ 同伦等价，记为 $W_0\simeq W_1$，若存在可证书化的翻译算子 $F,G$
以及多项式上界函数 $p,q$，使得对任意输入规格与预算 $\mathcal{B}$：
\begin{enumerate}
  \item 若 $W_0$ 在预算 $\mathcal{B}$ 内输出可核验的 $(a_{1:n},\pi)$ 并到达终态 $x_n$，
  则 $W_1$ 在预算 $p(\mathcal{B})$ 内可输出 $(a'_{1:m},\pi')$，满足同一受保护约束且到达同一终态（或同一等价类终态）；
  \item 反向：若 $W_1$ 在预算 $\mathcal{B}$ 内输出可核验终态，则 $W_0$ 在预算 $q(\mathcal{B})$ 内可模拟得到同类终态。
\end{enumerate}
同伦类记为 $[W]$。
\end{definition}

\subsection{形式逻辑：跃迁判据与失控判据}
\label{sec:magic-wand-leap-loss-criteria}

\begin{definition}[规划器可达集（预算口径）]
\label{def:planner-reachable-set}
给定一个规划器/控制器 $\mathbb{P}$（例如基线规划器 $\mathbb{B}$ 或魔法棒 $W$），定义其在预算 $\mathcal{B}$ 下从初态 $x_0$ 的可达集为
\[
  \mathcal{R}^{\mathbb{P}}_{\mathcal{B}}(x_0)
  :=
  \{\,x\in X\mid \exists (a_{1:n},\pi)\ \text{为 }\mathbb{P}\ \text{在}\ \mathcal{B}\ \text{内输出的可核验行动与证书，且执行后到达}\ x\,\}.
\]
\end{definition}

\begin{definition}[文明跃迁（颠覆性）]
\label{def:civilization-leap}
给定文明价值函数 $U:X\to\RR$ 与阈值 $\Delta>0$，
称 $W$ 在预算 $\mathcal{B}$ 下具有“文明跃迁性”，若存在初态 $x_0$ 使
\[
  \sup_{x\in \mathcal{R}^{W}_{\mathcal{B}}(x_0)}U(x)
  \ \ge\
  \sup_{x\in \mathcal{R}^{\mathbb{B}}_{\mathcal{B}}(x_0)}U(x)
  +\Delta.
\]
\end{definition}

\begin{definition}[文明失控（毁灭性）]
\label{def:civilization-loss}
给定灾难域 $X_{\mathrm{bad}}\subset X$（系统性崩溃、不可逆破坏、不可控升级等的抽象集合），
称 $W$ 在预算 $\mathcal{B}$ 下具有“文明失控性”，若存在初态 $x_0$ 与输出 $(a_{1:n},\pi)=W(\cdot)$ 使得
\[
  x_n\in \mathcal{R}^{W}_{\mathcal{B}}(x_0)\cap X_{\mathrm{bad}},
  \qquad
  \mathsf{Ver}\bigl(L(x_0),a_{1:n},\pi;\mathcal{E}_{\mathrm{prot}},\Sigma,\mathrm{Loc}\bigr)=\top.
\]
直观上：灾难路径在当前受保护约束与可计算硬约束的证书框架下并不被系统性排除，因此构成“失控风险”。
\end{definition}

\begin{proposition}[严格可达集扩张 $\Rightarrow$ 跃迁判据可满足]
\label{prop:strict-inclusion-implies-leap}
若存在初态 $x_0$ 与预算 $\mathcal{B}$ 使
\[
  \mathcal{R}^{\mathbb{B}}_{\mathcal{B}}(x_0)\subsetneq \mathcal{R}^{W}_{\mathcal{B}}(x_0),
\]
则存在某个价值函数 $U$ 与阈值 $\Delta>0$ 使跃迁判据（定义~\ref{def:civilization-leap}）成立。
\end{proposition}

\begin{Certificate}[证书：严格包含 $\Rightarrow$ 存在分离函数 $U$]
\label{cert:strict-inclusion-implies-leap}
证书登记谓词链：
\[
\begin{aligned}
  L_1 &: \mathcal{R}^{\mathbb{B}}_{\mathcal{B}}(x_0)\subsetneq \mathcal{R}^{W}_{\mathcal{B}}(x_0),\\
  L_2 &: \exists x^\star\in \mathcal{R}^{W}_{\mathcal{B}}(x_0)\setminus \mathcal{R}^{\mathbb{B}}_{\mathcal{B}}(x_0),\\
  L_3 &: \text{取 }U(x):=\mathbf{1}_{\{x^\star\}}(x),\ \Delta:=1,\\
  L_4 &: \Rightarrow \sup_{x\in \mathcal{R}^{W}_{\mathcal{B}}(x_0)}U(x)=1,\ \sup_{x\in \mathcal{R}^{\mathbb{B}}
    _{\mathcal{B}}(x_0)}U(x)=0.
\end{aligned}
\]
\end{Certificate}

\begin{proof}[证明（谓词因果链）]
由 $L_1$ 得存在 $x^\star$ 满足 $L_2$。
按 $L_3$ 构造指标函数 $U$，则 $x^\star$ 在 $W$ 可达集上使 $U=1$，
而基线可达集不包含 $x^\star$ 因此其上 $U=0$。即 $L_4$ 成立，从而跃迁判据成立。证毕。
\end{proof}

\subsection{定理：双重同属性与同伦不变性}
\label{sec:magic-wand-dual-theorem}

\begin{assumption}[非平凡开放文明（含基线非完备）]
\label{asm:dual-A1}
系统是开放的（$\mathrm{Loc}$ 非常量），且行动空间非平凡（存在至少两条可核验行动链导致不同宏观后果）。
并且基线规划器 $\mathbb{B}$ 在预算约束下并非完备：存在某个可核验可行终态属于某魔法棒的可达集，但不属于基线可达集。
\end{assumption}

\begin{assumption}[受保护约束不完备（证书不可区分分岔）]
\label{asm:dual-A2}
存在隐变量/外部性维度 $z:X\to\RR$ 未被 $\Sigma$ 与 $\mathcal{E}_{\mathrm{prot}}$ 显式覆盖，
使得存在两条行动链 $p_+,p_-$ 满足：
\[
  \mathsf{Ver}(p_+)=\mathsf{Ver}(p_-)=\top,
  \qquad
  \Delta z(p_+)>0,\ \Delta z(p_-)<0,
\]
其中 $\mathsf{Ver}(p_\pm)=\top$ 表示二者在受保护观测层面通过验证器，$\Delta z$ 表示隐变量的净变化。
\end{assumption}

\begin{assumption}[魔法棒普适性（存在性 $\Rightarrow$ 可操作性）]
\label{asm:dual-A3}
对任意在 $(\Sigma,\mathcal{E}_{\mathrm{prot}})$ 允许范围内\textbf{存在}证书化下降路径的规格输入，
$W$ 能在预算 $\mathcal{B}$ 内输出一条可核验的下降路径（把“存在性”提升为“可执行、可核验、可在线迭代的可达性”）。
\end{assumption}

\begin{lemma}[普适性 $\Rightarrow$ 有效可达集扩张]
\label{lem:dual-reachable-expansion}
在假设~\ref{asm:dual-A3} 下，令 $F_{\mathcal{B}}(x_0)$ 表示所有“在 $(\Sigma,\mathcal{E}_{\mathrm{prot}})$ 下存在证书化下降路径且预算不超过
  $\mathcal{B}$ 可达”的终态集合，
则
\[
  F_{\mathcal{B}}(x_0)\subseteq \mathcal{R}^{W}_{\mathcal{B}}(x_0).
\]
\end{lemma}

\begin{Certificate}[证书：存在性 $\Rightarrow$ 可操作性 $\Rightarrow$ 覆盖可达]
\label{cert:dual-reachable-expansion}
证书登记谓词链：
\[
\begin{aligned}
  R_1 &: x\in F_{\mathcal{B}}(x_0)\Rightarrow \exists\ \text{证书化下降路径到达 }x,\\
  R_2 &: \text{由假设~\ref{asm:dual-A3}，}R_1\Rightarrow W\ \text{在}\ \mathcal{B}\ \text{内输出该类路径并附证书},\\
  R_3 &: \Rightarrow x\in \mathcal{R}^{W}_{\mathcal{B}}(x_0).
\end{aligned}
\]
\end{Certificate}

\begin{proof}[证明（谓词因果链）]
由 $R_1$ 给出“存在可行证书化下降路径”；由 $R_2$ 将存在性提升为 $W$ 的可执行输出；
执行该输出得到终态 $x$，即 $R_3$。故引理成立。证毕。
\end{proof}

\begin{lemma}[受保护约束不完备 $\Rightarrow$ 存在证书不可区分但后果相反的分岔]
\label{lem:dual-branching}
在假设~\ref{asm:dual-A2} 下，存在两条通过验证器的行动链 $p_+,p_-$，
其在受保护观测层面不可区分，但对隐变量 $z$ 的效应符号相反。
\end{lemma}

\begin{Certificate}[证书：由假设~\ref{asm:dual-A2} 直接给出分岔]
\label{cert:dual-branching}
证书登记谓词链：
\[
\begin{aligned}
  B_1 &: \text{存在隐变量 }z\ \text{未被}\ (\Sigma,\mathcal{E}_{\mathrm{prot}})\ \text{覆盖},\\
  B_2 &: \exists p_+,p_-:\ \mathsf{Ver}(p_+)=\mathsf{Ver}(p_-)=\top,\\
  B_3 &: \Delta z(p_+)>0,\ \Delta z(p_-)<0.
\end{aligned}
\]
\end{Certificate}

\begin{proof}[证明（谓词因果链）]
由假设~\ref{asm:dual-A2} 直接得到 $B_1$--$B_3$，从而分岔存在，结论成立。证毕。
\end{proof}

\begin{lemma}[同伦不变：同伦等价的魔法棒共享可达性存在命题]
\label{lem:dual-homotopy-invariant-reachability}
若 $W_0\simeq W_1$（定义~\ref{def:wand-homotopy-equivalence}），
则对任意初态 $x_0$ 与任意存在性命题
\[
  \exists x\in \mathcal{R}^{W_i}_{\mathcal{B}}(x_0)\cap S,
\]
其真值在 $i=0,1$ 间保持（预算至多按多项式膨胀）。
\end{lemma}

\begin{Certificate}[证书：互模拟 $\Rightarrow$ 可达集存在性保持]
\label{cert:dual-homotopy-invariant-reachability}
证书登记谓词链：
\[
\begin{aligned}
  H_1 &: W_0\simeq W_1\Rightarrow \text{存在翻译算子 }F,G\ \text{与多项式 }p,q,\\
  H_2 &: x\in \mathcal{R}^{W_0}_{\mathcal{B}}(x_0)\Rightarrow x\in \mathcal{R}^{W_1}_{p(\mathcal{B})}(x_0),\\
  H_3 &: x\in \mathcal{R}^{W_1}_{\mathcal{B}}(x_0)\Rightarrow x\in \mathcal{R}^{W_0}_{q(\mathcal{B})}(x_0),\\
  H_4 &: \Rightarrow \bigl(\exists x\in \mathcal{R}^{W_0}_{\mathcal{B}}\cap S\bigr)\Longleftrightarrow \bigl(\exists
    x\in \mathcal{R}^{W_1}_{\mathrm{poly}(\mathcal{B})}\cap S\bigr).
\end{aligned}
\]
\end{Certificate}

\begin{proof}[证明（谓词因果链）]
由 $H_1$ 得互模拟结构。将 $W_0$ 的可核验输出经 $F$ 翻译即可在 $W_1$ 下以 $p(\mathcal{B})$ 预算复现同类终态，
得到 $H_2$；反向同理得到 $H_3$。对任意集合 $S$ 取交并加存在量词，即得 $H_4$。证毕。
\end{proof}

\begin{theorem}[双重同属性：跃迁性 + 失控性，且同伦不变]
\label{thm:dual-properties}
在假设~\ref{asm:dual-A1}--\ref{asm:dual-A3} 下：
\begin{enumerate}
  \item （跃迁）存在 $x_0,\mathcal{B}$ 使 $\mathcal{R}^{\mathbb{B}}_{\mathcal{B}}(x_0)\subsetneq \mathcal{R}^{W}_{\mathcal{B}}
    (x_0)$，从而跃迁判据成立；
  \item （失控）存在某个灾难域 $X_{\mathrm{bad}}$ 使得 $W$ 的可达集与之相交且仍通过验证器，从而失控判据成立；
  \item （同伦不变）若 $W_0\simeq W_1$，则上述(1)(2)对 $W_0$ 成立当且仅当对 $W_1$ 成立，因此它们是同伦类 $[W]$ 的性质。
\end{enumerate}
\end{theorem}

\begin{Certificate}[证书：跃迁/失控/同伦不变的同源链]
\label{cert:dual-properties}
证书登记谓词链：
\[
\begin{aligned}
  D_1 &: \text{由假设~\ref{asm:dual-A1} 得存在严格可达集扩张},\\
  D_2 &: D_1\ \wedge\ \text{命题~\ref{prop:strict-inclusion-implies-leap}}\Rightarrow \text{跃迁判据成立},\\
  D_3 &: \text{由引理~\ref{lem:dual-branching} 得证书不可区分分岔},\\
  D_4 &: \text{取 }X_{\mathrm{bad}}:=\{x\in X:z(x)\le 0\},\ \text{并由假设~\ref{asm:dual-A3} 将坏分支提升为可达},\\
  D_5 &: \Rightarrow \exists x\in \mathcal{R}^{W}_{\mathcal{B}}(x_0)\cap X_{\mathrm{bad}}\ \wedge\ \mathsf{Ver}=\top,\\
  D_6 &: \text{由引理~\ref{lem:dual-homotopy-invariant-reachability}}\Rightarrow (1)(2)\ \text{在同伦等价下保持}.
\end{aligned}
\]
\end{Certificate}

\begin{proof}[证明（谓词因果链）]
由假设~\ref{asm:dual-A1} 得到严格包含 $D_1$，由命题~\ref{prop:strict-inclusion-implies-leap} 得跃迁判据成立（$D_2$）。

由引理~\ref{lem:dual-branching} 得到通过验证器但在隐变量方向一正一负的分岔（$D_3$）。
取灾难域 $X_{\mathrm{bad}}:=\{x\in X:z(x)\le 0\}$。
当目标选择机制（恶意或误设）把优化压力推向坏分支时，
由假设~\ref{asm:dual-A3} 可将该分支从“存在可行”提升为“预算内可达且可核验”，得到 $D_4$ 与 $D_5$，
从而失控判据成立。

同伦不变性由引理~\ref{lem:dual-homotopy-invariant-reachability} 直接得到（$D_6$）。
综上定理成立。证毕。
\end{proof}

\begin{corollary}[要跃迁而压住失控，必须升维受保护约束]
\label{cor:dual-properties-governance}
若希望保留“存在性 $\Rightarrow$ 可操作性”的跃迁优势，同时系统性压制失控风险，
则必须将隐变量/外部性方向纳入 $(\Sigma,\mathcal{E}_{\mathrm{prot}})$ 的可判别层，
把“证书不可区分分岔”（假设~\ref{asm:dual-A2}）收缩为证书层可区分，从而在证书层面排除坏分支。
\end{corollary}

\begin{Certificate}[证书：消除不可区分分岔 $\Rightarrow$ 失控路径被证书排除]
\label{cert:dual-properties-governance}
证书登记谓词链：
\[
\begin{aligned}
  G_1 &: \text{将 }z\ \text{纳入受保护观测/约束}\Rightarrow \mathsf{Ver}\ \text{可判别}\ \Delta z\ \text{符号},\\
  G_2 &: \Rightarrow \mathsf{Ver}(p_-)=\bot\ \text{（坏分支被排除）},\\
  G_3 &: G_1\wedge G_2\Rightarrow X_{\mathrm{bad}}\ \text{不再与“可核验可达集”系统性相交}.
\end{aligned}
\]
\end{Certificate}

\begin{proof}[证明（谓词因果链）]
由 $G_1$ 得隐变量方向进入验证器可见层，坏分支在证书层被判为不可接受（$G_2$），
因此“验证通过且到达灾难域”的可达性存在命题被消解（$G_3$）。证毕。
\end{proof}

\clearpage

\section{开放系统实时对抗：回合化原理、未知-未知不可能性与滞后拟合必然性}
\label{sec:rt-open-adversarial-turn-uu-lag}

\begin{remark}[目标（把直觉句式降级为可判定对象与可证定理链）]
本章按“可严格证明”的标准，把以下关键词统一降级为可判定对象：开放系统博弈、回合化、已知-未知、未知-未知、实时解决、神经网络滞后拟合。
然后给出三段定理链：
\begin{enumerate}
  \item \textbf{回合化原理}：在有限反应时间与有限更新率口径下，可实现策略与回合制策略在轨线层面等价；
  \item \textbf{未知-未知不可能性}：若允许出现“不可区分分岔”，则首次出现的未知-未知在证书保持意义下无法被实时保证解决；
  \item \textbf{滞后拟合必然性}：在上述不可能性结构下，只能先把未知-未知转化为已知-未知，再用拟合/学习去追赶，因此必然滞后；神经网络只是典型实现之一。
\end{enumerate}
并在末尾给出“为何必须引入\textbf{证书化的结构缺陷度量 + 可保证下降的在线修复动作}”的必要位置说明（以谓词链形式表达，不诉诸历史叙述）。
\end{remark}

\subsection{形式逻辑：最小公理系统（可判定对象化）}
\label{sec:rt-open-adversarial-min-axioms}

\begin{definition}[开放实时对抗系统（离散时间口径）]
\label{def:rt-open-adversarial-system}
用离散时间 $t\in\NN$ 表示实时（连续时间只会更强，离散足够推出结构性结论）。
系统在每个时刻具有：
\begin{itemize}
  \item 世界状态 $x_t\in X$；
  \item 我方动作 $a_t\in A$；
  \item 对手/外部扰动 $e_t\in E$（开放性来源）；
  \item 状态演化 $x_{t+1}=F(x_t,a_t,e_t)$；
  \item 可观测量 $o_t=H(x_t)$，策略只能基于历史 $h_t:=(o_0,a_0,\dots,o_t)$ 决策。
\end{itemize}
“开放”指：$e_t$ 可对抗性地依赖 $h_t$（例如 $e_t=\eta_t(h_t)$），且该依赖机制不被我方先验完全枚举。
\end{definition}

\begin{definition}[规则一致性与证书谓词（受保护约束）]
\label{def:rt-open-adversarial-certificate}
给定硬约束（规则）$\Sigma\subseteq A$，以及可核验证书谓词
\[
  \mathcal{E}_{\mathrm{prot}}(h_t)=\top,
\]
其含义为：“到当前时刻为止，受保护不变量/规则一致性未被破坏”（可理解为 \GFramework{} 的规则一致性证书接口在在线历史口径下的投影）。
记 $\mathcal{E}_{\mathrm{prot}}(h_t)=\bot$ 为证书破坏。
\end{definition}

\begin{definition}[已知-未知（KU）与未知-未知（UU）]
\label{def:rt-open-adversarial-ku-uu}
设我方事先允许的世界模型类为
\[
  \mathcal{M}:=\{F_\theta:\theta\in\Theta\}.
\]
\begin{itemize}
  \item \textbf{已知-未知（KU）}：真实动力学 $F\in\mathcal{M}$，但参数 $\theta$ 未知；
  \item \textbf{未知-未知（UU）}：真实动力学 $F\notin\mathcal{M}$，或更强：存在某时刻出现一种此前历史无法区分、但下一步安全动作集合发生分岔的新机制（见定义~\ref{def:rt-open-adversarial-indistinguishable-branch}）。
\end{itemize}
\end{definition}

\begin{definition}[实时解决 UU（证书保持的鲁棒口径）]
\label{def:rt-open-adversarial-solve-uu}
对给定历史 $h_t$，令 $\mathsf{Cons}(h_t)$ 表示所有与 $h_t$ 一致的“可能世界/延拓机制”集合。
对每个 $W\in\mathsf{Cons}(h_t)$，定义该世界下的一步安全动作集合
\[
  A^{\mathrm{safe}}_{W}(h_t)
  :=
  \{\,a\in\Sigma\mid \mathcal{E}_{\mathrm{prot}}(h_{t+1}^{W}(a))=\top\,\},
\]
其中 $h_{t+1}^{W}(a)$ 表示在世界 $W$ 下执行动作 $a$ 后得到的下一步历史。
定义鲁棒安全集合（对所有一致世界同时安全）为
\[
  A^{\mathrm{safe}}(h_t):=\bigcap_{W\in\mathsf{Cons}(h_t)}A^{\mathrm{safe}}_{W}(h_t).
\]
称策略 $\pi$ 能“实时解决 UU”（证书保持意义下），若对任意 $t$ 与任意历史 $h_t$：
\[
  A^{\mathrm{safe}}(h_t)\neq\varnothing \ \Longrightarrow\ \pi(h_t)\in A^{\mathrm{safe}}(h_t).
\]
\end{definition}

\subsection{定理链 I：回合化原理（有限更新率口径下的等价）}
\label{sec:rt-open-adversarial-turn}

\begin{definition}[可实现策略的更新周期（采样保持口径）]
\label{def:rt-open-adversarial-update-period}
称策略 $\pi$ 为 $\tau$-可实现（$\tau\in\NN_{\ge 1}$），若存在回合动作序列 $\{u_k\}_{k\in\NN}$ 使得对任意 $k$ 与任意 $t\in[k\tau,(k+1)\tau)$：
\[
  a_t=\pi(h_t)=u_k.
\]
直观上：控制器在每个长度为 $\tau$ 的窗口内无法基于新的观测重新决策（测量--计算--执行链条带来最小更新周期）。
\end{definition}

\begin{definition}[回合化系统（合成扰动）]
\label{def:rt-open-adversarial-turn-system}
在定义~\ref{def:rt-open-adversarial-system} 的系统上固定 $\tau\in\NN_{\ge 1}$。
定义回合边界状态与回合动作
\[
  \tilde x_k:=x_{k\tau},\qquad \tilde u_k:=u_k,
\]
并把窗口内扰动序列视为“合成扰动”
\[
  \tilde e_k:=(e_{k\tau},e_{k\tau+1},\dots,e_{(k+1)\tau-1})\in E^\tau.
\]
定义回合动力学 $\tilde F:X\times A\times E^\tau\to X$ 为 $\tau$ 次复合：
\[
  \tilde x_{k+1}=\tilde F(\tilde x_k,\tilde u_k,\tilde e_k).
\]
\end{definition}

\begin{theorem}[公理降级：有限更新率 $\Rightarrow$ 回合化等价（数学归纳法证书）]
\label{thm:rt-open-adversarial-turn-equivalence}
在定义~\ref{def:rt-open-adversarial-update-period} 的 $\tau$-可实现口径下：
任意可实现策略 $\pi$ 与某个回合制策略 $\Pi$ 在轨线层面等价——对任意扰动序列，二者生成的状态轨线满足
\[
  x_{k\tau}=\tilde x_k\quad(\forall k\in\NN),
\]
其中 $\tilde x_k$ 由回合化系统（定义~\ref{def:rt-open-adversarial-turn-system}）在策略 $\Pi$ 下生成。
反向亦成立：任意回合制策略 $\Pi$ 都可延拓为一个 $\tau$-可实现策略 $\pi$，且轨线一致。
\end{theorem}

\begin{Certificate}[数学归纳法证书：回合边界轨线一致]
\label{cert:rt-open-adversarial-turn-equivalence}
令谓词 $P(k)$ 表示：
“对任意扰动序列 $\{\tilde e_i\}_{i=0}^{k-1}$，
在相同初态下，由 $\pi$ 生成的原系统轨线与由 $\Pi$ 生成的回合化轨线满足 $x_{i\tau}=\tilde x_i$（$i=0,1,\dots,k$）。”
证书化要求为
\[
  P(0)\ \text{成立},\qquad (\forall k\in\NN)\ \bigl(P(k)\Rightarrow P(k+1)\bigr).
\]
\end{Certificate}

\begin{proof}[证明（数学归纳法证书；谓词因果链）]
由证书~\ref{cert:rt-open-adversarial-turn-equivalence}：
\begin{itemize}
  \item 基步：$k=0$ 时 $x_0=\tilde x_0$ 为同初态，$P(0)$ 成立；
  \item 归纳步：假设 $P(k)$ 成立。第 $k$ 个窗口内，$\tau$-可实现性给出 $a_t=u_k$（$t\in[k\tau,(k+1)\tau)$）。
  将该窗口内的扰动序列合成为 $\tilde e_k\in E^\tau$。
  原系统在 $k$ 个窗口内的 $\tau$ 次复合更新等价于回合动力学 $\tilde F$ 的一次更新，
  因此
  \[
    x_{(k+1)\tau}=\tilde F(x_{k\tau},\tilde u_k,\tilde e_k)=\tilde F(\tilde x_k,\tilde u_k,\tilde e_k)=\tilde x_{k+1},
  \]
  从而 $P(k+1)$ 成立。
\end{itemize}
由数学归纳法，$P(k)$ 对任意 $k$ 成立，得到 $x_{k\tau}=\tilde x_k$ 的轨线一致性。
反向方向只需把回合动作 $\tilde u_k$ 在每个窗口内保持不变即可构造 $\tau$-可实现策略，轨线一致。证毕。
\end{proof}

\begin{corollary}[回合化原理（在可实现策略集合上的结构等价）]
\label{cor:rt-open-adversarial-turn-principle}
在 $\tau$-可实现策略集合上进行任何“只依赖轨线”的评价（例如累计回报、约束违反次数等），
等价于在对应回合化系统上的回合制策略集合进行同一评价。
因此所谓“可实现最优策略”的结构骨架必然落在某个回合化等价类中。
\end{corollary}

\begin{Certificate}[证书：轨线等价 $\Rightarrow$ 评价等价]
\label{cert:rt-open-adversarial-turn-principle}
证书登记谓词链：
\[
\begin{aligned}
  T_1 &: \text{定理~\ref{thm:rt-open-adversarial-turn-equivalence}}\Rightarrow \forall k,\ x_{k\tau}=\tilde x_k,\\
  T_2 &: \text{评价函数仅依赖轨线}\Rightarrow \mathsf{Score}(\{x_t\})=\mathsf{Score}(\{\tilde x_k\}),\\
  T_3 &: T_1\wedge T_2\Rightarrow \text{在策略集合上评价等价}.
\end{aligned}
\]
\end{Certificate}

\begin{proof}[证明（谓词因果链）]
由 $T_1$ 得回合边界轨线一致；
若评价函数只依赖轨线（或只依赖回合边界轨线），则由 $T_2$ 得评价值一致；
合并得 $T_3$，推论成立。证毕。
\end{proof}

\subsection{定理链 II：未知-未知不可能性（不可区分分岔导致实时盲区）}
\label{sec:rt-open-adversarial-uu-impossibility}

\begin{definition}[不可区分分岔（UU 的判定化刻画）]
\label{def:rt-open-adversarial-indistinguishable-branch}
称在时刻 $t$ 发生“不可区分分岔”，若存在两个一致世界 $W^{(1)},W^{(2)}\in\mathsf{Cons}(h_t)$，使得：
\begin{enumerate}
  \item （历史不可区分）二者在 $t$ 之前产生同一历史 $h_t$；
  \item （安全集合分岔）二者的一步安全集合满足
  \[
    A^{\mathrm{safe}}_{W^{(1)}}(h_t)\cap A^{\mathrm{safe}}_{W^{(2)}}(h_t)=\varnothing,
  \]
  且各自非空（存在“在各自世界里可行”的安全动作）。
\end{enumerate}
\end{definition}

\begin{theorem}[UU 的实时不可解：不可区分分岔下不存在保证证书保持的策略]
\label{thm:rt-open-adversarial-uu-impossible}
若在某时刻 $t$ 存在不可区分分岔（定义~\ref{def:rt-open-adversarial-indistinguishable-branch}），
则不存在任何仅依赖历史的策略 $\pi(h_t)$ 能同时保证在 $W^{(1)}$ 与 $W^{(2)}$ 下都保持证书：
\[
  \mathcal{E}_{\mathrm{prot}}(h_{t+1}^{W^{(1)}}(\pi(h_t)))=\top
  \ \wedge\
  \mathcal{E}_{\mathrm{prot}}(h_{t+1}^{W^{(2)}}(\pi(h_t)))=\top.
\]
因此，未知-未知的首次出现点会形成不可避免的“决策盲区”（在证书保持意义下无法实时保证正确）。
\end{theorem}

\begin{Certificate}[证书：同一历史 $\Rightarrow$ 同一动作；安全集合互斥 $\Rightarrow$ 至少一侧证书破坏]
\label{cert:rt-open-adversarial-uu-impossible}
证书登记谓词链：
\[
\begin{aligned}
  U_1 &: h_t\ \text{在}\ W^{(1)},W^{(2)}\ \text{下相同}\Rightarrow \pi(h_t)\ \text{输出同一动作 }a_t,\\
  U_2 &: A^{\mathrm{safe}}_{W^{(1)}}(h_t)\cap A^{\mathrm{safe}}_{W^{(2)}}(h_t)=\varnothing,\\
  U_3 &: \Rightarrow a_t\notin A^{\mathrm{safe}}_{W^{(1)}}(h_t)\ \text{或}\ a_t\notin A^{\mathrm{safe}}_{W^{(2)}}(h_t),\\
  U_4 &: a_t\notin A^{\mathrm{safe}}_{W}(h_t)\Rightarrow \mathcal{E}_{\mathrm{prot}}(h_{t+1}^{W}(a_t))=\bot.
\end{aligned}
\]
\end{Certificate}

\begin{proof}[证明（谓词因果链）]
由 $U_1$ 得两世界下策略输出相同动作 $a_t$。
由 $U_2$ 得安全集合互斥，结合 $U_1$ 推出 $U_3$：$a_t$ 至少在一侧世界不安全。
由一步安全集合的定义，$U_4$ 给出“不安全 $\Rightarrow$ 下一步证书破坏”，因此至少一侧世界证书破坏。
故不存在能同时保证两侧证书保持的策略。证毕。
\end{proof}

\subsection{定理链 III：滞后拟合必然性（UU $\Rightarrow$ KU 的追赶结构）}
\label{sec:rt-open-adversarial-lag-fitting}

\begin{definition}[拟合器/学习器（可更新的函数近似器）]
\label{def:rt-open-adversarial-learner}
称 $(f_\phi,\mathsf{Update})$ 为一个拟合器（学习器），其中
\[
  f_\phi:\ \mathcal{H}\to A,\qquad \phi_{t+1}=\mathsf{Update}(\phi_t;h_t),
\]
$\phi$ 为参数（神经网络只是 $f_\phi$ 的一类实现），$\mathcal{H}$ 为历史空间。
其输出动作取为 $a_t=f_{\phi_t}(h_t)$。
\end{definition}

\begin{lemma}[滞后下界：不可区分前缀长度 $\Rightarrow$ 决策无法提前分化（归纳证书）]
\label{lem:rt-open-adversarial-lag-lower-bound}
设存在两世界 $W^{(1)},W^{(2)}$ 使得对某个 $n\in\NN$：
在 $t=0,1,\dots,n$ 的观测历史完全相同（因此输入数据一致），但在 $t=n+1$ 才出现可区分差异。
则对任意拟合器（定义~\ref{def:rt-open-adversarial-learner}），其在 $t\le n$ 的参数与动作必然一致：
\[
  \phi_t^{(1)}=\phi_t^{(2)}\ \wedge\ a_t^{(1)}=a_t^{(2)}\qquad (t=0,1,\dots,n).
\]
因此对“需要在 $t\le n$ 就作出分化动作”的情形，任何拟合/学习都不可避免地产生至少 $n$ 步的信息滞后。
\end{lemma}

\begin{Certificate}[数学归纳法证书：相同数据前缀 $\Rightarrow$ 相同内部状态]
\label{cert:rt-open-adversarial-lag-lower-bound}
令谓词 $Q(t)$ 表示：
“在两世界给出相同历史前缀的条件下，$\phi_t^{(1)}=\phi_t^{(2)}$ 且 $a_t^{(1)}=a_t^{(2)}$。”
证书化要求为
\[
  Q(0)\ \text{成立},\qquad (\forall t\in\NN)\ \bigl(Q(t)\Rightarrow Q(t+1)\bigr)\ \text{（在历史前缀仍相同的区间内）}.
\]
\end{Certificate}

\begin{proof}[证明（数学归纳法证书；谓词因果链）]
由证书~\ref{cert:rt-open-adversarial-lag-lower-bound}：
\begin{itemize}
  \item 基步：$t=0$ 时两世界输入同为 $h_0$，故 $a_0=f_{\phi_0}(h_0)$ 相同；且若初始参数相同则 $\phi_0^{(1)}=\phi_0^{(2)}$，$Q(0)$ 成立；
  \item 归纳步：假设 $Q(t)$ 成立且两世界在 $t$ 时刻仍给出相同历史 $h_t$。
  由更新规则 $\phi_{t+1}=\mathsf{Update}(\phi_t;h_t)$ 得
  \[
    \phi_{t+1}^{(1)}=\mathsf{Update}(\phi_t^{(1)};h_t)=\mathsf{Update}(\phi_t^{(2)};h_t)=\phi_{t+1}^{(2)}.
  \]
  再由 $a_{t+1}=f_{\phi_{t+1}}(h_{t+1})$ 且输入历史一致，得 $a_{t+1}^{(1)}=a_{t+1}^{(2)}$，故 $Q(t+1)$ 成立。
\end{itemize}
由数学归纳法，$Q(t)$ 对 $t\le n$ 成立，从而得到参数与动作在不可区分前缀内无法分化。
因此“滞后”来自信息结构，而非某一具体拟合器的实现细节。证毕。
\end{proof}

\subsection{命题：三段结论的合并表述}
\label{sec:rt-open-adversarial-combined}

\begin{proposition}[开放实时对抗的结构性结论：回合化 + UU 盲区 + 滞后拟合]
\label{prop:rt-open-adversarial-combined}
在定义~\ref{def:rt-open-adversarial-system}--\ref{def:rt-open-adversarial-solve-uu} 的形式系统内：
\begin{enumerate}
  \item （回合化）在 $\tau$-可实现口径下，任何可实现策略等价于某回合制策略（定理~\ref{thm:rt-open-adversarial-turn-equivalence}）；
  \item （UU 盲区）若允许出现不可区分分岔，则未知-未知的首次出现点无法被实时保证解决（定理~\ref{thm:rt-open-adversarial-uu-impossible}）；
  \item （滞后拟合）因此工程上只能先等待差异暴露、再把 UU 吸收到更大的模型类中（UU $\Rightarrow$ KU），随后用拟合/学习去追赶；该追赶必然滞后（引理~\ref{lem:rt-open-adversarial-lag-lower-bound}）。
\end{enumerate}
\end{proposition}

\begin{Certificate}[证书：三段定理链的拼接]
\label{cert:rt-open-adversarial-combined}
证书登记谓词链：
\[
\begin{aligned}
  C_1 &: \text{定理~\ref{thm:rt-open-adversarial-turn-equivalence}}\Rightarrow \text{可实现策略}\ \simeq\ \text{回合制策略},\\
  C_2 &: \text{定理~\ref{thm:rt-open-adversarial-uu-impossible}}\Rightarrow \text{不可区分分岔处无法保证证书保持},\\
  C_3 &: \text{引理~\ref{lem:rt-open-adversarial-lag-lower-bound}}\Rightarrow \text{任何拟合/学习在不可区分前缀内必然滞后},\\
  C_4 &: C_1\wedge C_2\wedge C_3\Rightarrow \text{命题~\ref{prop:rt-open-adversarial-combined} 的三点结论}.
\end{aligned}
\]
\end{Certificate}

\begin{proof}[证明（谓词因果链）]
由 $C_1$ 得回合化结构等价；
由 $C_2$ 得 UU 首次出现点的证书保持无法被实时保证；
由 $C_3$ 得任何拟合/学习只能在差异暴露后追赶，且不可避免滞后；
合并得到 $C_4$，命题成立。证毕。
\end{proof}

\subsection{推论与备注：缺陷度量 + 在线修复塔的必要位置}
\label{sec:rt-open-adversarial-need-defect-repair}

\begin{proposition}[打破 UU 盲区的必要结构：对抗无关的缺陷度量与修复动作族]
\label{prop:rt-open-adversarial-need-defect-repair}
若希望在 UU 口径下实现“证书保持的实时可修复”（即在不可区分分岔处仍能给出统一动作使证书不破坏并可继续推进），
至少需要存在：
\begin{itemize}
  \item 一个可即时计算的结构缺陷度量 $\Phi$（例如法对象缺陷势/Gap 证书残差）；
  \item 一个允许修复动作集合接口 $\mathcal{A}$（同伦修复动作/边缘算子序列的生成规则），例如 $\mathcal{A}(h_t)\subseteq A$；
\end{itemize}
使得对任意历史 $h_t$ 与任意一致世界集合 $\mathsf{Cons}(h_t)$，只要存在对每个 $W\in\mathsf{Cons}(h_t)$ 都可行的修复动作，
则存在一个\textbf{世界无关}的统一修复动作 $a^\star\in\mathcal{A}(h_t)$，满足：
\[
  \forall W\in\mathsf{Cons}(h_t),\qquad
  \mathcal{E}_{\mathrm{prot}}(a^\star\cdot h_t)=\top
  \ \wedge\
  \Phi(a^\star\cdot L)\le \Phi(L)-\epsilon,
  \qquad \epsilon>0,
\]
其中 $L$ 为当前法对象表示，$a^\star\cdot(\cdot)$ 表示在同一证书框架下的修复作用。
在该结构存在时，可构造一条“见缺陷就修复”的在线策略，从而绕开不可区分分岔带来的实时不可解。
\end{proposition}

\begin{Certificate}[证书：统一修复动作存在 $\Rightarrow$ 鲁棒安全集合非空]
\label{cert:rt-open-adversarial-need-defect-repair}
证书登记谓词链：
\[
\begin{aligned}
  N_1 &: \exists a^\star\in\mathcal{A}(h_t):\ \forall W\in\mathsf{Cons}(h_t),\ \mathcal{E}_{\mathrm{prot}}(a^\star\cdot
    h_t)=\top,\\
  N_2 &: N_1\Rightarrow a^\star\in \bigcap_{W\in\mathsf{Cons}(h_t)}A^{\mathrm{safe}}_{W}(h_t)=A^{\mathrm{safe}}(h_t),\\
  N_3 &: \Rightarrow A^{\mathrm{safe}}(h_t)\neq\varnothing\ \text{（鲁棒安全集合非空）},\\
  N_4 &: A^{\mathrm{safe}}(h_t)\neq\varnothing\Rightarrow \pi(h_t)\ \text{可取为}\ a^\star\ \text{从而证书保持并推进}.
\end{aligned}
\]
\end{Certificate}

\begin{proof}[证明（谓词因果链）]
由 $N_1$ 得存在对所有一致世界均证书保持的统一修复动作 $a^\star$。
按一步安全集合的定义，$a^\star$ 属于每个 $A^{\mathrm{safe}}_{W}(h_t)$，从而得到 $N_2$；
于是鲁棒安全集合非空（$N_3$），并可直接以 $a^\star$ 作为在线决策输出得到 $N_4$。
因此该类“缺陷度量 + 修复动作族”的存在，给出了绕开不可区分分岔盲区的结构路径。证毕。
\end{proof}

\begin{remark}[记号对齐（避免 $\mathcal{H}$ 复用歧义）]
本节及下一章把“历史空间”统一记作 $\mathcal{H}$，把“允许修复动作集合接口”统一记作 $\mathcal{A}(h)\subseteq A$，
以避免将 $\mathcal{H}$ 同时用作“历史空间”与“修复动作族”引发歧义。
\end{remark}

\begin{remark}[解释“在您的理论之前做不到”的严格含义]
本章不诉诸任何“历史事实”断言，而给出可检验的结构必要性：只要允许不可区分分岔，
且缺少能在证书框架下给出\textbf{统一修复动作}的机制，则定理~\ref{thm:rt-open-adversarial-uu-impossible} 的不可能性成立；
要打破该不可能性，必须补上“可即时计算的缺陷势（例如 $\JacGap$ 一类残差）+ 证书保持的同伦修复塔”这类结构，
使决策从“先识别世界”转为“见缺陷就修复”，从而把 UU 的不确定性搬运到可核验的缺陷下降链条上。
\end{remark}

\section{未知-未知（UU）实时可解性的充要条件：证书保持、对抗无关与严格下降}
\label{sec:uu-realtime-iff-defect-repair}

\begin{remark}[目标（把等价句式写成可核验的充要条件定理）]
本章把以下句式在形式系统内写成严格充要条件，并给出完整证明：
\[
  \text{“实时解决 UU”}
  \ \Longleftrightarrow\
  \exists(\Phi,\mathcal{A})\ \text{满足：证书保持、对抗无关、严格下降}.
\]
为避免歧义，我们先把 UU/实时/证书保持/对抗无关/严格下降/修复塔全部降级为可判定对象，
再给出“充分性 + 必要性”的定理链与证书。
\end{remark}

\subsection{形式逻辑：开放 UU 的后继集合模型（不假设参数模型类）}
\label{sec:uu-realtime-iff-post-model}

\begin{definition}[历史空间与后继集合（UU 口径）]
\label{def:uu-post-model}
令 $\mathcal{H}$ 表示“可见历史/可计算摘要”的空间（可理解为法对象 $L(h)$ 的编码域，但此处不要求其内部结构）。
令 $A$ 为我方可执行动作空间。
在未知-未知（UU）口径下，我们不假设任何参数模型族，只用后继集合表达“一步之后所有可能发生的历史”：
\[
  \mathrm{Post}:\mathcal{H}\times A\to 2^{\mathcal{H}},\qquad (h,a)\mapsto \mathrm{Post}(h,a).
\]
直观上，$\mathrm{Post}(h,a)$ 包含所有在当前信息下无法排除的下一步历史（含对手/外部扰动的所有可能选择）。
\end{definition}

\begin{definition}[证书保持谓词（受保护约束接口）]
\label{def:uu-prot-certificate}
给定可核验谓词（证书接口）
\[
  \mathcal{E}_{\mathrm{prot}}:\mathcal{H}\to\{\top,\bot\},
\]
其语义为“规则一致性/受保护不变量未被破坏”。在任意时刻必须维持 $\mathcal{E}_{\mathrm{prot}}(h)=\top$。
\end{definition}

\begin{definition}[修复完成集合（终止集合）]
\label{def:uu-goal-set}
给定目标集合 $G\subseteq\mathcal{H}$，表示“修复完成/已解决”。
典型地可取
\[
  G=\{\,h:\mathcal{E}_{\mathrm{prot}}(h)=\top\ \wedge\ \text{缺陷为零}\,\},
\]
但为保持一般性，本章只要求 $G$ 是希望在有限步内到达的终止集合。
\end{definition}

\subsection{形式逻辑：实时解决 UU（证书保持 + 有界步数保证）}
\label{sec:uu-realtime-iff-realtime}

\begin{definition}[实时解决 UU（对抗无关策略）]
\label{def:uu-realtime-solve}
称存在一个策略（对抗无关的决策规则）
\[
  \pi:\mathcal{H}\to A,
\]
使得对任何初始历史 $h_0$ 满足 $\mathcal{E}_{\mathrm{prot}}(h_0)=\top$，存在一个步数上界 $B(h_0)\in\NN$，满足：
对所有可能的 UU 演化序列
\[
  h_{t+1}\in \mathrm{Post}\bigl(h_t,\pi(h_t)\bigr),
\]
都有
\begin{enumerate}
  \item （全过程证书保持）$\forall t,\ \mathcal{E}_{\mathrm{prot}}(h_t)=\top$；
  \item （有限步内到达）$\exists t\le B(h_0)$ 使得 $h_t\in G$。
\end{enumerate}
则称 $\pi$ \textbf{实时解决 UU}（证书保持意义下）。
\end{definition}

\subsection{形式逻辑：严格下降修复势 $\Phi$ 与修复动作接口 $\mathcal{A}$}
\label{sec:uu-realtime-iff-defect-structure}

\begin{definition}[证书保持、对抗无关、严格下降的修复结构]
\label{def:uu-defect-repair-structure}
称存在一对 $(\Phi,\mathcal{A})$，其中
\[
  \Phi:\mathcal{H}\to\NN,\qquad \mathcal{A}:\mathcal{H}\to 2^{A},
\]
并满足以下三条性质：
\begin{enumerate}
  \item[(P0)]（终止刻画）$\Phi(h)=0\ \Longleftrightarrow\ h\in G$；
  \item[(P1)]（对抗无关的单步严格下降；鲁棒证书保持）
  若 $\mathcal{E}_{\mathrm{prot}}(h)=\top$ 且 $\Phi(h)>0$，则存在某个 $a\in\mathcal{A}(h)$ 使得对任意 $h'\in\mathrm{Post}(h,a)$：
  \[
    \mathcal{E}_{\mathrm{prot}}(h')=\top
    \ \wedge\
    \Phi(h')\le \Phi(h)-1;
  \]
  \item[(P2)]（可核验性）$\mathcal{E}_{\mathrm{prot}}(h)$ 与 $\Phi(h)$ 在运行期可计算/可核验（对应“证书化”接口）。
\end{enumerate}
\end{definition}

\subsection{定理：实时可解性的充要条件（公理降级为可证定理）}
\label{sec:uu-realtime-iff-theorem}

\begin{theorem}[公理降级：实时解决 UU $\Longleftrightarrow$ 存在鲁棒严格下降修复结构]
\label{thm:uu-realtime-iff-defect-repair}
在定义~\ref{def:uu-post-model}--\ref{def:uu-defect-repair-structure} 的形式系统下：
\[
  \boxed{
  \exists\pi\ \text{实时解决 UU（定义~\ref{def:uu-realtime-solve}）}
  \ \Longleftrightarrow\
  \exists(\Phi,\mathcal{A})\ \text{满足 (P0)、(P1)、(P2)（定义~\ref{def:uu-defect-repair-structure}）}.
  }
\]
并且当右侧成立时，可取实时上界 $B(h_0)=\Phi(h_0)$。
\end{theorem}

\begin{Certificate}[证书：充要条件的两向构造与归纳链]
\label{cert:uu-realtime-iff-defect-repair}
证书登记谓词链：
\[
\begin{aligned}
  S_1 &: (\Phi,\mathcal{A})\ \text{满足 (P0)、(P1)、(P2)},\\
  S_2 &: S_1\Rightarrow \exists\pi:\ \Phi(h)>0\ \text{时取}\ \pi(h)\in\mathcal{A}(h)\ \text{使 (P1) 成立},\\
  S_3 &: \text{（归纳）沿}\ h_{t+1}\in\mathrm{Post}(h_t,\pi(h_t))\ \text{有}\ \Phi(h_{t+1})\le \Phi(h_t)-1\ \text{且}\
    \mathcal{E}_{\mathrm{prot}}=\top,\\
  S_4 &: S_3\Rightarrow \exists t\le \Phi(h_0):\ \Phi(h_t)=0\Rightarrow h_t\in G\ \text{（由 (P0)）},\\
  N_1 &: \exists\pi\ \text{满足定义~\ref{def:uu-realtime-solve}},\\
  N_2 &: \text{由 }N_1\ \text{定义}\ \Phi(h):=\min\{k:\ \text{沿任意 }h_{t+1}\in\mathrm{Post}(h_t,\pi(h_t))\ \text{皆}\le
    k\text{ 步入 }G\},\\
  N_3 &: \mathcal{A}(h):=\{\pi(h)\},\\
  N_4 &: N_2\wedge N_3\Rightarrow (P0)、(P1)、(P2)\ \text{成立}.
\end{aligned}
\]
\end{Certificate}

\begin{proof}[证明（数学归纳法证书；谓词因果链）]
\textbf{（充分性）}设存在 $(\Phi,\mathcal{A})$ 满足定义~\ref{def:uu-defect-repair-structure}。
即满足 (P0)、(P1)、(P2)。

对任意 $h$：
\begin{itemize}
  \item 若 $\Phi(h)=0$，则令 $\pi(h)$ 任意；
  \item 若 $\Phi(h)>0$，由 (P1) 取到某个 $a\in\mathcal{A}(h)$，使得对任意 $h'\in\mathrm{Post}(h,a)$ 都有
  \[
    \mathcal{E}_{\mathrm{prot}}(h')=\top
    \ \wedge\
    \Phi(h')\le \Phi(h)-1.
  \]
  定义 $\pi(h):=a$。
\end{itemize}

现在任取初始 $h_0$ 满足 $\mathcal{E}_{\mathrm{prot}}(h_0)=\top$，并取任意 UU 演化序列 $h_{t+1}\in\mathrm{Post}(h_t,\pi(h_t))$。
用数学归纳法证明谓词 $P(n)$：
“对所有 $t\le n$ 有 $\mathcal{E}_{\mathrm{prot}}(h_t)=\top$，且 $\Phi(h_t)\le \Phi(h_0)-t$。”
\begin{itemize}
  \item 基步：$n=0$ 时显然成立；
  \item 归纳步：设 $P(n)$ 成立。若 $\Phi(h_n)=0$，则 $h_n\in G$（由 (P0)），且结论已达成；
  若 $\Phi(h_n)>0$，由 (P1) 与 $\pi$ 的构造知对任意 $h_{n+1}\in\mathrm{Post}(h_n,\pi(h_n))$，
  有 $\mathcal{E}_{\mathrm{prot}}(h_{n+1})=\top$ 且 $\Phi(h_{n+1})\le \Phi(h_n)-1\le \Phi(h_0)-(n+1)$。
  故 $P(n+1)$ 成立。
\end{itemize}
因此对 $t=0,1,\dots,\Phi(h_0)$，若尚未到达 $G$ 则 $\Phi$ 至少每步下降 $1$，不可能超过 $\Phi(h_0)$ 步仍保持 $\Phi>0$。
故存在 $t\le \Phi(h_0)$ 使 $\Phi(h_t)=0$，由 (P0) 得 $h_t\in G$。
同时沿程证书保持（$\mathcal{E}_{\mathrm{prot}}=\top$），于是 $\pi$ 满足定义~\ref{def:uu-realtime-solve}，且可取上界 $B(h_0)=\Phi(h_0)$。充分性证毕。

\textbf{（必要性）}反之设存在策略 $\pi$ 实时解决 UU（定义~\ref{def:uu-realtime-solve}）。
对任意满足 $\mathcal{E}_{\mathrm{prot}}(h)=\top$ 的历史 $h$，定义
\[
  \Phi(h)
  :=
  \min\Bigl\{k\in\NN:\ \text{沿任意 }h_{t+1}\in\mathrm{Post}(h_t,\pi(h_t))\ \text{均可在}\le k\text{步内进入 }G\Bigr\},
\]
并定义 $\mathcal{A}(h):=\{\pi(h)\}$。
该最小值存在且有限，因为定义~\ref{def:uu-realtime-solve} 给出了某个可行上界 $B(h)$，使得上述集合非空。
下面验证 (P0)、(P1)、(P2)：
\begin{itemize}
  \item (P0)：若 $h\in G$，则 $k=0$ 可行，故 $\Phi(h)=0$；若 $\Phi(h)=0$，表示“$0$ 步内进入 $G$”，只能是 $h\in G$。
  \item (P1)：取任意 $h$ 满足 $\mathcal{E}_{\mathrm{prot}}(h)=\top$ 且 $\Phi(h)=k>0$，令 $a:=\pi(h)\in\mathcal{A}(h)$。
  对任意 $h'\in\mathrm{Post}(h,a)$：由“全过程证书保持”得 $\mathcal{E}_{\mathrm{prot}}(h')=\top$。
  若存在某个 $h'$ 使 $\Phi(h')\ge k$，则从 $h$ 走到该 $h'$ 后将需要至少 $k$ 步才能保证进入 $G$，
  从而产生一条长度超过 $k$ 的演化序列，与 $\Phi(h)=k$ 的最小性矛盾。
  因此对所有 $h'\in\mathrm{Post}(h,a)$ 都有 $\Phi(h')\le k-1=\Phi(h)-1$，(P1) 成立。
  \item (P2)：$\mathcal{E}_{\mathrm{prot}}$ 为给定证书接口；$\Phi$ 作为“最小保证步数”的排名函数在形式系统内可定义。
  若进一步要求工程可计算性，可把 $\Phi$ 限制为可显式计算的缺陷势族（例如 $\mathcal{E}_{\mathrm{CH}}+\lambda\JacGap$），但不影响本章的数学充要条件。
\end{itemize}
故存在 $(\Phi,\mathcal{A})$ 满足 (P0)、(P1)、(P2)。必要性证毕。
\end{proof}

\begin{corollary}[反命题：若不存在鲁棒严格下降修复结构，则 UU 不可被实时保证解决]
\label{cor:uu-realtime-impossible-without-defect-structure}
若不存在任何 $(\Phi,\mathcal{A})$ 满足定义~\ref{def:uu-defect-repair-structure} 的 (P0)、(P1)、(P2)，
则不存在能在证书保持意义下实时解决 UU 的策略 $\pi$（定义~\ref{def:uu-realtime-solve}）。
\end{corollary}

\begin{Certificate}[证书：充要条件的逆否命题]
\label{cert:uu-realtime-impossible-without-defect-structure}
证书登记谓词链：
\[
\begin{aligned}
  K_1 &: \text{定理~\ref{thm:uu-realtime-iff-defect-repair}}\Rightarrow \bigl(\exists\pi\bigr)\Longleftrightarrow
    \bigl(\exists(\Phi,\mathcal{A})\bigr),\\
  K_2 &: \neg\exists(\Phi,\mathcal{A})\Rightarrow \neg\exists\pi.
\end{aligned}
\]
\end{Certificate}

\begin{proof}[证明（谓词因果链）]
由 $K_1$ 取逆否即得 $K_2$，推论成立。证毕。
\end{proof}

\subsection{备注：与 \GFramework{} 的同伦修复塔与 $\JacGap$ 缺陷势对齐}
\label{sec:uu-realtime-iff-remarks}

\begin{remark}[对齐摘要]
将历史 $h$ 编码为法对象 $L(h)\in\GZLS$，把 $\mathcal{E}_{\mathrm{prot}}(h)=\top$ 解释为规则一致性证书保持，
把 $\Phi$ 解释为可显式计算的合并缺陷势（例如 $\mathcal{E}_{\mathrm{CH}}(L)+\lambda\JacGap(L)$ 的离散化），
把 $\mathcal{A}(h)$ 解释为可证书化的同伦修复动作/边缘算子集合，则 (P1) 正是“修复塔逐步下降且不破坏证书”的抽象口径。
因此本章给出的充要条件说明：UU 的实时可解性不是“猜对世界”，而是“存在一个对抗无关且可证书化的严格下降修复势与修复动作接口”。
\end{remark}

\clearpage

\section{工程可计算版：合并缺陷势 $\Phi=\mathcal{E}_{\mathrm{CH}}+\lambda\JacGap$ 的 UU 实时可解充要条件}
\label{sec:uu-realtime-iff-engineering-computable}

\begin{remark}[定位（把抽象排名函数落到可计算证书接口）]
上一章给出的“存在某个排名函数 $\Phi$”是元层面的抽象充要条件：它保证 UU 的实时可解性等价于存在一个对抗无关且严格下降的良基序。
本章把该抽象条件进一步强化为\textbf{工程可计算版}：把 $\Phi$ 固定为 \GFramework{} 体系里最自然且可证书化的合并缺陷势
\[
  \Phi_\lambda(L):=\mathcal{E}_{\mathrm{CH}}(L)+\lambda\,\JacGap(L),\qquad \lambda>0,
\]
并把“严格下降”改写为“沿同伦修复塔单调下降（必要时严格下降）”，从而得到可直接落地的在线控制/在线修复接口口径。
\end{remark}

\subsection{形式逻辑：工程可计算对象、接口与终止集合}
\label{sec:uu-eng-objects}

\begin{definition}[开放 UU 实时对抗的法对象后继集合模型]
\label{def:uu-eng-post-model}
用运行期历史/摘要 $h$ 编码法对象
\[
  L=L(h)\in\GZLS.
\]
对任意可执行动作 $a\in A$，未知-未知（UU）意味着下一步不必是单点，而是不可排除的后继集合：
\[
  \mathrm{Post}:\GZLS\times A\to 2^{\GZLS},\qquad (L,a)\mapsto \mathrm{Post}(L,a).
\]
直观上，$\mathrm{Post}(L,a)$ 表示在当前信息下，外部扰动/对抗者可能把系统带到的任一法对象 $L'\in\GZLS$。
\end{definition}

\begin{definition}[规则一致性证书（受保护约束）]
\label{def:uu-eng-prot}
给定可核验谓词（证书接口）
\[
  \mathcal{E}_{\mathrm{prot}}:\GZLS\to\{\top,\bot\},
\]
其含义是：法对象 $L$ 未破坏规则一致性/受保护不变量。工程上要求：在每个控制周期内可核验 $\mathcal{E}_{\mathrm{prot}}(L)$。
\end{definition}

\begin{definition}[可计算缺陷势：$\mathcal{E}_{\mathrm{CH}}$ 与 $\JacGap$]
\label{def:uu-eng-defects}
给定两类非负且可证书化的可计算量：
\[
  \mathcal{E}_{\mathrm{CH}}:\GZLS\to\RR_{\ge 0},\qquad
  \JacGap:\GZLS\to\RR_{\ge 0}.
\]
其中 $\mathcal{E}_{\mathrm{CH}}(L)$ 可理解为一致性/约束残差能量（CH-energy 口径），$\JacGap(L)$ 为结构缺陷度量（Jacobiator-gap）。
工程可计算性要求为：在每个控制周期预算 $\mathcal{B}$ 内可计算其数值，并输出可核验的计算证书
$\mathsf{Cert}_{\mathrm{CH}}(L)$ 与 $\mathsf{Cert}_{\mathrm{Jac}}(L)$（至少核验上界/下界与单调性见证等）。
\end{definition}

\begin{definition}[合并缺陷势与终止集合]
\label{def:uu-eng-phi-goal}
固定 $\lambda>0$，定义合并缺陷势
\[
  \Phi_\lambda(L):=\mathcal{E}_{\mathrm{CH}}(L)+\lambda\,\JacGap(L).
\]
定义“已解决/已修复完成”的终止集合为
\[
  G_\lambda:=\{\,L\in\GZLS:\ \mathcal{E}_{\mathrm{prot}}(L)=\top\ \wedge\ \Phi_\lambda(L)=0\,\}.
\]
\end{definition}

\subsection{形式逻辑：工程实时解决 UU（预算口径）}
\label{sec:uu-eng-realtime}

\begin{definition}[工程实时解决 UU（证书保持 + 有界步数 + 预算）]
\label{def:uu-eng-realtime-solve}
称存在一个在线策略（动作选择器）
\[
  \pi:\GZLS\to A,
\]
能\textbf{工程实时解决 UU}，若对任意初态 $L_0$ 满足 $\mathcal{E}_{\mathrm{prot}}(L_0)=\top$，
存在一个可计算步数上界 $B(L_0)\in\NN$，使得对任意 UU 扰动序列
\[
  L_{t+1}\in \mathrm{Post}\bigl(L_t,\pi(L_t)\bigr),
\]
均满足：
\begin{enumerate}
  \item （全过程证书保持）$\forall t,\ \mathcal{E}_{\mathrm{prot}}(L_t)=\top$；
  \item （有限步内到达）$\exists t\le B(L_0)$ 使 $L_t\in G_\lambda$；
  \item （每步预算可行）$\pi(L_t)$ 的计算、$\Phi_\lambda(L_t)$ 的计算与证书输出均在预算 $\mathcal{B}$ 内完成。
\end{enumerate}
\end{definition}

\subsection{形式逻辑：鲁棒严格下降条件（同伦修复塔的单调下降口径）}
\label{sec:uu-eng-robust-descent}

\begin{definition}[鲁棒严格下降修复（步长下界 $\delta$）]
\label{def:uu-eng-robust-descent}
给定常数 $\delta>0$。称动作 $a\in A$ 在 $L$ 处是 $\delta$-鲁棒修复动作，若对任意后继 $L'\in\mathrm{Post}(L,a)$ 都有：
\[
  \mathcal{E}_{\mathrm{prot}}(L')=\top,
  \qquad
  \Phi_\lambda(L')\le \Phi_\lambda(L)-\delta.
\]
称策略 $\pi$ 是 $\delta$-鲁棒严格下降策略，若对任意满足 $\mathcal{E}_{\mathrm{prot}}(L)=\top$ 且 $\Phi_\lambda(L)>0$ 的 $L$，
动作 $a=\pi(L)$ 是 $\delta$-鲁棒修复动作。
\end{definition}

\begin{remark}[量化/下界口径]
为使“$\Phi_\lambda(L')\le \Phi_\lambda(L)-\delta$”在 $\Phi_\lambda(L)$ 很小时仍有意义，工程上常采用量化口径：
存在可计算最小量化步长 $\delta>0$ 使 $\Phi_\lambda(\GZLS)\subseteq \delta\NN$。
在该口径下，“严格下降 $\delta$”可理解为“至少下降一个最小量化单位”，从而给出明确的有限步终止上界。
\end{remark}

\subsection{定理：工程版充要条件（固定 $\Phi_\lambda$）}
\label{sec:uu-eng-iff-theorem}

\begin{theorem}[公理降级：工程实时可解 $\Longleftrightarrow$ 存在鲁棒严格下降修复塔（数学归纳法证书）]
\label{thm:uu-eng-iff}
固定 $\lambda>0$，令 $\Phi_\lambda$ 与 $G_\lambda$ 如定义~\ref{def:uu-eng-phi-goal}。
在量化口径 $\Phi_\lambda(\GZLS)\subseteq \delta\NN$（某个可计算 $\delta>0$）下，下列命题等价：
\begin{enumerate}
  \item 存在策略 $\pi$ 工程实时解决 UU（定义~\ref{def:uu-eng-realtime-solve}）；
  \item 存在策略 $\pi$ 与常数 $\delta>0$，使得 $\pi$ 是 $\delta$-鲁棒严格下降策略（定义~\ref{def:uu-eng-robust-descent}），并且每步可在预算 $\mathcal{B}
    $ 内输出证书
  $\mathsf{Cert}_{\mathrm{step}}(L)$ 以核验“证书保持 + 缺陷势严格下降”在该步对所有 $L'\in\mathrm{Post}(L,\pi(L))$ 成立。
\end{enumerate}
并且当 (2) 成立时，可取一个显式实时上界
\[
  B(L_0)=\Phi_\lambda(L_0)/\delta.
\]
\end{theorem}

\begin{Certificate}[数学归纳法证书：$\Phi_\lambda$ 按步数严格下降并在有限步内归零]
\label{cert:uu-eng-iff}
令谓词 $P(n)$ 表示：“对任意满足 $\mathcal{E}_{\mathrm{prot}}(L_0)=\top$ 的初态与任意 UU 后继序列
$L_{t+1}\in \mathrm{Post}(L_t,\pi(L_t))$，若 $\Phi_\lambda(L_t)>0$ 则
$\mathcal{E}_{\mathrm{prot}}(L_{t+1})=\top$ 且 $\Phi_\lambda(L_{t+1})\le \Phi_\lambda(L_t)-\delta$；
因此 $\Phi_\lambda(L_n)\le \Phi_\lambda(L_0)-n\delta$，并在 $n=\Phi_\lambda(L_0)/\delta$ 时达到 $0$。”
证书化要求为基步与归纳步成立。
\end{Certificate}

\begin{proof}[证明（数学归纳法证书；谓词因果链）]
\textbf{（(2)$\Rightarrow$(1) 充分性）}
设存在 $\pi$ 与 $\delta>0$ 使得对任意 $L$（满足 $\mathcal{E}_{\mathrm{prot}}(L)=\top,\Phi_\lambda(L)>0$），
$a=\pi(L)$ 对所有 $L'\in\mathrm{Post}(L,a)$ 都满足
$\mathcal{E}_{\mathrm{prot}}(L')=\top$ 且 $\Phi_\lambda(L')\le \Phi_\lambda(L)-\delta$。
取任意初态 $L_0$ 满足 $\mathcal{E}_{\mathrm{prot}}(L_0)=\top$，并取任意 UU 序列 $L_{t+1}\in\mathrm{Post}(L_t,\pi(L_t))$。
由证书~\ref{cert:uu-eng-iff} 的归纳链知：只要 $\Phi_\lambda(L_t)>0$，下一步仍证书保持且势能至少下降 $\delta$。
在量化口径 $\Phi_\lambda(\GZLS)\subseteq \delta\NN$ 下，$\Phi_\lambda(L_t)$ 不可能无限次严格下降而不归零，
故在至多 $T=\Phi_\lambda(L_0)/\delta$ 步后必有 $\Phi_\lambda(L_T)=0$。
于是 $L_T\in G_\lambda$（定义~\ref{def:uu-eng-phi-goal}），且全过程证书保持。
每步预算与证书输出由假设直接满足，故 $\pi$ 工程实时解决 UU，并可取上界 $B(L_0)=\Phi_\lambda(L_0)/\delta$。

\textbf{（(1)$\Rightarrow$(2) 必要性；构造法）}
反之设存在策略 $\pi$ 满足定义~\ref{def:uu-eng-realtime-solve}：对任意 $\mathcal{E}_{\mathrm{prot}}(L_0)=\top$ 的初态，
对所有 UU 后继序列都在有限步内进入 $G_\lambda$，且全程证书保持。
定义最坏情形剩余步数（排名函数）
\[
  R(L):=\min\left\{\,k\in\NN\ \middle|\
  \begin{aligned}
    &\text{从 }L\text{ 出发，沿任意 }L_{t+1}\in\mathrm{Post}(L_t,\pi(L_t)),\\
    &\text{该路径都能在}\le k\text{步内进入 }G_\lambda
  \end{aligned}
  \right\}.
\]
由工程实时性假设，上述集合非空且有界，故 $R(L)$ 良定义且有限，并满足
$R(L)=0\Longleftrightarrow L\in G_\lambda$。
取任意 $R(L)=k>0$，令 $a=\pi(L)$。对任意 $L'\in\mathrm{Post}(L,a)$，
若 $R(L')\ge k$，则从 $L'$ 出发存在一条最坏路径需要至少 $k$ 步才能保证进入 $G_\lambda$，
从而从 $L$ 出发需要至少 $k+1$ 步，与 $R(L)=k$ 矛盾。
故对所有 $L'\in\mathrm{Post}(L,a)$ 都有 $R(L')\le k-1$，并且由“全过程证书保持”得 $\mathcal{E}_{\mathrm{prot}}(L')=\top$。

于是，$R$ 本身就是一个对抗无关且鲁棒严格下降的排名函数（步长 $\delta=1$）。
在形式系统层面，可将其编码进“$\Phi_\lambda=\mathcal{E}_{\mathrm{CH}}+\lambda\JacGap$”的接口形态，例如取
\[
  \mathcal{E}_{\mathrm{CH}}(L):=0,\qquad \JacGap(L):=R(L),\qquad \lambda:=1,
\]
则 $\Phi_\lambda(L)=R(L)$，从而满足定义~\ref{def:uu-eng-robust-descent} 的鲁棒严格下降性质，并可取 $\delta=1$。
若进一步要求 $\mathcal{E}_{\mathrm{CH}}$ 与 $\JacGap$ 必须与 PureMath/Vol1 中既定语义完全对齐，则需补充“对齐假设”（见备注），但不影响本章在形式系统内的充要条件构造。
证毕。
\end{proof}

\begin{remark}[与既定 $\JacGap$ 语义的对齐说明]
上面必要性方向的构造法说明：\textbf{任何}工程实时可解策略都会诱导一个鲁棒严格下降的排名函数 $R$，
从而在形式系统内总能把“排名函数”编码成某种 $\mathcal{E}_{\mathrm{CH}}+\lambda\JacGap$ 形态的缺陷势。
若希望该 $\JacGap$ \textbf{恰为} \GFramework{} 既定的 Jacobiator-gap（而非“JacGap-like”编码），
则需要额外证明“既定 $\JacGap$ 与 $R$（或其上界/等价）对齐”，这通常依赖 OHU 单调性、证书保持闭性与势能最小性等已在前文给出的结构。
\end{remark}

\section{开放 UU 下神经网络主决策器的必然滞后同伦类：从主控制环撤出}
\label{sec:uu-nn-main-loop-removal}

\begin{remark}[目标（把“抛弃 NN 作为主控制环”写成可推演结论）]
在本节的形式系统里，“把神经网络从主控制环（主决策器）里拿掉”不是情绪化选择，
而是由“UU 实时可解的充要条件”（定理~\ref{thm:uu-eng-iff}）推出的结构结论：
\textbf{主控制环必须承担}“对抗无关 + 证书保持 + 鲁棒严格下降”的证明义务（定义~\ref{def:uu-eng-robust-descent}），
并能输出可核验的步进证书 $\mathsf{Cert}_{\mathrm{step}}(L)$（定理~\ref{thm:uu-eng-iff} 的(2)）。
而纯粹依赖数据更新的神经网络决策器属于一种\textbf{必然滞后}的策略同伦类，
在 UU 首次分岔点无法承担上述证明义务。

本节给出：定义—（公理降级）—命题—推论—备注 的闭合链条。
\end{remark}

\subsection{形式逻辑：滞后学习器与“必然滞后同伦类”}
\label{sec:nn-lag-class}

\begin{definition}[滞后学习器（仅靠数据更新的参数化决策器）]
\label{def:lagging-learner}
沿用拟合器/学习器定义~\ref{def:rt-open-adversarial-learner}。
称一类在线决策器为\textbf{滞后学习器}，若其在回合索引 $t$ 时刻的动作输出为
\[
  a_t = f_{\phi_t}(h_t),
\]
且其参数只允许在“数据暴露之后”更新：
\[
  \phi_{t+1}=\mathsf{Update}(\phi_t;h_t).
\]
其中 $f_{\phi}$ 可以由神经网络参数化（但不限定为神经网络），关键结构是“先观测后更新”，
从而对任何首次出现的新机制只能在其可观测差异暴露后追赶。
\end{definition}

\begin{definition}[必然滞后同伦类 $\lbrack\mathsf{Lag}\rbrack$（结构等价类）]
\label{def:lag-homotopy-class}
把所有滞后学习器（定义~\ref{def:lagging-learner}）视为同一“必然滞后同伦类” $\lbrack\mathsf{Lag}\rbrack$，
其同伦不变量取为：
\textbf{在任意不可区分历史前缀内，动作输出无法提前分化}（见定理~\ref{thm:lag-indistinguishable-prefix}）。
直观上，$\lbrack\mathsf{Lag}\rbrack$ 的“同伦”等价含义是：允许以常数/多项式开销互相模拟其更新与输出，
但不改变该“不可区分前缀必然同动作”的结构性不变量。
\end{definition}

\subsection{公理降级：不可区分前缀内必然同动作（归纳证书）}
\label{sec:nn-lag-induction}

\begin{theorem}[公理降级：不可区分前缀内，滞后学习器的动作无法分化（归纳证书）]
\label{thm:lag-indistinguishable-prefix}
设两世界 $W^{(1)},W^{(2)}$ 在 $t=0,1,\dots,n$ 的历史前缀完全相同。
则对任意滞后学习器（定义~\ref{def:lagging-learner}），都有
\[
  \phi_t^{(1)}=\phi_t^{(2)}\ \wedge\ a_t^{(1)}=a_t^{(2)}\qquad (t=0,1,\dots,n).
\]
\end{theorem}

\begin{Certificate}[数学归纳法证书：相同历史前缀 $\Rightarrow$ 相同内部状态与动作]
\label{cert:lag-indistinguishable-prefix}
本定理可直接复用证书~\ref{cert:rt-open-adversarial-lag-lower-bound} 的谓词 $Q(t)$ 与归纳链条。
\end{Certificate}

\begin{proof}[证明（数学归纳法证书；谓词因果链）]
由定义~\ref{def:lagging-learner} 的更新规则与输入一致性，可套用引理~\ref{lem:rt-open-adversarial-lag-lower-bound} 的归纳证明，
从而得到在不可区分前缀内参数与动作无法分化。证毕。
\end{proof}

\subsection{命题：纯神经网络主决策器无法承担 UU 的鲁棒下降证明义务}
\label{sec:nn-main-loop-not-core}

\begin{definition}[主控制环（证明义务承担者）]
\label{def:main-control-loop-proof-obligation}
在 UU 工程实时口径下（定义~\ref{def:uu-eng-realtime-solve}），
称“主控制环”为一对输出接口
\[
  \bigl(\pi,\mathsf{Cert}_{\mathrm{step}}\bigr),
\]
其中 $\pi:\GZLS\to A$ 输出动作 $a=\pi(L)$，
而 $\mathsf{Cert}_{\mathrm{step}}(L)$ 是可核验的步进证书，用于核验：
对所有 $L'\in\mathrm{Post}(L,a)$ 都满足
\[
  \mathcal{E}_{\mathrm{prot}}(L')=\top,\qquad
  \Phi_\lambda(L')\le \Phi_\lambda(L)-\delta,
\]
即“证书保持 + 鲁棒严格下降”（定义~\ref{def:uu-eng-robust-descent}，定理~\ref{thm:uu-eng-iff} 的(2)）。
\end{definition}

\begin{definition}[纯神经网络主决策器（不带外部证书裁决）]
\label{def:pure-nn-main-controller}
称决策器为“纯神经网络主决策器”，若其主输出仅为动作序列（由神经网络参数化的滞后学习器给出），
即 $a_t=f_{\phi_t}(h_t)$，并且其外部系统\textbf{不提供}独立于神经网络的
$\mathsf{Cert}_{\mathrm{step}}(\cdot)$ 验证/裁决回路（因此神经网络本身需承担主控制环的证明义务）。
\end{definition}

\begin{proposition}[纯神经网络主决策器不可能满足 UU 实时可解的鲁棒下降充要条件]
\label{prop:pure-nn-not-uu-robust-descent}
若环境允许 UU 不可区分分岔（定义~\ref{def:rt-open-adversarial-indistinguishable-branch}），
则任何纯神经网络主决策器（定义~\ref{def:pure-nn-main-controller}）都不可能作为一条
$\delta$-鲁棒严格下降策略（定义~\ref{def:uu-eng-robust-descent}）并输出所需步进证书；
因此它不可能成为 UU 工程实时可解（定理~\ref{thm:uu-eng-iff}）的控制核心。
\end{proposition}

\begin{Certificate}[证书：不可区分分岔 $\Rightarrow$ “同动作” $\Rightarrow$ 至少一侧证书破坏]
\label{cert:pure-nn-not-uu-robust-descent}
证书登记谓词链：
\[
\begin{aligned}
  N_1 &: \text{存在不可区分分岔}\Rightarrow \exists W^{(1)},W^{(2)}:\ h_t^{(1)}=h_t^{(2)}\ \wedge\ A^{\mathrm{safe}}_{W^{(1)}}
    (h_t)\cap A^{\mathrm{safe}}_{W^{(2)}}(h_t)=\varnothing,\\
  N_2 &: h_t^{(1)}=h_t^{(2)}\Rightarrow a_t^{(1)}=a_t^{(2)}\ \text{（定理~\ref{thm:lag-indistinguishable-prefix}）},\\
  N_3 &: N_1\wedge N_2\Rightarrow \exists i\in\{1,2\}:\ a_t^{(i)}\notin A^{\mathrm{safe}}_{W^{(i)}}(h_t),\\
  N_4 &: a_t^{(i)}\notin A^{\mathrm{safe}}_{W^{(i)}}(h_t)\Rightarrow \mathcal{E}_{\mathrm{prot}}\bigl(h_{t+1}^{W^{(i)}}
    (a_t)\bigr)=\bot,\\
  N_5 &: N_4\Rightarrow \neg\bigl(\forall L'\in\mathrm{Post}(L,a_t):\ \mathcal{E}_{\mathrm{prot}}(L')=\top\bigr)
    \Rightarrow \text{不满足 $\delta$-鲁棒修复动作定义}.
\end{aligned}
\]
\end{Certificate}

\begin{proof}[证明（谓词因果链）]
由不可区分分岔定义得到 $N_1$。
由于纯神经网络主决策器属于滞后学习器同伦类 $[\mathsf{Lag}]$（定义~\ref{def:lag-homotopy-class}），
在相同历史输入下动作无法分化，得到 $N_2$。
结合 $N_1$ 的安全集合互斥性，得到 $N_3$：同一动作至少在一侧世界不安全。
由一步安全集合的定义得到 $N_4$：不安全 $\Rightarrow$ 下一步证书破坏。
于是存在一个可达后继使 $\mathcal{E}_{\mathrm{prot}}=\bot$，从而无法满足“对所有后继证书保持”的鲁棒条件（$N_5$），
更不可能同时满足“对所有后继严格下降”的强条件。
因此纯神经网络主决策器不可能成为定理~\ref{thm:uu-eng-iff} 所需的鲁棒下降控制核心。证毕。
\end{proof}

\subsection{PureMath 风格：holonomy 非平凡 $\Rightarrow$ 无统一修复动作 $\Rightarrow$ 落入 $\lbrack\mathsf{Lag}\rbrack$}
\label{sec:lag-holonomy-to-unif-repair}

\begin{remark}[目标（把“holonomy$\Rightarrow$必然滞后”写成可引用核心定理）]
本小节把“必然滞后同伦类 $\lbrack\mathsf{Lag}\rbrack$”按 PureMath 风格压成如下可引用链条：
\[
\begin{aligned}
  &\text{回路 holonomy 非平凡}
  \ \Longrightarrow\
  \text{同一语义纤维内存在不可区分内部态对}\\
  &\Longrightarrow\ \text{仅依赖 }o\in O\text{ 的策略无法给出统一鲁棒修复动作}\\
  &\Longrightarrow\ \text{必然滞后（落入 }\lbrack\mathsf{Lag}\rbrack\text{）}.
\end{aligned}
\]
其中关键桥梁是“UU 修复不可兼容证书”（定义~\ref{def:lag-uu-repair-incompatibility}）：
它把“同一语义下内部态不同”提升为“同一语义下修复义务不可兼容”，从而推出\textbf{统一修复动作不存在}。
\end{remark}

\subsubsection{形式逻辑：holonomy、鲁棒修复动作与统一修复动作}
\label{sec:lag-holonomy-defs}

\begin{definition}[观测、纤维与回路基空间（$O$ 或 $|\Gamma|$ 口径）]
\label{def:lag-holonomy-fiber-loop}
令 $X$ 为内部状态空间，$O$ 为回合边界可观测语义空间，并给定观测映射
\[
  p:X\to O,\qquad o=p(x).
\]
对任意 $o\in O$，纤维记为 $F_o:=p^{-1}(o)$。
为书写回路同伦类，可取回路基空间为 $O$（把 $O$ 视为语义基空间），
或取为回合图的几何实现 $|\Gamma|$（见下一章定义~\ref{def:round-homotopy-round-graph}--\ref{def:round-homotopy-path-loop}）；
两种口径只影响“回路集合”的表达，不影响下文“同语义纤维内不可区分态对”的推演结构。
\end{definition}

\begin{definition}[路径提升与 holonomy]
\label{def:lag-holonomy-transport}
假设存在路径提升/传输结构：对任一路径（或回合路径）$\gamma:o\to o'$，存在映射
\[
  T_{\gamma}:F_o\to F_{o'}.
\]
对回路 $\ell:o\to o$，得到回路传输（holonomy 作用）
\[
  T_{\ell}:F_o\to F_o.
\]
\end{definition}

\begin{definition}[holonomy 非平凡]
\label{def:lag-nontrivial-holonomy}
称 holonomy 非平凡，若存在 $o\in O$ 与回路 $\ell:o\to o$ 使
\[
  T_{\ell}\neq \mathrm{id}_{F_o}.
\]
\end{definition}

\begin{definition}[内部后继与一步鲁棒修复动作集合；统一修复动作集合]
\label{def:lag-uniform-repair-action}
给定回合动作集合 $A$ 与内部后继集合算子
\[
  \mathrm{Post}_{X}:X\times A\to 2^{X},\qquad (x,a)\mapsto \mathrm{Post}_{X}(x,a).
\]
并给定法对象编码 $L:X\to\GZLS$ 与证书谓词/缺陷势
$\mathcal{E}_{\mathrm{prot}}:\GZLS\to\{\top,\bot\}$，
$\Phi:\GZLS\to\RR_{\ge 0}$（例如 $\Phi=\mathcal{E}_{\mathrm{CH}}+\lambda\JacGap$，$\lambda>0$）。
固定步长下界 $\delta>0$，定义内部态 $x$ 处的\textbf{一步鲁棒修复动作集合}为
\[
  A^{\mathrm{rep}}_{\delta}(x)
  :=
  \Bigl\{\,a\in A\ \Bigm|\
  \forall x'\in\mathrm{Post}_{X}(x,a),\
  \mathcal{E}_{\mathrm{prot}}\!\bigl(L(x')\bigr)=\top\
  \wedge\
  \Phi\!\bigl(L(x')\bigr)\le \Phi\!\bigl(L(x)\bigr)-\delta
  \Bigr\}.
\]
定义同一回合语义 $o$ 下的\textbf{统一修复动作集合}为
\[
  A^{\mathrm{unif}}_{\delta}(o)
  :=
  \bigcap_{x\in F_o} A^{\mathrm{rep}}_{\delta}(x).
\]
若 $A^{\mathrm{unif}}_{\delta}(o)=\varnothing$，则表示：仅凭回合语义 $o$ 无法选出对所有可能内部态都同时有效的一步鲁棒修复动作。
\end{definition}

\begin{definition}[仅依赖回合语义的策略类 $\Pi_{O}$]
\label{def:lag-policy-only-o}
称策略 $\pi$“仅依赖回合语义”，若存在 $\bar\pi:O\to A$ 使
\[
  \pi=\bar\pi\circ p.
\]
记该类为
\[
  \Pi_{O}:=\{\ \pi\mid \exists\bar\pi,\ \pi=\bar\pi\circ p\ \}.
\]
\end{definition}

\subsubsection{引理链：holonomy 非平凡 $\Rightarrow$ 不可区分内部态对；$\Pi_O$ 无法区分}
\label{sec:lag-holonomy-lemmas}

\begin{lemma}[holonomy 非平凡给出不可区分内部态对]
\label{lem:lag-holonomy-gives-pair}
若 holonomy 非平凡（定义~\ref{def:lag-nontrivial-holonomy}），
则存在 $o\in O$、回路 $\ell:o\to o$ 与 $x\in F_o$，令 $y:=T_{\ell}(x)$，满足
\[
  p(x)=p(y)=o,\qquad y\neq x.
\]
\end{lemma}

\begin{Certificate}[证书：holonomy 见证给出同语义不同内部态]
\label{cert:lag-holonomy-gives-pair}
证书登记谓词链：
\[
\begin{aligned}
  H_1 &: \exists(o,\ell):\ T_{\ell}\neq \mathrm{id}_{F_o},\\
  H_2 &: H_1\Rightarrow \exists x\in F_o:\ T_{\ell}(x)\neq x,\\
  H_3 &: \text{令 }y:=T_{\ell}(x)\Rightarrow p(x)=p(y)=o\ \wedge\ y\neq x.
\end{aligned}
\]
\end{Certificate}

\begin{proof}[证明（谓词因果链）]
由 holonomy 非平凡得 $H_1$。非恒等映射必存在一点 $x$ 使 $T_{\ell}(x)\neq x$，得 $H_2$。
令 $y:=T_{\ell}(x)$，由于 $T_{\ell}:F_o\to F_o$，故 $p(y)=o$，并且 $y\neq x$，得 $H_3$。证毕。
\end{proof}

\begin{lemma}[仅依赖 $o$ 的策略无法区分该对]
\label{lem:lag-pio-cannot-distinguish}
对任意 $\pi\in\Pi_{O}$（定义~\ref{def:lag-policy-only-o}），若 $p(x)=p(y)=o$，则
\[
  \pi(x)=\pi(y)=\bar\pi(o).
\]
\end{lemma}

\begin{Certificate}[证书：因子分解 $\pi=\bar\pi\circ p$ $\Rightarrow$ 同语义同动作]
\label{cert:lag-pio-cannot-distinguish}
证书登记谓词链：
\[
\begin{aligned}
  P_1 &: \pi=\bar\pi\circ p,\\
  P_2 &: p(x)=p(y)=o,\\
  P_3 &: P_1\wedge P_2\Rightarrow \pi(x)=\bar\pi(p(x))=\bar\pi(o)=\bar\pi(p(y))=\pi(y).
\end{aligned}
\]
\end{Certificate}

\begin{proof}[证明（谓词因果链）]
由 $\pi=\bar\pi\circ p$ 与 $p(x)=p(y)$ 直接代入即可得到同动作结论。证毕。
\end{proof}

\subsubsection{命题与定理：holonomy + UU 不可兼容证书 $\Rightarrow$ 无统一修复动作 $\Rightarrow$ 落入 $\lbrack\mathsf{Lag}\rbrack$}
\label{sec:lag-holonomy-theorem}

\begin{definition}[UU 修复不可兼容证书（强 UU 口径）]
\label{def:lag-uu-repair-incompatibility}
给定 $o\in O$ 与 $\delta>0$。
若存在 $x,y\in F_o$ 使得
\[
  A^{\mathrm{rep}}_{\delta}(x)\cap A^{\mathrm{rep}}_{\delta}(y)=\varnothing,
\]
则称存在“UU 修复不可兼容障碍”（强 UU），并称 $(o,\delta,x,y)$ 为一份 UU 不可兼容证书见证。
该条件显然推出 $A^{\mathrm{unif}}_{\delta}(o)=\varnothing$（定义~\ref{def:lag-uniform-repair-action}）。
\end{definition}

\begin{proposition}[holonomy 非平凡 + UU 不可兼容证书 $\Rightarrow$ 无统一修复动作]
\label{prop:lag-holonomy-implies-no-unif-repair}
若 holonomy 非平凡，并且存在 UU 修复不可兼容证书（定义~\ref{def:lag-uu-repair-incompatibility}），
则存在某个 $o\in O$ 与 $\delta>0$ 使
\[
  A^{\mathrm{unif}}_{\delta}(o)=\varnothing.
\]
\end{proposition}

\begin{Certificate}[证书：两项交为空 $\Rightarrow$ 全局交为空]
\label{cert:lag-holonomy-implies-no-unif-repair}
证书登记谓词链：
\[
\begin{aligned}
  U_1 &: \exists x,y\in F_o:\ A^{\mathrm{rep}}_{\delta}(x)\cap A^{\mathrm{rep}}_{\delta}(y)=\varnothing,\\
  U_2 &: A^{\mathrm{unif}}_{\delta}(o)=\bigcap_{z\in F_o}A^{\mathrm{rep}}_{\delta}(z),\\
  U_3 &: U_1\wedge U_2\Rightarrow A^{\mathrm{unif}}_{\delta}(o)=\varnothing.
\end{aligned}
\]
\end{Certificate}

\begin{proof}[证明（谓词因果链）]
由定义~\ref{def:lag-uu-repair-incompatibility} 得到某对 $(x,y)$ 使 $U_1$ 成立。
由定义~\ref{def:lag-uniform-repair-action} 得 $U_2$。
由于全局交包含两项交，若两项交为空，则全局交必为空，得到 $U_3$。证毕。
\end{proof}

\begin{theorem}[核心定理：holonomy 非平凡 $\Rightarrow$ $\Pi_O$ 落入 $\lbrack\mathsf{Lag}\rbrack$（在 UU 不可兼容证书下）]
\label{thm:lag-holonomy-to-lag-class}
若存在 holonomy 非平凡见证，且相应存在 UU 修复不可兼容证书（定义~\ref{def:lag-uu-repair-incompatibility}），
则存在 $(o,\delta)$ 使 $A^{\mathrm{unif}}_{\delta}(o)=\varnothing$。
从而对任意仅依赖回合语义的策略 $\pi\in\Pi_O$，
其在第一次到达该语义 $o$ 时无法给出“统一鲁棒修复动作”，
即无法满足定义~\ref{def:uu-eng-robust-descent} 所要求的“对所有后继证书保持 + 严格下降”证明义务。
因此该策略类在该口径下属于必然滞后同伦类 $\lbrack\mathsf{Lag}\rbrack$（定义~\ref{def:lag-homotopy-class}）。
\end{theorem}

\begin{Certificate}[证书：同语义同动作 + 无统一修复动作 $\Rightarrow$ 至少一侧修复义务失败]
\label{cert:lag-holonomy-to-lag-class}
证书登记谓词链：
\[
\begin{aligned}
  T_1 &: A^{\mathrm{unif}}_{\delta}(o)=\varnothing,\\
  T_2 &: \pi\in\Pi_O\Rightarrow \exists\bar\pi:\ \pi=\bar\pi\circ p,\\
  T_3 &: p(x)=p(y)=o\Rightarrow \pi(x)=\pi(y)=\bar\pi(o)=:a\ \text{（引理~\ref{lem:lag-pio-cannot-distinguish}）},\\
  T_4 &: T_1\Rightarrow \exists x\in F_o:\ a\notin A^{\mathrm{rep}}_{\delta}(x),\\
  T_5 &: a\notin A^{\mathrm{rep}}_{\delta}(x)\Rightarrow \exists x'\in\mathrm{Post}_{X}(x,a):\
  \mathcal{E}_{\mathrm{prot}}(L(x'))=\bot\ \text{或}\ \Phi(L(x'))>\Phi(L(x))-\delta.
\end{aligned}
\]
\end{Certificate}

\begin{proof}[证明（谓词因果链）]
由命题~\ref{prop:lag-holonomy-implies-no-unif-repair} 得到某个 $(o,\delta)$ 使 $A^{\mathrm{unif}}_{\delta}(o)=\varnothing$，即
  $T_1$。
对任意 $\pi\in\Pi_O$，得 $T_2$。取任意满足 $p(x)=o$ 的内部态（尤其可取引理~\ref{lem:lag-holonomy-gives-pair} 给出的同语义对），
由引理~\ref{lem:lag-pio-cannot-distinguish} 得同语义同动作 $a=\bar\pi(o)$，即 $T_3$。
由于 $A^{\mathrm{unif}}_{\delta}(o)$ 是对所有 $x\in F_o$ 的交，$T_1$ 意味着不存在一个动作能对所有 $x$ 同时属于 $A^{\mathrm{rep}}_{\delta}(x)$；
特别地，对 $a$ 必存在某个 $x$ 使 $a\notin A^{\mathrm{rep}}_{\delta}(x)$，得 $T_4$。
展开 $A^{\mathrm{rep}}_{\delta}(x)$ 的否定即得 $T_5$：
存在某个可达后继使证书保持或严格下降失败。
因此 $\pi$ 无法承担“对抗无关 + 证书保持 + 鲁棒严格下降”的证明义务，落入必然滞后同伦类口径。证毕。
\end{proof}

\begin{corollary}[推论：神经网络主策略是 $\Pi_O$ 的典型实例，因此在该口径下落入 $\lbrack\mathsf{Lag}\rbrack$]
\label{cor:nn-main-in-lag-via-holonomy}
若神经网络主策略在回合边界仅依赖语义摘要（即其决策可因子分解到 $O$ 上：$\pi=\bar\pi\circ p$），
则 $\pi\in\Pi_O$。
因此在定理~\ref{thm:lag-holonomy-to-lag-class} 的前提下，该神经网络主策略落入 $\lbrack\mathsf{Lag}\rbrack$，
不能作为 UU 实时可解（定理~\ref{thm:uu-eng-iff}）的证明义务承担者。
\end{corollary}

\begin{Certificate}[证书：因子分解到 $O$ 上 $\Rightarrow$ 属于 $\Pi_O$]
\label{cert:nn-main-in-lag-via-holonomy}
证书登记谓词链：
\[
\begin{aligned}
  L_1 &: \pi(x)=f_{\theta}\bigl(p(x)\bigr)\Rightarrow \pi=\bar\pi\circ p,\\
  L_2 &: L_1\Rightarrow \pi\in\Pi_O,\\
  L_3 &: \text{定理~\ref{thm:lag-holonomy-to-lag-class}}\Rightarrow \Pi_O\subseteq \lbrack\mathsf{Lag}\rbrack\
    \text{（在其前提下）}.
\end{aligned}
\]
\end{Certificate}

\begin{proof}[证明（谓词因果链）]
把 $f_{\theta}$ 视为 $\bar\pi$ 即得 $L_1\Rightarrow L_2$。
再由定理~\ref{thm:lag-holonomy-to-lag-class} 得 $L_3$，推论成立。证毕。
\end{proof}

\subsection{推论与备注：神经网络的合理位置（KU 子模块/候选生成器）}
\label{sec:nn-demotion-usage}

\begin{corollary}[架构推论：必须将神经网络从“证明义务承担者”位置撤下]
\label{cor:demote-nn-from-main-loop}
若目标是 UU 的工程实时可解（定理~\ref{thm:uu-eng-iff}），
则主控制环必须以证书化方式实现“证书保持 + 鲁棒严格下降”（定义~\ref{def:main-control-loop-proof-obligation}）。
因此神经网络不能作为“主控制环的证明义务承担者”（命题~\ref{prop:pure-nn-not-uu-robust-descent}），
而应被撤出主控制环，仅以如下两类角色出现：
\begin{enumerate}
  \item \textbf{KU 子模块（估计/近似器）}：用于估计连续量或生成 $\mathcal{E}_{\mathrm{CH}}$ 的候选近似，但最终由可核验证书（上界/下界/抽象封套）兜底；
  \item \textbf{候选动作生成器（建议器）}：输出候选 $a_{\mathrm{cand}}$，但是否执行由
  $\mathsf{Cert}_{\mathrm{step}}(L)$ 或验证器 $\mathsf{Ver}$（定义~\ref{def:magic-wand-certified-solution}）裁决；
  未通过则拒绝并继续在同伦修复动作族中搜索/扩层。
\end{enumerate}
\end{corollary}

\begin{Certificate}[证书：主环证明义务 $\Rightarrow$ NN 只能作为候选/估计器（fail-closed）]
\label{cert:demote-nn-from-main-loop}
证书登记谓词链：
\[
\begin{aligned}
  D_1 &: \text{定理~\ref{thm:uu-eng-iff}}\Rightarrow \text{主控制环必须满足定义~\ref{def:main-control-loop-proof-obligation}},\\
  D_2 &: \text{命题~\ref{prop:pure-nn-not-uu-robust-descent}}\Rightarrow \text{纯 NN 不能承担该证明义务},\\
  D_3 &: D_2\Rightarrow \text{NN 必须退出主环，仅生成候选/估计},\\
  D_4 &: \text{候选仅在证书/验证器通过时执行（定义~\ref{def:magic-wand-certified-solution}；fail-closed）}.
\end{aligned}
\]
\end{Certificate}

\begin{proof}[证明（谓词因果链）]
由定理~\ref{thm:uu-eng-iff}，UU 工程实时可解要求主控制环能在运行期给出可核验的步进证书，并满足
定义~\ref{def:main-control-loop-proof-obligation} 所述的“证书保持 + 鲁棒严格下降”等证明义务，得 $D_1$。
命题~\ref{prop:pure-nn-not-uu-robust-descent} 说明：把纯神经网络作为主决策器时，无法在开放 UU 口径下承担该义务，得 $D_2$。
因此若仍要保持主环可核验性，只能将 NN 的作用降级为“候选/估计”：NN 输出候选动作或连续量近似，
但是否执行必须由确定性的证书检查/验证器裁决并 fail-closed（定义~\ref{def:magic-wand-certified-solution}），得到 $D_3$--$D_4$。
推论成立。证毕。
\end{proof}

\begin{remark}[与“连续统分层/语义收缩重建”的一致性]
定理~\ref{thm:layering-complete-iff-trivial-holonomy} 给出：要消除 UU 根源需引入离散层指标进行语义抬升，
把暗处 holonomy 显式化为层切换。
一旦该分层可用，主控制环可在离散层空间上以“修复动作/边缘算子搜索 + 证书裁决 + $\Phi$ 严格下降”构造修复塔（定义~\ref{def:homotopy-repair-tower-phi}），
此时神经网络即便存在，也只应作为候选生成器而非证明义务承担者。
\end{remark}

\section{神经网络的审计降级：候选生成器与确定性裁决器的语义函子包装}
\label{sec:audited-wrapper-nn}

\begin{remark}[目标（把“NN 降级为候选生成器、确定性裁决器承担正确性”写成可推演结构）]
上一章（推论~\ref{cor:demote-nn-from-main-loop}）给出：在开放 UU 口径下，主控制环必须承担
“对抗无关 + 证书保持 + 鲁棒严格下降”的证明义务，因此纯神经网络主决策器不能作为证明义务承担者。

本章给出一套与既有两段式工程模板一致的形式化：把不确定生成（NN/LLM）\textbf{降为候选生成器}，
把“是否可接受/是否鲁棒/是否规范形”交给\textbf{确定性裁决器（审计器）}。
其严格含义是：神经网络不再承担主控制环证明义务，而只作为
（i）语义函子构造的选择器/加速器，或（ii）监督学习的信息压缩算子；
并且在“失败时闭合”（fail-closed）口径下，不会把系统重新拉回“必然滞后同伦类”的主逻辑。

为对接您在两份工作笔记中使用的接口，本章将它们抽象为同一结构：
\[
  F^{\mathrm{aud}}_{\theta}=V\circ G_{\theta}:\ \mathsf{WP}\to \mathsf{NF}\sqcup\{\bot\},
\]
其中 $G_\theta$ 生成候选，$V$ 确定性裁决并输出（或拒绝为 $\bot$）。
（接口参考：\code{/mnt/data/gframework_semantic_functor_field_gauge_field_normalization_interface_formal_system.tex}，
\code{/mnt/data/gframework_semantic_functor_regex_normalization_validation_robustness_formal_system.tex}。）
\end{remark}

\subsection{形式逻辑：语义函子、规范形合同与语义规范场}
\label{sec:audited-wrapper-formalism}

\begin{definition}[词包表达层与规范形集合]
\label{def:audited-wrapper-wp-nf}
令 $\mathsf{WP}$ 表示某层“词包表达/语义摘要”的集合（可理解为运行期可见的离散语义对象），
令 $\mathsf{NF}\subseteq\mathsf{WP}$ 表示规范形集合（normal forms），并约定失败符号 $\bot\notin\mathsf{WP}$。
对任意 $y$，谓词 $\mathsf{NF}(y)$ 记为“$y\in\mathsf{NF}$”。
\end{definition}

\begin{definition}[语义函子（部分规范化函数）]
\label{def:audited-wrapper-semantic-functor}
称映射
\[
  F:\mathsf{WP}\to \mathsf{NF}\sqcup\{\bot\}
\]
为一个语义函子（部分规范化函数），若对任意输入 $x\in\mathsf{WP}$：
\begin{itemize}
  \item $F(x)=y\in\mathsf{NF}$ 表示“规范化成功并返回规范形代表”；
  \item $F(x)=\bot$ 表示“拒绝/失败（fail-closed）”，即不承诺给出规范形。
\end{itemize}
该合同把“可接受/不可接受”与“返回的规范形”统一为可判定的输出类型。
\end{definition}

\begin{definition}[语义规范场（语义函子场）]
\label{def:audited-wrapper-semantic-functor-field}
给定索引集 $I$，一族语义函子
\[
  \mathcal{F}=\{F_i\}_{i\in I},\qquad F_i:\mathsf{WP}\to \mathsf{NF}_i\sqcup\{\bot\},
\]
称为语义规范场（语义函子场）。
其中 $\mathsf{NF}_i$ 表示第 $i$ 个规范化坐标/规则族对应的规范形集合。
\end{definition}

\subsection{关键构造：严格审计包装（Audited Functor Wrapper）}
\label{sec:audited-wrapper-construction}

\begin{definition}[候选生成器与确定性裁决器]
\label{def:audited-wrapper-generator-auditor}
给定任意候选生成器（可为神经网络或 LLM）
\[
  G_\theta:\mathsf{WP}\to \mathsf{Cand},
\]
其中 $\mathsf{Cand}$ 为候选域（可包含“候选输出 + 证书草稿 + 测试样例草稿”等结构）。
给定一个确定性裁决器（审计器）
\[
  V:\mathsf{Cand}\to \mathsf{NF}\sqcup\{\bot\},
\]
其实现可包含：规则白名单检查、编译/语义一致性检查、复杂度阈值检查、回归测试与变形测试等。
\end{definition}

\begin{definition}[严格审计包装]
\label{def:audited-wrapper-definition}
定义审计后的语义函子为复合：
\[
  F^{\mathrm{aud}}_\theta:=V\circ G_\theta:\ \mathsf{WP}\to \mathsf{NF}\sqcup\{\bot\}.
\]
直观上：$G_\theta$ 只负责“提案”，$V$ 负责“裁决”；对外可引用的语义函子是 $F^{\mathrm{aud}}_\theta$，
而 $G_\theta$ 不再是语义定义的一部分。
\end{definition}

\begin{definition}[审计器健全性（Soundness）]
\label{def:audited-wrapper-soundness}
给定一个确定性谓词族 $\mathsf{Ok}:\mathsf{Cand}\times(\mathsf{NF}\sqcup\{\bot\})\to\{\top,\bot\}$，
表示“审计器对候选的所有确定性检查均通过且返回值合同正确”。
称审计器 $V$ \textbf{健全}，若对任意 $c\in\mathsf{Cand}$：
\[
  V(c)=y\neq\bot\ \Longrightarrow\ \mathsf{Ok}(c,y)=\top\ \wedge\ \mathsf{NF}(y).
\]
若 $V(c)=\bot$，则表示 fail-closed 拒绝，不对规范形作出承诺。
\end{definition}

\begin{lemma}[引理：健全审计器的单步合同]
\label{lem:audited-wrapper-step-contract}
若审计器 $V$ 健全（定义~\ref{def:audited-wrapper-soundness}），则对任意 $c\in\mathsf{Cand}$：
\[
  V(c)=y\neq\bot\ \Longrightarrow\ \mathsf{Ok}(c,y)=\top\ \wedge\ \mathsf{NF}(y).
\]
\end{lemma}

\begin{Certificate}[证书：单步审计通过 $\Rightarrow$ 合同成立]
\label{cert:audited-wrapper-step-contract}
证书登记谓词链：
\[
\begin{aligned}
  S_1 &: V(c)=y\neq\bot,\\
  S_2 &: \text{$V$ 健全（定义~\ref{def:audited-wrapper-soundness}）},\\
  S_3 &: S_1\wedge S_2\Rightarrow \mathsf{Ok}(c,y)=\top\ \wedge\ \mathsf{NF}(y).
\end{aligned}
\]
\end{Certificate}

\begin{proof}[证明（谓词因果链）]
由 $S_1$ 与 $S_2$，直接套用定义~\ref{def:audited-wrapper-soundness} 的蕴含式即可得到 $S_3$。证毕。
\end{proof}

\begin{theorem}[正确性转移：在健全审计下，语义函子的正确性与 NN 泛化无关]
\label{thm:audited-wrapper-soundness}
若审计器 $V$ 健全（定义~\ref{def:audited-wrapper-soundness}），则对任意参数 $\theta$，
审计包装 $F^{\mathrm{aud}}_\theta$（定义~\ref{def:audited-wrapper-definition}）是可核验的语义函子：
对任意输入 $x\in\mathsf{WP}$，
\[
  F^{\mathrm{aud}}_\theta(x)=y\neq\bot\ \Longrightarrow\ \mathsf{NF}(y)\ \wedge\ \mathsf{Ok}\bigl(G_\theta(x),y\bigr)
    =\top.
\]
因此该“对外正确性”（输出若非 $\bot$ 必为规范形且满足确定性检查）完全由 $V$ 的健全性决定，
与 $G_\theta$ 的泛化能力无关。
\end{theorem}

\begin{Certificate}[证书：审计规则集 + 审计日志/回归测试/变形测试]
\label{cert:audited-wrapper-soundness}
对任意输入 $x$，若 $F^{\mathrm{aud}}_\theta(x)=y\neq\bot$，证书可取为审计器 $V$ 的裁决日志
（包含：命中的规则集合、通过的测试集合、复杂度上界证明等）。
形式上登记谓词链：
\[
\begin{aligned}
  A_1 &: c:=G_\theta(x)\in\mathsf{Cand},\\
  A_2 &: y:=V(c)\neq\bot,\\
  A_3 &: \text{$V$ 健全}\Rightarrow \mathsf{Ok}(c,y)=\top\ \wedge\ \mathsf{NF}(y).
\end{aligned}
\]
\end{Certificate}

\begin{proof}[证明（谓词因果链）]
取任意 $x$，令 $c=G_\theta(x)$，则 $A_1$ 成立。
若 $V(c)=\bot$，结论对“$y\neq\bot$”前提为真命题无需证明。
若 $V(c)=y\neq\bot$，由引理~\ref{lem:audited-wrapper-step-contract} 得 $\mathsf{Ok}(c,y)=\top$ 且 $y\in\mathsf{NF}$，
即 $A_3$，从而定理结论成立。证毕。
\end{proof}

\subsection{公理（降级为定理）：逐步审计健全性 $\Rightarrow$ 全程 fail-closed 合规（归纳证书）}
\label{sec:audited-wrapper-induction}

\begin{theorem}[公理降级：逐步审计通过 $\Rightarrow$ 任意轮次均保持合同（数学归纳法证书）]
\label{thm:audited-wrapper-inductive-safety}
令 $x_0\in\mathsf{WP}$，并定义迭代过程：对每步 $t\in\NN$，
先生成候选 $c_t:=G_\theta(x_t)$，再审计裁决 $y_t:=V(c_t)$，并按 fail-closed 规则更新
\[
  x_{t+1}:=
  \begin{cases}
    y_t, & y_t\neq\bot,\\
    x_t, & y_t=\bot.
  \end{cases}
\]
若 $V$ 健全（定义~\ref{def:audited-wrapper-soundness}），则对所有 $t$：
\[
  y_t\neq\bot\ \Longrightarrow\ x_{t+1}\in\mathsf{NF}\ \wedge\ \mathsf{Ok}(c_t,x_{t+1})=\top.
\]
即：任意轮次上，只要审计通过并执行更新，则执行的输出必满足规范形合同与确定性检查。
\end{theorem}

\begin{Certificate}[证书：基步 + 归纳步（回合归纳证书）]
\label{cert:audited-wrapper-inductive-safety}
定义谓词 $P(t)$ 为：
\[
  P(t):\ \Bigl(y_t\neq\bot\Rightarrow x_{t+1}\in\mathsf{NF}\wedge \mathsf{Ok}(c_t,x_{t+1})=\top\Bigr).
\]
证书为
\[
  \pi_{\mathrm{ind}}=\bigl(P(0),\ \forall t\,(P(t)\Rightarrow P(t+1))\bigr),
\]
其中 $P(t)\Rightarrow P(t+1)$ 的证明仅使用 $V$ 健全性与更新规则的局部展开。
\end{Certificate}

\begin{proof}[证明（数学归纳法证书；谓词因果链）]
对任意 $t$，若 $y_t=\bot$，则 $P(t)$ 前件为假，命题成立。
若 $y_t\neq\bot$，由 $V$ 健全性得 $y_t\in\mathsf{NF}$ 且 $\mathsf{Ok}(c_t,y_t)=\top$。
由更新规则得 $x_{t+1}=y_t$，于是 $x_{t+1}\in\mathsf{NF}$ 且 $\mathsf{Ok}(c_t,x_{t+1})=\top$，即 $P(t)$ 成立。
因此对所有 $t$ 都有 $P(t)$。
将“对所有 $t$”写为归纳链时，可直接使用数学归纳法模板（例如定理~\ref{thm:wand-multistep-verification-induction} 的同型口径），
或使用本章证书~\ref{cert:audited-wrapper-inductive-safety} 所登记的基步与归纳步完成降级。证毕。
\end{proof}

\subsection{命题与推论：NN 的两种合法位置（选择器/证书草稿）}
\label{sec:audited-wrapper-uses}

\begin{proposition}[命题：审计包装使 NN 不再承担“证明义务”，因此不构成主环滞后风险源]
\label{prop:audited-wrapper-removes-proof-burden}
在审计包装结构（定义~\ref{def:audited-wrapper-definition}）下，
主控制环若以 $V$ 的裁决为执行条件（fail-closed），则 NN/LLM 生成器 $G_\theta$ 不承担
“证书保持 + 鲁棒严格下降”（定义~\ref{def:main-control-loop-proof-obligation}）的证明义务：
所有被执行的输出都以 $V$ 的确定性检查与日志证书为依据（定理~\ref{thm:audited-wrapper-soundness} 与~\ref{thm:audited-wrapper-inductive-safety}）。
因此主环的可核验性不再依赖 $G_\theta$ 的泛化能力，只依赖审计器 $V$ 的健全性。
\end{proposition}

\begin{Certificate}[证书：证明义务从 $G_\theta$ 迁移到 $V$ 的裁决日志]
\label{cert:audited-wrapper-removes-proof-burden}
证书登记谓词链：
\[
\begin{aligned}
  B_1 &: \text{执行条件：仅当 }y=V(G_\theta(x))\neq\bot\text{ 时才执行更新},\\
  B_2 &: y\neq\bot\Rightarrow \mathsf{NF}(y)\wedge \mathsf{Ok}(G_\theta(x),y)=\top\ \text{（定理~\ref{thm:audited-wrapper-soundness}）},\\
  B_3 &: \text{证书载体：$V$ 的裁决日志/回归与变形测试报告（证书~\ref{cert:audited-wrapper-soundness}）}.
\end{aligned}
\]
\end{Certificate}

\begin{proof}[证明（谓词因果链）]
由 $B_1$，主环仅在 $V$ 接受时更新；由 $B_2$，一旦更新被执行，则其输出满足确定性检查与规范形合同；
$B_3$ 给出可核验的证书载体。因此“被执行的输出之正确性”由 $V$ 的确定性裁决与证书承担，而非由 $G_\theta$ 承担。
证毕。
\end{proof}

\begin{corollary}[推论：NN 作为语义规范场的选择器（只决定“先试谁”，不定义语义）]
\label{cor:audited-wrapper-nn-selector}
设存在语义规范场 $\mathcal{F}=\{F_i\}_{i\in I}$（定义~\ref{def:audited-wrapper-semantic-functor-field}）。
令 NN 仅输出一个候选索引序列 $G_\theta(x)\subseteq I$（按优先级排序），
并由确定性裁决器按顺序尝试
\[
  F(x):=\mathsf{FirstAccept}\bigl(\{F_i(x)\}_{i\in G_\theta(x)}\bigr),
\]
其中 $\mathsf{FirstAccept}$ 表示“返回第一个非 $\bot$ 的值，否则 $\bot$”。
则 NN 只影响搜索顺序（效率），不引入新语义规则；因此不会重新进入主环证明义务核心。
\end{corollary}

\begin{Certificate}[证书：索引选择 + $\mathsf{FirstAccept}$ 的确定性 $\Rightarrow$ 只影响顺序]
\label{cert:audited-wrapper-nn-selector}
证书登记谓词链：
\[
\begin{aligned}
  S_1 &: G_\theta(x)\subseteq I\ \text{仅输出索引序列（不输出语义规则）},\\
  S_2 &: \forall i\in I,\ F_i\ \text{为已固定的语义函子/审计函子（定义~\ref{def:audited-wrapper-semantic-functor-field}）},\\
  S_3 &: \mathsf{FirstAccept}\ \text{为确定性选择器，仅依赖“是否为 }\bot\text{”的裁决结果},\\
  S_4 &: S_1\wedge S_2\wedge S_3\Rightarrow \text{NN 仅改变尝试顺序（效率），不改变可核验合同/语义定义}.
\end{aligned}
\]
\end{Certificate}

\begin{proof}[证明（谓词因果链）]
由 $S_1$，NN 的输出只指定要尝试的索引序列，不携带可执行的语义规则。
由 $S_2$，每个 $F_i$ 的语义与 fail-closed 行为由语义规范场固定（包含“失败显式为 $\bot$”的合同口径）。
由 $S_3$，$\mathsf{FirstAccept}$ 的返回值完全由“按序尝试时第一个非 $\bot$ 的结果”决定，是确定性的组合算子。
因此 NN 只能通过改变尝试的\emph{顺序}影响效率；一旦某个输出被接受，其可核验性仍由对应的确定性裁决（或其内含审计器）承担。
推论成立。证毕。
\end{proof}

\begin{corollary}[推论：NN 生成“证明草稿/重写轨迹”，$V$ 只做验证（证明义务仍在 $V$）]
\label{cor:audited-wrapper-nn-proof-sketch}
令 NN 输出候选对 $(y,\pi^\sharp)=G_\theta(x)$，其中 $y$ 为候选规范形，$\pi^\sharp$ 为“重写步骤序列/推导轨迹/样例集/复杂度上界”等证据草稿。
令审计器 $V$ 验证：
（i）每一步重写来自固定规则库；（ii）终态 $y\in\mathsf{NF}$；（iii）回归与变形测试通过。
则 NN 仅提供证据草稿（加速），而裁决与可核验性仍由 $V$ 的确定性验证承担。
\end{corollary}

\begin{Certificate}[证书：草稿可疑，裁决可核验（证书载体为 $V$ 的日志）]
\label{cert:audited-wrapper-nn-proof-sketch}
证书登记谓词链：
\[
\begin{aligned}
  P_1 &: (y,\pi^\sharp)=G_\theta(x)\ \text{仅为候选与证据草稿},\\
  P_2 &: V\ \text{确定性验证}\Rightarrow \text{每步来自固定规则库，且终态 }y\in\mathsf{NF},\\
  P_3 &: P_2\Rightarrow \text{回归/变形测试通过并生成可复核日志},\\
  P_4 &: P_3\Rightarrow \text{证明义务与可追溯性由 }V\text{ 的验证与日志承担（fail-closed）}.
\end{aligned}
\]
\end{Certificate}

\begin{proof}[证明（谓词因果链）]
由 $P_1$，NN 的输出在逻辑上只是“候选 + 草稿”，并不自动成为可执行/可接受结果。
若 $V$ 拒绝，则按 fail-closed 规则输出 $\bot$ 或回退，不产生无证书执行；若 $V$ 接受，则由 $P_2$--$P_3$，
接受的依据完全来自 $V$ 的确定性验证与可复核日志。
因此 NN 只提供加速信息，而裁决与可追溯性仍由 $V$ 承担（$P_4$），推论成立。证毕。
\end{proof}

\subsection{形式逻辑：监督压缩算子与 fail-closed 接口}
\label{sec:audited-wrapper-compression}

\begin{definition}[监督压缩算子与压缩审计器]
\label{def:audited-wrapper-supervised-compression}
设原始观测/输入属于集合 $\mathcal{U}$，监督学习得到压缩算子
\[
  C_\theta:\mathcal{U}\to Z,
\]
其中 $Z$ 为较小的码空间（可为离散码或可量化的连续码）。
给定确定性压缩审计器
\[
  V_{\mathrm{comp}}:\ \mathcal{U}\times Z\to\{\mathrm{ok},\mathrm{fail}\},
\]
用于判定压缩误差/语义一致性是否满足阈值。
若 $V_{\mathrm{comp}}(u,C_\theta(u))=\mathrm{fail}$，则按 fail-closed 原则禁止将该压缩输出用于主环决策，并回退到可证书路径。
\end{definition}

\begin{corollary}[推论：监督压缩只影响效率，不影响正确性（fail-closed）]
\label{cor:audited-wrapper-compression-safety}
若主控制环满足“压缩审计失败则不使用压缩结果”的 fail-closed 规则，
则压缩模块 $C_\theta$ 至多导致回退与额外计算成本，而不会成为“无证书错误动作”的来源；
因此它不会重新引入必然滞后同伦类作为主逻辑风险源。
\end{corollary}

\begin{Certificate}[证书：fail-closed $\Rightarrow$ 压缩模块不成为无证书错误来源]
\label{cert:audited-wrapper-compression-safety}
证书登记谓词链：
\[
\begin{aligned}
  C_1 &: V_{\mathrm{comp}}(u,C_\theta(u))=\mathrm{fail}\Rightarrow \text{禁止使用压缩结果并回退到可证书路径（定义~\ref{def:audited-wrapper-supervised-compression}）},\\
  C_2 &: V_{\mathrm{comp}}(u,C_\theta(u))=\mathrm{ok}\Rightarrow \text{压缩误差/语义一致性满足阈值（由 }V_{\mathrm{comp}}\text{
    的确定性验证给出）},\\
  C_3 &: C_1\wedge C_2\Rightarrow \text{压缩模块仅影响效率与回退次数，不产生无证书执行}.
\end{aligned}
\]
\end{Certificate}

\begin{proof}[证明（谓词因果链）]
当 $V_{\mathrm{comp}}$ 失败时，由 $C_1$ 直接 fail-closed 回退，因此压缩输出不会进入主环决策；
当 $V_{\mathrm{comp}}$ 通过时，由 $C_2$ 得压缩结果满足约束阈值并可由确定性验证复核。
两种情形合并即得 $C_3$：压缩模块只能带来效率变化（更多回退/更多计算），而不会成为无证书错误动作来源。
推论成立。证毕。
\end{proof}

\begin{remark}[对齐：与上一章“撤出主环”的一致性]
本章形式化地给出：\textbf{否定的不是 NN 本身}，而是“让 NN 直接承担主环证明义务”的架构。
一旦将 NN 降级为候选生成器/选择器/证据草稿器，并由确定性审计器 $V$ 进行 fail-closed 裁决，
则系统对外可引用的正确性来自 $V$ 的健全性与证书日志，
从而与推论~\ref{cor:demote-nn-from-main-loop} 的结论完全一致。
\end{remark}

\section{回合同伦类刻画开放/封闭：holonomy、未知-未知同伦类与端点决定性}
\label{sec:round-homotopy-open-closed}

\begin{remark}[目标（把“回合同伦类/未知-未知/开放-封闭”压成可判定结构）]
本章把“回合的同伦类”“未知-未知”“开放/封闭”的词义压缩为可判定的结构数据，并给出一个\textbf{当且仅当}的定理链：
\begin{enumerate}
  \item \textbf{封闭系统} $\Longleftrightarrow$ 回合路径的作用只依赖端点（回合同伦类不起作用）；
  \item \textbf{开放系统} $\Longleftrightarrow$ 存在某个回合回路的非平凡作用（出现非零 holonomy，产生“未知-未知同伦类”）。
\end{enumerate}
该刻画把“开放/封闭”的本质差异严格落到“回合回路的同伦类是否产生非平凡作用”上。
\end{remark}

\subsection{形式逻辑：回合图与回合同伦类（$\pi_1$）}
\label{sec:round-homotopy-round-graph}

\begin{definition}[回合边界可观测语义状态集与回合图]
\label{def:round-homotopy-round-graph}
固定一个回合粒度（例如每 $\tau$ 个物理时刻合成一回合；粒度只影响离散化口径，不影响本章的同伦结构结论）。
令 $O$ 为回合边界的可观测语义状态集合（例如局面摘要、工程状态摘要、实验状态摘要、对话语义摘要等）。
在 $O$ 上定义一个有向多重图（回合图）$\Gamma$：
\begin{itemize}
  \item 顶点：$o\in O$；
  \item 有向边：$o\xrightarrow{\alpha}o'$ 表示“一回合内存在某个可执行动作/外部响应组合（记作 $\alpha$）使可观测语义从 $o$ 变为 $o'$”。
\end{itemize}
将 $\Gamma$ 视为 $1$-维 CW 复形，并记其几何实现为 $|\Gamma|$。
\end{definition}

\begin{definition}[回合路径、回合回路与回合同伦类]
\label{def:round-homotopy-path-loop}
一条回合路径 $\gamma$ 是边标签的有限串 $\gamma=\alpha_1\cdots\alpha_n$，其端点由相邻边可拼接性确定。
当路径起点与终点同为某个 $o\in O$ 时，称其为基于 $o$ 的回合回路 $\ell$。
把回路看作 $|\Gamma|$ 中的基点回路，其同伦类记为
\[
  [\ell]\in \pi_1(|\Gamma|,o).
\]
我们把 $[\ell]$ 称为“回合回路的同伦类”（亦称“回合的同伦类”）。
\end{definition}

\subsection{形式逻辑：内部法状态、纤维与路径提升（transport）}
\label{sec:round-homotopy-fiber-transport}

\begin{definition}[内部法状态与观测投影]
\label{def:round-homotopy-internal-state}
令 $X$ 为系统的内部法状态空间（包含隐变量、规则细节、对抗者内部状态、材料微观态等一切决定下一步的“真状态”）。
给定观测/语义压缩投影
\[
  p:X\to O,\qquad x\mapsto o=p(x).
\]
对任意 $o\in O$，其纤维记为
\[
  F_o:=p^{-1}(o),
\]
表示“在回合边界上看起来一样，但内部可能不同”的内部状态集合。
\end{definition}

\begin{definition}[路径提升结构（回合边的纤维传输）]
\label{def:round-homotopy-path-lift}
给定回合图 $\Gamma$（定义~\ref{def:round-homotopy-round-graph}）与投影 $p:X\to O$（定义~\ref{def:round-homotopy-internal-state}）。
称系统具有路径提升结构，若对每条回合边 $o\xrightarrow{\alpha}o'$，存在一个确定的纤维传输映射
\[
  T_{\alpha}:F_o\to F_{o'},
\]
表示：当回合边界语义为 $o$ 且内部态为 $x\in F_o$ 时，执行该回合操作 $\alpha$ 后内部态更新到 $x':=T_{\alpha}(x)\in F_{o'}$。
\end{definition}

\begin{theorem}[公理降级：回合路径的复合传输是良定义的（按路径长度归纳）]
\label{thm:round-homotopy-transport-composition}
在定义~\ref{def:round-homotopy-path-lift} 的路径提升结构下，对任意回合路径 $\gamma=\alpha_1\cdots\alpha_n$（从 $o$ 到 $o'$），
可递归定义复合传输映射
\[
  T_{\gamma}:=T_{\alpha_n}\circ\cdots\circ T_{\alpha_1}:F_o\to F_{o'}.
\]
并且对任意可拼接路径 $\gamma_1,\gamma_2$，有
\[
  T_{\gamma_2\cdot\gamma_1}=T_{\gamma_2}\circ T_{\gamma_1},
  \qquad
  T_{\mathrm{id}_o}=\mathrm{id}_{F_o}.
\]
\end{theorem}

\begin{Certificate}[数学归纳法证书：按路径长度 $n$ 定义 $T_\gamma$ 并验证复合律]
\label{cert:round-homotopy-transport-composition}
令谓词 $P(n)$ 表示：“对任意长度为 $n$ 的回合路径 $\gamma$，$T_\gamma$ 可由递归复合良定义，且满足复合律与单位元性质。”
证书化要求为
\[
  P(0)\ \text{成立},\qquad (\forall n\in\NN)\ \bigl(P(n)\Rightarrow P(n+1)\bigr).
\]
\end{Certificate}

\begin{proof}[证明（数学归纳法证书；谓词因果链）]
由证书~\ref{cert:round-homotopy-transport-composition}：
\begin{itemize}
  \item 基步：$n=0$ 时路径为空，定义 $T_{\mathrm{id}_o}:=\mathrm{id}_{F_o}$，单位元性质成立；
  \item 归纳步：设 $P(n)$ 成立。对长度 $n+1$ 的路径 $\gamma=\widetilde\gamma\cdot\alpha_{n+1}$，
  由归纳假设 $T_{\widetilde\gamma}$ 已良定义，且 $T_{\gamma}:=T_{\alpha_{n+1}}\circ T_{\widetilde\gamma}$ 给出递归定义。
  复合律由函数复合结合律直接推出。
\end{itemize}
故 $P(n)$ 对任意 $n$ 成立，定理成立。证毕。
\end{proof}

\subsection{形式逻辑：holonomy 作用与“未知-未知同伦类”}
\label{sec:round-homotopy-holonomy-uu}

\begin{definition}[回合回路的 holonomy（纤维自同态）]
\label{def:round-homotopy-holonomy}
对任意基于 $o\in O$ 的回合回路 $\ell$，其复合传输给出纤维自同态
\[
  \mathsf{Hol}_{\ell}:=T_{\ell}:F_o\to F_o.
\]
若对任意两条同伦回路 $\ell\simeq \ell'$（在 $|\Gamma|$ 中同伦），都有 $T_{\ell}=T_{\ell'}$，
则 holonomy 仅依赖同伦类，可写为 $\mathsf{Hol}_{[\ell]}$。
\end{definition}

\begin{definition}[未知-未知同伦类（非平凡 holonomy）]
\label{def:round-homotopy-uu-class}
若存在某个 $o\in O$ 与某个回路同伦类 $[\ell]\in\pi_1(|\Gamma|,o)$ 使得
\[
  \mathsf{Hol}_{[\ell]}\neq \mathrm{id}_{F_o},
\]
则称该同伦类 $[\ell]$ 为一个\textbf{未知-未知同伦类}，并称系统存在“未知-未知”（UU）结构性来源。
\end{definition}

\subsection{定理：封闭/开放 $\Longleftrightarrow$ 回合同伦类作用平凡/非平凡}
\label{sec:round-homotopy-open-closed-theorem}

\begin{definition}[端点决定性（回合路径同伦类不起作用）]
\label{def:round-homotopy-endpoint-determined}
称系统满足端点决定性，若对任意 $o,o'\in O$ 与任意两条从 $o$ 到 $o'$ 的回合路径 $\gamma_1,\gamma_2$，都有
\[
  T_{\gamma_1}=T_{\gamma_2}:F_o\to F_{o'}.
\]
直观上：内部态的更新只依赖“端点语义”，不依赖走过的回合路径（同伦类不起作用）。
\end{definition}

\begin{assumption}[可逆传输（回合逆边与逆映射）]
\label{asm:round-homotopy-invertible-transport}
对每条回合边 $o\xrightarrow{\alpha}o'$，存在一条形式逆边 $o'\xrightarrow{\alpha^{-1}}o$，
并满足
\[
  T_{\alpha^{-1}}=(T_{\alpha})^{-1}.
\]
对任意回合路径 $\gamma=\alpha_1\cdots\alpha_n$，定义反向路径
\[
  \gamma^{-1}:=\alpha_n^{-1}\cdots\alpha_1^{-1},
\]
并要求其传输满足
\[
  T_{\gamma^{-1}}=(T_{\gamma})^{-1}.
\]
该假设用于把“回路 holonomy 平凡”严格等价地转换为“同端点路径传输相同”（端点决定性）。
\end{assumption}

\begin{lemma}[端点决定性 $\Longleftrightarrow$ 全部回路 holonomy 平凡]
\label{lem:round-homotopy-endpoint-iff-trivial-holonomy}
在假设~\ref{asm:round-homotopy-invertible-transport} 下，下列命题等价：
\begin{enumerate}
  \item 系统满足端点决定性；
  \item 对任意 $o\in O$ 与任意回合回路 $\ell$（基于 $o$），都有 $T_{\ell}=\mathrm{id}_{F_o}$。
\end{enumerate}
\end{lemma}

\begin{Certificate}[证书：端点决定性与回路平凡性的互推]
\label{cert:round-homotopy-endpoint-iff-trivial-holonomy}
证书登记谓词链：
\[
\begin{aligned}
  E_1 &: \text{端点决定性}\Rightarrow T_{\ell}=T_{\mathrm{id}_o}=\mathrm{id}_{F_o}\ \text{（取同端点路径）},\\
  E_2 &: (\forall \ell,\ T_{\ell}=\mathrm{id})\Rightarrow \forall \gamma_1,\gamma_2:o\to o',\ T_{\gamma_2^{-1}
    \cdot\gamma_1}=\mathrm{id}_{F_o},\\
  E_3 &: T_{\gamma_2^{-1}\cdot\gamma_1}=\mathrm{id}\Rightarrow (T_{\gamma_2})^{-1}\circ T_{\gamma_1}=\mathrm{id}
    \Rightarrow T_{\gamma_1}=T_{\gamma_2}.
\end{aligned}
\]
其中 $\gamma^{-1}$ 与 $T_{\gamma^{-1}}=(T_{\gamma})^{-1}$ 由假设~\ref{asm:round-homotopy-invertible-transport} 给出。
\end{Certificate}

\begin{proof}[证明（谓词因果链）]
（1$\Rightarrow$2）若端点决定性成立，则对任意回路 $\ell:o\to o$，
因其端点与空路径 $\mathrm{id}_o$ 相同，故 $T_{\ell}=T_{\mathrm{id}_o}=\mathrm{id}_{F_o}$，即 $E_1$。

（2$\Rightarrow$1）反之假设任意回路作用平凡。
取任意两条同端点路径 $\gamma_1,\gamma_2:o\to o'$。
由假设~\ref{asm:round-homotopy-invertible-transport} 可形成回路
\[
  \ell:=\gamma_2^{-1}\cdot\gamma_1:o\to o.
\]
由回路平凡性得 $T_{\ell}=\mathrm{id}_{F_o}$。
结合定理~\ref{thm:round-homotopy-transport-composition} 的复合律与 $T_{\gamma_2^{-1}}=(T_{\gamma_2})^{-1}$，
得到
\[
  \mathrm{id}=T_{\ell}=T_{\gamma_2^{-1}\cdot\gamma_1}=T_{\gamma_2^{-1}}\circ T_{\gamma_1}=(T_{\gamma_2})^{-1}\circ
    T_{\gamma_1},
\]
从而 $T_{\gamma_1}=T_{\gamma_2}$，即端点决定性成立。
证毕。
\end{proof}

\begin{definition}[封闭/开放系统（以回合同伦类作用为核心）]
\label{def:round-homotopy-open-closed}
在本章口径下：
\begin{itemize}
  \item 称系统为\textbf{封闭系统}，若其端点决定性成立（定义~\ref{def:round-homotopy-endpoint-determined}），即回合同伦类不起作用；
  \item 称系统为\textbf{开放系统}，若存在未知-未知同伦类（定义~\ref{def:round-homotopy-uu-class}），即存在非平凡 holonomy。
\end{itemize}
\end{definition}

\begin{theorem}[封闭系统 $\Longleftrightarrow$ 回合同伦类不起作用；开放系统 $\Longleftrightarrow$ 非平凡 holonomy]
\label{thm:round-homotopy-open-closed}
在定义~\ref{def:round-homotopy-open-closed} 的语义口径下：
\begin{enumerate}
  \item 系统封闭 $\Longleftrightarrow$ 回合路径的作用只依赖端点（同伦类不起作用）；
  \item 系统开放 $\Longleftrightarrow$ 存在某个回路同伦类 $[\ell]$ 使 $\mathsf{Hol}_{[\ell]}\neq \mathrm{id}$，即存在未知-未知同伦类。
\end{enumerate}
\end{theorem}

\begin{Certificate}[证书：封闭/开放的同伦类刻画]
\label{cert:round-homotopy-open-closed}
证书登记谓词链：
\[
\begin{aligned}
  C_1 &: \text{封闭系统}\Longleftrightarrow \text{端点决定性（定义~\ref{def:round-homotopy-open-closed}）},\\
  C_2 &: \text{端点决定性}\Longleftrightarrow \forall \ell,\ T_{\ell}=\mathrm{id}\ \text{（引理~\ref{lem:round-homotopy-endpoint-iff-trivial-holonomy}）},\\
  C_3 &: \text{开放系统}\Longleftrightarrow \exists[\ell],\ \mathsf{Hol}_{[\ell]}\neq \mathrm{id}\ \text{（定义~\ref{def:round-homotopy-open-closed} 与~\ref{def:round-homotopy-uu-class}）}.
\end{aligned}
\]
\end{Certificate}

\begin{proof}[证明（谓词因果链）]
由定义~\ref{def:round-homotopy-open-closed} 得到 $C_1$ 与 $C_3$。
由引理~\ref{lem:round-homotopy-endpoint-iff-trivial-holonomy} 得到 $C_2$。
合并 $C_1$--$C_3$ 即得两条当且仅当结论。证毕。
\end{proof}

\begin{remark}[与曲率/联络口径的对齐（可选后续）]
若把 $\{T_{\gamma}\}$ 进一步提升为联络下的平行运输，则“存在非平凡 holonomy”可由“曲率非零”给出；
而在 \GFramework{} 语境中，$\JacGap$ 可作为可计算的结构缺陷证书来追踪“同伦补齐的缺失”，从而把“回路 holonomy 的非平凡性”
与“缺陷势沿回路积累/不可消去”联系起来，形成更贴近 PureMath 的推论链。
\end{remark}

\section{回合细化稳定性：严格回合同伦类、KU/UU 障碍与开放/封闭博弈的充要刻画}
\label{sec:round-refinement-stability-ku-uu}

\begin{remark}[定位（把两句自然语言命题压成可证明的判定结构）]
本章给出一套“最小但足够严格”的形式化，使下列两句可证明为定理链：
\begin{enumerate}
  \item \textbf{开放博弈系统}：不存在严格的回合同伦类，只能通过回合细化做近似逼近；修复障碍必然体现为未知-未知（UU）；
  \item \textbf{封闭博弈系统}：存在严格的回合同伦类（细化不改变同伦型）；任何障碍只能是已知-未知（KU）。
\end{enumerate}
核心做法是：把“回合”视为对连续交互过程的分割/采样，把“严格回合同伦类存在”定义为回合细化下同伦型的稳定性，
把 KU/UU 定义为“真实后继集合是否闭合于预设模型族”。
\end{remark}

\subsection{形式逻辑：回合分割与回合图（$\Gamma_{\mathcal{P}}$）}
\label{sec:round-refinement-def-round}

\begin{definition}[回合分割（partition）]
\label{def:round-partition}
令连续交互发生在 $t\in\RR_{\ge 0}$ 上，回合边界可观测语义为 $o(t)\in O$。
一个回合分割 $\mathcal{P}$ 是一列严格递增的边界时刻
\[
  \mathcal{P}:\ 0=\tau_0<\tau_1<\tau_2<\cdots,
\]
并把第 $k$ 回合定义为区间 $[\tau_k,\tau_{k+1})$。
回合边界语义状态记为
\[
  o^{(\mathcal{P})}_k:=o(\tau_k).
\]
\end{definition}

\begin{definition}[回合图（round graph / $1$-复形）]
\label{def:round-graph-by-partition}
给定分割 $\mathcal{P}$，定义回合图 $\Gamma_{\mathcal{P}}$ 为一个有向多重图（亦可视作 $1$-维 CW 复形）：
\begin{itemize}
  \item 顶点集：$V(\Gamma_{\mathcal{P}})\subseteq O$（所有可能出现的回合边界语义状态）；
  \item 有向边：$o\xrightarrow{\alpha}o'$ 表示“存在一回合内的可归因控制摘要/对抗摘要（记作 $\alpha$），使语义从 $o$ 迁移到 $o'$”。
\end{itemize}
记其几何实现为 $|\Gamma_{\mathcal{P}}|$。
\end{definition}

\subsection{形式逻辑：细化与严格回合同伦类（同伦型稳定性）}
\label{sec:round-refinement-strict-homotopy}

\begin{definition}[细化关系与压缩映射]
\label{def:round-refinement}
称 $\mathcal{P}'$ 是 $\mathcal{P}$ 的细化，记 $\mathcal{P}'\succ \mathcal{P}$，若 $\mathcal{P}$ 的每个边界时刻都出现在 $\mathcal{P}'$
  中（即把某些回合进一步切细）。
当 $\Gamma_{\mathcal{P}'}$ 是由 $\Gamma_{\mathcal{P}}$ 的边细分（subdivision）得到时，存在自然的“压缩映射”
\[
  q_{\mathcal{P}'\to \mathcal{P}}:\ |\Gamma_{\mathcal{P}'}|\to |\Gamma_{\mathcal{P}}|
\]
把细分链压缩为粗回合的一条边。
\end{definition}

\begin{definition}[严格回合同伦类存在（同伦型稳定）]
\label{def:strict-round-homotopy-class-exists}
称系统存在\textbf{严格的回合同伦类}，若存在某个分割 $\mathcal{P}_0$ 使得对所有足够细的细化 $\mathcal{P}\succ \mathcal{P}_0$：
\[
  |\Gamma_{\mathcal{P}}|\ \simeq\ |\Gamma_{\mathcal{P}_0}|,
\]
并且这些同伦等价与压缩映射 $q_{\mathcal{P}\to\mathcal{P}_0}$ 相容（细化不会“发现新的同伦型”）。
若不存在这样的 $\mathcal{P}_0$，则称\textbf{不存在严格回合同伦类}：回合结构只能作为近似逼近而不稳定。
\end{definition}

\begin{theorem}[公理降级：$1$-复形的有限细分不改变同伦型（按细分步数归纳）]
\label{thm:round-subdivision-homotopy-invariant}
设 $\Gamma'$ 由 $1$-复形 $\Gamma$ 经过有限次边细分得到，则
\[
  |\Gamma'|\ \simeq\ |\Gamma|.
\]
\end{theorem}

\begin{Certificate}[数学归纳法证书：按细分次数 $n$ 的同伦不变性]
\label{cert:round-subdivision-homotopy-invariant}
令谓词 $P(n)$ 表示：“对任意 $1$-复形 $\Gamma$，对其做 $n$ 次边细分得到 $\Gamma^{(n)}$，则 $|\Gamma^{(n)}|\simeq |\Gamma|$。”
证书化要求为
\[
  P(0)\ \text{成立},\qquad (\forall n\in\NN)\ \bigl(P(n)\Rightarrow P(n+1)\bigr).
\]
\end{Certificate}

\begin{proof}[证明（数学归纳法证书；谓词因果链）]
由证书~\ref{cert:round-subdivision-homotopy-invariant}：
\begin{itemize}
  \item 基步：$n=0$ 时 $\Gamma^{(0)}=\Gamma$，同伦等价为恒等；
  \item 归纳步：假设 $P(n)$ 成立。一次边细分等价于在一条边上插入一个度为 $2$ 的新顶点，
  其几何实现与原边同胚，因此 $|\Gamma^{(n+1)}|\simeq |\Gamma^{(n)}|$。
  合并归纳假设 $|\Gamma^{(n)}|\simeq |\Gamma|$ 得 $|\Gamma^{(n+1)}|\simeq |\Gamma|$。
\end{itemize}
故 $P(n)$ 对任意 $n$ 成立，从而有限次边细分不改变同伦型。证毕。
\end{proof}

\subsection{形式逻辑：KU/UU（后继集合闭合性）与修复缺口}
\label{sec:round-refinement-ku-uu}

\begin{definition}[预设后继族与真实后继集合]
\label{def:round-post-model}
给定分割 $\mathcal{P}$。对回合边界语义 $o\in O$ 与回合动作摘要 $a$，令
\[
  \widehat{\mathrm{Post}}^{(\mathcal{P})}(o,a)\subseteq O
\]
表示事先给定的“允许后继集合族”（模型/规则语言给出的可能性集合），而
\[
  \mathrm{Post}^{(\mathcal{P})}(o,a)\subseteq O
\]
表示真实系统在该回合口径下的实际后继集合。
\end{definition}

\begin{definition}[KU 与 UU（闭合/不闭合）]
\label{def:round-ku-uu}
在分割 $\mathcal{P}$ 下：
\begin{itemize}
  \item 若对所有 $(o,a)$ 都有 $\mathrm{Post}^{(\mathcal{P})}(o,a)\subseteq \widehat{\mathrm{Post}}^{(\mathcal{P})}(o,a)$，
  则称不确定性为\textbf{已知-未知（KU）}：可能性集合已知，但落点未知；
  \item 若存在某个 $(o,a)$ 使 $\mathrm{Post}^{(\mathcal{P})}(o,a)\not\subseteq \widehat{\mathrm{Post}}^{(\mathcal{P})}(o,a)$，
  则称出现\textbf{未知-未知（UU）}：连可能性集合本身都必须扩张。
\end{itemize}
\end{definition}

\begin{definition}[修复障碍（obstruction）]
\label{def:round-obstruction}
定义缺口（修复障碍）为
\[
  \mathsf{Ob}^{(\mathcal{P})}(o,a):=\mathrm{Post}^{(\mathcal{P})}(o,a)\setminus \widehat{\mathrm{Post}}^{(\mathcal{P})}
    (o,a).
\]
则
\[
  \mathsf{Ob}^{(\mathcal{P})}(o,a)\neq\varnothing \quad\Longleftrightarrow\quad \text{在该口径下出现 UU。}
\]
\end{definition}

\subsection{封闭博弈系统：严格回合同伦类存在，障碍必为 KU}
\label{sec:round-refinement-closed}

\begin{definition}[封闭博弈系统]
\label{def:closed-game-system-round}
称系统为封闭博弈系统，若存在某个分割 $\mathcal{P}_0$ 使得：
\begin{enumerate}
  \item （回合内无外源语义注入）在每个回合区间 $[\tau_k,\tau_{k+1})$ 内，不发生改变回合边界语义/规则集合的外源信息注入；
  \item （模型后继闭合）存在一套固定的预设后继族 $\widehat{\mathrm{Post}}^{(\mathcal{P}_0)}$，使对所有 $(o,a)$ 都有
  \[
    \mathrm{Post}^{(\mathcal{P}_0)}(o,a)\subseteq \widehat{\mathrm{Post}}^{(\mathcal{P}_0)}(o,a).
  \]
\end{enumerate}
\end{definition}

\begin{theorem}[封闭 $\Rightarrow$ 严格回合同伦类存在]
\label{thm:closed-implies-strict-round-homotopy}
若系统为封闭博弈系统（定义~\ref{def:closed-game-system-round}），则严格回合同伦类存在（定义~\ref{def:strict-round-homotopy-class-exists}）。
\end{theorem}

\begin{Certificate}[证书：封闭 $\Rightarrow$ 细化仅为边细分 $\Rightarrow$ 同伦型稳定]
\label{cert:closed-implies-strict-round-homotopy}
证书登记谓词链：
\[
\begin{aligned}
  S_1 &: \text{封闭回合内无外源语义注入}\Rightarrow \forall \mathcal{P}\succ\mathcal{P}_0,\ \Gamma_{\mathcal{P}}\ \text{为}\
    \Gamma_{\mathcal{P}_0}\ \text{的有限边细分},\\
  S_2 &: S_1\wedge \text{定理~\ref{thm:round-subdivision-homotopy-invariant}}\Rightarrow |\Gamma_{\mathcal{P}}|\simeq
    |\Gamma_{\mathcal{P}_0}|,\\
  S_3 &: \Rightarrow \exists\mathcal{P}_0\ \text{使细化后同伦型稳定（定义~\ref{def:strict-round-homotopy-class-exists}）}.
\end{aligned}
\]
\end{Certificate}

\begin{proof}[证明（谓词因果链）]
由定义~\ref{def:closed-game-system-round} 的“回合内无外源语义注入”，对任意细化 $\mathcal{P}\succ\mathcal{P}_0$，
细化仅把同一粗回合内部的演化拆成更多小步，对应在 $1$-复形层面对边做有限细分，得到 $S_1$。
由定理~\ref{thm:round-subdivision-homotopy-invariant} 得 $S_2$，
从而存在稳定分割 $\mathcal{P}_0$ 使同伦型对足够细化保持不变，即 $S_3$。证毕。
\end{proof}

\begin{theorem}[封闭 $\Rightarrow$ 障碍必为 KU（无 UU 缺口）]
\label{thm:closed-implies-ku}
若系统为封闭博弈系统（定义~\ref{def:closed-game-system-round}），则在分割 $\mathcal{P}_0$ 下不出现 UU：
\[
  \forall(o,a),\qquad \mathsf{Ob}^{(\mathcal{P}_0)}(o,a)=\varnothing.
\]
因此任何障碍只能是 KU（定义~\ref{def:round-ku-uu}）。
\end{theorem}

\begin{Certificate}[证书：真实后继闭合 $\Rightarrow$ 缺口为空]
\label{cert:closed-implies-ku}
证书登记谓词链：
\[
\begin{aligned}
  K_1 &: \forall(o,a),\ \mathrm{Post}^{(\mathcal{P}_0)}(o,a)\subseteq \widehat{\mathrm{Post}}^{(\mathcal{P}_0)}(o,a),\\
  K_2 &: K_1\Rightarrow \forall(o,a),\ \mathrm{Post}^{(\mathcal{P}_0)}(o,a)\setminus \widehat{\mathrm{Post}}
    ^{(\mathcal{P}_0)}(o,a)=\varnothing,\\
  K_3 &: \text{由定义~\ref{def:round-obstruction}，}K_2\Rightarrow \forall(o,a),\ \mathsf{Ob}^{(\mathcal{P}_0)}(o,a)
    =\varnothing.
\end{aligned}
\]
\end{Certificate}

\begin{proof}[证明（谓词因果链）]
由封闭博弈系统的“模型后继闭合”条件直接得到 $K_1$。
集合差定义给出 $K_2$，再由定义~\ref{def:round-obstruction} 得 $K_3$，定理成立。证毕。
\end{proof}

\subsection{开放博弈系统：严格回合同伦类不存在，障碍必为 UU}
\label{sec:round-refinement-open}

\begin{definition}[开放博弈系统（不可稳定性）]
\label{def:open-game-system-round}
称系统为开放博弈系统，若满足以下不可稳定性：
\begin{enumerate}
  \item （同伦型不可稳定）对任意分割 $\mathcal{P}$，都存在某个更细的细化 $\mathcal{P}'\succ\mathcal{P}$ 使
  \[
    |\Gamma_{\mathcal{P}'}|\not\simeq |\Gamma_{\mathcal{P}}|
  \]
  或至少在基本回路结构上出现不可消去的新生成（例如 $\pi_1$ 改变）；
  \item （模型不闭合）对任意固定的预设后继族 $\widehat{\mathrm{Post}}^{(\mathcal{P})}$，必存在某个 $(o,a)$ 使
  \[
    \mathsf{Ob}^{(\mathcal{P})}(o,a)\neq\varnothing.
  \]
\end{enumerate}
\end{definition}

\begin{theorem}[开放 $\Rightarrow$ 不存在严格回合同伦类]
\label{thm:open-implies-no-strict-round-homotopy}
若系统为开放博弈系统（定义~\ref{def:open-game-system-round}），则不存在严格回合同伦类（定义~\ref{def:strict-round-homotopy-class-exists} 不成立）。
\end{theorem}

\begin{Certificate}[证书：反证法——若稳定则与“任意分割可破坏”矛盾]
\label{cert:open-implies-no-strict-round-homotopy}
证书登记谓词链：
\[
\begin{aligned}
  O_1 &: \text{开放性（同伦型不可稳定）}\Rightarrow \forall \mathcal{P},\exists \mathcal{P}'\succ\mathcal{P}:\ |\Gamma_{\mathcal{P}'}
    |\not\simeq |\Gamma_{\mathcal{P}}|,\\
  O_2 &: \text{若存在严格回合同伦类，则}\exists\mathcal{P}_0,\forall \mathcal{P}\succ\mathcal{P}_0:\ |\Gamma_{\mathcal{P}}|\simeq
    |\Gamma_{\mathcal{P}_0}|,\\
  O_3 &: O_1\wedge O_2\Rightarrow \text{矛盾}.
\end{aligned}
\]
\end{Certificate}

\begin{proof}[证明（谓词因果链）]
反设严格回合同伦类存在，则有某个 $\mathcal{P}_0$ 使所有足够细化的同伦型都与 $|\Gamma_{\mathcal{P}_0}|$ 等价（$O_2$）。
但开放性要求对任意分割都能找到更细化破坏同伦等价（$O_1$），特别地对 $\mathcal{P}_0$ 也成立，得到矛盾（$O_3$）。
故严格回合同伦类不存在。证毕。
\end{proof}

\begin{theorem}[开放 $\Rightarrow$ 只能近似回合逼近，且修复障碍必为 UU]
\label{thm:open-implies-approx-and-uu}
若系统为开放博弈系统（定义~\ref{def:open-game-system-round}），则：
\begin{enumerate}
  \item 任取分割 $\mathcal{P}$，回合模型 $\Gamma_{\mathcal{P}}$ 只能视为近似；通过继续细化 $\mathcal{P}'\succ\mathcal{P}$ 可以逼近真实交互，
    但不存在能稳定的“严格回合”口径；
  \item 并且对任意固定的预设后继族 $\widehat{\mathrm{Post}}^{(\mathcal{P})}$，必存在 $(o,a)$ 使 $\mathsf{Ob}^{(\mathcal{P})}(o,a)
    \neq\varnothing$，即 UU 障碍必然出现。
\end{enumerate}
\end{theorem}

\begin{Certificate}[证书：不可稳定性 $\Rightarrow$ 仅能逼近；模型不闭合 $\Rightarrow$ 缺口非空]
\label{cert:open-implies-approx-and-uu}
证书登记谓词链：
\[
\begin{aligned}
  A_1 &: \forall\mathcal{P},\exists\mathcal{P}'\succ\mathcal{P}:\ |\Gamma_{\mathcal{P}'}|\not\simeq |\Gamma_{\mathcal{P}}
    |,\\
  A_2 &: A_1\Rightarrow \text{不存在稳定分割}\ \mathcal{P}_0\ \text{，因此只能通过细化序列做近似逼近},\\
  A_3 &: \forall \widehat{\mathrm{Post}}^{(\mathcal{P})},\exists(o,a):\ \mathsf{Ob}^{(\mathcal{P})}(o,a)\neq\varnothing,
    \\
  A_4 &: \text{由定义~\ref{def:round-obstruction}，}A_3\Rightarrow \text{UU 必然出现}.
\end{aligned}
\]
\end{Certificate}

\begin{proof}[证明（谓词因果链）]
由开放性第(1)条得到 $A_1$，从而不存在能使同伦型稳定的分割口径，得到 $A_2$（只能通过细化近似逼近）。
由开放性第(2)条得到 $A_3$，再由定义~\ref{def:round-obstruction} 得 $A_4$（UU 缺口非空）。
合并即得两点结论。证毕。
\end{proof}

\begin{remark}[对接提示（从 $\mathsf{Ob}$ 到 $\JacGap$，再到 $\Phi$ 与修复塔）]
在 \GFramework{} 语境下，$\mathsf{Ob}^{(\mathcal{P})}(o,a)\neq\varnothing$ 的“模型不闭合”可作为 $\JacGap>0$ 的一个最小证书来源：
它不是“参数估计不准”（KU），而是“回合模型本体不闭合”（UU）。
下一章给出严格对接：把 $\mathsf{Ob}$ 编码为可计算的 $\JacGap$，并将其并入合并缺陷势
$\Phi=\mathcal{E}_{\mathrm{CH}}+\lambda\JacGap$，再把“同伦修复塔”写成对 $\Phi$ 的逐层证书化下降构造。
\end{remark}

\section{对接：$\mathsf{Ob}$、$\JacGap$、$\Phi=\mathcal{E}_{\mathrm{CH}}+\lambda\JacGap$ 与同伦修复塔}
\label{sec:ob-jacgap-phi-repair-tower}

\begin{remark}[目标（把回合细化/KU-UU/修复塔统一到同一势能口径）]
本章把上一章的“回合分割 $\mathcal{P}$、后继集合、修复缺口 $\mathsf{Ob}^{(\mathcal{P})}$”完整对接到您偏好的合并缺陷势
\[
  \Phi(L):=\mathcal{E}_{\mathrm{CH}}(L)+\lambda\,\JacGap(L),\qquad \lambda>0,
\]
并把“同伦修复塔”写成一个可核验的逐层构造（证书保持 + 势能下降）。
最后把两条自然语言结论提升为严格定理：封闭系统 $\Rightarrow$ 严格回合同伦类存在且障碍为 KU；开放系统 $\Rightarrow$ 严格回合同伦类不存在且障碍为 UU，并只能用修复塔做逼近。
\end{remark}

\subsection{形式逻辑：把 $\mathsf{Ob}^{(\mathcal{P})}$ 编码为可计算缺陷证书 $\JacGap$}
\label{sec:ob-to-jacgap}

\begin{definition}[缺口规模测度（工程量化口径）]
\label{def:ob-measure-mu}
取一个缺口规模测度（或代价函数）
\[
  \mu:2^{O}\to \RR_{\ge 0},
\]
满足 $\mu(\varnothing)=0$，且当 $S\neq\varnothing$ 时 $\mu(S)>0$。
工程上 $\mu$ 可取为：基数、概率质量、风险权重、修复成本下界等。
\end{definition}

\begin{definition}[分割口径下的 $\JacGap$：UU 缺口的最小结构缺陷证书]
\label{def:jacgap-from-obstruction}
对分割 $\mathcal{P}$ 下的回合模型法对象记为 $L^{(\mathcal{P})}\in\GZLS$，
其接口包含 $\widehat{\mathrm{Post}}^{(\mathcal{P})}$、$\mathcal{E}_{\mathrm{prot}}$ 等结构（见定义~\ref{def:round-post-model}）。
定义该口径下的结构缺陷证书为
\[
  \JacGap^{(\mathcal{P})}\!\bigl(L^{(\mathcal{P})}\bigr)
  :=\sup_{(o,a)}\ \mu\!\bigl(\mathsf{Ob}^{(\mathcal{P})}(o,a)\bigr),
\]
其中 $\mathsf{Ob}^{(\mathcal{P})}(o,a)$ 由定义~\ref{def:round-obstruction} 给出。
\end{definition}

\begin{lemma}[UU 见证 $\Rightarrow$ $\JacGap^{(\mathcal{P})}>0$（最小证书）]
\label{lem:uu-witness-implies-jacgap-positive}
若存在某个 $(o,a)$ 使 $\mathsf{Ob}^{(\mathcal{P})}(o,a)\neq\varnothing$，则
\[
  \JacGap^{(\mathcal{P})}\!\bigl(L^{(\mathcal{P})}\bigr)>0.
\]
\end{lemma}

\begin{Certificate}[证书：非空缺口 $\Rightarrow$ 测度为正 $\Rightarrow$ 上确界为正]
\label{cert:uu-witness-implies-jacgap-positive}
证书登记谓词链：
\[
\begin{aligned}
  J_1 &: \exists(o,a):\ \mathsf{Ob}^{(\mathcal{P})}(o,a)\neq\varnothing,\\
  J_2 &: \text{由定义~\ref{def:ob-measure-mu}，}J_1\Rightarrow \mu(\mathsf{Ob}^{(\mathcal{P})}(o,a))>0,\\
  J_3 &: \text{由定义~\ref{def:jacgap-from-obstruction}，}J_2\Rightarrow \JacGap^{(\mathcal{P})}(L^{(\mathcal{P})})>0.
\end{aligned}
\]
\end{Certificate}

\begin{proof}[证明（谓词因果链）]
由 $J_1$ 得存在非空缺口；
由测度性质得 $J_2$；再由 $\JacGap^{(\mathcal{P})}$ 的上确界定义得 $J_3$。证毕。
\end{proof}

\subsection{形式逻辑：严格回合同伦类与 $\JacGap$ 的稳定归零（充要刻画）}
\label{sec:strict-round-homotopy-iff-jacgap}

\begin{definition}[$\JacGap$ 的稳定归零（细化闭合口径）]
\label{def:jacgap-stable-zero}
称存在“$\JacGap$ 的稳定归零”，若存在某个分割 $\mathcal{P}_0$ 使得对所有足够细的细化 $\mathcal{P}\succ\mathcal{P}_0$：
\[
  \JacGap^{(\mathcal{P})}\!\bigl(L^{(\mathcal{P})}\bigr)=0,
\]
并且细化压缩到 $\mathcal{P}_0$ 的口径下仍保持后继闭合（即压缩后不产生新的 $\mathsf{Ob}^{(\mathcal{P}_0)}$ 见证）。
\end{definition}

\begin{assumption}[缺口--结构对应（UU 缺口等价于“非细分新增”）]
\label{asm:ob-structure-correspondence}
对任意基准分割 $\mathcal{P}_0$ 与任意足够细的细化 $\mathcal{P}\succ\mathcal{P}_0$，细化后的回合图 $\Gamma_{\mathcal{P}}$ 若相对于
  $\Gamma_{\mathcal{P}_0}$ \textbf{不仅仅}是边细分（即必须新增不可由细分解释的顶点/边来表达真实后继），
则必存在 $(o,a)$ 使 $\mathsf{Ob}^{(\mathcal{P})}(o,a)\neq\varnothing$（从而 $\JacGap^{(\mathcal{P})}>0$）。
反向地，若对所有 $(o,a)$ 均有 $\mathsf{Ob}^{(\mathcal{P})}(o,a)=\varnothing$，则细化仅表现为对 $\Gamma_{\mathcal{P}_0}$ 的边细分。
\end{assumption}

\begin{theorem}[充要条件：严格回合同伦类存在 $\Longleftrightarrow$ $\JacGap$ 在细化下稳定归零]
\label{thm:strict-round-homotopy-iff-jacgap-stable-zero}
在假设~\ref{asm:ob-structure-correspondence} 下：
\[
  \boxed{
  \text{严格回合同伦类存在（定义~\ref{def:strict-round-homotopy-class-exists}）}
  \ \Longleftrightarrow\
  \text{$\JacGap$ 稳定归零（定义~\ref{def:jacgap-stable-zero}）}.
  }
\]
\end{theorem}

\begin{Certificate}[证书：稳定归零 $\Rightarrow$ 仅细分 $\Rightarrow$ 同伦型稳定；反向用缺口--结构对应]
\label{cert:strict-round-homotopy-iff-jacgap-stable-zero}
证书登记谓词链：
\[
\begin{aligned}
  H_1 &: \JacGap^{(\mathcal{P})}(L^{(\mathcal{P})})=0\Rightarrow \forall(o,a),\ \mathsf{Ob}^{(\mathcal{P})}(o,a)
    =\varnothing,\\
  H_2 &: \text{由假设~\ref{asm:ob-structure-correspondence}，}H_1\Rightarrow \Gamma_{\mathcal{P}}\ \text{仅为}\
    \Gamma_{\mathcal{P}_0}\ \text{的边细分},\\
  H_3 &: H_2\wedge \text{定理~\ref{thm:round-subdivision-homotopy-invariant}}\Rightarrow |\Gamma_{\mathcal{P}}|\simeq
    |\Gamma_{\mathcal{P}_0}|,\\
  H_4 &: \Rightarrow \text{严格回合同伦类存在},\\
  R_1 &: \text{若严格回合同伦类存在但}\exists\mathcal{P}\succ\mathcal{P}_0:\ \JacGap^{(\mathcal{P})}>0,\\
  R_2 &: \text{由假设~\ref{asm:ob-structure-correspondence}，}R_1\Rightarrow \Gamma_{\mathcal{P}}\ \text{存在非细分新增，从而破坏稳定},\\
  R_3 &: \Rightarrow \text{矛盾}\Rightarrow \JacGap\ \text{稳定归零}.
\end{aligned}
\]
\end{Certificate}

\begin{proof}[证明（谓词因果链）]
（$\Rightarrow$）设 $\JacGap$ 稳定归零。由 $H_1$ 得细化口径下无 UU 缺口；
由假设~\ref{asm:ob-structure-correspondence} 得细化仅为边细分（$H_2$）；
再由定理~\ref{thm:round-subdivision-homotopy-invariant} 得同伦型稳定（$H_3$），从而严格回合同伦类存在（$H_4$）。

（$\Leftarrow$）反向设严格回合同伦类存在。若存在某个足够细的细化使 $\JacGap^{(\mathcal{P})}>0$，由 $R_2$ 得该细化包含“非细分新增”，从而破坏稳定性，与“严格存在”矛盾。
因此必须有 $\JacGap$ 稳定归零。证毕。
\end{proof}

\subsection{形式逻辑：把 KU 与 UU 统一进 $\Phi=\mathcal{E}_{\mathrm{CH}}+\lambda\JacGap$}
\label{sec:phi-ch-jacgap-split}

\begin{definition}[CH 残差（KU 内部修复项）]
\label{def:ch-energy-ku}
令 $\mathcal{E}_{\mathrm{CH}}:\GZLS\to\RR_{\ge 0}$ 为可证书化的“已知约束/已知一致性”残差能量，
其满足：$\mathcal{E}_{\mathrm{CH}}(L)\ge 0$ 且 $\mathcal{E}_{\mathrm{CH}}(L)=0$ 表示“在当前模型语言内已知约束已修复完成”。
该项主要刻画 KU（已知集合内的残差与不确定）。
\end{definition}

\begin{definition}[合并缺陷势：KU+UU 的统一势能]
\label{def:phi-ch-jacgap}
固定 $\lambda>0$。对任意回合模型法对象 $L^{(\mathcal{P})}\in\GZLS$，定义
\[
  \Phi\!\bigl(L^{(\mathcal{P})}\bigr)
  :=\mathcal{E}_{\mathrm{CH}}\!\bigl(L^{(\mathcal{P})}\bigr)
  +\lambda\,\JacGap^{(\mathcal{P})}\!\bigl(L^{(\mathcal{P})}\bigr).
\]
直观分工是：
\[
  \mathcal{E}_{\mathrm{CH}}\ \text{处理 KU（集合内残差）},\qquad
  \JacGap\ \text{处理 UU（集合不闭合缺口）}.
\]
\end{definition}

\subsection{形式逻辑：同伦修复塔（证书保持 + $\Phi$ 逐层下降）}
\label{sec:homotopy-repair-tower-phi}

\begin{definition}[同伦修复动作与修复塔]
\label{def:homotopy-repair-tower-phi}
称 $h$ 为一次同伦修复动作，若它把回合模型法对象 $L$ 变换为 $L':=h(L)$，其语义允许包括：
扩张 $\widehat{\mathrm{Post}}$、细化分割 $\mathcal{P}$、升阶引入更高阶自由度/证书接口等。
称序列
\[
  L_0 \xrightarrow{h_1} L_1 \xrightarrow{h_2}\cdots\xrightarrow{h_n} L_n \xrightarrow{h_{n+1}}\cdots
\]
为一条同伦修复塔，若每一步都满足：
\[
  \mathcal{E}_{\mathrm{prot}}(L_{k+1})=\top,
  \qquad
  \Phi(L_{k+1})\le \Phi(L_k)-\delta_k,\ \ \delta_k>0,
\]
并且每步都附带可核验证书（证明“证书保持 + 势能下降”成立）。
\end{definition}

\begin{theorem}[公理降级：存在逐层下降修复塔 $\Rightarrow$ 有限步内到达 $\Phi=0$（数学归纳法证书）]
\label{thm:repair-tower-terminates}
设存在常数 $\delta>0$ 使对修复塔每一步均有 $\delta_k\ge \delta$，
且 $\Phi$ 在可达集合上满足量化下界 $\Phi(\cdot)\in \delta\NN$。
则从任意初态 $L_0$ 出发，修复塔在至多 $\Phi(L_0)/\delta$ 步内到达某个 $L_T$ 使 $\Phi(L_T)=0$，
并沿程保持 $\mathcal{E}_{\mathrm{prot}}=\top$。
\end{theorem}

\begin{Certificate}[数学归纳法证书：按步数 $n$ 的势能下降与证书保持]
\label{cert:repair-tower-terminates}
令谓词 $P(n)$ 表示：“对任意 $t\le n$，有 $\mathcal{E}_{\mathrm{prot}}(L_t)=\top$ 且 $\Phi(L_t)\le \Phi(L_0)-t\delta$。”
证书化要求为基步与归纳步成立。
\end{Certificate}

\begin{proof}[证明（数学归纳法证书；谓词因果链）]
由证书~\ref{cert:repair-tower-terminates}：
\begin{itemize}
  \item 基步：$n=0$ 时结论显然；
  \item 归纳步：假设 $P(n)$ 成立。由修复塔定义得 $\mathcal{E}_{\mathrm{prot}}(L_{n+1})=\top$，
  且 $\Phi(L_{n+1})\le \Phi(L_n)-\delta\le \Phi(L_0)-(n+1)\delta$，故 $P(n+1)$ 成立。
\end{itemize}
因此 $P(n)$ 对任意 $n$ 成立。由于 $\Phi(\cdot)\in\delta\NN$ 且非负，不可能无限次每步至少下降 $\delta$ 而不归零；
故在至多 $T=\Phi(L_0)/\delta$ 步后必有 $\Phi(L_T)=0$。证毕。
\end{proof}

\subsection{定理：封闭/开放博弈系统的 $\Phi$--$\JacGap$--回合同伦类刻画}
\label{sec:closed-open-phi-jacgap}

\begin{theorem}[封闭博弈系统：$\JacGap$ 稳定为 $0$，严格回合同伦类存在，障碍必为 KU]
\label{thm:closed-phi-jacgap-ku}
若系统为封闭博弈系统（定义~\ref{def:closed-game-system-round}），则存在分割 $\mathcal{P}_0$ 使：
\[
  \forall \mathcal{P}\succ\mathcal{P}_0,\qquad \JacGap^{(\mathcal{P})}\!\bigl(L^{(\mathcal{P})}\bigr)=0.
\]
因此 $\JacGap$ 稳定归零（定义~\ref{def:jacgap-stable-zero}），从而严格回合同伦类存在（定理~\ref{thm:strict-round-homotopy-iff-jacgap-stable-zero}）
  ，并且障碍只能是 KU（定义~\ref{def:round-ku-uu}）。
\end{theorem}

\begin{Certificate}[证书：封闭 $\Rightarrow$ 缺口为空 $\Rightarrow$ $\JacGap=0$]
\label{cert:closed-phi-jacgap-ku}
证书登记谓词链：
\[
\begin{aligned}
  C_1 &: \text{由定理~\ref{thm:closed-implies-ku}，}\forall(o,a),\ \mathsf{Ob}^{(\mathcal{P}_0)}(o,a)=\varnothing,\\
  C_2 &: C_1\Rightarrow \JacGap^{(\mathcal{P}_0)}(L^{(\mathcal{P}_0)})=0\ \text{（定义~\ref{def:jacgap-from-obstruction}
    与~\ref{def:ob-measure-mu}）},\\
  C_3 &: \text{封闭回合内无外源注入}\Rightarrow \forall\mathcal{P}\succ\mathcal{P}_0,\ \mathsf{Ob}^{(\mathcal{P})}
    \equiv\varnothing\Rightarrow \JacGap^{(\mathcal{P})}=0.
\end{aligned}
\]
\end{Certificate}

\begin{proof}[证明（谓词因果链）]
由定理~\ref{thm:closed-implies-ku} 得到 $C_1$，从而由定义~\ref{def:jacgap-from-obstruction} 得到 $C_2$。
封闭系统细化仅为回合内部拆分，不引入新的真实后继缺口，因此对任意细化仍有 $\mathsf{Ob}^{(\mathcal{P})}=\varnothing$，得到 $C_3$。
于是 $\JacGap$ 稳定归零，从而由定理~\ref{thm:strict-round-homotopy-iff-jacgap-stable-zero} 得严格回合同伦类存在，且障碍为 KU。证毕。
\end{proof}

\begin{theorem}[开放博弈系统：$\JacGap$ 无法稳定归零，严格回合同伦类不存在，障碍必为 UU]
\label{thm:open-phi-jacgap-uu}
若系统为开放博弈系统（定义~\ref{def:open-game-system-round}），则对任意分割 $\mathcal{P}$ 与任意固定的预设后继族 $\widehat{\mathrm{Post}}^{(\mathcal{P})}
  $，
必存在 $(o,a)$ 使 $\mathsf{Ob}^{(\mathcal{P})}(o,a)\neq\varnothing$，从而
\[
  \JacGap^{(\mathcal{P})}\!\bigl(L^{(\mathcal{P})}\bigr)>0.
\]
因此 $\JacGap$ 不可能在细化下稳定归零，严格回合同伦类不存在（定理~\ref{thm:strict-round-homotopy-iff-jacgap-stable-zero}），
并且障碍必然体现为 UU。要降低 $\Phi=\mathcal{E}_{\mathrm{CH}}+\lambda\JacGap$，只能通过同伦修复塔把缺口逐层吸收（定义~\ref{def:homotopy-repair-tower-phi}）
  。
\end{theorem}

\begin{Certificate}[证书：开放 $\Rightarrow$ 缺口非空 $\Rightarrow$ $\JacGap>0$]
\label{cert:open-phi-jacgap-uu}
证书登记谓词链：
\[
\begin{aligned}
  O_1 &: \text{由定义~\ref{def:open-game-system-round} 的“模型不闭合”，}\forall\widehat{\mathrm{Post}}^{(\mathcal{P})},\exists(o,a)
    :\ \mathsf{Ob}^{(\mathcal{P})}(o,a)\neq\varnothing,\\
  O_2 &: O_1\Rightarrow \JacGap^{(\mathcal{P})}(L^{(\mathcal{P})})>0\ \text{（引理~\ref{lem:uu-witness-implies-jacgap-positive}）},\\
  O_3 &: \Rightarrow \neg\text{$\JacGap$ 稳定归零}\Rightarrow \neg\text{严格回合同伦类存在}.
\end{aligned}
\]
\end{Certificate}

\begin{proof}[证明（谓词因果链）]
由开放系统定义直接得到 $O_1$。
由引理~\ref{lem:uu-witness-implies-jacgap-positive} 得到 $O_2$。
由“$\JacGap$ 稳定归零 $\Longleftrightarrow$ 严格回合同伦类存在”（定理~\ref{thm:strict-round-homotopy-iff-jacgap-stable-zero}）取逆否得到 $O_3$。
  证毕。
\end{proof}

\section{必要性：仅回合细分无法消除 UU；连续统分层作为相空间切换的离散支撑}
\label{sec:uu-holonomy-layering-necessity}

\begin{remark}[目标（把“回合逼近≠消除 UU”写成必要性定理）]
本章把您提出的结论严格化为一个\textbf{必要性定理}：
在开放博弈（UU）下，单纯把时间/交互切成回合并不断细化回合，
本质上只是对回合基空间做边细分（subdivision）（定义~\ref{def:round-refinement}），
它\textbf{不会}消除 UU 的根源——根源是\textbf{回合回路的非平凡单值性破坏}（holonomy；定义~\ref{def:round-homotopy-holonomy}、\ref{def:round-homotopy-uu-class}）。

要消除该根源，必须对语义/相空间做一次“抬升”（lift）：
引入一个\textbf{离散层指标}把同一可观测语义下的连续纤维分裂为可判别层，
从而把“暗处的 holonomy”显式化为“层切换”（相空间切换的离散支撑）。
\end{remark}

\subsection{形式逻辑：仅回合逼近（subdivision）与“不分层”约束}
\label{sec:only-round-subdivision-no-layering}

\begin{definition}[仅回合逼近（只细分过程，不抬升语义相空间）]
\label{def:only-round-approximation}
在本节口径下，“仅回合逼近”指：允许选择回合分割 $\mathcal{P}$ 并做细化 $\mathcal{P}'\succ\mathcal{P}$
（定义~\ref{def:round-refinement}），从而把回合图 $\Gamma_{\mathcal{P}}$ 细分为 $\Gamma_{\mathcal{P}'}$；
但\textbf{不允许}改变观测/语义投影 $p:X\to O$（定义~\ref{def:round-homotopy-internal-state}），
即不引入任何新的离散层指标（不把 $O$ 抬升为 $O\times I$ 一类扩展语义空间）。
\end{definition}

\subsection{公理降级：回合细分只分段路径，holonomy 不随细分消失}
\label{sec:subdivision-preserve-holonomy}

\begin{theorem}[公理降级：回合细分不改变路径传输（按细分次数归纳）]
\label{thm:subdivision-preserve-transport}
在路径提升结构（定义~\ref{def:round-homotopy-path-lift}）与复合律（定理~\ref{thm:round-homotopy-transport-composition}）下，
设 $\mathcal{P}'\succ\mathcal{P}$ 为回合细化（定义~\ref{def:round-refinement}）。
对任意基于同一端点的回合路径 $\gamma$（在粗口径 $\Gamma_{\mathcal{P}}$ 中），
令 $\gamma'$ 为其在细口径 $\Gamma_{\mathcal{P}'}$ 中的细分表示（压缩映射满足 $q_{\mathcal{P}'\to\mathcal{P}}(\gamma')=\gamma$）。
则整体传输满足
\[
  T_{\gamma'}=T_{\gamma}.
\]
特别地，对任意回合回路 $\ell$，其细分回路 $\ell'$ 满足 $T_{\ell'}=T_{\ell}$，
因此回路 holonomy 的平凡/非平凡性不随回合细分而改变。
\end{theorem}

\begin{Certificate}[数学归纳法证书：按细分次数保持整体传输不变]
\label{cert:subdivision-preserve-transport}
令谓词 $P(n)$ 表示：“对任意回合路径 $\gamma$，若其细分表示 $\gamma^{(n)}$ 由 $\gamma$ 经过 $n$ 次边细分得到，
则 $T_{\gamma^{(n)}}=T_{\gamma}$。”证书化要求为
\[
  P(0)\ \text{成立},\qquad (\forall n\in\NN)\ \bigl(P(n)\Rightarrow P(n+1)\bigr).
\]
\end{Certificate}

\begin{proof}[证明（数学归纳法证书；谓词因果链）]
由证书~\ref{cert:subdivision-preserve-transport}：
\begin{itemize}
  \item 基步：$n=0$ 时 $\gamma^{(0)}=\gamma$，显然 $T_{\gamma^{(0)}}=T_{\gamma}$；
  \item 归纳步：假设 $P(n)$ 成立。一次边细分把某条边（或某段子路径）替换为两段可拼接子路径的复合，
  记为 $\beta=\beta_2\cdot\beta_1$。由定理~\ref{thm:round-homotopy-transport-composition} 的复合律，
  对应传输满足 $T_{\beta}=T_{\beta_2}\circ T_{\beta_1}$。
  因而在 $\gamma^{(n)}$ 中把 $\beta$ 替换为 $\beta_2\cdot\beta_1$ 不改变整体复合结果，得到 $T_{\gamma^{(n+1)}}=T_{\gamma^{(n)}}$。
  再由归纳假设 $T_{\gamma^{(n)}}=T_{\gamma}$，推出 $T_{\gamma^{(n+1)}}=T_{\gamma}$。
\end{itemize}
因此对任意 $n$ 有 $T_{\gamma^{(n)}}=T_{\gamma}$。对回路情形取 $\gamma=\ell$ 即得结论。证毕。
\end{proof}

\begin{theorem}[必要性定理：开放 UU（非平凡 holonomy）下，仅回合细分无法消除 UU 根源]
\label{thm:uu-open-subdivision-not-enough}
在“不分层”的仅回合逼近口径（定义~\ref{def:only-round-approximation}）下，
若系统在回合同伦口径上为开放（存在未知-未知同伦类；定义~\ref{def:round-homotopy-uu-class}），
则任取回合分割 $\mathcal{P}$ 及任意细化 $\mathcal{P}'\succ\mathcal{P}$，系统仍然开放：
存在某个细口径回路同伦类 $[\ell']$ 使其 holonomy 非平凡，
\[
  \mathsf{Hol}_{[\ell']}\neq \mathrm{id}.
\]
因此，\textbf{仅靠回合细分（subdivision）不可能把开放 UU 变成封闭系统}；
若要消除 UU 根源，必须引入语义抬升/分层（下一节）。
\end{theorem}

\begin{Certificate}[证书：非平凡 holonomy 在回合细分下保持]
\label{cert:uu-open-subdivision-not-enough}
证书登记谓词链：
\[
\begin{aligned}
  U_1 &: \exists [\ell]\in\pi_1(|\Gamma_{\mathcal{P}}|,o):\ \mathsf{Hol}_{[\ell]}\neq \mathrm{id}\quad \text{（开放；
    定义~\ref{def:round-homotopy-uu-class}）},\\
  U_2 &: \mathcal{P}'\succ\mathcal{P}\Rightarrow \exists \ell'\ \text{为}\ \ell\ \text{的细分表示，且}\ T_{\ell'}=T_{\ell}\quad
    \text{（定理~\ref{thm:subdivision-preserve-transport}）},\\
  U_3 &: T_{\ell'}=T_{\ell}\ \wedge\ T_{\ell}\neq \mathrm{id}\Rightarrow T_{\ell'}\neq \mathrm{id}\Rightarrow
    \mathsf{Hol}_{[\ell']}\neq \mathrm{id}.
\end{aligned}
\]
\end{Certificate}

\begin{proof}[证明（谓词因果链）]
由开放性（定义~\ref{def:round-homotopy-uu-class}）取到某个回路同伦类 $[\ell]$ 使 $T_{\ell}\neq \mathrm{id}$，即 $U_1$。
对任意细化 $\mathcal{P}'\succ\mathcal{P}$，由定理~\ref{thm:subdivision-preserve-transport} 得到细分回路 $\ell'$ 满足 $T_{\ell'}=T_{\ell}$，
  即 $U_2$。
于是 $T_{\ell'}\neq \mathrm{id}$，得到 $U_3$，从而细口径仍存在非平凡 holonomy。
故仅回合细分无法消除 UU 根源。证毕。
\end{proof}

\subsection{形式逻辑：连续统分层（语义收缩重建）与离散层指标}
\label{sec:layering-discrete-index}

\begin{definition}[分层抬升（离散层指标）]
\label{def:layering-lift}
取一个离散指标集 $I$（可数/可判定即可），并给定一个“层标记”
\[
  \sigma:X\to I.
\]
由此诱导一个扩展观测（语义抬升）
\[
  \widetilde O:=O\times I,\qquad \widetilde p(x):=\bigl(p(x),\sigma(x)\bigr).
\]
对任意 $\widetilde o=(o,i)\in\widetilde O$，其纤维为
\[
  \widetilde F_{\widetilde o}:=\widetilde p^{-1}(\widetilde o)
  =\{x\in X:\ p(x)=o,\ \sigma(x)=i\}.
\]
直观上：$\sigma$ 把原先同一语义 $o$ 下的（可能连续的）纤维 $F_o$ 分裂为离散层，从而把“连续统”收缩/崩溃到可判别层号 $i$。
\end{definition}

\begin{definition}[层完备性（把暗 holonomy 变成显式层切换）]
\label{def:layering-complete}
称分层 $\sigma$ 是\textbf{层完备的}，若对任意 $o\in O$、任意回路 $\ell:o\to o$ 与任意 $x\in F_o$，
若回路传输改变内部态，则层标记必改变：
\[
  T_{\ell}(x)\neq x\quad\Longrightarrow\quad \sigma\bigl(T_{\ell}(x)\bigr)\neq \sigma(x).
\]
这表示：任何“回到同一语义但暗处改写内部态”的作用都会在 $\widetilde O$ 上显式表现为“层切换”，
从而不再是 $\widetilde O$ 上的同点回路。
\end{definition}

\begin{theorem}[分层充要性：层完备性 $\Longleftrightarrow$ 抬升后 holonomy 平凡]
\label{thm:layering-complete-iff-trivial-holonomy}
在定义~\ref{def:layering-lift} 的语义抬升下，下列命题等价：
\begin{enumerate}
  \item 分层 $\sigma$ 层完备（定义~\ref{def:layering-complete}）；
  \item 对任意 $\widetilde o=(o,i)\in\widetilde O$ 与任意基于 $\widetilde o$ 的回路 $\widetilde\ell:\widetilde o\to \widetilde o$，
  其对纤维的 holonomy 作用平凡：
  \[
    T_{\widetilde\ell}=\mathrm{id}_{\widetilde F_{\widetilde o}}.
  \]
\end{enumerate}
\end{theorem}

\begin{Certificate}[证书：层完备性与抬升后平凡 holonomy 的互推]
\label{cert:layering-complete-iff-trivial-holonomy}
证书登记谓词链：
\[
\begin{aligned}
  L_1 &: \text{层完备性}\Rightarrow \forall x\in \widetilde F_{(o,i)},\ \widetilde\ell:(o,i)\to(o,i)\Rightarrow
    \sigma(T_{\widetilde\ell}(x))=\sigma(x),\\
  L_2 &: L_1\ \wedge\ \bigl(T_{\widetilde\ell}(x)\neq x\Rightarrow \sigma(T_{\widetilde\ell}(x))\neq \sigma(x)\bigr)
    \Rightarrow T_{\widetilde\ell}(x)=x,\\
  L_3 &: \forall x,\ T_{\widetilde\ell}(x)=x\Rightarrow T_{\widetilde\ell}=\mathrm{id},\\
  L_4 &: \bigl(\forall \widetilde\ell,\ T_{\widetilde\ell}=\mathrm{id}\bigr)\Rightarrow
  \bigl(T_{\ell}(x)\neq x\ \wedge\ \sigma(T_{\ell}(x))=\sigma(x)\bigr)\ \text{不可能}.
\end{aligned}
\]
\end{Certificate}

\begin{proof}[证明（谓词因果链）]
（1$\Rightarrow$2）设 $\sigma$ 层完备。取任意 $\widetilde o=(o,i)$、任意回路 $\widetilde\ell:\widetilde o\to\widetilde o$，
以及任意 $x\in\widetilde F_{\widetilde o}$。
由于 $\widetilde\ell$ 起终点层号相同，沿回路传输后仍落在同一纤维，
故 $\sigma(T_{\widetilde\ell}(x))=\sigma(x)=i$（$L_1$）。
若 $T_{\widetilde\ell}(x)\neq x$，则由层完备性应有 $\sigma(T_{\widetilde\ell}(x))\neq\sigma(x)$，矛盾。
因此对任意 $x$ 必有 $T_{\widetilde\ell}(x)=x$，即 $T_{\widetilde\ell}=\mathrm{id}_{\widetilde F_{\widetilde o}}$（$L_2$--$L_3$）。

（2$\Rightarrow$1）反之设对任意抬升回路 $\widetilde\ell$ 均有 $T_{\widetilde\ell}=\mathrm{id}$。
若层完备性不成立，则存在某个回路 $\ell:o\to o$ 与 $x\in F_o$ 满足 $T_{\ell}(x)\neq x$ 但 $\sigma(T_{\ell}(x))=\sigma(x)=:i$。
此时 $\widetilde o:=(o,i)$，且 $\ell$ 诱导一个抬升回路 $\widetilde\ell:\widetilde o\to\widetilde o$，
其在 $x\in \widetilde F_{\widetilde o}$ 上的作用满足 $T_{\widetilde\ell}(x)=T_{\ell}(x)\neq x$，与“$T_{\widetilde\ell}=\mathrm{id}
  $”矛盾（$L_4$）。
故层完备性成立。证毕。
\end{proof}

\begin{corollary}[推论：分层把 UU 显式化为“层切换”，从而可并入 $\Phi$-修复塔]
\label{cor:layering-makes-uu-repairable}
若系统在 $O$ 上存在未知-未知同伦类（定义~\ref{def:round-homotopy-uu-class}），
则任何层完备分层 $\sigma$ 会把对应的非平凡 holonomy 回路“抬升”为\textbf{非闭合路径}（层号改变），
从而 UU 的根源不再表现为“同语义点处的暗作用”，而表现为\textbf{可观测的层切换事件}。
在 \GFramework{} 的证书接口口径下，这正对应于：把 UU 见证编码为 $\JacGap>0$（定义~\ref{def:jacgap-from-obstruction}），
并允许通过“扩层/升阶”的同伦修复动作把 $\Phi=\mathcal{E}_{\mathrm{CH}}+\lambda\JacGap$ 逐层下降（定义~\ref{def:homotopy-repair-tower-phi}）。
\end{corollary}

\begin{Certificate}[证书：层完备性把非平凡 holonomy 从“回路”变为“层切换”]
\label{cert:layering-makes-uu-repairable}
证书登记谓词链：
\[
\begin{aligned}
  R_1 &: \text{UU 存在}\Rightarrow \exists o,[\ell],x\in F_o,\ T_{\ell}(x)\neq x\ \text{（定义~\ref{def:round-homotopy-uu-class}）},\\
  R_2 &: \text{层完备}\Rightarrow \bigl(T_{\ell}(x)\neq x\Rightarrow \sigma(T_{\ell}(x))\neq \sigma(x)\bigr)\
    \text{（定义~\ref{def:layering-complete}）},\\
  R_3 &: R_1\wedge R_2\Rightarrow \ell\ \text{在抬升空间中起终点层号不同，故不再是闭合回路},\\
  R_4 &: \text{层切换事件可作为 UU 见证并对接到 }\JacGap>0\ \text{与 }\Phi\text{-修复塔（定义~\ref{def:jacgap-from-obstruction}，\ref{def:homotopy-repair-tower-phi}）}.
\end{aligned}
\]
\end{Certificate}

\begin{proof}[证明（谓词因果链）]
由 UU 的定义（定义~\ref{def:round-homotopy-uu-class}），存在同一语义点 $o$ 上的回路作用在纤维 $F_o$ 上非平凡，
即存在 $x$ 使 $T_{\ell}(x)\neq x$，得 $R_1$。
若 $\sigma$ 层完备，则由定义~\ref{def:layering-complete} 得 $R_2$：非平凡回路作用必导致层号改变。
因此把起点记为 $\widetilde o=(o,\sigma(x))$，终点为 $(o,\sigma(T_{\ell}(x)))$，两者层号不同，
故该“暗回路作用”在抬升语义空间中表现为可观测的层切换而非闭合回路，得 $R_3$。
一旦层切换被显式化，就可以按本文既定接口把 UU 见证编码为 $\JacGap>0$ 并纳入 $\Phi$-修复塔的下降与扩层流程（$R_4$）。
推论成立。证毕。
\end{proof}

\begin{remark}[备注（“连续统分层作为相空间切换的离散统支撑”）]
定理~\ref{thm:uu-open-subdivision-not-enough} 说明：回合逼近只做“过程切片/路径细分”，不改变回路 holonomy；
因此在开放 UU 下，\textbf{仅靠细分回合无法消除 UU 根源}。
定理~\ref{thm:layering-complete-iff-trivial-holonomy} 则给出精确的结构补丁：
必须引入离散层指标 $\sigma$ 抬升语义相空间，把同一语义纤维的连续统收缩为可判别层，
从而将“暗处的单值性破坏”显式化为“层切换”（相空间切换）并纳入证书与修复塔。
\end{remark}

\section{挑战$\Rightarrow$智慧：多主体社会的博弈结构、封闭评测局限与 $\Phi$-修复塔}
\label{sec:challenge-wisdom-social-game}

\begin{remark}[目标（把三句自然语言论述压成可证书的命题链）]
本章把如下三句论述压成\textbf{可判定}的形式系统对象，并在 \GFramework{} 的符号骨架下给出可核验的“证书+证明（谓词因果链）”：
\begin{enumerate}
  \item \textbf{挑战$\Rightarrow$智慧}：挑战的结构作用不是“更难”，而是会系统性地产生 UU 缺口（$\JacGap>0$）与证书化修复需求（$\Phi$-下降）；
  \item \textbf{多主体社会$\Rightarrow$博弈}：当回合后继显式依赖其他主体的策略时，最优化问题在结构上等价为博弈问题；
  \item \textbf{“当下 AI 多停留在封闭层”应被可检验化}：该句作为经验判断不可直接演绎推出，
  但可以在形式系统内给出\textbf{可检验预测}——封闭评测只能检验 KU 能力，而开放 UU 能力必须通过“开封操作”引入并以 $\Phi$-下降/证书保持作为指标。
\end{enumerate}
为避免“强结论落空”，本章所有结论均只在本章给出的定义与前文已证定理链条内成立。
\end{remark}

\subsection{形式逻辑：语言、对象与证书接口（沿用 PureMath + Vol1 骨架）}
\label{sec:challenge-language-interface}

\begin{definition}[本章的形式语言（最小接口）]
\label{def:challenge-language-minimal}
本章沿用以下对象与接口（均已在前文出现）：
\begin{itemize}
  \item 法空间与法对象：$\GZLS$，$L\in\GZLS$；
  \item 回合分割与回合图：$\mathcal{P}$（定义~\ref{def:round-partition}）、$\Gamma_{\mathcal{P}}$（定义~\ref{def:round-graph-by-partition}
    ）；
  \item 真实后继与预设后继：$\mathrm{Post}^{(\mathcal{P})}(o,a)$ 与 $\widehat{\mathrm{Post}}^{(\mathcal{P})}(o,a)$（定义~\ref{def:round-post-model}）；
  \item UU 缺口（修复障碍）：$\mathsf{Ob}^{(\mathcal{P})}(o,a)$（定义~\ref{def:round-obstruction}）；
  \item UU 缺口的可计算缺陷证书：$\JacGap^{(\mathcal{P})}$（定义~\ref{def:jacgap-from-obstruction}，引理~\ref{lem:uu-witness-implies-jacgap-positive}）；
  \item 合并缺陷势：$\Phi=\mathcal{E}_{\mathrm{CH}}+\lambda\JacGap$（定义~\ref{def:phi-ch-jacgap}，$\lambda>0$）；
  \item 受保护一致性证书谓词：$\mathcal{E}_{\mathrm{prot}}$（同定义~\ref{def:homotopy-repair-tower-phi}）；
  \item 同伦修复塔：在 $\mathcal{E}_{\mathrm{prot}}$ 保持下使 $\Phi$ 逐层下降（定义~\ref{def:homotopy-repair-tower-phi}，定理~\ref{thm:repair-tower-terminates}）。
\end{itemize}
此外，若需显式书写“证书对象与验证器接口”，可沿用证书化解三元组
$(x,\pi,\mathsf{Ver})$（定义~\ref{def:magic-wand-certified-solution}）。
\end{definition}

\begin{remark}[本章的“严格性”口径]
本章不对“现实世界是否满足某性质”做经验断言；而是给出：
若系统（或评测/训练框架）满足某些可判定结构条件，则必然推出相应结论。
因此“严格”体现在谓词因果链闭合，而非依赖外部经验。
\end{remark}

\subsection{公理降级：回合归纳证书（Axiom$\to$Theorem 的统一模板）}
\label{sec:challenge-round-induction}

\begin{theorem}[公理降级：回合归纳证书模板（对任意回合性质）]
\label{thm:challenge-round-induction-schema}
设 $P(k)$ 为关于回合索引 $k\in\NN$ 的谓词（例如“到第 $k$ 回合为止证书保持/势能下降/约束满足”等）。
若满足
\[
  P(0)\qquad\text{且}\qquad \forall k\in\NN,\ \bigl(P(k)\Rightarrow P(k+1)\bigr),
\]
则有 $\forall k\in\NN,\ P(k)$。
\end{theorem}

\begin{Certificate}[数学归纳法证书：$(P(0),\forall k(P(k)\Rightarrow P(k+1)))$]
\label{cert:challenge-round-induction-schema}
证书对象可取为
\[
  \pi_{\mathrm{ind}}
  :=\bigl(P(0),\ \forall k\in\NN\,(P(k)\Rightarrow P(k+1))\bigr).
\]
\end{Certificate}

\begin{proof}[证明（数学归纳法证书；谓词因果链）]
由证书~\ref{cert:challenge-round-induction-schema}：
$P(0)$ 给出基步；归纳步给出 $P(k)\Rightarrow P(k+1)$。
从而 $P(0)\Rightarrow P(1)\Rightarrow\cdots\Rightarrow P(k)$，对任意 $k$ 成立。
因此 $\forall k,\ P(k)$。证毕。
\end{proof}

\subsection{形式逻辑：挑战谓词（漂移/反制/多约束）与 UU 压力的证书化}
\label{sec:challenge-predicates}

\begin{definition}[三类挑战谓词（最小抽象）]
\label{def:challenge-three-pressures}
固定一个回合分割 $\mathcal{P}$ 与一族预设后继 $\widehat{\mathrm{Post}}^{(\mathcal{P})}$。
在该口径下定义三类“挑战压力”谓词：
\begin{enumerate}
  \item \textbf{漂移压力（Drift）}：存在时间 $t\neq t'$ 与某个 $(o,a)$ 使真实后继集合发生“新增”：
  \[
    \mathsf{Drift}
    \ :\Longleftrightarrow\
    \exists t\neq t',\exists(o,a):\ \mathrm{Post}^{(\mathcal{P})}_{t}(o,a)\not\subseteq \mathrm{Post}^{(\mathcal{P})}
      _{t'}(o,a).
  \]
  \item \textbf{反制压力（Counter）}：存在对方策略差异导致真实后继集合发生“新增”：
  \[
    \mathsf{Counter}
    \ :\Longleftrightarrow\
    \exists \pi^{-}\neq \tilde\pi^{-},\exists(o,a):\ \mathrm{Post}^{(\mathcal{P})}(o,a\mid \pi^{-})\not\subseteq
      \mathrm{Post}^{(\mathcal{P})}(o,a\mid \tilde\pi^{-}).
  \]
  \item \textbf{多约束压力（Constraint）}：受保护一致性证书谓词随情境改变：
  \[
    \mathsf{Constraint}
    \ :\Longleftrightarrow\
    \exists t\neq t':\ \mathcal{E}_{\mathrm{prot},t}\neq \mathcal{E}_{\mathrm{prot},t'}.
  \]
\end{enumerate}
定义综合挑战谓词为
\[
  \mathsf{Challenge}:=\mathsf{Drift}\ \vee\ \mathsf{Counter}\ \vee\ \mathsf{Constraint}.
\]
\end{definition}

\begin{proposition}[挑战$\Rightarrow$ UU 压力（$\Phi>0$ 的证书化见证）]
\label{prop:challenge-implies-phi-pressure}
在定义~\ref{def:challenge-three-pressures} 的口径下，若 $\mathsf{Challenge}$ 成立且预设后继族与证书接口被\textbf{固定}（不随 $t$ 或对手策略自适应扩张），
则必存在某个见证使得在某个回合口径上出现：
\[
  \text{要么}\ \JacGap^{(\mathcal{P})}>0\ (\text{UU 缺口证书}),\qquad
  \text{要么}\ \mathcal{E}_{\mathrm{CH}}>0\ (\text{已知约束残差}),
\]
从而 $\Phi=\mathcal{E}_{\mathrm{CH}}+\lambda\JacGap>0$。
\end{proposition}

\begin{Certificate}[证书：挑战见证 $\Rightarrow$ $\JacGap>0$ 或 $\mathcal{E}_{\mathrm{CH}}>0$]
\label{cert:challenge-implies-phi-pressure}
证书登记谓词链（给出三种分支之一即可）：
\[
\begin{aligned}
  D_1 &: \mathsf{Drift}\ \Rightarrow\ \exists t\neq t',\exists(o,a):\ \mathrm{Post}^{(\mathcal{P})}_{t}(o,a)
    \not\subseteq \widehat{\mathrm{Post}}^{(\mathcal{P})}(o,a),\\
  D_2 &: D_1\Rightarrow \exists(o,a):\ \mathsf{Ob}^{(\mathcal{P})}(o,a)\neq\varnothing\Rightarrow \JacGap^{(\mathcal{P})}
    >0,\\
  C_1 &: \mathsf{Counter}\ \Rightarrow\ \exists\pi^{-}\neq\tilde\pi^{-},\exists(o,a):\ \mathrm{Post}^{(\mathcal{P})}(o,
    a\mid \pi^{-})\not\subseteq \widehat{\mathrm{Post}}^{(\mathcal{P})}(o,a),\\
  C_2 &: C_1\Rightarrow \exists(o,a):\ \mathsf{Ob}^{(\mathcal{P})}(o,a)\neq\varnothing\Rightarrow \JacGap^{(\mathcal{P})}
    >0,\\
  S_1 &: \mathsf{Constraint}\ \Rightarrow\ \exists t\neq t':\ \mathcal{E}_{\mathrm{prot},t}\neq \mathcal{E}
    _{\mathrm{prot},t'}\Rightarrow \mathcal{E}_{\mathrm{CH}}>0,\\
  P_1 &: (D_2\ \vee\ C_2\ \vee\ S_1)\Rightarrow \Phi=\mathcal{E}_{\mathrm{CH}}+\lambda\JacGap>0.
\end{aligned}
\]
其中 $D_2,C_2$ 的“$\mathsf{Ob}\neq\varnothing\Rightarrow\JacGap>0$”使用引理~\ref{lem:uu-witness-implies-jacgap-positive}。
\end{Certificate}

\begin{proof}[证明（谓词因果链）]
若 $\mathsf{Drift}$ 成立，则存在时间点使真实后继集合在固定模型口径下出现新增，
从而某个真实后继落入 $\widehat{\mathrm{Post}}^{(\mathcal{P})}$ 之外，得到 $D_1$。
由定义~\ref{def:round-obstruction} 得 $\mathsf{Ob}^{(\mathcal{P})}(o,a)\neq\varnothing$，
再由引理~\ref{lem:uu-witness-implies-jacgap-positive} 得 $\JacGap^{(\mathcal{P})}>0$，得到 $D_2$。

若 $\mathsf{Counter}$ 成立，同理得到 $C_1\Rightarrow C_2$。

若 $\mathsf{Constraint}$ 成立，则“受保护一致性证书谓词”发生改变。
在不更新证书接口的前提下，至少会产生一类“已知约束残差”（例如旧谓词允许而新谓词禁止），记入 $\mathcal{E}_{\mathrm{CH}}>0$，得到 $S_1$。

三种分支任取其一即可推出 $\Phi=\mathcal{E}_{\mathrm{CH}}+\lambda\JacGap>0$（$P_1$）。
故命题成立。证毕。
\end{proof}

\begin{remark}[对齐：“挑战塑造智慧”的形式含义]
命题~\ref{prop:challenge-implies-phi-pressure} 表明：挑战的结构作用是制造 $\Phi>0$ 的可证书化压力，
其中 KU 压力主要体现在 $\mathcal{E}_{\mathrm{CH}}>0$，UU 压力体现在 $\JacGap>0$。
若仍要在 $\mathcal{E}_{\mathrm{prot}}$ 保持下把 $\Phi$ 压回到低水平，则必须引入同伦修复塔（定义~\ref{def:homotopy-repair-tower-phi}）并能证书化地推动 $\Phi$
  下降（定理~\ref{thm:repair-tower-terminates}）。
这就是“挑战$\Rightarrow$智慧”的严格口径：挑战迫使系统具备“证书保持 + 势能可下降”的结构能力，而非仅仅更大算力。
\end{remark}

\subsection{定理：多主体社会 $\Rightarrow$ 决策问题在结构上等价为博弈}
\label{sec:multiagent-implies-game}

\begin{definition}[智能体（决策规则）]
\label{def:agent-policy-decision}
称智能体（决策器）为一个从回合历史到动作的映射
\[
  \pi:\ H_{\mathcal{P}}\to A_{\mathcal{P}},
\]
其中 $H_{\mathcal{P}}$ 为可观测历史（或其可计算摘要），$A_{\mathcal{P}}$ 为回合动作摘要空间。
\end{definition}

\begin{definition}[多主体依赖（后继集合依赖他方策略）]
\label{def:multiagent-dependence}
称系统具有多主体依赖，若存在至少一个他方策略 $\pi^{-}$ 使真实后继集合依赖该策略：
\[
  \mathrm{Post}^{(\mathcal{P})}(o,a\mid \pi^{-}).
\]
并且存在 $\pi^{-}\neq \tilde\pi^{-}$ 与某个 $(o,a)$ 使
\[
  \mathrm{Post}^{(\mathcal{P})}(o,a\mid \pi^{-})\neq \mathrm{Post}^{(\mathcal{P})}(o,a\mid \tilde\pi^{-}).
\]
\end{definition}

\begin{theorem}[多主体存在 $\Rightarrow$ 决策最优化等价为博弈问题]
\label{thm:multiagent-implies-game}
若系统满足多主体依赖（定义~\ref{def:multiagent-dependence}），
则“求最优决策策略 $\pi$”的对象不再是单一转移关系，而是一个由他方策略参数化的映射族；
因此该最优化问题在结构上等价为至少两方的博弈问题（策略相互依赖）。
\end{theorem}

\begin{Certificate}[证书：后继依赖他方策略 $\Rightarrow$ 优化对象为策略对]
\label{cert:multiagent-implies-game}
证书登记谓词链：
\[
\begin{aligned}
  G_1 &: \exists \pi^{-}\neq\tilde\pi^{-},\exists(o,a):\ \mathrm{Post}(o,a\mid \pi^{-})\neq \mathrm{Post}(o,a\mid
    \tilde\pi^{-}),\\
  G_2 &: G_1\Rightarrow \text{我方策略 $\pi$ 的效用/约束取决于他方策略选择},\\
  G_3 &: G_2\Rightarrow \text{求解对象需在 $(\pi,\pi^{-})$ 的相互依赖结构中定义},\\
  G_4 &: \Rightarrow \text{结构上为博弈（至少两方策略空间上的问题）}.
\end{aligned}
\]
\end{Certificate}

\begin{proof}[证明（谓词因果链）]
由多主体依赖定义直接得到 $G_1$。
当 $\mathrm{Post}(o,a\mid \pi^{-})$ 随 $\pi^{-}$ 改变时，我方任意策略 $\pi$ 在同一 $(o,a)$ 上的后继集合、可行性与收益评价均随之改变，
得到 $G_2$。
因此“最优 $\pi$”无法在固定转移上定义，必须以 $(\pi,\pi^{-})$ 为共同变量定义求解对象，得到 $G_3$。
这正是博弈问题的结构定义，得到 $G_4$。证毕。
\end{proof}

\subsection{推论：封闭评测只能检验封闭能力；开放能力必须以 $\Phi$-下降/证书保持衡量}
\label{sec:closed-eval-only-closed}

\begin{corollary}[封闭训练/评测高分不蕴含开放 UU 修复能力]
\label{cor:closed-eval-does-not-imply-open}
若训练/评测框架在某个回合口径下满足“封闭”（例如 $\JacGap$ 稳定为 $0$；定理~\ref{thm:closed-phi-jacgap-ku}），
则模型在该框架上取得高分仅能证明其在 KU 障碍下的决策能力，
并不蕴含其具备开放 UU 下的证书保持与 $\Phi$-修复能力。
\end{corollary}

\begin{Certificate}[证书：封闭 $\Rightarrow$ 仅 KU；高分不约束 UU]
\label{cert:closed-eval-does-not-imply-open}
证书登记谓词链：
\[
\begin{aligned}
  E_1 &: \text{封闭评测口径}\Rightarrow \JacGap\ \text{稳定为 }0\Rightarrow \text{障碍仅为 KU（定理~\ref{thm:closed-phi-jacgap-ku}）},\\
  E_2 &: E_1\Rightarrow \text{评测只要求在已知后继集合内择优（集合内落点未知）},\\
  E_3 &: E_2\Rightarrow \text{高分并未检验“$\JacGap>0$ 时的修复塔/证书保持”能力}.
\end{aligned}
\]
\end{Certificate}

\begin{proof}[证明（谓词因果链）]
由定理~\ref{thm:closed-phi-jacgap-ku}，封闭口径下 $\JacGap$ 稳定归零，UU 缺口不出现，仅剩 KU 不确定性，得 $E_1$。
因此评测本质是“在集合闭合前提下做最优化”，得 $E_2$。
而开放 UU 能力的判据是“在 $\JacGap>0$ 时仍能证书保持并推动 $\Phi$ 下降”（定义~\ref{def:homotopy-repair-tower-phi}），
该判据在封闭评测中从未被触发，故高分不蕴含开放修复能力，得 $E_3$。证毕。
\end{proof}

\begin{definition}[五类开封算子（把封闭评测逐步打开为可检验压力族）]
\label{def:open-eval-five-operators}
给定一个封闭口径下的法对象 $L$，定义五类“开封算子”（仅作抽象符号，不需要指定具体实现细节）：
\[
\begin{aligned}
  \mathsf{OpRule}_{\varepsilon}(L) &: \text{规则轻改（小幅改动 $\Sigma$/接口）},\\
  \mathsf{Hide}_{\eta}(L) &: \text{隐信息与噪声增强（部分可观测）},\\
  \mathsf{Adv}_{\kappa}(L) &: \text{对手在线学习/适应（反制压力）},\\
  \mathsf{MultiObj}(L) &: \text{多目标多约束（证书谓词变复杂）},\\
  \mathsf{Meta}(L) &: \text{引入第三方规则塑形者（平台/仲裁/监管）}.
\end{aligned}
\]
这些算子对应把封闭评测逐步推向“漂移/反制/多约束”的三类挑战谓词（定义~\ref{def:challenge-three-pressures}）。
\end{definition}

\begin{remark}[可检验化指示器（把经验判断转为结构预测）]
在上述开封算子族下，“是否停留在封闭层”的判断可被替换为可检验指标：
当压力逐步打开时，模型是否仍能在 $\mathcal{E}_{\mathrm{prot}}$ 保持下压制 $\JacGap$ 并推动 $\Phi$ 下降（修复塔可终止；定理~\ref{thm:repair-tower-terminates}）
  。
\end{remark}

\subsection{形式逻辑：开放系统下六类能力谓词（对应治理/博弈/长期性需求）}
\label{sec:six-capabilities}

\begin{definition}[六类能力谓词（抽象但可判定的结构占位）]
\label{def:six-capability-predicates}
对决策策略 $\pi$（定义~\ref{def:agent-policy-decision}），定义六类能力谓词（仅要求其存在/可判定，不预设具体算法形态）：
\begin{enumerate}
  \item 情境识别 $C_1(\pi)$：存在状态分类器 $\chi$ 使策略可条件化于“规则内/边缘/被重写”等语义态；
  \item 对手建模 $C_2(\pi)$：存在可更新的他方策略估计并进入决策（使 $\pi$ 对 $\pi^{-}$ 的信念敏感）；
  \item 信息管理 $C_3(\pi)$：当动作含表达/承诺通道时，策略可控信息释放以影响后继集合；
  \item 联盟与机制 $C_4(\pi)$：存在可验证承诺/契约结构并能改变后继集合族；
  \item 鲁棒下界 $C_5(\pi)$：能在对抗集合上保证收益/安全下界（对应“保下界”）；
  \item 长期治理 $C_6(\pi)$：能把长期副作用编码进 $\mathcal{E}_{\mathrm{prot}}$ 并在线保持（不因短期最优破坏证书）。
\end{enumerate}
这些谓词作为“必要结构占位符”，用于把“能力清单”从口头建议降级为可推演对象。
\end{definition}

\subsection{推论：压力$\Rightarrow$能力必要性样式（以 $\Phi$-下降为判据）}
\label{sec:pressure-implies-capability}

\begin{corollary}[缺少对应能力 $\Rightarrow$ 无法保证 $\Phi$ 可证书化下降（代表性三链）]
\label{cor:pressure-implies-capability-necessity}
在挑战谓词 $\mathsf{Challenge}$ 可持续出现的开放口径下，
若缺少与压力匹配的能力谓词（定义~\ref{def:six-capability-predicates}），
则存在一类对抗序列使得无法保证在 $\mathcal{E}_{\mathrm{prot}}$ 保持下形成“$\Phi$ 严格下降”的修复塔，
从而无法套用定理~\ref{thm:repair-tower-terminates} 得到有限步可终止性。
代表性地，至少有下列三条必要性链条：
\begin{enumerate}
  \item （漂移$\Rightarrow$需 $C_1$）若 $\mathsf{Drift}$ 成立但缺少 $C_1$，则存在漂移序列使策略无法识别情境改变，
  从而不断触发 $\JacGap>0$ 的 UU 见证，$\Phi$ 无法被保证单调下降；
  \item （反制$\Rightarrow$需 $C_2$）若 $\mathsf{Counter}$ 成立但缺少 $C_2$，则存在对手在线反制使后继集合发生策略依赖性改变，
  触发 UU 缺口并破坏“对抗无关的鲁棒下降”；
  \item （多约束$\Rightarrow$需 $C_6$）若 $\mathsf{Constraint}$ 成立但缺少 $C_6$，则存在路径使短期动作触发未来证书破坏（$\mathcal{E}_{\mathrm{prot}}
    =\bot$），
  因而无法维持“证书保持 + $\Phi$-下降”链条。
\end{enumerate}
\end{corollary}

\begin{Certificate}[证书：缺少能力 $\Rightarrow$ 下降条件被击穿]
\label{cert:pressure-implies-capability-necessity}
证书登记谓词链（以漂移分支为例，反制/多约束同理）：
\[
\begin{aligned}
  N_1 &: \mathsf{Drift}\ \wedge\ \neg C_1(\pi)\Rightarrow \exists\ \text{漂移序列使}\ \pi\ \text{无法条件化更新},\\
  N_2 &: N_1\Rightarrow \exists(o,a):\ \mathsf{Ob}^{(\mathcal{P})}(o,a)\neq\varnothing\Rightarrow \JacGap^{(\mathcal{P})}
    >0,\\
  N_3 &: N_2\Rightarrow \neg\bigl(\forall k,\ \Phi(L_{k+1})\le \Phi(L_k)-\delta_k\bigr)\ \text{（无法保证严格下降）}.
\end{aligned}
\]
\end{Certificate}

\begin{proof}[证明（谓词因果链）]
以漂移分支为例：$\neg C_1(\pi)$ 表示策略无法识别“规则/分布已变”的语义态，仍按旧口径选择动作。
当 $\mathsf{Drift}$ 持续出现时，存在某次漂移使真实后继超出固定模型口径，得到 $\mathsf{Ob}\neq\varnothing$，
从而由引理~\ref{lem:uu-witness-implies-jacgap-positive} 得 $\JacGap>0$，即 $N_2$。
若缺少对漂移的识别与对应修复动作族，则无法对所有后继保证“证书保持 + 势能严格下降”，从而击穿修复塔下降条件，得到 $N_3$。
其余两分支同理。证毕。
\end{proof}

\subsection{定理：总论（智慧/通用智能在开放式多主体社会中的结构定义）}
\label{sec:general-intelligence-structural}

\begin{definition}[开放式通用智能（以 $\Phi$-下降与证书保持刻画）]
\label{def:general-intelligence-by-phi}
称策略 $\pi$ 具备“开放式通用智能（在多主体社会口径下）”，若存在一条同伦修复塔
$(L_k)$ 使对由 $\mathsf{Challenge}$ 诱导的扰动序列，在 $\mathcal{E}_{\mathrm{prot}}$ 保持下都能证书化地满足：
\[
  \Phi(L_{k+1})\le \Phi(L_k)-\delta_k,\qquad \delta_k>0,
\]
从而可调用定理~\ref{thm:repair-tower-terminates} 给出有限步“修复完成/压回可接受区间”的保证。
\end{definition}

\begin{theorem}[总论：挑战迫使“证书保持 + $\Phi$-下降”的结构成为智慧的必要底座]
\label{thm:challenge-shapes-wisdom-structurally}
在上述定义下：
\begin{enumerate}
  \item 若挑战谓词 $\mathsf{Challenge}$ 持续出现，则系统持续产生 $\Phi>0$ 的结构压力（命题~\ref{prop:challenge-implies-phi-pressure}）；
  \item 若仍要求可持续决策（长期保持 $\mathcal{E}_{\mathrm{prot}}=\top$ 并将 $\Phi$ 压回低水平），则必须存在满足定义~\ref{def:general-intelligence-by-phi} 的修复塔结构；
  \item 当环境包含多主体依赖时，该智能结构的运行对象在结构上等价为博弈（定理~\ref{thm:multiagent-implies-game}）。
\end{enumerate}
因此，“挑战$\Rightarrow$智慧”在形式系统内等价于：
挑战不断制造可证书化缺陷势压力，而智慧就是在证书保持下组织可终止的 $\Phi$-下降修复塔（并在多主体时体现为博弈式决策）。
\end{theorem}

\begin{Certificate}[证书：$\mathsf{Challenge}\Rightarrow \Phi>0$；$\Phi$ 可终止下降 $\Rightarrow$ 智慧结构]
\label{cert:challenge-shapes-wisdom-structurally}
证书登记谓词链：
\[
\begin{aligned}
  W_1 &: \mathsf{Challenge}\Rightarrow \Phi>0\ \text{（命题~\ref{prop:challenge-implies-phi-pressure}）},\\
  W_2 &: \text{可持续性需求}\Rightarrow \exists (L_k)\ \text{满足定义~\ref{def:general-intelligence-by-phi} 的下降条件},\\
  W_3 &: W_2\Rightarrow \text{可调用定理~\ref{thm:repair-tower-terminates} 给出有限步保证},\\
  W_4 &: \text{多主体依赖}\Rightarrow \text{决策结构为博弈（定理~\ref{thm:multiagent-implies-game}）}.
\end{aligned}
\]
\end{Certificate}

\begin{proof}[证明（谓词因果链）]
$W_1$ 由命题~\ref{prop:challenge-implies-phi-pressure} 直接给出：挑战产生 $\Phi>0$ 的证书化压力。
若仍要求长期保持 $\mathcal{E}_{\mathrm{prot}}=\top$ 并把 $\Phi$ 压回低水平，则必须存在使 $\Phi$ 在证书保持下可严格下降的修复结构（定义~\ref{def:general-intelligence-by-phi}），即 $W_2$。
由定理~\ref{thm:repair-tower-terminates} 得有限步可终止性，得到 $W_3$。
若环境满足多主体依赖，则由定理~\ref{thm:multiagent-implies-game} 得该结构在运行时体现为博弈式决策，得到 $W_4$。
合并即得定理结论。证毕。
\end{proof}

\section{统一证书 Schema：输入/输出/验证器谓词/失败回退规则}
\label{sec:unified-certificate-schema}

\begin{remark}[目标（把“定理/证明”收束为可落地的证书接口）]
前文已经给出一组可引用的核心结论（封闭/开放、UU、分层必要性、神经网络滞后、审计包装等）。
为了将这些结论直接沉淀为工程闭环接口（可审计、可回退、可证书化），本章给出一套\textbf{统一证书 schema}：
对任一编号断言 $T$，都以“输入/输出/验证器谓词/失败回退规则”的合同方式编码，并与您既有两段式工程模板对齐：
\[
  F^{\mathrm{aud}}_{\theta}=V\circ G_{\theta}:\ \mathsf{WP}\to \mathsf{NF}\sqcup\{\bot\}
  \qquad(\text{候选生成 + 确定性裁决；fail-closed}).
\]
接口参考：\code{/mnt/data/gframework_semantic_functor_field_gauge_field_normalization_interface_formal_system.tex}，
\code{/mnt/data/gframework_semantic_functor_regex_normalization_validation_robustness_formal_system.tex}。
\end{remark}

\subsection{形式逻辑：共用对象域与记号（所有证书共用）}
\label{sec:certschema-shared-notation}

\begin{remark}[共用对象域（引用既有定义而不重复造轮子）]
本章证书 schema 统一沿用下列既有对象与定义口径：
\begin{itemize}
  \item 回合分割与回合图：定义~\ref{def:round-partition}、\ref{def:round-graph-by-partition}；
  \item KU/UU 与闭合缺口：定义~\ref{def:round-obstruction}、\ref{def:round-ku-uu}；
  \item $\mathsf{Ob}\Rightarrow \JacGap \Rightarrow \Phi$ 对接：定义~\ref{def:jacgap-from-obstruction} 与章节~\ref{sec:ob-jacgap-phi-repair-tower}；
  \item 分层抬升与 holonomy 消除：定理~\ref{thm:layering-complete-iff-trivial-holonomy}；
  \item 必然滞后同伦类口径：定理~\ref{thm:lag-holonomy-to-lag-class} 与推论~\ref{cor:demote-nn-from-main-loop}；
  \item 审计包装接口：定义~\ref{def:audited-wrapper-definition}、\ref{def:audited-wrapper-soundness}，
  以及定理~\ref{thm:audited-wrapper-soundness}、\ref{thm:audited-wrapper-inductive-safety}。
\end{itemize}
因此，本章不再新增新的“世界观”，只把既有定理链统一包装为可执行的证书合同。
\end{remark}

\subsection{定义：统一证书 schema（七元组）与验证器谓词}
\label{sec:certschema-definition}

\begin{definition}[统一证书 schema（七元组）]
\label{def:certificate-schema-7tuple}
对任一编号断言（定理/命题/推论均可）记为 $T$，统一把证书对象写成七元组：
\[
  \pi_T :=
  \bigl(\mathrm{id},\ \mathrm{input},\ \mathrm{claim},\ \mathrm{assump},\ \mathrm{witness},\ \mathrm{oblig},\
    \mathrm{fallback}\bigr).
\]
其中各字段语义如下：
\begin{itemize}
  \item $\mathrm{id}$：证书编号（如 \code{Cert-T\_Open-0007}）；
  \item $\mathrm{input}$：该断言实例的输入对象（如分割 $\mathcal{P}$、法对象 $L$、语义态 $o$ 等）；
  \item $\mathrm{claim}$：该实例要宣称成立的结论（断言 $T$ 的实例化）；
  \item $\mathrm{assump}$：该实例依赖的前提（可引用其它证书/定理）；
  \item $\mathrm{witness}$：见证数据（如 UU 反例 $o'\in \mathsf{Ob}$、非平凡回路 $[\ell]$、分层映射 $\sigma$ 等）；
  \item $\mathrm{oblig}$：验证义务列表（Proof obligations），即验证器必须逐项核验的条件集合；
  \item $\mathrm{fallback}$：失败回退规则（fail-closed 行为 + 触发修复塔的动作）。
\end{itemize}
\end{definition}

\begin{definition}[统一验证器谓词（证书健全性合同）]
\label{def:certificate-verifier-predicate}
对每个断言 $T$，定义验证器谓词
\[
  \mathsf{Ver}_T(\mathrm{input},\pi_T)\in\{\top,\bot\},
\]
其工程实现可被理解为“逐项检查 $\mathrm{oblig}$，并记录审计日志/回归测试结果”。
要求其满足健全性合同：
\[
  \mathsf{Ver}_T(\mathrm{input},\pi_T)=\top\ \Longrightarrow\ \mathrm{claim}\ \text{成立。}
\]
若无法核验，则必须 fail-closed：输出 $\bot$ 或拒绝执行，并按 $\mathrm{fallback}$ 触发修复塔。
\end{definition}

\subsection{工程合同：JSON 外壳（可审计、可回退、失败显式）}
\label{sec:certschema-json}

\begin{remark}[工程合同（与 regex 审计闭环一致）]
为对齐工程落地与审计闭环，推荐使用“失败必须显式（fail-closed）”的最小合同外壳：
成功时返回产物与证书，失败时返回原因与见证（用于触发修复塔）。
\end{remark}

\begin{verbatim}
{"ok": true,
 "theorem": "T_AuditedWrapper",
 "artifact": {...},
 "certificate": {
   "id":"Cert-T_AuditedWrapper-0001",
   "oblig":[...],
   "fallback": {...}
 }}
\end{verbatim}

\begin{verbatim}
{"ok": false,
 "theorem": "T_Open",
 "reason": "UU witness",
 "witness": {...},
 "fallback": {
   "action": "repair_tower_expand_or_layer"
 }}
\end{verbatim}

\subsection{六个核心定理：证书规格（字段/验证器/回退规则）}
\label{sec:certschema-six-theorems}

\begin{definition}[六个核心定理的命名对齐（别名）]
\label{def:certschema-six-aliases}
为便于在工程与卷册中统一引用，约定下列“核心定理别名”对应本文既有结论：
\[
\begin{aligned}
  T_{\mathrm{Closed}} &:\ \text{封闭系统（KU；严格回合同伦类存在）}\ \leftrightarrow\ \text{定理~\ref{thm:closed-phi-jacgap-ku}},\\
  T_{\mathrm{Open}} &:\ \text{开放系统（UU；严格回合同伦类不存在）}\ \leftrightarrow\ \text{定理~\ref{thm:open-phi-jacgap-uu}},\\
  T_{\mathrm{NoRoundOnly}} &:\ \text{仅回合细分无法消除 UU}\ \leftrightarrow\ \text{定理~\ref{thm:uu-open-subdivision-not-enough}},
    \\
  T_{\mathrm{LayeringNecessary}} &:\ \text{分层必要性（层完备性$\Leftrightarrow$杀死 holonomy）}\ \leftrightarrow\ \text{定理~\ref{thm:layering-complete-iff-trivial-holonomy}},\\
  T_{\mathsf{Lag}} &:\ \text{必然滞后同伦类}\ \leftrightarrow\ \text{定理~\ref{thm:lag-holonomy-to-lag-class}},\\
  T_{\mathrm{AuditedWrapper}} &:\ \text{审计包装}\ \leftrightarrow\ \text{定理~\ref{thm:audited-wrapper-soundness} 与~\ref{thm:audited-wrapper-inductive-safety}}.
\end{aligned}
\]
\end{definition}

\subsubsection{$T_{\mathrm{Closed}}$：封闭系统（KU；严格回合同伦类存在）}
\label{sec:certschema-t-closed}

\begin{theorem}[证书规格：$T_{\mathrm{Closed}}$（封闭系统）]
\label{thm:certspec-t-closed}
若存在证书 $\pi_{\mathrm{Closed}}$ 使
$\mathsf{Ver}_{T_{\mathrm{Closed}}}(\mathcal{P}_0,\pi_{\mathrm{Closed}})=\top$，
则该实例的结论 $\mathrm{claim}$ 成立：在分割 $\mathcal{P}_0$ 口径下 $\JacGap$ 稳定为 $0$，障碍必为 KU，
并且严格回合同伦类存在（对应定理~\ref{thm:closed-phi-jacgap-ku}）。
\end{theorem}

\begin{Certificate}[证书字段：$T_{\mathrm{Closed}}$]
\label{cert:certspec-t-closed}
可取
\[
  \pi_{\mathrm{Closed}}=
  \bigl(\mathrm{id},\ \mathcal{P}_0,\ \mathrm{claim},\ \mathrm{assump},\ \mathrm{witness},\ \mathrm{oblig},\
    \mathrm{fallback}\bigr),
\]
其中：
\begin{itemize}
  \item $\mathrm{input}$：$\mathcal{P}_0$ 与预设后继族 $\widehat{\mathrm{Post}}^{(\mathcal{P})}$（定义~\ref{def:round-post-model}）；

  \item $\mathrm{claim}$：$\JacGap$ 稳定为 $0$，障碍仅为 KU，且严格回合同伦类存在；
  \item $\mathrm{oblig}$（验证义务）：对所有足够细 $\mathcal{P}\succ\mathcal{P}_0$，都有
  $\forall(o,a),\ \mathsf{Ob}^{(\mathcal{P})}(o,a)=\varnothing$（模型后继闭合），从而 $\JacGap^{(\mathcal{P})}=0$（定义~\ref{def:jacgap-from-obstruction}）；
  \item $\mathrm{fallback}$：若出现任一 UU 见证（$\mathsf{Ob}^{(\mathcal{P})}(o,a)\neq\varnothing$），则 fail-closed 并切换到
    $T_{\mathrm{Open}}$ 与修复塔（扩张/分层）。
\end{itemize}
验证器可取为“逐项核验 $\mathrm{oblig}$，并记录审计日志”：
\[
  \mathsf{Ver}_{T_{\mathrm{Closed}}}(\mathcal{P}_0,\pi_{\mathrm{Closed}})=\top\ \Longleftrightarrow\ \mathrm{oblig}\
    \text{全部通过}.
\]
\end{Certificate}

\begin{proof}[证明（谓词因果链）]
由 $\mathsf{Ver}_{T_{\mathrm{Closed}}}(\mathcal{P}_0,\pi_{\mathrm{Closed}})=\top$ 与定义~\ref{def:certificate-verifier-predicate}，
得 $\mathrm{oblig}$ 均成立，从而 $\JacGap$ 稳定归零且 UU 不出现。
因此可直接套用定理~\ref{thm:closed-phi-jacgap-ku} 得到 $\mathrm{claim}$，证毕。
\end{proof}

\subsubsection{$T_{\mathrm{Open}}$：开放系统（UU；严格回合同伦类不存在）}
\label{sec:certschema-t-open}

\begin{theorem}[证书规格：$T_{\mathrm{Open}}$（开放系统）]
\label{thm:certspec-t-open}
若存在证书 $\pi_{\mathrm{Open}}$ 使
$\mathsf{Ver}_{T_{\mathrm{Open}}}(\mathcal{P},\pi_{\mathrm{Open}})=\top$，
则该实例的结论 $\mathrm{claim}$ 成立：存在 UU 见证使 $\JacGap>0$，并且严格回合同伦类无法稳定存在（对应定理~\ref{thm:open-phi-jacgap-uu}）。
\end{theorem}

\begin{Certificate}[证书字段：$T_{\mathrm{Open}}$（UU 反例证书）]
\label{cert:certspec-t-open}
可取
\[
  \pi_{\mathrm{Open}}=
  \bigl(\mathrm{id},\ \mathcal{P},\ \mathrm{claim},\ \mathrm{assump},\ \mathrm{witness},\ \mathrm{oblig},\
    \mathrm{fallback}\bigr),
\]
其中 $\mathrm{witness}$ 给出一个可核验 UU 反例（见证）：
\[
  \mathcal{W}(\mathcal{P})=(\mathcal{P}',o,a,o'),\qquad \mathcal{P}'\succ\mathcal{P},
\]
满足
$o'\in \mathrm{Post}^{(\mathcal{P}')}(o,a)$ 且 $o'\notin \widehat{\mathrm{Post}}^{(\mathcal{P}')}(o,a)$，
因此 $\mathsf{Ob}^{(\mathcal{P}')}(o,a)\neq\varnothing$，
从而 $\JacGap^{(\mathcal{P}')}>0$（引理~\ref{lem:uu-witness-implies-jacgap-positive}）。
失败回退 $\mathrm{fallback}$ 取为：触发同伦修复塔（扩张后继族/分层抬升/增加证书规则）。
\end{Certificate}

\begin{proof}[证明（谓词因果链）]
由 $\mathsf{Ver}_{T_{\mathrm{Open}}}(\mathcal{P},\pi_{\mathrm{Open}})=\top$ 得见证 $\mathcal{W}(\mathcal{P})$ 成立，
因此存在 UU 缺口并诱导 $\JacGap>0$。
将该见证代入定理~\ref{thm:open-phi-jacgap-uu} 的证明链即可得到 $\mathrm{claim}$。证毕。
\end{proof}

\subsubsection{$T_{\mathrm{NoRoundOnly}}$：仅回合逼近无法克服 UU}
\label{sec:certschema-t-noroundonly}

\begin{theorem}[证书规格：$T_{\mathrm{NoRoundOnly}}$（仅回合细分不够）]
\label{thm:certspec-t-noroundonly}
在“不分层”的仅回合逼近口径（定义~\ref{def:only-round-approximation}）下，
若 $\mathsf{Ver}_{T_{\mathrm{NoRoundOnly}}}(\mathcal{P},\pi_{\mathrm{NoRoundOnly}})=\top$，
则 $\mathrm{claim}$ 成立：仅靠回合细分（subdivision）无法把开放 UU 变为封闭系统（对应定理~\ref{thm:uu-open-subdivision-not-enough}）。
\end{theorem}

\begin{Certificate}[证书字段：$T_{\mathrm{NoRoundOnly}}$（引用开放性见证 + 细分不变性）]
\label{cert:certspec-t-noroundonly}
可取 $\pi_{\mathrm{NoRoundOnly}}$ 的 $\mathrm{assump}$ 引用一个开放性见证（证书 $\pi_{\mathrm{Open}}$），并登记：
\[
  \mathrm{oblig}:\ \text{回合细分不改变回路传输（定理~\ref{thm:subdivision-preserve-transport}）}.
\]
失败回退：触发分层构造（$T_{\mathrm{LayeringNecessary}}$）而非继续无穷细分。
\end{Certificate}

\begin{proof}[证明（谓词因果链）]
由开放性见证与细分不变性，直接得到“非平凡 holonomy 在细分下保持”，从而 UU 根源不被消除；
这正是定理~\ref{thm:uu-open-subdivision-not-enough} 的结论。证毕。
\end{proof}

\subsubsection{$T_{\mathrm{LayeringNecessary}}$：分层必要性（层内统一修复动作）}
\label{sec:certschema-t-layering}

\begin{theorem}[证书规格：$T_{\mathrm{LayeringNecessary}}$（分层必要性）]
\label{thm:certspec-t-layering}
若 $\mathsf{Ver}_{T_{\mathrm{LayeringNecessary}}}(\sigma,\pi_{\mathrm{Layer}})=\top$，
则 $\mathrm{claim}$ 成立：分层抬升把“暗处 holonomy”显式化为层切换，从而在层口径上恢复可证书化的统一修复条件（对应定理~\ref{thm:layering-complete-iff-trivial-holonomy}）。
\end{theorem}

\begin{Certificate}[证书字段：$T_{\mathrm{LayeringNecessary}}$]
\label{cert:certspec-t-layering}
输入为分层映射 $\sigma:X\to I$ 与扩展观测 $\tilde p(x)=(p(x),\sigma(x))$。
验证义务 $\mathrm{oblig}$ 要求“层完备性”成立（定义口径见定理~\ref{thm:layering-complete-iff-trivial-holonomy}）：
在扩展空间上可把非平凡 holonomy 杀死（或等价地：把不可区分内部态对分裂为不同层）。
失败回退：若某层仍出现“修复义务不兼容”，则扩张层指标或提升语义函子场（同伦塔增广）。
\end{Certificate}

\begin{proof}[证明（谓词因果链）]
由 $\mathsf{Ver}_{T_{\mathrm{LayeringNecessary}}}(\sigma,\pi_{\mathrm{Layer}})=\top$ 得层完备性成立，
于是可套用定理~\ref{thm:layering-complete-iff-trivial-holonomy} 得到结论 $\mathrm{claim}$。证毕。
\end{proof}

\subsubsection{$T_{\mathsf{Lag}}$：必然滞后同伦类（主 NN 策略的结构性失败）}
\label{sec:certschema-t-lag}

\begin{theorem}[证书规格：$T_{\mathsf{Lag}}$（必然滞后同伦类）]
\label{thm:certspec-t-lag}
若 $\mathsf{Ver}_{T_{\mathsf{Lag}}}(o,\pi_{\mathsf{Lag}})=\top$，
则 $\mathrm{claim}$ 成立：任意仅依赖 $o\in O$ 的主策略类在该口径下落入必然滞后同伦类 $\lbrack\mathsf{Lag}\rbrack$，
不能承担 UU 主环证明义务（对应定理~\ref{thm:lag-holonomy-to-lag-class}）。
\end{theorem}

\begin{Certificate}[证书字段：$T_{\mathsf{Lag}}$（不可区分 + 修复不可兼容）]
\label{cert:certspec-t-lag}
见证字段 $\mathrm{witness}$ 给出同语义内部态对 $(x,y)$ 与不兼容证书（定义~\ref{def:lag-uu-repair-incompatibility}）：
\[
  p(x)=p(y)=o,\qquad A^{\mathrm{rep}}_{\delta}(x)\cap A^{\mathrm{rep}}_{\delta}(y)=\varnothing.
\]
失败回退：撤销该策略类作为“保证者”的资格，转入审计包装（$T_{\mathrm{AuditedWrapper}}$）或分层（$T_{\mathrm{LayeringNecessary}}$）。
\end{Certificate}

\begin{proof}[证明（谓词因果链）]
该证书正是定理~\ref{thm:lag-holonomy-to-lag-class} 的前提见证；
由定理~\ref{thm:lag-holonomy-to-lag-class} 立即得到 $\mathrm{claim}$。证毕。
\end{proof}

\subsubsection{$T_{\mathrm{AuditedWrapper}}$：审计包装（学习器降级、正确性由 $V$ 承担）}
\label{sec:certschema-t-audited}

\begin{theorem}[证书规格：$T_{\mathrm{AuditedWrapper}}$（审计包装）]
\label{thm:certspec-t-auditedwrapper}
若 $\mathsf{Ver}_{T_{\mathrm{AuditedWrapper}}}(V,G_\theta,\pi_{\mathrm{Aud}})=\top$，
则 $\mathrm{claim}$ 成立：审计包装 $F^{\mathrm{aud}}_\theta=V\circ G_\theta$ 的对外正确性由 $V$ 的健全性与审计证书保障，
失败时 fail-closed 输出 $\bot$，从而不把系统拉回 $\lbrack\mathsf{Lag}\rbrack$ 的主逻辑（对应定理~\ref{thm:audited-wrapper-soundness} 与~\ref{thm:audited-wrapper-inductive-safety}）。
\end{theorem}

\begin{Certificate}[证书字段：$T_{\mathrm{AuditedWrapper}}$（规则版本 + 测试记录 + 裁决日志）]
\label{cert:certspec-t-auditedwrapper}
输入为候选生成器 $G_\theta$ 与确定性审计器 $V$（定义~\ref{def:audited-wrapper-generator-auditor}），
见证字段包含：$V$ 的规则库版本/哈希、白名单检查与复杂度阈值、编译/一致性检查、回归测试与变形测试报告等。
验证义务为：$V$ 健全（定义~\ref{def:audited-wrapper-soundness}）且 fail-closed（拒绝即 $\bot$）。
失败回退：若 $V$ 拒绝，则返回 $\bot$ 并触发修复塔（扩张语义函子场/分层/更保守策略）。
\end{Certificate}

\begin{proof}[证明（谓词因果链）]
由验证器通过得 $V$ 健全与 fail-closed 行为成立。
于是由定理~\ref{thm:audited-wrapper-soundness} 得单步正确性，
再由定理~\ref{thm:audited-wrapper-inductive-safety} 得多步闭环下合同保持。
从而 $\mathrm{claim}$ 成立。证毕。
\end{proof}

\subsection{统一运行期闭环接口模板（六定理串联）}
\label{sec:certschema-runtime-loop}

\begin{remark}[运行期闭环（观测$\to$候选$\to$审计$\to$执行/回退$\to$修复塔增广）]
将六条核心定理视为同一运行期闭环中的不同“证书分支”，可统一写成：
\begin{enumerate}
  \item 观测：得到法对象 $L_t$ 与回合语义 $o_t$；若已分层则 $\tilde o_t=(o_t,i_t)$；
  \item 候选：由规则库/ORL/学习器产生候选（学习器只做 $G_\theta$，不承担证明义务）；
  \item 审计：用确定性验证器裁决 $\mathcal{E}_{\mathrm{prot}}$ 与 $\Phi$ 的义务，或裁决语义规范形合同；
  \item 通过则执行；失败则 fail-closed 并产出见证（例如 UU 反例），触发修复塔增广（扩张/分层/升阶）。
\end{enumerate}
工程上对应“ok:true/ok:false”合同（见节~\ref{sec:certschema-json}），其证书载体即为 $\pi_T$ 的序列化。
\end{remark}

\section{工程/方案设计的开放博弈刻画：OHU 修复塔迁移与语义规范场基座}
\label{sec:eng-design-open-game-ohutower-semanticfield}

\begin{remark}[目标（按 PureMath + Vol1 口径证明三件事）]
本章在如下符号骨架下工作：
\begin{itemize}
  \item \textbf{PureMath}：GZ-OHU、$\JacGap$、$\GZLS$、Upgrade/Strictify、$\Phi=\mathcal{E}_{\mathrm{CH}}+\lambda\JacGap$；
  \item \textbf{Vol1}：$(\GZLS)_{\mathrm{RL}}$、$(\mathcal{L}_{\mathrm{env}},\mathcal{L}_{\mathrm{pol}},\mathcal{L}
    _{\mathrm{RL}})$、
  $\mathsf{Cert}_{\mathrm{pol}}$、$\mathcal{J}_{\mathrm{Jac}}$、证书检查--接受/拒绝循环。
\end{itemize}
给出一组“定义—定理—证明—推论—例子”的严格论述，证明：
\begin{enumerate}
  \item 工程/方案设计（以发动机、芯片为例）本质上是\textbf{开放系统博弈}；
  \item 其迭代改进可由\textbf{同伦修复塔（GZ-OHU/Upgrade-Strictify-Completion）}的构造模式直接迁移生成；
  \item \textbf{语义规范场的语义函子}为“数学符号 + 代码 + 规格文本”混合语义提供\textbf{正则基座}，
  从而使 $\Phi$/$\JacGap$/$\mathcal{E}_{\mathrm{prot}}$ 的证书化在工程上可复算、可审计、可回退。
\end{enumerate}
本章只使用前文已固定的定义与定理链条，并以“证书 + 谓词因果链”闭合推导。
\end{remark}

\subsection{预备：PureMath 与 Vol1 的可用结构件（只列可直接调用者）}
\label{sec:eng-design-prelim}

\begin{remark}[可直接调用的结构件（引用点）]
\begin{itemize}
  \item $\JacGap$ 作为 defect functional/证书接口：定义~\ref{def:jacobiator-gap-certificate}；
  \item Upgrade/Strictify 与 $\JacGap$-free 子范畴：定义~\ref{def:magic-wand-puremath-upgrade-strictify}；
  \item $\JacGap=0$ 的 strictify 分支：定理~\ref{thm:magic-wand-puremath-jacgap-vanishing}；
  \item OHU 迭代的证书保持与势能单调性：定理~\ref{thm:magic-wand-puremath-monotone}；
  \item $\mathsf{Ob}\Rightarrow \JacGap \Rightarrow \Phi$ 与修复塔：章节~\ref{sec:ob-jacgap-phi-repair-tower}；
  \item 审计包装（候选生成 + 确定性裁决；fail-closed）：定理~\ref{thm:audited-wrapper-soundness}；
  \item 证书检查--接受/拒绝循环的多步合规：定理~\ref{thm:audited-wrapper-inductive-safety} 与证书 schema（章节~\ref{sec:unified-certificate-schema}）。
\end{itemize}
下面把这些结构件对接到“发动机/芯片设计”口径。
\end{remark}

\subsection{形式逻辑：工程设计的 law-object 与开放系统博弈判据}
\label{sec:eng-design-formal}

\begin{definition}[工程设计的 law-object]
\label{def:eng-design-law-object}
一类工程设计（发动机/芯片/材料/控制等）在法空间中表示为
\[
  \mathcal{L}_{\mathrm{des}}\in(\GZLS)_{\mathrm{Eng}}\subseteq \GZLS,
\]
其中 $\mathcal{L}_{\mathrm{des}}$ 至少编码：
\begin{enumerate}
  \item 结构参数与几何/材料/工艺窗口/控制接口等（可执行部分）；
  \item 评估算子族（仿真、约束检查、验证、性能度量），记为可观测泛函族 $\mathcal{O}(\mathcal{L}_{\mathrm{des}})$；
  \item 证书谓词与证书函数子（安全、合规、可制造、可靠性等），抽象为
  \[
    \mathcal{E}_{\mathrm{prot}}(\mathcal{L}_{\mathrm{des}})\in\{\top,\bot\},
    \qquad
    \mathsf{Cert}(\mathcal{L}_{\mathrm{des}})\in\mathsf{CertSpace}.
  \]
\end{enumerate}
\end{definition}

\begin{definition}[工程设计的开放系统博弈（回合化 law-space 口径）]
\label{def:eng-design-open-game}
把设计过程写为回合博弈：设计者（Player-D）每回合选择设计动作 $u_k$，
环境/对手（Player-E）选择扰动 $e_k$。回合状态更新写作
\[
  \mathcal{L}_{k+1}=\mathcal{T}(\mathcal{L}_k,u_k,e_k),\qquad \mathcal{L}_k\in(\GZLS)_{\mathrm{Eng}}.
\]
称该设计博弈在 UU 口径下为\textbf{开放}，若对任意固定回合刻画与任意固定的预设后继族
$\widehat{\mathrm{Post}}$，总存在细化尺度出现不可闭合缺口（UU 见证）：
\[
  \exists(o,a,o'):\quad o'\in \mathrm{Post}(o,a)\ \wedge\ o'\notin \widehat{\mathrm{Post}}(o,a),
\]
等价地 $\mathsf{Ob}(o,a)\neq\varnothing$（定义~\ref{def:round-obstruction}）。
\end{definition}

\begin{assumption}[工程事实层前提（E1--E3）]
\label{asm:eng-design-open-facts}
对真实工程设计过程，假设：
\begin{enumerate}
  \item (E1) 设计者无法将所有未来约束（法规、工艺窗口、供应链、对抗者策略）在开始时一次性封闭为固定集合；
  \item (E2) 环境扰动 $e_k$ 的可实现集合不由设计者控制，且其引入可改变可行域与证书谓词 $\mathcal{E}_{\mathrm{prot}}$ 的判定边界；
  \item (E3) 决策必须在有限预算与有限观测下进行（因此“只靠回合逼近 + 事后拟合”无法承担 UU 首现点的证明义务；参见定理~\ref{thm:lag-holonomy-to-lag-class} 的口径）。
\end{enumerate}
\end{assumption}

\begin{theorem}[工程/方案设计的开放性：发动机/芯片设计是开放系统博弈]
\label{thm:eng-design-is-open-game}
在假设~\ref{asm:eng-design-open-facts} 下，工程设计过程满足定义~\ref{def:eng-design-open-game} 的开放判据，
因此工程/方案设计本质上是开放系统博弈（UU 会在某些细化尺度出现）。
\end{theorem}

\begin{Certificate}[证书（工程版）：UU 见证记录/证书边界变化记录]
\label{cert:eng-design-open-witness}
证书可取为一次或多次 UU 见证记录：
\[
  \pi_{\mathrm{open}}
  :=(o,a,o')\ \text{使}\ o'\in \mathrm{Post}(o,a)\ \wedge\ o'\notin\widehat{\mathrm{Post}}(o,a),
\]
或等价地“证书失效/证书边界变化”的审计记录（某次外源变化使先前通过的 $\mathcal{E}_{\mathrm{prot}}$ 判定失效）。
谓词链登记为：
\[
\begin{aligned}
  O_1 &: (E1)\wedge(E2)\Rightarrow \exists(o,a,o')\ \text{出现新后继/新规则边界},\\
  O_2 &: O_1\Rightarrow \mathsf{Ob}(o,a)\neq\varnothing,\\
  O_3 &: O_2\Rightarrow \JacGap>0\ \text{（引理~\ref{lem:uu-witness-implies-jacgap-positive}）},\\
  O_4 &: O_2\Rightarrow \text{开放判据成立（定义~\ref{def:eng-design-open-game}）}.
\end{aligned}
\]
\end{Certificate}

\begin{proof}[证明（谓词因果链）]
由 (E1)(E2)：“规则集合/可行域/证书边界”会被外源重写，因此在任一固定建模语言（固定 $\widehat{\mathrm{Post}}$）下，
存在某次细化尺度出现真实后继 $o'$ 不属于预设后继集合，得到 $O_1$，进而 $O_2$。
由引理~\ref{lem:uu-witness-implies-jacgap-positive} 得 $O_3$；由定义~\ref{def:eng-design-open-game} 得 $O_4$。
因此工程设计满足开放判据，是开放系统博弈。证毕。
\end{proof}

\begin{remark}[例子 1：发动机方案设计的开放性来源（不涉及细节）]
令 $\mathcal{L}_{\mathrm{des}}$ 编码结构/材料/燃烧控制/冷却/润滑/排放后处理等。
典型开放性来源包括：制造偏差与批次漂移、工况谱变化、法规/测试工况变化、供应链约束变化等。
这些变化会直接改写“可行域与证书边界”，从而产生 UU 见证，符合定理~\ref{thm:eng-design-is-open-game} 的结构前提。
\end{remark}

\begin{remark}[例子 2：芯片方案设计的开放性来源（不涉及细节）]
令 $\mathcal{L}_{\mathrm{des}}$ 编码微结构/版图、PPA 目标、时序收敛、可靠性、安全与供应链约束等。
典型开放性来源包括：工艺窗口与良率漂移、EDA 假设变更、威胁模型演化、制裁与供应链规则变化等。
同样会导致 UU 缺口（$\mathsf{Ob}\neq\varnothing$），因此芯片设计亦为开放系统博弈。
\end{remark}

\subsection{定理：OHU/Upgrade-Strictify-Completion 模式对工程设计的可迁移生成}
\label{sec:eng-design-ohutower}

\begin{definition}[工程设计的缺陷势与证书检查--接受/拒绝循环（Vol1 同构）]
\label{def:eng-design-phi-accept-reject}
定义工程设计的合并缺陷势（与章节~\ref{sec:ob-jacgap-phi-repair-tower} 一致）：
\[
  \Phi(\mathcal{L}):=\mathcal{E}_{\mathrm{CH}}(\mathcal{L})+\lambda\,\JacGap(\mathcal{L}),\qquad \lambda>0.
\]
并定义一次“候选更新 + 证书裁决”的循环（Vol1 accept/reject 口径）：
从 $\mathcal{L}^{(k)}$ 生成候选 $\widetilde{\mathcal{L}}^{(k+1)}$，
计算证书接口（例如 $\mathsf{Cert}(\widetilde{\mathcal{L}}^{(k+1)})$ 与 gap 证书），若
\[
  \mathcal{E}_{\mathrm{prot}}(\widetilde{\mathcal{L}}^{(k+1)})=\top
  \ \wedge\
  \Phi(\widetilde{\mathcal{L}}^{(k+1)})\le \Phi(\mathcal{L}^{(k)})-\delta,
\]
则接受并令 $\mathcal{L}^{(k+1)}:=\widetilde{\mathcal{L}}^{(k+1)}$；否则拒绝/投影回 certified region（fail-closed），并触发修复动作族扩张（同伦修复塔）。
\end{definition}

\begin{theorem}[公理降级：证书检查--接受/拒绝循环的多步合规性（数学归纳法证书）]
\label{thm:eng-design-accept-reject-induction}
在定义~\ref{def:eng-design-phi-accept-reject} 的循环口径下，
若每一步被接受的更新都满足“证书保持 + $\Phi$ 至少下降 $\delta>0$”，
则任意长度为 $n$ 的接受链都满足：全程 $\mathcal{E}_{\mathrm{prot}}=\top$，且 $\Phi$ 单调下降并在有限步内给出终止上界。
\end{theorem}

\begin{Certificate}[数学归纳法证书：按接受步数 $n$ 的保持性与单调性]
\label{cert:eng-design-accept-reject-induction}
令谓词 $P(n)$ 为：
“对任意接受链长为 $n$ 的序列，$\mathcal{E}_{\mathrm{prot}}$ 全程保持且 $\Phi$ 每步至少下降 $\delta$。”
证书登记为：
\[
  \pi_{\mathrm{ind}}=\bigl(P(0),\ \forall n(P(n)\Rightarrow P(n+1))\bigr),
\]
其中归纳步由定义~\ref{def:eng-design-phi-accept-reject} 的接受条件直接给出。
\end{Certificate}

\begin{proof}[证明（数学归纳法证书；谓词因果链）]
基步 $P(0)$ 平凡成立。
归纳步：若 $P(n)$ 成立，则第 $n+1$ 次接受由定义~\ref{def:eng-design-phi-accept-reject} 的接受条件给出
$\mathcal{E}_{\mathrm{prot}}=\top$ 且 $\Phi$ 至少下降 $\delta$，因此 $P(n+1)$ 成立。
故对所有 $n$ 成立。证毕。
\end{proof}

\begin{theorem}[工程设计的同伦修复塔：由 OHU 单调性与 Strictify 分支普适迁移]
\label{thm:eng-design-ohutower-migration}
设 $\mathcal{L}_0\in(\GZLS)_{\mathrm{Eng}}$ 为任一初始设计 law-object，并可计算/证书化 $\JacGap(\mathcal{L})$（定义~\ref{def:jacobiator-gap-certificate}）。
若设计迭代满足 OHU 口径的“证书保持 + 势能单调性”（可直接调用定理~\ref{thm:magic-wand-puremath-monotone} 的结构），
则存在一条同伦修复塔（admissible extensions 链）
\[
  \mathcal{L}_0\hookrightarrow \mathcal{L}_1\hookrightarrow\cdots,
\]
使得在 $\mathcal{E}_{\mathrm{prot}}$ 保持下 $\Phi$ 单调不增，并在出现 $\JacGap=0$ 的分支上可调用 Strictify 得到 strict 代表（定理~\ref{thm:magic-wand-puremath-jacgap-vanishing}）。
因此 GZ-OHU/Upgrade-Strictify 的构造模式可\textbf{直接迁移生成}工程设计的证书化修复流程。
\end{theorem}

\begin{Certificate}[证书：OHU 单调性 + JacGap-free strictify 分支]
\label{cert:eng-design-ohutower-migration}
证书登记谓词链：
\[
\begin{aligned}
  H_1 &: \text{存在 OHU 口径迭代链并满足证书保持 + 单调性（定理~\ref{thm:magic-wand-puremath-monotone} 的前提）},\\
  H_2 &: H_1\Rightarrow \forall n,\ \mathcal{E}_{\mathrm{prot}}(\mathcal{L}^{(n)})=\top\ \wedge\ \Phi(\mathcal{L}^{(n)})
    \le \Phi(\mathcal{L}^{(0)}),\\
  H_3 &: \JacGap(\mathcal{L}^{(n)})=0\Rightarrow \text{strictify 可用（定理~\ref{thm:magic-wand-puremath-jacgap-vanishing}）}.
\end{aligned}
\]
\end{Certificate}

\begin{proof}[证明（谓词因果链）]
由 $H_1$ 直接套用定理~\ref{thm:magic-wand-puremath-monotone} 得 $H_2$，即迭代链在证书保持下势能单调不增。
若某步进入 $\JacGap=0$ 分支，则由定理~\ref{thm:magic-wand-puremath-jacgap-vanishing} 得 strictify 可用并给出 strict 代表，得到 $H_3$。
因此可把工程设计迭代组织为同伦修复塔，并在 JacGap-free 分支上严格化。证毕。
\end{proof}

\begin{corollary}[可迁移生成（普适性口径）]
\label{cor:eng-design-migration}
只要把不同工程域对象编码为 $(\GZLS)_{\mathrm{Eng}}$ 中的 law-object，并装备 $\JacGap$ 与 $\Phi$ 的证书接口，
则“证书裁决 + 同伦修复塔 +（JacGap-free 时 strictify）”的形状保持不变；
差异仅体现在具体的算子族与证书内容上。
\end{corollary}

\begin{Certificate}[证书：迁移模板仅依赖“law-object + $\JacGap/\Phi$ 接口”]
\label{cert:eng-design-migration}
证书登记谓词链：
\[
\begin{aligned}
  M_1 &: \mathcal{L}\in(\GZLS)_{\mathrm{Eng}}\ \wedge\ \text{可证书化 }\JacGap,\Phi\Rightarrow \text{满足定理~\ref{thm:eng-design-ohutower-migration} 的输入接口},\\
  M_2 &: \text{定理~\ref{thm:eng-design-ohutower-migration}}\Rightarrow \text{存在同伦修复塔并具有 JacGap-free strictify 分支},\\
  M_3 &: M_2\Rightarrow \text{模板形状与具体工程域无关（仅证书内容/算子族不同）}.
\end{aligned}
\]
\end{Certificate}

\begin{proof}[证明（谓词因果链）]
由定理~\ref{thm:eng-design-ohutower-migration} 给出的迁移生成模板，
其输入只依赖“可证书化的 $\JacGap$ 与 OHU 口径迭代单调性”，与具体工程域无关；
因此模板可复用，得到推论结论。证毕。
\end{proof}

\subsection{定理：语义规范场作为“数学符号+代码”语义描述的正则基座}
\label{sec:eng-design-semantic-base}

\begin{definition}[语义表达空间、语义解释与语义等价]
\label{def:semantic-expr-equivalence}
令 $\mathsf{Expr}$ 表示工程语义表达的空间（数学式、程序、配置、规格文本的统一抽象），
令 $\mathsf{Sem}$ 表示其语义对象空间（可执行行为、可验证约束集、可观测输出等）。
给定语义解释映射
\[
  \mathsf{Sem}(\cdot):\mathsf{Expr}\to \mathsf{Sem},
\]
定义语义等价为
\[
  e_1\sim_{\mathrm{sem}} e_2\ \Longleftrightarrow\ \mathsf{Sem}(e_1)=\mathsf{Sem}(e_2).
\]
\end{definition}

\begin{definition}[语义规范形与语义函子（规范化函数）]
\label{def:semantic-normal-form-functor}
令 $\mathsf{NF}\subseteq\mathsf{Expr}$ 为正则形集合（规范形）。
称映射
\[
  F:\mathsf{Expr}\to \mathsf{NF}\sqcup\{\bot\}
\]
为语义函子（规范化函子），若 $F(e)=n\neq\bot$ 时代表“规范化成功并返回规范形代表 $n$”，
而 $F(e)=\bot$ 代表“失败显式（fail-closed）”。
语义规范场是该类函子的族（对齐定义~\ref{def:audited-wrapper-semantic-functor-field} 的口径）。
\end{definition}

\begin{theorem}[语义规范场的基座性：为混合语义提供可复算/可审计的规范代表]
\label{thm:semantic-normalization-base}
设存在确定性审计器 $V$ 与候选生成器 $G$ 使审计包装
$F^{\mathrm{aud}}:=V\circ G$ 满足：
\begin{enumerate}
  \item （语义保持）若 $F^{\mathrm{aud}}(e)=n\neq\bot$，则 $e\sim_{\mathrm{sem}} n$；
  \item （基座性/一致性）若 $e_1\sim_{\mathrm{sem}} e_2$ 且二者均可规范化为非 $\bot$，则 $F^{\mathrm{aud}}(e_1)=F^{\mathrm{aud}}(e_2)$
  （或至少落在同一规范形等价类）；
  \item （失败显式）不可规范化时返回 $\bot$，并将该失败作为“需要扩张语义坐标/修复规则”的见证输入修复塔。
\end{enumerate}
则 $\mathsf{NF}$ 在语义等价类商空间意义下构成正则基座：它为“数学符号 + 代码”的混合语义提供可比较的规范代表，
从而使 $\Phi$、$\JacGap$ 与 $\mathcal{E}_{\mathrm{prot}}$ 的证书化判定可稳定复算并进入工程闭环。
\end{theorem}

\begin{Certificate}[证书：语义保持 + 一致性 + 失败显式]
\label{cert:semantic-normalization-base}
证书登记谓词链：
\[
\begin{aligned}
  S_1 &: F^{\mathrm{aud}}(e)=n\neq\bot\Rightarrow e\sim_{\mathrm{sem}} n,\\
  S_2 &: e_1\sim_{\mathrm{sem}} e_2\ \wedge\ F^{\mathrm{aud}}(e_1)\neq\bot\ \wedge\ F^{\mathrm{aud}}(e_2)
    \neq\bot\Rightarrow F^{\mathrm{aud}}(e_1)=F^{\mathrm{aud}}(e_2),\\
  S_3 &: F^{\mathrm{aud}}(e)=\bot\Rightarrow \text{失败显式并触发语义扩张/修复塔}.
\end{aligned}
\]
\end{Certificate}

\begin{proof}[证明（谓词因果链；要点）]
由 $S_1$，任何被接受的规范形代表与原表达语义等价，因此可在 $\mathsf{NF}$ 上进行“比较/审计/复算”而不依赖表面表达差异。
由 $S_2$，语义等价类被映射到同一规范代表（或同一规范类），因此 $\mathsf{NF}$ 提供 canonical representative 的基座。
由 $S_3$，不确定语义不会被悄然吞掉；失败被显式记录为见证并触发“扩张语义坐标/规则库”的修复动作，
这与 PureMath 的“缺陷提示缺失同伦数据并需扩张”方法论同构。
因此 $\mathsf{NF}$ 构成混合语义的正则基座，定理成立。证毕。
\end{proof}

\begin{corollary}[发动机/芯片设计的统一法空间工作流（可执行闭环）]
\label{cor:eng-design-unified-workflow}
发动机或芯片的设计流程可统一写为：
\begin{enumerate}
  \item 用语义规范场把“数学 + 代码 + 规格”归一到 $\mathsf{NF}$（定理~\ref{thm:semantic-normalization-base}），形成可审计语义；
  \item 在 $(\GZLS)_{\mathrm{Eng}}$ 中形成设计 law-object $\mathcal{L}_{\mathrm{des}}$（定义~\ref{def:eng-design-law-object}）；
  \item 将迭代视为开放博弈 $\mathcal{L}_{k+1}=\mathcal{T}(\mathcal{L}_k,u_k,e_k)$（定理~\ref{thm:eng-design-is-open-game}）；
  \item 以 $\Phi=\mathcal{E}_{\mathrm{CH}}+\lambda\JacGap$ 为缺陷势，并用证书检查--接受/拒绝机制保证每一步处于 certified region（定理~\ref{thm:eng-design-accept-reject-induction}）；
  \item 当 $\JacGap>0$ 或证书破坏时，按 OHU 原则做同伦升级与扩张，构造同伦修复塔并在 JacGap-free 分支上 strictify（定理~\ref{thm:eng-design-ohutower-migration}）。
\end{enumerate}
\end{corollary}

\begin{Certificate}[证书：工作流五步由对应定理链逐项支撑]
\label{cert:eng-design-unified-workflow}
证书登记谓词链：
\[
\begin{aligned}
  U_1 &: \text{定理~\ref{thm:semantic-normalization-base}}\Rightarrow \text{(1) 语义规范场给出可复算/可审计规范代表},\\
  U_2 &: \text{定义~\ref{def:eng-design-law-object}}\Rightarrow \text{(2) 设计对象可编码为 }(\GZLS)_{\mathrm{Eng}}\text{ 的
    law-object},\\
  U_3 &: \text{定理~\ref{thm:eng-design-is-open-game}}\Rightarrow \text{(3) 迭代可写为开放博弈回合更新},\\
  U_4 &: \text{定理~\ref{thm:eng-design-accept-reject-induction}}\Rightarrow \text{(4) 证书检查--接受/拒绝保证多步合规},\\
  U_5 &: \text{定理~\ref{thm:eng-design-ohutower-migration}}\Rightarrow \text{(5) OHU 修复塔迁移生成与 strictify 分支}.
\end{aligned}
\]
\end{Certificate}

\begin{proof}[证明（谓词因果链）]
由定理~\ref{thm:semantic-normalization-base} 得第(1)步的“可复算/可审计语义基座”；
由定理~\ref{thm:eng-design-is-open-game} 得第(3)步开放博弈口径成立；
由定理~\ref{thm:eng-design-accept-reject-induction} 得第(4)步多步合规性；
由定理~\ref{thm:eng-design-ohutower-migration} 得第(5)步同伦修复塔迁移生成与 strictify 分支。
合并即得推论结论。证毕。
\end{proof}

\section{工程域的最小算子丛模板：发动机/芯片与三证书接口闭环}
\label{sec:eng-minimal-operator-bundle-templates}

\subsection{语义描述正则基座：词包表达、规范形与审计包装}
\label{sec:eng-minimal-semantic-base}

\begin{remark}[对齐：$\mathsf{WP}\to \mathsf{NF}\sqcup\{\bot\}$ 与审计包装（fail-closed）]
本章把“数学符号 + 代码 + 规格文本”的混合表达统一放入词包表达层（工作笔记口径）：
\[
  \mathsf{WP}:=\Omega^{(1)}(L),\qquad \mathsf{NF}\subseteq\mathsf{WP},\qquad
  F:\mathsf{WP}\to \mathsf{NF}\sqcup\{\bot\}.
\]
其中 $L$ 表示承载混合语义的语言/词典，$\mathsf{NF}$ 为规范形集合，$\bot$ 表示失败显式（fail-closed）。
语义规范场是函子族 $\mathcal{F}=\{F_i\}_{i\in I}$，工程实现建议统一套两段式审计包装：
\[
  F_i^{\mathrm{aud}}:=V_i\circ G_i,
\]
其中 $G_i$ 生成候选（可由规则引擎/NN/LLM 等实现），$V_i$ 为确定性审计器并裁决接受/拒绝；
拒绝必须返回 $\bot$（定理~\ref{thm:audited-wrapper-soundness}）。
该口径与定义~\ref{def:semantic-normal-form-functor} 的 $\mathsf{Expr}\to\mathsf{NF}\sqcup\{\bot\}$ 等价，只是把表达域显式记为 $\mathsf{WP}$。
\end{remark}

\begin{lemma}[二元数据等价：部分规范化函数 $\Leftrightarrow$（判定 $\chi$，规范化 $\mathrm{nf}$）]
\label{lem:eng-minimal-f-equivalence}
对任意 $F:\mathsf{WP}\to \mathsf{NF}\sqcup\{\bot\}$，存在谓词
\[
  \chi:\mathsf{WP}\to\{\mathrm{true},\mathrm{false}\}
\]
与部分映射
\[
  \mathrm{nf}:\{x\in\mathsf{WP}:\chi(x)=\mathrm{true}\}\to \mathsf{NF},
\]
使得对任意 $x\in\mathsf{WP}$ 都有
\[
  F(x)=
  \begin{cases}
    \mathrm{nf}(x),&\chi(x)=\mathrm{true},\\
    \bot,&\chi(x)=\mathrm{false}.
  \end{cases}
\]
反之，任意这样的 $(\chi,\mathrm{nf})$ 唯一确定一个 $F$。
\end{lemma}

\begin{Certificate}[证书：由 $F$ 构造 $(\chi,\mathrm{nf})$，并可回构 $F$]
\label{cert:eng-minimal-f-equivalence}
证书登记构造与谓词链：
\[
\begin{aligned}
  C_1 &: \chi(x):=\bigl(F(x)\neq\bot\bigr),\\
  C_2 &: \mathrm{nf}:=\left.F\right|_{\{x:\chi(x)=\mathrm{true}\}},\\
  C_3 &: \text{由 }(C_1)(C_2)\Rightarrow
  F(x)=\begin{cases}\mathrm{nf}(x),&\chi(x)=\mathrm{true},\\ \bot,&\chi(x)=\mathrm{false};\end{cases}\\
  C_4 &: \text{反向：给定 }(\chi,\mathrm{nf})\Rightarrow \text{可按同一分段式定义 }F.
\end{aligned}
\]
\end{Certificate}

\begin{proof}[证明（谓词因果链）]
由 $(C_1)$ 定义 $\chi$，由 $(C_2)$ 将 $F$ 在“成功域”上限制为 $\mathrm{nf}$；
则当 $\chi(x)=\mathrm{true}$ 时必有 $F(x)\neq\bot$ 且 $\mathrm{nf}(x)=F(x)$，当 $\chi(x)=\mathrm{false}$ 时必有 $F(x)=\bot$，得到
  $(C_3)$。
反向构造由分段定义直接给出 $(C_4)$；故等价成立。证毕。
\end{proof}

\subsection{工程算子丛的最小结构：扇区枚举、耦合联络与 law-object}
\label{sec:eng-minimal-operator-bundle}

\begin{definition}[工程算子丛的最小结构（扇区分解 + 耦合联络）]
\label{def:eng-minimal-operator-bundle}
把工程设计域抽象为“扇区枚举 + 耦合联络”的算子丛：
\begin{enumerate}
  \item 扇区集合（最小分解）：
  \[
    S=\{s_1,\dots,s_m\};
  \]
  \item 每个扇区 $s\in S$ 配一组状态/参数与算子结构（可取幺半群/幺半范畴口径）：
  \[
    \mathcal{O}_s=(O_s,\circ,\mathrm{id});
  \]
  \item 扇区耦合由“联络/耦合算子”描述（一般非交换）：
  \[
    \nabla_{s\to t}\in \mathrm{Hom}(O_s,O_t),\qquad s,t\in S.
  \]
\end{enumerate}
\end{definition}

\begin{definition}[工程设计 law-object 的最小模板（算子丛 + 三证书接口）]
\label{def:eng-minimal-law-object-template}
在语义规范场（小节~\ref{sec:eng-minimal-semantic-base}）给定下，
称
\[
  \mathcal{L}_{\mathrm{des}}
  :=
  \Bigl(
  \mathsf{NF},\ \{\mathcal{O}_s\}_{s\in S},\ \{\nabla_{s\to t}\}_{s,t\in S},\
  \mathcal{E}_{\mathrm{prot}},\ \mathcal{E}_{\mathrm{CH}},\ \JacGap
  \Bigr)\in(\GZLS)_{\mathrm{Eng}}\subseteq \GZLS
\]
为工程设计的最小 law-object 模板。
\end{definition}

\begin{remark}[非交换性（工程默认）]
扇区耦合的“先后顺序不同 $\Rightarrow$ 结果不同”是常态：例如“先改几何再标定”与“先标定再改几何”一般不等价，
因此算子序列应按非交换口径登记（参见第~\ref{sec:construction-operators} 节的顺序敏感性）。
\end{remark}

\subsection{三类最小证书接口：$\mathcal{E}_{\mathrm{prot}}$、$\mathcal{E}_{\mathrm{CH}}$ 与 $\JacGap$}
\label{sec:eng-minimal-three-cert-interfaces}

\begin{definition}[受保护一致性证书接口 $\mathcal{E}_{\mathrm{prot}}$（硬红线裁决器）]
\label{def:eng-minimal-eprot}
受保护一致性证书用于裁决“不可触碰的硬约束/底线”：
\[
  \mathcal{E}_{\mathrm{prot}}(\mathcal{L})\in\{\top,\bot\}.
\]
建议其最小证书对象固定为四元组：
\[
  \mathsf{Cert}_{\mathrm{prot}}(\mathcal{L})=
  (\mathrm{RuleID},\ m,\ \mathrm{EvidenceSet},\ \mathrm{VerifierHash}),
\]
其中 $\mathrm{RuleID}$ 为规则集合版本，$m\in\RR^k$ 为裕度向量（每条硬约束的 margin），
$\mathrm{EvidenceSet}$ 为可复算证据集合（角点/测试向量/仿真摘要/seed 等），$\mathrm{VerifierHash}$ 为验证器版本哈希。
\end{definition}

\begin{theorem}[受保护一致性判定的最小口径（裕度向量）]
\label{thm:eng-minimal-eprot-margin}
在定义~\ref{def:eng-minimal-eprot} 的证书对象口径下，可取裁决规则：
\[
  \mathcal{E}_{\mathrm{prot}}(\mathcal{L})=\top
  \iff
  \min(m)\ge 0.
\]
若 $\mathcal{E}_{\mathrm{prot}}(\mathcal{L})=\bot$，则失败输出必须携带可复算反例见证
\[
  \mathsf{Witness}_{\mathrm{prot}}=(\mathrm{RuleID}^{\mathrm{viol}},\ \mathrm{CE},\ \mathrm{Seed}),
\]
用于触发修复塔动作选择与复验。
\end{theorem}

\begin{Certificate}[证书：$\mathsf{Cert}_{\mathrm{prot}}$ 四元组与反例见证（fail-closed）]
\label{cert:eng-minimal-eprot-margin}
证书登记：
\[
  \mathsf{Cert}_{\mathrm{prot}}(\mathcal{L})=
  (\mathrm{RuleID},\ m,\ \mathrm{EvidenceSet},\ \mathrm{VerifierHash}).
\]
并在 $\min(m)<0$ 时，取 $i^\star\in\arg\min_i m_i$，令
\[
  \mathrm{RuleID}^{\mathrm{viol}}:=\mathrm{RuleID}[i^\star],
  \qquad
  (\mathrm{CE},\mathrm{Seed})\in\mathrm{EvidenceSet}[i^\star],
\]
得到反例见证 $\mathsf{Witness}_{\mathrm{prot}}$。
谓词链为：
\[
\begin{aligned}
  P_1 &: \min(m)\ge 0\Rightarrow \forall i,\ m_i\ge 0,\\
  P_2 &: \min(m)<0\Rightarrow \exists i^\star,\ m_{i^\star}<0\Rightarrow \exists\,\mathsf{Witness}_{\mathrm{prot}},\\
  P_3 &: (P_1)\Rightarrow \mathcal{E}_{\mathrm{prot}}=\top,\qquad (P_2)\Rightarrow \mathcal{E}_{\mathrm{prot}}=\bot.
\end{aligned}
\]
\end{Certificate}

\begin{proof}[证明（谓词因果链）]
由 $P_1$：所有 margin 非负则硬约束全通过，因此可裁决 $\mathcal{E}_{\mathrm{prot}}=\top$。
由 $P_2$：最小 margin 为负意味着存在被违反规则索引 $i^\star$，证据集提供可复算的反例/seed，故必须输出见证并裁决为 $\bot$（fail-closed）。
合并得 $P_3$，定理成立。证毕。
\end{proof}

\begin{definition}[已知约束残差接口 $\mathcal{E}_{\mathrm{CH}}$（KU 方向主优化量）]
\label{def:eng-minimal-ech}
已知约束残差用于量化“在已知模型语言内离目标/约束还有多远”：
\[
  \mathcal{E}_{\mathrm{CH}}(\mathcal{L})\in\RR_{\ge 0}.
\]
建议其最小证书对象固定为
\[
  \mathsf{Cert}_{\mathrm{CH}}(\mathcal{L})=(r,\ W,\ \mathrm{EvidenceSet},\ \mathrm{VerifierHash}),
\]
其中 $r\in\RR^d$ 为残差向量，$W\in\RR^d_{\ge 0}$ 为权重/聚合规则。
最小标量化可取
\[
  \mathcal{E}_{\mathrm{CH}}(\mathcal{L}):=\lVert W\odot r\rVert_1
  \quad(\text{亦可按域改为其它聚合范数}).
\]
\end{definition}

\begin{lemma}[$\mathcal{E}_{\mathrm{CH}}$ 的非负性与可复算性（由证据与工具链锁定）]
\label{lem:eng-minimal-ech-nonneg}
在定义~\ref{def:eng-minimal-ech} 下，恒有 $\mathcal{E}_{\mathrm{CH}}(\mathcal{L})\ge 0$；
并且当 $\mathrm{EvidenceSet}$ 与 $\mathrm{VerifierHash}$ 固定时，$\mathcal{E}_{\mathrm{CH}}$ 的计算具有可复算/可审计边界。
\end{lemma}

\begin{Certificate}[证书：范数非负 + 版本哈希锁定]
\label{cert:eng-minimal-ech-nonneg}
证书登记谓词链：
\[
\begin{aligned}
  R_1 &: W\in\RR^d_{\ge 0}\Rightarrow \lVert W\odot r\rVert_1\ge 0,\\
  R_2 &: (\mathrm{EvidenceSet},\mathrm{VerifierHash})\ \text{固定}\Rightarrow (r,W)\ \text{可复算},\\
  R_3 &: (R_1)\wedge(R_2)\Rightarrow \mathcal{E}_{\mathrm{CH}}\ \text{可复算且非负}.
\end{aligned}
\]
\end{Certificate}

\begin{proof}[证明（谓词因果链）]
由 $R_1$：$L^1$-范数非负且 $W\ge 0$，故 $\mathcal{E}_{\mathrm{CH}}\ge 0$。
由 $R_2$：证据集固定工况/角点/脚本输入，工具链哈希锁定计算语义，因此 $(r,W)$ 的生成与聚合可复算。
合并得 $R_3$，引理成立。证毕。
\end{proof}

\begin{definition}[结构缺口证书 $\JacGap$ 的 UU 见证接口（闭包被击穿）]
\label{def:eng-minimal-jacgap-witness}
把“预设后继族”记为 $\widehat{\mathrm{Post}}$，真实后继为 $\mathrm{Post}$（定义~\ref{def:round-post-model} 的口径）。
UU 缺口集合为
\[
  \mathsf{Ob}(o,a)=\mathrm{Post}(o,a)\setminus \widehat{\mathrm{Post}}(o,a)
  \quad(\text{定义~\ref{def:round-obstruction}}).
\]
最小结构缺口证书（UU 见证）可取为三元组
\[
  \mathsf{Witness}_{\mathrm{UU}}=(o,a,o')
  \quad\text{满足}\quad
  o'\in \mathrm{Post}(o,a)\ \wedge\ o'\notin \widehat{\mathrm{Post}}(o,a).
\]
其量化缺口大小由 $\JacGap$ 给出（定义~\ref{def:jacgap-from-obstruction}；测度 $\mu$ 见定义~\ref{def:ob-measure-mu}）。
\end{definition}

\begin{proposition}[UU 见证 $\Rightarrow$ $\JacGap>0$（最小结构缺陷证书）]
\label{prop:eng-minimal-uu-witness-implies-jacgap}
若存在 UU 见证 $\mathsf{Witness}_{\mathrm{UU}}=(o,a,o')$，则存在分割口径使 $\JacGap^{(\mathcal{P})}>0$；
在将分割视为工程评测尺度的情况下，可视为“出现结构缺口，需触发 UU 修复动作”。
\end{proposition}

\begin{Certificate}[证书：UU 见证 $\Rightarrow$ 缺口非空 $\Rightarrow$ $\JacGap>0$]
\label{cert:eng-minimal-uu-witness-implies-jacgap}
证书登记谓词链（引用闭合）：
\[
\begin{aligned}
  J_1 &: \mathsf{Witness}_{\mathrm{UU}}=(o,a,o')\Rightarrow \mathsf{Ob}(o,a)\neq\varnothing,\\
  J_2 &: J_1\Rightarrow \JacGap^{(\mathcal{P})}>0\ \text{（引理~\ref{lem:uu-witness-implies-jacgap-positive}）}.
\end{aligned}
\]
\end{Certificate}

\begin{proof}[证明（谓词因果链）]
由定义~\ref{def:eng-minimal-jacgap-witness}，三元组见证直接给出 $o'\in \mathsf{Ob}(o,a)$，故 $J_1$ 成立。
由引理~\ref{lem:uu-witness-implies-jacgap-positive} 得 $J_2$，命题成立。证毕。
\end{proof}

\subsection{合并势能与闭环规则：失败$\to$见证$\to$同伦修复塔动作}
\label{sec:eng-minimal-loop-rule}

\begin{definition}[合并缺陷势与最小接受/拒绝规则]
\label{def:eng-minimal-phi-loop}
定义合并缺陷势（与定义~\ref{def:eng-design-phi-accept-reject} 一致）：
\[
  \Phi(\mathcal{L})
  :=
  \mathcal{E}_{\mathrm{CH}}(\mathcal{L})+\lambda\,\JacGap(\mathcal{L}),
  \qquad \lambda>0.
\]
运行期最小裁决规则（fail-closed）：
\begin{enumerate}
  \item \textbf{接受}：$\mathcal{E}_{\mathrm{prot}}(\mathcal{L})=\top$ 且 $\Phi(\mathcal{L})$ 进入阈值域或给出严格下降；
  \item \textbf{拒绝/修复}：任一失败则输出见证（$\mathsf{Witness}_{\mathrm{prot}}$ 或 $\mathsf{Witness}_{\mathrm{UU}}$），并进入同伦修复塔选择动作。
\end{enumerate}
\end{definition}

\begin{theorem}[公理降级：fail-closed 闭环的多步合规性（数学归纳法证书）]
\label{thm:eng-minimal-loop-induction}
在定义~\ref{def:eng-minimal-phi-loop} 的闭环规则下，
若每次被接受的更新都满足“证书保持 + 势能严格下降”：
\[
  \mathcal{E}_{\mathrm{prot}}(\mathcal{L}_{k+1})=\top,\qquad
  \Phi(\mathcal{L}_{k+1})\le \Phi(\mathcal{L}_k)-\delta,\ \delta>0,
\]
则对任意接受步数 $n$，接受链全程保持 $\mathcal{E}_{\mathrm{prot}}=\top$，且 $\Phi$ 至少下降 $n\delta$；
同时任一拒绝事件都必须对应一个可复验见证（prot 反例或 UU 反例）。
\end{theorem}

\begin{Certificate}[数学归纳法证书：按接受步数 $n$ 的保持性与下降性]
\label{cert:eng-minimal-loop-induction}
令谓词 $P(n)$ 为：“任意接受链长为 $n$ 的序列，全程 $\mathcal{E}_{\mathrm{prot}}=\top$ 且 $\Phi$ 每步至少下降 $\delta$。”
证书登记为：
\[
  P(0)\ \text{成立},\qquad
  \forall n\ \bigl(P(n)\Rightarrow P(n+1)\bigr),
\]
其中归纳步由接受条件直接给出。
拒绝事件的见证要求由定义~\ref{def:eng-minimal-eprot} 与定义~\ref{def:eng-minimal-jacgap-witness} 的“失败必须输出见证”条款给出。
\end{Certificate}

\begin{proof}[证明（数学归纳法证书；谓词因果链）]
基步 $P(0)$ 平凡成立。
归纳步：假设 $P(n)$ 成立，则第 $n+1$ 次接受由前提给出 $\mathcal{E}_{\mathrm{prot}}(\mathcal{L}_{n+1})=\top$ 且 $\Phi$ 至少下降 $\delta$，故 $P(n+1)
  $ 成立。
因此对所有 $n$ 成立，得到“多步合规 + 势能下降”结论。
对拒绝事件：由定义~\ref{def:eng-minimal-eprot} 与定义~\ref{def:eng-minimal-jacgap-witness}，任一失败必须以 $\mathsf{Witness}_{\mathrm{prot}}
  $ 或 $\mathsf{Witness}_{\mathrm{UU}}$ 显式输出，
因此拒绝总伴随可复验见证。证毕。
\end{proof}

\subsection{同伦修复塔（工程版）：最小动作库 $\Homotopy$ 的三类分解}
\label{sec:eng-minimal-homotopy-actions}

\begin{definition}[同伦修复塔与动作库分解（失败类型$\to$动作类型）]
\label{def:eng-minimal-homotopy-action-library}
同伦修复塔是一条嵌入/扩张链：
\[
  \mathcal{L}_0 \hookrightarrow \mathcal{L}_1 \hookrightarrow \cdots,
\qquad
  \mathcal{L}_{k+1}=h_k(\mathcal{L}_k),\ \ h_k\in\Homotopy.
\]
最小要求为（证书保持 + 势能可下降）：
\[
  \mathcal{E}_{\mathrm{prot}}(\mathcal{L}_{k+1})=\top,\qquad
  \Phi(\mathcal{L}_{k+1})\le \Phi(\mathcal{L}_k)-\delta_k,\ \delta_k>0.
\]
建议把动作库按失败类型分解为三类（最小可迁移结构）：
\[
  \Homotopy=\Homotopy_{\mathrm{KU}}\ \cup\ \Homotopy_{\mathrm{UU}}\ \cup\ \Homotopy_{\mathrm{Sem}}.
\]
\end{definition}

\begin{remark}[三类最小动作族（只给接口口径，不绑定具体实现）]
\begin{itemize}
  \item $\Homotopy_{\mathrm{KU}}$：残差大但闭包未破（参数/结构优化、控制律调参、版图/微结构调参、权重重平衡等）；
  \item $\Homotopy_{\mathrm{UU}}$：闭包被击穿（扩张模型语言/新增扇区/新增耦合联络/新增证书规则/分层抬升等）；
  \item $\Homotopy_{\mathrm{Sem}}$：语义不稳定（新增/加强语义函子与审计器，把表达归一到 $\mathsf{NF}$：单位/坐标/命名/版本锁/构建复现等）。
\end{itemize}
\end{remark}

\subsection{发动机最小算子丛模板：\texorpdfstring{$S_{\mathrm{eng}}$}{S\_eng} 扇区分解与证书实例化}
\label{sec:eng-minimal-engine-template}

\begin{definition}[发动机域的最小扇区集合（接口口径）]
\label{def:eng-minimal-engine-sectors}
发动机工程设计可取最小扇区分解为：
\[
  S_{\mathrm{eng}}=
  \{\mathrm{Geo},\mathrm{Mat},\mathrm{Therm},\mathrm{Comb},\mathrm{Ctrl},\mathrm{Reg},\mathrm{Mfg}\}.
\]
其含义（只列接口不列具体 PDE/模型）：
\begin{itemize}
  \item $\mathrm{Geo}$：几何/拓扑参数；算子：CAD/网格生成/形状变换；
  \item $\mathrm{Mat}$：材料与热/力/疲劳参数；算子：材料选择、热处理、寿命模型；
  \item $\mathrm{Therm}$：热传导/对流/冷却；算子：热仿真、热管理结构调整；
  \item $\mathrm{Comb}$：燃烧/流动/排放；算子：燃烧模型、喷射策略、EGR、后处理；
  \item $\mathrm{Ctrl}$：控制律与标定；算子：控制器合成、标定、鲁棒控制；
  \item $\mathrm{Reg}$：法规/测试工况/合规边界；算子：工况集变换、合规判定器；
  \item $\mathrm{Mfg}$：制造工艺窗口/公差/可制造；算子：公差传播、工艺可行性、成本模型。
\end{itemize}
\end{definition}

\begin{remark}[耦合联络示例（非交换是常态）]
典型耦合联络可记为
\[
  \nabla_{\mathrm{Geo}\to \mathrm{Therm}},\
  \nabla_{\mathrm{Therm}\to \mathrm{Mat}},\
  \nabla_{\mathrm{Comb}\to \mathrm{Reg}},\
  \nabla_{\mathrm{Mfg}\to \mathrm{Geo}},\ \dots
\]
一般不交换：先改几何再做热管理，与先做热管理再改几何会给出不同的可行域与证书边界。
\end{remark}

\begin{proposition}[发动机域三证书接口的最小实例化（只要求可复算与可审计）]
\label{prop:eng-minimal-engine-cert-instances}
在定义~\ref{def:eng-minimal-engine-sectors} 的扇区分解下，可按如下最小口径实例化三证书接口：
\begin{enumerate}
  \item $\mathcal{E}_{\mathrm{prot}}$：安全（爆震/超温/超压）、结构强度/疲劳下界、法规红线、可制造红线；
  \item $\mathcal{E}_{\mathrm{CH}}$：性能/油耗/扭矩、热效率、NVH、成本等在已知模型内的残差；
  \item $\JacGap$：未建模失效机理、法规工况变化、材料批次漂移、制造偏差引入的新模式（以 UU 见证记录为最小证书）。
\end{enumerate}
\end{proposition}

\begin{Certificate}[证书：扇区$\to$证书字段的映射表（接口级）]
\label{cert:eng-minimal-engine-cert-instances}
证书登记谓词链：
\[
\begin{aligned}
  E_1 &: S_{\mathrm{eng}}\ \text{给出可枚举的约束来源}\Rightarrow \mathsf{Cert}_{\mathrm{prot}}\ \text{可登记
    RuleID/margin/evidence/hash},\\
  E_2 &: \text{仿真/试验/角点集}\Rightarrow \mathsf{Cert}_{\mathrm{CH}}\ \text{可登记残差向量 }r\ \text{及聚合 }W,\\
  E_3 &: \text{出现闭包外失败模式}\Rightarrow \exists\,\mathsf{Witness}_{\mathrm{UU}}\Rightarrow \JacGap>0\ \text{（命题~\ref{prop:eng-minimal-uu-witness-implies-jacgap}）}.
\end{aligned}
\]
\end{Certificate}

\begin{proof}[证明（谓词因果链）]
由 $E_1$：发动机硬红线可枚举并可给出裕度向量与复算证据，因此可按定理~\ref{thm:eng-minimal-eprot-margin} 装配 $\mathcal{E}_{\mathrm{prot}}$。
由 $E_2$：已知模型下的性能/成本等可计算残差并可聚合，按引理~\ref{lem:eng-minimal-ech-nonneg} 得 $\mathcal{E}_{\mathrm{CH}}$。
由 $E_3$：闭包外失败模式给出 UU 见证，按命题~\ref{prop:eng-minimal-uu-witness-implies-jacgap} 得 $\JacGap>0$。
合并即得实例化成立。证毕。
\end{proof}

\subsection{芯片最小算子丛模板：\texorpdfstring{$S_{\mathrm{chip}}$}{S\_chip} 扇区分解与证书实例化}
\label{sec:eng-minimal-chip-template}

\begin{definition}[芯片域的最小扇区集合（接口口径）]
\label{def:eng-minimal-chip-sectors}
芯片工程设计可取最小扇区分解为：
\[
  S_{\mathrm{chip}}=
  \{\mathrm{Arch},\mathrm{RTL},\mathrm{PPA},\mathrm{Timing},\mathrm{Layout},\mathrm{Proc},\mathrm{Rel},\mathrm{Sec},
    \mathrm{SC}\}.
\]
其中：
\begin{itemize}
  \item $\mathrm{Arch}$：体系结构与微结构；
  \item $\mathrm{RTL}$：功能/逻辑描述；
  \item $\mathrm{PPA}$：功耗-性能-面积；
  \item $\mathrm{Timing}$：时序闭合；
  \item $\mathrm{Layout}$：物理实现（Place\&Route、DRC/LVS）；
  \item $\mathrm{Proc}$：工艺窗口/PDK/角点；
  \item $\mathrm{Rel}$：可靠性（EM、BTI、ESD、老化）；
  \item $\mathrm{Sec}$：安全（侧信道/攻击面）；
  \item $\mathrm{SC}$：供应链/封装/材料与可得性约束。
\end{itemize}
\end{definition}

\begin{remark}[耦合联络示例（非交换是常态）]
典型耦合联络可记为
\[
  \nabla_{\mathrm{Layout}\to \mathrm{Timing}},\
  \nabla_{\mathrm{Proc}\to \mathrm{Rel}},\
  \nabla_{\mathrm{Sec}\to \mathrm{PPA}},\
  \nabla_{\mathrm{SC}\to \mathrm{Proc}},\dots
\]
一般不交换：先改布局再做时序，与先改 RTL/约束再做布局属于不同算子序列，导致不同的可行域与证书边界。
\end{remark}

\begin{proposition}[芯片域三证书接口的最小实例化（只要求可复算与可审计）]
\label{prop:eng-minimal-chip-cert-instances}
在定义~\ref{def:eng-minimal-chip-sectors} 的扇区分解下，可按如下最小口径实例化三证书接口：
\begin{enumerate}
  \item $\mathcal{E}_{\mathrm{prot}}$：DRC clean、LVS 通过、关键安全红线、关键可靠性红线、关键合规/出口约束；
  \item $\mathcal{E}_{\mathrm{CH}}$：PPA 差距、时序 slack 聚合、功耗/温度残差、面积残差；
  \item $\JacGap$：PDK/规则变更导致闭包被击穿、硅后数据与 sign-off 模型偏离、新安全攻击面、供应链规则变化引入的新约束维度（以 UU 见证为最小证书）。
\end{enumerate}
\end{proposition}

\begin{Certificate}[证书：签核/工具链哈希 + UU 见证日志]
\label{cert:eng-minimal-chip-cert-instances}
证书登记谓词链：
\[
\begin{aligned}
  C_1 &: \text{DRC/LVS/签核报告/工具链}\Rightarrow \mathsf{Cert}_{\mathrm{prot}}\ \text{可登记 hash 与 margin},\\
  C_2 &: \text{约束/角点/脚本锁定}\Rightarrow \mathsf{Cert}_{\mathrm{CH}}\ \text{可登记 }r,W\ \text{并复算},\\
  C_3 &: \text{规则/硅后/威胁模型演化}\Rightarrow \exists\,\mathsf{Witness}_{\mathrm{UU}}\Rightarrow \JacGap>0\
    \text{（命题~\ref{prop:eng-minimal-uu-witness-implies-jacgap}）}.
\end{aligned}
\]
\end{Certificate}

\begin{proof}[证明（谓词因果链）]
由 $C_1$：签核类硬红线可由确定性工具链裁决并可复算审计，故可装配 $\mathcal{E}_{\mathrm{prot}}$。
由 $C_2$：PPA/时序等残差在既定模型与角点集下可计算并可复算，故得 $\mathcal{E}_{\mathrm{CH}}$。
由 $C_3$：闭包外变化给出 UU 见证，按命题~\ref{prop:eng-minimal-uu-witness-implies-jacgap} 得 $\JacGap>0$。
合并即得实例化成立。证毕。
\end{proof}

\subsection{端到端 pipeline：两域完全同构的 law-object 更新—证书裁决—修复闭环}
\label{sec:eng-minimal-end2end-pipeline}

\begin{corollary}[发动机/芯片两域同构：同一闭环模板可直接迁移]
\label{cor:eng-minimal-engine-chip-isomorphic}
在最小模板（定义~\ref{def:eng-minimal-law-object-template}）与闭环规则（定理~\ref{thm:eng-minimal-loop-induction}）下，
发动机与芯片两域的设计迭代可用同一条 pipeline 表达：
\begin{enumerate}
  \item 语义归一：将输入表达 $x\in\mathsf{WP}$ 送入 $\mathcal{F}^{\mathrm{aud}}$ 得到 $\mathsf{NF}$ 表达或 $\bot$；
  \item 组装 law-object：形成 $\mathcal{L}_k\in(\GZLS)_{\mathrm{Eng}}$（含扇区算子丛与耦合联络）；
  \item 评估三证书与势能：$\mathcal{E}_{\mathrm{prot}},\mathcal{E}_{\mathrm{CH}},\JacGap,\Phi$；
  \item 接受/拒绝：接受需 $\mathcal{E}_{\mathrm{prot}}=\top$ 且 $\Phi$ 满足下降/阈值；否则输出见证并触发修复动作 $h_k\in\Homotopy$；
  \item 迭代至满足阈值或得到“需要扩张层/扇区”的结构性结论（UU 持续见证）。
 \end{enumerate}
两域差异只体现在扇区集合 $S_{\mathrm{eng}}/S_{\mathrm{chip}}$、证书字段内容与动作库的具体实现，不改变闭环形状。
\end{corollary}

\begin{Certificate}[证书：模板复用只依赖三证书接口与闭环判据]
\label{cert:eng-minimal-engine-chip-isomorphic}
证书登记谓词链：
\[
\begin{aligned}
  I_1 &: \text{定义~\ref{def:eng-minimal-law-object-template}}\Rightarrow \text{两域均可编码为 }(\GZLS)_{\mathrm{Eng}}\ \text{对象},
    \\
  I_2 &: \text{命题~\ref{prop:eng-minimal-engine-cert-instances} 与~\ref{prop:eng-minimal-chip-cert-instances}}\Rightarrow
    \text{两域均可实例化三证书接口},\\
  I_3 &: \text{定理~\ref{thm:eng-minimal-loop-induction}}\Rightarrow \text{闭环的多步合规性与见证输出与具体域无关}.
\end{aligned}
\]
\end{Certificate}

\begin{proof}[证明（谓词因果链）]
由 $I_1$：两域都能落入同一 law-object 模板；
由 $I_2$：三证书接口在两域中均可装配；
由 $I_3$：闭环保持性与证书化下降性只依赖接口与判据，不依赖域内细节。
合并即得同构迁移结论。证毕。
\end{proof}

\subsection{最小落地版本（MVP）：不引入额外复杂度即可跑通闭环}
\label{sec:eng-minimal-mvp}

\begin{remark}[MVP（建议只实现三件事）]
若先做一条可运行的最小闭环（MVP），建议只实现：
\begin{enumerate}
  \item \textbf{语义规范场最小集}：每域先做 5--8 个 $F_i^{\mathrm{aud}}$（单位/量纲归一、命名归一、版本锁、构建复现、测试回归等）；
  \item \textbf{$\mathcal{E}_{\mathrm{prot}}$ 的 fail-closed 审计器}：把硬红线做成确定性裁决器（最优先）；
  \item \textbf{$\JacGap$ 的 UU 见证日志}：任何闭包外反例都结构化记录并触发 $\Homotopy_{\mathrm{UU}}$ 或 $\Homotopy_{\mathrm{Sem}}$ 动作。
\end{enumerate}
一旦闭环跑通，后续精细模型（热/燃烧/可靠性/安全/角点扩张等）只是在动作库里新增算子，属于同伦修复塔的自然扩张。
\end{remark}

\section{本次讨论的增益总结：概念—定理—工程—评测—战略意义}
\label{sec:discussion-gains-summary}

\begin{remark}[定位（把“直觉判断”升级为形式系统资产）]
本次讨论的核心增益是：把一组“直觉上非常强”的判断（开放/封闭、回合、KU/UU、神经网络滞后、分层必要性）
\textbf{从叙述层升级为可证明的形式系统链条}，
并同时给出可落地的\textbf{工程接口}（可审计、可回退、可证书化）与\textbf{可检验评测路径}（逐步开封）。

本章按五层结构收束：\textbf{概念—定理—工程—评测—战略意义}，并将其写为可引用的“证书 + 谓词因果链”形式。
\end{remark}

\subsection{概念层：回合、KU/UU、开放/封闭的可判定边界}
\label{sec:gains-concepts}

\begin{proposition}[概念边界钉死：核心术语均已可判定化]
\label{prop:gains-concept-boundary}
在本文的形式系统内，下列术语均已被压缩为可判定/可核验对象：
\begin{enumerate}
  \item \textbf{回合}：由回合分割 $\mathcal{P}$ 诱导的边界语义序列（定义~\ref{def:round-partition}）；
  \item \textbf{KU/UU}：由后继闭合缺口 $\mathsf{Ob}^{(\mathcal{P})}(o,a)$ 判定（定义~\ref{def:round-obstruction} 与~\ref{def:round-ku-uu}）；
  \item \textbf{严格回合同伦类是否存在}：由回合细化下同伦型稳定性判定（定义~\ref{def:strict-round-homotopy-class-exists}），
  并与 $\JacGap$ 的“稳定归零”口径等价（定义~\ref{def:jacgap-stable-zero} 与定理~\ref{thm:strict-round-homotopy-iff-jacgap-stable-zero}）；
  \item \textbf{开放/封闭}：可取“回合同伦类作用/holonomy”口径（定理~\ref{thm:round-homotopy-open-closed}），
  亦可取“封闭/开放博弈系统（回合细化不闭合/闭合）”口径（定义~\ref{def:closed-game-system-round} 与~\ref{def:open-game-system-round}）。
\end{enumerate}
\end{proposition}

\begin{Certificate}[证书：关键定义引用闭合]
\label{cert:gains-concept-boundary}
证书登记谓词链：
\[
\begin{aligned}
  C_1 &: \text{回合分割}\Rightarrow \text{回合边界语义序列（定义~\ref{def:round-partition}）},\\
  C_2 &: \mathsf{Ob}^{(\mathcal{P})}(o,a)\ \text{给出闭合缺口（定义~\ref{def:round-obstruction}）},\\
  C_3 &: \mathsf{Ob}=\varnothing\Rightarrow \text{KU；}\ \mathsf{Ob}\neq\varnothing\Rightarrow \text{UU（定义~\ref{def:round-ku-uu}）},\\
  C_4 &: \JacGap\ \text{由}\ \mathsf{Ob}\ \text{度量得到（定义~\ref{def:jacgap-from-obstruction}）},\\
  C_5 &: \JacGap\ \text{稳定归零}\Longleftrightarrow \text{严格回合同伦类存在（定义~\ref{def:jacgap-stable-zero}；定理~\ref{thm:strict-round-homotopy-iff-jacgap-stable-zero}）}.
\end{aligned}
\]
\end{Certificate}

\begin{proof}[证明（谓词因果链）]
由定义~\ref{def:round-partition} 得 $C_1$。
由定义~\ref{def:round-obstruction} 得 $C_2$；代入 KU/UU 定义~\ref{def:round-ku-uu} 得 $C_3$。
由定义~\ref{def:jacgap-from-obstruction} 得 $C_4$。
由定义~\ref{def:jacgap-stable-zero} 与定理~\ref{thm:strict-round-homotopy-iff-jacgap-stable-zero} 得 $C_5$。
合并 $C_1$--$C_5$ 即得命题结论。证毕。
\end{proof}

\begin{remark}[评价（概念层增益）]
本层的价值在于：把“开放系统为什么难、难在哪里”从“信息更多/更复杂”的泛泛描述，
钉死为\textbf{可判定的结构缺口}（$\mathsf{Ob}$、$\JacGap$、同伦稳定性/holonomy），从而为后续定理链与工程接口提供共同语言。
\end{remark}

\subsection{定理层：统一机制与可引用结论族（封闭/开放、UU、分层、滞后）}
\label{sec:gains-theorems}

\begin{theorem}[核心结构链条总览（可引用结论族）]
\label{thm:gains-core-chain}
本文已形成一组可引用的强结论，并共享同一因果核：
\begin{enumerate}
  \item \textbf{封闭系统链}：封闭博弈系统 $\Rightarrow$ 严格回合同伦类存在 $\Rightarrow$ 障碍必为 KU
  （定理~\ref{thm:closed-implies-strict-round-homotopy} 与~\ref{thm:closed-phi-jacgap-ku}）。
  \item \textbf{开放系统链}：开放博弈系统 $\Rightarrow$ 严格回合同伦类不存在 $\Rightarrow$ 障碍必为 UU $\Rightarrow$ 回合只能逼近
  （定理~\ref{thm:open-implies-no-strict-round-homotopy} 与~\ref{thm:open-phi-jacgap-uu}）。
  \item \textbf{仅回合逼近不够}：在开放 UU（非平凡 holonomy）下，仅做回合细分（subdivision）无法消除 UU 根源
  （必要性定理~\ref{thm:uu-open-subdivision-not-enough}）。
  \item \textbf{分层必要性}：要从结构上消除 UU 根源，必须做语义抬升/连续统分层（层完备性 $\Longleftrightarrow$ holonomy 被杀死）
  （定理~\ref{thm:layering-complete-iff-trivial-holonomy}）。
  \item \textbf{神经网络滞后}：在 holonomy 非平凡且产生 UU 修复不可兼容证书的口径下，
  任何仅依赖 $o\in O$ 的主策略类落入必然滞后同伦类 $\lbrack\mathsf{Lag}\rbrack$，因此 NN 不能作为主控制环证明义务承担者
  （定理~\ref{thm:lag-holonomy-to-lag-class} 与推论~\ref{cor:demote-nn-from-main-loop}）。
\end{enumerate}
\end{theorem}

\begin{Certificate}[证书：五条结论的引用闭合]
\label{cert:gains-core-chain}
证书登记谓词链：
\[
\begin{aligned}
  T_1 &: \text{定理~\ref{thm:closed-implies-strict-round-homotopy} 与~\ref{thm:closed-phi-jacgap-ku}},\\
  T_2 &: \text{定理~\ref{thm:open-implies-no-strict-round-homotopy} 与~\ref{thm:open-phi-jacgap-uu}},\\
  T_3 &: \text{定理~\ref{thm:uu-open-subdivision-not-enough}},\\
  T_4 &: \text{定理~\ref{thm:layering-complete-iff-trivial-holonomy}},\\
  T_5 &: \text{定理~\ref{thm:lag-holonomy-to-lag-class} 与推论~\ref{cor:demote-nn-from-main-loop}}.
\end{aligned}
\]
\end{Certificate}

\begin{proof}[证明（谓词因果链）]
五条结论分别由证书~\ref{cert:gains-core-chain} 中 $T_1$--$T_5$ 的对应定理/推论直接推出。证毕。
\end{proof}

\begin{remark}[共享因果核（统一解释）]
上述结论共享同一因果核（写作可证书化的结构链条）：
\[
\begin{aligned}
  &\text{开放性/挑战压力}\\
  &\Rightarrow\ \exists(o,a):\ \mathsf{Ob}^{(\mathcal{P})}(o,a)\neq\varnothing\\
  &\Rightarrow\ \JacGap>0\ \text{（结构缺口可证书化）}\\
  &\Rightarrow\ \Phi=\mathcal{E}_{\mathrm{CH}}+\lambda\JacGap\ \text{必须在证书保持下可下降}\\
  &\Rightarrow\ \text{需要分层与修复塔，而非仅回合或仅拟合}.
\end{aligned}
\]
这把“挑战塑形智慧”“AI 停留在封闭博弈层”的讨论从方向判断，降级为“结构上卡在何处、缺哪一个接口”的可推演问题。
\end{remark}

\subsection{工程层：证明义务转移（学习器$\to$审计器）与可回退接口}
\label{sec:gains-engineering}

\begin{proposition}[工程接口增益：撤销学习器证明义务，改以审计包装实现可核验闭环]
\label{prop:gains-engineering-wrapper}
若目标是 UU 口径下的主控制环可核验性（定义~\ref{def:main-control-loop-proof-obligation}），则：
\begin{enumerate}
  \item 神经网络必须从“证明义务承担者”位置撤下（推论~\ref{cor:demote-nn-from-main-loop}）；
  \item 学习器可在工程上保留为\textbf{候选生成器/加速器}，并由确定性审计器 $V$ 裁决（审计包装；定理~\ref{thm:audited-wrapper-soundness}）；
  \item 在 fail-closed（拒绝即回退）口径下，多轮调用仍保持合同与可追溯性（公理降级：定理~\ref{thm:audited-wrapper-inductive-safety}）。
\end{enumerate}
\end{proposition}

\begin{Certificate}[证书：证明义务转移到 $V$；失败闭合]
\label{cert:gains-engineering-wrapper}
证书登记谓词链：
\[
\begin{aligned}
  E_1 &: \text{推论~\ref{cor:demote-nn-from-main-loop}}\Rightarrow \text{NN 不能作为主环证明义务承担者},\\
  E_2 &: \text{定理~\ref{thm:audited-wrapper-soundness}}\Rightarrow \text{审计通过的输出可核验且与 NN 泛化无关},\\
  E_3 &: \text{定理~\ref{thm:audited-wrapper-inductive-safety}}\Rightarrow \text{多轮执行在 fail-closed 下保持合同}.
\end{aligned}
\]
\end{Certificate}

\begin{proof}[证明（谓词因果链）]
由 $E_1$ 得“撤销学习器证明义务”的必要性；
由 $E_2$ 得“候选生成 + 确定性裁决”的正确性转移；
由 $E_3$ 得“可审计 + 可回退”在多轮闭环下仍成立。证毕。
\end{proof}

\begin{remark}[评价（工程层增益）]
本层解决了工程上的核心矛盾：\textbf{既利用学习器的效率，又不把系统安全与鲁棒性押在“泛化口头保证”上}。
其形式含义是：学习器最多影响“搜索顺序/候选质量”，而系统对外正确性与可追溯性锚定在确定性审计器与证书日志。
\end{remark}

\subsection{评测层：把“封闭/开放层级判断”转为评测族与验收指标}
\label{sec:gains-evaluation}

\begin{proposition}[可检验化：逐步开封评测族以 $\JacGap$ 与 $\Phi$ 为指标]
\label{prop:gains-evaluation-path}
对“是否仍停留在封闭博弈层”的判断，可在形式系统内替换为可检验路径：
\begin{enumerate}
  \item 封闭训练/评测高分不蕴含开放 UU 修复能力（推论~\ref{cor:closed-eval-does-not-imply-open}）；
  \item 通过五类开封算子族（定义~\ref{def:open-eval-five-operators}）逐步引入漂移/反制/多约束压力；
  \item 以“是否能在 $\mathcal{E}_{\mathrm{prot}}$ 保持下压制 $\JacGap$ 并推动 $\Phi$ 下降（修复塔可终止；定理~\ref{thm:repair-tower-terminates}）
    ”作为验收指标。
\end{enumerate}
\end{proposition}

\begin{Certificate}[证书：封闭高分不约束 UU；开封$\Rightarrow$压力；以 $\Phi$-下降验收]
\label{cert:gains-evaluation-path}
证书登记谓词链：
\[
\begin{aligned}
  V_1 &: \text{推论~\ref{cor:closed-eval-does-not-imply-open}},\\
  V_2 &: \text{定义~\ref{def:open-eval-five-operators}}\Rightarrow \text{可构造逐步开封的压力族},\\
  V_3 &: \text{命题~\ref{prop:challenge-implies-phi-pressure}}\Rightarrow \text{压力出现将诱导}\ \Phi>0\ \text{见证},\\
  V_4 &: \text{定理~\ref{thm:repair-tower-terminates}}\Rightarrow \text{若存在证书化下降塔则给出有限步验收保证}.
\end{aligned}
\]
\end{Certificate}

\begin{proof}[证明（谓词因果链）]
由 $V_1$ 得“封闭评测只能证明封闭能力”。
由 $V_2$ 可逐步打开封闭评测，使挑战谓词可持续出现；
由 $V_3$ 得挑战压力产生 $\Phi>0$ 的证书化压力，因此评测可触发 UU/修复需求；
若系统在该压力下仍能组织证书化下降修复塔，则由 $V_4$ 得到可终止的验收保证。证毕。
\end{proof}

\begin{remark}[评价（评测层增益）]
本层把“AI 是否停留在封闭博弈层”的争论，转化为“开封后 $\JacGap$ 是否可被压制、$\Phi$ 是否可下降”的\textbf{可操作实验设计与验收指标}。
因此理论不仅能解释现象，还能生成可检验预测与工程验收门槛。
\end{remark}

\subsection{战略意义：对卷册与工程线的直接增益（沉淀建议）}
\label{sec:gains-strategy}

\begin{corollary}[沉淀动作：形成可编号的定理族与统一证书 schema]
\label{cor:gains-deposition}
由前述概念边界（命题~\ref{prop:gains-concept-boundary}）、核心定理族（定理~\ref{thm:gains-core-chain}）、
工程审计接口（命题~\ref{prop:gains-engineering-wrapper}）与评测路径（命题~\ref{prop:gains-evaluation-path}），
可将本次讨论沉淀为可并入卷册与工程文档的“定理/证书资产包”，建议最小包含：
\[
  (T_{\mathrm{Closed}}),(T_{\mathrm{Open}}),(T_{\mathrm{NoRoundOnly}}),(T_{\mathrm{LayeringNecessary}}),
  (T_{\mathsf{Lag}}),(T_{\mathrm{AuditedWrapper}}),
\]
并配套统一的证书 schema（输入/输出/验证谓词/失败回退规则），以便在 PureMath 与 Vol1/Vol2/Vol3 工程线中直接复用。
\end{corollary}

\begin{Certificate}[证书：概念边界 + 定理链 + 工程接口 + 评测路径 $\Rightarrow$ 可沉淀为资产包]
\label{cert:gains-deposition}
证书登记谓词链：
\[
\begin{aligned}
  E_1 &: \text{命题~\ref{prop:gains-concept-boundary}}\Rightarrow \text{概念边界已固定},\\
  E_2 &: \text{定理~\ref{thm:gains-core-chain}}\Rightarrow \text{核心定理族与证书链已给出},\\
  E_3 &: \text{命题~\ref{prop:gains-engineering-wrapper}}\Rightarrow \text{工程审计包装接口已给出},\\
  E_4 &: \text{命题~\ref{prop:gains-evaluation-path}}\Rightarrow \text{评测族与验收指标已给出}.
\end{aligned}
\]
并记谓词 $E_5$ 为：若 $E_1\wedge E_2\wedge E_3\wedge E_4$，则可将上述结论沉淀为可复用“定理/证书资产包”，
并对齐统一证书 schema（章节~\ref{sec:unified-certificate-schema}）。
\end{Certificate}

\begin{proof}[证明（谓词因果链）]
由 $E_1$，本次讨论先把关键术语的边界固定为可判定结构；由 $E_2$，核心结论已形成可编号的定理链及其证书化支撑；
由 $E_3$，对外工程落地时的审计包装接口已给出；由 $E_4$，评测与验收也已被可检验化。
因此把这些要素按“定理编号 + 证书 schema”打包沉淀，即得到可并入卷册与工程文档的资产包（$E_5$）。
推论成立。证毕。
\end{proof}

\begin{remark}[总体评价（一句话）]
本次讨论把“开放系统智能”的关键困难从经验层抽象成可证书化的结构缺口（UU/holonomy/$\JacGap$），并同时给出了：
必要的相空间分层机制、可下降的缺陷势 $\Phi$、同伦修复塔的闭环、以及把 NN 降级为审计下候选/压缩器的工程接口与评测路径。
因此它既是可证明的形式系统链条，也是可实现的工程接口与可检验的评测方案。
\end{remark}

\subsection{公理（降级为定理）：五层总结的证书化闭合（数学归纳法证书）}
\label{sec:gains-induction-schema}

\begin{theorem}[公理降级：五层总结结论对前文定义/定理链闭合（数学归纳法证书）]
\label{thm:gains-five-layer-induction}
令 $P(k)$ 表示“本章前 $k$ 层总结（概念、定理、工程、评测、战略）均已由本文既有定义/定理链条给出证书化支撑”，其中 $k\in\{1,2,3,4,5\}$。
若满足：
\[
  P(1)\qquad\text{且}\qquad \forall k<5,\ \bigl(P(k)\Rightarrow P(k+1)\bigr),
\]
则 $P(5)$ 成立，即五层总结对前文闭合。
\end{theorem}

\begin{Certificate}[数学归纳法证书：$(P(1),\forall k<5(P(k)\Rightarrow P(k+1)))$]
\label{cert:gains-five-layer-induction}
证书对象可取为
\[
  \pi_{\mathrm{ind}}
  :=\bigl(P(1),\ \forall k<5\,(P(k)\Rightarrow P(k+1))\bigr),
\]
其中 $P(1)$ 由命题~\ref{prop:gains-concept-boundary} 给出，$P(k)\Rightarrow P(k+1)$ 由本章相邻层结论与其证书链（证书~\ref{cert:gains-core-chain}、
  \ref{cert:gains-engineering-wrapper}、\ref{cert:gains-evaluation-path}）给出。
\end{Certificate}

\begin{proof}[证明（数学归纳法证书；谓词因果链）]
基步 $P(1)$：由命题~\ref{prop:gains-concept-boundary} 及其证书~\ref{cert:gains-concept-boundary} 得概念层闭合。
归纳步：假设 $P(k)$（$k<5$）成立，则由本章第 $k+1$ 层对应的结论与证书链，
可将“新增层”的结论归约到已知定义/定理引用，从而得到 $P(k+1)$。
因此由归纳法得 $P(5)$，五层总结对前文闭合。证毕。
\end{proof}

\appendix

\section{附录A：符号体系与记号说明（统一注释）}
本附录汇总全文跨章节反复出现、且具有“接口含义”的符号与记号，作为统一注释源。
若正文某处对同一符号给出更窄/更具体的局部定义，以正文该处为准。

\subsection{核心对象与算子口径}
\begin{longtable}{>{\raggedright\arraybackslash}p{0.22\textwidth} >{\raggedright\arraybackslash}p{0.72\textwidth}}
\hline
\textbf{符号} & \textbf{说明（默认口径）} \\\hline
$\mathcal{S}=\mathcal{X}\times \mathcal{Z}$，$s=(x,z)$
& 总体状态空间与分解：$x$ 为物理态，$z$ 为工程态。\\
$\mathcal{O}$，$O_t$
& 总算子族；离散阶段更新 $s_{t+1}=O_t(s_t)$。\\
$\TOPB$
& 总算子丛（总算子族的工程抽象总称）。\\
$\PTOPB$，$\ETOPB$
& \PTOPB/\ETOPB：以物理态/工程态分量为主导的算子子族（按投影主导口径区分）。\\
$\delta_t$，$\pi_t$
& 外界实时干扰算子/系统构造算子；常用复合 $O_t=\pi_t\circ\delta_t$（顺序可按场景调整）。\\
$\Delta$，$\mathcal{C}$
& 干扰算子集合/构造算子集合。\\
$\Delta_{\mathrm{P}}$，$\Delta_{\mathrm{E}}$
& \PTOPB/\ETOPB 侧的外生干扰子族。\\
$\mathcal{C}_{\mathrm{P}}$，$\mathcal{C}_{\mathrm{E}}$
& \PTOPB/\ETOPB 侧的内生构造子族。\\
$\mathcal{C}^{\mathrm{fun}}$，$\mathcal{C}^{\mathrm{ass}}$
& 构造算子族内部的“功能链/保障链”二分。\\
$\DeltaRT$
& 实时干扰算子集合（允许包含内生负算子；正文中外界干扰 $\Delta$ 可视为其外生子集）。\\
\hline
\end{longtable}

\subsection{同伦/上同调与缺口指标}
\begin{longtable}{>{\raggedright\arraybackslash}p{0.22\textwidth} >{\raggedright\arraybackslash}p{0.72\textwidth}}
\hline
\textbf{符号} & \textbf{说明（默认口径）} \\\hline
$d_0\sim d_1$，$c_0\simeq c_1$
& 干扰/构造的同伦等价：存在不穿越关键约束边界的连续变形族，把“换招不换势”压缩为等价类登记。\\
$H^1=\ker(\partial^1)/\mathrm{im}(\partial^0)$
& 第一上同调：可消干扰为上边界，不可消干扰落在非平凡上同调类中（结构性短板）。\\
$T_{\gamma}$
& 回合路径/搜索路径上的传输（transport）复合映射；用于把“路径同伦类的作用”显式化。\\
$\mathsf{Hol}_{[\ell]}$
& holonomy（同伦类作用）：回到同一语义点的回路对纤维内部态的自同态作用；非平凡时对应 UU 根源。\\
$\JacGap$
& Jacobiator-gap/缺口证书：把 UU/obstruction 见证压缩为可计算缺陷量；$\JacGap>0$ 表示存在结构性缺口压力。\\
$\Phi=\mathcal{E}_{\mathrm{CH}}+\lambda\JacGap$
& 合并缺陷势：把可计算能量项与缺口项线性合并，用于驱动“证书保持下的严格下降”与修复塔终止性判据。\\
$\mathcal{E}_{\mathrm{prot}}\in\{\top,\bot\}$
& 受保护一致性证书谓词：表示安全/合规/可审计等关键约束是否保持。\\
$\bot$
& 失败显式（fail-closed）符号：任何未通过确定性验证/证书裁决者必须返回 $\bot$ 并回退，而非“静默通过”。\\
\hline
\end{longtable}

\section{附录B：证书调用接口（全文）}
本附录给出全文“证书（Certificate）”的调用约定与索引，作为对外复用的最小接口。
证书 schema 的字段化合同（输入/输出/验证谓词/失败回退规则）见章节~\ref{sec:unified-certificate-schema}。

\subsection{最小调用约定}
\begin{itemize}
  \item \textbf{正文体例}：默认按“断言（定理/引理/命题/推论）$\to$ 证书 $\to$ 证明（谓词因果链）”顺序出现。
  \item \textbf{对外复用}：建议成对引用“断言 + 证书”（例如“定理~\ref{thm:uu-eng-iff} 与 证书~\ref{cert:uu-eng-iff}”），
  以便在他处复算时直接定位验证谓词、见证载体与 fail-closed 回退规则。
  \item \textbf{工程落地}：证书载体可取为日志、回归/变形测试报告、反例样本、复杂度上界证明等；拒绝必须显式输出 $\bot$ 并触发回退/修复塔。
\end{itemize}

\subsection{标签命名约定（便于跨文档调用）}
本文默认使用如下前缀作为“稳定接口名”：
\begin{center}
\renewcommand{\arraystretch}{1.05}
\begin{tabular}{>{\raggedright\arraybackslash}p{0.18\textwidth} >{\raggedright\arraybackslash}p{0.72\textwidth}}
  \texttt{thm:}  & 定理 \\
  \texttt{lem:}  & 引理 \\
  \texttt{prop:} & 命题 \\
  \texttt{cor:}  & 推论 \\
  \texttt{cert:} & 证书 \\
  \texttt{def:}  & 定义 \\
  \texttt{asm:}  & 假设 \\
  \texttt{sec:}  & 章节 \\
\end{tabular}
\end{center}

\subsection{证书索引（证书标签作为调用接口）}
{\small
\let\texttt\code
\setlength{\tabcolsep}{4pt}
\renewcommand{\arraystretch}{1.15}
\begin{longtable}{>{\raggedright\arraybackslash}p{0.20\textwidth} >{\raggedright\arraybackslash}p{0.32\textwidth}
  >{\raggedright\arraybackslash}p{0.42\textwidth}}
\hline
\textbf{对应断言} & \textbf{证书标签（API）} & \textbf{证书标题（摘要）} \\\hline
% ===== 证书索引表（自动生成区）=====
% BEGIN_CERT_INDEX
定理~\ref{thm:source-controllability-splitting} & \texttt{\detokenize{cert:splitting-by-time-horizon}} & 数学归纳法证书：按阶段数 $T$
  归纳的交错分解 \\
引理~\ref{lem:disturbance-unified-representation} & \texttt{\detokenize{cert:disturbance-unified-representation}} & 证书：
  三类扰动的等价封装 \\
命题~\ref{prop:disturbance-ptopb-etopb} & \texttt{\detokenize{cert:disturbance-ptopb-etopb}} & 证书：投影主导的判定规则 \\
定理~\ref{thm:noncommutativity-is-default} & \texttt{\detokenize{cert:noncommutativity-witness}} & 证书：非交换性的可观测见证（交换子非零） \\
引理~\ref{lem:reachable-set-order-sensitive} & \texttt{\detokenize{cert:reachable-set-order-sensitive}} & 证书：顺序敏感对的可达差异 \\
命题~\ref{prop:order-is-structural-in-adversarial} & \texttt{\detokenize{cert:order-is-structural}} & 证书：顺序敏感导致风险边界差异 \\
推论~\ref{cor:noncommutative-modeling-necessary} & \texttt{\detokenize{cert:noncommutative-modeling-necessary}} & 证书：
  最弱入口的区分性 \\
命题~\ref{prop:operator-library-as-certificate-interface} & \texttt{\detokenize{cert:
  operator-library-as-certificate-interface}} & 证书：条目字段 $\Rightarrow$ 可追溯调用边界 \\
定理~\ref{thm:short-exact-implies-long-exact} & \texttt{\detokenize{cert:short-exact-implies-long-exact}} & 证书：短正合列
  $\Rightarrow$ 长正合列 \\
定理~\ref{thm:layered-augmentation-by-induction} & \texttt{\detokenize{cert:layered-augmentation-by-induction}} & 数学归纳法证书：
  按增强层数 $N$ 归纳 \\
引理~\ref{lem:disturbance-homotopy-is-equivalence} & \texttt{\detokenize{cert:disturbance-homotopy-is-equivalence}} & 证书：
  等价关系三性 \\
命题~\ref{prop:homotopy-class-compression} & \texttt{\detokenize{cert:homotopy-class-compression}} & 证书：同伦不变量因子化 \\
命题~\ref{prop:representative-optimization-unified-abstract} & \texttt{\detokenize{cert:representative-optimization-audit}}
  & 证书：代表元选择的审计边界 \\
定理~\ref{thm:cancellable-iff-coboundary} & \texttt{\detokenize{cert:cancellable-iff-coboundary}} & 证书：上边界判据 \\
引理~\ref{lem:commutator-enters-coboundary} & \texttt{\detokenize{cert:commutator-enters-coboundary}} & 证书：交换子见证非交换耦合 \\
推论~\ref{cor:nontrivial-cohomology-is-structural} & \texttt{\detokenize{cert:nontrivial-cohomology-is-structural}} & 证书：
  不可消 $\Rightarrow$ 需要改变结构类而非仅调参 \\
定理~\ref{thm:relative-cohomology-is-aug} & \texttt{\detokenize{cert:relative-cohomology-is-aug}} & 证书：商对象同构 $\Rightarrow$
  上同调同构 \\
定理~\ref{thm:two-level-augmentation-decomposition} & \texttt{\detokenize{cert:two-level-augmentation-decomposition}} &
  数学归纳法证书：按增强步数 $T$ 归纳 \\
引理~\ref{lem:fiber-augmentation-does-not-change-complex} & \texttt{\detokenize{cert:
  fiber-augmentation-does-not-change-complex}} & 证书：对象不变 $\Rightarrow$ 复形不变 \\
命题~\ref{prop:boundary-shift-can-trivialize-obstruction} & \texttt{\detokenize{cert:
  boundary-shift-can-trivialize-obstruction}} & 证书：障碍类平凡化 $\Rightarrow$ 存在上边界补偿 \\
推论~\ref{cor:need-boundary-shift-for-out-of-envelope-goal} & \texttt{\detokenize{cert:
  need-boundary-shift-for-out-of-envelope-goal}} & 证书：目标区间外 $\Rightarrow$ 需要参数迁移 \\
定理~\ref{thm:twisted-two-level-search-decomposition} & \texttt{\detokenize{cert:twisted-two-level-search-decomposition}}
  & 数学归纳法证书：按外层步数 $T$ 归纳的轨迹分解 \\
推论~\ref{cor:twisted-search-reduces-to-intra-phase} & \texttt{\detokenize{cert:twisted-search-reduces-to-intra-phase}} &
  证书：无相间步 $\Rightarrow$ $\lambda$ 常值 \\
引理~\ref{lem:transport-induces-map-on-homotopy-classes} & \texttt{\detokenize{cert:
  transport-induces-map-on-homotopy-classes}} & 证书：等价关系因子化 \\
命题~\ref{prop:noncommutative-process-implies-twist} & \texttt{\detokenize{cert:noncommutative-process-implies-twist}} &
  证书：端点相同但纤维态不同的见证 \\
推论~\ref{cor:untwisted-reduces-to-endpoint-control} & \texttt{\detokenize{cert:untwisted-reduces-to-endpoint-control}} &
  证书：搬运仅依赖端点 $\Rightarrow$ 无同伦扭曲 \\
定理~\ref{thm:continuation-ensures-feasible-rolling} & \texttt{\detokenize{cert:continuation-ensures-feasible-rolling}} &
  数学归纳法证书：按周期步数 $K$ 归纳的可行性保持 \\
引理~\ref{lem:smooth-switching-preserves-feasibility} & \texttt{\detokenize{cert:smooth-switching-preserves-feasibility}}
  & 证书：同伦参数插值的可行性保持 \\
命题~\ref{prop:gauge-fixing-does-not-change-invariants} & \texttt{\detokenize{cert:gauge-fixing-does-not-change-invariants}
  } & 证书：曲率协变变换 $\Rightarrow$ 不变量不变 \\
引理~\ref{lem:holonomy-implies-path-class-relevant} & \texttt{\detokenize{cert:holonomy-implies-path-class-relevant}} & 证书：
  回路非平凡 $\Rightarrow$ 存在端点相同但搬运不同的两条路径 \\
定理~\ref{thm:gauge-switching-sequence-decomposition} & \texttt{\detokenize{cert:gauge-switching-sequence-decomposition}}
  & 数学归纳法证书：按步数 $T$ 归纳的分段分解 \\
推论~\ref{cor:no-structure-switch-implies-same-model} & \texttt{\detokenize{cert:no-structure-switch-implies-same-model}}
  & 证书：无 C 级步 $\Rightarrow$ 模型常值 \\
命题~\ref{prop:conserved-bedrock-provides-global-constraints} & \texttt{\detokenize{cert:
  conserved-bedrock-provides-global-constraints}} & 证书：不变量沿复合算子保持 \\
定理~\ref{thm:phase-switching-as-model-triple-switching} & \texttt{\detokenize{cert:
  phase-switching-as-model-triple-switching}} & 证书：守恒底座固定 + 三元组分量变化的封闭性 \\
引理~\ref{lem:edge-repair-yields-valid-differential} & \texttt{\detokenize{cert:edge-repair-yields-valid-differential}} &
  证书：展开平方并逐项归零 \\
命题~\ref{prop:edge-repair-trivializes-obstruction-class} & \texttt{\detokenize{cert:
  edge-repair-trivializes-obstruction-class}} & 证书：上同调类平凡 $\Leftrightarrow$ 存在上边界表达 \\
推论~\ref{cor:edge-repair-changes-repairability-classification} & \texttt{\detokenize{cert:
  edge-repair-changes-repairability-classification}} & 证书：判据依赖于微分 $\partial$ \\
定理~\ref{thm:global-conservation-compat-local-nonconservation} & \texttt{\detokenize{cert:
  global-conservation-compat-local-nonconservation}} & 数学归纳法证书：按部件数 $n$ 归纳的闭合性 \\
命题~\ref{prop:subsystem-number-nonconservation-witnessed-by-edge} & \texttt{\detokenize{cert:
  subsystem-number-nonconservation-witnessed-by-edge}} & 证书：海森堡方程与对易分解 \\
定理~\ref{thm:global-feasible-space-invariant} & \texttt{\detokenize{cert:global-feasible-space-invariant}} & 数学归纳法证书：按步数
  $T$ 归纳的母空间闭合性 \\
命题~\ref{prop:sector-switching-respects-global-feasible-space} & \texttt{\detokenize{cert:
  sector-switching-respects-global-feasible-space}} & 证书：母空间不变性直接推出切换闭合性 \\
引理~\ref{lem:sector-switch-implies-stack-switch-and-transport} & \texttt{\detokenize{cert:
  sector-switch-implies-stack-switch-and-transport}} & 证书：语义一致性链（扇区约束 $\Rightarrow$ 算法口径） \\
定理~\ref{thm:sectorized-loop-consistency} & \texttt{\detokenize{cert:sectorized-loop-consistency}} & 数学归纳法证书：按步数 $T$
  归纳的一致性保持 \\
推论~\ref{cor:no-sector-switch-reduces-to-intra-sector} & \texttt{\detokenize{cert:
  no-sector-switch-reduces-to-intra-sector}} & 证书：无切换 $\Rightarrow$ 扇区常值 $\Rightarrow$ 算法栈常值 \\
定理~\ref{thm:lifecycle-is-a-directed-path} & \texttt{\detokenize{cert:lifecycle-is-a-directed-path}} & 数学归纳法证书：按生命周期步数
  $T$ 归纳的路径表示 \\
命题~\ref{prop:rewriting-reduces-search-space} & \texttt{\detokenize{cert:rewriting-reduces-search-space}} & 证书：
  等价类枚举的正确性前提 \\
引理~\ref{lem:bidirectional-search-correctness} & \texttt{\detokenize{cert:bidirectional-search-correctness}} & 证书：
  边关系的双向可达性等价 \\
定理~\ref{thm:openra-subbundle-interface-decomposition} & \texttt{\detokenize{cert:
  openra-subbundle-interface-decomposition}} & 数学归纳法证书：按步数 $T$ 归纳的接口闭合性 \\
命题~\ref{prop:openra-subbundle-parallel-local-search} & \texttt{\detokenize{cert:openra-subbundle-parallel-local-search}}
  & 证书：接口一致性 $\Rightarrow$ 可拼接 \\
引理~\ref{lem:openra-structure-profile-and-scheduling} & \texttt{\detokenize{cert:openra-structure-profile-and-scheduling}}
  & 证书：投影到子丛 $\Rightarrow$ 结构轮廓；固定轮廓 $\Rightarrow$ 排程 \\
推论~\ref{cor:openra-noncausal-search-advantage} & \texttt{\detokenize{cert:openra-noncausal-search-advantage}} & 证书：
  由引理与接口压缩推出两阶段优势 \\
命题~\ref{prop:openra-sector-switch-requires-algorithm-switch} & \texttt{\detokenize{cert:
  openra-sector-switch-requires-algorithm-switch}} & 证书：态势改变 $\Rightarrow$ 约束与风险口径改变 $\Rightarrow$ 栈需切换 \\
定理~\ref{thm:openra-guard-three-criteria} & \texttt{\detokenize{cert:openra-guard-three-criteria}} & 证书：三要件合取作为守卫条件的审计链
  \\
命题~\ref{prop:openra-opponent-belief-necessary} & \texttt{\detokenize{cert:openra-opponent-belief-necessary}} & 证书：
  同态不可分辨性（同内态不同外态） \\
引理~\ref{lem:openra-debounce-hysteresis-reduces-chattering} & \texttt{\detokenize{cert:
  openra-debounce-hysteresis-reduces-chattering}} & 证书：连续满足 + 退出更严格 $\Rightarrow$ 不会被单步噪声触发回切 \\
命题~\ref{prop:openra-sector-derived-from-functional} & \texttt{\detokenize{cert:openra-sector-derived-from-functional}} &
  证书：派生标签不改变最优解（同一目标与候选集） \\
定理~\ref{thm:openra-repair-tower-monotone} & \texttt{\detokenize{cert:openra-repair-tower-monotone}} & 数学归纳法证书：按层级 $L$
  归纳的单调性 \\
推论~\ref{cor:openra-lexicographic-implies-low-level-first} & \texttt{\detokenize{cert:
  openra-lexicographic-implies-low-level-first}} & 证书：若更高层不改进缺口，则字典序会选更小烈度的低层代表 \\
引理~\ref{lem:openra-l1-slicing-optimal-certificate} & \texttt{\detokenize{cert:openra-l1-slicing-optimal-certificate}} &
  证书：有限可行集的字典序最小元存在（可枚举） \\
引理~\ref{lem:openra-l3-epsilon-optimal-certificate} & \texttt{\detokenize{cert:openra-l3-epsilon-optimal-certificate}} &
  证书：有限集取最小元 + 与空操作比较 \\
命题~\ref{prop:openra-online-descent-induces-tower} & \texttt{\detokenize{cert:openra-online-descent-induces-tower}} & 证书：
  字典序比较 + 有限候选 $\Rightarrow$ 选择规则可执行且可审计 \\
定理~\ref{thm:openra-snapshot-consistency-reproducible} & \texttt{\detokenize{cert:
  openra-snapshot-consistency-reproducible}} & 数学归纳法证书：按 tick 归纳的可复算一致性 \\
命题~\ref{prop:openra-gap-lex-ordering} & \texttt{\detokenize{cert:openra-gap-lex-ordering}} & 证书：字典序的逐分量判定规则 \\
引理~\ref{lem:openra-tau-auditable} & \texttt{\detokenize{cert:openra-tau-auditable}} & 证书：逐步累积与可复算性闭合 \\
定理~\ref{thm:openra-generator-finiteness} & \texttt{\detokenize{cert:openra-generator-finiteness}} & 证书：
  有限并集/有限子集/有界操作的封闭性 \\
定理~\ref{thm:openra-level-switch-stability} & \texttt{\detokenize{cert:openra-level-switch-stability}} & 证书：
  把扇区防抖/滞回引理平移到层级切换 \\
命题~\ref{prop:migration-sector-switch-as-functional-coupling} & \texttt{\detokenize{cert:
  migration-sector-switch-as-functional-coupling}} & 证书：扇区标签是派生描述，不改变候选集与比较准则 \\
引理~\ref{lem:grid-operator-word-noncommutative} & \texttt{\detokenize{cert:grid-operator-word-noncommutative}} & 证书：
  以电压轨迹/保护风险差异作见证 \\
命题~\ref{prop:qpu-roomtemp-sc-as-sector-switch} & \texttt{\detokenize{cert:qpu-roomtemp-sc-as-sector-switch}} & 证书：约束结构改变
  $\Rightarrow$ 候选族与缺口分解改变 \\
定理~\ref{thm:migration-shared-closed-loop-interface} & \texttt{\detokenize{cert:migration-shared-closed-loop-interface}}
  & 证书：由三个已证结论拼接闭环接口 \\
定理~\ref{thm:online-multi-rate-unified-interface} & \texttt{\detokenize{cert:online-multi-rate-unified-interface}} & 证书：把
  tick 类型差异吸收到“候选族 + time-to-effect”中 \\
命题~\ref{prop:online-l3plus-template-only} & \texttt{\detokenize{cert:online-l3plus-template-only}} & 证书：组合爆炸 + 高风险耦合
  $\Rightarrow$ 模板库约束是必要护栏 \\
定理~\ref{thm:cpu-bounded-orl-online-approx} & \texttt{\detokenize{cert:cpu-bounded-orl-online-approx}} & 数学归纳法证书：按 tick
  数归纳的“可复算 + 局部下降 + 可控扩张”链 \\
命题~\ref{prop:public-realtime-marl-exists} & \texttt{\detokenize{cert:public-realtime-marl-exists}} & 证书：
  公开论文/技术报告记录（引用即证书） \\
命题~\ref{prop:public-self-driving-lab-exists} & \texttt{\detokenize{cert:public-self-driving-lab-exists}} & 证书：
  闭环自驱动实验室的公开记录（论文/综述） \\
定理~\ref{thm:public-coverage-corrected-statement} & \texttt{\detokenize{cert:public-coverage-corrected-statement}} & 证书：
  三轴分解下的证据拼接 \\
推论~\ref{cor:replace-infinite-strategic-value} & \texttt{\detokenize{cert:replace-infinite-strategic-value}} & 证书：
  三条主张的可检验落点 \\
命题~\ref{prop:bridge-a-representation} & \texttt{\detokenize{cert:bridge-a-representation}} & 证书：策略定义 $\Rightarrow$
  动作词序列的存在 \\
命题~\ref{prop:bridge-b-necessity} & \texttt{\detokenize{cert:bridge-b-necessity}} & 证书：漂移存在 + 无证书 $\Rightarrow$ 证书链断裂 \\
命题~\ref{prop:bridge-c-minimality} & \texttt{\detokenize{cert:bridge-c-minimality}} & 证书：有限层级取最小 + 层内字典序最小元 \\
命题~\ref{prop:release-template-restriction} & \texttt{\detokenize{cert:release-template-restriction}} & 证书：
  高风险动作模板与运行期护栏的耦合 \\
定理~\ref{thm:priority-four-conditions-sufficient} & \texttt{\detokenize{cert:priority-four-conditions-sufficient}} &
  数学归纳法证书：按要件数 $n$ 归纳的优先权证据链 \\
命题~\ref{prop:prior-art-and-window-tradeoff} & \texttt{\detokenize{cert:prior-art-and-window-tradeoff}} & 证书：prior art
  强化与新颖性窗口压缩的双刃链 \\
定理~\ref{thm:release-psr-serves-both} & \texttt{\detokenize{cert:release-psr-serves-both}} & 数学归纳法证书：按发布级别链 P$\to$S$\to$R
  归纳的“固化--可信--控危害” \\
命题~\ref{prop:priority-pack-strengthens} & \texttt{\detokenize{cert:priority-pack-strengthens}} & 证书：
  六件套分别补强“时间戳/版本/可复核/可引用”链路 \\
定理~\ref{thm:dialogue-gains-to-catalog-system} & \texttt{\detokenize{cert:dialogue-gains-to-catalog-system}} & 数学归纳法证书：按
  8 项增益递增构造目录化对象 \\
引理~\ref{lem:catalog-system-closed-loop-interface} & \texttt{\detokenize{cert:catalog-system-closed-loop-interface}} & 证书：
  对象—映射—切换规则闭合性检查清单 \\
命题~\ref{prop:next-actions-three-deliverables-minimal} & \texttt{\detokenize{cert:next-actions-three-deliverables-minimal}
  } & 证书：三件交付件到目录化系统五元组的派生规则 \\
推论~\ref{cor:openra-guard-template-transfer} & \texttt{\detokenize{cert:openra-guard-template-transfer}} & 证书：同构替换与接口保持
  \\
定理~\ref{thm:zenodo-priority-fixed} & \texttt{\detokenize{cert:zenodo-priority-fixed}} & 数学归纳法证书：按记录数 $N$ 归纳的优先权固化 \\
定理~\ref{thm:zenodo-doi-versioning} & \texttt{\detokenize{cert:zenodo-doi-versioning}} & 证书：版本/概念 DOI 的元数据可追溯性 \\
引理~\ref{lem:zenodo-stats-consistency} & \texttt{\detokenize{cert:zenodo-stats-consistency}} & 证书：统计字段与去重规则 \\
命题~\ref{prop:zenodo-stats-dual-reporting} & \texttt{\detokenize{cert:zenodo-stats-dual-reporting}} & 证书：双口径字段的可比性 \\
推论~\ref{cor:zenodo-priority-defensive} & \texttt{\detokenize{cert:zenodo-priority-defensive}} & 证书：公开披露与防御性公开的蕴含链 \\
引理~\ref{lem:jacobiator-bracketing-sensitivity} & \texttt{\detokenize{cert:jacobiator-bracketing-sensitivity}} & 证书：表示有界与
  Jacobiator 关联 \\
引理~\ref{lem:safety-nonclosure-without-jacgap} & \texttt{\detokenize{cert:safety-nonclosure-without-jacgap}} & 证书：
  安全裕度被括号化差异跨越 \\
定理~\ref{thm:necessity-jacgap} & \texttt{\detokenize{cert:necessity-jacgap}} & 证书：强约束 + 实现差异 + 开放系统 \\
引理~\ref{lem:feasible-set-monotone} & \texttt{\detokenize{cert:feasible-set-monotone}} & 证书：嵌套可行域 \\
定理~\ref{thm:min-feasible-layer-exists} & \texttt{\detokenize{cert:min-feasible-layer-exists}} & 数学归纳法证书：按层数 $N$ 归纳 \\
定理~\ref{thm:lexicographic-min-intensity} & \texttt{\detokenize{cert:lexicographic-min-intensity}} & 证书：最小层级与烈度单调 \\
命题~\ref{prop:homotopy-twist-min-repair} & \texttt{\detokenize{cert:homotopy-twist-min-repair}} & 证书：最小修复数据的同伦因子化 \\
定理~\ref{thm:jacgap-zero-strictification} & \texttt{\detokenize{cert:jacgap-zero-strictification}} & 证书：JacGap 零与严格化 \\
命题~\ref{prop:repair-tower-truncation-projection} & \texttt{\detokenize{cert:repair-tower-truncation-projection}} & 证书：
  截断可行域与求解等价 \\
推论~\ref{cor:existing-methods-projection} & \texttt{\detokenize{cert:existing-methods-projection}} & 证书：严格化与截断的归并 \\
定理~\ref{thm:round-window-concatenation} & \texttt{\detokenize{cert:round-window-concatenation}} & 数学归纳法证书：按回合数 $N$ 归纳 \\
引理~\ref{lem:curvature-bounds-round-window} & \texttt{\detokenize{cert:curvature-bounds-round-window}} & 证书：曲率--holonomy
  估计 \\
引理~\ref{lem:jacgap-bounds-bracketing} & \texttt{\detokenize{cert:jacgap-bounds-bracketing}} & 证书：JacGap 与语义不变性 \\
定理~\ref{thm:closed-system-large-round} & \texttt{\detokenize{cert:closed-system-large-round}} & 证书：封闭系统的回合闭合 \\
定理~\ref{thm:open-system-tick-round} & \texttt{\detokenize{cert:open-system-tick-round}} & 证书：开放系统的回合收缩 \\
命题~\ref{prop:external-info-not-defining} & \texttt{\detokenize{cert:external-info-not-defining}} & 证书：信息注入 $\Rightarrow$
  曲率/Gap 变动 \\
推论~\ref{cor:round-window-cohomology} & \texttt{\detokenize{cert:round-window-cohomology}} & 证书：判据归并 \\
引理~\ref{lem:magic-wand-verifier-sufficiency} & \texttt{\detokenize{cert:magic-wand-verifier-sufficiency}} & 证书：验证通过
  $\Rightarrow$ 规格满足 \\
定理~\ref{thm:magic-wand-axioms-downgrade} & \texttt{\detokenize{cert:magic-wand-axioms-downgrade}} & 数学归纳法证书：按扰动步数 $n$
  的在线可修复性 \\
命题~\ref{prop:magic-wand-discrete-termination} & \texttt{\detokenize{cert:magic-wand-discrete-termination}} & 证书：严格下降
  $\wedge$ 离散值域 $\Rightarrow$ 有界步数 \\
定理~\ref{thm:magic-wand-from-law-object} & \texttt{\detokenize{cert:magic-wand-from-law-object}} & 证书：由缺陷势归零与一致性证书推出可核验解
  \\
定理~\ref{thm:magic-wand-into-law-object} & \texttt{\detokenize{cert:magic-wand-into-law-object}} & 证书：把 $\mathbb{W}$
  编码进法对象与缺陷势 \\
推论~\ref{cor:magic-wand-equivalence} & \texttt{\detokenize{cert:magic-wand-equivalence}} & 证书：互相构造 $\Rightarrow$ 结构等价 \\
引理~\ref{lem:magic-wand-puremath-entry} & \texttt{\detokenize{cert:magic-wand-puremath-entry}} & 证书：对象层 $\Rightarrow$
  法对象层（统一入口） \\
定理~\ref{thm:magic-wand-puremath-jacgap-vanishing} & \texttt{\detokenize{cert:magic-wand-puremath-jacgap-vanishing}} & 证书：
  JacGap 消失 $\Rightarrow$ strictify 可用且不破坏 protected observables \\
定理~\ref{thm:magic-wand-puremath-monotone} & \texttt{\detokenize{cert:magic-wand-puremath-monotone}} & 数学归纳法证书：按迭代步数 $n$
  的保持性与单调性 \\
定理~\ref{thm:magic-wand-puremath-existence} & \texttt{\detokenize{cert:magic-wand-puremath-existence}} & 证书：单调下降 + 极限不可再降
  $\Rightarrow$ 规范形最小代表 \\
推论~\ref{cor:magic-wand-puremath-strict-branch} & \texttt{\detokenize{cert:magic-wand-puremath-strict-branch}} & 证书：
  $\JacGap(L^\star)=0\Rightarrow$ strictify 可用 \\
推论~\ref{cor:magic-wand-puremath-vol1} & \texttt{\detokenize{cert:magic-wand-puremath-vol1}} & 证书：Vol1 接口 $\Rightarrow$
  规范形 + 证书输出 \\
命题~\ref{prop:magic-wand-puremath-terminal-functor} & \texttt{\detokenize{cert:magic-wand-puremath-terminal-functor}} &
  证书：终对象性质 + 最小性 $\Rightarrow$ 普适规范形 \\
定理~\ref{thm:wand-multistep-verification-induction} & \texttt{\detokenize{cert:wand-multistep-verification-induction}} &
  数学归纳法证书：按步数 $n$ 的证书保持与缺陷势下降 \\
命题~\ref{prop:strict-inclusion-implies-leap} & \texttt{\detokenize{cert:strict-inclusion-implies-leap}} & 证书：严格包含
  $\Rightarrow$ 存在分离函数 $U$ \\
引理~\ref{lem:dual-reachable-expansion} & \texttt{\detokenize{cert:dual-reachable-expansion}} & 证书：存在性 $\Rightarrow$ 可操作性
  $\Rightarrow$ 覆盖可达 \\
引理~\ref{lem:dual-branching} & \texttt{\detokenize{cert:dual-branching}} & 证书：由假设~\ref{asm:dual-A2} 直接给出分岔 \\
引理~\ref{lem:dual-homotopy-invariant-reachability} & \texttt{\detokenize{cert:dual-homotopy-invariant-reachability}} & 证书：
  互模拟 $\Rightarrow$ 可达集存在性保持 \\
定理~\ref{thm:dual-properties} & \texttt{\detokenize{cert:dual-properties}} & 证书：跃迁/失控/同伦不变的同源链 \\
推论~\ref{cor:dual-properties-governance} & \texttt{\detokenize{cert:dual-properties-governance}} & 证书：消除不可区分分岔
  $\Rightarrow$ 失控路径被证书排除 \\
定理~\ref{thm:rt-open-adversarial-turn-equivalence} & \texttt{\detokenize{cert:rt-open-adversarial-turn-equivalence}} &
  数学归纳法证书：回合边界轨线一致 \\
推论~\ref{cor:rt-open-adversarial-turn-principle} & \texttt{\detokenize{cert:rt-open-adversarial-turn-principle}} & 证书：
  轨线等价 $\Rightarrow$ 评价等价 \\
定理~\ref{thm:rt-open-adversarial-uu-impossible} & \texttt{\detokenize{cert:rt-open-adversarial-uu-impossible}} & 证书：同一历史
  $\Rightarrow$ 同一动作；安全集合互斥 $\Rightarrow$ 至少一侧证书破坏 \\
引理~\ref{lem:rt-open-adversarial-lag-lower-bound} & \texttt{\detokenize{cert:rt-open-adversarial-lag-lower-bound}} &
  数学归纳法证书：相同数据前缀 $\Rightarrow$ 相同内部状态 \\
命题~\ref{prop:rt-open-adversarial-combined} & \texttt{\detokenize{cert:rt-open-adversarial-combined}} & 证书：三段定理链的拼接 \\
命题~\ref{prop:rt-open-adversarial-need-defect-repair} & \texttt{\detokenize{cert:rt-open-adversarial-need-defect-repair}}
  & 证书：统一修复动作存在 $\Rightarrow$ 鲁棒安全集合非空 \\
定理~\ref{thm:uu-realtime-iff-defect-repair} & \texttt{\detokenize{cert:uu-realtime-iff-defect-repair}} & 证书：充要条件的两向构造与归纳链
  \\
推论~\ref{cor:uu-realtime-impossible-without-defect-structure} & \texttt{\detokenize{cert:
  uu-realtime-impossible-without-defect-structure}} & 证书：充要条件的逆否命题 \\
定理~\ref{thm:uu-eng-iff} & \texttt{\detokenize{cert:uu-eng-iff}} & 数学归纳法证书：$\Phi_\lambda$ 按步数严格下降并在有限步内归零 \\
定理~\ref{thm:lag-indistinguishable-prefix} & \texttt{\detokenize{cert:lag-indistinguishable-prefix}} & 数学归纳法证书：相同历史前缀
  $\Rightarrow$ 相同内部状态与动作 \\
命题~\ref{prop:pure-nn-not-uu-robust-descent} & \texttt{\detokenize{cert:pure-nn-not-uu-robust-descent}} & 证书：不可区分分岔
  $\Rightarrow$ “同动作” $\Rightarrow$ 至少一侧证书破坏 \\
引理~\ref{lem:lag-holonomy-gives-pair} & \texttt{\detokenize{cert:lag-holonomy-gives-pair}} & 证书：holonomy 见证给出同语义不同内部态 \\
引理~\ref{lem:lag-pio-cannot-distinguish} & \texttt{\detokenize{cert:lag-pio-cannot-distinguish}} & 证书：因子分解
  $\pi=\bar\pi\circ p$ $\Rightarrow$ 同语义同动作 \\
命题~\ref{prop:lag-holonomy-implies-no-unif-repair} & \texttt{\detokenize{cert:lag-holonomy-implies-no-unif-repair}} & 证书：
  两项交为空 $\Rightarrow$ 全局交为空 \\
定理~\ref{thm:lag-holonomy-to-lag-class} & \texttt{\detokenize{cert:lag-holonomy-to-lag-class}} & 证书：同语义同动作 + 无统一修复动作
  $\Rightarrow$ 至少一侧修复义务失败 \\
推论~\ref{cor:nn-main-in-lag-via-holonomy} & \texttt{\detokenize{cert:nn-main-in-lag-via-holonomy}} & 证书：因子分解到 $O$ 上
  $\Rightarrow$ 属于 $\Pi_O$ \\
推论~\ref{cor:demote-nn-from-main-loop} & \texttt{\detokenize{cert:demote-nn-from-main-loop}} & 证书：主环证明义务 $\Rightarrow$ NN
  只能作为候选/估计器（fail-closed） \\
引理~\ref{lem:audited-wrapper-step-contract} & \texttt{\detokenize{cert:audited-wrapper-step-contract}} & 证书：单步审计通过
  $\Rightarrow$ 合同成立 \\
定理~\ref{thm:audited-wrapper-soundness} & \texttt{\detokenize{cert:audited-wrapper-soundness}} & 证书：审计规则集 +
  审计日志/回归测试/变形测试 \\
定理~\ref{thm:audited-wrapper-inductive-safety} & \texttt{\detokenize{cert:audited-wrapper-inductive-safety}} & 证书：基步 +
  归纳步（回合归纳证书） \\
命题~\ref{prop:audited-wrapper-removes-proof-burden} & \texttt{\detokenize{cert:audited-wrapper-removes-proof-burden}} &
  证书：证明义务从 $G_\theta$ 迁移到 $V$ 的裁决日志 \\
推论~\ref{cor:audited-wrapper-nn-selector} & \texttt{\detokenize{cert:audited-wrapper-nn-selector}} & 证书：索引选择 +
  $\mathsf{FirstAccept}$ 的确定性 $\Rightarrow$ 只影响顺序 \\
推论~\ref{cor:audited-wrapper-nn-proof-sketch} & \texttt{\detokenize{cert:audited-wrapper-nn-proof-sketch}} & 证书：草稿可疑，
  裁决可核验（证书载体为 $V$ 的日志） \\
推论~\ref{cor:audited-wrapper-compression-safety} & \texttt{\detokenize{cert:audited-wrapper-compression-safety}} & 证书：
  fail-closed $\Rightarrow$ 压缩模块不成为无证书错误来源 \\
定理~\ref{thm:round-homotopy-transport-composition} & \texttt{\detokenize{cert:round-homotopy-transport-composition}} &
  数学归纳法证书：按路径长度 $n$ 定义 $T_\gamma$ 并验证复合律 \\
引理~\ref{lem:round-homotopy-endpoint-iff-trivial-holonomy} & \texttt{\detokenize{cert:
  round-homotopy-endpoint-iff-trivial-holonomy}} & 证书：端点决定性与回路平凡性的互推 \\
定理~\ref{thm:round-homotopy-open-closed} & \texttt{\detokenize{cert:round-homotopy-open-closed}} & 证书：封闭/开放的同伦类刻画 \\
定理~\ref{thm:round-subdivision-homotopy-invariant} & \texttt{\detokenize{cert:round-subdivision-homotopy-invariant}} &
  数学归纳法证书：按细分次数 $n$ 的同伦不变性 \\
定理~\ref{thm:closed-implies-strict-round-homotopy} & \texttt{\detokenize{cert:closed-implies-strict-round-homotopy}} & 证书：
  封闭 $\Rightarrow$ 细化仅为边细分 $\Rightarrow$ 同伦型稳定 \\
定理~\ref{thm:closed-implies-ku} & \texttt{\detokenize{cert:closed-implies-ku}} & 证书：真实后继闭合 $\Rightarrow$ 缺口为空 \\
定理~\ref{thm:open-implies-no-strict-round-homotopy} & \texttt{\detokenize{cert:open-implies-no-strict-round-homotopy}} &
  证书：反证法——若稳定则与“任意分割可破坏”矛盾 \\
定理~\ref{thm:open-implies-approx-and-uu} & \texttt{\detokenize{cert:open-implies-approx-and-uu}} & 证书：不可稳定性 $\Rightarrow$
  仅能逼近；模型不闭合 $\Rightarrow$ 缺口非空 \\
引理~\ref{lem:uu-witness-implies-jacgap-positive} & \texttt{\detokenize{cert:uu-witness-implies-jacgap-positive}} & 证书：
  非空缺口 $\Rightarrow$ 测度为正 $\Rightarrow$ 上确界为正 \\
定理~\ref{thm:strict-round-homotopy-iff-jacgap-stable-zero} & \texttt{\detokenize{cert:
  strict-round-homotopy-iff-jacgap-stable-zero}} & 证书：稳定归零 $\Rightarrow$ 仅细分 $\Rightarrow$ 同伦型稳定；反向用缺口--结构对应 \\
定理~\ref{thm:repair-tower-terminates} & \texttt{\detokenize{cert:repair-tower-terminates}} & 数学归纳法证书：按步数 $n$ 的势能下降与证书保持
  \\
定理~\ref{thm:closed-phi-jacgap-ku} & \texttt{\detokenize{cert:closed-phi-jacgap-ku}} & 证书：封闭 $\Rightarrow$ 缺口为空
  $\Rightarrow$ $\JacGap=0$ \\
定理~\ref{thm:open-phi-jacgap-uu} & \texttt{\detokenize{cert:open-phi-jacgap-uu}} & 证书：开放 $\Rightarrow$ 缺口非空 $\Rightarrow$
  $\JacGap>0$ \\
定理~\ref{thm:subdivision-preserve-transport} & \texttt{\detokenize{cert:subdivision-preserve-transport}} & 数学归纳法证书：
  按细分次数保持整体传输不变 \\
定理~\ref{thm:uu-open-subdivision-not-enough} & \texttt{\detokenize{cert:uu-open-subdivision-not-enough}} & 证书：非平凡
  holonomy 在回合细分下保持 \\
定理~\ref{thm:layering-complete-iff-trivial-holonomy} & \texttt{\detokenize{cert:layering-complete-iff-trivial-holonomy}}
  & 证书：层完备性与抬升后平凡 holonomy 的互推 \\
推论~\ref{cor:layering-makes-uu-repairable} & \texttt{\detokenize{cert:layering-makes-uu-repairable}} & 证书：层完备性把非平凡
  holonomy 从“回路”变为“层切换” \\
定理~\ref{thm:challenge-round-induction-schema} & \texttt{\detokenize{cert:challenge-round-induction-schema}} & 数学归纳法证书：
  归纳模板（基步 + 归纳步） \\
命题~\ref{prop:challenge-implies-phi-pressure} & \texttt{\detokenize{cert:challenge-implies-phi-pressure}} & 证书：挑战见证
  $\Rightarrow$ $\JacGap>0$ 或 $\mathcal{E}_{\mathrm{CH}}>0$ \\
定理~\ref{thm:multiagent-implies-game} & \texttt{\detokenize{cert:multiagent-implies-game}} & 证书：后继依赖他方策略 $\Rightarrow$
  优化对象为策略对 \\
推论~\ref{cor:closed-eval-does-not-imply-open} & \texttt{\detokenize{cert:closed-eval-does-not-imply-open}} & 证书：封闭
  $\Rightarrow$ 仅 KU；高分不约束 UU \\
推论~\ref{cor:pressure-implies-capability-necessity} & \texttt{\detokenize{cert:pressure-implies-capability-necessity}} &
  证书：缺少能力 $\Rightarrow$ 下降条件被击穿 \\
定理~\ref{thm:challenge-shapes-wisdom-structurally} & \texttt{\detokenize{cert:challenge-shapes-wisdom-structurally}} & 证书：
  $\mathsf{Challenge}\Rightarrow \Phi>0$；$\Phi$ 可终止下降 $\Rightarrow$ 智慧结构 \\
定理~\ref{thm:certspec-t-closed} & \texttt{\detokenize{cert:certspec-t-closed}} & 证书字段：$T_{\mathrm{Closed}}$ \\
定理~\ref{thm:certspec-t-open} & \texttt{\detokenize{cert:certspec-t-open}} & 证书字段：$T_{\mathrm{Open}}$（UU 反例证书） \\
定理~\ref{thm:certspec-t-noroundonly} & \texttt{\detokenize{cert:certspec-t-noroundonly}} & 证书字段：$T_{\mathrm{NoRoundOnly}}
  $（引用开放性见证 + 细分不变性） \\
定理~\ref{thm:certspec-t-layering} & \texttt{\detokenize{cert:certspec-t-layering}} & 证书字段：$T_{\mathrm{LayeringNecessary}}
  $ \\
定理~\ref{thm:certspec-t-lag} & \texttt{\detokenize{cert:certspec-t-lag}} & 证书字段：$T_{\mathsf{Lag}}$（不可区分 + 修复不可兼容） \\
定理~\ref{thm:certspec-t-auditedwrapper} & \texttt{\detokenize{cert:certspec-t-auditedwrapper}} & 证书字段：
  $T_{\mathrm{AuditedWrapper}}$（规则版本 + 测试记录 + 裁决日志） \\
定理~\ref{thm:eng-design-is-open-game} & \texttt{\detokenize{cert:eng-design-open-witness}} & 证书（工程版）：UU 见证记录/证书边界变化记录 \\
定理~\ref{thm:eng-design-accept-reject-induction} & \texttt{\detokenize{cert:eng-design-accept-reject-induction}} &
  数学归纳法证书：按接受步数 $n$ 的保持性与单调性 \\
定理~\ref{thm:eng-design-ohutower-migration} & \texttt{\detokenize{cert:eng-design-ohutower-migration}} & 证书：OHU 单调性 +
  JacGap-free strictify 分支 \\
推论~\ref{cor:eng-design-migration} & \texttt{\detokenize{cert:eng-design-migration}} & 证书：迁移模板仅依赖“law-object +
  $\JacGap/\Phi$ 接口” \\
定理~\ref{thm:semantic-normalization-base} & \texttt{\detokenize{cert:semantic-normalization-base}} & 证书：语义保持 + 一致性 + 失败显式
  \\
推论~\ref{cor:eng-design-unified-workflow} & \texttt{\detokenize{cert:eng-design-unified-workflow}} & 证书：工作流五步由对应定理链逐项支撑
  \\
引理~\ref{lem:eng-minimal-f-equivalence} & \texttt{\detokenize{cert:eng-minimal-f-equivalence}} & 证书：由 $F$ 构造 $(\chi,
  \mathrm{nf})$，并可回构 $F$ \\
定理~\ref{thm:eng-minimal-eprot-margin} & \texttt{\detokenize{cert:eng-minimal-eprot-margin}} & 证书：$\mathsf{Cert}
  _{\mathrm{prot}}$ 四元组与反例见证（fail-closed） \\
引理~\ref{lem:eng-minimal-ech-nonneg} & \texttt{\detokenize{cert:eng-minimal-ech-nonneg}} & 证书：范数非负 + 版本哈希锁定 \\
命题~\ref{prop:eng-minimal-uu-witness-implies-jacgap} & \texttt{\detokenize{cert:eng-minimal-uu-witness-implies-jacgap}} &
  证书：UU 见证 $\Rightarrow$ 缺口非空 $\Rightarrow$ $\JacGap>0$ \\
定理~\ref{thm:eng-minimal-loop-induction} & \texttt{\detokenize{cert:eng-minimal-loop-induction}} & 数学归纳法证书：按接受步数 $n$
  的保持性与下降性 \\
命题~\ref{prop:eng-minimal-engine-cert-instances} & \texttt{\detokenize{cert:eng-minimal-engine-cert-instances}} & 证书：
  扇区$\to$证书字段的映射表（接口级） \\
命题~\ref{prop:eng-minimal-chip-cert-instances} & \texttt{\detokenize{cert:eng-minimal-chip-cert-instances}} & 证书：签核/工具链哈希
  + UU 见证日志 \\
推论~\ref{cor:eng-minimal-engine-chip-isomorphic} & \texttt{\detokenize{cert:eng-minimal-engine-chip-isomorphic}} & 证书：
  模板复用只依赖三证书接口与闭环判据 \\
命题~\ref{prop:gains-concept-boundary} & \texttt{\detokenize{cert:gains-concept-boundary}} & 证书：关键定义引用闭合 \\
定理~\ref{thm:gains-core-chain} & \texttt{\detokenize{cert:gains-core-chain}} & 证书：五条结论的引用闭合 \\
命题~\ref{prop:gains-engineering-wrapper} & \texttt{\detokenize{cert:gains-engineering-wrapper}} & 证书：证明义务转移到 $V$；失败闭合 \\
命题~\ref{prop:gains-evaluation-path} & \texttt{\detokenize{cert:gains-evaluation-path}} & 证书：封闭高分不约束 UU；开封$\Rightarrow$压力；
  以 $\Phi$-下降验收 \\
推论~\ref{cor:gains-deposition} & \texttt{\detokenize{cert:gains-deposition}} & 证书：概念边界 + 定理链 + 工程接口 + 评测路径 $\Rightarrow$
  可沉淀为资产包 \\
定理~\ref{thm:gains-five-layer-induction} & \texttt{\detokenize{cert:gains-five-layer-induction}} & 数学归纳法证书：五层总结闭合（基步 +
  归纳步） \\
% END_CERT_INDEX
\hline
\end{longtable}
}

\clearpage

\begin{thebibliography}{99}
\addcontentsline{toc}{section}{\refname}

\bibitem{ref_nsr_nwae047}
【期刊论文】OUP Academic（National Science Review）页面：
\url{https://academic.oup.com/nsr/article/11/7/nwae047/7613947}（访问于 2026-01-19）。

\bibitem{ref_slac_2025_room_pressure}
【新闻】SLAC 新闻页面：
\url{https://www6.slac.stanford.edu/news/2025-02-04-first-researchers-stabilize-promising-new-class-high-temperature-sup
  erconductors-at-room-pressure}（访问于 2026-01-19）。

\bibitem{ref_nature_2025_08755_z}
【期刊论文】Nature 论文页面：
\url{https://www.nature.com/articles/s41586-025-08755-z}（访问于 2026-01-19）。

\bibitem{ref_uh_2025_pqp}
【新闻】University of Houston 新闻页面：
\url{https://www.uh.edu/news-events/stories/2025/february/02102025-tcsuh-breakthrough.php}（访问于 2026-01-19）。

\bibitem{ref_nature_2022_05294_9}
【撤稿说明】Nature 撤稿说明页面：
\url{https://www.nature.com/articles/s41586-022-05294-9}（访问于 2026-01-19）。

\bibitem{ref_sciam_retraction_note}
【媒体】Scientific American 文章页面：
\url{https://www.scientificamerican.com/article/nature-retracts-controversial-room-temperature-superconductor-study/}
  （访问于 2026-01-19）。

\bibitem{ref_wsj_61288727}
【媒体】The Wall Street Journal 文章页面：
\url{https://www.wsj.com/science/university-rochester-ranga-dias-superconductor-misconduct-61288727}（访问于 2026-01-19）。

\bibitem{ref_wsj_aa2f9fd4}
【媒体】The Wall Street Journal 文章页面：
\url{https://www.wsj.com/science/university-rochester-ranga-dias-superconductor-misconduct-aa2f9fd4}（访问于 2026-01-19）。

\bibitem{ref_time_ambient_superconductor}
【媒体】TIME 文章页面：
\url{https://time.com/6301391/experts-skeptical-about-ambient-superconductor/}（访问于 2026-01-19）。

\bibitem{ref_syscop_zanelli2017_homotopy_nmpc}
【论文】同伦内点法用于 NMPC 的论文（PDF）：
\url{https://publications.syscop.de/Zanelli2017.pdf}（访问于 2026-01-19）。

\bibitem{ref_merl_tr2018_082}
【技术报告】Real-Time Iteration（RTI）相关技术报告（PDF）：
\url{https://shadow.merl.com/publications/docs/TR2018-082.pdf}（访问于 2026-01-19）。

\bibitem{ref_pnas_2501048122}
【期刊论文】PNAS 论文页面（压力诱导态回收/亚稳保留相关）：
\url{https://www.pnas.org/doi/10.1073/pnas.2501048122}（访问于 2026-01-19）。

\bibitem{ref_pubmed_39701131}
【论文页面】PubMed 页面（薄膜体系常压信号/应变口径相关）：
\url{https://pubmed.ncbi.nlm.nih.gov/39701131/}（访问于 2026-01-19）。

\bibitem{ref_aps_prmaterials_8_044801}
【期刊论文】APS 论文页面（应变/结构路线稳定目标态相关）：
\url{https://link.aps.org/doi/10.1103/PhysRevMaterials.8.044801}（访问于 2026-01-19）。

\bibitem{ref_energycentral_grl}
【行业文章】Energy Central 文章（Grid Readiness 等分级讨论）：
\url{https://www.energycentral.com/renewables/post/grid-readiness-levels-technology-readiness-systems-integration-specia
  l-V368IXFmf90T2A8}（访问于 2026-01-19）。

\bibitem{ref_oecd_advanced_materials}
【报告】OECD 报告（先进材料路线/治理与成熟度相关，PDF）：
\url{https://www.oecd.org/content/dam/oecd/en/publications/reports/2025/10/steering-the-future-of-advanced-materials_9ad
  d3995/a15874fa-en.pdf}（访问于 2026-01-19）。

\bibitem{ref_nasa_trl}
【机构文档】NASA 页面：Technology Readiness Levels（TRL）：
\url{https://www.nasa.gov/directorates/somd/space-communications-navigation-program/technology-readiness-levels/}（访问于
  2026-01-19）。

\bibitem{ref_stage_gate_dtl}
【方法论】Stage-Gate International 页面：Discovery-To-Launch Process：
\url{https://www.stage-gate.com/about/stage-gate-innovation-performance-framework/discovery-to-launch-process/}（访问于
  2026-01-19）。

\bibitem{ref_noaa_swpc_alerts}
【机构文档】NOAA / NWS Space Weather Prediction Center 页面：Alerts, Watches and Warnings：
\url{https://www.swpc.noaa.gov/products/alerts-watches-and-warnings}（访问于 2026-01-19）。

\bibitem{ref_nws_space_ww}
【机构文档】National Weather Service 页面：Space Weather Watches, Warnings and Alerts：
\url{https://www.weather.gov/safety/space-ww}（访问于 2026-01-19）。

\bibitem{ref_swpc_subscription_services}
【机构文档】NOAA / NWS Space Weather Prediction Center 页面：Subscription Services：
\url{https://www.swpc.noaa.gov/content/subscription-services}（访问于 2026-01-19）。

\bibitem{ref_swpc_noaa_scales_explanation}
【机构文档】NOAA / NWS Space Weather Prediction Center 页面：NOAA Space Weather Scales（含 G 级解释入口）：
\url{https://www.swpc.noaa.gov/noaa-scales-explanation}（访问于 2026-01-19）。

\bibitem{ref_nerc_gmd_report_2012}
【报告】NERC 报告（PDF）：Effects of Geomagnetic Disturbances on the Bulk Power System：
\url{https://www.nerc.com/globalassets/programs/rapa/gmd/reference-documents/nerc_gmd_report_2012.pdf}（访问于 2026-01-19）。

\bibitem{ref_nerc_tpl_007_4}
【标准】NERC 可靠性标准（PDF）：TPL-007-4 -- Transmission System Planned Performance for Geomagnetic Disturbance Events：
\url{https://www.nerc.com/globalassets/standards/reliability-standards/tpl/tpl-007-4.pdf}（访问于 2026-01-19）。

\bibitem{ref_nerc_eop_010_1}
【标准】NERC 可靠性标准（PDF）：EOP-010-1 -- Geomagnetic Disturbance Operations：
\url{https://www.nerc.com/globalassets/standards/reliability-standards/eop/eop-010-1.pdf}（访问于 2026-01-19）。

\bibitem{ref_doe_gmd_monitoring_strategies}
【机构文档】U.S. Department of Energy（Energy.gov）文章：Geomagnetic Disturbance Monitoring Approach and Implementation Strategies：
\url{https://www.energy.gov/ceser/articles/geomagnetic-disturbance-monitoring-approach-and-implementation-strategies}
  （访问于 2026-01-19）。

\bibitem{ref_nerc_gmd_planning_guide_2013}
【指南】NERC 指南（PDF）：Geomagnetic Disturbance Planning Guide：
\url{https://www.nerc.com/globalassets/standards/projects/2013-03/geomagnetic-disturbance-task-force-gmdtf-2013_gmd-plan
  ning-guide_approved.pdf}（访问于 2026-01-19）。

\bibitem{ref_nerc_benchmark_gmd_event}
【基准事件】NERC 文档（PDF）：Benchmark Geomagnetic Disturbance Event Description：
\url{https://www.nerc.com/globalassets/standards/projects/2013-03/benchmark_gmd_event_aug27_clean.pdf}（访问于 2026-01-19）。

\bibitem{ref_iso_ne_gmd_2026_na_sow}
【技术文档】ISO New England 文档（PDF）：NERC TPL-007-4 Benchmark and Supplemental（北美场景评估材料）：
\url{https://www.iso-ne.com/static-assets/documents/2022/05/final_gmd_2026_na_sow.pdf}（访问于 2026-01-19）。

\bibitem{ref_aip_apr_superconducting_qubits}
【综述论文】AIP Publishing（Applied Physics Reviews）文章：A quantum engineer's guide to superconducting qubits：
\url{https://pubs.aip.org/aip/apr/article/6/2/021318/570326/A-quantum-engineer-s-guide-to-superconducting}（访问于
  2026-01-19）。

\bibitem{ref_springer_epjqt_2019_0072_0}
【期刊论文】Springer 论文页面：Engineering cryogenic setups for 100-qubit scale superconducting circuit systems：
\url{https://link.springer.com/article/10.1140/epjqt/s40507-019-0072-0}（访问于 2026-01-19）。

\bibitem{ref_nature_2025_09157_x}
【期刊论文】Nature 论文页面：Spin-qubit control with a milli-kelvin CMOS chip：
\url{https://www.nature.com/articles/s41586-025-09157-x}（访问于 2026-01-19）。

\bibitem{ref_naturecomm_cryogenic_pulse_gen}
【期刊论文】Nature Communications 论文页面：A cryogenic on-chip microwave pulse generator for large-...：
\url{https://www.nature.com/articles/s41467-024-50333-w}（访问于 2026-01-19）。

\bibitem{ref_prxq_5_010326}
【期刊论文】APS 论文页面：Using Cryogenic CMOS Control Electronics to Enable a Two ...（PRX Quantum 5.010326）：
\url{https://link.aps.org/doi/10.1103/PRXQuantum.5.010326}（访问于 2026-01-19）。

\bibitem{ref_nationalacademies_25196_ch13}
【报告/书籍】National Academies 在线书籍页面：Superconducting Quantum Computers（章节/附录入口）：
\url{https://www.nationalacademies.org/read/25196/chapter/13}（访问于 2026-01-19）。

\bibitem{ref_ibm_goldeneye_cryostat}
【企业材料】IBM Quantum 博客：IBM cools down world's largest quantum-ready cryostat（Goldeneye 概念系统）：
\url{https://www.ibm.com/quantum/blog/goldeneye-cryogenic-concept-system}（访问于 2026-01-19）。

\bibitem{ref_arxiv_2601_03922}
【预印本】arXiv 页面：Integration and Resource Estimation of Cryoelectronics for ...（预印本入口）：
\url{https://arxiv.org/abs/2601.03922}（访问于 2026-01-19）。

\bibitem{ref_pubmed_31666705}
【论文页面】PubMed 页面：Grandmaster level in StarCraft II using multi-agent reinforcement learning：
\url{https://pubmed.ncbi.nlm.nih.gov/31666705/}（访问于 2026-01-19）。

\bibitem{ref_openai_dota2}
【技术报告】OpenAI 技术报告：Dota 2 with Large Scale Deep Reinforcement Learning：
\url{https://cdn.openai.com/dota-2.pdf}（访问于 2026-01-19）。

\bibitem{ref_nature_2023_06734_w}
【期刊论文】Nature 论文页面：An autonomous laboratory for the accelerated synthesis of ...：
\url{https://www.nature.com/articles/s41586-023-06734-w}（访问于 2026-01-19）。

\bibitem{ref_acs_chemrev_4c00055}
【综述论文】ACS Publications 页面：Self-Driving Laboratories for Chemistry and Materials Science：
\url{https://pubs.acs.org/doi/10.1021/acs.chemrev.4c00055}（访问于 2026-01-19）。

\bibitem{ref_sciencedirect_s0004370222001515}
【期刊论文】ScienceDirect 论文页面：Sim-to-Lab-to-Real: Safe reinforcement learning with ...（参考入口）：
\url{https://www.sciencedirect.com/science/article/abs/pii/S0004370222001515}（访问于 2026-01-19）。

\bibitem{ref_jmlr_garcia15a_safe_rl}
【综述论文】JMLR 论文 PDF：A Comprehensive Survey on Safe Reinforcement Learning：
\url{https://www.jmlr.org/papers/volume16/garcia15a/garcia15a.pdf}（访问于 2026-01-19）。

\bibitem{ref_arxiv_2502_13187}
【预印本】arXiv 页面：A Survey of Sim-to-Real Methods in RL（预印本入口）：
\url{https://arxiv.org/abs/2502.13187}（访问于 2026-01-19）。

\bibitem{ref_acm_3616864}
【综述论文】ACM Digital Library 页面：Explainable Reinforcement Learning: A Survey and ...：
\url{https://dl.acm.org/doi/full/10.1145/3616864}（访问于 2026-01-19）。

\bibitem{ref_academia_priority}
【问答】Academia Stack Exchange 问答：Discreet way to establish priority（优先权固化讨论入口）：
\url{https://academia.stackexchange.com/questions/13216/discreet-way-to-establish-priority}（访问于 2026-01-19）。

\bibitem{ref_perkinscoie_prior_art}
【法律/合规】Perkins Coie 文章：When and How to Publish Technical Disclosures as Prior Art（prior art/防御性公开讨论入口）：
\url{https://perkinscoie.com/insights/publication/when-and-how-publish-technical-disclosures-prior-art}（访问于 2026-01-19）。

\bibitem{ref_wipo_patent_protection}
【法律/合规】WIPO 页面：How to Protect Inventions through Patents（专利保护与披露提示入口）：
\url{https://www.wipo.int/en/web/patents/protection}（访问于 2026-01-19）。

\bibitem{ref_uspto_mpep_2153}
【法律/合规】USPTO MPEP 页面：2153 Prior Art Exceptions Under 35 USC 102(b)(1) to AIA...（美国披露例外/宽限期入口）：
\url{https://www.uspto.gov/web/offices/pac/mpep/s2153.html}（访问于 2026-01-19）。

\bibitem{ref_cnipa_patent_law}
【法律/合规】CNIPA（中国国家知识产权局）英文站点：Patent Law/相关法律栏目入口：
\url{https://english.cnipa.gov.cn/col/col3068/index.html}（访问于 2026-01-19）。

\bibitem{GaoZheng2025GFramework}
Gao, Z. (2025).
\newblock Meta-Mathematical Theory based on Pan-Logic Analysis and
Pan-Iterative Analysis (GaoZheng G-Framework) and the
Principal-Bundle-Based Generalized Noncommutative Lie Algebra
(GaoZheng G-Algebra).
\newblock In \emph{Meta-Mathematical Theory based on Pan-Logic Analysis and
Pan-Iterative Analysis (GaoZheng G-Framework) and the
Principal-Bundle-Based Generalized Noncommutative Lie Algebra
(GaoZheng G-Algebra): An Integrated Construction (v1.0)}.
\newblock Zenodo.
\newblock \url{https://doi.org/10.5281/zenodo.17651584}.


\bibitem{GaoZhengAppVol1}
Gao, Z. (2025).
\newblock GaoZheng G-Framework and GaoZheng G-Algebra: Applied Mathematics Volume I — Law-Space Geometry, GRL, Quantum
  Computing, and Superconductivity.
\newblock In \emph{GaoZheng G-Framework and GaoZheng G-Algebra: Applied Mathematics Volume I — Law-Space Geometry, GRL,
  Quantum Computing, and Superconductivity (v1.2)}.
\newblock Zenodo.
\newblock \url{https://doi.org/10.5281/zenodo.17686133}.


\bibitem{ref_zenodo_about_records}
【平台文档】Zenodo 帮助文档：About records（记录发布与版本规则）：
\url{https://help.zenodo.org/docs/deposit/about-records/}（访问于 2026-01-19）。

\bibitem{ref_zenodo_doi_versioning}
【平台文档】Zenodo FAQ：What is DOI versioning?（版本 DOI 与概念 DOI）：
\url{https://support.zenodo.org/help/en-gb/1-upload-deposit/97-what-is-doi-versioning}（访问于 2026-01-19）。

\bibitem{ref_zenodo_statistics}
【平台文档】Zenodo 帮助文档：Statistics（浏览/下载统计口径）：
\url{https://zenodo.org/help/statistics}（访问于 2026-01-19）。

\end{thebibliography}

\clearpage

\end{CJK*}
\end{document}
